\documentclass[10pt]{article}
\usepackage[T1]{fontenc}
\usepackage[utf8]{inputenc}
\usepackage[english]{babel}
\usepackage{lmodern}
\usepackage{amsmath,amssymb,array,longtable,graphicx}
\usepackage{geometry}
\usepackage{microtype}
\geometry{paperwidth=6.5in,paperheight=9in,margin=0.52in}
\pagestyle{empty}
\setlength{\parindent}{1em}
\setlength{\parskip}{0.16em}
\setlength{\abovedisplayskip}{0.45em}
\setlength{\belowdisplayskip}{0.45em}
\emergencystretch=2em
\allowdisplaybreaks
\hbadness=10000
\hfuzz=2pt
\newcommand{\err}[4]{#1 & #2 & #3 & #4\\}
\newcommand{\tocentry}[3]{\S\,#1 & #2 & #3\\}
\newcommand{\tocsec}[2]{\multicolumn{3}{c}{\bfseries #1}\\[-0.15em]\multicolumn{3}{c}{\bfseries #2}\\[0.25em]}
\begin{document}
\small

% English translation of Clebsch and Gordan, Theorie der Abelschen Functionen, 1866.
% Translation scope: source PDF pp. 001--016; non-authorial scan pages are in the source witness/provenance only.

\begin{center}
{\Large\bfseries THEORY}\par\vspace{1.2em}
{\small\bfseries OF}\par\vspace{1.5em}
{\Large\bfseries ABELIAN FUNCTIONS.}
\end{center}

\clearpage

\begin{center}
{\Large\bfseries THEORY}\par\vspace{1.3em}
{\small\bfseries OF}\par\vspace{1.4em}
{\LARGE\bfseries ABELIAN FUNCTIONS}\par\vspace{2.2em}
{\small\bfseries BY}\par\vspace{1.8em}
{\large\bfseries A. CLEBSCH \quad AND \quad P. GORDAN,}\par
{\scriptsize\bfseries PROFESSORS AT THE UNIVERSITY OF GIESSEN.}\par\vfill
{\Large\bfseries LEIPZIG.}\par
{\bfseries PRINTED AND PUBLISHED BY B. G. TEUBNER.}\par
{\large 1866.}
\end{center}

\clearpage

\begin{center}
{\Large\bfseries Preface.}
\end{center}

Almost nine years have now passed since Riemann, torn so early from the sciences, created the theory of Abelian functions through his immortal works and unexpectedly enlarged the boundaries of analysis. This theory has since then rightly not lacked admirers, whereas understanding and advancement could not everywhere be achieved in equal measure. Only very recently Mr. Neumann, in his excellent work on this subject, presented a large part of Riemann's investigations in detail, gave their application to the hyperelliptic functions, and thus made this field accessible to study.

The difficulties of entering the theory of Abelian functions along Riemann's path rest exactly upon the same circumstances that are most suited to heighten admiration for the inventor. We mean, on the one hand, the foundation taken from the most general, and yet again cautiously restricted, concepts of function; on the other hand, the indirect procedure by which the functions to be investigated are constructed more synthetically than analytically developed. Precisely connected with this is the circumstance emphasized by Mr. Neumann, that it has hitherto not been possible to introduce the $\Theta$-function naturally as a necessary member of the theory.

When the authors were occupied with the thought of grounding this theory in a new way, the original source of the whole discipline in particular seemed to them to deserve renewed attention: the Jacobi inversion problem, upon which Mr. Weierstrass so successfully founded the theory of hyperelliptic functions. It seemed possible from the outset to pass directly from this problem to a formula analogous to the Jacobi formula
\[
\frac{\Theta(u)}{\Theta(0)}=e^{\int_0^u Z(u)\,du},
\]
and which therefore, as in Mr. Weierstrass's theory, would contain the solution of the inversion problem. In fact, only a few geometrical considerations and a few theorems from the theory of integration through the imaginary were needed in order to carry out this direct path completely. In this way one avoided in particular all considerations concerning functions in general, which are always precarious because they attach to a wholly indeterminate concept containing unknown possibilities. By its nature, the theory of Abelian functions must be capable of being settled by theorems that refer only to known functions, especially algebraic functions. When in a few places we appeal to theorems given in the work of Messrs. Briot and Bouquet, we understand these only within the range in which they apply to algebraic functions and to certain functions derived from them; within this range they are correct, while in further extension they are in part more than doubtful, as for years we have never neglected to explain in our lectures held at this university.

Since what is first at issue is a new fixing of the general outlines of a very broad theory, it will not be held against us that we have nowhere entered into special cases. Such special cases, which are always bound to special properties of the algebraic fundamental equation, also stand out especially in the hyperelliptic functions. We believed we could all the more dispense with a special mention of the latter: first, because it is very easy to derive their theory along our path as well; second, because their extremely special properties are only too well suited to divert attention from the points most essential to the general theory.

Among the parts of the theory treated by Riemann we believed we had to exclude the question of the number of moduli of a class of Abelian functions. This question has become an object of controversy through Mr. Cayley's acute considerations; indeed it is probably to be decided definitively, at first, only through deep algebraic investigations, for whose difficulties the methods now known no longer seem sufficient. And such investigations are in fact those from which, in our conviction, the further development of this discipline is to be expected, and in which the true deeper difficulties of the theory of Abelian functions must be seen. This discipline too finally ends in those branches of modern algebra which seem destined to be the center of all modern mathematical development.

\vspace{0.6em}
\noindent Giessen, 18 August 1866.

\begin{flushright}
{\bfseries The Authors.}
\end{flushright}

\vspace{1.0em}

\begin{center}
{\Large\bfseries Misprints.}
\end{center}

\begin{tabular}{@{}r@{\ }l@{\ }l@{\ }l@{}}
p. & 7. & line 9 from top, for $s_x$ & read $s_t$\\
--- & & line 8 from bottom, for $m+n-3$ & read $m+n-5$\\
p. & 11. & line 2 from top, for $s_x-s_x$ & read $s_x-s_x^\circ$\\
--- & & line 3 from top, for $\dfrac{\partial F}{\partial s}$ & read $\dfrac{\partial F}{\partial s_x}$\\
p. & 17. & line 10 from top, for $\gamma(x-x^\circ)$ & read $\beta(x-x^\circ)$\\
p. & 28. & line 6 from bottom, for $\lim S_{\xi\xi}$ & read $\lim S_{\xi\eta}$\\
p. & 29. & line 2 from top, for $\eta$ & read $\xi$\\
p. & 30. & line 14 from top to p. 31, line 6 from top, for $\eta$ & read $\xi$\\
p. & 39. & line 3 from bottom, for $Z_ns^n$ & read $Z_ms^m$\\
p. & 40. & line 8 from bottom, for $S_nz^n$ & read $S_mz^m$\\
p. & 74. & line 8 from top, for $w_{ik}$ & read $w'_{ik}$\\
p. & 89. & line 4 from bottom, for $s_\beta$ & read $s_\alpha$\\
p. & 93. & line 9 from bottom, for left & read right\\
--- & & line 8 from bottom, for right & read left\\
p. & 119. & line 15 from top, for (5) & read (6)\\
p. & 122. & line 15 from bottom, for $s_x, x$ & read for $s_x, x$\\
p. & 123. & line 2 from bottom, for (1) & read (3)\\
p. & 131. & line 4 from bottom, for $Q(q-1)$ & read $Q(q-2)$\\
p. & 133. & line 3 from bottom, for two & read two consecutive\\
p. & 157. & line 9 from top, for $p$ & read $p-1$\\
p. & 158. & line 6 from bottom, for $u-1$ & read $n-2$\\
\end{tabular}

\vspace{0.7em}
\noindent p. 159. Lines 6, 7 and 9, 10: the expressions
\[
\eta_1\frac{\partial f(\xi)}{\partial \xi_1}
+\eta_2\frac{\partial f(\xi)}{\partial \xi_2}
+\eta_3\frac{\partial f(\xi)}{\partial \xi_3}
\quad\text{and}\quad
\xi_1\frac{\partial f(\eta)}{\partial \eta}
+\xi_2\frac{\partial f(\eta)}{\partial \eta}
+\xi_3\frac{\partial f(\eta)}{\partial \eta}
\]
are to be interchanged.



% --- Continuation: source PDF pp. 017--038. ---
\clearpage


% English translation of Clebsch and Gordan, Theorie der Abelschen Functionen, 1866.
% Source witness checked against Google scan PDF theoriederabels00clebgoog.pdf, source PDF pp. 017--038.
% Source PDF p. 022 is blank/scan whitespace and is represented in the source witness and ledger only.

\begin{center}
{\Large\bfseries Contents.}
\end{center}

{\scriptsize
\begin{longtable}{@{}r@{\ }p{0.74\linewidth}@{\ }r@{}}
\tocsec{First Section.}{The algebraic integrals.}
\tocentry{1.}{Representation of the differential expressions to be investigated}{1}
\tocentry{2.}{Classification of the integrals}{5}
\tocentry{3.}{Discontinuities. Double and return points}{10}
\tocentry{4.}{Integrals of the first kind}{14}
\tocentry{5.}{Integrals of the third kind}{16}
\tocentry{6.}{Formation of the normal integrals of the third kind}{20}
\tocentry{7.}{Properties of the normal integrals of the third kind}{22}
\tocentry{8.}{Integrals of the second kind}{27}
\tocsec{Second Section.}{Abel's theorem.}
\tocentry{9.}{Systems of intersection points}{34}
\tocentry{10.}{Transformation of the sums of integrals which correspond to a system of intersection points}{36}
\tocentry{11.}{Jacobi's theorem}{39}
\tocentry{12.}{Abel's theorem}{44}
\tocentry{13.}{Special cases}{47}
\tocsec{Third Section.}{The one-valued transformations.}
\tocentry{14.}{Transformation of the algebraic integrals}{50}
\tocentry{15.}{Concept of the one-valued transformations. Preservation of the number $p$}{52}
\tocentry{16.}{Direct proof for the preservation of the number $p$}{54}
\tocentry{17.}{Proof of the auxiliary theorem}{60}
\tocentry{18.}{Lowering of the order of the curves. Normal curves}{63}
\tocentry{19.}{Treatment of the case $p=0$. Normal curve: a straight line}{67}
\tocentry{20.}{Treatment of the case $p=1$. Normal curve: a general curve of the third order}{69}
\tocentry{21.}{Treatment of the case $p=2$. Normal curve: a curve of the fourth order with a double point}{77}
\tocsec{Fourth Section.}{Periodicity.}
\tocentry{22.}{Branch points}{80}
\tocentry{23.}{Cycles and circuits}{82}
\tocentry{24.}{Reciprocity between two classes of circuits and between two classes of fundamental points}{85}
\tocentry{25.}{Integration along different paths. Periodicity moduli of the integrals of the first kind}{88}
\tocentry{26.}{Relations between the periodicity moduli of the integrals of the first kind}{90}
\tocentry{27.}{Form of the relation found and properties of its integral coefficients}{95}
\tocentry{28.}{Example for setting up the relation between the periodicity moduli of the integrals of the first kind}{99}
\tocentry{29.}{Transformation of the relation}{103}
\tocentry{30.}{Normal form of the integrals of the first kind}{107}
\tocentry{31.}{Execution of the transformation for the example treated in \S\ 28}{109}
\tocsec{Fifth Section.}{Interchange of parameter and argument.}
\tocentry{32.}{Closer determination of the integral of the third kind}{112}
\tocentry{33.}{Interchange of parameter and argument}{114}
\tocentry{34.}{Finite normal form of the integral of the third kind. Simplest form of the interchange theorem. Theorems which follow from it}{116}
\tocentry{35.}{Theorems on integrals of the second kind which follow from the interchange theorem}{119}
\tocentry{36.}{Reduction of the differential quotients of integrals of the third kind with respect to a parameter to integrals of the second and third kinds}{123}
\tocentry{37.}{Application of the methods used above for integrals of the third kind to algebraic functions and logarithms of such. Abel's theorem}{124}
\tocsec{Sixth Section.}{The inversion problem.}
\tocentry{38.}{Auxiliary theorem for judging periodic functions}{129}
\tocentry{39.}{Functions with several periods}{134}
\tocentry{40.}{Concept of Jacobi's inversion problem}{136}
\tocentry{41.}{Formulation of the inversion problem. Abelian functions and transcendentals}{138}
\tocentry{42.}{Reduction of the inversion problem to the investigation of the transcendental $T_{\xi\eta}$}{141}
\tocentry{43.}{Extended inversion problem}{143}
\tocentry{44.}{Periods of the functions of the extended inversion problem. Modification when integrals of the second kind occur}{148}
\tocsec{Seventh Section.}{The function $T_{\xi\eta}$.}
\tocentry{45.}{Decomposition of the function $T_{\xi\eta}$ into the difference of two functions, each of which depends only on $p+2$ points}{152}
\tocentry{46.}{Proof of the decomposability of the function $\frac12 T_{\xi\eta}(x)$ into two functions depending on only $p+1$ points}{155}
\tocentry{47.}{On the character of the special transcendental $T_{\xi\eta}(x)$}{160}
\tocentry{48.}{Theory of the adjoint transcendentals}{165}
\tocentry{49.}{The function $U$. The transcendental $T_{\xi\eta}$ expressed through it}{170}
\tocentry{50.}{Discontinuities of the function $U$. The function $V=e^{-U}$}{173}
\tocentry{51.}{Properties of the function $V=e^{-U}$}{175}
\tocsec{Eighth Section.}{Monodromy.}
\tocentry{52.}{Differential of an Abelian function. Value systems of the $v$ for which the system of the $x$ is undetermined}{181}
\tocentry{53.}{Proof that the function $V$ is a synektic function of the $v$}{187}
\tocentry{54.}{Periodicity moduli of $-U=\log V$}{190}
\tocentry{55.}{Series development}{193}
\tocentry{56.}{The arguments of the function $\Theta$}{195}
\tocentry{57.}{Solution of the inversion problem by $\Theta$-functions}{198}
\tocentry{58.}{An integral of the third kind and the double ratio $\left(\frac{\varphi}{\psi}\right)_{\xi}:\left(\frac{\varphi}{\psi}\right)_{\eta}$ represented by $\Theta$-functions}{200}
\tocentry{59.}{Quotients of alternating functions represented by $\Theta$-functions}{202}
\tocsec{Ninth Section.}{The function $\Theta$.}
\tocentry{60.}{Vanishing of the $\Theta$-function. Indeterminate case of the inversion problem}{206}
\tocentry{61.}{Special cases in which the $\Theta$-function vanishes}{211}
\tocentry{62.}{The inversion of Abel's theorem}{216}
\tocentry{63.}{Products of $\Theta$-functions reduced to other and to algebraic functions}{218}
\tocentry{64.}{On the representation of the algebraic expressions contained in the foregoing}{222}
\tocentry{65.}{Representation of quotients of $\Theta$-products by algebraic functions}{224}
\tocentry{66.}{Relations between $\Theta$-functions}{226}
\tocsec{Tenth Section.}{Division.}
\tocentry{67.}{Multiplication}{230}
\tocentry{68.}{Problem of general division}{231}
\tocentry{69.}{The problem of special division. Reduction of general division to special division}{235}
\tocentry{70.}{Solution of the general division equation}{239}
\tocentry{71.}{Reduction of the special division equation to the case in which $m$ is a power of a prime number}{241}
\tocentry{72.}{Solution of the special division equation for a prime number and powers of such a number}{244}
\tocentry{73.}{Algebraically representable $\Theta$-quotients}{248}
\tocentry{74.}{Algebraic determination of $\dfrac{\Theta\left(\frac{P}{m}\right)}{\Theta(0)}$}{256}
\tocentry{75.}{Determination of $\dfrac{\Theta\left(\frac{P}{2}\right)}{\Theta(0)}$. Curves of order $(n-3)$ which touch in $p$ points and pass through the double and return points}{259}
\tocentry{76.}{Reduction of the division problems for $m=2$}{265}
\tocsec{Eleventh Section.}{The extended inversion problem.}
\tocentry{77.}{Reduction of the extended inversion problem to a function $U^{(q)}$ depending on only $(p+1)$ points}{270}
\tocentry{78.}{Discontinuities of $U^{(q)}$. The synektic function $V^{(q)}=e^{-U^{(q)}}$}{274}
\tocentry{79.}{Series development of the function $V^{(q)}$ from its periodicity properties}{277}
\tocentry{80.}{Determination of the coefficients $A$ of the function $V^{(q)}$}{281}
\tocentry{81.}{The function $\Theta^{(q)}$, with suitable arrangement, is a function of its arguments which is even or odd with $q$}{287}
\tocentry{82.}{Division for the extended inversion problem}{290}
\tocentry{83.}{Division for $m=2$}{294}
\tocentry{84.}{Influence of return points}{296}
\tocsec{Twelfth Section.}{The infinitely many forms of the function $\Theta$.}
\tocentry{85.}{Canonical determinants}{300}
\tocentry{86.}{Reduction of the canonical determinants to simple ones}{304}
\tocentry{87.}{Transformation of the integrals of the first and third kinds when the system of normal periods is changed}{308}
\tocentry{88.}{Transformation of the $\Theta$-function by means of canonical substitutions}{311}
\tocentry{89.}{Contact curves which occur in the lower limits}{314}
\tocentry{90.}{Reduction of the constant occurring in the transformation formula for $\Theta$ to a numerical one}{318}
\tocentry{91.}{The numerical constant in the composition of transformations}{321}
\tocentry{92.}{Determination of the numerical constant for the simple transformations}{326}
\tocentry{93.}{Calculation of the $a_{ik}$ for the different kinds of $\Theta$-functions. Direct establishment of the series developments}{331}
\end{longtable}}

\begin{center}
{\large\bfseries First Section.}\par
{\Large\bfseries The algebraic integrals.}
\end{center}

\begin{center}
{\bfseries \S\ 1. Representation of the differential expressions to be investigated.}
\end{center}

Let us denote by
\[
(1)\ldots\quad f(x_1,x_2,x_3)=0
\]
the equation of an algebraic curve of order $n$, referred to triangular coordinates. Because of the two equations which follow from it,
\[
\frac{\partial f}{\partial x_1}x_1+\frac{\partial f}{\partial x_2}x_2+\frac{\partial f}{\partial x_3}x_3=nf=0
\]
and
\[
\frac{\partial f}{\partial x_1}dx_1+\frac{\partial f}{\partial x_2}dx_2+\frac{\partial f}{\partial x_3}dx_3=0,
\]
one has the proportion
\[
x_2dx_3-x_3dx_2:x_3dx_1-x_1dx_3:x_1dx_2-x_2dx_1
=\frac{\partial f}{\partial x_1}:\frac{\partial f}{\partial x_2}:\frac{\partial f}{\partial x_3}.
\]
Therefore, in the differential expression\footnote{This symmetric representation of such differentials is due to Mr. Aronhold (cf. the monthly reports of the Berlin Academy, session of 25 April 1861).}
\[
\frac{\Sigma\pm c_1x_2dx_3}{c_1\frac{\partial f}{\partial x_1}+c_2\frac{\partial f}{\partial x_2}+c_3\frac{\partial f}{\partial x_3}},
\]
the terms multiplied by the individual $c$ in numerator and denominator are proportional, and the whole expression is independent of the $c$.

If one further multiplies it by a function $\Phi$ which is rational and homogeneous in the $x$, and whose numerator is of order $n-3$ higher than the denominator, then
\[
\frac{\Phi\,\Sigma\pm c_1x_2dx_3}{c_1\frac{\partial f}{\partial x_1}+c_2\frac{\partial f}{\partial x_2}+c_3\frac{\partial f}{\partial x_3}}
\]
is a differential expression which depends on only one variable, for it is homogeneous of order zero in the $x$, does not change if one simultaneously multiplies the $x$ by a common factor, and hence depends only on the two ratios of the quantities $x$. But between these there exists equation (1); hence these two ratios are algebraically connected, and one can be regarded as a function of the other. Thus in the end only one variable remains, while the differential expression, because equation (1) has been used, indeed becomes irrational; and one has before one an irrational algebraic differential whose irrationalities are determined exclusively by equation (1). By integrating this differential from one point of the curve (value system of the $x$) to another, one arrives at the irrational algebraic integral
\[
(2)\ldots\quad V=\int\frac{\Phi\cdot\Sigma\pm c_1x_2dx_3}{c_1\frac{\partial f}{\partial x_1}+c_2\frac{\partial f}{\partial x_2}+c_3\frac{\partial f}{\partial x_3}},
\]
which forms the basis of the following investigations. This integral is to be taken in such a way that in the course of integration the variables $x$ are always connected by equation (1).

The integral (2) is the most general integral of a rational function of $x_1,x_2,x_3$ which one can form. Namely, suppose some differential expression
\[
dV=A_1dx_1+A_2dx_2+A_3dx_3,
\]
in which the $A$ denote homogeneous functions of order $\mu$, is to pass, by means of $f=0$, into an expression with one variable. Then it may depend only on the ratios of the $x$; it may therefore not change when $kx_1,kx_2,kx_3$ are put in place of $x_1,x_2,x_3$. By this one obtains
\[
dV=[(A_1x_1+A_2x_2+A_3x_3)dk+k(A_1dx_1+A_2dx_2+A_3dx_3)]k^\mu.
\]
Therefore first $k^{\mu+1}=1$, $\mu=-1$, must hold, that is, the $A$ must be of dimension $-1$; but then also
\[
A_1x_1+A_2x_2+A_3x_3=nM\cdot f,
\]
or
\[
\left(A_1-M\frac{\partial f}{\partial x_1}\right)x_1+
\left(A_2-M\frac{\partial f}{\partial x_2}\right)x_2+
\left(A_3-M\frac{\partial f}{\partial x_3}\right)x_3=0.
\]
This equation is satisfied in the most general way if one sets
\[
A_1=M\frac{\partial f}{\partial x_1}+C_2x_3-C_3x_2,
\]
\[
A_2=M\frac{\partial f}{\partial x_2}+C_3x_1-C_1x_3,
\]
\[
A_3=M\frac{\partial f}{\partial x_3}+C_1x_2-C_2x_1.
\]
And our differential expression, omitting the vanishing $df$, passes into
\[
dV=\Sigma\pm C_1x_2dx_3,
\]
where now the $C$ are of dimension $-2$. These functions can have a common factor $N$. Let
\[
C_i=Nc_i
\]
and
\[
\Phi=N\left(c_1\frac{\partial f}{\partial x_1}+c_2\frac{\partial f}{\partial x_2}+c_3\frac{\partial f}{\partial x_3}\right);
\]
then immediately $dV$ passes into the form established by us:
\[
V=\int\frac{\Phi\,\Sigma\pm c_1x_2dx_3}{c_1\frac{\partial f}{\partial x_1}+c_2\frac{\partial f}{\partial x_2}+c_3\frac{\partial f}{\partial x_3}},
\]
where now $\Phi$ is of dimension $(n-3)$. The $c$, insofar as they occur here outside $\Phi$, have no influence on the value of the integral and can therefore, outside the function $\Phi$, be replaced by entirely arbitrarily chosen quantities. We shall in general regard them as constants, that is, as coordinates of a certain point $c$.

The integral $V$ can be brought in many ways into the form of an ordinary integral with one variable. We shall first do this here by means of a linear transformation. Let $a,b$ be two arbitrarily chosen fixed points which, together with $c$, determine a new triangle. Introduce three new variables $x,s_x,t$ with the help of the formulas
\[
(3)\ldots\quad
\begin{aligned}
\mu x_1&=a_1t+b_1x+c_1s_x,\\
\mu x_2&=a_2t+b_2x+c_2s_x,\\
\mu x_3&=a_3t+b_3x+c_3s_x,
\end{aligned}
\]
where $\mu$ is an entirely arbitrary quantity. Then $t,x,s_x$ are the triangular coordinates of the point $x$ referred to the new triangle. If one substitutes these expressions for the $x_i$ in $V$, then $\mu$ cancels out completely according to what precedes; namely, one has
\[
\mu^n f(x_1,x_2,x_3)=F(t,x,s_x),
\]
\[
\mu^{n-3}\Phi(x_1,x_2,x_3)=\Psi(t,x,s_x),
\]
\[
\begin{aligned}
c_1\frac{\partial f}{\partial x_1}+c_2\frac{\partial f}{\partial x_2}+c_3\frac{\partial f}{\partial x_3}
&=\mu\left(\frac{\partial f}{\partial x_1}\frac{\partial x_1}{\partial s_x}+\frac{\partial f}{\partial x_2}\frac{\partial x_2}{\partial s_x}+\frac{\partial f}{\partial x_3}\frac{\partial x_3}{\partial s_x}\right)\\
&=\mu\frac{\partial f}{\partial s_x}=\frac{1}{\mu^{n-1}}\frac{\partial F}{\partial s_x};
\end{aligned}
\]
and finally
\[
\begin{aligned}
\mu^2\Sigma\pm c_1x_2dx_3&=\left|\begin{array}{ccc}
c_1&a_1t+b_1x+c_1s_x&a_1dt+b_1dx+c_1ds_x\\
c_2&a_2t+b_2x+c_2s_x&a_2dt+b_2dx+c_2ds_x\\
c_3&a_3t+b_3x+c_3s_x&a_3dt+b_3dx+c_3ds_x
\end{array}\right|\\
&=\Sigma\pm a_1b_2c_3\cdot(t\,dx-x\,dt);
\end{aligned}
\]
hence, if one lets $\Sigma\pm a_1b_2c_3$ enter into $\Psi$,
\[
V=\int\frac{\Psi(t\,dx-x\,dt)}{\frac{\partial F}{\partial s_x}},
\]
while between $s_x,x,t$ there exists the equation $F=0$, which, like the former one, is of degree $n$.

The new expression for $V$ is also homogeneous of order zero; hence one may set $t=1$ without impairing generality, and obtains
\[
(4)\ldots\quad V=\int\frac{\Psi\,dx}{\frac{\partial F}{\partial s_x}},
\]
where $F=0$ is now an equation between $x,s_x$, which defines $s_x$ as a function of $x$ and rises to dimension $n$.

The geometrical meaning of $x,s_x$ for $t=1$ is very simple. From (3) the equations follow:
\[
0=\Sigma\pm a_1b_2x_3+s_x\Sigma\pm c_1b_2x_3,
\]
\[
0=\Sigma\pm a_1c_2x_3+x\Sigma\pm b_1c_2x_3.
\]
These are the equations of two pencils of rays, the first of which has its vertex at the intersection of the two lines
\[
\Sigma\pm a_1b_2x_3=0,
\qquad
\Sigma\pm c_1b_2x_3=0,
\]
that is, at $b$, while the other has its vertex at the intersection of the lines
\[
\Sigma\pm a_1c_2x_3=0,
\qquad
\Sigma\pm b_1c_2x_3=0,
\]
that is, at $c$; their parameters are respectively $s_x$ and $x$.

We shall use the integral $V$, according to circumstances, in the form (2) or (4). Let us note only that of course $\Psi$ can be any rational function whatever of $s_x$ and $x$, and that therefore (4) also embraces all algebraic integrals of rational functions of $s_x$ and $x$ in which $s_x,x$ are connected by the equation
\[
F(s_x,x)=0.
\]

\begin{center}
{\bfseries \S\ 2. Classification of the integrals.}
\end{center}

The integral $V$ contains the fractional function $\Psi$, whose denominator may be of order $m$, and whose numerator then rises to order $(m+n-3)$\footnote{The difference between the orders in numerator and denominator can apparently become larger, by the denominator containing a power of $t$ as factor which is set equal to $1$. This, however, can evidently always be avoided by an altered choice of the points $c,b$, whose joining line is $t=0$.}. Set
\[
\Psi=\frac{M}{N}.
\]
Then we can apply to this fraction a kind of partial fraction decomposition, and thereby decompose the integral $V$ into a number of simpler integrals.

The curves $F=0$, $N=0$ cut one another in $mn$ points. From the point $b$ we draw $mn$ rays through these points, to which just as many values of $x$ correspond. The product of the equations of all these rays is an equation
\[
X=0
\]
of degree $mn$. It arises when one seeks the values of $x$ which, with corresponding values of $s_x$, satisfy the equations $F=0$, $N=0$; that is, when one eliminates $s_x$ from the equations $F=0$, $N=0$. Let, arranged according to $s_x$,
\[
F=f_0s_x^n+f_1s_x^{n-1}+\cdots+f_n,
\]
\[
N=n_0s_x^m+n_1s_x^{m-1}+\cdots+n_m.
\]
Then $X$ is the determinant of $m+n$ rows:
\[
X=\left|\begin{array}{r|cccc}
1)&f_0&f_1&f_2&\cdots\\
2)&0&f_0&f_1&\cdots\\
3)&0&0&f_0&\cdots\\
&\cdots&\cdots&\cdots&\\
m+1)&n_0&n_1&n_2&\cdots\\
m+2)&0&n_0&n_1&\cdots\\
m+3)&0&0&n_0&\cdots\\
&\cdots&\cdots&\cdots&
\end{array}\right|.
\]
If one imagines the vertical columns of $X$ multiplied successively by
\[
s_x^{m+n-1},\quad s_x^{m+n-2},\ldots,1,
\]
and adds all columns so multiplied to the last one, then the terms of the last vertical column become
\[
s_x^{m-1}F,\\quad s_x^{m-2}F,\ldots,F,\quad s_x^{n-1}N,\quad s_x^{n-2}N,\ldots,N.
\]
Thus, arranging according to this column, one obtains $X$ in the form
\[
(1)\ldots\quad X=AF+BN,
\]
where $A$ and $B$ are determinants which agree completely with $X$, except that the last vertical column for $A$ is
\[
s_x^{m-1},\ s_x^{m-2},\ldots,1,\quad 0,\quad 0,\ldots,0,
\]
but for $B$ is
\[
0,\quad 0,\ldots,0,\quad s_x^{n-1},\ s_x^{n-2},\ldots,1.
\]
One can therefore multiply the left sides of the equations of two other curves $A=0$, $B=0$ by the curves of orders $n$ and $m$, $F=0$, $N=0$, so that the sum of the products gives the equation $X=0$ as result. The geometrical meaning of this is evident. The straight lines of the pencil $X=0$ cut $F=0$ in a further $mn(n-1)$ points which do not lie on $N=0$, and these lie on a curve $B=0$ of order $m(n-1)$. Likewise those straight lines cut the curve $N=0$ in a further $mn(m-1)$ points which do not lie on $F=0$, and these lie on a curve $A=0$ of order $n(m-1)$.

The curves $A=0$, $B=0$ are not completely determined by equation (1), for that equation remains unchanged if one replaces $A$ and $B$, respectively, by
\[
A+CN,\qquad B-CF,
\]
where $C$ denotes an arbitrary curve of order $(mn-m-n)$. But it is very important to note that the special curves $A=0$, $B=0$ obtained above have special important properties. Namely, $A$ contains the quantity $s_x$ only to the $(m-1)$st, and $B$ contains it only to the $(n-1)$st power. If one therefore imagines $A,B$ made homogeneous again by introducing the variable $t$ which above was set equal to $1$, then all terms in $A$ and $B$ are, in $s_x$ and $x$, at least of order $(n-1)(m-1)$; the two curves $A$ and $B$ are thus here determined so that both still have at the point $c$ an $(n-1)(m-1)$-fold point.

Thus multiply the fraction
\[
\Psi=\frac{M}{N}
\]
in numerator and denominator by $B$. Then because of the equation $F=0$ one may also put $X$ for the denominator, and one obtains
\[
\Psi=\frac{MB}{X},
\]
where now the denominator is a function of order $mn$ of $x$ alone, while the numerator is a function of order $(mn+n-3)$ which, however, because of the properties of $B$, contains the quantity $s_x$ only up to the power $(m+2n-4)$. Thus $MB$ is of the form
\[
MB=s_x^{m+2n-4}X_{(m-1)(n-1)}+s_x^{m+2n-5}X_{(m-1)(n-1)+1}+\cdots+X_{mn+n-3},
\]
where the $X$ are functions of $x$ alone whose order is indicated by the appended index. If in the same way $C$ is a function of order $(mn-3)$ which contains $s_x$ only to the power $(m+n-4)$,
\[
C=s_x^{m+n-4}\xi_{(m-1)(n-1)}+s_x^{m+n-3}\xi_{(m-1)(n-1)+1}+\cdots+\xi_{mn-3},
\]
then one can determine the functions $\xi$ so that
\[
MB+CF
\]
contains only powers of $s_x$ which lie below the $n$th. Namely, if only $f_0$ is not equal to zero, which always occurs when the point $c$ ($t=0$, $x=0$) does not lie on the curve itself, then the following equations must hold:
\[
0=X_{(m-1)(n-1)}+\xi_{(m-1)(n-1)}f_0,
\]
\[
0=X_{(m-1)(n-1)+1}+\xi_{(m-1)(n-1)+1}f_0+\xi_{(m-1)(n-1)}f_1,
\]
\[
0=X_{(m-1)(n-1)+2}+\xi_{(m-1)(n-1)+2}f_0+\xi_{(m-1)(n-1)+1}f_1+\xi_{(m-1)(n-1)}f_2,
\]
\[
\cdots
\]
\[
0=X_{mn-3}+\xi_{mn-3}f_0+\xi_{mn-4}f_1+\cdots+\xi_{mn-n-3}f_n,
\]
whereby the functions $\xi$, since $f_0$ is a non-zero constant, are completely and uniquely determined.

After this determination of the $\xi$, the highest power of $s_x$ in $MB+CF$ becomes the $(n-1)$st, and so
\[
MB+CF=X'_{mn-2}s_x^{n-1}+X'_{mn-1}s_x^{n-2}+\cdots+X'_{mn+n-3},
\]
where the $X'$ are again functions of $x$ alone whose degree is indicated by the index. Since moreover
\[
\frac{\partial F}{\partial s_x}=nf_0s_x^{n-1}+(n-1)f_1s_x^{n-2}+\cdots,
\]
one can destroy the first term of the preceding series by subtracting
\[
\frac{1}{nf_0}\cdot X'_{mn-2}\frac{\partial F}{\partial s_x}=D\frac{\partial F}{\partial s_x},
\]
and obtains an expression of the form
\[
MB+CF-D\frac{\partial F}{\partial s_x}=X''_{mn-1}s_x^{n-2}+X''_{mn}s_x^{n-3}+\cdots+X''_{mn+n-3}.
\]
With regard to the equation $F=0$, one therefore now has
\[
\frac{M}{N}=\frac{MB-D\frac{\partial F}{\partial s_x}}{X}+\frac{D}{X}\frac{\partial F}{\partial s_x}
\]
\[
=\frac{1}{X}\left\{X''_{mn-1}s_x^{n-2}+X''_{mn}s_x^{n-3}+\cdots+X''_{mn+n-3}\right\}+\frac{D}{X}\frac{\partial F}{\partial s_x}.
\]
Our integral
\[
V=\int\frac{M}{N}\frac{dx}{F'_{s_x}}
\]
accordingly decomposes into two parts. The first, which comes from the last term in the above expression for $\frac{M}{N}$, is
\[
\int\frac{D}{X}\,dx,
\]
and hence the integral of a rational function; it can therefore be integrated by logarithms and rational functions.

It is enough to restrict $V$ to the part which comes from the remaining terms of the expression for $\frac{M}{N}$.

The first part of $\frac{M}{N}$ can first be decomposed by division into an entire function which, with respect to $s_x$ and $x$, rises to dimension $(n-3)$, and into a remainder which with respect to $x$ becomes of order $(mn-1)$ and must be treated by partial fraction decomposition. In this last process, in which we first suppose the roots of $X=0$ to be distinct, there arise $mn$ different terms whose numerators are entire functions of degree $(n-2)$ in $s_x$, while the denominator always has the form $x-a$.

We can therefore reduce every integral of the form considered here to aggregates of the following kinds of integrals:
\begin{enumerate}
\item Integrals of rational functions of $x$.
\item Integrals of the form
\[
\int\frac{Q\,dx}{F'_{s_x}},
\]
where $Q$ is an entire function which rises, in $s_x$ and $x$, to dimension $(n-3)$.
\item Integrals of the form
\[
\int\frac{Q\,dx}{(x-a)F'_{s_x}},
\]
where $Q$ rises to dimension $(n-2)$.
\end{enumerate}
Finally, if in the above the equation $X=0$ had equal roots, there can also occur integrals which have higher powers of $x-a$ in the denominator. These, however, can be allowed to arise by a limiting transition from those cited under 3). According to the theory of partial fractions, the corresponding terms of $\frac{M}{N}$, if $\hat P$ denotes the denominator left over after removal of the equal factors $(x-a)$ from $X$, are of the form
\[
\frac{\left(\frac{R}{P}\right)_{x=a}}{(x-a)^\mu}
+\frac{1}{1}\frac{\left(\frac{\partial}{\partial x}\frac{R}{P}\right)_{x=a}}{(x-a)^{\mu-1}}
+\frac{1}{1\cdot2}\frac{\left(\frac{\partial^2}{\partial x^2}\frac{R}{P}\right)_{x=a}}{(x-a)^{\mu-2}}+\cdots.
\]
The integrals arising from this are obtained, therefore, when one differentiates certain integrals contained in form 3) with respect to the parameter $a$\footnote{This kind of reduction to integral 3) is given here only for preliminary orientation. Later (cf. Section 5) an entirely different path will be shown for the representation of such integrals. Likewise we shall learn to express the logarithmic and rational functions by integrals of the forms 1)--3).}.

The integrals contained under the forms 2) and 3) are now to be examined more closely, and in particular with respect to the discontinuities which can occur in them and which determine their character.

\begin{center}
{\bfseries \S\ 3. Discontinuities. Double and return points.}
\end{center}

The integrals of form 2), expressed in $x$ and $s_x$, are
\[
V=\int\frac{Q\,dx}{F'_{s_x}},
\]
where $Q$ rises to dimension $(n-3)$. In the $x_i$ they therefore have the form
\[
V=\int\frac{Q\cdot\Sigma\pm c_1x_2dx_3}{c_1\frac{\partial f}{\partial x_1}+c_2\frac{\partial f}{\partial x_2}+c_3\frac{\partial f}{\partial x_3}},
\]
where $Q$ is a homogeneous function of order $(n-3)$. The function to be integrated can here become discontinuous only when the denominator vanishes, that is, when $F'(s_x)=0$ or
\[
c_1\frac{\partial f}{\partial x_1}+c_2\frac{\partial f}{\partial x_2}+c_3\frac{\partial f}{\partial x_3}=0.
\]
This is the equation of the polar of the point $c$ with respect to the curve $f=0$; the two curves cut one another in $n(n-1)$ points, at which the tangents drawn from the arbitrarily chosen point $c$ to the curve $f=0$ touch it. At one of these points let $s_x^\circ,x^\circ$ be the values of $s_x$ and $x$. Then one has
\[
F(x^\circ,s_x^\circ)=0,\qquad \frac{\partial F(x^\circ,s_x^\circ)}{\partial s_x^\circ}=0;
\]
hence, near such a point, when $s_x-s_x^\circ$ and $x-x^\circ$ are very small quantities, and when, for brevity, $F^0$ is written for $F(s_x^\circ,x^\circ)$,
\[
F(s_x,x)=0=\frac{\partial F^0}{\partial x^0}(x-x^0)+\frac12\frac{\partial^2F^0}{\partial s_x^{0\,2}}(s_x-s_x^0)^2+\cdots
\]
and
\[
\frac{\partial F}{\partial s_x}=\frac{\partial^2F^0}{\partial s_x^{0\,2}}(s_x-s_x^0)+\cdots,
\]
so that, by the first equation, $s_x-s_x^0$ becomes of the same order as $\sqrt{x-x^0}$, and thus also $\frac{\partial F}{\partial s}$, the denominator of $V$, is proportional to $\sqrt{x-x^0}$. Thus the integral $V$, if the integration moves near these points, becomes proportional to
\[
\int\frac{dx}{\sqrt{x-x^0}}=2\sqrt{x-x^0},
\]
that is, it remains finite, although the function to be integrated becomes infinite.

For the point of contact of a tangent drawn from $c$ to the curve, the integral $V$ therefore does not become discontinuous.

But if the curve $f=0$ has some of the ordinary singularities (and we shall restrict ourselves to these)\footnote{It may be noted that an $r$-fold point can always, in all relations, be thought of as replaced by $\frac{r\cdot r-1}{2}$ double points, provided only that all its tangents are distinct. The function $f$ in hyperelliptic functions, for example, can always be reduced to this case, so that these make no further difficulties. But special considerations are required when several of the tangents of a multiple point coincide.}, if it has double and return points, then these considerations require some modifications. As is known, each double point diminishes by two the number of tangents which can be drawn from a point to the curve, each return point by three. Let $d$ denote the number of double points and $r$ the number of return points; then the number of proper tangents which can be drawn from $c$ to the curve, or the number of movable intersection points of curve and polar, is
\[
n(n-1)-2d-3r;
\]
the polar of every point, hence also that of $c$, passes through every double or return point of the curve, and touches the tangent of each of the latter, so that among the intersection points of the polar with the curve, two are united and fixed at each double point, and three at each return point.

A double point is known to be characterized by the fact that in it
\[
\frac{\partial f}{\partial x_1}=0,
\qquad \frac{\partial f}{\partial x_2}=0,
\qquad \frac{\partial f}{\partial x_3}=0,
\]
or, if one uses the equation $F(s_x,x)=0$, that simultaneously
\[
F=0,\qquad \frac{\partial F}{\partial s_x}=0,\qquad \frac{\partial F}{\partial x}=0.
\]
If one therefore expands as above, near a double point $s_x^\circ,x^\circ$, in powers of $s_x-s_x^\circ$, $x-x^\circ$, one has:
\[
\begin{aligned}
F(s_x,x)=0={}&\frac12\left\{\frac{\partial^2F^\circ}{\partial s_x^{\circ\,2}}(s_x-s_x^\circ)^2
+2\frac{\partial^2F^\circ}{\partial s_x^\circ\partial x^\circ}(s_x-s_x^\circ)(x-x^\circ)
+\frac{\partial^2F^\circ}{\partial x^{\circ\,2}}(x-x^\circ)^2\right\}+\cdots,
\end{aligned}
\]
\[
\frac{\partial F}{\partial s_x}=\frac{\partial^2F^\circ}{\partial s_x^{\circ\,2}}(s_x-s_x^\circ)+\frac{\partial^2F^\circ}{\partial x^\circ\partial s_x^\circ}(x-x^\circ)+\cdots.
\]
From the first of these equations one sees that here again $s_x-s_x^0$ becomes proportional in the limit to $x-x^0$ itself, not as before to $\sqrt{x-x^0}$, and indeed for $\frac{s_x-s_x^0}{x-x^0}$ two values are obtained from the quadratic equation
\[
\left(\frac{s_x-s_x^0}{x-x^0}\right)^2\frac{\partial^2F^0}{\partial s_x^{0\,2}}
+2\frac{s_x-s_x^0}{x-x^0}\frac{\partial^2F^0}{\partial s_x^0\partial x^0}
+\frac{\partial^2F^0}{\partial x^{0\,2}}=0.
\]
Therefore $\frac{\partial F}{\partial s_x}$ also becomes proportional to $x-x^\circ$ near the double point, and the integral $V$, when the integration moves near $s_x^\circ,x^\circ$, becomes proportional to
\[
\int\frac{dx}{x-x^0}=\log(x-x^0),
\]
that is, the integral becomes logarithmically infinite. Excepted from this is the case in which the numerator $Q$ itself vanishes for $s_x=s_x^\circ$, $x=x^\circ$. If this is the case, then near these values $Q$ does not approach a finite constant value, as otherwise, but one has
\[
Q=\frac{\partial Q^\circ}{\partial s_x}(s_x-s_x^\circ)+\frac{\partial Q^\circ}{\partial x}(x-x^\circ)+\cdots;
\]
hence $Q$ is also proportional to $x-x^\circ$. In this case $\dfrac{Q}{\frac{\partial F}{\partial s_x}}$ thus approaches a finite limit, and the integral $V$ remains perfectly finite near the double point.

If $s_x^\circ,x^\circ$ is a return point, then in addition to the conditions of the double point there is the condition that the preceding quadratic equation have two equal roots, that is, that
\[
\frac{\partial^2F^0}{\partial s_x^{0\,2}}\cdot\frac{\partial^2F^0}{\partial x^{0\,2}}-
\left(\frac{\partial^2F^0}{\partial s_x^0\partial x^0}\right)^2=0.
\]
Then one can determine two quantities $a,b$ so that
\[
\frac{\partial^2F^0}{\partial s_x^{0\,2}}=a^2,
\qquad
\frac{\partial^2F^0}{\partial x^{0\,2}}=b^2,
\qquad
\frac{\partial^2F^0}{\partial s_x^0\partial x^0}=ab.
\]
Taking into account in the expansion of $F(s_x,x)$ also the terms of third order, one has
\[
\begin{aligned}
F(s_x,x)=0={}&\frac12\left[a(s_x-s_x^0)+b(x-x^0)\right]^2\\
&+\frac16\left\{\frac{\partial^3F^0}{\partial s_x^{0\,3}}(s_x-s_x^0)^3
+3\frac{\partial^3F^0}{\partial s_x^{0\,2}\partial x^0}(s_x-s_x^0)^2(x-x^0)\right.\\
&\left.\quad+3\frac{\partial^3F^0}{\partial s_x^0\partial x^{0\,2}}(s_x-s_x^0)(x-x^0)^2
+\frac{\partial^3F^0}{\partial x^{0\,3}}(x-x^0)^3\right\}+\cdots,
\end{aligned}
\]
while
\[
\frac{\partial F}{\partial s_x}=a\left[a(s_x-s_x^0)+b(x-x^0)\right]+\cdots.
\]
The first equation shows that the combination
\[
a(s_x-s_x^0)+b(x-x^0)
\]
is of higher than first order, namely of order $\frac32$; hence in the terms of third order of $F(s_x,x)=0$ we can use the equation $a(s_x-s_x^0)+b(x-x^0)=0$, or set
\[
s_x-s_x^0=-\frac{b}{a}(x-x^0),
\]
whereby the expression
\[
a(s_x-s_x^0)+b(x-x^0)
\]
becomes proportional to the $\frac32$ power of $x-x^0$. Thus also $\frac{\partial F}{\partial s_x}$, which differs from this expression only by the factor $a$, is proportional to $(x-x^0)^{\frac32}$, and the integral $V$, as long as the integration moves near the return point, since the coefficient of $\frac{1}{x-x^0}$ under the integral sign vanishes, is of the form
\[
\int\left(\frac{A}{\sqrt{(x-x^0)^3}}+\frac{B}{\sqrt{x-x^0}}\right)dx=
\frac{2}{3}\frac{A}{\sqrt{x-x^0}}+2B\sqrt{x-x^0},
\]
hence algebraically infinite, and indeed proportional to the $-\frac12$ power of $x-x^0$. Excepted here again is the case in which the numerator $Q$ found under the integral sign itself vanishes for the return point, where then $V$ is proportional to
\[
\int\frac{dx}{\sqrt{x-x^0}}=2\sqrt{x-x^0},
\]
and thus remains finite. One sees that the return point behaves like the union of a double point with a point of contact, so that after the double point has been removed the properties of the point of contact still remain.

\begin{center}
{\bfseries \S\ 4. Integrals of the first kind.}
\end{center}

From what has been said one sees that there can be integrals which never become infinite at all. They must have the property that their numerator $Q$ or $\Theta$ vanishes for all double and return points. This is in fact always possible if $n>2$; for $n=1$ or $n=2$ there are no integrals of the kind considered here at all, since there is no longer an entire function $Q$ of order $(n-3)$.

The highest number of double and return points which a proper curve of order $n$ can possess is
\[
\frac{n-1\cdot n-2}{2}.
\]
For if one more double or return point were present, then through the
\[
\frac{n-1\cdot n-2}{2}+1
\]
points of this kind and through another $2n-3$ points chosen on the curve, hence in all through
\[
\frac{(n-1)(n+2)}{2}
\]
points, one could draw a curve of order $(n-1)$, which would thereby be just determined. This curve would then cut the given one, since each double or return point is to be counted as a double intersection point, in
\[
(n-1)(n-2)+2+2n-3=n(n-1)+1
\]
points, hence in one intersection point more than can occur at all for a curve of order $n$ and a curve of order $(n-1)$. Therefore a part of both curves would have to coincide, and thus the curve of order $n$ itself would decompose into parts and cease to be a proper curve of order $n$.

Thus $d+r$ is at most equal to
\[
\frac{n-1\cdot n-2}{2}.
\]
The latter number, however, agrees exactly with the number of coefficients in the curve of order $(n-3)$, $\Theta=0$, which forms the numerator under the integral sign of $V$. The curve $\Theta=0$ can therefore be determined so that it passes through the $d+r$ double and return points, that is, so that it satisfies the $d+r$ linear equations for its coefficients which state that $\Theta$ vanishes for the coordinates of one of these points. For the highest number of double and return points, indeed, one obtains as many homogeneous equations as unknowns, and these must all vanish; for
\[
d+r=\frac{n-1\cdot n-2}{2}
\]
the number of everywhere finite integrals is therefore zero. In general, however, the number of coefficients still left undetermined in $\Theta$ is
\[
p=\frac{n-1\cdot n-2}{2}-d-r,
\]
and the most general everywhere finite integral is therefore composed of just as many special integrals; or, as one may express it, there are
\[
p=\frac{n-1\cdot n-2}{2}-d-r
\]
different integrals of the first kind, if by this latter expression one denotes integrals which remain finite everywhere. Such integrals we shall in what follows generally call integrals $I$.

All integrals of the form
\[
V=\int\frac{\Theta\Sigma\pm c_1x_2dx_3}{c_1\frac{\partial f}{\partial x_1}+c_2\frac{\partial f}{\partial x_2}+c_3\frac{\partial f}{\partial x_3}}
=\int\frac{Q\,dx}{\frac{\partial F}{\partial s_x}}
\]
can be composed from $\frac{n-1\cdot n-2}{2}$ special integrals. As such special integrals, besides the everywhere finite integrals just treated, one can choose those which become infinite only at one each of the $d+r$ exceptional points. Indeed, let $H=0$ be an arbitrary curve of order $(n-3)$, and let $\Theta_1=0$, $\Theta_2=0\ldots\Theta_{d+r}=0$ be curves of order $(n-3)$ which pass through the $d+r$ exceptional points, in each case with one excepted. Then one can always determine constants $\alpha_1,\alpha_2\ldots\alpha_{d+r}$ so that the curve
\[
(1)\ldots\quad H-\alpha_1\Theta_1-\alpha_2\Theta_2\cdots-\alpha_{d+r}\Theta_{d+r}=\Theta=0
\]
passes through all $d+r$ exceptional points, and therefore
\[
\begin{aligned}
\int\frac{H\Sigma\pm c_1x_2dx_3}{c_1\frac{\partial f}{\partial x_1}+c_2\frac{\partial f}{\partial x_2}+c_3\frac{\partial f}{\partial x_3}}
={}&\int\frac{\Theta\Sigma\pm c_1x_2dx_3}{c_1\frac{\partial f}{\partial x_1}+c_2\frac{\partial f}{\partial x_2}+c_3\frac{\partial f}{\partial x_3}}\\
&+\Sigma_i\alpha_i\int\frac{\Theta_i\Sigma\pm c_1x_2dx_3}{c_1\frac{\partial f}{\partial x_1}+c_2\frac{\partial f}{\partial x_2}+c_3\frac{\partial f}{\partial x_3}}
\end{aligned}
\]
is composed in the indicated way. Namely, if one denotes by upper indices the values which the functions $H,\Theta_1,\Theta_2\ldots$ assume for the exceptional points, then, by setting up the conditions that the curve (1) pass through all exceptional points, one has
\[
H^{(1)}-\alpha_1\Theta_1^{(1)}=0,
\]
\[
H^{(2)}-\alpha_2\Theta_2^{(2)}=0,
\]
\[
\cdots
\]
\[
H^{(d+r)}-\alpha_{d+r}\Theta_{d+r}^{(d+r)}=0.
\]
All the remaining $\Theta_i^{(k)}$ vanish by assumption; the preceding equations, however, determine the $\alpha$ completely and uniquely. The representation in question is therefore possible in only one way once the curves $\Theta_i$ have been chosen. We shall return to these at a later place.

\begin{center}
{\bfseries \S\ 5. Integrals of the third kind.}
\end{center}

Instead of the integrals listed in \S\ 2 under 3),
\[
V=\int\frac{Q\,dx}{(x-a)\frac{\partial F}{\partial s_x}},
\]
where $Q$ is a function of degree $(n-2)$, we shall consider the somewhat more general form
\[
V=\int\frac{Q\,dx}{(\alpha s_x+\beta x+\gamma)\frac{\partial F}{\partial s_x}}
=\int\frac{\Omega\Sigma\pm c_1x_2dx_3}{(\alpha_1x_1+\alpha_2x_2+\alpha_3x_3)\left(c_1\frac{\partial f}{\partial x_1}+c_2\frac{\partial f}{\partial x_2}+c_3\frac{\partial f}{\partial x_3}\right)}.
\]

% Continuation package Abel03: source pp. 039--055, with Toronto backfill for missing Google printed p. 30.
% --- Continuation: source PDF pp. 039--055. ---
% English translation of the authorial German text on printed pp. 17--33.

in which an arbitrary linear function appears in the denominator. The present integral behaves, with respect to the points of contact of the tangents drawn from $c$ to $f=0$, exactly like the preceding one; in general it also becomes discontinuous at the double points and cusps. But it also becomes discontinuous at those $n$ points at which the curve $F=0$ is cut by the line
\[
\alpha s_x+\beta x+\gamma=0.
\]
If $s_x^\circ$, $x^\circ$ is one of these points, then in its neighbourhood
\[
\begin{aligned}
\alpha s_x+\beta x+\gamma&=\alpha(s_x-s_x^\circ)+\gamma(x-x^\circ),\\
F(s_x,x)=0&=\frac{\partial F^\circ}{\partial s_x^\circ}(s_x-s_x^\circ)+\frac{\partial F^\circ}{\partial x^\circ}(x-x^\circ)+\cdots .
\end{aligned}
\]
The last equation shows that $s_x-s_x^\circ$ is proportional to $x-x^\circ$; the first then shows that the denominator under the integral sign of $V$ is likewise proportional to $x-x^\circ$, and hence that $V$ becomes logarithmically infinite. In this respect the intersection points of the line with $F=0$ behave like the double points of the curve.

Since the function $Q$ or $\Omega$ is of order $(n-2)$, it contains
\[
\frac{n(n-1)}{2}
\]
arbitrary coefficients, and one may represent it as a combination of the same number of special functions of the same degree, or, what is the same thing, represent the integral $V$ as an aggregate of
\[
\frac{n(n-1)}{2}
\]
special integrals of the same kind. Now one can evidently place special curves $\Omega=0$ so that they pass through $n-2$ intersection points of the line with $f=0$. More is impossible, unless $\Omega$ splits into the line and a curve of order $(n-3)$, so that $V$ becomes an integral of the first kind; for if a curve $\Omega=0$ of order $(n-2)$ were to pass through $n-1$ of these intersection points, it would have $n-1$ points in common with the line and would therefore contain the line itself. But such a special curve can also be made to pass through the $d+r$ exceptional points of $f=0$; for since their number is at most
\[
\frac{(n-1)(n-2)}{2},
\]
the number of points thereby prescribed for $\Omega=0$ is at most
\[
\frac{(n-1)(n-2)}{2}+n-2=\frac{(n+1)(n-2)}{2},
\]
which is exactly sufficient to determine a curve of order $(n-2)$.

Let then $\Omega_2=0$, $\Omega_3=0\ldots\Omega_n=0$ be special curves which all contain the exceptional points and, besides these, contain all intersection points of the line with $f=0$, respectively excepting, for $\Omega_2=0$, the first and second, for $\Omega_3=0$, the first and third, $\ldots$, and for $\Omega_n=0$, the first and $n$th. Then in the equation
\[
\Omega-k_2\Omega_2-k_3\Omega_3\cdots-k_n\Omega_n=0
\]
one can determine the constants $k$ so that the curve thus represented passes through $n-1$ intersection points of the line with $f=0$, and hence contains the whole line. Namely, if upper indices denote that the coordinates of the corresponding intersection point are to be inserted, one has only the equations
\[
\Omega^{(2)}-k_2\Omega_2^{(2)}=0,
\qquad
\Omega^{(3)}-k_3\Omega_3^{(3)}=0,
\]
\[
\cdots
\qquad
\Omega^{(n)}-k_n\Omega_n^{(n)}=0
\]
to satisfy, in order that the second, third, $\ldots$, $n$th intersection point should lie on the above curve. When the $k$ are determined from these equations, the above curve decomposes into the line and a curve $\Theta=0$ of order $(n-3)$; hence one has
\[
\Omega=k_2\Omega_2+k_3\Omega_3\cdots+k_n\Omega_n+(\alpha_1x_1+\alpha_2x_2+\alpha_3x_3)\Theta.
\]
Thus one obtains the following representation for $V$:
\[
\begin{aligned}
V={}&\sum_{i=2}^{i=n}k_i\int
\frac{\Omega_i\Sigma\pm c_1x_2dx_3}{(\alpha_1x_1+\alpha_2x_2+\alpha_3x_3)\left(c_1\frac{\partial f}{\partial x_1}+c_2\frac{\partial f}{\partial x_2}+c_3\frac{\partial f}{\partial x_3}\right)}\\
&+\int\frac{\Theta\Sigma\pm c_1x_2dx_3}{c_1\frac{\partial f}{\partial x_1}+c_2\frac{\partial f}{\partial x_2}+c_3\frac{\partial f}{\partial x_3}}.
\end{aligned}
\]
This representation shows that the integral $V$ can be reduced to integrals already treated earlier and to such integrals as become logarithmically infinite at only two points. Indeed, an integral
\[
\int\frac{\Omega_i\Sigma\pm c_1x_2dx_3}{(\alpha_1x_1+\alpha_2x_2+\alpha_3x_3)\left(c_1\frac{\partial f}{\partial x_1}+c_2\frac{\partial f}{\partial x_2}+c_3\frac{\partial f}{\partial x_3}\right)}
\]
then remains discontinuous only at the two intersection points of the line with the curve $f=0$ through which the curve $\Omega_i=0$ does not pass. Such integrals, discontinuous at only two points, are called integrals of the third kind.

Not only all the integrals listed under \S\ 3 can be reduced to these integrals, together with integrals of the form treated in the preceding section, but also those integrals treated above,
\[
\int\frac{\Theta_i\Sigma\pm c_1x_2dx_3}{c_1\frac{\partial f}{\partial x_1}+c_2\frac{\partial f}{\partial x_2}+c_3\frac{\partial f}{\partial x_3}},
\]
which became infinite at only one double point. For if, in the integral
\[
\int\frac{\Omega\Sigma\pm c_1x_2dx_3}{(\alpha_1x_1+\alpha_2x_2+\alpha_3x_3)\left(c_1\frac{\partial f}{\partial x_1}+c_2\frac{\partial f}{\partial x_2}+c_3\frac{\partial f}{\partial x_3}\right)},
\]
one lets the line $\alpha$ pass through a double point, and determines $\Omega$, as was taught earlier, so that it passes through all double points and cusps and, in addition, through the remaining $n-2$ points which the line $\alpha$ has in common with the curve $f=0$, then one plainly has an integral of the character treated above. But $\Omega$ then has $n-1$ points in common with the line $\alpha$ and therefore contains it entirely; hence, when
\[
\Omega=(\alpha_1x_1+\alpha_2x_2+\alpha_3x_3)\Theta
\]
the above integral passes into the form
\[
\int\frac{\Theta\Sigma\pm c_1x_2dx_3}{c_1\frac{\partial f}{\partial x_1}+c_2\frac{\partial f}{\partial x_2}+c_3\frac{\partial f}{\partial x_3}},
\]
where $\Theta$ passes through all exceptional points except the one lying on $\alpha$. Thus this integral also supplies the scheme for the not-everywhere-finite integrals listed ad 2) in \S\ 2, insofar as they become discontinuous at double points. These, too, we therefore count among the integrals of the third kind. And since these must always be discontinuous at two points, one may say that the former are discontinuous at the two points which lie united in the double point under consideration.

How these considerations must be modified when the discontinuity point becomes a cusp is to be discussed below.

Since integrals of the third kind, as well as their differential quotients with respect to constants, and finally the algebraic or logarithmic functions separated above, in short all integrals occurring here except the integrals of the first kind, become discontinuous at some points, the following proposition follows as a corollary:

Every integral of the present form which is not discontinuous at any point of $f=0$ is an integral of the first kind.

\begin{center}
{\bfseries \S\ 6. Formation of the normal integrals of the third kind.}
\end{center}

Concerning the formation of an integral of the third kind we can now give the following more precise account.

Let $\xi$, $\eta$ be the two discontinuity points of the integral, so that, since both lie on $f=0$, the equations
\[
f(\xi_1,\xi_2,\xi_3)=0,
\qquad
f(\eta_1,\eta_2,\eta_3)=0
\]
hold. Then the equation of the line $\alpha$ is the following:
\[
\Sigma\pm x_1\xi_2\eta_3=0,
\]
and the integral becomes
\[
V=\int\frac{\Omega\Sigma\pm c_1x_2dx_3}{\Sigma\pm x_1\xi_2\eta_3\left(c_1\frac{\partial f}{\partial x_1}+c_2\frac{\partial f}{\partial x_2}+c_3\frac{\partial f}{\partial x_3}\right)}.
\]
The curve $\Omega=0$ of order $(n-2)$ is now to be determined so that it passes through all exceptional points of $f=0$, and through the $n-2$ points in which the line $\xi\eta$ still cuts the curve $f=0$. Thus the number of all points fixed for $\Omega=0$ is $d+r+n-2$, so that one may still arbitrarily add
\[
\frac{(n-1)(n-2)}{2}-d-r=p
\]
points through which $\Omega$ is to be forced to pass.

The coordinates of a point of the line $\xi$, $\eta$ are known to be
\[
\xi_1+\lambda\eta_1,
\qquad
\xi_2+\lambda\eta_2,
\qquad
\xi_3+\lambda\eta_3,
\]
where $\lambda$ is a variable parameter.

The values of $\lambda$ which correspond to the other $n-2$ intersection points of the line $\xi\eta$ with $f=0$ are obtained by substituting these coordinate expressions into $f=0$. Ordered by powers of $\lambda$, this equation assumes the form
\[
u_0+\lambda u_1+\lambda^2u_2\cdots+\lambda^nu_n=0,
\]
and so gives an equation of degree $n$ in $\lambda$ for determining the $n$ intersection points of our line with $f=0$. But we already know two roots: $\lambda=0$, which gives the point $\xi$, and $\lambda=\infty$, which gives the point $\eta$. Thus in the above equation one must have $u_0=0$, $u_n=0$; and in fact
\[
u_0=f(\xi_1,\xi_2,\xi_3),
\qquad
u_n=f(\eta_1,\eta_2,\eta_3).
\]
After omission of the factor $\lambda$, the above equation therefore passes into the equation of degree $(n-2)$
\[
(1)\ldots\quad u_1+\lambda u_2+\lambda^2u_3\cdots+\lambda^{n-2}u_{n-1}=0,
\]
which serves to determine the remaining $n-2$ intersection points.

These intersection points are also to lie on the curve $\Omega=0$. This means that, if in $\Omega$ one substitutes the quantities $\xi_i+\lambda\eta_i$ for the coordinates, the same equation in $\lambda$ must result as the one just obtained. Thus, if one sets
\[
\Omega(\xi_1+\lambda\eta_1,\xi_2+\lambda\eta_2,\xi_3+\lambda\eta_3)
=v_1+\lambda v_2+\lambda^2v_3\cdots+\lambda^{n-2}v_{n-1},
\]
the $v$ can differ from the corresponding quantities $u$ only by a common factor. Hence one has the equations
\[
(2)\ldots\quad
\begin{cases}
 v_1=\mu u_1,\\
 v_2=\mu u_2,\\
 \cdots\\
 v_{n-1}=\mu u_{n-1},
\end{cases}
\]
where $\mu$ denotes this undetermined common factor.

The equations (2) are $n-1$ in number, linear in the coefficients of $\Omega$ occurring in the $v$, and in the quantity $\mu$, which, without loss of generality, may be put equal to 1; this will be done in what follows. If we add the
\[
\frac{(n-1)(n-2)}{2}
\]
other equations which state that $\Omega=0$ passes through the same number of given points, then we have
\[
\frac{n(n-1)}{2}
\]
linear equations with just as many unknowns, namely the coefficients of $\Omega$, and these equations determine them completely. Substituting the values obtained, $\Omega$ becomes the negative quotient of two determinants. The numerator determinant contains, in its first vertical row, the elements
\[
0, u_1, u_2\ldots u_{n-1}, 0,0\ldots0;
\]
while in the following vertical rows it contains the terms of the functions $\Omega$, $v_1$, $v_2\ldots v_{n-1}$ multiplied by the coefficients of $\Omega$, as well as those which form the left-hand sides of the remaining linear equations. The denominator determinant is the determinant of the linear equations themselves, or the subdeterminant of the preceding one taken with respect to its first element\footnote{This denominator has, moreover, a very remarkable property. It decomposes into the product of $\frac{(n-1)(n-2)}{2}$ three-term determinants and a larger determinant. Setting one of the former equal to zero says that one of the $\frac{(n-1)(n-2)}{2}$ fixed points lies on the line $\xi\eta$; the vanishing of the latter says that those points themselves lie on a curve of order $(n-3)$, a case which must always be avoided.}.

This mode of representing the function $\Omega$ will be adopted in what follows as the basis; the integral of the third kind so formed will be denoted by $S$, or by $S_{\xi\eta}$, or by $S_{\xi\eta}(x)$. It has some important properties, which are to be developed here.

\begin{center}
{\bfseries \S.\ 7. Properties of the normal integrals of the third kind.}
\end{center}

Let us first consider the behaviour of $\Omega$ and $S$ near the discontinuity points. Near the point $\xi$, the $x$ may be replaced by the $\xi$ wherever no vanishingly small factors are involved. Now the first two horizontal rows of the numerator of $\Omega$--the first from $\Omega$ itself, the second from the first equation (2)--are the following:
\[
\begin{array}{cccccc}
0 & x_1^{n-2} & nx_1^{n-3}x_2 & \cdots & x_3^{n-2},\\
u_1 & \xi_1^{n-2} & n\xi_1^{n-3}\xi_2 & \cdots & \xi_3^{n-2}.
\end{array}
\]
When $x$ approaches $\xi$, these rows differ only in the first term; if the second row is subtracted from the first, then the first term of the first row becomes $-u_1$, and the remaining terms vanish. The numerator of $\Omega$ is therefore reduced to $-u_1$ multiplied by the subdeterminant of the first term; and since this is exactly equal to the denominator of $\Omega$, one obtains
\[
\Omega=u_1,
\]
or, if one takes account of the origin of $u_1$ as the coefficient of $\lambda$ in the expansion of
\[
f(\xi_1+\lambda\eta_1,\xi_2+\lambda\eta_2,\xi_3+\lambda\eta_3),
\]
\[
\Omega=\eta_1\frac{\partial f(\xi)}{\partial \xi_1}+\eta_2\frac{\partial f(\xi)}{\partial \xi_2}+\zeta_3\frac{\partial f(\xi)}{\partial \xi_3}.
\]
Thus as long as the integration moves in the neighbourhood of the point $\xi$, one may put for $S$:
\[
S=\frac{\eta_1\frac{\partial f(\xi)}{\partial \xi_1}+\eta_2\frac{\partial f(\xi)}{\partial \xi_2}+\zeta_3\frac{\partial f(\xi)}{\partial \xi_3}}{c_1\frac{\partial f(\xi)}{\partial \xi_1}+c_2\frac{\partial f(\xi)}{\partial \xi_2}+c_3\frac{\partial f(\xi)}{\partial \xi_3}}
\cdot
\int\frac{\Sigma\pm c_1\xi_2dx_3}{\Sigma\pm x_1\xi_2\eta_3}.
\]
If here we put the $\eta$ for the wholly arbitrary quantities $c$, one immediately obtains
\[
S=-\int\frac{\Sigma\pm dx_1\xi_2\eta_3}{\Sigma\pm x_1\xi_2\eta_3}=-\log\Sigma\pm x_1\xi_2\eta_3.
\]
Thus near the discontinuity point $\xi$ the normal integral $S$ behaves like $-\log\Sigma\pm x_1\xi_2\eta_3$; consequently near the discontinuity point $\eta$ it behaves like $+\log\Sigma\pm x_1\xi_2\eta_3$\footnote{This does not change even if a discontinuity point falls at the contact point of one of the tangents to be drawn from $c$ to the curve. Such a coincidence is in general to be regarded as completely immaterial.}.

Assume further three distinct points on the curve $f=0$, which we denote by $\xi$, $\eta$, $\zeta$, and form the sum
\[
S_{\xi\eta}+S_{\eta\zeta}+S_{\zeta\xi}.
\]
Then, apart from a common factor, the integrand contains a sum of the form
\[
\frac{\Omega}{\Sigma\pm x_1\xi_2\eta_3}
+\frac{\Omega'}{\Sigma\pm x_1\eta_2\zeta_3}
+\frac{\Omega''}{\Sigma\pm x_1\zeta_2\xi_3},
\]
where $\Omega'$, $\Omega''$ denote constructions for the point-pairs $\eta\zeta$, $\zeta\xi$ analogous to $\Omega$ for the pair $\xi\eta$. Reducing this expression to a common denominator, the numerator is the function of order $n$
\[
\begin{aligned}
\varphi={}&\Omega\,\Sigma\pm x_1\eta_2\zeta_3\cdot\Sigma\pm x_1\zeta_2\xi_3\\
&+\Omega'\,\Sigma\pm x_1\zeta_2\xi_3\cdot\Sigma\pm x_1\xi_2\eta_3\\
&+\Omega''\,\Sigma\pm x_1\xi_2\eta_3\cdot\Sigma\pm x_1\eta_2\zeta_3.
\end{aligned}
\]
Set equal to zero, this is the equation of a curve of order $n$. It passes through all intersection points of $f=0$ with the lines $\xi\eta$, $\eta\zeta$, $\zeta\xi$. For as regards, for example, the line $\xi\eta$, the left-hand side of its equation is directly contained in two of the above terms. In the first, however, $\Omega$ vanishes at the $n-2$ intersection points which the line has in common with $f=0$ besides $\xi$ and $\eta$, while $\Sigma\pm x\eta\zeta$ vanishes at $\eta$ and $\Sigma\pm x\zeta\xi$ at $\xi$. This property is not changed if, instead of $\varphi$, one considers the expression $\varphi+kf$, where $k$ is a constant. One can always choose this constant so that the curve $\varphi+kf=0$ has the three points $\xi$, $\eta$, $\zeta$ as double points. We set down, for the point $\xi$, the conditions contained in this. They consist in the vanishing of the differential quotients of $\varphi+kf$ with respect to the $x$ for $x=\xi$. For $x=\xi$, omitting the terms whose vanishing for $x=\xi$ is immediately seen, one obtains
\[
\begin{aligned}
\frac{\partial(\varphi+kf)}{\partial x_1}
={}&\bigl[\Omega(\xi_2\zeta_3-\zeta_2\xi_3)+\Omega''(\xi_2\eta_3-\eta_2\xi_3)\bigr]\Sigma\pm \xi_1\eta_2\zeta_3+k\frac{\partial f(\xi)}{\partial \xi_1},\\
\frac{\partial(\varphi+kf)}{\partial x_2}
={}&\bigl[\Omega(\zeta_3\xi_1-\xi_3\zeta_1)+\Omega''(\eta_3\xi_1-\xi_3\eta_1)\bigr]\Sigma\pm \xi_1\eta_2\zeta_3+k\frac{\partial f(\xi)}{\partial \xi_2},\\
\frac{\partial(\varphi+kf)}{\partial x_3}
={}&\bigl[\Omega(\xi_1\zeta_2-\zeta_1\xi_2)+\Omega''(\xi_1\eta_2-\eta_1\xi_2)\bigr]\Sigma\pm \xi_1\eta_2\zeta_3+k\frac{\partial f(\xi)}{\partial \xi_3}.
\end{aligned}
\]
In these expressions, by what precedes,
\[
\Omega=-\left(\eta_1\frac{\partial f(\xi)}{\partial \xi_1}+\eta_2\frac{\partial f(\xi)}{\partial \xi_2}+\zeta_3\frac{\partial f(\xi)}{\partial \xi_3}\right),
\]
\[
\Omega''=\left(\zeta_1\frac{\partial f(\xi)}{\partial \xi_1}+\zeta_2\frac{\partial f(\xi)}{\partial \xi_2}+\zeta_3\frac{\partial f(\xi)}{\partial \xi_3}\right).
\]
Observing also that
\[
0=\xi_1\frac{\partial f(\xi)}{\partial \xi_1}+\xi_2\frac{\partial f(\xi)}{\partial \xi_2}+\xi_3\frac{\partial f(\xi)}{\partial \xi_3}=nf(\xi),
\]
one may write the expression bracketed in the first of the above equations as
\[
\left|
\begin{array}{ccc}
\xi_1\frac{\partial f(\xi)}{\partial \xi_1}+\xi_2\frac{\partial f(\xi)}{\partial \xi_2}+\xi_3\frac{\partial f(\xi)}{\partial \xi_3} & \xi_2 & \xi_3\\
\eta_1\frac{\partial f(\xi)}{\partial \xi_1}+\eta_2\frac{\partial f(\xi)}{\partial \xi_2}+\zeta_3\frac{\partial f(\xi)}{\partial \xi_3} & \eta_2 & \eta_3\\
\zeta_1\frac{\partial f(\xi)}{\partial \xi_1}+\zeta_2\frac{\partial f(\xi)}{\partial \xi_2}+\zeta_3\frac{\partial f(\xi)}{\partial \xi_3} & \zeta_2 & \zeta_3
\end{array}
\right|
=\frac{\partial f(\xi)}{\partial \xi_1}\Sigma\pm \xi_1\eta_2\zeta_3.
\]
Treating the other expressions in the same way, one finally finds
\[
\frac{\partial(\varphi+kf)}{\partial x_i}=\frac{\partial f(\xi)}{\partial \xi_i}\left[k+(\Sigma\pm\xi_1\eta_2\zeta_3)^2\right].
\]
Hence $\xi$ becomes a double point of the curve $\varphi+kf=0$ if
\[
k=-(\Sigma\pm\xi_1\eta_2\zeta_3)^2
\]
is put. But if one determines $k$ in the same way so that $\eta$ or $\zeta$ becomes a double point, one evidently obtains always the same value of $k$, since this value is not changed by interchange of $\xi$, $\eta$, $\zeta$. Thus the curve $\varphi+kf=0$ has all three points as double points. It therefore has with each of the lines $\xi\eta$, $\eta\zeta$, $\zeta\xi$ not only the $n-2$ remaining intersections of the line with $f=0$, but also two double points, thus $n+2$ points in all; hence it must contain these lines entirely. Therefore
\[
\varphi+kf=\Sigma\pm x_1\xi_2\eta_3\cdot\Sigma\pm x_1\eta_2\zeta_3\cdot\Sigma\pm x_1\zeta_2\xi_3\cdot\Theta,
\]
where $\Theta$ is an expression of order $(n-3)$.

It follows, taking account of the equation $f=0$ holding during the integration, that
\[
\frac{\Omega}{\Sigma\pm x_1\xi_2\eta_3}
+\frac{\Omega'}{\Sigma\pm x_1\eta_2\zeta_3}
+\frac{\Omega''}{\Sigma\pm x_1\zeta_2\xi_3}=\Theta,
\]
and therefore
\[
S_{\xi\eta}+S_{\eta\zeta}+S_{\zeta\xi}
=\int\frac{\Theta\Sigma\pm c_1x_2dx_3}{c_1\frac{\partial f}{\partial x_1}+c_2\frac{\partial f}{\partial x_2}+c_3\frac{\partial f}{\partial x_3}}.
\]
Since, by assumption, the $\Omega$ vanish at the double points and cusps, the same holds for $\Theta$; the right-hand side of this equation is therefore an integral of the first kind, and we have the theorem:

The sum of three normal integrals of the third kind, formed for each two of three points of the curve $f=0$ as discontinuity points, is an integral of the first kind.

This equation, which we may write more briefly in the form
\[
(1)\ldots\quad S_{\xi\eta}+S_{\eta\zeta}+S_{\zeta\xi}=I,
\]
is of great importance for the theory of integrals of the third kind. If we put for $\zeta$ any arbitrarily chosen constant point of the curve $f=0$, and observe that by the mode of formation of the function $S$
\[
S_{\eta\xi}=-S_{\xi\eta},
\qquad
S_{\xi\zeta}=-S_{\zeta\xi},
\]
then from (1) one obtains
\[
S_{\xi\eta}=S_{\xi\zeta}-S_{\eta\zeta}+I,
\]
which can be stated in the following theorem:

The normal integral of the third kind $S_{\xi\eta}$ is, up to a term which becomes equal to an integral of the first kind, the difference of two functions formed in exactly the same way, one depending only on the $\xi$, the other only on the $\eta$, and differing only by these parameters.

In this way the integrals $S$ depending on two points $\xi$, $\eta$ have been reduced to functions depending on only one of the discontinuity points.

A further theorem concerns the circumstance that $\Omega$, and therefore also $S$, were not completely determined, but that in forming $\Omega$ a series of wholly arbitrary points had been used. If we replace them by others, we obtain a function $\Omega'$ and an integral $S'$, which agree with $\Omega$ and $S$ in all essential points. The curve $\Omega-\Omega'=0$ of order $(n-2)$
\[
\Omega-\Omega'=0
\]
passes through the $n-2$ remaining intersection points of the line $\xi\eta$ with $f=0$, since at these points $\Omega$ and $\Omega'$ vanish separately. But since, by what precedes, $\Omega$ is reduced for $x=\xi$ to a quite definite value, wholly independent of the arbitrary points chosen, $\Omega'$ assumes the same value for $x=\xi$, and $\Omega-\Omega'$ vanishes for $x=\xi$, and likewise for $x=\eta$. Thus the curve $\Omega-\Omega'=0$ has $n$ points in common with the line $\xi\eta$ and therefore contains it entirely, so that
\[
\Omega-\Omega'=\Theta\cdot\Sigma\pm x_1\xi_2\eta_3,
\]
where $\Theta=0$ is a curve of order $(n-3)$ which, like $\Omega=0$ and $\Omega'=0$, passes through the double points and cusps of $f=0$. Hence
\[
S-S'=\int\frac{\Theta\Sigma\pm c_1x_2dx_3}{c_1\frac{\partial f}{\partial x_1}+c_2\frac{\partial f}{\partial x_2}+c_3\frac{\partial f}{\partial x_3}}=I,
\]
i.e.:

If one forms the normal integral $S$ in two different ways for the same discontinuity points $\xi$, $\eta$, then the two formations differ only by an integral of the first kind.

For that form of the integrals of the third kind in which $x$, $s_x$ are introduced, and $s_x$ is regarded as a function of $x$, we similarly replace the points $\xi$, $\eta$ by value-pairs $s_\xi$, $\xi$, $s_\eta$, $\eta$, which, according to equations 3 of \S\ 1, are connected with the coordinates of those points. Then, by what precedes, the integral of the third kind approaches, near $\xi$, the expression $-\log\Sigma\pm x_1\xi_2\eta_3$, which here first becomes
\[
-\log\left|
\begin{array}{ccc}
1&x&s_x\\
1&\xi&s_\xi\\
1&\eta&s_\eta
\end{array}
\right|
=-\log\bigl[(x-\xi)(s_\xi-s_\eta)-(s_x-s_\xi)(\xi-\eta)\bigr].
\]
If one now observes that near $\xi$
\[
0=(x-\xi)\frac{\partial F(\xi,s_\xi)}{\partial \xi}+(s_x-s_\xi)\frac{\partial F(\xi,s_\xi)}{\partial s_\xi},
\]
then, up to an additive constant, the above expression becomes $-\log(x-\xi)$, and the integral near this discontinuity point therefore behaves like the logarithm of the difference between the variable and the parameter.

\begin{center}
{\bfseries \S\ 8. Integrals of the second kind.}
\end{center}

An integral of the second kind arises when the two discontinuity points of an integral of the third kind come infinitely close together, and hence the line $\xi\eta$ becomes the tangent to the curve $f=0$. Such an integral can be formed directly in the form
\[
\int
\frac{\Omega\Sigma\pm c_1x_2dx_3}{\left(x_1\frac{\partial f(\xi)}{\partial \xi_1}+x_2\frac{\partial f(\xi)}{\partial \xi_2}+x_3\frac{\partial f(\xi)}{\partial \xi_3}\right)
\left(c_1\frac{\partial f}{\partial x_1}+c_2\frac{\partial f}{\partial x_2}+c_3\frac{\partial f}{\partial x_3}\right)},
\]
where an expression now appears in the denominator which, set equal to zero, gives the equation of the tangent at the point $\xi$. One then has to determine the function $\Omega$ of order $(n-2)$ so that $\Omega=0$ passes through the double points and cusps of $f=0$ and through the $n-2$ points at which the tangent at the point $\xi$ still cuts the curve $f=0$.

It follows at once that an integral infinite only at a cusp is an integral of the second kind; for every line drawn through the cusp connects in it two infinitely close points, unlike a line drawn through a double point, which is to be regarded as the joining line of two wholly distinct points. The same consideration which above taught us to recognize an integral infinite at a double point as an integral of the third kind thus characterizes the integral infinite at a cusp--which, moreover, likewise has the form of an integral of the first kind--as discontinuous at two infinitely close points, that is, as an integral of the second kind.

But the representation of an integral of the second kind just indicated generally leads to considerable algebraic difficulties. To avoid these, one may proceed in the following way. If in the equation
\[
S_{\xi\eta}=S_{\xi\zeta}-S_{\eta\zeta}+I
\]
one lets the points $\xi$, $\eta$ approach one another, putting for instance
\[
\eta_1=\xi_1-\varepsilon\alpha_1,
\qquad
\eta_2=\xi_2-\varepsilon\alpha_2,
\qquad
\eta_3=\xi_3-\varepsilon\alpha_3,
\]
and letting $\varepsilon$ tend to zero, one obtains
\[
\lim S_{\xi\eta}=\varepsilon\left(\alpha_1\frac{\partial S_{\xi\zeta}}{\partial \xi_1}+\alpha_2\frac{\partial S_{\xi\zeta}}{\partial \xi_2}+\alpha_3\frac{\partial S_{\xi\zeta}}{\partial \xi_3}\right)+\lim I.
\]
Since $I$ contains the $\xi$, $\eta$, $\zeta$ only in the coefficients, the limit to which it tends still remains an integral of the first kind. Abstracting from it, we may therefore take as an integral of the second kind the expression
\[
E_\xi=\alpha_1\frac{\partial S_{\xi\zeta}}{\partial \xi_1}+\alpha_2\frac{\partial S_{\xi\zeta}}{\partial \xi_2}+\alpha_3\frac{\partial S_{\xi\zeta}}{\partial \xi_3}.
\]
Since, however, $S_{\xi\eta}$ is also equal to the difference of a function of $\xi$ and the same function of $\eta$, apart from an integral of the first kind, the above expression too can contain $\zeta$ only in a term composed of integrals of the first kind; or one can determine an integral $I$ so that
\[
\alpha_1\frac{\partial S_{\xi\zeta}}{\partial \xi_1}+\alpha_2\frac{\partial S_{\xi\zeta}}{\partial \xi_2}+\alpha_3\frac{\partial S_{\xi\zeta}}{\partial \xi_3}-I
\]
is independent of $\zeta$.

The quantities $\alpha$ are not completely arbitrary, but must be chosen so that $\eta$ remains on the curve, or so that
\[
0=f(\eta)=f(\xi)-\varepsilon\left(\alpha_1\frac{\partial f}{\partial \xi_1}+\alpha_2\frac{\partial f}{\partial \xi_2}+\alpha_3\frac{\partial f}{\partial \xi_3}\right)+\cdots .
\]
Since the first term vanishes already, there remains the condition
\[
\alpha_1\frac{\partial f}{\partial \xi_1}+\alpha_2\frac{\partial f}{\partial \xi_2}+\alpha_3\frac{\partial f}{\partial \xi_3}=0.
\]
It can be satisfied in the most general way by setting
\[
\alpha_1=\beta_3\frac{\partial f}{\partial \xi_2}-\beta_2\frac{\partial f}{\partial \xi_3},
\quad
\alpha_2=\beta_1\frac{\partial f}{\partial \xi_3}-\beta_3\frac{\partial f}{\partial \xi_1},
\quad
\alpha_3=\beta_2\frac{\partial f}{\partial \xi_1}-\beta_1\frac{\partial f}{\partial \xi_2},
\]
where the $\beta$ are wholly arbitrary.

Substituting this, one obtains
\[
(1)\ldots\quad
\alpha_1\frac{\partial S_{\xi\zeta}}{\partial \xi_1}+\alpha_2\frac{\partial S_{\xi\zeta}}{\partial \xi_2}+\alpha_3\frac{\partial S_{\xi\zeta}}{\partial \xi_3}
=\left|
\begin{array}{ccc}
\frac{\partial S_{\xi\zeta}}{\partial \xi_1}&\frac{\partial f}{\partial \xi_1}&\beta_1\\
\frac{\partial S_{\xi\zeta}}{\partial \xi_2}&\frac{\partial f}{\partial \xi_2}&\beta_2\\
\frac{\partial S_{\xi\zeta}}{\partial \xi_3}&\frac{\partial f}{\partial \xi_3}&\beta_3
\end{array}
\right|.
\]
But $S_{\xi\zeta}$ is a function depending only on the ratios of the $\xi$, as a counting shows; therefore identically
\[
\xi_1\frac{\partial S_{\xi\zeta}}{\partial \xi_1}+\xi_2\frac{\partial S_{\xi\zeta}}{\partial \xi_2}+\xi_3\frac{\partial S_{\xi\zeta}}{\partial \xi_3}=0.
\]
Moreover,
\[
\xi_1\frac{\partial f}{\partial \xi_1}+\xi_2\frac{\partial f}{\partial \xi_2}+\xi_3\frac{\partial f}{\partial \xi_3}=0.
\]
From the two equations it follows that
\[
\begin{aligned}
\frac{\partial S_{\xi\zeta}}{\partial \xi_2}\frac{\partial f}{\partial \xi_3}
-\frac{\partial S_{\xi\zeta}}{\partial \xi_3}\frac{\partial f}{\partial \xi_2}&=M\cdot\xi_1,\\
\frac{\partial S_{\xi\zeta}}{\partial \xi_3}\frac{\partial f}{\partial \xi_1}
-\frac{\partial S_{\xi\zeta}}{\partial \xi_1}\frac{\partial f}{\partial \xi_3}&=M\cdot\xi_2,\\
\frac{\partial S_{\xi\zeta}}{\partial \xi_1}\frac{\partial f}{\partial \xi_2}
-\frac{\partial S_{\xi\zeta}}{\partial \xi_2}\frac{\partial f}{\partial \xi_1}&=M\cdot\xi_3,
\end{aligned}
\]
and therefore also
\[
\alpha_1\frac{\partial S_{\xi\zeta}}{\partial \xi_1}
+\alpha_2\frac{\partial S_{\xi\zeta}}{\partial \xi_2}
+\alpha_3\frac{\partial S_{\xi\zeta}}{\partial \xi_3}
=M(\beta_1\xi_1+\beta_2\xi_2+\beta_3\xi_3).
\]
The quantity $M$ has become wholly free of the accidental values of the constants $\beta$.

It can also be represented easily. For one obtains
\[
\frac{E_\xi}{\beta_1\xi_1+\beta_2\xi_2+\beta_3\xi_3}
=M=\int\frac{N\cdot\Sigma\pm c_1x_2dx_3}{c_1\frac{\partial f}{\partial x_1}+c_2\frac{\partial f}{\partial x_2}+c_3\frac{\partial f}{\partial x_3}},
\]
where
\[
\begin{aligned}
(\beta_1\xi_1+\beta_2\xi_2+\beta_3\xi_3)N
={}&\alpha_1\frac{\partial}{\partial \xi_1}\frac{\Omega}{\Sigma\pm x_1\xi_2\eta_3}
+\alpha_2\frac{\partial}{\partial \xi_2}\frac{\Omega}{\Sigma\pm x_1\xi_2\eta_3}\\
&+\alpha_3\frac{\partial}{\partial \xi_3}\frac{\Omega}{\Sigma\pm x_1\xi_2\eta_3}.
\end{aligned}
\]
This gives
\[
\begin{aligned}
(\beta_1\xi_1+\beta_2\xi_2+\beta_3\xi_3)N
={}&\frac{1}{(\Sigma\pm x_1\xi_2\eta_3)^2}
\left\{
\left|\begin{array}{ccc}
\frac{\partial f}{\partial \xi_1}&\beta_1&\frac{\partial\Omega}{\partial \xi_1}\\
\frac{\partial f}{\partial \xi_2}&\beta_2&\frac{\partial\Omega}{\partial \xi_2}\\
\frac{\partial f}{\partial \xi_3}&\beta_3&\frac{\partial\Omega}{\partial \xi_3}
\end{array}\right|
\left|\begin{array}{ccc}
x_1&\xi_1&\eta_1\\
x_2&\xi_2&\eta_2\\
x_3&\xi_3&\eta_3
\end{array}\right|
\right.\\
&\left.
-\Omega
\left|\begin{array}{ccc}
\frac{\partial f}{\partial \xi_1}&\beta_1&\eta_2x_3-\eta_3x_2\\
\frac{\partial f}{\partial \xi_2}&\beta_2&\eta_3x_1-\eta_1x_3\\
\frac{\partial f}{\partial \xi_3}&\beta_3&\eta_1x_2-\eta_2x_1
\end{array}\right|
\right\}.
\end{aligned}
\]
Or, if one carries out the first determinant product and transforms the second determinant in the known way,
\[
\begin{aligned}
(\beta_1\xi_1+\beta_2\xi_2+\beta_3\xi_3)N
={}&\frac{1}{(\Sigma\pm x_1\xi_2\eta_3)^2}
\left\{
\left|\begin{array}{ccc}
\sum x_i\frac{\partial f}{\partial \xi_i}&0&\sum \eta_i\frac{\partial f}{\partial \xi_i}\\
\sum x_i\beta_i&\sum \xi_i\beta_i&\sum \eta_i\beta_i\\
\sum x_i\frac{\partial\Omega}{\partial \xi_i}&\Omega&\sum \eta_i\frac{\partial\Omega}{\partial \xi_i}
\end{array}\right|
\right.\\
&\left.
-\Omega
\left|\begin{array}{cc}
\sum \eta_i\frac{\partial f}{\partial \xi_i}&\sum x_i\frac{\partial f}{\partial \xi_i}\\
\sum \eta_i\beta_i&\sum x_i\beta_i
\end{array}\right|
\right\}\\
={}&(\sum\beta_i\xi_i)
\frac{\sum x_i\frac{\partial f}{\partial \xi_i}\sum \eta_i\frac{\partial\Omega}{\partial \xi_i}
-\sum x_i\frac{\partial\Omega}{\partial \xi_i}\sum \eta_i\frac{\partial f}{\partial \xi_i}}{(\Sigma\pm x_1\xi_2\eta_3)^2}.
\end{aligned}
\]
One has therefore
\[
N=-\frac{\sum x_i\frac{\partial f}{\partial \xi_i}\sum \eta_i\frac{\partial\Omega}{\partial \xi_i}
-\sum x_i\frac{\partial\Omega}{\partial \xi_i}\sum \eta_i\frac{\partial f}{\partial \xi_i}}{(\Sigma\pm x_1\xi_2\eta_3)^2},
\]
and hence for the integral of the second kind the form freed from the arbitrary quantities $\alpha$, $\beta$:
\[
\frac{E_\xi}{\beta_1\xi_1+\beta_2\xi_2+\beta_3\xi_3}=\int
\frac{\sum x_i\frac{\partial f}{\partial \xi_i}\sum \eta_i\frac{\partial\Omega}{\partial \xi_i}
-\sum x_i\frac{\partial\Omega}{\partial \xi_i}\sum \eta_i\frac{\partial f}{\partial \xi_i}}{(\Sigma\pm x_1\xi_2\eta_3)^2}
\cdot
\frac{\Sigma\pm c_1x_2dx_3}{c_1\frac{\partial f}{\partial x_1}+c_2\frac{\partial f}{\partial x_2}+c_3\frac{\partial f}{\partial x_3}}.
\]
Since $S_{\xi\zeta}$ is, up to integrals of the first kind, the difference of two functions, one depending only on $\xi$ and the other only on $\zeta$, it is clear that $E_\xi$ is independent of $\zeta$, up to terms which are integrals of the first kind. The point $\zeta$ is no longer a discontinuity point of $E_\xi$; for in the neighbourhood of this point $S_{\xi\zeta}$ passes into
\[
\log\Sigma\pm x_1\zeta_2\xi_3,
\]
which in $E_\xi$ gives the term
\[
\frac{\Sigma\pm x_1\zeta_2\alpha_3}{\Sigma\pm x_1\zeta_2\xi_3}.
\]
Putting here
\[
x_i=\zeta_i+\varrho\xi_i+\tau\alpha_i,
\]
where $\varrho$, $\tau$ are very small, one obtains $-\frac{\varrho}{\tau}$. But at the same time
\[
0=f(x)=\varrho\left(\xi_1\frac{\partial f(\zeta)}{\partial \zeta_1}+\xi_2\frac{\partial f(\zeta)}{\partial \zeta_2}+\xi_3\frac{\partial f(\zeta)}{\partial \zeta_3}\right)
+\tau\left(\alpha_1\frac{\partial f(\zeta)}{\partial \zeta_1}+\alpha_2\frac{\partial f(\zeta)}{\partial \zeta_2}+\alpha_3\frac{\partial f(\zeta)}{\partial \zeta_3}\right),
\]
and therefore
\[
\lim\frac{\Sigma\pm x_1\zeta_2\alpha_3}{\Sigma\pm x_1\zeta_2\xi_3}
=-\frac{\varrho}{\tau}
=\frac{\alpha_1\frac{\partial f(\zeta)}{\partial \zeta_1}+\alpha_2\frac{\partial f(\zeta)}{\partial \zeta_2}+\alpha_3\frac{\partial f(\zeta)}{\partial \zeta_3}}
{\xi_1\frac{\partial f(\zeta)}{\partial \zeta_1}+\xi_2\frac{\partial f(\zeta)}{\partial \zeta_2}+\xi_3\frac{\partial f(\zeta)}{\partial \zeta_3}},
\]
so it does not become infinite.

On the other hand, in the neighbourhood of $\xi$,
\[
\lim S_{\xi\zeta}=-\log\Sigma\pm x_1\xi_2\zeta_3,
\]
and therefore
\[
\lim E_\xi=-\frac{\Sigma\pm x_1\alpha_2\zeta_3}{\Sigma\pm x_1\xi_2\zeta_3}.
\]
If now
\[
x_i=\xi_i+\varrho\zeta_i+\tau\alpha_i
\]
is put, where again $\varrho$, $\tau$ are very small quantities, then
\[
\lim E_\xi=\frac{1}{\tau}.
\]
Thus $E_\xi$ becomes algebraically infinite near the point $\xi$, and indeed of the first order. Observing again that
\[
0=f(x)=\varrho\left(\zeta_1\frac{\partial f(\xi)}{\partial \xi_1}+\zeta_2\frac{\partial f(\xi)}{\partial \xi_2}+\zeta_3\frac{\partial f(\xi)}{\partial \xi_3}\right)
+\tau\left(\alpha_1\frac{\partial f(\xi)}{\partial \xi_1}+\alpha_2\frac{\partial f(\xi)}{\partial \xi_2}+\alpha_3\frac{\partial f(\xi)}{\partial \xi_3}\right),
\]
and that the coefficient of $\tau$ vanishes here because of the condition imposed on the $\alpha$, one sees that $\varrho$ becomes of higher order, and therefore more briefly one can put
\[
x_i=\xi_i+\tau\alpha_i.
\]
From this one obtains the following expressions, containing only the ratios of the $x$:
\[
\lim E_\xi=\frac{1}{\tau}
=-\frac{x_2\alpha_3-x_3\alpha_2}{x_2\xi_3-x_3\xi_2}
=-\frac{x_3\alpha_1-x_1\alpha_3}{x_3\xi_1-x_1\xi_3}
=-\frac{x_1\alpha_2-x_2\alpha_1}{x_1\xi_2-x_2\xi_1},
\]
which, as must be the case, are independent of $\zeta$.

Using the equations
\[
\nu\xi_i=a_i+b_i\xi+c_is_\xi,
\qquad
\nu_1\zeta_i=a_i+b_i\zeta+c_is_\zeta,
\]
one introduces into $S_{\xi\zeta}$ the quantities $\xi$, $s_\xi$, $\zeta$, $s_\zeta$, just as $s_x$, $x$ are introduced in place of the $x$, where then
\[
F(\xi,s_\xi)=0,
\qquad
F(\zeta,s_\zeta)=0,
\]
and at the same time sets
\[
\alpha_i=\frac{1}{\nu}\left(b_i+c_i\frac{ds_\xi}{d\xi}\right).
\]
Then for $E_\xi$ one obtains the simple form
\[
E_\xi=\frac{\partial S_{\xi\zeta}}{\partial \xi}+\frac{\partial S_{\xi\zeta}}{\partial s_\xi}\frac{ds_\xi}{d\xi}
=\frac{dS_{\xi\zeta}}{d\xi}.
\]
Thus in this representation the integral of the second kind may be replaced by the differential quotient of the integral of the third kind with respect to the parameter $\xi$.

The system of values assigned above to the $\alpha$ is in fact possible; for the equation
\[
\alpha_1\frac{\partial f(\xi)}{\partial \xi_1}+\alpha_2\frac{\partial f(\xi)}{\partial \xi_2}+\alpha_3\frac{\partial f(\xi)}{\partial \xi_3}=0
\]
passes into
\[
\frac{\partial F(s_\xi,\xi)}{\partial \xi}
+\frac{\partial F(s_\xi,\xi)}{\partial s_\xi}\frac{ds_\xi}{d\xi}=0,
\]
which is an identity.

Since near $\xi$ the integral $S$ passed into $-\log(x-\xi)$, at that place
\[
\lim E_\xi=\frac{1}{x-\xi};
\]
and at the same time one may remark that because $E$ is obtained by differentiating a function which becomes logarithmically infinite at only two points and has at all other places the character of an algebraic function, no terms becoming logarithmically infinite can occur in $E$\footnote{For the simplest of the cases treated in this section $(p=0,n=2)$ compare the excellent work of Herr Aronhold, Crelle's Journal, vol. 61, p. 95 ff.}.

% --- Continuation: source PDF pp. 056--071. ---
\clearpage

\begin{center}
{\large\bfseries Second Section.}\par\vspace{0.8em}
{\Large\bfseries Abel's Theorem.}
\end{center}

\begin{center}
{\bfseries \S\ 9. Systems of intersections.}
\end{center}

The given curve of order $n$, $f=0$, is cut by another curve of order $m$, $\varphi=0$, in $mn$ points. These $mn$ points, when $\varphi=0$ is itself given, are completely determined, namely by means of an equation of degree $mn$. But when the curve $\varphi=0$ is not itself given, a certain number of intersection-points suffices in order to determine the remaining ones\footnote{Jacobi, de relationibus etc. Crelle's Journal, vol. 15.}. If $m<n$, then
\[
\frac{m\cdot(m+3)}{2}
\]
intersection-points are required as given; by these the curve of order $m$, $\varphi=0$, is completely determined, and therefore also the remaining intersection-points. If, on the other hand, $m\geq n$, then, if only the intersection-points of the two curves are in question, one may put in place of the equation $\varphi=0$ the equation
\[
\varphi'=\varphi+\mu f=0,
\]
where $\mu$ is an expression of order $(m-n)$; and one can choose the
\[
\frac{(m-n+1)(m-n+2)}{2}
\]
coefficients of $\mu$ so that the same number of coefficients of the reduced function $\varphi'$ vanish. Thus, without changing the system of intersections, instead of a general curve $\varphi=0$ one may put a special one
\[
\varphi'=\varphi+\mu f=0
\]
which depends only on
\[
\frac{(m+1)(m+2)}{2}-\frac{(m-n+1)(m-n+2)}{2}-1
=mn-\frac{(n-1)(n-2)}{2}
\]
coefficients. This curve is then determined by
\[
mn-\frac{(n-1)(n-2)}{2}
\]
points; if therefore that many intersection-points of $f=0$, $\varphi=0$ are given, then the remaining
\[
\frac{(n-1)(n-2)}{2}
\]
are determined by them.

The two numbers
\[
mn-\frac{m(m+3)}{2}
\quad\text{and}\quad
\frac{(n-1)(n-2)}{2},
\]
which state how many intersection-points of the system are determined by the others, still agree for $m=n-1$ and $m=n-2$. For $m>n-3$ one may therefore always let the number $\frac{(n-1)(n-2)}{2}$ stand, which, as one sees, is wholly independent of $m$. For smaller values of $m$, however, the number of points which are given by the others is smaller.

These considerations require a small modification when $\varphi$ passes through some double or return points of $f$. These points count as determining elements of $\varphi=0$ simply, but as intersection-points doubly; so that the number of intersection-points determined by the others is diminished by the number of those points which fall into these exceptional points and hence are known from the outset. Thus one has the theorem:

Of the intersection-points of a given curve of order $n$, $f=0$, with a curve of order $m$ which is to pass through $\delta$ exceptional points of $f=0$, the following are determined by the others:
\[
\begin{array}{ll}
1) & \text{if }m\geq n-2,\quad \dfrac{(n-1)(n-2)}{2}-\delta;\\[0.4em]
2) & \text{if }m<n-2,\quad mn-\dfrac{m(m+3)}{2}-\delta.
\end{array}
\]
This theorem suffers an exception only when the given intersection-points themselves are not independent of one another. Namely, if the curve $\varphi=0$ decomposes into several curves of lower degree, then the whole system of intersections likewise decomposes into several, and in each of these systems of intersections only as many points may occur among the given ones as suffice for such a curve of lower degree, by means of the above theorem, to determine the remaining ones.

It is easy to find the equation on which these remaining intersection-points depend. If one eliminates one of the variables from $\varphi=0$, $f=0$, then an equation of degree $mn$ remains; this equation is satisfied by those values of the unknown which correspond to the given intersection-points. If, then, one denotes the unknown, say, by $z$, and the values of it corresponding to the given intersection-points by $z_1,z_2,\ldots,z_{\varrho}$, and denotes the resultant just mentioned by $Z=0$, then
\[
\frac{Z}{(z-z_1)(z-z_2)\cdots(z-z_{\varrho})}=0
\]
is the equation which supplies the remaining intersection-points.

\begin{center}
{\bfseries \S\ 10. Transformation of integral sums which correspond to a system of intersections.}
\end{center}

Abel's theorem makes known transcendental relations which arise from the equations existing among the coordinates of a system of intersections. It is concerned, namely, with the determination of a sum of $mn$ integrals of the same kind of a function which contains no irrationalities other than those conditioned by $f=0$; here the lower limits are the intersection-points of a curve of order $m$, $\varphi=0$, with $f=0$, and the upper limits are the intersection-points of another curve $\psi=0$ with $f=0$. According to our earlier considerations, the question is therefore the determination of a sum of integrals of the form
\[
\Sigma\int F(x_1,x_2,x_3)\,\Sigma\pm c_1x_2\,dx_3,
\]
where $F$ is a homogeneous rational function of order $(-2)$, and where the sum extends over the $mn$ pairs of points which, belonging to the one and the other system of intersections, are to become lower and upper limits of the integrals.

But since we have learned to decompose every function $F$ into such parts that its integral is one of our normal integrals, or at least is composed from differential quotients of such integrals with respect to constants (apart from certain parts that can be integrated immediately), it will suffice to understand by $F$ the form which is found under the sign of our normal integral of the third kind, that is, to consider the sum
\[
V=\Sigma\int
\frac{\Omega_{n-2}\cdot \Sigma\pm c_1x_2\,dx_3}
{\Sigma\pm x_1\xi_2\eta_3\cdot(c_1 f'_{x_1}+c_2 f'_{x_2}+c_3 f'_{x_3})}.
\]
The different integrals arise as the individual intersection-points move from their initial positions to their final positions. One can make this appear more pregnantly by considering a pencil of curves $\varphi+\lambda\psi=0$ and inserting, between the systems of intersections which $f=0$ has in common with $\varphi=0$ $(\lambda=0)$ and with $\psi=0$ $(\lambda=\infty)$, those which $f=0$ has in common with $\varphi+\lambda\psi=0$ (for values of $\lambda$ lying between $0$ and $\infty$). These intermediate systems connect the limiting systems in such a way that, while $\lambda$ moves from $0$ to $\infty$, the systems of points move from the lower limits of the integrals to the upper limits. Thus, if at each instant one considers the $mn$ intersection-points of $f=0$ with a curve of the system $\varphi+\lambda\psi=0$, the coordinates of these points become functions of $\lambda$ alone, and one obtains
\[
V=\int_0^\infty d\lambda\cdot \Sigma
\frac{\Omega_{n-2}\,\Sigma\pm c_1x_2\dfrac{\partial x_3}{\partial\lambda}}
{\Sigma\pm x_1\xi_2\eta_3\cdot(c_1 f'_{x_1}+c_2 f'_{x_2}+c_3 f'_{x_3})},
\]
where the sum now refers only to the algebraic function under the integral sign, and therefore forms a symmetric function of the coordinates of the $mn$ points which are determined by the intersection of $f=0$ with $\varphi+\lambda\psi=0$.

We first transform the expression under the summation sign by multiplying numerator and denominator by the functional determinant
\[
\begin{vmatrix}
\dfrac{\partial f}{\partial x_1}&\dfrac{\partial f}{\partial x_2}&\dfrac{\partial f}{\partial x_3}\\
\dfrac{\partial\varphi}{\partial x_1}&\dfrac{\partial\varphi}{\partial x_2}&\dfrac{\partial\varphi}{\partial x_3}\\
\dfrac{\partial\psi}{\partial x_1}&\dfrac{\partial\psi}{\partial x_2}&\dfrac{\partial\psi}{\partial x_3}
\end{vmatrix}.
\]
This, multiplied by
\[
\begin{vmatrix}
c_1&c_2&c_3\\
x_1&x_2&x_3\\
\dfrac{\partial x_1}{\partial\lambda}&\dfrac{\partial x_2}{\partial\lambda}&\dfrac{\partial x_3}{\partial\lambda}
\end{vmatrix},
\]
gives, by known theorems on homogeneous functions,
\[
\begin{vmatrix}
c_1\dfrac{\partial f}{\partial x_1}+c_2\dfrac{\partial f}{\partial x_2}+c_3\dfrac{\partial f}{\partial x_3} & nf & \dfrac{\partial f}{\partial\lambda}\\
c_1\dfrac{\partial\varphi}{\partial x_1}+c_2\dfrac{\partial\varphi}{\partial x_2}+c_3\dfrac{\partial\varphi}{\partial x_3} & m\varphi & \dfrac{\partial\varphi}{\partial\lambda}\\
c_1\dfrac{\partial\psi}{\partial x_1}+c_2\dfrac{\partial\psi}{\partial x_2}+c_3\dfrac{\partial\psi}{\partial x_3} & m\psi & \dfrac{\partial\psi}{\partial\lambda}
\end{vmatrix}.
\]
For all value-systems of the $x$ occurring in the integral one always has $f=0$, and hence also $\frac{\partial f}{\partial\lambda}=0$; further, $\varphi+\lambda\psi=0$, and hence also
\[
\frac{\partial\varphi}{\partial\lambda}+\lambda\frac{\partial\psi}{\partial\lambda}+\psi=0,
\]
or, if one multiplies by $\psi$ and puts $-\varphi$ for $\lambda\psi$,
\[
\psi\frac{\partial\varphi}{\partial\lambda}-\varphi\frac{\partial\psi}{\partial\lambda}+\psi^2=0.
\]
Therefore the determinant product standing in our numerator becomes
\[
m\left(c_1\frac{\partial f}{\partial x_1}+c_2\frac{\partial f}{\partial x_2}+c_3\frac{\partial f}{\partial x_3}\right)\psi^2.
\]
In the denominator the functional determinant is to be multiplied by
\[
\begin{vmatrix}
x_1&x_2&x_3\\
\xi_1&\xi_2&\xi_3\\
\eta_1&\eta_2&\eta_3
\end{vmatrix};
\]
the product becomes
\[
\begin{vmatrix}
nf&\xi_1\dfrac{\partial f}{\partial x_1}+\xi_2\dfrac{\partial f}{\partial x_2}+\xi_3\dfrac{\partial f}{\partial x_3}&\eta_1\dfrac{\partial f}{\partial x_1}+\eta_2\dfrac{\partial f}{\partial x_2}+\eta_3\dfrac{\partial f}{\partial x_3}\\
m\varphi&\xi_1\dfrac{\partial\varphi}{\partial x_1}+\xi_2\dfrac{\partial\varphi}{\partial x_2}+\xi_3\dfrac{\partial\varphi}{\partial x_3}&\eta_1\dfrac{\partial\varphi}{\partial x_1}+\eta_2\dfrac{\partial\varphi}{\partial x_2}+\eta_3\dfrac{\partial\varphi}{\partial x_3}\\
m\psi&\xi_1\dfrac{\partial\psi}{\partial x_1}+\xi_2\dfrac{\partial\psi}{\partial x_2}+\xi_3\dfrac{\partial\psi}{\partial x_3}&\eta_1\dfrac{\partial\psi}{\partial x_1}+\eta_2\dfrac{\partial\psi}{\partial x_2}+\eta_3\dfrac{\partial\psi}{\partial x_3}
\end{vmatrix}
\]
\[
\begin{aligned}
&=m\psi\Bigl[\Bigl(\xi_1\frac{\partial f}{\partial x_1}+\xi_2\frac{\partial f}{\partial x_2}+\xi_3\frac{\partial f}{\partial x_3}\Bigr)\Bigl(\eta_1\frac{\partial(\varphi+\lambda\psi)}{\partial x_1}+\eta_2\frac{\partial(\varphi+\lambda\psi)}{\partial x_2}+\eta_3\frac{\partial(\varphi+\lambda\psi)}{\partial x_3}\Bigr)\\
&\quad-\Bigl(\eta_1\frac{\partial f}{\partial x_1}+\eta_2\frac{\partial f}{\partial x_2}+\eta_3\frac{\partial f}{\partial x_3}\Bigr)\Bigl(\xi_1\frac{\partial(\varphi+\lambda\psi)}{\partial x_1}+\xi_2\frac{\partial(\varphi+\lambda\psi)}{\partial x_2}+\xi_3\frac{\partial(\varphi+\lambda\psi)}{\partial x_3}\Bigr)\Bigr].
\end{aligned}
\]
And if one denotes the last bracket for the moment by $R$, then, after introducing the results obtained, one has for $V$ the transformed expression
\[
V=\int_0^\infty d\lambda\cdot\Sigma\frac{\Omega\psi}{R}.
\]

It is now natural to introduce, in place of $x_1,x_2,x_3$, new variables $z,s,t$ which are connected with the $x_i$ by the linear relations
\[
x_1=\xi_1s+\eta_1z+\zeta_1t,\quad
x_2=\xi_2s+\eta_2z+\zeta_2t,\quad
x_3=\xi_3s+\eta_3z+\zeta_3t.
\]
Then $s,z,t$ are triangle-coordinates, namely referred to a coordinate triangle whose vertices are the points of discontinuity $\xi$, $\eta$ of the integral and an arbitrary third point $\zeta$. In these new variables one has, for every function $f$,
\[
\frac{\partial f}{\partial s}=\xi_1\frac{\partial f}{\partial x_1}+\xi_2\frac{\partial f}{\partial x_2}+\xi_3\frac{\partial f}{\partial x_3},\qquad
\frac{\partial f}{\partial z}=\eta_1\frac{\partial f}{\partial x_1}+\eta_2\frac{\partial f}{\partial x_2}+\eta_3\frac{\partial f}{\partial x_3},
\]
and therefore $R$ assumes the simple expression
\[
R=\frac{\partial f}{\partial s}\frac{\partial(\varphi+\lambda\psi)}{\partial z}
-\frac{\partial f}{\partial z}\frac{\partial(\varphi+\lambda\psi)}{\partial s}.
\]
In order now to carry out the sum $\frac{\Omega\psi}{R}$, which is to be extended over all value-systems that simultaneously satisfy the equations $f=0$ and $\varphi+\lambda\psi=0$, it is necessary to go back to an algebraic theorem due to Jacobi\footnote{Jacobi, theoremata nova algebraica, Crelle's Journal vol. 14.}.

\begin{center}
{\bfseries \S\ 11. Jacobi's theorem.}
\end{center}

If two equations (curves), $f=0$, $F=0$, of orders $n$ and $m$ are given, homogeneous with respect to the three variables $s,z,t$, then one can eliminate $s$ or $z$ from the two equations and obtains two equations $Z=0$, $S=0$, which are of degree $mn$, respectively in $z,t$ and in $s,t$. If, ordered according to powers of $s$,
\[
f=z_0+z_1s+z_2s^2\cdots+z_ns^n,
\qquad
F=Z_0+Z_1s+Z_2s^2\cdots+Z_ns^n,
\]
then
\[
Z=\begin{vmatrix}
z_0&z_1&z_2&\cdots\\
0&z_0&z_1&\cdots\\
0&0&z_0&\cdots\\
\cdots&\cdots&\cdots&\cdots\\
Z_0&Z_1&Z_2&\cdots\\
0&Z_0&Z_1&\cdots\\
\cdots&\cdots&\cdots&\cdots
\end{vmatrix};
\]
and if one multiplies the second vertical row by $s$, the third by $s^2$, and so on, and adds them to the first row, then one finds
\[
Z=\begin{vmatrix}
f&z_1&z_2&\cdots\\
sf&z_0&z_1&\cdots\\
s^2f&0&z_0&\cdots\\
\cdots&\cdots&\cdots&\cdots\\
F&Z_1&Z_2&\cdots\\
sF&Z_0&Z_1&\cdots\\
\cdots&\cdots&\cdots&\cdots
\end{vmatrix}=Af+BF,
\]
where $A$, $B$ are the determinants
\[
A=\begin{vmatrix}
1&z_1&z_2&\cdots\\
s&z_0&z_1&\cdots\\
s^2&0&z_0&\cdots\\
\cdots&\cdots&\cdots&\cdots\\
0&Z_1&Z_2&\cdots\\
0&Z_0&Z_1&\cdots\\
\cdots&\cdots&\cdots&\cdots
\end{vmatrix},\quad
B=\begin{vmatrix}
0&z_1&z_2&\cdots\\
0&z_0&z_1&\cdots\\
0&0&z_0&\cdots\\
\cdots&\cdots&\cdots&\cdots\\
1&Z_1&Z_2&\cdots\\
s&Z_0&Z_1&\cdots\\
\cdots&\cdots&\cdots&\cdots
\end{vmatrix}.
\]
Quite similarly one has, ordered according to $z$,
\[
f=s_0+s_1z+s_2z^2\cdots+s_nz^n,
\qquad
F=S_0+S_1z+S_2z^2\cdots+S_nz^n,
\]
and
\[
S=\begin{vmatrix}
s_0&s_1&s_2&\cdots\\
0&s_0&s_1&\cdots\\
0&0&s_0&\cdots\\
\cdots&\cdots&\cdots&\cdots\\
S_0&S_1&S_2&\cdots\\
0&S_0&S_1&\cdots\\
\cdots&\cdots&\cdots&\cdots
\end{vmatrix}=Cf+DF,
\]
where
\[
C=\begin{vmatrix}
1&s_1&s_2&\cdots\\
z&s_0&s_1&\cdots\\
z^2&0&s_0&\cdots\\
\cdots&\cdots&\cdots&\cdots\\
0&S_1&S_2&\cdots\\
0&S_0&S_1&\cdots\\
\cdots&\cdots&\cdots&\cdots
\end{vmatrix},\quad
D=\begin{vmatrix}
0&s_1&s_2&\cdots\\
0&s_0&s_1&\cdots\\
0&0&s_0&\cdots\\
\cdots&\cdots&\cdots&\cdots\\
1&S_1&S_2&\cdots\\
z&S_0&S_1&\cdots\\
\cdots&\cdots&\cdots&\cdots
\end{vmatrix}.
\]
The degrees up to which the functions $A,B,C,D$ rise with respect to the variables $s,z,t$ are the following:
\[
\begin{array}{c|ccc}
& s&z&t\\ \hline
A&m-1&mn-n&mn-n\\
B&n-1&mn-m&mn-m\\
C&mn-n&m-1&mn-n\\
D&mn-m&n-1&mn-m
\end{array}
\]
The equation $Z=0$ supplies $mn$ values of $\frac{z}{t}$ corresponding to the $mn$ intersection-points of $f=0$, $F=0$; likewise $S=0$ supplies $mn$ values of $\frac{s}{t}$; these values belong together pairwise. Let us imagine the roots of $Z=0$, and likewise those of $S=0$, ordered, and the corresponding roots denoted with the same indices.

If, in the equations
\[
Z=Af+BF,
\qquad
S=Cf+DF,
\]
one substitutes for $\frac{z}{t}$, $\frac{s}{t}$ two pairs of values not belonging together, then $S$ and $Z$ vanish, but not $f$ and $F$. For these values the determinant
\[
\Delta=AD-BC
\]
must therefore vanish.

But if one substitutes into the above equations a pair of values belonging together, then $f$, $F$ also vanish; by differentiating with respect to $s$ and $z$, omitting as vanishing the terms multiplied by $f$, $F$, one obtains
\[
\frac{\partial Z}{\partial z}=A\frac{\partial f}{\partial z}+B\frac{\partial F}{\partial z},\qquad
0=C\frac{\partial f}{\partial z}+D\frac{\partial F}{\partial z},
\]
\[
0=A\frac{\partial f}{\partial s}+B\frac{\partial F}{\partial s},\qquad
\frac{\partial S}{\partial s}=C\frac{\partial f}{\partial s}+D\frac{\partial F}{\partial s},
\]
and determinant formation gives
\[
\begin{vmatrix}
\dfrac{\partial Z}{\partial z}&0\\[0.2em]
0&\dfrac{\partial S}{\partial s}
\end{vmatrix}
=
\begin{vmatrix}A&B\\C&D\end{vmatrix}
\cdot
\begin{vmatrix}
\dfrac{\partial f}{\partial z}&\dfrac{\partial F}{\partial z}\\[0.2em]
\dfrac{\partial f}{\partial s}&\dfrac{\partial F}{\partial s}
\end{vmatrix}.
\]
The last determinant, analogously to the earlier one, is denoted by $-R$; from this equation one then finds, as the value of $\Delta$ for a corresponding pair of values,
\[
\Delta_i=-\left(\frac{\dfrac{\partial Z}{\partial z}\dfrac{\partial S}{\partial s}}{R}\right)_i,
\]
where the appended index is that of the pair of values in question.

Let us consider $\Delta$ more closely. By the preceding table $\Delta$ is, in all, of degree $2mn-m-n$; but it contains
\[
s\text{ only up to degree }mn-1,
\qquad
z\text{ only up to degree }mn-1,
\]
or, geometrically expressed, the curve $\Delta=0$ has two $(m-1)(n-1)$-fold points, one at $s=0$, $t=0$, the other at $z=0$, $t=0$. Owing to this property of $\Delta$, the homogeneous expression of dimension zero
\[
\frac{\Delta t^{m+n}}{SZ},
\]
which contains only the quotients $\frac{s}{t}$, $\frac{z}{t}$, can be decomposed into partial fractions with respect to $\frac{s}{t}$ and $\frac{z}{t}$, and indeed without the production of an integral function. The decomposition with respect to $\frac{s}{t}$ gives in the numerator the value of $\Delta$ after putting for $\frac{s}{t}$ the corresponding roots of $S=0$. And if one then also carries out the decomposition with respect to $\frac{z}{t}$, the numerators become just those values, after the corresponding roots of $Z=0$ have also been put for $\frac{z}{t}$. Hence all terms whose denominators do not use corresponding roots of $S$, $Z$ vanish from the outset, and only the simple sum remains:
\[
\frac{\Delta t^{m+n}}{SZ}
=t^2\sum_{i=1}^{i=mn}
\frac{\Delta_i\,t_i^{m+n}}
{\left(\dfrac{\partial S}{\partial s}\right)_i\left(\dfrac{\partial Z}{\partial z}\right)_i(st_i-ts_i)(zt_i-tz_i)},
\]
in which only the negative powers of $t$, $t_i$ found in the denominator have been multiplied back up; or, putting the value of $\Delta$, one has
\[
\frac{\Delta t^{m+n}}{SZ}
=-t^2\sum_{i=1}^{i=mn}
\frac{t_i^{m+n}}{R_i(st_i-ts_i)(zt_i-tz_i)}.
\]
If, in this partial-fraction decomposition, instead of $t^{m+n}$ one had had in the numerator an arbitrary function $U$ of order $(m+n-2)$ multiplied by $t^2$, then terms would indeed have been added which, with respect to $z$ or $s$, or to both quantities, would have had the character of an integral function. But, apart from these terms (which the stroke through the equality sign is to indicate), one would have obtained
\[
\frac{\Delta U}{SZ}\not=
\sum_{i=1}^{i=mn}
\frac{U_i t_i^2}{R_i(st_i-ts_i)(zt_i-tz_i)},
\]
where the factor $t^2$ has merely been omitted on both sides.

If one expands both sides of this equation according to ascending powers of $t$, then the coefficients of like powers on the two sides must be equal to one another, excepting such terms as would also occur in the expansion of the omitted terms. Let us consider in particular the terms furnished with the factor $\frac{t^0}{sz}$. Since the omitted terms contain only either $z$ or $s$ in the denominator, no term can appear in their expansion that simultaneously contains $z$ and $s$ in the denominator; hence these terms need not be taken into account. But instead of determining the relevant coefficients from the above equation by expansion, one can simply put $t=0$ on both sides, and only then on the left seek the coefficient of $\frac{1}{sz}$. In this way $S$, $Z$ reduce to their highest terms, $Z$ to $\alpha z^{mn}$, $S$ to $\beta s^{mn}$; and $\alpha$, $\beta$ are nothing other than the left-hand sides of the equations which result when in $f=0$, $F=0$ one puts $t=0$ from the beginning and then eliminates once $s$, once $z$. But in this operation the other variable always drops out of itself; one therefore sees that $\beta=\pm\alpha$, and that, if for $t=0$
\[
f=a_0s^n+a_1s^{n-1}z+\cdots+a_nz^n=f_0,
\qquad
F=b_0s^m+b_1s^{m-1}z+\cdots+b_mz^m=F_0,
\]
then $\alpha$ is nothing other than the determinant
\[
\alpha=\begin{vmatrix}
a_0&a_1&a_2&\cdots\\
0&a_0&a_1&\cdots\\
0&0&a_0&\cdots\\
\cdots&\cdots&\cdots&\cdots\\
b_0&b_1&b_2&\cdots\\
0&b_0&b_1&\cdots\\
0&0&b_0&\cdots\\
\cdots&\cdots&\cdots&\cdots
\end{vmatrix}.
\]
If, in order to form $\beta$, one writes the coefficients $a$ and the coefficients $b$ in reverse order, then there arises a determinant
\[
\beta=(-1)^{mn}\alpha.
\]
And so one has the theorem established by Jacobi:

If $f$, $F$ are two homogeneous functions of orders $n$ and $m$ in $s,z,t$, and the sum extends over the $mn$ value-systems which are common to the equations $f=0$, $F=0$; if, further,
\[
Z=Af+BF=0,
\qquad
S=Cf+DF=0
\]
are the equations which arise from $f=0$, $F=0$ by elimination of $s$ or $z$, $\alpha$ is the common coefficient of $z^{mn}$ in $Z$ and of $s^{mn}$ in $S$, $\Delta=AD-BC$, and finally $U$ is an arbitrary integral function of order $m+n-2$, while $U_0$, $\Delta_0$ are the values of $U$, $\Delta$ for $t=0$, then
\[
\sum_{i=1}^{i=mn}
\frac{U_i}{\left(\dfrac{\partial f}{\partial z}\dfrac{\partial F}{\partial s}-\dfrac{\partial f}{\partial s}\dfrac{\partial F}{\partial z}\right)_i}
=\frac{(-1)^{mn}}{\alpha^2}\left[\frac{\Delta_0U_0}{(sz)^{mn}}\right]_{s^{-1}z^{-1}}.
\]

\begin{center}
{\bfseries \S\ 12. Abel's theorem.}
\end{center}

Let us now apply this theorem to our case, by setting $F=\varphi+\lambda\psi$. First one has to note that in $V$ the product $\Omega\psi$ occurs in place of $U$. Now $t=0$ is the line joining the discontinuity-points $\xi\eta$ of our integral of the third kind. But the curve $\Omega=0$ was so determined that (compare \S\ 6), for the points of this line, it passes into $\frac{f}{sz}$, which for $t=0$ becomes an integral function of order $(n-2)$, since the points $s=0$, $t=0$ and $z=0$, $t=0$ lie on $f=0$. Thus one can put
\[
U_0=\frac{f_0}{sz}\cdot\psi_0.
\]

Besides, from
\[
Z=Af+BF,
\qquad
S=Cf+DF
\]
one finds by combination
\[
DZ-BS=\Delta f,
\]
or
\[
\frac{\Delta f}{ZS}=\frac{D}{S}-\frac{B}{Z}.
\]
The development coefficient
\[
\left[\frac{\Delta U}{SZ}\right]_{s^{-1}z^{-1}}
\]
therefore passes into
\[
\left[\left(\frac{D}{S}-\frac{B}{Z}\right)\frac{\psi}{sz}\right]_{s^{-1}z^{-1}}
\]
for $t=0$, or into
\[
\frac{(-1)^{mn}}{\alpha}\left[\frac{D_0\psi_0}{s^{mn+1}}\right]_{s^{-1}}
-\frac{1}{\alpha}\left[\frac{B_0\psi_0}{z^{mn+1}}\right]_{z^{-1}},
\]
and one sees that, in the first development coefficient, one may put not only $t$ but also $z$ equal to zero, while in the second not only $t$ but also $s$. Let us see what values $D_0$, $B_0$ then assume.

Since $Z=\alpha z^{mn}$, we can form $B_0$ by decomposing $z^{mn}$ into $z^{(m-1)(n-1)}$ and $z^{m+n-1}$, multiplying the last vertical row of $\alpha$ by the latter, and adding the preceding vertical rows multiplied by $sz^{m+n-2}$, $s^2z^{m+n-3}\ldots s^{m+n-1}$. The last vertical row in $\alpha$ will then consist of the following elements:
\[
s^{m-1}f_0,\\ s^{m-2}zf_0,\ldots,z^{m-1}f_0,
\quad s^{n-1}F_0,
\quad s^{n-2}zF_0,\ldots,z^{n-1}F_0.
\]
If we now omit the terms with $f_0$ and take the factor $F_0$ out of the rest, then the remaining determinant, multiplied by $z^{(m-1)(n-1)}$, is $B_0$. But we need $B_0$ only insofar as $s=0$ has also been put in it, that is, by reducing the last vertical row to its last member. Thus $B_0$ is to be put equal to $z^{m(n-1)}$ multiplied by the subdeterminant of $\alpha$ taken after the last element of the last row.

One finds $D_0$ in the same way. Since $S=\beta s^{mn}$, we divide $s^{mn}$ into $s^{(m-1)(n-1)}$ and into $s^{m+n-1}$, multiply the last vertical row of $\alpha$ by the latter, add the remaining rows multiplied by
\[
s^{m+n-2}\cdot z,
\quad s^{m+n-3}\cdot z^2,
\ldots,
 z^{m+n-1},
\]
and then find as last vertical row
\[
s^{m-1}f_0,
\quad s^{m-2}zf_0,\ldots,z^{m-1}f_0,
\quad s^{n-1}F_0,
\quad s^{n-2}zF_0,\ldots,z^{n-1}F_0.
\]
If here we omit the terms with $f_0$, take the factor $F_0$ out of the remainder, and multiply the determinant thereby formed by $s^{(m-1)(n-1)}$, we have $D_0$. But we need $D_0$ only for $z=0$. Hence we can restrict the last vertical row to the member $s^{n-1}F_0$, and can therefore set $D_0$ equal to $s^{m(n-1)}$ multiplied by the subdeterminant of $\beta$ after the last element $(b_0)$ of the last vertical row, or, what is the same, to $(-1)^{mn}\,s^{m(n-1)}$ multiplied by the subdeterminant of $\alpha$ after the $(m+1)$st element $b_0$ of the first vertical row.

Meanwhile we have not yet made use of the fact that in $f_0$ itself the first and last terms ($a_0$ and $a_n$) vanish, since $\xi$ ($z=0$, $t=0$) and $\eta$ ($s=0$, $t=0$) lie on the curve $f=0$. Let us therefore denote by $\gamma$ the second subdeterminant of $\alpha$, taken after the element $b_0$ of the first vertical row and after the last element $b_n$ of the last vertical row, these being the only non-vanishing elements in those rows; then
\[
\alpha=b_0b_n\gamma;
\]
and for $B_0$ we have to put
\[
z^{m(n-1)}\,b_0\gamma,
\]
for $D_0$
\[
(-1)^{mn}s^{m(n-1)}\,b_n\gamma.
\]
And hence our development coefficient becomes, after $\gamma$ cancels,
\[
\frac{1}{b_0}\left[\frac{\psi_0}{s^{m+1}}\right]_{s^{-1}}
-\frac{1}{b_n}\left[\frac{\psi_0}{z^{m+1}}\right]_{z^{-1}}.
\]
In the first bracket we must put for $\psi_0$ its highest term in $s$. But $\psi_0$ was $\psi$, after putting for the $x$ therein, since $t=0$,
\[
x_i=s\xi_i+z\eta_i.
\]
Therefore the highest term in $s$ is plainly
\[
s^m\psi(\xi).
\]
Likewise the highest term in $z$, which is to be taken in the second bracket, becomes
\[
z^m\psi(\eta).
\]
Finally one must note that $b_0$, $b_n$ are in exactly the same way the coefficients of the highest terms in $s$ and $z$ which arise if one makes the same substitution in $F=\varphi+\lambda\psi$, so that
\[
b_0=\varphi(\xi)+\lambda\psi(\xi),
\qquad
b_n=\varphi(\eta)+\lambda\psi(\eta).
\]
Our development coefficient is therefore
\[
\frac{\psi(\xi)}{\varphi(\xi)+\lambda\psi(\xi)}
-\frac{\psi(\eta)}{\varphi(\eta)+\lambda\psi(\eta)}.
\]
Thus our integral becomes
\[
V=-\int_0^\infty \frac{\psi(\xi)\,d\lambda}{\varphi(\xi)+\lambda\psi(\xi)}
+\int_0^\infty \frac{\psi(\eta)\,d\lambda}{\varphi(\eta)+\lambda\psi(\eta)}
\]
\[
=\left[\log\frac{\varphi(\eta)+\lambda\psi(\eta)}{\varphi(\xi)+\lambda\psi(\xi)}\right]_0^\infty
=\log\frac{\psi(\eta)\cdot\varphi(\xi)}{\varphi(\eta)\cdot\psi(\xi)}.
\]
And we have the theorem:

The sum of $mn$ normal integrals of the third kind of the same kind with the discontinuity-points $\xi$, $\eta$, whose lower limits are the $mn$ value-systems corresponding to the $mn$ intersection-points of $f=0$ with a curve of order $m$, $\varphi=0$, and whose upper limits indicate the $mn$ intersection-points of $f=0$ with a curve of order $m$, $\psi=0$, has the value
\[
\log\frac{\psi(\eta)\cdot\varphi(\xi)}{\varphi(\eta)\cdot\psi(\xi)}.
\]

\begin{center}
{\bfseries \S\ 13. Special cases.}
\end{center}

To this general theorem we now have to add the discussion of some special cases.

1) If the given integral is to pass into an integral of the first kind, or at least into such an integral of the third kind as becomes discontinuous only at one exceptional point of $f=0$, then one can express this by saying that $\Omega$ is to contain the factor $\Sigma\pm x_1\xi_2\eta_3$. This factor vanishes with $t$, for the triangle side $t=0$ is completely identical with the line $\Sigma\pm x_1\xi_2\eta_3=0$. Therefore, at the place of our development where $t=0$ was put, the development coefficient, and hence $V$, vanishes. Thus one has the theorem:

The sum of integrals of the first kind of the same kind, or of such integrals of the third kind as become discontinuous only at one exceptional point of $f=0$, extended from the system of intersections of the curve $f=0$ with a curve of order $m$, $\varphi=0$, to the system of intersections of $f=0$ with a curve of order $m$, $\psi=0$, is always equal to zero.

2) If the curves $\varphi=0$, $\psi=0$ have points in common with the curve $f=0$, the corresponding integrals simply drop out, because their upper and lower limits coincide.

3) If the point $\xi$ itself is one of the fixed intersection-points of the curves $\varphi$, $\psi$, then the equation of the theorem becomes illusory. Thus, while in general $\frac{(n-1)(n-2)}{2}$ integral sums vanish, corresponding to the $\frac{(n-1)(n-2)}{2}$ coefficients of a function of order $(n-3)$, this number is reduced by $\delta$ when the curves $\varphi$, $\psi$ pass through $\delta$ double or return points of $f=0$.

4) In general there therefore exist $\frac{(n-1)(n-2)}{2}-\delta$ independent transcendental conditions which express that just as many intersection-points of a curve of order $m$ with $f=0$ are given when this curve is to pass through $\delta$ double or return points of $f=0$. This number agrees exactly with the number of algebraic conditions which we set up above, and the transcendental conditions of Abel's theorem are only another form of the latter.

But if $m\leq n-3$, then the number of algebraic conditions is smaller, namely only $mn-\frac{m(m+3)}{2}-\delta$. Therefore the transcendental conditions cannot then be independent of one another, but must reduce to the same number of mutually independent conditions. Thus, for $m=n-3$, one obtains $\frac{n(n-3)}{2}-\delta+1$ transcendental conditions, while only $\frac{n(n-3)}{2}-\delta$ algebraic conditions are present.

5) If one writes the equation of the theorem for the integrals of the third kind in the form
\[
(1)\quad \Sigma S_{\xi\eta}(x)-\Sigma S_{\xi\eta}(y)
=\log\frac{\varphi(\xi)}{\psi(\xi)}-\log\frac{\varphi(\eta)}{\psi(\eta)},
\]
where the $x$ represent the one system of intersections and the $y$ the other, then from it one obtains the equation for the integrals of the second kind by putting $\eta_i=\xi_i-\varepsilon\alpha_i$ (compare \S\ 8) and letting $\varepsilon$ converge to zero. In place of $S_{\xi\eta}$ there then appears the integral $E_\xi$, augmented by integrals of the first kind, which however vanish from the above sum by the theorems just proved\footnote{The theorem 1) can also be derived from formula (1). Namely, if one forms $S_{\xi\eta}$ in two ways (compare p. 27), the equation (1) holds each time. But the difference of corresponding $S$ is, by the theorem cited at the place just mentioned, an integral of the first kind $I(x)$, and therefore, if one subtracts the two forms of equation (1) from one another, one has $\Sigma I(x)-\Sigma I(y)=0$, which is again theorem 1).}, and it remains:
\[
\Sigma E_\xi(x)-\Sigma E_\xi(y)
=\alpha_1\frac{\partial\log\frac{\varphi(\xi)}{\psi(\xi)}}{\partial\xi_1}
+\alpha_2\frac{\partial\log\frac{\varphi(\xi)}{\psi(\xi)}}{\partial\xi_2}
+\alpha_3\frac{\partial\log\frac{\varphi(\xi)}{\psi(\xi)}}{\partial\xi_3},
\]
which again becomes illusory when the point $\xi$ belongs to both curves $\varphi=0$, $\psi=0$.

6) Cases occur in which, in the general theorem on the normal integrals of the third kind, the logarithmic function vanishes. Namely, if $\xi$ and $\eta$ lie with the base points of the pencil $\varphi+\lambda\psi=0$ on a curve of order $m$, then it is the same curve of this pencil which passes through $\xi$ and through $\eta$. If $\varrho$ is its parameter, then one has
\[
\varphi(\xi)+\varrho\psi(\xi)=0,
\qquad
\varphi(\eta)+\varrho\psi(\eta)=0,
\]
and hence
\[
\frac{\varphi(\xi)}{\psi(\xi)}\cdot\frac{\psi(\eta)}{\varphi(\eta)}=1,
\]
so that the equation of Abel's theorem passes into
\[
\Sigma S_{\xi\eta}(x)-\Sigma S_{\xi\eta}(y)=0.
\]



% --- Continuation source pp. 072--088 ---

\begin{center}
{\large Third Section.}\\[0.35em]
{\Large The one-to-one transformations.}
\end{center}

\subsection*{\S\,14. Transformation of the algebraic integrals.}

The integral which, like those considered by us, refers to a curve
\[
f(x_1,x_2,x_3)=0
\]
and contains no other irrationalities than those grounded in this equation itself, can be transformed in infinitely many ways into other integrals of a similar nature, which refer to another curve
\[
F(y_1,y_2,y_3)=0.
\]
For if, in place of the \(x\), we put in our integral arbitrary rational functions of the \(y\), then on the one hand the integral goes over into a quite similar integral in the \(y\), and on the other \(f=0\) goes over into another equation \(F=0\). In order to prove the first point one only has to show that the expression
\[
\frac{\Omega\cdot \Sigma\pm c_1x_2\,dx_3}
 {c_1f'_{x_1}+c_2f'_{x_2}+c_3f'_{x_3}},
\]
where \(\Omega\) denotes a rational, not necessarily integral, function of \((n-3)\)-rd dimension, is transformed into a similar one. Let the transformation formulae be
\[
\tag{1}
\begin{aligned}
x_1&=\varphi_1(y_1,y_2,y_3),\\
x_2&=\varphi_2(y_1,y_2,y_3),\\
x_3&=\varphi_3(y_1,y_2,y_3),
\end{aligned}
\]
where on the left one may, if desired, still add an undetermined factor. Then \(f=0\) will pass into \(MF=0\), where \(M\) is a factor which may possibly separate off, but may also be \(1\). Since
\[
f=MF,
\]
one also has
\[
\frac{\partial f}{\partial x_i}
=M\left(
\frac{\partial F}{\partial y_1}\frac{\partial y_1}{\partial x_i}
+\frac{\partial F}{\partial y_2}\frac{\partial y_2}{\partial x_i}
+\frac{\partial F}{\partial y_3}\frac{\partial y_3}{\partial x_i}
\right)
\]
for all systems of values for which, as is to be assumed here, \(F\) vanishes; hence
\[
c_1f'_{x_1}+c_2f'_{x_2}+c_3f'_{x_3}
=M\left[k_1F'_{y_1}+k_2F'_{y_2}+k_3F'_{y_3}\right],
\]
where
\[
k_h=c_1\frac{\partial y_h}{\partial x_1}
+c_2\frac{\partial y_h}{\partial x_2}
+c_3\frac{\partial y_h}{\partial x_3}.
\]
We further multiply numerator and denominator of the differential expression by the determinant
\[
\Delta=\Sigma\pm
\frac{\partial y_1}{\partial x_1}
\frac{\partial y_2}{\partial x_2}
\frac{\partial y_3}{\partial x_3};
\]
then
\[
\Delta\cdot\Sigma\pm c_1x_2\,dx_3
=r\begin{vmatrix}
k_1&k_2&k_3\\
y_1&y_2&y_3\\
dy_1&dy_2&dy_3
\end{vmatrix},
\]
where \(r\) gives the order of the functions \(\varphi\). Thus, finally, if for \(\frac1\Delta\) one writes the determinant
\[
D=\Sigma\pm
\frac{\partial x_1}{\partial y_1}
\frac{\partial x_2}{\partial y_2}
\frac{\partial x_3}{\partial y_3},
\]
an integral function of the \(y\), one obtains
\[
\frac{\Omega\cdot\Sigma\pm c_1x_2\,dx_3}
 {c_1f'_{x_1}+c_2f'_{x_2}+c_3f'_{x_3}}
=
\frac{r\Omega D\cdot\Sigma\pm k_1y_2\,dy_3}
 {M\cdot (k_1F'_{y_1}+k_2F'_{y_2}+k_3F'_{y_3})}.
\]
Here the expression on the right is indeed of a nature quite similar to that on the left, except that the function \(\frac{r\Omega D}{M}\) has taken the place of \(\Omega\), and that the equally arbitrary quantities \(k\) appear in place of the arbitrary quantities \(c\).

It is, however, very important that, if the original integral is an integral of the first kind, so that \(\Omega\) was an integral function of \((n-3)\)-rd order, the transformed integral is again an integral of the first kind; for then the numerator \(\Omega D\) can always, with the aid of \(F=0\), be transformed so that it is divisible by \(M\), and so that the remainder vanishes at the double points and cusps of \(F=0\).

To prove this, consider the equations
\[
\begin{aligned}
M\frac{\partial F}{\partial y_1}=0
&=\frac{\partial f}{\partial x_1}\frac{\partial x_1}{\partial y_1}
+\frac{\partial f}{\partial x_2}\frac{\partial x_2}{\partial y_1}
+\frac{\partial f}{\partial x_3}\frac{\partial x_3}{\partial y_1},\\
M\frac{\partial F}{\partial y_2}=0
&=\frac{\partial f}{\partial x_1}\frac{\partial x_1}{\partial y_2}
+\frac{\partial f}{\partial x_2}\frac{\partial x_2}{\partial y_2}
+\frac{\partial f}{\partial x_3}\frac{\partial x_3}{\partial y_2},\\
M\frac{\partial F}{\partial y_3}=0
&=\frac{\partial f}{\partial x_1}\frac{\partial x_1}{\partial y_3}
+\frac{\partial f}{\partial x_2}\frac{\partial x_2}{\partial y_3}
+\frac{\partial f}{\partial x_3}\frac{\partial x_3}{\partial y_3}.
\end{aligned}
\]
These equations must hold simultaneously for all possible discontinuity points of the given integral. It must be shown that, when the expressions here set equal to zero vanish, the numerator \(\Omega D\) becomes zero to an equally high order. It follows that the new integral never becomes discontinuous and, by \S\,5, must be an integral of the first kind. The equations may hold together either because the determinant of the right-hand parts, namely \(D\), vanishes, or because
\[
\frac{\partial f}{\partial x_1}=0,\qquad
\frac{\partial f}{\partial x_2}=0,\qquad
\frac{\partial f}{\partial x_3}=0.
\]
In the latter case either \(f\) has a double point at \(x\), and then \(\Omega\) vanishes because the given integral was to be of the first kind; or all the \(x\) vanish, since for the value-system of the \(y\) in question the three functions \(\varphi\) vanish simultaneously. In this case the expressions on the right become zero of order \(n-1\), since the \(\frac{\partial f}{\partial x}\) become zero of order \(n-1\) as soon as all \(x\) vanish; but at the same time \(D\) (see below) becomes zero of the first order, and \(\Omega\), as a function of order \(n-2\) in the \(x\), becomes zero of order \(n-2\). Thus the product \(\Omega D\) becomes zero of order \(n-1\), like the denominator.

Thus one sees in fact that, as soon as the denominator vanishes, the numerator vanishes in the same way, and therefore that one has an integral of the first kind.

\subsection*{\S\,15. Concept of one-to-one transformations. Preservation of the number \(p\).}

Among the various transformations, those are for the moment of special interest in which the \(x\) can be expressed rationally by the \(y\), and conversely the \(y\) rationally by the \(x\). This is conceivable in so far as, by combining equations (1) with \(F=0\), one can return by elimination of the \(y\) to the equation \(f=0\), and in the course of the elimination process find the \(y\) proportional to integral rational functions of the \(x\). This will occur, however, only when the factor \(M\) assumed above really separates off, and \(F\) is therefore not a function of \(\varphi_1,\varphi_2,\varphi_3\). In that latter case the elimination process would give the \(y\) as functions of the \(x\) only by solving equations (1) themselves, and hence in general with irrationalities attached.

Such transformations, in which, with a certain equation \(f=0\) or \(F=0\) assumed, the \(x\) can be expressed rationally by the \(y\), and these, or at least their ratios, which alone matter, rationally by the \(x\), we call one-to-one transformations. The two curves \(f=0\), \(F=0\) then correspond point for point one-to-one: in general to every point \(y\) of the second curve there corresponds from (1) a single point \(x\) of the first, and conversely.

In such a case every integral of the first kind to be formed for the \(x\) yields one such integral in the \(y\), and every integral of the first kind to be formed in the \(y\) yields one for the \(x\). It follows that the number of integrals of the first kind must be the same in both cases; or, as we may also express it:

\begin{quote}
For curves which correspond to one another point for point, the number \(p\) has the same value.\footnote{Cf. Riemann, Crelle's Journal, vol. 54, p. 133.}
\end{quote}

If \(n,\nu\) are the orders, and \(d+r,\delta+\varrho\) are the numbers of double points and cusps for the two curves, one always has
\[
\frac{(n-1)(n-2)}{2}-d-r
=
\frac{(\nu-1)(\nu-2)}{2}-\delta-\varrho.\footnote{An interesting application of this theorem is to the theory of reciprocal curves, or to the comparison of the singularities which an algebraic curve exhibits in point- and line-coordinates. The equations which connect the coordinates \(y\) of the tangent with the coordinates \(x\) of the point of contact unquestionably belong to the class of transformation equations treated here, since in general to every point there corresponds only one tangent, and to every tangent only one point of contact. In passing from point- to line-coordinates the order \(n\) corresponds to the class \(k\), the number of double points \(d\) to the number of bitangents \(t\), and the number of cusps \(r\) to the number of inflectional tangents \(w\). Equating the two values of \(p\) therefore gives
\[
\frac{(n-1)(n-2)}{2}-d-r
=
\frac{(k-1)(k-2)}{2}-t-w;
\]
an equation which can be verified with the aid of the known Plücker formulae, but which at the same time, since the expressions for \(k,w\) are easily formed, gives the most direct derivation of the number of bitangents of a curve.}
\]

\subsection*{\S\,16. Direct proof of the preservation of the number \(p\).}

This theorem seems important enough to make one wish for another, directly algebraic proof. Such a proof can be carried out as follows.

We start from the equation \(f=0\) and from the inverse formulae (1):
\[
\tag{2}
\begin{aligned}
\mu y_1&=\Phi_1(x_1,x_2,x_3),\\
\mu y_2&=\Phi_2(x_1,x_2,x_3),\\
\mu y_3&=\Phi_3(x_1,x_2,x_3),
\end{aligned}
\]
where \(\mu\) denotes an undetermined factor.

Let the functions \(\Phi\) be of degree \(s\). The equations \(\Phi_1=0\), \(\Phi_2=0\), \(\Phi_3=0\) represent three curves of order \(s\). For generality suppose that there exists a series of points common to these three curves and at the same time to the curve \(f=0\), and denote the number of these common points by \(\sigma+\tau\), where \(\sigma\) is the number of those points which are simple points on \(f=0\), and \(\tau\) the number of those which are perhaps double points or cusps of \(f=0\).

Eliminating the \(x\) from equations (2) and from \(f=0\), one obtains an equation \(F=0\) of degree \(\nu\) in the \(y\); it is the equation of the transformed curve. We first determine the order of this curve.

If we cut \(F=0\) by an arbitrary line
\[
\alpha_1y_1+\alpha_2y_2+\alpha_3y_3=0,
\]
we obtain the order of \(F\) as the number of solutions admitted by this problem. Instead of the still unknown equation \(F=0\), however, put the equation \(f=0\) in connection with equations (2). Substituting the values of the \(y\) in the equation of the line, we obtain, for the determination of the points \(x\) corresponding to the intersection points of \(f=0\) with the line, the equations
\[
\tag{3}
\begin{cases}
f=0,\\
\alpha_1\Phi_1+\alpha_2\Phi_2+\alpha_3\Phi_3=0.
\end{cases}
\]
Each solution of these equations gives in general a point \(y\), so that altogether \(ns\) intersection points would have to exist. But from these one must exclude such solutions as satisfy equations (3) independently of the special values of the \(\alpha\), and which therefore have nothing to do with the proposed problem. For these, \(f\) and the \(\Phi\) must vanish simultaneously. Such extraneous solutions therefore exist, and they lead to the \(\sigma+\tau\) points \(x\) mentioned above; the first \(\sigma\) of them appear as simple intersections of the two curves (3), whereas the other \(\tau\) appear as double, since they are double points or cusps of the first curve. Thus the number \(ns\) of solutions of (3) is to be diminished by \(\sigma+2\tau\), and therefore the order of \(F\) is found from
\[
\nu=ns-\sigma-2\tau.
\]

We now determine the number of double points and cusps of the transformed curve. These may in part correspond to the double points and cusps of the given curve. Conversely, just as double points of the first curve may have split into pairs of points of the second, or cusps into two infinitely near points, so also pairs of points of the first curve may have united into double points of the second, and thus new double points, or, if those points were infinitely near, cusps, may have arisen on the transformed curve.

Let \(\delta',\varrho'\) denote the number of double points and cusps of \(F\) to which double points and cusps of \(f\) correspond, and let \(\delta'',\varrho''\) denote the number of those to which pairs of points on \(f\) correspond. Then
\[
\delta+\varrho=\delta'+\varrho'+\delta''+\varrho''.
\]

On differentiating equations (2), the differentials of the \(y\) are expressed by the differentials of the \(x\) by means of the formulae
\[
\tag{4}
\begin{aligned}
\mu\,dy_1+y_1\,d\mu
&=\frac{\partial\Phi_1}{\partial x_1}dx_1+
  \frac{\partial\Phi_1}{\partial x_2}dx_2+
  \frac{\partial\Phi_1}{\partial x_3}dx_3,\\
\mu\,dy_2+y_2\,d\mu
&=\frac{\partial\Phi_2}{\partial x_1}dx_1+
  \frac{\partial\Phi_2}{\partial x_2}dx_2+
  \frac{\partial\Phi_2}{\partial x_3}dx_3,\\
\mu\,dy_3+y_3\,d\mu
&=\frac{\partial\Phi_3}{\partial x_1}dx_1+
  \frac{\partial\Phi_3}{\partial x_2}dx_2+
  \frac{\partial\Phi_3}{\partial x_3}dx_3.
\end{aligned}
\]
It follows from these equations that, if \(x\) is such a double point or cusp of \(f=0\) for which not all \(\Phi\), and hence not \(\mu\), vanish, then to every direction of progress \(dx\) on the first curve there corresponds a direction of progress \(dy\) on the second; hence to the double point or cusp \(x\) there corresponds a double point or cusp \(y\). If, on the other hand, \(x\) is one of those double points or cusps for which the \(\Phi\) vanish, then \(\mu\) also vanishes. One of the \(dx\) may then be set arbitrarily equal to zero; between the remaining ones there is the equation of second degree
\[
\frac{\partial^2 f}{\partial x_1^2}dx_1^2
+2\frac{\partial^2 f}{\partial x_1\partial x_2}dx_1dx_2+\cdots=0.
\]
Determining from this equation the ratios of the \(dx\), equations (4), from which the \(dy\) have disappeared, give two value-systems of the \(y\), hence two points of \(F=0\), corresponding to such a double point of \(f=0\).\footnote{This is not invalidated by the circumstance that the two points \(y\) may move infinitely close to one another when the double point on \(f=0\) passes into a cusp.} The \(\tau\) double points or cusps of \(f=0\) which also lie on the curves \(\Phi=0\) therefore give no double points or cusps of \(F=0\). Hence
\[
\delta'+\varrho'=d+r-\tau.
\]

The remaining double points and cusps of \(F=0\) arise because the same point \(y\) corresponds to two points \(x\), that is, because there exist value-systems of \(x,x'\) such that
\[
\tag{5}
\begin{gathered}
\Phi_1(x)=\Phi_1(x'),\qquad
\Phi_2(x)=\Phi_2(x'),\qquad
\Phi_3(x)=\Phi_3(x'),\\
f(x)=0,\qquad f(x')=0;
\end{gathered}
\]
for brevity only one argument is written instead of three. The number of solutions admitted by these equations determines \(\delta''+\varrho''\).

But these equations contain the obvious and unusable solution
\[
x_1=x_1',\qquad x_2=x_2',\qquad x_3=x_3'.
\]
To avoid it, we first put in place of equations (5) the following:
\[
\tag{6}
\begin{gathered}
\Phi_1(x)+\varepsilon\Psi_1(x)=\Phi_1(x'),\\
\Phi_2(x)+\varepsilon\Psi_2(x)=\Phi_2(x'),\qquad f(x)=0,\quad f(x')=0,\\
\Phi_3(x)+\varepsilon\Psi_3(x)=\Phi_3(x'),
\end{gathered}
\]
where the \(\Psi\) are arbitrary functions of order \(s\), and eliminate the \(x'\) from these equations. Letting \(\varepsilon\) converge to zero in the result, we obtain the required equation which, together with \(f=0\), yields the pairs of points corresponding to the newly arising double points and cusps.

Since, however, by assumption the elimination of the \(x\) from the equations
\[
\Phi_1(x')=\mu y_1,\qquad
\Phi_2(x')=\mu y_2,\qquad
\Phi_3(x')=\mu y_3,\qquad
f(x')=0
\]
gives the equation
\[
F(y)=0,
\]
and therefore, for all arbitrary values of \(x'\),
\[
F[\Phi_1(x'),\Phi_2(x'),\Phi_3(x')]=0,
\]
the elimination of \(x'\) from equations (6) gives
\[
F[\Phi_1(x)+\varepsilon\Psi_1(x),\Phi_2(x)+\varepsilon\Psi_2(x),
\Phi_3(x)+\varepsilon\Psi_3(x)]=0
\]
\[
=F(\Phi)+\varepsilon\left[
\Psi_1(x)F'_{\Phi_1}+\Psi_2(x)F'_{\Phi_2}+\Psi_3(x)F'_{\Phi_3}
\right]+\cdots.
\]
The first term of this expansion vanishes; hence, after division by \(\varepsilon\), the required resultant remains in the form
\[
\tag{7}
R=\Psi_1(x)F'_{\Phi_1}+\Psi_2(x)F'_{\Phi_2}+\Psi_3(x)F'_{\Phi_3}=0.
\]
This equation, however, still depends on the entirely accidental functions \(\Psi\); one must separate them off in an inessential factor.

Now observe that equations (6), for infinitely small values of \(\varepsilon\), can still be satisfied whenever the quantities \(x\), besides satisfying \(f(x)=0\), satisfy a second equation depending on the \(\Psi\), and at the same time the \(x'\) differ infinitely little from them, if these infinitely small differences are suitably chosen. Put
\[
x_i'=x_i+\varepsilon\xi_i.
\]
Expanding by Taylor's theorem, equations (6), after omission of the factor \(\varepsilon\), become
\[
\begin{aligned}
\Psi_1(x)&=\xi_1\Phi'_1(x_1)+\xi_2\Phi'_1(x_2)+\xi_3\Phi'_1(x_3),\\
\Psi_2(x)&=\xi_1\Phi'_2(x_1)+\xi_2\Phi'_2(x_2)+\xi_3\Phi'_2(x_3),\\
\Psi_3(x)&=\xi_1\Phi'_3(x_1)+\xi_2\Phi'_3(x_2)+\xi_3\Phi'_3(x_3),\\
0&=\xi_1f'(x_1)+\xi_2f'(x_2)+\xi_3f'(x_3).
\end{aligned}
\]
After eliminating the \(\xi\), this gives the sole condition for the \(x\):
\[
S=0=
\begin{vmatrix}
\Psi_1(x)&\Phi'_1(x_1)&\Phi'_1(x_2)&\Phi'_1(x_3)\\
\Psi_2(x)&\Phi'_2(x_1)&\Phi'_2(x_2)&\Phi'_2(x_3)\\
\Psi_3(x)&\Phi'_3(x_1)&\Phi'_3(x_2)&\Phi'_3(x_3)\\
0&f'(x_1)&f'(x_2)&f'(x_3)
\end{vmatrix}.
\]
The value-systems of \(x\) which satisfy \(f=0\), \(S=0\) also satisfy \(R=0\); but the former correspond to no double points or cusps, whereas double points or cusps give only such \(x\) as satisfy \(\frac{R}{S}=0\). In general the division on the left cannot be carried out, even with the aid of \(f=0\).

Since the functions \(\Psi\) are completely arbitrary, one must have
\[
\frac{R}{S}
=\frac{F'_{\Phi_1}}{\Sigma\pm\Phi'_2(x_1)\Phi'_3(x_2)f'(x_3)}
=\frac{F'_{\Phi_2}}{\Sigma\pm\Phi'_3(x_1)\Phi'_1(x_2)f'(x_3)}
=\frac{F'_{\Phi_3}}{\Sigma\pm\Phi'_1(x_1)\Phi'_2(x_2)f'(x_3)}.
\]
It is easy to give a fourth value of this ratio, which is symmetric. Namely, if
\[
F(\Phi_1,\Phi_2,\Phi_3)=M\cdot f,
\]
then, since \(f=0\), one also has
\[
F'_{\Phi_1}\frac{\partial\Phi_1}{\partial x_i}
+F'_{\Phi_2}\frac{\partial\Phi_2}{\partial x_i}
+F'_{\Phi_3}\frac{\partial\Phi_3}{\partial x_i}
-M\frac{\partial f}{\partial x_i}=0.
\]
If in the three equations thus represented one regards \(F'_{\Phi_1},F'_{\Phi_2},F'_{\Phi_3},-M\) as the unknowns and expresses their ratios by minors, one finds the equality of the ratios found above to be \(\frac{R}{S}\), and in addition the fourth value
\[
-\frac{M}{\Sigma\pm
\frac{\partial\Phi_1}{\partial x_1}
\frac{\partial\Phi_2}{\partial x_2}
\frac{\partial\Phi_3}{\partial x_3}}.
\]

The curve \(R=0\) is of degree \(\nu s\), and the curve \(S=0\) of degree \(n+3s-3\). The number of intersections of \(R=0\), \(f=0\) which still exist besides those of \(S=0\), \(f=0\) is therefore
\[
n[\nu s-n-3s+3].
\]
Nevertheless not all these intersections yield points \(x\) corresponding to double points or cusps; the value-systems which coincide with the coordinates of the common intersection points of the curves \(\Phi=0\) with \(f=0\) must still be separated off, since for these equations (5) say nothing at all. For the points \(\sigma,\tau\), \(R\) plainly becomes zero of order \(\nu-1\); moreover one can show that \(S\) becomes zero to second order at the points \(\sigma\), and to third order at the points \(\tau\). It follows that \(\frac{R}{S}\) vanishes to order \(\nu-3\) for the former and to order \(\nu-4\) for the latter. Adding that the former are simple points and the latter double points of \(f=0\), one sees that each of the former absorbs \(\nu-3\), and each of the latter \(2(\nu-4)\), intersections. Thus from the number given above one still has to subtract
\[
(\nu-3)\sigma+2(\nu-4)\tau,
\]
and obtains
\[
n[\nu s-n-3s+3]-(\nu-3)\sigma-2(\nu-4)\tau
=(\nu-1)(\nu-2)-(n-1)(n-2)+2\tau
\]
usable intersections. Of these, two together always correspond to one double point or cusp of \(F=0\); therefore
\[
\delta''+\varrho''
=
\frac{(\nu-1)(\nu-2)}{2}
-\frac{(n-1)(n-2)}{2}
+\tau,
\]
\[
\delta+\varrho
=\delta'+\varrho'+\delta''+\varrho''
=
\frac{(\nu-1)(\nu-2)}{2}
-\frac{(n-1)(n-2)}{2}
+d+r,
\]
or finally
\[
\frac{(\nu-1)(\nu-2)}{2}-\delta-\varrho
=
\frac{(n-1)(n-2)}{2}-d-r,
\]
as was to be proved.

\subsection*{\S\,17. Proof of the auxiliary theorem.}

It remains to prove the propositions which were used here, and which amount to the following:

\begin{quote}
If one has three functions \(\varphi,\psi,f\), of which the first two are of the same order, then their functional determinant vanishes to second order at every point common to the curves \(\varphi=0\), \(\psi=0\), \(f=0\); and to third order if this point is a double point or cusp of \(f=0\).
\end{quote}

Geometrically this theorem may be expressed as follows, if the curve
\[
\Sigma\pm
\frac{\partial\varphi}{\partial x_1}
\frac{\partial\psi}{\partial x_2}
\frac{\partial f}{\partial x_3}
\]
is called Jacobi's curve:\footnote{Cremona, Introd. ad una teor. geom. delle curve piane.}

\begin{quote}
Common intersection points of \(\varphi=0\), \(\psi=0\), \(f=0\) also lie on Jacobi's curve; this curve touches the curve \(f=0\) at them whenever \(\varphi,\psi\) are of the same degree, and if \(f\) has such a point as a double point, it has it likewise as a double point, so that its tangents coincide with those of \(f=0\).
\end{quote}

Let \(\Delta=0\) be the equation of Jacobi's curve, \(x\) one of the points under consideration, and substitute near this point \(x_i+\lambda\xi_i\) for \(x_i\), where \(\lambda\) is very small. Then instead of \(\Delta=0\) one obtains
\[
\Delta+\lambda\sum \xi_i\frac{\partial\Delta}{\partial x_i}
+\frac{\lambda^2}{1\cdot2}\sum\sum \xi_i\xi_k
\frac{\partial^2\Delta}{\partial x_i\partial x_k}
+\cdots=0,
\]
and in place of \(f=0\)
\[
f+\lambda\sum \xi_i\frac{\partial f}{\partial x_i}
+\frac{\lambda^2}{1\cdot2}\sum\sum \xi_i\xi_k
\frac{\partial^2f}{\partial x_i\partial x_k}
+\cdots=0.
\]
If now \(\Delta=0\), \(f=0\), the two equations have the intersection point \(x\) in common. If the terms of first order also agree, so that
\[
\frac{\partial\Delta}{\partial x_i}
=m\frac{\partial f}{\partial x_i},
\]
then, with the aid of the second equation, one can destroy the first-order terms in the first; the first equation begins with \(\lambda^2\), and therefore indicates two intersections coinciding at \(x\). Finally, if both curves have double points at \(x\), then
\[
\frac{\partial\Delta}{\partial x_i},\qquad
\frac{\partial f}{\partial x_i}
\]
are zero. If one then also has
\[
\frac{\partial^2\Delta}{\partial x_i\partial x_k}
=m\frac{\partial^2f}{\partial x_i\partial x_k},
\]
so that the tangent directions of the double point coincide for the two curves, then one can reduce the first equation to terms of third order, while the second begins with terms of second order, and the double point therefore absorbs 6 intersections.

To prove the theorem above, let \(\varphi,\psi\) be of order \(\varrho\), and \(f\) of order \(n\). Further let \(c_1,c_2,c_3\) be arbitrary constants. Then
\[
\tag{1}
\begin{aligned}
&\Delta\,(c_1x_1+c_2x_2+c_3x_3)\\
&=
\begin{vmatrix}
\frac{\partial\varphi}{\partial x_1}&\frac{\partial\varphi}{\partial x_2}&\frac{\partial\varphi}{\partial x_3}&0\\
\frac{\partial\psi}{\partial x_1}&\frac{\partial\psi}{\partial x_2}&\frac{\partial\psi}{\partial x_3}&0\\
\frac{\partial f}{\partial x_1}&\frac{\partial f}{\partial x_2}&\frac{\partial f}{\partial x_3}&0\\
c_1&c_2&c_3&c_1x_1+c_2x_2+c_3x_3
\end{vmatrix}\\
&=
\begin{vmatrix}
\frac{\partial\varphi}{\partial x_1}&\frac{\partial\varphi}{\partial x_2}&\frac{\partial\varphi}{\partial x_3}&-\varrho\varphi\\
\frac{\partial\psi}{\partial x_1}&\frac{\partial\psi}{\partial x_2}&\frac{\partial\psi}{\partial x_3}&-\varrho\psi\\
\frac{\partial f}{\partial x_1}&\frac{\partial f}{\partial x_2}&\frac{\partial f}{\partial x_3}&-nf\\
c_1&c_2&c_3&0
\end{vmatrix}.
\end{aligned}
\]
It follows that \(\Delta\) vanishes together with \(\varphi,\psi,f\). Using this and differentiating with respect to \(x_i\), one has for those values of \(x\) at which \(\varphi,\psi,f,\Delta\) are zero:
\[
\begin{aligned}
&\frac{\partial\Delta}{\partial x_i}(c_1x_1+c_2x_2+c_3x_3)\\
&=
\begin{vmatrix}
\frac{\partial\varphi}{\partial x_1}&\frac{\partial\varphi}{\partial x_2}&\frac{\partial\varphi}{\partial x_3}&-\varrho\frac{\partial\varphi}{\partial x_i}\\
\frac{\partial\psi}{\partial x_1}&\frac{\partial\psi}{\partial x_2}&\frac{\partial\psi}{\partial x_3}&-\varrho\frac{\partial\psi}{\partial x_i}\\
\frac{\partial f}{\partial x_1}&\frac{\partial f}{\partial x_2}&\frac{\partial f}{\partial x_3}&-n\frac{\partial f}{\partial x_i}\\
c_1&c_2&c_3&0
\end{vmatrix}\\
&=
\begin{vmatrix}
\frac{\partial\varphi}{\partial x_1}&\frac{\partial\varphi}{\partial x_2}&\frac{\partial\varphi}{\partial x_3}&0\\
\frac{\partial\psi}{\partial x_1}&\frac{\partial\psi}{\partial x_2}&\frac{\partial\psi}{\partial x_3}&0\\
\frac{\partial f}{\partial x_1}&\frac{\partial f}{\partial x_2}&\frac{\partial f}{\partial x_3}&(\varrho-n)\frac{\partial f}{\partial x_i}\\
c_1&c_2&c_3&\varrho c_i
\end{vmatrix}\\
&=-(\varrho-n)\frac{\partial f}{\partial x_i}\,
 \Sigma\pm
 \frac{\partial\varphi}{\partial x_1}
 \frac{\partial\psi}{\partial x_2}c_3.
\end{aligned}
\]
Thus \(\Delta=0\) and \(f=0\) touch.

If, however, \(x\) is a double point of \(f=0\), the right side, and therefore also the left side, vanishes; hence \(\Delta=0\) likewise has the point \(x\) as a double point. Differentiating (1) twice, with respect to \(x_i\) and \(x_k\), one may for simplification note that, since here at the same time
\[
(n-1)\frac{\partial f}{\partial x_i}
=x_1\frac{\partial^2 f}{\partial x_1\partial x_i}
+x_2\frac{\partial^2 f}{\partial x_2\partial x_i}
+x_3\frac{\partial^2 f}{\partial x_3\partial x_i}=0,
\]
\[
\varrho\varphi
=x_1\frac{\partial\varphi}{\partial x_1}
+x_2\frac{\partial\varphi}{\partial x_2}
+x_3\frac{\partial\varphi}{\partial x_3}=0,
\]
\[
\varrho\psi
=x_1\frac{\partial\psi}{\partial x_1}
+x_2\frac{\partial\psi}{\partial x_2}
+x_3\frac{\partial\psi}{\partial x_3}=0,
\]
the determinant
\[
\begin{vmatrix}
\frac{\partial^2 f}{\partial x_1\partial x_i}&
\frac{\partial^2 f}{\partial x_2\partial x_i}&
\frac{\partial^2 f}{\partial x_3\partial x_i}\\
\frac{\partial\varphi}{\partial x_1}&
\frac{\partial\varphi}{\partial x_2}&
\frac{\partial\varphi}{\partial x_3}\\
\frac{\partial\psi}{\partial x_1}&
\frac{\partial\psi}{\partial x_2}&
\frac{\partial\psi}{\partial x_3}
\end{vmatrix}
\]
also vanishes. In the transformed determinant (1) one need not consider the terms in which the third horizontal row is not differentiated, since it then vanishes entirely. If one differentiates the third row once and the first or second also once, determinants arise in which the last vertical row contains only one non-zero term. Destroying this with the aid of the other vertical rows gives only the last term of the last vertical row as non-zero, and this term is multiplied by one of the expressions whose vanishing we have just noted. Hence, differentiating only the third horizontal row twice, there remains
\[
\frac{\partial^2\Delta}{\partial x_i\partial x_k}
(c_1x_1+c_2x_2+c_3x_3)
=-(\varrho-n)\frac{\partial^2f}{\partial x_i\partial x_k}
\,\Sigma\pm
\frac{\partial\varphi}{\partial x_1}
\frac{\partial\psi}{\partial x_2}c_3.
\]
Thus the tangents of the double point of \(\Delta\) coincide with those of the double point of \(f\).

\subsection*{\S\,18. Reduction of the order of the curve. Normal curves.}

We now wish to treat the question of the lowest degree \(\nu\) of a curve \(F=0\) which can be obtained by transformation from a given \(f=0\). First suppose that the order \(s\) of the expressions \(\Phi\), which in the transformation are set equal to the new variables, is greater than \(n-3\). In this case the points at which a curve \(\Phi=0\) cuts the curve \(f=0\), and which alone matter since \(\Phi\) occurs only in connection with \(f=0\), are given, by \S\,9, by
\[
ns-\frac{(n-1)(n-2)}{2}=ns-p-d-r
\]
of them. The number \(\sigma+\tau\) of points common to all \(\Phi\) with \(f=0\) must therefore be at least 2 less than this number. If it were only 1 less, then, with the aid of \(f=0\), one could first destroy in one \(\Phi\) all coefficients except \(ns-p-d-r+1\), and the remaining ones would then have to satisfy the \(ns-p-d-r-1\) linear equations carried along by the given points. The function \(\Phi\) would thereby be determined up to a single coefficient \(\lambda\); the three \(\Phi\) would all have the forms \(u+\lambda_1v\), \(u+\lambda_2v\), \(u+\lambda_3v\), and there would be a linear relation among them, which is not allowed. One must therefore have
\[
\sigma+\tau\le ns-p-d-r-2.
\]
Adding the equation
\[
\nu=ns-\sigma-2\tau,
\]
one obtains
\[
\nu\ge p+d+r+2-\tau.
\]

Thus the smallest value which the order \(\nu\) of the transformed curve can assume occurs when \(\tau=d+r\) and \(\nu=p+2\). This value, as one sees, is entirely independent of the order of the transformation equations; all curves \(\Phi=0\) pass through all double points of the curve \(f=0\), and in addition through
\[
\sigma=ns-\nu-2\tau=ns-p-2-2(d+r)
\]
\[
=ns-\frac{(n-1)(n-2)}{2}-2-d-r
\]
fixed common points. For the possibility of the transformation it is still necessary that \(\sigma\) be zero or positive. Now \(s>n-3\), and therefore also
\[
\sigma>\frac{n(n-3)}{2}-3-d-r=p-4.
\]
Thus \(\sigma\) is indeed zero or positive as soon as \(p>2\). Only the cases \(p=0\), \(p=1\), \(p=2\) require special treatment.

We see that with the aid of substitutions whose order exceeds \(n-3\) we obtain, at the lowest, curves of order \(p+2\). We now examine substitutions of lower order.

If the order of the substitutions is
\[
s=n-3-k,\qquad (k=0,1\ldots,n-4),
\]
then first again
\[
\nu=ns-\sigma-2\tau=n(n-3-k)-\sigma-2\tau.
\]
Further, the number \(\sigma+\tau\) of common points must be at least 2 less than the number of points which determine a curve of order \(s\); that is, one must have
\[
\sigma+\tau\le \frac{(n-3-k)(n-k)}{2}-2.
\]
Together with
\[
\sigma+2\tau=n(n-3-k)-\nu
\]
this gives the inequality
\[
\tau\ge \frac{(n-3-k)(n+k)}{2}-\nu+2,
\]
or
\[
\tau\ge p+1+d+r-\frac{k\cdot k+3}{2}-\nu,
\]
or finally
\[
\nu\ge p+1+(d+r-\tau)-\frac{k\cdot k+3}{2}.
\]
It follows that, for a given \(k\) or \(s\), the most favorable case occurs when
\[
\nu=p+1+(d+r-\tau)-\frac{k\cdot k+3}{2},
\]
and at the same time
\[
\sigma=p-3-kn+(d+r-\tau)+\frac{k\cdot k+3}{2}.
\]

It is therefore always possible to reduce the curve \(f=0\) to a curve of order \(p+1\); this is achieved by a transformation of order \(n-3\) (\(k=0\)), in which all curves \(\Phi\) pass through all double points and cusps \((\tau=d+r)\), and in addition through \(p-3\) points \(\sigma\) chosen arbitrarily on \(f=0\).

Here also, of course, the cases \(p=0,1,2\) are to be excluded and treated separately.

This transformation is generally applicable, since it occurs no matter how large the number of double points contained in the curve.

From the inequalities set up on p. 64 it follows that one always has
\[
\tau\le \frac{(n-3-k)(n-k)}{2}-2,
\]
\[
\tau\ge \frac{(n-3-k)(n+k)}{2}+2-\nu.
\]
It follows that
\[
\frac{(n-3-k)(n-k)}{2}-2
\ge
\frac{(n-3-k)(n+k)}{2}+2-\nu,
\]
or
\[
\nu\ge 4+k(n-k-3),
\]
or finally
\[
\nu-n\ge (k-1)(n-k-4).
\]
Since \(k\) can take only the values \(0,1,\ldots,n-4\), the right side of this inequality is zero or positive, so that \(\nu\ge n\), except when \(k=0\). Hence a reduction of the order of the curve by a transformation can arise in general, that is, without special assumptions on the position of the double points and cusps, only for \(k=0\).

Therefore, if the degree of the transformed curve is to be lower than \(n\) at all, this must occur through the general transformation mentioned above, which is of order \(n-3\) and leads to the degree \(p+1\). One must then have
\[
p+1<n,
\]
and hence for the number of double points the determination
\[
d+r>\frac{(n-1)(n-4)}{2}.
\]
Since the transformed curve becomes of order \(p+1\) when one applies a transformation of order \(n-3\) in the indicated way, and since the number \(p\) is preserved, the new curve has
\[
\frac{p(p-1)}{2}-p=\frac{p(p-3)}{2}
\]
double points. The curves of order \(p+1\) with \(\frac{p(p-3)}{2}\) double points therefore form, as it were, the normal form for all curves belonging to a given number \(p\), since all such curves can be transformed into them. Conversely, however, by what precedes, these curves are in general not transformable into lower ones, and all curves of lower order to which the same \(p\) belongs can arise only from special curves of order \(p+1\) with \(\frac{p(p-3)}{2}\) double points, which by their special nature exceptionally allow the order to be lowered. Our normal curves, enumerated more precisely from \(p=3\) onward, are the following:
\[
\begin{array}{ll}
p=3:& \text{curves of order }4\text{ with }0\text{ double points},\\
p=4:& \text{curves of order }5\text{ with }2\text{ double points},\\
p=5:& \text{curves of order }6\text{ with }5\text{ double points},\\
p=6:& \text{curves of order }7\text{ with }9\text{ double points},\\
p=7:& \text{curves of order }8\text{ with }14\text{ double points},\\
&\text{and so on.}
\end{array}
\]
The general curve of order \(5\) appears thereby as the transformation of a special curve of order \(7\) with 9 double points; the curve of order \(5\) with one double point as the transformation of a special curve of order \(6\) with 5 double points; etc.

% Continuation: source PDF pp. 089--106.

\subsection*{\S\,19. Treatment of the case \(p=0\). Normal curve: a straight line.}

The cases which still require a separate investigation are those in which \(p=0,1,2\). They shall be treated in order. The simplest curves of this kind are respectively the straight line, the curve of order \(3\) without a double point, and the curve of order \(4\) with one double point. It is to be shown how the most general curves for which \(p=0,1,2\) can respectively be transformed into these curves.

If for a curve of order \(n\) one has \(p=0\), then it possesses
\[
\frac{(n-1)(n-2)}{2}
\]
double and return points. Consider these, together with \(2n-3\) further points chosen on the curve, as base-points of a pencil of curves of order \(n-1\). This pencil is thereby completely determined; for through the
\[
\frac{(n-1)(n+2)}{2}-1
\]
fixed points the constants of a curve of order \(n-1\) are determined up to one. Let this one, left arbitrary, be \(\lambda\); then the equation of the pencil is of the form
\[
(1)\ldots\quad u+\lambda v=0,
\]
where \(u\) and \(v\) are functions of order \(n-1\) which, set equal to zero, represent any two curves of the pencil.

If we combine equation (1) with the equation \(f=0\) of the given curve, we obtain the \(n\cdot(n-1)\) points in which a curve of the pencil cuts the given curve. But these consist of the \(\frac{(n-1)(n-2)}{2}\) double and return points, to be counted twice, of the \(2n-3\) arbitrarily chosen fixed points, and of
\[
n(n-1)-(n-1)(n-2)-(2n-3)=1
\]
moving point of intersection, that is, one dependent on \(\lambda\). If we combine the equations
\[
f=0,\qquad u+\lambda v=0
\]
with the equation
\[
\alpha=\alpha_1x_1+\alpha_2x_2+\alpha_3x_3=0,
\]
in which the \(\alpha\)'s denote the coordinates of an arbitrary line, and eliminate the \(x\)'s from these equations, there arises an equation for the \(\alpha\)'s which says that the line \(\alpha\) passes through a point of intersection of the two curves, that is, there arises the product of the \(n\cdot(n-1)\) equations of the individual intersection points. Since the degrees of the given equations in the variables are represented by the table
\[
\begin{array}{c|ccc}
 & x & \alpha & \lambda\\ \hline
f=0 & n & 0 & 0\\
u+\lambda v=0 & n-1 & 0 & 1\\
\alpha=0 & 1 & 1 & 0
\end{array}
\]
the final equation of the elimination is in fact of degree \(n\cdot(n-1)\) in the \(\alpha\)'s and also of degree \(n\) in \(\lambda\). This final equation must split into \(n\cdot(n-1)\) linear factors, of which only one contains \(\lambda\), namely the one which gives the equation of the moving point of intersection. Hence a factor of degree \([n(n-1)-1]\) in the \(\alpha\)'s, not containing \(\lambda\), must separate off. After separating it, the elimination equation has the form
\[
A_1\alpha_1+A_2\alpha_2+A_3\alpha_3=0,
\]
where the \(A\)'s are functions of degree \(n\) in \(\lambda\). This is the equation of the moving point of intersection; its coordinates therefore have the same ratios as the functions \(A\), and thus one has represented the coordinates \(x\) of the moving point of intersection as rational functions of order \(n\) of a parameter \(\lambda\), by means of the equations
\[
(1)\ldots\quad
\begin{gathered}
\mu x_1=A_1,\\
\mu x_2=A_2,\\
\mu x_3=A_3,
\end{gathered}
\]
where \(\mu\) is an arbitrary factor.

If we put \(\lambda=\frac{y_1}{y_2}\), then the \(x\)'s are represented as homogeneous functions of order \(n\) of \(y_1,y_2\). If now \(y_1,y_2,y_3\) are regarded as coordinates of a point of the transformed curve, then the \(x\)'s may be regarded as functions of the \(y\)'s for which
\[
y_3=0;
\]
this is therefore the equation of the transformed curve, that is, the curve is a straight line, as was required.

Thus every curve of order \(n\) with \(\frac{(n-1)(n-2)}{2}\) double and return points can be transformed into a straight line by means of the transformation (1), which is of order \(n\).

\subsection*{\S\,20. Treatment of the case \(p=1\). Normal curve: a general curve of third order.}

We pass to curves of order \(n\) with \(\frac{n(n-3)}{2}\) double and return points, for which \(p=1\). Consider the double and return points, together with \(2n-2\) further fixed points chosen arbitrarily on the curve, as the
\[
\frac{(n-1)(n+2)}{2}-1
\]
base-points of a pencil \(u+\lambda v=0\) of curves of order \(n-1\). Proceeding again exactly as before, one has to form the eliminant resultant from the equations
\[
f=0,\qquad u+\lambda v=0,\qquad \alpha=0.
\]
As before, it is of degree \(n(n-1)\) in the \(\alpha\)'s and of degree \(n\) in \(\lambda\), and represents the product of \(n(n-1)\) linear factors in the \(\alpha\)'s, the equations of the \(n(n-1)\) intersection points of \(f=0\) with \(u+\lambda v=0\). Of these, \(n(n-3)\) fall into the double and return points, and \(2n-2\) into the arbitrarily chosen fixed points; there therefore remain
\[
n(n-1)-n(n-3)-(2n-2)=2
\]
moving intersection points. After separation of the factor of degree \([n(n-1)-2]\), not containing \(\lambda\), the resultant therefore passes into the form
\[
(1)\ldots\quad
0=w_{11}\alpha_1^2+w_{22}\alpha_2^2+w_{33}\alpha_3^2
+2w_{23}\alpha_2\alpha_3+2w_{31}\alpha_3\alpha_1+2w_{12}\alpha_1\alpha_2,
\]
where the \(w\)'s are functions of degree \(n\) in \(\lambda\). Since this equation represents the product of the equations of the two moving intersection points, it must be resolvable into two linear factors; that is, independently of the value of \(\lambda\), one must have
\[
(2)\ldots\quad
W=\begin{vmatrix}
w_{11}&w_{12}&w_{13}\\
w_{21}&w_{22}&w_{23}\\
w_{31}&w_{32}&w_{33}
\end{vmatrix}=0.
\]
Furthermore, if \(x_1,x_2,x_3\) and \(\xi_1,\xi_2,\xi_3\) are the coordinates of the moving intersection points, then the above equation must be equivalent to
\[
(3)\ldots\quad
0=(x_1\alpha_1+x_2\alpha_2+x_3\alpha_3)(\xi_1\alpha_1+\xi_2\alpha_2+\xi_3\alpha_3).
\]
Thus, to determine \(x,\xi\), one has the equations
\[
(4)\ldots\quad
\begin{array}{lll}
x_1\xi_1=w_{11},&
x_2\xi_3+x_3\xi_2=2w_{23},\\
x_2\xi_2=w_{22},&
x_3\xi_1+x_1\xi_3=2w_{31},\\
x_3\xi_3=w_{33},&
x_1\xi_2+x_2\xi_1=2w_{12}.
\end{array}
\]
If the coordinates of the line joining the points \(x,\xi\) are denoted by \(p_1,p_2,p_3\), so that
\[
(5)\ldots\quad
\begin{gathered}
p_1=x_2\xi_3-x_3\xi_2,\\
p_2=x_3\xi_1-x_1\xi_3,\\
p_3=x_1\xi_2-x_2\xi_1,
\end{gathered}
\]
then from (4) one immediately finds
\[
p_ip_k=-4W_{ik},
\]
where the \(W_{ik}\) are the minors of \(W\). The minors to be formed from the \(W_{ik}\) are therefore all identically zero, so that for every value of \(\lambda\):
\[
\begin{array}{lll}
W_{22}W_{33}=W_{23}^{2},& W_{12}W_{13}=W_{11}W_{23},\\
W_{33}W_{11}=W_{31}^{2},& W_{23}W_{21}=W_{22}W_{31},\\
W_{11}W_{22}=W_{12}^{2},& W_{31}W_{32}=W_{33}W_{12}.
\end{array}
\]
This is possible only if the \(W_{ik}\), apart from a common factor which we shall denote by \(-Q\), become equal to the squares and products of rational functions \(l_1,l_2,l_3\), that is, if
\[
W_{ik}=-l_il_kQ.
\]
It then follows that
\[
p_1=2l_1\sqrt Q,\qquad p_2=2l_2\sqrt Q,\qquad p_3=2l_3\sqrt Q,
\]
and further, from equations (4) and (5) together,
\[
(6)\ldots\quad
\begin{array}{lll}
x_1\xi_1=w_{11},&x_2\xi_1=w_{21}-l_3\sqrt Q,&x_3\xi_1=w_{31}+l_2\sqrt Q,\\
x_1\xi_2=w_{12}+l_3\sqrt Q,&x_2\xi_2=w_{22},&x_3\xi_2=w_{32}-l_1\sqrt Q,\\
x_1\xi_3=w_{13}-l_2\sqrt Q,&x_2\xi_3=w_{23}+l_1\sqrt Q,&x_3\xi_3=w_{33}.
\end{array}
\]
Instead of these equations one may also write the following ones. They arise when these rows are multiplied by three arbitrary quantities \(\gamma_1,\gamma_2,\gamma_3\), added, and \(\varrho=\xi_1\gamma_1+\xi_2\gamma_2+\xi_3\gamma_3\) is regarded as an arbitrary factor:
\[
(7)\ldots\quad
\begin{aligned}
\varrho x_1&=w_{11}\gamma_1+w_{12}\gamma_2+w_{13}\gamma_3+(\gamma_2l_3-\gamma_3l_2)\sqrt Q,\\
\varrho x_2&=w_{21}\gamma_1+w_{22}\gamma_2+w_{23}\gamma_3+(\gamma_3l_1-\gamma_1l_3)\sqrt Q,\\
\varrho x_3&=w_{31}\gamma_1+w_{32}\gamma_2+w_{33}\gamma_3+(\gamma_1l_2-\gamma_2l_1)\sqrt Q.
\end{aligned}
\]
These expressions depend on the quantities \(\gamma\) only formally, as follows from equations (6). If the sign of \(\sqrt Q\) is left undetermined, then these equations give both the \(x\)'s and the \(\xi\)'s, since one merely has to give the root the positive sign once and the negative sign once.

Everything now depends on determining the degree of the function \(Q\). This is done by the following consideration. There are curves of the pencil \(u+\lambda v=0\) which touch the curve \(f=0\), so that the two moving intersection points coincide. For these curves \(Q\) must plainly vanish. We therefore investigate how many such curves there are. Since \(f=0\) is to be touched by a curve \(u+\lambda v=0\), one has, for the point of contact,
\[
\varrho\frac{\partial f}{\partial x_1}=\frac{\partial u}{\partial x_1}+\lambda\frac{\partial v}{\partial x_1},\qquad
\varrho\frac{\partial f}{\partial x_2}=\frac{\partial u}{\partial x_2}+\lambda\frac{\partial v}{\partial x_2},\qquad
\varrho\frac{\partial f}{\partial x_3}=\frac{\partial u}{\partial x_3}+\lambda\frac{\partial v}{\partial x_3}.
\]
Hence
\[
0=\Sigma\pm \frac{\partial f}{\partial x_1}\frac{\partial u}{\partial x_2}\frac{\partial v}{\partial x_3};
\]
the points of contact are thus intersection points of Jacobi's curve with \(f=0\). But not all these intersection points are essential. For, by what preceded, in each simple point which \(u,v\) have in common with \(f\) there fall 2 intersection points of Jacobi's curve with \(f=0\), and even 6 if this point was a double point of \(f\). Now the degree of Jacobi's curve is \(3n-5\), \(u,v\) have \(\frac{n(n-3)}{2}\) double points and \(2n-2\) simple points in common with \(f\); hence the number of essential intersection points, that is, the degree of \(Q\), is
\[
n(3n-5)-2(2n-2)-3n(n-3)=4.
\]
The function \(Q\) is therefore of the fourth degree, while the \(l\)'s become of degree \(n-2\) in \(\lambda\).

We now write \(Q\) in the form
\[
Q=(\lambda-a_1)(\lambda-a_2)(\lambda-a_3)(\lambda-a_4)
\]
and put
\[
y_1=(\lambda-a_1)(a_3-a_2),\qquad
y_2=(\lambda-a_2)(a_3-a_1),
\]
\[
(8)\ldots\quad y_1y_3=\frac{\sqrt Q}{m},\qquad
m^2=\frac{a_4-a_2}{(a_3-a_1)(a_1-a_2)^2(a_3-a_2)^2}.
\]
Then \(w,l\) can be represented as functions \(w',l'\) of \(y_1,y_2\), whose degrees are respectively equal to the orders of \(w,l\); \(\sqrt Q\), up to a constant which may enter into the \(l\)'s, becomes equal to \(y_1y_3\), and one finds the \(x\)'s expressed by the formulas
\[
(9)\ldots\quad
\begin{aligned}
\varrho x_1&=\gamma_1w'_{11}+\gamma_2w'_{12}+\gamma_3w'_{13}
+m(\gamma_2l'_3-\gamma_3l'_2)y_1y_3,\\
\varrho x_2&=\gamma_1w'_{21}+\gamma_2w'_{22}+\gamma_3w'_{23}
+m(\gamma_3l'_1-\gamma_1l'_3)y_1y_3,\\
\varrho x_3&=\gamma_1w'_{31}+\gamma_2w'_{32}+\gamma_3w'_{33}
+m(\gamma_1l'_2-\gamma_2l'_1)y_1y_3.
\end{aligned}
\]
Thus the \(x\)'s are represented as integral functions of order \(n\) of the \(y\)'s; conversely, the \(y\)'s can also be expressed rationally through the \(x\)'s, since \(\lambda=-\frac{u}{v}\) is a rational function of the \(x\)'s, whereby \(y_1,y_2\) are immediately so represented, while \(y_3\) occurs only linearly in equations (9). Finally, from (8), the following equation holds between the \(y\)'s:
\[
(10)\ldots\quad
F=y_3^2y_1-y_2(y_1-y_2)(y_1-k^2y_2)=0,\qquad
k^2=\frac{\frac{a_3-a_2}{a_3-a_1}}{\frac{a_4-a_2}{a_4-a_1}}.
\]
This is the equation of the transformed curve. It is of the third order, and the coordinate system of the \(y\)'s has, with respect to it, a very noteworthy position. For \(y_1=0\) the equation of the curve becomes \(y_2^3=0\); the axis \(y_1=0\) therefore cuts the curve in three infinitely near points, that is, at \(y_1=0,\ y_2=0\) there is a point of inflexion and \(y_1=0\) is its tangent. Through it pass the three lines
\[
y_2=0,\qquad y_1-y_2=0,\qquad y_1-k^2y_2=0,
\]
for which the equation of the curve always passes into
\[
y_3^2y_1=0.
\]
Those three lines are therefore the tangents that can be drawn from the point of inflexion to the curve, and \(y_3=0\) is the line on which the points of contact of the three tangents lie, the harmonic line belonging to the point of inflexion. At the same time \(k^2\) is the cross-ratio of the tangents drawn from the point of inflexion to the curve, including the inflexional tangent itself.

The curves of order \(n\) which arise when the \(x\)'s are set equal to zero must have with curve (10) a certain number \(\sigma\) of fixed points in common. The latter has no double points; hence the degree of the curve \(f=0\), which, as we know, is \(n\), and which again arises from (10) through the substitutions of order \(n\), (9), must be equal to \(3n-\sigma\), or \(\sigma=2n\). These \(\sigma\) points are easy to detect. Take any linear combination of the right-hand sides of (9) and form the equation
\[
(11)\ldots\quad
\Sigma\gamma_i\beta_kw'_{ik}+y_1y_3\Sigma\pm\gamma_1\beta_2l'_3=0.
\]
The curve represented by it must likewise have those \(2n\) constant points of intersection with (10). But near \(y_1=0,\ y_2=0\), equation (11) passes into the pencil of \(n-1\) lines
\[
y_1\cdot \Sigma\pm\gamma_1\beta_2l'_3=0.
\]
Curve (11) therefore has, at the point \(y_1=0,\ y_2=0\), an \((n-1)\)-fold point, one of whose tangents, \(y_1=0\), touches curve (10). Thus \(n\) intersection points of the two curves fall into this point.

Furthermore equation (11) gives
\[
(12)\ldots\quad
(\Sigma\gamma_i\beta_kw'_{ik})^2
=y_1^2y_3^2(\Sigma\pm\gamma_1\beta_2l'_3)^2,
\]
where \(y_3^2\) can now be eliminated by means of (10). But then \(l'_il'_ky_1^2y_3^2\), by the preceding, passes into the minor \(W'_{ik}\) of the \(w'_{ik}\), multiplied by \(-1\), and the right-hand side of (12) into the determinant
\[
-\begin{vmatrix}
W'_{11}&W'_{12}&W'_{13}&\gamma_1&\beta_1\\
W'_{21}&W'_{22}&W'_{23}&\gamma_2&\beta_2\\
W'_{31}&W'_{32}&W'_{33}&\gamma_3&\beta_3\\
\gamma_1&\gamma_2&\gamma_3&0&0\\
\beta_1&\beta_2&\beta_3&0&0
\end{vmatrix}
=(\Sigma w'_{ik}\gamma_i\beta_k)^2
-(\Sigma w'_{ik}\gamma_i\gamma_k)(\Sigma w'_{ik}\beta_i\beta_k).
\]
Equation (12) therefore splits into the two equations
\[
\Sigma w'_{ik}\gamma_i\gamma_k=0,\qquad
\Sigma w'_{ik}\beta_i\beta_k=0.
\]
Only the second of these depends on the special constants \(\beta\) of curve (11); the first supplies \(n\) constant ratios \(y_1:y_2\), and to each of them, from (11), there corresponds a value of \(y_3\), likewise independent of the \(\beta\)'s. Thus from
\[
(13)\ldots\quad \Sigma\gamma_i\gamma_kw'_{ik}=0
\]
in combination with (10), one finds \(n\) further constant points of intersection. As for the value of \(y_3\), if the value found from (11) is really to be independent of the \(\beta\)'s, it is necessary that
\[
y_3=
\frac{w'_{11}\gamma_1+w'_{12}\gamma_2+w'_{13}\gamma_3}{\gamma_3l'_2-\gamma_2l'_3}
=
\frac{w'_{21}\gamma_1+w'_{22}\gamma_2+w'_{23}\gamma_3}{\gamma_1l'_3-\gamma_3l'_1}
=
\frac{w'_{31}\gamma_1+w'_{32}\gamma_2+w'_{33}\gamma_3}{\gamma_2l'_1-\gamma_1l'_2}.
\]
Equating these quotients pairwise, one obtains the three conditions
\[
\begin{aligned}
\gamma_1\Sigma w'_{ik}\gamma_il'_k-l'_1\Sigma w'_{ik}\gamma_i\gamma_k&=0,\\
\gamma_2\Sigma w'_{ik}\gamma_il'_k-l'_2\Sigma w'_{ik}\gamma_i\gamma_k&=0,\\
\gamma_3\Sigma w'_{ik}\gamma_il'_k-l'_3\Sigma w'_{ik}\gamma_i\gamma_k&=0,
\end{aligned}
\]
which, because of (13), reduce to the single equation
\[
\Sigma w'_{ik}\gamma_il'_k=0.
\]
But, multiplied by \(y_1y_3l'_h\), this passes into
\[
\sum_{ik}w'_{ik}\gamma_iW'_{kh}=W\cdot\gamma_h=0,
\]
and is therefore identically satisfied.

The question now arises whether the modulus \(k^2\) is an essential constant, or whether its value can be changed by applying some other transformation.

To answer this question,\footnote{We owe this proof to a written communication from Mr. Cayley. On it, as well as on various matters relating to the subjects treated here, compare Mr. Cayley's paper in the Proceedings of the London Mathematical Society, Oct. 16, 1865.} let us suppose that the curve \(f=0\) is transformed, by an arbitrary other transformation, into the curve of third order
\[
(14)\ldots\quad z_3^2z_1=z_2(z_1-z_2)(z_1-\lambda^2z_2).
\]
Since the \(x\)'s correspond one-to-one to the \(y\)'s as well as to the \(z\)'s, these latter also correspond one-to-one to each other, and the new transformation may therefore be composed of our original one and a transformation which changes the curve of third order (10) into the curve of third order (14). The question is whether the value of \(\lambda^2\) can thereby be made arbitrarily large.

The quantity \(\lambda^2\) represents the cross-ratio of the tangents drawn from the point \(z_1=0,\ z_2=0\) to curve (14), just as \(k^2\) represents the corresponding cross-ratio for the curve \(F=0\). By a known theorem from the theory of curves of third order, this cross-ratio is known to remain unchanged from whichever point of the curve the pencil of tangents is drawn.

If the transformation leading from (10) to (14) is
\[
z_i=\varphi_i(y_1,y_2,y_3)
\]
of order \(\varrho\), then the curves \(\varphi=0\) must pass through \(3\varrho-3\) fixed points of curve (10), in order that the final result be of the third order. A line
\[
(15)\ldots\quad \alpha_1z_1+\alpha_2z_2+\alpha_3z_3=0
\]
corresponds, under this transformation, to a curve of order \(\varrho\),
\[
(16)\ldots\quad \alpha_1\varphi_1+\alpha_2\varphi_2+\alpha_3\varphi_3=0,
\]
in such a way that the moving intersection points of the one curve with (14) correspond to those of the other with (10). If line (15) is allowed to rotate about a point, so that one of the moving intersection points is fixed, then one obtains the line pencil
\[
(17)\ldots\quad Z+\lambda Z'=0,
\]
and to it, from (16), there corresponds a pencil of curves of order \(\varrho\),
\[
(18)\ldots\quad Y+\lambda Y'=0,
\]
which arises by fixing the corresponding moving point of intersection of curve (16). The forms (17) and (18) are projective, in so far as the same ranging element \(\lambda\) occurs in both and connects the corresponding elements.

Since the two remaining intersection points of a line of the pencil (17) always correspond to the two remaining intersection points of a curve of the pencil (18), both pairs of points coincide simultaneously. Thus to the 4 tangents which can be drawn from the fixed point on (14) to this curve there correspond 4 curves of the pencil (18) which touch the curve \(F=0\); and since the quantities \(\lambda\) of the corresponding curves and lines are the same, the cross-ratio of the four contact curves is equal to that of the four tangents. The latter remained unchanged when the point fixed on (14) is allowed to move over the curve. Thus the cross-ratio of the four contact curves likewise remains unchanged when the fourth point moves over the curve \(F=0\). And what holds of this point must likewise hold of the other points of the pencil. Hence one has the theorem:

\emph{If through any \(3s-2\) points of a curve of third order one lays a pencil of curves of order \(s\)}\footnote{The pencil is not thereby in general completely determined, but the families of intersection points which its curves have in common with the curve of third order are.}\emph{, then in it there are 4 curves which touch the curve of third order, and their cross-ratio is independent of the position of those \(3s-2\) points.}

If now \(3s-3\) of those points are so placed that they form the complete intersection of a curve of order \(s-1\) with the curve of third order, then the pencil of curves of order \(s\) passes into a fixed curve of order \(s-1\) (respectively its system of intersections with the curve of third order) and into a pencil of lines whose vertex is the last fixed point. It follows further:

\emph{That cross-ratio is also independent of the number \(s\), and is equal to the cross-ratio of the tangents which can be drawn from a point of the curve of third order to the curve itself.}

Thus \(k^2\) really has an absolutely invariant value, and \(\lambda^2\) must either be equal to \(k^2\), or else equal to one of the five other cross-ratios
\[
\left(\frac1{k^2},\ 1-k^2,\ \frac1{1-k^2},\ \frac{k^2-1}{k^2},\ \frac{k^2}{k^2-1}\right),
\]
which can be formed from the same elements.

Thus the curves for which \(p=1\) possess an absolute invariant.

\subsection*{\S\,21. Treatment of the case \(p=2\). Normal curve: a curve of fourth order with one double point.}

The reduction of curves with \(\frac{(n-1)(n-2)}{2}-2\) double and return points \((p=2)\) proceeds in a quite similar way to the case \(p=1\). Take a pencil of curves of order \(n-1\), passing through the double and return points and through \(2n\) further arbitrarily chosen points on the curve \(f=0\), and denote it by
\[
u+\lambda v=0.
\]
Then the curves of this pencil again have only two moving intersection points in common with \(f=0\). If, therefore, one eliminates the \(x\)'s from the equations
\[
f=0,\qquad u+\lambda v=0,\qquad \alpha_1x_1+\alpha_2x_2+\alpha_3x_3=0,
\]
one again obtains, after disregarding the factors independent of \(\lambda\), the product of the equations of the two moving intersection points in the form
\[
\Sigma w_{ik}\alpha_i\alpha_k=0,
\]
where the \(w\)'s are integral functions of order \(n\) of \(\lambda\), whose determinant vanishes. Hence, exactly as before, the following expressions in \(\lambda\) hold for the \(x\)'s:
\[
\begin{aligned}
\varrho x_1&=w_{11}\gamma_1+w_{12}\gamma_2+w_{13}\gamma_3+(\gamma_2l_3-\gamma_3l_2)\sqrt Q,\\
\varrho x_2&=w_{21}\gamma_1+w_{22}\gamma_2+w_{23}\gamma_3+(\gamma_3l_1-\gamma_1l_3)\sqrt Q,\\
\varrho x_3&=w_{31}\gamma_1+w_{32}\gamma_2+w_{33}\gamma_3+(\gamma_1l_2-\gamma_2l_1)\sqrt Q,
\end{aligned}
\]
where the \(\gamma\)'s are arbitrary quantities and \(Q,l_1,l_2,l_3\) are rational functions of \(\lambda\), so that
\[
l_il_kQ=-W_{ik}.
\]
But here the degree of the function \(Q\) is different from before. It is obtained by diminishing the number of intersection points of \(f=0\) with Jacobi's curve
\[
\Sigma\pm\frac{\partial f}{\partial x_1}\frac{\partial u}{\partial x_2}\frac{\partial v}{\partial x_3}=0
\]
by twice the number of fixed simple points which \(u,v\) have in common with \(f\), and by six times the number of double points; the degree of \(Q\) is therefore
\[
n(3n-5)-4n-6\left(\frac{(n-1)(n-2)}{2}-2\right)=6.
\]
Set
\[
Q=(\lambda-a_1)(\lambda-a_2)(\lambda-a_3)(\lambda-a_4)(\lambda-a_5)(\lambda-a_6)
\]
and then
\[
y_1=(\lambda-a_1)(a_3-a_2),\qquad y_2=(\lambda-a_2)(a_3-a_1),
\]
\[
y_1y_2y_3=\frac{1}{m}\sqrt Q,\qquad
m^2=-\frac{(a_4-a_2)(a_5-a_2)(a_6-a_2)}
{(a_3-a_1)(a_3-a_1)^4(a_1-a_2)^4}.
\]
Then the functions \(w,l\) can be represented as homogeneous functions \(w',l'\) of \(y_1\) and \(y_2\), and for the \(x\)'s one obtains the rational expressions in the \(y\)'s
\[
(1)\ldots\quad
\begin{aligned}
\varrho x_1&=\gamma_1w'_{11}+\gamma_2w'_{12}+\gamma_3w'_{13}
+m(\gamma_2l'_3-\gamma_3l'_2)y_1y_2y_3,\\
\varrho x_2&=\gamma_1w'_{21}+\gamma_2w'_{22}+\gamma_3w'_{23}
+m(\gamma_3l'_1-\gamma_1l'_3)y_1y_2y_3,\\
\varrho x_3&=\gamma_1w'_{31}+\gamma_2w'_{32}+\gamma_3w'_{33}
+m(\gamma_1l'_2-\gamma_2l'_1)y_1y_2y_3,
\end{aligned}
\]
while between the \(y\)'s themselves there is the equation
\[
(2)\ldots\quad
y_1y_2y_3^2=(y_1-y_2)(y_1-k^2y_2)(y_1-\lambda^2y_2)(y_1-\mu^2y_2),
\]
where
\[
k^2=\frac{\frac{a_3-a_2}{a_3-a_1}}{\frac{a_4-a_2}{a_4-a_1}},
\qquad
\lambda^2=\frac{\frac{a_3-a_2}{a_3-a_1}}{\frac{a_5-a_2}{a_5-a_1}},
\qquad
\mu^2=\frac{\frac{a_3-a_2}{a_3-a_1}}{\frac{a_6-a_2}{a_6-a_1}}.
\]
Equation (2) is the equation of the transformed curve; specifically, it is the equation of a curve of order \(4\) with the double point \(y_1=0,\ y_2=0\). The axes \(y_1=0,\ y_2=0\) themselves are the tangents of the double point, while the 4 lines
\[
y_1-y_2=0,\qquad y_1-k^2y_2=0,\qquad
y_1-\lambda^2y_2=0,\qquad y_1-\mu^2y_2=0
\]
are the four tangents that can be drawn from the double point to the curve; their points of contact lie on the line \(y_3=0\).

Instead of the system of curves of order \(n-1\), one may also use a system of order \(n-2\), passing through all double and return points and through \(n\) fixed points of \(f=0\).

This system of curves arises from the preceding one when \(n\) of the fixed points lie in a straight line, so that the system of order \(n-1\) passes into this fixed line itself and into a system of order \(n-2\). Finally, one may use a system of order \(n-3\), which passes through all double and return points except one. This arises when, in the system of order \(n-1\), \(2n\) of the fixed points lie on a conic; thereby the system of order \(n-1\) passes into this conic and into a system of order \(n-3\) having all double and return points as base-points.

\begin{center}
{\large Fourth Section.}\\[0.35em]
{\Large Periodicity.}
\end{center}

\subsection*{\S\,22. Branch points.}

Represent a value of the independent variable \(x\) by a point of the \(XY\)-plane, whose coordinates are the real and imaginary parts of the quantity \(x=X+Y\sqrt{-1}\). To each point of the plane there correspond \(n\) values of \(s_x\), obtained from the equation \(F(s_x,x)=0\). These values are all distinct, unless \(x\) is one of the values for which simultaneously
\[
F=0,\qquad \frac{\partial F}{\partial s_x}=0.
\]
The points corresponding to such values are called branch points;\footnote{Riemann in Crelle's Journal, vol. 54.} in them, in general, two values of \(s_x\) coincide. We shall confine ourselves to the consideration of this case, since all others follow from it as limiting cases.\footnote{A precise investigation of all possible cases is found in Puiseux, Recherches sur les fonctions algébriques, in Liouville's Journal, vol. 15.}

If at an arbitrary place in the \(XY\)-plane one chooses one of the values of \(s_x\), and lets it vary continuously on a closed curve, then, as Puiseux has shown, one always returns to the original value of \(s_x\), provided that the closed curve encloses no branch point. If branch points are enclosed, however, then under some circumstances one comes to another value of \(s_x\). This value, in turn, remains unchanged when the closed curve is varied so that, during the deformation of the curve, no branch point enters or leaves.

If only one branch point is enclosed, one may arrive at another root \(s_x\). We may then contract the closed curve into a path \(A\) leading to the branch point, an infinitely small circle around it, and the path \(A\) in the reverse direction. Suppose we begin with the value \(s'_x\), and suppose that, while traversing path \(A\) toward the branch point, it has become \(s_x^\circ\), the value which is to correspond to the branch point as the associated double root. If we break off path \(A\) at distance \(\varepsilon\) from the branch point in order to traverse the small circle, then on it \(s_x\) differs very little from \(s_x^\circ\), and \(x\) very little from the branch value \(x^\circ\); and if by \(\varphi\) one denotes the angle which a radius of the circle makes with the \(X\)-axis, one has
\[
x-x^\circ=\varepsilon e^{i\varphi}.
\]
Furthermore, by what was said earlier, \(F(s_x,x)=0\) passes into
\[
\frac{\partial F^\circ}{\partial x^\circ}(x-x^\circ)
+\frac{1}{2}\frac{\partial^2F^\circ}{\partial (s_x^\circ)^2}(s_x-s_x^\circ)^2=0,
\]
so that
\[
s_x-s_x^\circ=\operatorname{Const.}\sqrt{x-x^\circ}
=\operatorname{Const.}\sqrt\varepsilon\, e^{\frac{i\varphi}{2}}.
\]
When the circle is traversed, \(\varphi\) increases by \(2i\pi\), and one obtains, when one has returned to path \(A\),
\[
s_x-s_x^\circ=-\operatorname{Const.}\sqrt\varepsilon\, e^{\frac{i\varphi}{2}}.
\]
Thus the value of \(s_x\) has changed during the circuit around the circle; and therefore, by following path \(A\) back, one cannot arrive at the original value \(s'_x\), but must arrive at another value \(s''_x\). Conversely, starting from \(s''_x\), one returns to \(s'_x\).

Now fix in the \(XY\)-plane an arbitrary point \(O\), and from it draw straight or curved, mutually non-intersecting lines \(A\) to the branch points. Then every branch point is assigned to two of the values of \(s_x\) that belong to point \(O\), to be denoted by \(s^{(1)},s^{(2)},\ldots,s^{(n)}\); namely, to those which pass into one another by describing a circuit around the branch point. On the other hand, if a circuit around a branch point is begun with a value of \(s_x\) not assigned to it, the branch point behaves completely neutrally, and one returns to the same value of \(s_x\). If the curves \(A\) are altered, the kind of assignment changes as soon as one of the curves crosses a branch point during the alteration.

This behavior of the branch points rests essentially on the fact that \(s_x-s_x^\circ\) near such a point becomes proportional to \(\sqrt{x-x^\circ}\). If, instead of a branch point corresponding to a point of contact (\S\,3), one takes a point corresponding to a double point of the curve \(F(s_x,x)=0\), then near such a point \(s_x-s_x^\circ\) is proportional to \(x-x^\circ\) itself, and the characteristic property of the branch point just described is lost. We shall therefore not count such points among the branch points; but we shall count those corresponding to return points, in which \(s_x-s_x^\circ\) is proportional to \(\sqrt{x-x^\circ}^{\,3}\). Since the number of contact points is \(n(n-1)-2d-3r\), the number of branch points is therefore
\[
\nu=n(n-1)-2d-2r=2p+2(n-1).
\]

The path toward a branch point, around it in a small circle, and back again shall, for brevity, be called a loop. Each loop connects two definite roots belonging to point \(O\) with one another.

\subsection*{\S\,23. Cycles and circuits.}

The equation \(F(s_x,x)=0\), considered as an equation for \(s_x\), is necessarily irreducible unless the entire curve is to split into parts. Hence one may state the theorem:

\emph{The position of the wholly arbitrary point \(O\) may always be supposed chosen so that, by a suitable series of loops, one can pass from any arbitrarily prescribed value of \(s_x\) occurring at \(O\) to any other arbitrarily prescribed one.}

For if this were not the case, there would have to be a cycle of roots \(s_x^{(\alpha)},s_x^{(\beta)},\ldots\) which always pass only into one another, whatever loops one traverses. Every symmetric function of \(s_x^{(\alpha)},s_x^{(\beta)},\ldots\), if allowed to vary continuously along any path, or along the loops into which it can be contracted, would then have the same value at every arbitrarily chosen point \(O\), and so would possess only a single value everywhere in the \(XY\)-plane; consequently it would be a rational function of \(x\).\footnote{Briot et Bouquet, Théorie des fonctions doublement périodiques.} One could therefore set up an equation whose coefficients would be rational functions of \(x\), and whose roots would be only certain of the \(n\) roots \(s_x\). This equation would then have to be a factor of \(F(s_x,x)\), which is excluded.

Thus if we single out any branch point, assigned for example to the roots \(s^{(1)},s^{(2)}\), there will always be a root \(s^{(3)}\) connected by a loop either with \(s^{(1)}\) or with \(s^{(2)}\); then a root \(s^{(4)}\) connected with \(s^{(1)},s^{(2)}\), or \(s^{(3)}\); and so on. We choose \(n-1\) branch points as fundamental points of the first kind, having the property that one leads, say, from \(s^{(1)}\) to \(s^{(2)}\), the second from \(s^{(1)}\) or \(s^{(2)}\) to \(s^{(3)}\), the third from \(s^{(1)},s^{(2)}\), or \(s^{(3)}\) to \(s^{(4)}\), and so forth. Every other branch point forms a cycle with some of these, in so far as the loop of any other branch point, leading say from \(s^{(i)}\) to \(s^{(k)}\), can be connected, in one and only one way, with suitably chosen loops of the fundamental points, so that by traversing the whole path thus formed one returns to the initial value of \(s\). To make this clear, we first combine the fundamental loops into paths \(W_2,W_3,\ldots,W_n\) such that \(W_2\) leads from \(s^{(1)}\) to \(s^{(2)}\), \(W_3\) from \(s^{(1)}\) to \(s^{(3)}\), and so on. The loop leading from \(s^{(i)}\) to \(s^{(k)}\) then forms a cycle with \(W_i\) and \(W_k\), the path \(W_1\), corresponding to \(s^{(1)}\), being omitted.

From the fundamental points of the first kind themselves, no cycle can evidently be composed. Any cycle whatever, however, can be resolved into such cycles as contain only one other branch point. For by supplementing each loop of the given cycle, except fundamental loops, with fundamental loops so as to produce such simple cycles, one obtains the given cycle as equal to the sum of the new ones, increased or diminished by certain fundamental loops. But the aggregate of these would itself have to form a cycle, which can never occur; hence these fundamental points must simply drop out, since the loop of each of them occurs in the positive and in the negative sense. Thus the given cycle becomes equal to a sum of simple cycles.

Among the cycles there are some which are especially distinguished: namely those which arise from complete circuits. By a complete circuit we understand a path which comprises all loops in order; and we think of the loops as ordered in the way in which the tangents of the paths drawn from point \(O\) follow one another at \(O\) itself, when one moves in a small circle about \(O\). The whole path of a circuit can then be contracted through infinity to the point \(O\), without crossing a branch point, whence it follows that every circuit returns to the initial value of \(s_x\). There are therefore \(n\) separate circuits, each beginning with one of the values of \(s_x\) occurring at \(O\) and ending with the same one. In such a circuit some branch points are essential, that is, they lead from one root to another; others are inessential. In two different circuits the same loop cannot be traversed with the same root-values, for otherwise all preceding and following loops, and hence also the whole circuits, would have to agree. Thus each branch point can be essential in only two circuits, by leading in the one from one of its assigned roots \(s^{(i)}\) to the other \(s^{(k)}\), and in the other circuit conversely from \(s^{(k)}\) to \(s^{(i)}\). The cycles of the \(n\) circuits therefore consist of the loops of their essential branch points, and every circuit can be replaced by a sum of simple cycles. Since the whole circuit, by


% --- Continuation: source PDF pp. 107--121. ---
% English translation corresponding to German source PDF pp. 107--121.

continuous change to a point, and thereby its effect becomes identically zero, one may also say that one of these simple cycles can be expressed by the remaining ones of the same circuit. It could therefore seem as though one might express \(n\) cycles which occur in different circuits by the others; in fact this is possible for only \(n-1\) of them. Namely, the path which consists of all circuits contains each loop twice essentially, once leading from \(s^{(i)}\) to \(s^{(k)}\), and the other time from \(s^{(k)}\) to \(s^{(i)}\). The different loops of this total circuit therefore directly cancel each other, without needing to call on a displacement of the path. Hence the last circuit is nothing other than a combination of the earlier ones, and therefore can give occasion to no new relation between the cycles. Since we have had \(2p+n-1\) simple cycles, and since \(n-1\) relations exist between them, we have the theorem:

\begin{center}
\emph{There are \(2p\) simple cycles independent of one another.}
\end{center}

\subsection*{\S\,24. Reciprocity between two classes of circuits and between two classes of fundamental points.}

Between the \(n\) circuits on the one hand and the \(n\) values of \(s_x\) occurring at the point \(O\) on the other hand there is now a remarkable reciprocity, which can be represented as follows.

Connect the branch points in the order given above by a broken line, which begins at the first and ends at the last. The loops \(a_1,a_2\ldots a_{\nu}\) drawn to the several branch points then all lie on the same side of this broken line. Now describe further paths \(A_1,A_2\ldots A_{\nu}\), which, starting from \(O\), always exclude only the one branch point indicated by the index, but enclose all the others. Such a path is always assigned to two circuits, namely to the two in which the excluded branch point occurs as an essential one. If the loop belonging to the branch point connects the roots \(s^{(i)}\) and \(s^{(k)}\), and if the point is essential for the circuits beginning with \(s^{(\alpha)}\) and \(s^{(\beta)}\), then first the corresponding path \(A\), when it is not begun with one of the roots \(s^{(\alpha)},s^{(\beta)}\), always returns to the initial root, since the excluded branch point is then immaterial, and the path \(A\) can therefore be replaced by a complete circuit. But if one begins the path \(A\) with the root \(s^{(\alpha)}\), then up to the excluded loop everything agrees exactly with the \(\alpha^{\text{th}}\) circuit, up to the place where that loop would have carried \(s^{(i)}\) into \(s^{(k)}\). Instead, one proceeds with \(s^{(i)}\) to the next loop, and so behaves as in the \(\beta^{\text{th}}\) circuit, in which conversely the excluded loop would have ended with \(s^{(i)}\). Thus, when begun with \(s^{(\alpha)}\), the path \(A\) consists of the beginning of the \(\alpha^{\text{th}}\) and the end of the \(\beta^{\text{th}}\) circuit; and similarly, when it is begun with \(s^{(\beta)}\), it consists of the beginning of the \(\beta^{\text{th}}\) and the end of the \(\alpha^{\text{th}}\) circuit.

This path \(A\) is therefore assigned to the circuits \(\alpha\) and \(\beta\), just as the loop corresponding to the branch point is assigned to the roots \(i\) and \(k\). But the path \(A\) can be brought by continuous change into a loop which behaves similarly to the loop \(a\), but lies on the other side of the broken line. \emph{The system of the new loops \(A\) is therefore completely analogous to the system of the original loops \(a\); and the two are conjugate with respect to the broken line in such a way that the loop of one system connects the roots \(s^{(\alpha)},s^{(\beta)}\) in whose corresponding circuits the corresponding loop of the other system is essential, and conversely.}

To the theorem stated above, that one can pass by loops of the first kind from every root \(s^{(i)}\) to every other \(s^{(k)}\), corresponds therefore the theorem that by loops of the second kind one can pass from every circuit into every other one; and consequently also by loops of the first kind, since every loop of one kind can be represented by loops of the other kind. To the fundamental points of the first kind there correspond also \(n-1\) fundamental points of the second kind, which stand in the same relation to the second system of loops as the former stand to the first, and which therefore have to the first system the relation that their branch points connect all circuits of that first system with one another.

But one can further prove the remarkable theorem \emph{that the two sequences of fundamental points can always be chosen so that no fundamental point of one kind belongs at the same time to the other kind.} This is shown in the following way.

In a circuit in which the first fundamental point of the first kind occurs, since the points of the circuit form a cycle and therefore cannot be only fundamental points of the first kind, at least one other essential branch point must be present. We choose it as the first fundamental point of the second kind. This point connects the first circuit with a second. One next proves, secondly, that in the first and second circuits together there must be at least one point which is not a fundamental point of the first kind, and which occurs in only one of the two circuits. This point, as the second fundamental point of the second kind, connects the first two circuits with a third. Then one shows, thirdly, that in all three circuits together there occurs at least one point which appears only once and is not a fundamental point of the first kind. This point, as the third fundamental point of the second kind, connects the first three circuits with a fourth. Continuing in this way, one arrives at a complete system of fundamental points of the second kind which has no point in common with the first system. Here it remains only to prove the following auxiliary theorem:

\begin{center}
\emph{If \(i\) circuits are already connected by fundamental points of the second kind, then in them, besides fundamental points of the first kind, there is at least one point which occurs in only one of these circuits.}
\end{center}

The proof of this theorem is very simple. Take a symmetric function of the roots \(s^{(1)},s^{(2)}\ldots s^{(i)}\) with which these \(i\) circuits begin, and let the \(s\)'s vary continuously while we go around all branch points. After describing the circuit we arrive at the original values of the \(s\)'s, and therefore also back at the value of the function. But for this circuit and this symmetric function all branch points that occur essentially in two of the \(i\) circuits are inessential, since they only simultaneously carry \(s^{(k)}\) into \(s^{(h)}\) and \(s^{(h)}\) into \(s^{(k)}\), and thus do not change the symmetric function. The circuit formed for the symmetric function, however, must possess essential branch points, because otherwise the function would be rational, contrary to what was shown earlier. Hence there must be branch points occurring singly in the \(i\) circuits. But these necessarily form a sequence of cycles\footnote{Namely, one can state the following theorems: \emph{Two cyclic paths in which a branch point is common, and whose loop is traversed in opposite directions in the two paths, can be united into one cyclic path in which that doubled branch point has disappeared.} \emph{A cyclic path in which a branch point occurs twice, so that its loop is traversed the two times in opposite directions, can be divided into two cyclic paths in which that loop no longer occurs.} With the aid of these two theorems one can plainly always reduce any number of circuits to cycles from which the branch points common to the circuits have been removed. The first theorem is proved as follows. Let the first path lead from \(s^{(\alpha)}\) in some way to \(s^{(i)}\), where the common loop begins, and, after traversing that loop, from \(s^{(k)}\) back to \(s^{(\alpha)}\). Let the other lead from \(s^{(\beta)}\) to \(s^{(k)}\), and after traversing the common loop from \(s^{(i)}\) back to \(s^{(\beta)}\). The united path, in which that loop drops out, then goes like the first from \(s^{(k)}\) through \(s^{(\alpha)}\) to \(s^{(i)}\), and from there, attaching itself to the second, from \(s^{(i)}\) through \(s^{(\beta)}\) back to \(s^{(k)}\). Similarly, in the second theorem let the path lead from \(s^{(\alpha)}\) in some way to \(s^{(i)}\), after traversing the loop from \(s^{(k)}\) through \(s^{(\beta)}\) to \(s^{(k)}\), and finally after traversing the loop once more from \(s^{(i)}\) back to \(s^{(\beta)}\). One cyclic path then goes from \(s^{(k)}\) through \(s^{(\beta)}\) back to \(s^{(k)}\), the other from \(s^{(i)}\) through \(s^{(\alpha)}\) back to \(s^{(i)}\), thereby accomplishing the required decomposition.}. Since fundamental points of the first kind cannot form a cycle, at least one other point must be present in this remainder, as was to be proved.

\subsection*{\S\,25. Integration along different paths. Periodicity moduli of the integrals of the first kind.}

We shall now examine the paths considered above in connection with the continuous change of the function \(s_x\), taking them as paths of integration of an integral of the first kind. Such an integral is completely determined when the initial and final values of \(s_x\) and \(x\), together with the path by which one passes from one value system to the other, are specified. As the initial value of \(x\) we can always take the one which belongs to the point \(O\), and as initial value of \(s_x\) one of the corresponding values of this quantity. Two paths of integration which can be transformed into one another through the finite or the infinite without crossing a branch point always yield the same result as the value of the integral. But if they enclose branch points, and neither all nor none of them, so that a direct transformation of one path of integration into the other is not possible, then in general one has two different values of the integral, although the pairs of values \(x,s_x\) for the upper and lower limits are the same. One can then put the first path in place of the second, augmented by a closed curve which contains the second path of integration in the direct direction and the first in the opposite direction; thus one can think of the second path of integration as arising from the first by adding a cyclic path. Now, as was shown above, all cyclic paths can be composed from \(2p\) of them, corresponding to the \(2p\) independent simple cycles. The most general path of integration which can supply new values of the integral is therefore to be thought of as arising from some one path by adding the simple cyclic paths arbitrarily often, positively or negatively. Thus, if we call the value of the integral extended over a cyclic path a \emph{periodicity modulus} of the integral, it first follows \emph{that every periodicity modulus is expressible linearly with integral coefficients by \(2p\) simple periodicity moduli, or that an integral of the first kind has \(2p\) independent periodicity moduli.} Further, if we call \(I\) the integral taken along an arbitrary path, and \(I',I^{(2)}\ldots I^{(2p)}\) its periodicity moduli, then
\[
I+m'I'+m^{(2)}I^{(2)}\cdots+m^{(2p)}I^{(2p)}
\]
is the most general value which the integral can obtain by change of the path of integration, where \(m',m^{(2)},\ldots,m^{(2p)}\) are positive or negative whole numbers.

We shall denote the \(p\) independent integrals of the first kind, taken along any common path between the same limits, by \(I_1,I_2,\ldots,I_p\). If we then change the path of integration, the values of all the integrals change simultaneously; the numbers \(m\), which depend only on the nature of the new path, remain the same for all integrals, while the periodicity moduli have individual values for each integral. The general value system of the integrals \(I\) along any path is therefore represented by the schema
\[
\begin{gathered}
I_1+m'I_1'+m^{(2)}I_1^{(2)}\cdots+m^{(2p)}I_1^{(2p)},\\
I_2+m'I_2'+m^{(2)}I_2^{(2)}\cdots+m^{(2p)}I_2^{(2p)},\\
\vdots\\
I_p+m'I_p'+m^{(2)}I_p^{(2)}\cdots+m^{(2p)}I_p^{(2p)}.
\end{gathered}
\]
The periodicity moduli form a system of \(2p^2\) constants, in which the \(2p\) quantities of the same horizontal row refer to the same integral and to different cyclic paths, while the \(p\) quantities of the same vertical column refer to the same cyclic path and to different integrals.

Algebraic relations exist between the quantities \(I_h^{(k)}\); in particular, between the quantities \(I_h^{(k)}\) of two horizontal rows there are equations which are linear in the \(I_h^{(k)}\) of each of the two rows. These are now to be developed.

\subsection*{\S\,26. Relations between the periodicity moduli of the integrals of the first kind.\footnotemark}\footnotetext{The relations to be established here by no means comprise all existing ones, but contain only those which are necessary for the course of the theory.}

Let \(I\) and \(H\) denote any two of the \(p\) integrals of the first kind, and carry the integral \(\int I\,dH\) over all \(n\) complete circuits. The function to be integrated, \(I\dfrac{\partial H}{\partial x}\), is nowhere discontinuous and has no branch points except those of \(F(s_x,x)=0\). Therefore the integral taken over one circuit is equal to zero, and so also is the integral which we are considering. It remains to express the left hand side of the equation
\[
\int I\,dH=0
\]
in another way. This integral first splits into \(n\) parts, corresponding to the \(n\) circuits; each part again into \(\nu\) parts, the integrals over the individual loops. One immediately sees that each loop gives the value zero as its contribution to the integral, whenever the loop was inessential in the relevant circuit. For if we consider the outward and return path which must be made between \(O\) and the branch point for such a loop, then at the same place \(I\) has both times the same value, whereas \(dH\) has opposite values. The elements of the outward and return paths thus exactly cancel one another. Therefore only the essential loops of the circuits remain to be considered, so that each loop is to be taken into account only twice. In both cases the same values of \(s_x\) and \(x\) are traversed only in the opposite order, since the loop leads once from \(s^{(i)}\) to \(s^{(k)}\), and the other time from \(s^{(k)}\) to \(s^{(i)}\). If we again compare two elements of such integrals corresponding to the same loop, for which \(s_x\) and \(x\) have the same values, then \(dH\) has equal and opposite values for the two; hence, if we denote the two values of \(I\) by \(I\) and \(I'\), the element of our integral is
\[
(I-I')\,dH.
\]
Now \(I\) is the integral taken from the beginning of the first circuit in which the branch point under consideration is essential, first up to the value \(s^{(i)}\) with which the loop of the branch point begins, and then from there to the variable point of the loop. Similarly, \(I'\) is the integral taken from the beginning of the second circuit in which the branch point is essential, first up to \(s^{(k)}\), with which the loop begins in this circuit, and then on to the variable point. In place of \(-I'\), one can therefore put the remainder of this circuit, that is, first the integral from the variable point to the value \(s^{(i)}\) with which the loop ends in this circuit, and then from there further to the end of the circuit. Thus \(I-I'\), since the integrals from the variable point to \(s^{(i)}\) cancel, becomes an integral over a path which encloses all branch points except the one under consideration; a path which in the preceding section we transformed into a loop of the second kind.

Let the values of the integral \(I\), taken over the loops so that one passes from a root \(s\) with smaller index to a root \(s\) with larger index, be denoted by \(a_1,a_2\ldots a_{\nu}\), and let the corresponding values taken over loops of the second kind be denoted by \(A_1,A_2\ldots A_{\nu}\), in such a way that, when \(A_i\) is expressed by the \(a\)'s as above, the path of \(A\) begins with that circuit in which the loop \(+a_i\) occurs. Then the element just examined of the \(i^{\text{th}}\) loop becomes
\[
(I-I')\,dH=A_i\,dH;
\]
hence the contribution of this loop to the integral is \(A_i\int dH\), and if the integrals taken for \(H\), corresponding to \(a_i,A_i\), are denoted by \(b_i,B_i\), this contribution is \(A_i b_i\). The above equation therefore becomes
\[
(1)\quad A_1b_1+A_2b_2+\cdots+A_{\nu}b_{\nu}=0.
\]
We may also write this equation in the form
\[
(2)\quad B_1a_1+B_2a_2+\cdots+B_{\nu}a_{\nu}=0.
\]
For if we add both equations and return to their derivation, we have
\[
\Sigma(A_i b_i+B_i a_i)=\int(I\,dH+H\,dI).
\]
The indefinite integral of the right hand side of this expression gives \(H.I\); and since both \(H\) and \(I\) have the same value at the beginning and at the end of each circuit, the integral on the right vanishes over each single circuit, and therefore also over all circuits taken together. Thus one sees that the expressions (1), (2) must be transformable into one another by the sole use of the linear relations arising between the \(a\)'s, the \(b\)'s, the \(A\)'s, and the \(B\)'s, which state that the integral \(I\) or \(H\), extended over all loops of the same circuit, has the value zero.

Equations (1), (2) hold between integrals extended over the separate loops. They can be transformed so that they contain only integrals extended over cyclic paths, that is, periodicity moduli. This is achieved by three successive transformations.

In place of the quantities \(b_1,b_2\ldots\), we first introduce the integrals \((b_1),(b_2)\ldots\), corresponding to the relevant branch points and extended over the simple cyclic paths. Form, as combinations of the integrals over fundamental loops of the first kind, first \(n-1\) quantities \(f_2,f_3\ldots f_n\), whose paths lead respectively from \(s^{(1)}\) to \(s^{(2)}\), from \(s^{(1)}\) to \(s^{(3)}\), \(\ldots\), from \(s^{(1)}\) to \(s^{(n-1)}\) (compare \S\,23). Then, if \(b_i\) leads from \(s^{(\alpha)}\) to \(s^{(\beta)}\), the integral taken over the corresponding simple cyclic path, the periodicity modulus, is
\[
(b_i)=b_i+f_{\alpha}-f_{\beta};
\]
here \(f_1\) is always to be replaced by zero. If in (1) we therefore put
\[
b_i=(b_i)-f_{\alpha}+f_{\beta},
\]
then on the left hand side of (1) each quantity \(f_{\alpha}\) is multiplied by all those \(A\)'s for whose branch point \(s^{(\alpha)}\) is one of the assigned roots; and the corresponding \(A\) is negative when the positive loop of the first kind of the branch point begins with \(s^{(\alpha)}\), and positive in the opposite case. The loops of the second kind belonging to these \(A\)'s therefore form a complete circuit of the second kind, or their sum vanishes. Thus one obtains the transformation
\[
(3)\quad \Sigma A_i b_i=\Sigma A_i(b_i).
\]
The sum on the left contains \(n-1\) terms fewer than the sum on the right, since all those \((b_i)\) vanish identically which themselves belong to the fundamental points of the first kind.

Quite similarly, the linear equations existing between the \(a\)'s or \(b\)'s, which say that an integral taken over a circuit vanishes, have the property that in them one can replace the \(a\)'s or \(b\)'s by the \((a)\)'s or \((b)\)'s. This follows already from the proof given in \S\,23, that in a circuit one can omit the fundamental points of the first kind and replace the other points by the corresponding simple cycles of the first kind.

In order to carry out the second transformation, suppose now that the quantities \(A\) have been expressed as linear functions of the \(a\)'s, according to how the paths of the \(A\)'s were originally composed of loops of the first kind. If in these expressions the corresponding periods \((a)\) are substituted for the \(a\)'s, let the \(A\)'s pass into expressions \((A)\). Let the combinations of the \(a\)'s corresponding to the combinations of the \(b\)'s denoted above by \(f\) be denoted by \(g\), and if the circuit \(A_i\) begins with \(s^{(\alpha)}\) and ends with \(s^{(\beta)}\), then
\[
(A_i)=A_i+g_{\alpha}-g_{\beta}.
\]
Thus, if in the right hand side of equation (3) one puts
\[
A_i=(A_i)-g_{\alpha}+g_{\beta},
\]
then \(g_{\alpha}\) is multiplied there by all those \((b_i)\) whose branch points belong to the \(\alpha^{\text{th}}\) circuit; namely negatively when the corresponding loop of the first kind occurs in that circuit positively, and positively when the opposite occurs. Therefore all \((b_i)\)'s multiplied by \(g_{\alpha}\) form a complete circuit; their algebraic sum vanishes by the preceding result, and one has
\[
(4)\quad \Sigma A_i b_i=\Sigma(A_i)(b_i).
\]
We can now transform the expression (4) for a third time by eliminating, with the help of the conditional equations holding between the \((b_i)\), \(n-1\) of these quantities. According to the theory established above concerning the fundamental points of the second kind, the corresponding quantities \((b_i)\) can always be determined from the linear conditional equations. For in each following circuit there was at least one new fundamental point of the second kind, and therefore also a new quantity \((b_i)\), whose coefficient is \(\pm1\), and which therefore can be represented as a linear function of the remaining \((b_i)\)'s with integral coefficients. If these determinations are supposed to have been carried out, one finally has
\[
(5)\quad \Sigma A_i b_i=\Sigma[A_i](b_i),
\]
where the sum extends only over \(2p\) branch points which are fundamental points neither of the first nor of the second kind. The \(2p\) quantities \((b_i)\) are now no longer subject to any linear relation, since the conditional equations were used to eliminate the remaining \(n-1\) quantities \((b_i)\). The quantities \([A_i]\) must be examined and characterized more precisely.

\subsection*{\S\,27. Form of the relation found and properties of its integral coefficients.}

From the \(n-1\) periods \((A_i)\) corresponding to the \(n-1\) fundamental points of the second kind, form combinations \(h_2,h_3\ldots h_n\) such that the paths correspondingly composed from the \(A_i\) lead from the first circuit respectively into the second, third, and so on. If the \(i^{\text{th}}\) branch point belongs to the \(\alpha^{\text{th}}\) circuit with positively traversed loop and to the \(\beta^{\text{th}}\) with negatively traversed loop, then
\[
{}[A_i]=(A_i)+h_{\alpha}-h_{\beta},
\]
where \(h_1\) is always to be replaced by zero. Indeed, if from this one substitutes
\[
(A_i)=[A_i]-h_{\alpha}+h_{\beta}
\]
into the right hand side of equation (4), the following happens. The \((A_i)\)'s belonging to the fundamental points of the first kind do not occur in that expression at all. The \([A_i]\)'s corresponding to the fundamental points of the second kind vanish identically, because the combination \(h_{\alpha}-h_{\beta}\) then necessarily reduces to \(-(A_i)\), since fundamental points of the second kind can never generate a cycle of the second kind. Finally, everything multiplied by \(h_{\alpha},h_{\beta}\) also vanishes. For \(h_{\alpha}\) is multiplied by all those \((b_i)\)'s whose branch points lie on the \(\alpha^{\text{th}}\) circuit, negatively when the loop of the point occurs positively on this circuit, positively when the opposite occurs. The coefficient of \(h_{\alpha}\) is thus exactly one of the expressions which vanish because of the linear conditional equations. One therefore comes back to the form (5), and the new definition of the \([A_i]\)'s must agree with the one given there, since no relations remain between the surviving \((b_i)\)'s.

The quantities \([A_i]\) are periodicity moduli of such a kind that the branch points used in their formation constitute not only cycles of the first kind, but also cycles of the second kind, that is, cycles formed with regard to the loops of the second kind, or cycles which in continuous succession always connect one circuit of the first kind with another and finally return to the circuit used at the beginning, just as the cycles of the first kind lead continuously from one root to another and finally back to the root used at the beginning.

In their formation the quantities \(A_i\) appeared as linear functions of the \(a_i\)'s with integral coefficients. Since the \((a_i)\)'s corresponding to the fundamental points of the second kind can be expressed in the same way by the remaining \(2p\), one finally has the \(A_i\)'s, and hence also the \([A_i]\)'s composed from them, as linear functions of the latter \(2p\) quantities \((a_i)\). Thus put
\[
{}[A_i]=c_{i1}a_1+c_{i2}a_2\cdots+c_{i,2p}a_{2p},
\]
where, for convenience, we now use the indices \(1,2\ldots2p\) for the remaining branch points; then the condition existing between the \((a)\)'s and the \((b)\)'s is
\[
(6)\quad \sum_{i,k=1}^{2p} c_{ik}b_i a_k=0.
\]
We shall now prove two main properties of the integral coefficients \(c_{ik}\) of this equation, which is fundamental for what follows.

Since the equation
\[
\Sigma A_i b_i+\Sigma B_i a_i=0
\]
became an identity with the help of the equations existing between the \(a\)'s alone and between the \(b\)'s alone, the equation arising from applying the above transformation to both parts,
\[
\Sigma[A_i]b_i+\Sigma[B_k]a_k=0,
\]
in which no such relations between the \((a_i)\)'s and \((b_i)\)'s hold any longer, must be a complete identity. But one also has
\[
[B_k]=c_{k1}b_1+c_{k2}b_2\cdots+c_{k,2p}b_{2p},
\]
and the above identity therefore becomes
\[
\sum c_{ik}b_i a_k+\sum c_{ki}b_i a_k=0,
\]
which resolves into the \(\dfrac{2p\cdot(2p+1)}{2}\) equations
\[
c_{ik}+c_{ki}=0.
\]
All products \((a_i)(b_i)\) therefore disappear from expression (6), and the coefficients of \((a_i)(b_k)\) and \((b_i)(a_k)\) are equal and opposite. \emph{The coefficients \(c_{ik}\) therefore form an alternating determinant.}

It follows already from this that the determinant of the \(c\)'s must be the square of an integer, since every alternating determinant of even order is the square of an expression rationally formed from its elements. But the following important theorem can further be proved:

\begin{center}
\emph{The determinant of the \(c\)'s is equal to 1.}
\end{center}

This theorem, which is in a certain sense the pivot of this theory, is proved as follows.

The \([A_i]\)'s can be represented as linear functions of the remaining \(2p\) \((a_i)\)'s with integral coefficients. Let the value of the integral determinant of the \(c\)'s be \(C\). If, analogously, the \((a_i)\)'s can also be represented as integral linear functions of the \([A_i]\)'s, and if the determinant of the latter expressions is \(\Gamma\), then \(C,\Gamma\) are adjugate determinants, so that \(C\cdot\Gamma=1\), hence \(C=\Gamma=\pm1\). And since \(C\) is necessarily positive, \(C=1\), as was to be proved.

In order now to express the \(2p\) quantities \((a_i)\), which do not correspond to fundamental points, by the \(2p\) corresponding \([A]\)'s, we consider the following sequence of equations:

\begin{center}
\scriptsize
\begin{tabular}{p{0.43\linewidth}|p{0.43\linewidth}}
\centering I. & \centering II.\tabularnewline
The \(A\)'s are expressed by the \(a\)'s by forming circuits of the first kind which each exclude only one branch point. &
The \(a\)'s are expressed by the \(A\)'s by forming circuits of the second kind which each exclude only one branch point.\tabularnewline[0.6em]
\centering III. & \centering IV.\tabularnewline
By describing complete circuits of the first kind one forms equations between the \(a\)'s. &
By describing complete circuits of the second kind one forms equations between the \(A\)'s.\tabularnewline[0.6em]
\centering V. & \centering VI.\tabularnewline
From the \(a\)'s one forms the \((a)\)'s by adding \(a\)'s which correspond to the fundamental points of the first kind. In doing so, the \(a\)'s belonging to the latter vanish. &
From the \(A\)'s one forms the \((A)\)'s by taking the \(A\)'s from I and in these expressions completing the \(a\)'s to \((a)\)'s. In general none of the \(A\)'s vanishes.\tabularnewline[0.6em]
\centering VII. & \centering VIII.\tabularnewline
Equations I therefore continue to hold if the \(a,A\) occurring in them are replaced by \((a),A\). &
Equations II also continue to hold if \(a,A\) are replaced by \((a),(A)\). For if a trace remained from the \(a\)'s corresponding to fundamental points of the first kind, then integrals over non-cyclic paths would in general have to reduce to cyclic ones, which is impossible.\tabularnewline[0.6em]
\centering IX. & \centering X.\tabularnewline
\multicolumn{2}{c}{The same holds for the conditions III. IV.}\tabularnewline[0.6em]
\centering XI. & \centering XII.\tabularnewline
From the \((A)\)'s one forms the \([A]\)'s by adding \(A\)'s which correspond to the fundamental points of the second kind. &
From the \((a)\)'s one forms the \([a]\)'s by defining the latter in such a way that equations VIII are to continue to hold also for \([a],[A]\).
\end{tabular}
\end{center}

The quantity \([a_i]\) is therefore expressed by the \((a_i)\)'s, augmented by those combinations \(h_{\alpha},h_{\beta}\) of the \((A)\)'s corresponding to fundamental points of the second kind which, formed from the \(A\)'s, would close the incomplete circuit of the \(A\)'s described in forming the \(a\)'s (with \(h_1\) replaced by zero). Thus \(\alpha,\beta\) are the indices of the complete circuits of the second kind in which the \(i^{\text{th}}\) branch point is essential, or, because of the reciprocity existing between the two systems of loops, the indices of those roots which the \(i^{\text{th}}\) loop of the first kind connects with one another. But some of the \((a_i)\)'s are zero, namely those corresponding to the fundamental points of the first kind. Hence their associated \([a_i]\)'s are simply equal to the differences of the corresponding \(h\)'s. And since the corresponding \(n-1\) quantities \(a_i\), as belonging to fundamental points of the first kind, can be combined so that they lead from \(s_1\) to any other \(s\), the \(n-1\) \([a_i]\)'s can be combined in the same way, and then one of the quantities \(h_{\alpha}\) is obtained directly expressed by these \([a_i]\)'s. Thus the quantities \((a_i)\) can now also be expressed by the \([a_i]\)'s, that is, linearly and integrally by the \([A_i]\)'s, whose number is \(2p+n-1\), since the \([A_i]\)'s corresponding to the fundamental points of the second kind vanish.

But equations IV. X plainly still remain valid when one replaces \(A,(A)\) by \([A]\). For in passing from \((A)\) to \([A]\), one only ever has to add the differences of the \(k\)'s which, formed from the \(A\)'s and added to the \(A\)'s, would produce a simple cycle of the second kind. The indices of two \(k\)'s occurring in such a difference indicate the roots which a loop of the second kind connects. Since each following difference begins with the \(k\) with which the preceding one ended, and the last closes with the \(k\) with which the first began, all \(k\)'s really cancel, and only equations IV. X remain, written with the \([A]\)'s.

But just as we proved earlier that from equations I the \(a\)'s corresponding to the first fundamental points can be expressed integrally and linearly by the others, so reciprocally the theorem holds that from equations II the \(A\)'s corresponding to the second fundamental points can be expressed integrally and linearly by the others. If therefore in those equations we put \([A]\) in place of \(A\), we obtain the \(n-1\) \([A]\)'s corresponding to fundamental points of the second kind expressed by the remaining \(2p\). Consequently the \((a_i)\)'s which above were expressed by all the \([A_i]\)'s can also be expressed integrally and linearly by the latter \(2p\), as was to be proved.

\subsection*{\S\,28. Example for setting up the relation between the periodicity moduli of the integrals of the first kind.}

For explanation, in the case \(n=4\) and for 12 branch points, we insert the following schema, in which the operations indicated above are carried out.



% Source PDF pp. 122--133 / printed pp. 100--111.
% Continuation of § 28 after the introductory sentence already present in the cumulative text.

\begin{center}
\scriptsize
\begin{tabular}{c@{\qquad}c@{\qquad}c}
\textbf{Index of}\\[-0.25em]
\textbf{the branch}\\[-0.25em]
\textbf{points.}
&
\begin{tabular}{c}
\textbf{Their loops of the \(1^{\mathrm{st}}\) kind}\\
\textbf{connect the roots, or they lie}\\
\textbf{on the circuits}\\
\textbf{of the \(2^{\mathrm{nd}}\) kind.}
\end{tabular}
&
\begin{tabular}{c}
\textbf{Their loops of the \(2^{\mathrm{nd}}\) kind}\\
\textbf{connect the roots, or}\\
\textbf{lie on the circuits}\\
\textbf{of the \(1^{\mathrm{st}}\) kind.}
\end{tabular}
\\[0.6em]\hline\\[-0.6em]
1. & \(1-2\) & \(1-2\)\\
2. & \(2-3\) & \(1-3\)\\
3. & \(1-4\) & \(2-4\)\\
4. & \(2-4\) & \(3-2\)\\
5. & \(3-4\) & \(1-3\)\\
6. & \(1-3\) & \(4-3\)\\
7. & \(2-4\) & \(2-1\)\\
8. & \(1-2\) & \(3-1\)\\
9. & \(2-3\) & \(3-4\)\\
10. & \(2-4\) & \(4-2\)\\
11. & \(3-4\) & \(3-4\)\\
12. & \(3-4\) & \(4-3\)
\end{tabular}
\end{center}

\noindent Relations between the \(a\)'s from circuits of the \(1^{\mathrm{st}}\) kind:
\[
\mathrm{I.}\quad
\begin{array}{r@{\quad}l}
1&0=a_1+a_2+a_3-a_7-a_8,\\
2&0=-a_1+a_3-a_4+a_7-a_{10},\\
3&0=-a_2+a_4-a_5-a_6+a_8+a_9+a_{11}-a_{12},\\
4&0=-a_3+a_5-a_9+a_{10}-a_{11}+a_{12}.
\end{array}
\]

\noindent Relations between the \(A\)'s from circuits of the \(2^{\mathrm{nd}}\) kind:
\[
\mathrm{II.}\quad
\begin{array}{r@{\quad}l}
1&0=A_1+A_3+A_6+A_8,\\
2&0=-A_1+A_2+A_4+A_7-A_8+A_9+A_{10},\\
3&0=-A_2+A_5-A_6-A_9+A_{11}+A_{12},\\
4&0=-A_3-A_4-A_5-A_7-A_{10}-A_{11}-A_{12}.
\end{array}
\]

\noindent Expressions of the \(A\)'s by the \(a\)'s:
\[
\mathrm{III.}\quad
\begin{array}{rcl}
A_1&=&a_3-a_4+a_7-a_{10}\\
A_2&=&a_1+a_4-a_5-a_6+a_8+a_9+a_{11}-a_{12}\\
A_3&=&-a_1+a_6-a_9+a_{10}-a_{11}+a_{12}\\
A_4&=&-a_2+a_7-a_{10}\\
A_5&=&a_1+a_2-a_6+a_8+a_9+a_{11}-a_{12}\\
A_6&=&-a_3+a_8+a_9+a_{11}-a_{12}\\
A_7&=&-a_1+a_3-a_4-a_5\\
A_8&=&-a_2+a_4-a_5-a_6\\
A_9&=&-a_2+a_4-a_5-a_6+a_8+a_{10}-a_{11}+a_{12}\\
A_{10}&=&-a_3+a_6-a_9\\
A_{11}&=&-a_2+a_4-a_5-a_6+a_8+a_9+a_{12}\\
A_{12}&=&-a_3+a_6-a_9+a_{10}-a_{11}.
\end{array}
\]

\noindent Expressions of the \(a\)'s by the \(A\)'s:
\[
\mathrm{IV.}\quad
\begin{array}{rcl}
a_1&=&A_2+A_4+A_7-A_8+A_9+A_{10}\\
a_2&=&-A_1+A_5-A_6-A_9+A_{11}+A_{12}\\
a_3&=&A_1-A_4-A_5-A_7-A_{10}-A_{11}-A_{12}\\
a_4&=&-A_1+A_2-A_5-A_7-A_{10}-A_{11}-A_{12}\\
a_5&=&-A_2-A_7-A_{10}-A_{11}-A_{12}\\
a_6&=&A_1+A_3-A_9+A_{11}+A_{12}\\
a_7&=&-A_1+A_2+A_4-A_{10}-A_{11}-A_{12}\\
a_8&=&A_1+A_3+A_6+A_9+A_{10}\\
a_9&=&-A_1+A_2+A_4+A_7-A_8+A_{11}+A_{12}\\
a_{10}&=&-A_1+A_2+A_4+A_7-A_8+A_9-A_{11}-A_{12}\\
a_{11}&=&-A_2+A_5-A_6-A_9-A_{12}\\
a_{12}&=&-A_2+A_5-A_6-A_9+A_{11}.
\end{array}
\]

\noindent Fundamental points of the \(1^{\mathrm{st}}\) kind:\quad \(1.\quad2.\quad3.\)
\hfill Fundamental points of the \(2^{\mathrm{nd}}\) kind:\quad \(4.\quad5.\quad6.\)

\noindent Combinations which in the first system lead from the root \(1\) to \(2,3,4\):
\[
g_2=a_1,\qquad g_3=a_1+a_2,\qquad g_4=a_3.
\]
Combinations which in the second system lead from the root \(1\) to \(2,3,4\):
\[
k_2=A_4+A_5,\qquad k_3=A_5,\qquad k_4=A_5-A_6.
\]

\noindent Formation of the round-bracketed quantities:
\[
\begin{array}{rcl@{\qquad\qquad}rcl}
(a_1)&=&0&(A_1)&=&A_1-a_1\\
(a_2)&=&0&(A_2)&=&A_2-a_1-a_2\\
(a_3)&=&0&(A_3)&=&A_3-a_3+a_1\\
(a_4)&=&a_4-a_3+a_1&(A_4)&=&A_4+a_2\\
(a_5)&=&a_5-a_3+a_1+a_2&(A_5)&=&A_5-a_1-a_2\\
(a_6)&=&a_6-a_1-a_2&(A_6)&=&A_6+a_3-a_1-a_2\\
(a_7)&=&a_7-a_3+a_1&(A_7)&=&A_7+a_1\\
(a_8)&=&a_8-a_1&(A_8)&=&A_8+a_1+a_2\\
(a_9)&=&a_9-a_2&(A_9)&=&A_9-a_3+a_1+a_2\\
(a_{10})&=&a_{10}-a_3+a_1&(A_{10})&=&A_{10}+a_3-a_1\\
(a_{11})&=&a_{11}-a_3+a_1+a_2&(A_{11})&=&A_{11}-a_3+a_1+a_2\\
(a_{12})&=&a_{12}-a_3+a_1+a_2&(A_{12})&=&A_{12}+a_3-a_1-a_2.
\end{array}
\]

\noindent Formation of the square-bracketed quantities:
\[
\begin{array}{rcl@{\qquad\qquad}rcl}
{}[A_1]&=&(A_1)-(A_4)-(A_5)&[a_1]&=&-(A_4)-(A_5)\\
{}[A_2]&=&(A_2)-(A_5)&[a_2]&=&(A_4)\\
{}[A_3]&=&(A_3)+(A_4)+(A_6)&[a_3]&=&-(A_5)+(A_6)\\
{}[A_4]&=&0&[a_4]&=&(a_4)+(A_4)+(A_6)\\
{}[A_5]&=&0&[a_5]&=&(a_5)+(A_6)\\
{}[A_6]&=&0&[a_6]&=&(a_6)-(A_5)\\
{}[A_7]&=&(A_7)+(A_4)+(A_5)&[a_7]&=&(a_7)+(A_4)+(A_6)\\
{}[A_8]&=&(A_8)+(A_5)&[a_8]&=&(a_8)-(A_4)-(A_5)\\
{}[A_9]&=&(A_9)+(A_6)&[a_9]&=&(a_9)+(A_4)\\
{}[A_{10}]&=&(A_{10})-(A_4)-(A_6)&[a_{10}]&=&(a_{10})+(A_4)+(A_6)\\
{}[A_{11}]&=&(A_{11})+(A_6)&[a_{11}]&=&(a_{11})+(A_6)\\
{}[A_{12}]&=&(A_{12})-(A_6)&[a_{12}]&=&(a_{12})+(A_6).
\end{array}
\]

\noindent Expressions of the \([A]\)'s by the \((a)\)'s and conversely:
\[
\begin{array}{rcl}
{}[A_7]&=&-(a_4)-(a_6)+(a_7)+(a_9)-(a_{10})+(a_{11})-(a_{12})\\
{}[A_8]&=&(a_4)-(a_5)-2(a_6)+(a_8)+(a_9)+(a_{11})-(a_{12})\\
{}[A_9]&=&(a_4)-(a_5)-(a_6)+2(a_8)+(a_9)+(a_{10})\\
{}[A_{10}]&=&(a_6)-(a_7)-(a_8)-2(a_9)+(a_{10})-(a_{11})+(a_{12})\\
{}[A_{11}]&=&(a_4)-(a_5)-(a_6)+2(a_8)+2(a_9)+(a_{11})\\
{}[A_{12}]&=&(a_6)-(a_8)-2(a_9)+(a_{10})-2(a_{11})+(a_{12})\\[0.3em]
(a_7)&=&-2[A_1]+2[A_2]+2[A_7]-[A_8]+[A_9]+[A_{10}]\\
(a_8)&=&[A_1]-[A_2]+[A_3]-[A_7]+[A_8]\\
(a_9)&=&[A_2]+[A_7]-[A_8]+[A_9]\\
(a_{10})&=&-2[A_1]+2[A_2]+3[A_7]-2[A_8]+2[A_9]+2[A_{10}]\\
(a_{11})&=&-2[A_1]+2[A_7]-[A_8]-[A_9]+2[A_{10}]+2[A_{11}]+[A_{12}]\\
(a_{12})&=&-2[A_1]+2[A_7]-[A_8]-[A_9]+2[A_{10}]+3[A_{11}]+2[A_{12}].
\end{array}
\]

\noindent Expressions of \((a_4),(a_5),(a_6),[A_1],[A_2],[A_3]\) from relations I. II.:
\[
\begin{array}{rcl}
(a_4)&=&(a_7)-(a_{10})\\
(a_5)&=&(a_7)+(a_8)\\
(a_6)&=&(a_9)-(a_{10})+(a_{11})-(a_{12})\\
{}[A_1]&=&[A_7]-[A_8]+[A_{10}]+[A_{11}]+[A_{12}]\\
{}[A_2]&=&-[A_9]+[A_{11}]+[A_{12}]\\
{}[A_3]&=&-[A_7]-[A_{10}]-[A_{11}]-[A_{12}].
\end{array}
\]

\noindent Definitive expressions of the quantities \((a_i),[A_i]\), which correspond to no fundamental points:
\[
\begin{array}{rcl@{\qquad}rcl}
{}[A_7]&=&(a_{10})&(a_7)&=&[A_8]-[A_9]-[A_{10}]\\
{}[A_8]&=&-(a_9)+(a_{10})-(a_{11})+(a_{12})&(a_8)&=&-[A_7]+[A_9]-[A_{11}]-[A_{12}]\\
{}[A_9]&=&(a_8)+(a_{10})-(a_{11})+(a_{12})&(a_9)&=&[A_7]-[A_8]+[A_{11}]+[A_{12}]\\
{}[A_{10}]&=&-(a_7)-(a_8)-(a_9)&(a_{10})&=&[A_7]\\
{}[A_{11}]&=&(a_8)+(a_9)+(a_{12})&(a_{11})&=&[A_8]-[A_9]-[A_{12}]\\
{}[A_{12}]&=&-(a_8)-(a_9)-(a_{11})&(a_{12})&=&[A_8]-[A_9]+[A_{11}].
\end{array}
\]
The determinant of the \(c_{ik}\), with rows/columns suppressed, is therefore
\[
\left|\begin{array}{rrrrrr}
 0& 0& 0& 1& 0& 0\\
 0& 0&-1& 1&-1& 1\\
 0& 1& 0& 1&-1& 1\\
-1&-1&-1& 0& 0& 0\\
 0& 1& 1& 0& 0& 1\\
 0&-1&-1& 0&-1& 0
\end{array}\right|.
\]

\subsection*{\S\,29. Transformation of the relation.}

The quantities \((a_i)\) which correspond to no fundamental points may be designated as fundamental periodicity moduli; by them all other periodicity moduli of the same integral can be expressed linearly and integrally. Between the fundamental periodicity moduli of any two different integrals of the first kind there is the equation
\[
\tag{1}\sum c_{ik}(a_i)(b_k)=0,
\]
where the \(c\)'s each time denote the same numbers.

Instead of the periodicity moduli \((a_i)\), however, one may use, in order to express all possible periodicity moduli linearly and integrally, any \(2p\) integral linear combinations of them whose determinant is \(1\). Thus we may introduce, in place of the \((a_i)\), the \(2p\) new periodicity moduli
\[
\tag{2}M^{(h)}=\alpha_1^{(h)}(a_1)+\alpha_2^{(h)}(a_2)+\cdots+\alpha_{2p}^{(h)}(a_{2p}),\quad(h=1,2..2p)
\]
and therefore at the same time in place of the \((b)\)'s the \(2p\) new periodicity moduli
\[
\tag{3}N^{(h)}=\alpha_1^{(h)}(b_1)+\alpha_2^{(h)}(b_2)+\cdots+\alpha_{2p}^{(h)}(b_{2p}),\quad(h=1,2..2p).
\]
As soon as the determinant of the integers \(\alpha\) is equal to \(1\), one may express all periodicity moduli of the first integral by the \(M\)'s, and all those of the second by the \(N\)'s. If the solutions of the systems (2), (3) are
\[
\tag{4}
\left\{\begin{array}{rcl}
(a_i)&=&\beta_i^{(1)}M^{(1)}+\beta_i^{(2)}M^{(2)}+\cdots+\beta_i^{(2p)}M^{(2p)}\\
(b_i)&=&\beta_i^{(1)}N^{(1)}+\beta_i^{(2)}N^{(2)}+\cdots+\beta_i^{(2p)}N^{(2p)},
\end{array}\right.
\]
then the \(\beta\)'s are again integers whose determinant is \(1\), and equation (1) is transformed into
\[
\tag{5}\sum C^{(h,m)}M^{(h)}N^{(m)}=0,
\]
where
\[
\tag{6}C^{(h,m)}=\sum c_{ik}\beta_i^{(h)}\beta_k^{(m)}.
\]
From this one sees, first, that always again
\[
C^{(h,m)}=-C^{(m,h)}.
\]
For if \(h\) is interchanged with \(m\), and at the same time, which is irrelevant, \(i\) with \(k\), then the right side of (6) merely changes sign. Secondly, by known determinant theorems the determinant of \(C\) is equal to the determinant of \(c\) multiplied by the square of the determinant of \(\beta\), hence is equal to \(1\). Therefore the determinant of \(C\) has exactly the properties belonging to the determinant of \(c\).

It follows further that these properties of the coefficients continue to be preserved through a sequence of linear substitutions in which, successively, other fundamental periodicity moduli occur in place of the \((a),(b)\).

Among the linear substitutions there are now such as have on the determinant of the \(c\)'s the same effect as if in it one vertical row, multiplied by a factor, were added to another, and at the same time the same were done with the corresponding horizontal rows. Namely, let
\[
(a_i)=M^{(i)},
\]
and, deviating from this only for a certain value \(i=\rho\),
\[
(a_\rho)=M^{(\rho)}+mM^{(\sigma)}\quad(\sigma\ \text{different from}\ \rho).
\]
The determinant of the \(\beta\)'s then contains in the diagonal only \(1\)'s, otherwise zeros, except in a single place in which the element \(m\) occurs. Therefore the determinant of the \(\beta\)'s really has the value \(1\). The left side of equation (1) passes into
\[
\sum c_{ik}M^{(i)}N^{(k)}+
m\left[M^{(\sigma)}\sum_k c_{\rho k}N^{(k)}+
N^{(\sigma)}\sum_i c_{i\rho}M^{(i)}\right].
\]
Thus, in general,
\[
C^{(ik)}=c_{ik},
\]
but in particular
\[
C^{(i\sigma)}=c_{i\sigma}+mc_{i\rho},\qquad
C^{(\sigma i)}=c_{\sigma i}+mc_{\rho i},
\]
and
\[
C^{(\sigma\sigma)}=c_{\sigma\sigma}+m(c_{\sigma\rho}+c_{\rho\sigma})=0.
\]
Thus the determinant of \(C\) arises from the determinant of \(c\) when the \(\rho^{\mathrm{th}}\) vertical row, multiplied by \(m\), is added to the \(\sigma^{\mathrm{th}}\), and the same operation is carried out with the horizontal rows. This operation is equivalent to the linear transformation
\[
\begin{array}{rcl@{\qquad}rcl}
(a_\rho)&=&M^{(\rho)}+mM^{(\sigma)},&(a_i)&=&M^{(i)}\\
M^{(\rho)}&=&(a_\rho)-m(a_\sigma),&M^{(i)}&=&(a_i)
\end{array}
\quad(i=1,2..\rho-1,\rho+1..2p).
\]

We now see how we can bring the determinant of the \(c\)'s, by a sequence of such operations, into as simple a form as possible.

Since the determinant has the value \(1\), the elements of a row, for example of the first horizontal row, have no common factor. Hence, by suitable addition of vertical rows, we can make one number of the first horizontal row equal to \(1\); and then at once we can destroy the other terms of the same horizontal and vertical row with the help of this term. We can also bring the \(1\) to any desired place of the first vertical row and then destroy the \(1\) at its former place. It can therefore always be arranged that the second place of the first horizontal row is \(1\), all the others of the same horizontal row and of the second vertical row consisting of zeros; and since we always perform the corresponding operations with the vertical rows, the first vertical row will then contain \(-1\) in the second place and \(0\) in all the others, and the second horizontal row will contain \(-1\) in the first place and \(0\) in all the others. One may now proceed further with the third and fourth rows exactly as just with the first and second, then likewise with the fifth and sixth, and so on, so that finally the determinant of the \(c\)'s assumes the form
\[
\left|\begin{array}{rrrrrrr}
0&1&0&0&\cdots&0&0\\
-1&0&0&0&\cdots&0&0\\
0&0&0&1&\cdots&0&0\\
0&0&-1&0&\cdots&0&0\\
\multicolumn{7}{c}{\cdots\qquad\qquad\cdots}\\
0&0&0&0&\cdots&0&1\\
0&0&0&0&\cdots&-1&0
\end{array}\right|.
\]
Thus, finally, by successive transformation such a transformation has arisen for which in the transformed equation (1) all \(c_{ik}\) vanish except
\[
c_{12},\ c_{34},\ c_{56}\ldots c_{2p-1,2p},
\]
which are equal to \(1\). Or, in other words: one can always set up integral linear transformation formulae
\[
\begin{array}{rcl}
M^{(h)}&=&\alpha_1^{(h)}(a_1)+\alpha_2^{(h)}(a_2)+\cdots+\alpha_{2p}^{(h)}(a_{2p})\\
N^{(h)}&=&\alpha_1^{(h)}(b_1)+\alpha_2^{(h)}(b_2)+\cdots+\alpha_{2p}^{(h)}(b_{2p})
\end{array}
\]
by means of which identically
\[
\begin{split}
\sum c_{ik}(a_i)(b_k)
&=(M^{(1)}N^{(2)}-N^{(1)}M^{(2)})+(M^{(3)}N^{(4)}-N^{(3)}M^{(4)})\\
&\quad+\cdots+(M^{(2p-1)}N^{(2p)}-N^{(2p-1)}M^{(2p)})
\end{split}
\]
is obtained.

These combinations of the \((a),(b)\), to which the transformation of equation (1) has here led, will be called normal periodicity moduli; there are infinitely many systems of normal periodicity moduli. In such a system the normal periodicity moduli of an integral are pairwise conjugate, inasmuch as the quantities bearing the indices \(1\) and \(2\), \(3\) and \(4,\ldots,2p-1\) and \(2p\) stand in special relations to one another; and between the normal periodicity moduli of two integrals of the first kind there is the relation
\[
\tag{7}\sum_{i=1}^{i=p}(M^{(2i-1)}N^{(2i)}-N^{(2i-1)}M^{(2i)})=0.\footnotemark
\]
\footnotetext{In this form Rosenhain and Weierstrass gave the relations occurring between the periodicity moduli for the hyperelliptic functions, and Riemann for the Abelian functions.}

\subsection*{\S\,30. Normal form of the integrals of the first kind.}

These considerations finally lead to a definite normal form of the integrals of the first kind.

Hitherto, by the \(p\) integrals of the first kind we understood expressions which were linearly composed from any \(p\) integrals of the first kind. Each such integral therefore still contains \(p\) arbitrary constants. If \(H_1,H_2\ldots H_p\) are any mutually independent integrals of the first kind, that is, such that no relation with constant coefficients exists between them, then the most general \(p\) integrals of the first kind which we can take as a basis are
\[
\begin{array}{rcl}
I_1&=&m_{11}H_1+m_{12}H_2+\cdots+m_{1p}H_p\\
I_2&=&m_{21}H_1+m_{22}H_2+\cdots+m_{2p}H_p\\
&\vdots&\\
I_p&=&m_{p1}H_1+m_{p2}H_2+\cdots+m_{pp}H_p,
\end{array}
\]
where the determinant of the arbitrary constants \(m\), so that the \(I\)'s too are independent of one another, must not vanish.

If for these \(I\)'s we form a system of normal periodicity moduli, then the \(H\)'s are transformed into integrals between fixed limits, hence into given constants. We can therefore determine the \(m\)'s so that \(p\) of the normal periodicity moduli of each integral assume definite values. If we denote the normal periodicity moduli of \(I_k\) by \(I_k^{(1)}, I_k^{(2)}\ldots I_k^{(2p)}\), and the corresponding ones of \(H_k\) by \(H_k^{(1)}, H_k^{(2)}\ldots H_k^{(2p)}\), then we choose the \(m_{ik}\) so that all \(I_k^{(h)}\) with an odd upper index vanish, except \(I_k^{(2k-1)}\), which are to assume the common value \(2i\pi\). For determining \(m_{i1},m_{i2}\ldots m_{ip}\) we then have the following system of linear equations:
\[
\begin{array}{rcl}
m_{i1}H_1^{(1)}+m_{i2}H_2^{(1)}+\cdots+m_{ip}H_p^{(1)}&=&0\\
m_{i1}H_1^{(3)}+m_{i2}H_2^{(3)}+\cdots+m_{ip}H_p^{(3)}&=&0\\
\multicolumn{3}{c}{\cdots}\\
m_{i1}H_1^{(2i-1)}+m_{i2}H_2^{(2i-1)}+\cdots+m_{ip}H_p^{(2i-1)}&=&2i\pi\\
\multicolumn{3}{c}{\cdots}\\
m_{i1}H_1^{(2p-1)}+m_{i2}H_2^{(2p-1)}+\cdots+m_{ip}H_p^{(2p-1)}&=&0,
\end{array}
\]
and therefore, if \(\Delta\) denotes the determinant of the \(H\)'s with odd upper index and \(\Delta_k^{(2m-1)}\) one of its minors,
\[
m_{ik}=\frac{2\pi\sqrt{-1}}{\Delta}\Delta_k^{(2i-1)}.
\]
The normal integrals of the first kind which arise in this way from \(I_1,I_2\ldots I_p\) will henceforth be denoted by \(u_1,u_2\ldots u_p\). The first half of the system of normal periodicity moduli of the normal integrals is given by the schema
\[
\begin{array}{c@{\quad}cccc}
u_1&2i\pi&0&\cdots&0\\
u_2&0&2i\pi&\cdots&0\\
\vdots&&\ddots&&\vdots\\
u_p&0&0&\cdots&2i\pi.
\end{array}
\]
Returning now to equation (7), in order to investigate what relations occur between the remaining normal periodicity moduli, we can first write that equation, applied to two integrals \(I_h,I_k\), in the form
\[
\sum_{i=1}^{i=p}\left(I_h^{(2i-1)}I_k^{(2i)}-I_k^{(2i-1)}I_h^{(2i)}\right)=0.
\]
But since \(I_h^{(2i-1)}\) and \(I_k^{(2i-1)}\) vanish except for \(I_h^{(2h-1)}\) and \(I_k^{(2k-1)}\), which become equal to \(2i\pi\), this equation passes into the simple form
\[
I_k^{(2h)}=I_h^{(2k)}.
\]
Thus, if we put
\[
I_k^{(2h)}=I_h^{(2k)}=a_{hk},
\]
where \(a_{hk}\) is not to change when the indices are interchanged, then the equations flowing from (7) are satisfied, and one has the theorem:

The second half of the normal periodicity moduli of normal integrals of the first kind forms a system of symmetric determinant:
\[
\begin{array}{c@{\quad}cccc}
u_1&a_{11}&a_{12}&\cdots&a_{1p}\\
u_2&a_{21}&a_{22}&\cdots&a_{2p}\\
\vdots&&\cdots&&\\
u_p&a_{p1}&a_{p2}&\cdots&a_{pp}.
\end{array}
\]
And the most general values which the normal integrals \(u_1,u_2\ldots u_p\) can assume through a change of the path of integration are
\[
\begin{array}{rcl}
u_1'&=&u_1+2i\pi m_1+a_{11}q_1+a_{12}q_2+\cdots+a_{1p}q_p\\
u_2'&=&u_2+2i\pi m_2+a_{21}q_1+a_{22}q_2+\cdots+a_{2p}q_p\\
&\vdots&\\
u_p'&=&u_p+2i\pi m_p+a_{p1}q_1+a_{p2}q_2+\cdots+a_{pp}q_p,
\end{array}
\]
where \(m,q\) denote arbitrary positive or negative integers (cf. Riemann).

\subsection*{\S\,31. Execution of the transformation for the example treated in \S\,28.}

We will carry out, for the example treated above, the search for the normal periodicity moduli. Our determinant was:
\[
\left|\begin{array}{rrrrrr}
0&0&0&1&0&0\\
0&0&-1&1&-1&1\\
0&1&0&1&-1&1\\
-1&-1&-1&0&0&0\\
0&1&1&0&0&1\\
0&-1&-1&0&-1&0
\end{array}\right|.
\]
Subtracting the third horizontal, respectively vertical, row from the second gives:
\[
\mathrm{I.}\quad
\left|\begin{array}{rrrrrr}
0&0&0&1&0&0\\
0&0&-1&0&0&0\\
0&1&0&1&-1&1\\
-1&0&-1&0&0&0\\
0&0&1&0&0&1\\
0&0&-1&0&-1&0
\end{array}\right|.
\]
Subtracting the first row from the third:
\[
\mathrm{II.}\quad
\left|\begin{array}{rrrrrr}
0&0&0&1&0&0\\
0&0&-1&0&0&0\\
0&1&0&0&-1&1\\
-1&0&0&0&0&0\\
0&0&1&0&0&1\\
0&0&-1&0&-1&0
\end{array}\right|.
\]
Adding the second row to the fifth:
\[
\mathrm{III.}\quad
\left|\begin{array}{rrrrrr}
0&0&0&1&0&0\\
0&0&-1&0&0&0\\
0&1&0&0&0&1\\
-1&0&0&0&0&0\\
0&0&0&0&0&1\\
0&0&-1&0&-1&0
\end{array}\right|.
\]
Subtracting the second row from the sixth:
\[
\mathrm{IV.}\quad
\left|\begin{array}{rrrrrr}
0&0&0&1&0&0\\
0&0&-1&0&0&0\\
0&1&0&0&0&0\\
-1&0&0&0&0&0\\
0&0&0&0&0&1\\
0&0&0&0&-1&0
\end{array}\right|.
\]
This determinant has the desired form, apart from the order of the rows, which ultimately is immaterial.

The four transformations of the determinant give the following substitutions:
\[
\begin{array}{c|c|c|c}
\mathrm{I.}&\mathrm{II.}&\mathrm{III.}&\mathrm{IV.}\\
(a_7)=a_7'&a_7'=a_7''-a_9''&a_7''=a_7'''&a_7'''=M_7\\
(a_8)=a_8'&a_8'=a_8''&a_8''=a_8'''+a_{11}'''&a_8'''=M_8-M_{12}\\
(a_9)=a_9'-a_8'&a_9'=a_9''&a_9''=a_9'''&a_9'''=M_9\\
(a_{10})=a_{10}'&a_{10}'=a_{10}''&a_{10}''=a_{10}'''&a_{10}'''=M_{10}\\
(a_{11})=a_{11}'&a_{11}'=a_{11}''&a_{11}''=a_{11}'''&a_{11}'''=M_{11}\\
(a_{12})=a_{12}'&a_{12}'=a_{12}''&a_{12}''=a_{12}'''&a_{12}'''=M_{12}
\end{array}
\]
Between the first and last variables there therefore exist the equations
\[
\begin{array}{rcl@{\qquad}rcl}
(a_7)&=&M_7-M_9&M_7&=&(a_7)+(a_9)+(a_8)\\
(a_8)&=&M_8+M_{11}-M_{12}&M_8&=&(a_8)-(a_{11})+(a_{12})\\
(a_9)&=&-M_8+M_9-M_{11}+M_{12}&M_9&=&(a_8)+(a_9)\\
(a_{10})&=&M_{10}&M_{10}&=&(a_{10})\\
(a_{11})&=&M_{11}&M_{11}&=&(a_{11})\\
(a_{12})&=&M_{12}&M_{12}&=&(a_{12}).
\end{array}
\]
and if one similarly puts
\[
\begin{array}{rcl@{\qquad}rcl}
(b_7)&=&N_7-N_9&N_7&=&(b_7)+(b_9)+(b_8)\\
(b_8)&=&N_8+N_{11}-N_{12}&N_8&=&(b_8)-(b_{11})+(b_{12})\\
(b_9)&=&-N_8+N_9-N_{11}+N_{12}&N_9&=&(b_8)+(b_9)\\
(b_{10})&=&N_{10}&N_{10}&=&(b_{10})\\
(b_{11})&=&N_{11}&N_{11}&=&(b_{11})\\
(b_{12})&=&N_{12}&N_{12}&=&(b_{12}),
\end{array}
\]
then the equation existing between the \((a),(b)\) passes into
\[
0=(M_{10}N_7-M_7N_{10})+(M_8N_9-M_9N_8)+(M_{12}N_{11}-M_{11}N_{12}).
\]
Thus if the periodicity moduli corresponding to the three different integrals of the first kind are denoted by \(M,N,P\), one has the three equations
\[
\begin{array}{l}
(M_{10}N_7-M_7N_{10})+(M_8N_9-M_9N_8)+(M_{12}N_{11}-M_{11}N_{12})=0\\
(N_{10}P_7-N_7P_{10})+(N_8P_9-N_9P_8)+(N_{12}P_{11}-N_{11}P_{12})=0\\
(P_{10}M_7-P_7M_{10})+(P_8M_9-P_9M_8)+(P_{12}M_{11}-P_{11}M_{12})=0.
\end{array}
\]
If therefore the constants of the integrals are determined so that
\[
\begin{array}{lll}
M_7=2i\pi,&N_7=0,&P_7=0,\\
M_9=0,&N_9=2i\pi,&P_9=0,\\
M_{11}=0,&N_{11}=0,&P_{11}=2i\pi,
\end{array}
\]
the preceding equations pass into
\[
M_8=N_{10},\qquad N_{12}=P_8,\qquad P_{10}=M_{12},
\]
and the complete system of the normal periodicity moduli becomes
\[
\begin{array}{llllll}
M_7=2i\pi,&M_9=0,&M_{11}=0,&M_{10}=a_{11},&M_8=a_{12},&M_{12}=a_{13},\\
N_7=0,&N_9=2i\pi,&N_{11}=0,&N_{10}=a_{12},&N_8=a_{22},&N_{12}=a_{23},\\
P_7=0,&P_9=0,&P_{11}=2i\pi,&P_{10}=a_{13},&P_8=a_{23},&P_{12}=a_{33}.
\end{array}
\]


% Source PDF pp. 134--149 / printed pp. 112--127.
\section*{Fifth Section.}
\subsection*{Interchange of Parameter and Argument.}

\subsubsection*{\S\,32. More precise determination of the integral of the third kind.}

An integral of the third kind behaves, with respect to the branch-points, quite similarly to an integral of the first kind; in general it changes its value when the path over which the integral is carried crosses a branch-point during its displacement. But in addition the value of an integral of the third kind also changes when, during its displacement, the path crosses one of the two points of discontinuity. One can therefore derive the most general value of the integral from one of them by first adjoining cyclic paths leading around branch-points, and then a further path which encloses one or the other of the two points of discontinuity.

In the first section, \S\,7, it was shown that, near the first point of discontinuity \(\xi\), our normal integral of the third kind behaves as \(-\log(x-\xi)\), while near the second it behaves as \(\log(x-\eta)\). Hence, in a circuit in a small circle, in the direction of the motion of the hands of a clock, if the small circle has the first point of discontinuity as its centre, the integral of the third kind changes by \(2\pi\sqrt{-1}\), but by \(-2\pi\sqrt{-1}\) if the second point of discontinuity is the centre. Thus a circuit around the two points of discontinuity simultaneously does not change the integral at all, while circuits around one of them change the integral only by \(2m\pi\sqrt{-1}\), where \(m\) is a positive or negative integer or zero.

We now introduce here, as with the integrals of the first kind, the concept of the periodicity modulus in such a way that it denotes the integral over a cyclic path; and if \(S_{\xi\eta}^{(1)}\), \(S_{\xi\eta}^{(2)}\), \(\ldots\), \(S_{\xi\eta}^{(2p)}\) are mutually independent periodicity moduli of the integral \(S_{\xi\eta}\), then the most general value which this integral can assume by alteration of the path of integration is:
\[
\Sigma_{\xi\eta}=S_{\xi\eta}+2m\pi\sqrt{-1}+m^{(1)}S_{\xi\eta}^{(1)}+m^{(2)}S_{\xi\eta}^{(2)}+\cdots+m^{(2p)}S_{\xi\eta}^{(2p)},
\]
where the \(m\)'s are integers.

In what follows we shall restrict the value of \(S_{\xi\eta}\) to this extent: we wish to connect the points of discontinuity \(\xi,\eta\) by a curve which leads from \(\xi\), with the corresponding value \(s_\xi\), to \(\eta\), with the corresponding value \(s_\eta\); and which cuts none of the loops leading to the branch-points in such a way that at the point of intersection the values of \(s_x\) also coincide. We agree never to cross this curve. The complete integrals \(S_{\xi\eta}^{(1)}\), \(S_{\xi\eta}^{(2)}\), \(\ldots\), which at first we supposed to have been led over a number of loops, can then indeed be led over equivalent paths by extending the systems of loops to closed curves; but these too may never be shifted over one of the points \(\xi,\eta\), so that those integrals acquire genuinely fully determined values. At the same time, the term \(2m\pi\sqrt{-1}\) drops out of the above formula when \(S_{\xi\eta}^{(1)}\), \(S_{\xi\eta}^{(2)}\), etc. are understood in the indicated way.

We shall further extend the curve which connects the two points of discontinuity to a loop which goes from \(0\) to a point of that curve, along it to the point of discontinuity \(\xi\), around this point in a small circle, and indeed in the direction of the hands of a clock, then along the curve to the other point of discontinuity, around this point, again along a part of the curve, and back to \(0\) on the same path running from \(0\) to the curve. This loop, which may cut none of those led to the branch-points in such a way that at the point of intersection the values of \(s_x\) also coincide, is to be added to the earlier system of loops. It is assigned to one definite root-value occurring at \(0\); for only from one such value will one be led at the points of discontinuity to the corresponding values of \(s_x\), whereas for other initial values that loop does not exist.

But the integral carried over this new loop is identically zero; for by known theorems the loop, which is a closed path for the values of \(s_x\), can be reduced to two circles described in the same sense around \(\xi\) and \(\eta\); and since the value of the integral taken over the two circles is equal and opposite, \((2\pi\sqrt{-1}\) and \(-2\pi\sqrt{-1})\), the sum of the resulting values is zero.

If, then, the integral \(S_{\xi\eta}\) is formed around one of the closed circuits mentioned in the preceding section, and these circuits are now enlarged only by the new loop, one obtains zero equal to a sum of integrals over individual loops, which, as before, can be completed to periodicity moduli; but the new loop has no share in these equations. Thus the linear relations between complete integrals of the third kind are exactly the same as the relations between the corresponding periodicity moduli of integrals of the first kind.

\subsubsection*{\S\,33. Interchange of parameter and argument.}

Let us now consider two integrals of the third kind, \(S_{\xi\eta}\), \(S_{\alpha\beta}\), whose points of discontinuity are denoted by the appended indices. For each of these integrals we form the loop encircling its points of discontinuity, and insert these loops among the circuits; here we arrange the position of the loops so that they cut neither one another nor the loops directed toward the branch-points in the forbidden manner described above.

The integral \(\int S_{\xi\eta}\,dS_{\alpha\beta}\), carried over a complete circuit, then gives zero, because it can be contracted to a point. The same integral, carried out successively over all \(n\) circuits, therefore gives rise to an equation quite similar to that developed in the preceding section. Since the value of \(S_{\xi\eta}\), or of \(S_{\alpha\beta}\), changes in no way after a circuit around a loop encircling the points of discontinuity, these loops have no influence on the function-values of the following ones. Thus, insofar as the result arises from these remaining loops, it has exactly the form previously obtained for integrals of the first kind, especially since the conditions existing between the complete integrals are exactly the same as before. If we denote the mutually independent periodicity moduli of the two integrals \(S_{\xi\eta}\) and \(S_{\alpha\beta}\) by \(A_i,B_i\) \((i=1,2,\ldots,2p)\), then the part of our integral arising from the loops of the branch-points is
\[
\sum c_{ik}B_iA_k,
\]
where the integers \(c\) have the values assigned to them in the preceding section.

As for the share of the integral which comes from the loops of the points of discontinuity, it is plainly necessary to consider only those parts of them which connect \(\xi\) with \(\eta\), respectively \(\alpha\) with \(\beta\), backwards and forwards. The path lying between \(\xi\) and \(\eta\) is traversed once in the direction from \(\xi\) to \(\eta\), and another time in the opposite direction; at corresponding places, therefore, \(dS_{\alpha\beta}\) has the same but opposite value. At corresponding places of the two paths, however, the value of the integral \(S_{\xi\eta}\) differs by \(2\pi\sqrt{-1}\), since in the meantime one of the points of discontinuity, \(\eta\), has been encircled. Thus the value of the element \(S_{\xi\eta}dS_{\alpha\beta}\) is the first time \(S_{\xi\eta}dS_{\alpha\beta}\), the second time \(-(S_{\xi\eta}-2\pi\sqrt{-1})dS_{\alpha\beta}\), so that together they yield the value \(2\pi\sqrt{-1}\,dS_{\alpha\beta}\). Hence the constituent of our integral arising from the first loop is
\[
2\pi\sqrt{-1}\int_{\xi}^{\eta}dS_{\alpha\beta}.\footnotemark
\]
\footnotetext{We understand the notation of the limits so that \(\xi,\eta\) are not only the variables, but also the points \(\xi,\eta\) of the curve \(f=0\); thus not merely one coordinate, but both coordinates of the same points are assumed to be given.}

As for the loop leading around \(\alpha,\beta\), we can contract it to two small circles described around \(\alpha,\beta\), where \(S_{\xi\eta}\) then respectively receives the constant values which belong to it at the points \(\alpha,\beta\). Since the loop was to move around the points \(\alpha,\beta\) in the direction of the hands of a clock, the integration of \(dS_{\alpha\beta}\) gives at \(\alpha\) the value \(2\pi\sqrt{-1}\), at \(\beta\) the value \(-2\pi\sqrt{-1}\); and the constituent arising from this loop is thus the sum of \(2\pi\sqrt{-1}\) multiplied by the value of \(S_{\xi\eta}\) at \(\alpha\), and \(-2\pi\sqrt{-1}\) multiplied by the value of \(S_{\xi\eta}\) at \(\beta\). Or it may be represented by
\[
2\pi\sqrt{-1}\int_{\beta}^{\alpha}dS_{\xi\eta}.
\]

Thus finally we have obtained the equation
\[
\tag{1}
2\pi\sqrt{-1}\left[\int_{\alpha}^{\beta}dS_{\xi\eta}-\int_{\xi}^{\eta}dS_{\alpha\beta}\right]=\sum c_{ik}B_iA_k.
\]

\noindent\textbf{I.} If, therefore, one forms two integrals of the third kind in which the limits of one are the parameters of the other, and conversely, their difference can be represented by the periodicity moduli of both integrals.

\subsubsection*{\S\,34. Finite normal form of the integrals of the third kind. Simplest form of the interchange theorem. Theorems which follow from it.}

This important theorem can be stated more precisely if we determine the normal form of our integrals of the third kind somewhat more exactly. We have already seen in the first section that the integrals of the third kind, as we set them up, were determined only up to additive integrals of the first kind, that is, up to linear expressions with \(p\) constants. Thus, without changing anything, instead of an integral \(S_{\xi\eta}\) we may put the expression
\[
S_{\xi\eta}+c_1u_1+c_2u_2+\cdots+c_pu_p,
\]
where we have deliberately expressed the linear function by our normal integrals of the first kind. If we now form the first half of those combinations of the periodicity moduli which above we called normal for integrals of the first kind, and which here too we shall call normal for integrals of the third kind, we obtain the values
\[
S_{\xi\eta}^{(1)}+2\pi\sqrt{-1}\,c_1,\quad
S_{\xi\eta}^{(2)}+2\pi\sqrt{-1}\,c_2,\quad\ldots,\quad
S_{\xi\eta}^{(p)}+2\pi\sqrt{-1}\,c_p,
\]
where \(S_{\xi\eta}^{(1)}\), \(S_{\xi\eta}^{(2)}\) \(\ldots\) denote the parts originating from the somehow fixed integral \(S_{\xi\eta}\). Plainly we can now determine the constants \(c\) so that all these expressions vanish. Then the expression to be put in place of \(S_{\xi\eta}\),
\[
\tag{2}
\Pi_{\xi\eta}=S_{\xi\eta}-\frac{1}{2\pi\sqrt{-1}}\left\{S_{\xi\eta}^{(1)}u_1+S_{\xi\eta}^{(2)}u_2+\cdots+S_{\xi\eta}^{(p)}u_p\right\},
\]
is completely determined; and henceforth we shall designate this expression exclusively as the normal integral of the third kind. Thus the normal integral of the third kind is determined so that the first half of its normal periodicity moduli vanish.

But in the preceding chapter it was shown that the expression \(\sum c_{ik}a_ib_k\) can be transformed so that each term is the product of a periodicity modulus of one half with one of the other half. If, therefore, we apply the same transformation to the right-hand side of equation (1), it must vanish, since we have now normalized the integrals of the third kind so that the whole of one half of their periodicity moduli vanishes. Equation (1) therefore becomes
\[
\tag{3}
\int_{\xi}^{\eta}d\Pi_{\alpha\beta}=\int_{\alpha}^{\beta}d\Pi_{\xi\eta},
\]
and one has the following theorem for the interchange of parameter and argument:

\noindent\textbf{II.} The normal integral of the third kind \(\Pi\) does not change if one interchanges the limits with the points of discontinuity.

This theorem is capable of very many applications. Recall, for example, the theorem established in the first section that the sum of three integrals of the third kind, two of which always have one point of discontinuity in common, is an integral of the first kind. We can write this sum as
\[
\int_{\alpha}^{\beta}\left(d\Pi_{\xi\eta}+d\Pi_{\eta\zeta}+d\Pi_{\zeta\xi}\right),
\]
and, with the aid of our theorem, transform it into
\[
\int_{\xi}^{\eta}d\Pi_{\alpha\beta}+\int_{\eta}^{\zeta}d\Pi_{\alpha\beta}+\int_{\zeta}^{\xi}d\Pi_{\alpha\beta},
\]
which is identically zero if the closed curve composed of the paths \(\xi\eta\), \(\eta\zeta\), \(\zeta\xi\) can be contracted to a point. Thus, provided that the four integrals with points of discontinuity \(\xi\eta\), \(\eta\zeta\), \(\zeta\xi\), \(\alpha\beta\) satisfy the conditions laid down above in the geometric representation--namely that the loops used for their definition, led around the points of discontinuity, do not cut one another and the loops of the branch-points--one has the theorem:

\noindent\textbf{III.} The sum of three normal integrals of the third kind, any two of which have a point of discontinuity in common, is identically zero.

If we take for \(\zeta\) any numerical value, the following theorem further follows:

\noindent\textbf{IV.} The normal integral of the third kind is the difference of two functions of the same kind, each of which depends on only one of its parameters.

If such functions are denoted by \(\log \Phi_\xi\), \(\log \Phi_\eta\), where \(\Phi\) is by no means a completely determined function, then one has
\[
\tag{4}
\Pi_{\xi\eta}=\log\frac{\Phi_\eta}{\Phi_\xi}.
\]

Theorem III also serves to determine the integral of the first kind to which the sum of the corresponding \(S\)'s is equal. For if, using (2), one puts the value of \(\Pi\) in terms of \(S\), one obtains
\[
\tag{5}
\int_{\alpha}^{\beta}\left(dS_{\xi\eta}+dS_{\eta\zeta}+dS_{\zeta\xi}\right)
=\frac{1}{2\pi\sqrt{-1}}\sum_{h=1}^{h=p}
\left(S_{\xi\eta}^{(h)}+S_{\eta\zeta}^{(h)}+S_{\zeta\xi}^{(h)}\right)
\int_{\alpha}^{\beta}du_h.
\]

It was determined above that the first half of the periodicity moduli for the integral of the third kind \(\Pi_{\xi\eta}\) should vanish. It remains to determine the other half of the periodicity moduli for this integral. We obtain them by letting, in equation (1), the one integral \(S_{\alpha\beta}\) pass into an integral of the first kind \(u_h\). Then \(\int_{\alpha}^{\beta}dS_{\xi\eta}\) vanishes; and if we transform \(S\) into \(\Pi\), the half of the terms on the right vanish because they are multiplied by the vanishing periodicity moduli of \(\Pi_{\xi\eta}\). The corresponding periodicity moduli of \(u\), however, do not all vanish; rather, the \(h^{\mathrm{th}}\) remains, which is equal to \(2\pi\sqrt{-1}\). This is now multiplied, in the transformed expression of the right-hand side, by the negative periodicity modulus of \(S_{\xi\eta}\) which is furnished with the index \(h\) in the second half of the periodicity moduli. If we denote it by \(\Pi_{\xi\eta}^{(h)}\), then we have the equation
\[
\tag{6}
\Pi_{\xi\eta}^{(h)}=-\int_{\xi}^{\eta}du_h.
\]

\noindent\textbf{V.} The second half of the periodicity moduli of the integral \(\Pi\) is therefore expressed by the integrals of the first kind taken between the points of discontinuity as limits.

If in (5) one substitutes for \(\Pi\) its expression in \(S\), then the second series \(\Sigma^{(h)}\) of the periodicity moduli of \(S\) is obtained from the formula
\[
\tag{7}
\Sigma_{\xi\eta}^{(h)}=-\int_{\xi}^{\eta}du_h+\frac{1}{2\pi\sqrt{-1}}\sum_{k=1}^{k=p}S_{\xi\eta}^{(k)}a_{kh}.
\]

\subsubsection*{\S\,35. Theorems on integrals of the second kind which follow from the interchange theorem.}

We arrive at a further series of theorems if, in the preceding discussion, we let one or both integrals of the third kind pass into integrals of the second kind, or, what is the same, if we differentiate with respect to a parameter.

We have set above, p. 32,
\[
E_\xi=\frac{\partial S_{\xi\zeta}}{\partial \xi}.
\]
This integral \(E_\xi\) is not independent of the parameter \(\xi\), but it contains \(\xi\) only in a part which is an aggregate of integrals of the first kind. In fact, if we put for \(S\) its expression through \(\Pi\) and through the first \(p\) periodicity moduli of \(S\), we have
\[
\tag{8}
E_\xi=\frac{\partial \Pi_{\xi\zeta}}{\partial \xi}
+\frac{1}{2\pi\sqrt{-1}}\sum_{h=1}^{h=p}
\frac{\partial S_{\xi\zeta}^{(h)}}{\partial \xi}\,u_h.
\]

Let us first consider the expression depending only on the one parameter \(\xi\),
\[
\tag{9}
Z_\xi=\frac{\partial \Pi_{\xi\zeta}}{\partial \xi}=\frac{\partial \log\Phi_\xi}{\partial \xi},
\]
and enclose it between the limits \(\alpha,\beta\). Then by Theorem II we find
\[
\tag{10}
\int_{\alpha}^{\beta}dZ_\xi
=\frac{\partial}{\partial \xi}\int_{\xi}^{\zeta}d\Pi_{\alpha\beta}
=-\left(\frac{\partial\Pi_{\alpha\beta}}{\partial x}\right)_{x=\xi}
=-\left(\frac{\partial\Pi_{\alpha\beta}}{\partial x}\right)_\xi.
\]
The right-hand side, however, is an algebraic function of the values of \(x,s_x\) belonging to the point \(\xi\). Thus one has the theorem:

\noindent\textbf{VI.} The integral of the second kind \(Z\), taken between definite limits and arising from differentiation of \(\Pi\) with respect to a parameter, is an algebraic function of the parameter with respect to which it is differentiated.

Putting for \(\Pi\) its value in \(S\), we obtain the formula
\[
\tag{11}
\int_{\alpha}^{\beta}dZ_\xi
=-\left(\frac{\partial S_{\alpha\beta}}{\partial x}\right)_\xi
+\frac{1}{2\pi\sqrt{-1}}\sum_{h=1}^{h=p}S_{\alpha\beta}^{(h)}
\left(\frac{\partial u_h}{\partial x}\right)_\xi,
\]
and therefore the theorem:

\noindent\textbf{VII.} The integral of the second kind \(Z\), taken between definite limits, is expressed by an algebraic function and by a linearly occurring series of transcendent functions of its limits, which are the first series of periodicity moduli of that integral of the third kind whose points of discontinuity are those limits.

Finally, for \(E\) one obtains from (8) the formula
\[
\tag{12}
\begin{aligned}
\int_{\alpha}^{\beta}dE_\xi
&=-\left(\frac{\partial S_{\alpha\beta}}{\partial x}\right)_\xi\\
&\quad+\frac{1}{2\pi\sqrt{-1}}\sum_{h=1}^{h=p}\left[
S_{\alpha\beta}^{(h)}\left(\frac{\partial u_h}{\partial x}\right)_\xi
+\frac{\partial S_{\xi\zeta}^{(h)}}{\partial \xi}\int_{\alpha}^{\beta}du_h\right],
\end{aligned}
\]
from which the theorem follows:

\noindent\textbf{VIII.} The integral of the second kind \(E\), taken between definite limits, is expressed by algebraic functions and by integrals of the first kind, and moreover linearly

1) by the first series of periodicity moduli of that integral of the third kind \(S\) whose points of discontinuity are the limits,

2) by the periodicity moduli of the first series
\[
\tag{13}
E_\xi^{(h)}=\frac{\partial S_{\xi\zeta}^{(h)}}{\partial \xi}
\]
of the integral \(E\) itself.

From equations (8), (9), (13) we further find the combination
\[
\tag{14}
Z_\xi=E_\xi-\frac{1}{2\pi\sqrt{-1}}\sum_{h=1}^{h=p}E_\xi^{(h)}u_h,
\]
which is entirely analogous to equation (2) between \(\Pi\) and \(S\). Hence it follows immediately:

\noindent\textbf{IX.} The first series of periodicity moduli of \(Z_\xi\) vanishes.

If equation (6) is then differentiated with respect to \(\xi\), one obtains the second series of periodicity moduli of \(Z\):
\[
\tag{15}
Z_\xi^{(h)}=\left(\frac{\partial u_h}{\partial x}\right)_\xi.
\]
Thus:

\noindent\textbf{X.} The second series of periodicity moduli of the function \(Z\) consists of algebraic functions of the parameter.

And the second series of the periodicity moduli of \(E_\xi\) is composed, as follows from (14), or by differentiating (7), from the first series and from the periodicity moduli of \(S_{\xi\zeta}^{(h)}\), with the aid of the formula
\[
\tag{16}
H_\xi^{(h)}=\left(\frac{\partial u_h}{\partial x}\right)_\xi
+\frac{1}{2\pi\sqrt{-1}}\sum_{k=1}^{k=p}E_\xi^{(k)}a_{kh}.
\]

From the above theorems one sees that the integrals of the second kind \(E_\xi\) can always be expressed by algebraic functions and integrals of the first kind, by their periodicity moduli used in forming the normal integrals of the first kind, by the first series of periodicity moduli of the integrals of the third kind, and finally by their differential quotients taken with respect to a parameter, which form the first series of the periodicity moduli of the integral of the second kind.

For the latter quantities one can now set up further relations which prove the theorem:

\noindent\textbf{XI.} The first series of periodicity moduli of the integrals \(E_\xi\), with arbitrary point of discontinuity \(\xi\), can always be expressed by the corresponding periodicity moduli of such \(p\) integrals \(E_\alpha,E_\beta\ldots E_\pi\) as become discontinuous at definite, arbitrarily chosen points \(\alpha,\beta\ldots\pi\).

A further property of the function \(Z\) leads to the proof of this theorem. Namely, if the appended index indicates that in the function for \(x,s_x\) the values corresponding to the point indicated by the index are to be set, then
\[
\left(\frac{\partial Z_\xi}{\partial x}\right)_\alpha
=-\frac{\partial}{\partial\alpha}\int_{\alpha}^{\beta}dZ_\xi
=-\frac{\partial}{\partial\alpha}\int_{\alpha}^{\beta}d\frac{\partial\Pi_{\xi\zeta}}{\partial\xi}
=-\frac{\partial^2}{\partial\alpha\partial\xi}\int_{\alpha}^{\beta}d\Pi_{\xi\zeta}.
\]
Since this expression does not change if one interchanges \(\xi\) with \(\alpha\), and \(\zeta\) with \(\beta\), one has the theorem:

\noindent\textbf{XII.} If one differentiates \(Z_\xi\) with respect to the argument and then substitutes for \(s_x,x\) the values which these variables receive for an arbitrary point \(\alpha\), the result is the same as if, in the differential quotient of \(Z_\alpha\) with respect to the argument \(s_x,x\), one substitutes the values corresponding to the point \(\xi\).

This theorem can also be expressed in the following form:

\noindent\textbf{XIII.} If one has \(q\) points and forms sums \(M_1,M_2,\ldots,M_q\) so that each \(M\) is the sum of \(q-1\) functions \(Z\), whose arguments correspond to \(q-1\) of the above points, while the parameter corresponds to the \(q^{\mathrm{th}}\) point, and if \(x_1,x_2\ldots x_q\) are the values of \(x\) for these \(q\) points, then always
\[
M_1dx_1+M_2dx_2+\cdots+M_qdx_q
\]
is a complete differential.

But equation
\[
\tag{17}
\left(\frac{\partial Z_\xi}{\partial x}\right)_\alpha=
\left(\frac{\partial Z_\alpha}{\partial x}\right)_\xi
\]
gives, if the functions \(Z\) are expressed by the \(E\)'s,
\[
\tag{18}
\begin{aligned}
&\left(\frac{\partial E_\xi}{\partial x}\right)_\alpha
-\frac{1}{2\pi\sqrt{-1}}\sum_{h=1}^{h=p}E_\xi^{(h)}\left(\frac{\partial u_h}{\partial x}\right)_\alpha\\
&\qquad=\left(\frac{\partial E_\alpha}{\partial x}\right)_\xi
-\frac{1}{2\pi\sqrt{-1}}\sum_{h=1}^{h=p}E_\alpha^{(h)}\left(\frac{\partial u_h}{\partial x}\right)_\xi.
\end{aligned}
\]

If one now forms this equation by using, in place of \(\alpha\), successively \(p\) arbitrarily chosen constant points \(\alpha,\beta,\ldots,\pi\), and observes that the bracketed quantities are purely algebraic expressions, then one has before oneself a system of linear equations by virtue of which the \(E_\xi^{(h)}\) can be expressed by \(E_\alpha^{(h)}\), \(E_\beta^{(h)}\ldots E_\pi^{(h)}\), and by algebraic functions of \(\xi\), as Theorem XI requires.

If one similarly forms, from equations (12), by putting for \(\xi\) successively \(a,b,c\ldots\), \(p\) different equations, then the complete integrals \(S_{\alpha\beta}^{(h)}\) are found, by means of a system of linear equations, expressed by purely algebraic integrals of the second kind, by algebraic functions, and by the quantities \(E_a^{(h)}\), \(E_b^{(h)}\), etc. Thus one also has the theorem:

\noindent\textbf{XIV.} The complete integrals \(S_{\alpha\beta}^{(h)}\) can be expressed by purely algebraic integrals of the first and second kinds and by algebraic functions, where only complete integrals of the second kind with absolutely fixed points of discontinuity occur in the coefficients.

\subsubsection*{\S\,36. Reduction of the differential quotients of integrals of the third kind with respect to a parameter to integrals of the second and third kinds.}

In the first section it was proved that every integral of a rational function of \(x\) and \(s_x\) can be expressed by algebraic and logarithmic functions, by integrals of the first and third kinds, and by the differential quotients of the latter taken with respect to the points of discontinuity. We now wish to show how all higher differential quotients of the integrals of the third kind can also be reduced to algebraic functions, to integrals of the first kind and their periodicity moduli, to integrals of the third kind and their periodicity moduli, and to the first series of the periodicity moduli of \(p\) integrals of the second kind, whose points of discontinuity are chosen arbitrarily.

Namely, differentiate equation (1) \(\mu\) times with respect to \(\xi\), \(\nu\) times with respect to \(\eta\). One then obtains
\[
\frac{\partial^{\mu+\nu}}{\partial\xi^\mu\partial\eta^\nu}\int_{\alpha}^{\beta}d\Pi_{\xi\eta}
=
\frac{\partial^{\mu+\nu}}{\partial\xi^\mu\partial\eta^\nu}\int_{\xi}^{\eta}d\Pi_{\alpha\beta}.
\]
If now the \(\Pi\)'s are expressed by the integrals \(S\), this becomes
\[
\tag{19}
\begin{aligned}
&\frac{\partial^{\mu+\nu}}{\partial\xi^\mu\partial\eta^\nu}\int_{\alpha}^{\beta}dS_{\xi\eta}
-\frac{1}{2\pi\sqrt{-1}}\sum_{h=1}^{h=p}
\frac{\partial^{\mu+\nu}S_{\xi\eta}^{(h)}}{\partial\xi^\mu\partial\eta^\nu}\int_{\alpha}^{\beta}du_h\\
&\quad=\frac{\partial^{\mu+\nu}}{\partial\xi^\mu\partial\eta^\nu}\int_{\xi}^{\eta}dS_{\alpha\beta}
-\frac{1}{2\pi\sqrt{-1}}\sum_{h=1}^{h=p}S_{\alpha\beta}^{(h)}
\frac{\partial^{\mu+\nu}}{\partial\xi^\mu\partial\eta^\nu}\int_{\xi}^{\eta}du_h.
\end{aligned}
\]
This equation in fact already expresses the required integral in the desired way, except that the quantities
\[
\frac{\partial^{\mu+\nu}S_{\xi\eta}^{(h)}}{\partial\xi^\mu\partial\eta^\nu},
\]
the higher differential quotients of the first series of periodicity moduli of the integral \(S_{\xi\eta}\), still occur in this formula. But if now in (18) we put \(\xi,\eta\) in place of \(\alpha,\beta\), and put for \(\xi\) successively \(p\) arbitrarily chosen points, and if we then differentiate \(\mu\) times with respect to \(\xi\) and \(\nu\) times with respect to \(\eta\), we have before us \(p\) equations which are linear in the sought quantities
\[
\frac{\partial^{\mu+\nu}S_{\xi\eta}^{(h)}}{\partial\xi^\mu\partial\eta^\nu},
\]
and which, besides, contain only algebraic expressions and the first series of the periodicity moduli, corresponding to those \(p\) points, of integrals of the second kind. By these, therefore, and by algebraic expressions, we can represent the sought quantities; and by entering this into (19), we obtain the required formula.

\subsubsection*{\S\,37. Application of the methods used above for integrals of the third kind to algebraic functions and logarithms of such functions. Abel's theorem.}

We have thus far derived all results which arise from the investigation of the integral \(\int I\,dH\), when \(I,H\) are integrals of the first, second, or third kind. We now wish to investigate the results which arise when for one of the functions \(I,H\) we introduce an algebraic function or the logarithm of such a function.

To treat the integral \(\int \log\frac{\varphi}{\psi}\,dS_{\alpha\beta}\), in which \(\varphi,\psi\) are arbitrary entire functions of order \(m\) in \(s_x,x\), we first construct for ourselves the loop for \(S_{\alpha\beta}\) as before. The functions \(\varphi,\psi\) each vanish at \(mn\) points, corresponding to the intersection-points of the curve \(\varphi=0\), \(f=0\) on the one hand, and \(\psi=0\), \(f=0\) on the other. If some of the first coincide with some of the last, then \(\log\frac{\varphi}{\psi}\) remains finite for them, while for the others \(\log\frac{\varphi}{\psi}\) becomes infinitely large. We shall always connect a point \(s_x,x\) for which \(\varphi\) vanishes with a point \(s_x,x\) for which \(\psi\) vanishes; from these curves we construct loops which lead around each pair of such points and to the point \(0\), as those corresponding earlier to integrals of the third kind were drawn, and we lay them so that no two different loops have a value-pair \(s_x,x\) in common. How one connects the points arising from \(\varphi=0\) with those arising from \(\psi=0\) is immaterial for the first pairs of points; but the last ones we combine so that, if the curve \(\varphi=0\) is carried into \(\psi=0\) by continuous variation, and if this happens in such a way that, of the points first combined, each always passes into its corresponding one, then also, of the later pairs, one point is carried into the other.

The function \(\log\frac{\varphi}{\psi}\) does not change on a circuit around such a loop; for on a circuit around the one point of discontinuity it always changes by \(-2\pi\sqrt{-1}\), and on a circuit around the other by \(+2\pi\sqrt{-1}\), and therefore remains unchanged when both are traversed simultaneously. If we now form \(\int d\log\frac{\varphi}{\psi}\) around a closed circuit, the inserted loops have no influence on the value of this integral, and one obtains among the periodicity moduli of \(\log\frac{\varphi}{\psi}\) the same relations as hold for those of integrals of the first kind. Moreover, all these periodicity moduli have the value \(2m\pi\sqrt{-1}\), where the \(m\)'s may be arbitrary integers; for this is the only increment which \(\log\frac{\varphi}{\psi}\) can receive at all on a circuit around a cyclic path.

The integral \(\int\log\frac{\varphi}{\psi}\,dS_{\alpha\beta}\), taken around all circuits, consists, exactly as before, of one part arising from the integrals led around the branch-points and composed from the periodicity moduli of \(S_{\alpha\beta}\) and \(\log\frac{\varphi}{\psi}\); of a second part arising from the loop laid around \(\alpha,\beta\) and, in the manner explained above, having the value
\[
2\pi\sqrt{-1}\int_{\beta}^{\alpha}d\log\frac{\varphi}{\psi};
\]
and finally of the terms which arise from the points of discontinuity of the function \(\log\frac{\varphi}{\psi}\). These terms can be treated in exactly the same way as those above which arose from \(S_{\xi\eta}\), and one finds
\[
-2\pi\sqrt{-1}\sum_{\xi}\int_{\xi}^{\eta}dS_{\alpha\beta},
\]
where the \(\xi\)'s refer to the points at which \(\varphi=0\), and the \(\eta\)'s to those at which \(\psi=0\).

Dividing now by \(2\pi\sqrt{-1}\), one obtains, since the periodicity moduli of \(\log\frac{\varphi}{\psi}\) are thereby reduced to integers,
\[
\sum_{\xi}\int_{\xi}^{\eta}dS_{\alpha\beta}-\int_{\beta}^{\alpha}d\log\frac{\varphi}{\psi}=K,
\]
where \(K\) is an integral combination of the periods of \(S_{\alpha\beta}\). It can depend on the position of the loops led around the points \(\xi,\eta\) relative to the branch-loops, but not on the special position of the points \(\xi,\eta\). Therefore, to determine \(K\), one may let the curves \(\varphi,\psi\) approach one another infinitely closely; and since by assumption the pairwise assigned points \(\xi,\eta\) are then also to coincide, one has \(K=0\), and hence
\[
\log\left(\frac{\varphi}{\psi}\right)_\alpha-\log\left(\frac{\varphi}{\psi}\right)_\beta
=\sum_{\xi}\int_{\xi}^{\eta}dS_{\alpha\beta}.
\]
This equation is Abel's theorem, and indeed in a form in which both the integrals on the right and the logarithms are fully determined. In it one can, in accordance with the derivation of the formulas, put \(\Pi\) in place of \(S\); and this agrees with the theorem that \(\sum_{\xi}\int_{\xi}^{\eta}du_h\) must vanish, which forms a part of Abel's theorem. In this representation of Abel's theorem it is assumed, first, that the path on which one passes from \(\log(\frac{\varphi}{\psi})_\beta\) to \(\log(\frac{\varphi}{\psi})_\alpha\) is cut by none of the integration paths on the right, and second, that the upper and lower limits of these paths are assigned to one another in such a way that they become infinitely small if one lets the curves \(\varphi=0\), \(\psi=0\) approach one another infinitely closely.

But if one puts \(\Pi\) in place of \(S\), and applies the theorem on the interchange of parameter and argument, one finds:
\[
\mathrm{I.}\quad
\log\left(\frac{\varphi}{\psi}\right)_\alpha-\log\left(\frac{\varphi}{\psi}\right)_\beta
=\sum_{\xi}\int_{\alpha}^{\beta}d\Pi_{\xi\eta}.
\]
This formula contains the reduction of the logarithm of an algebraic function to a sum of normal integrals of the third kind. It does not cease to hold when some of the points \(\xi\), or of the points \(\eta\), are made to approach one another very closely, whereby only individual parts of the loops laid around the points of discontinuity approach one another infinitely closely; or also whole loops, as soon as two pairs of assigned points approach one another very closely, that is, when corresponding points of the curves \(\varphi=0\), \(\psi=0\) become points of contact of these curves with \(f=0\). ---

If we now pass to the investigation of the integral \(\int \frac{\varphi}{\psi}\,dS_{\alpha\beta}\), we first lay down a system of loops, which becomes necessary because of the discontinuity points of \(\frac{\varphi}{\psi}\), in the following manner. Each loop goes out from the point \(0\) to one of the discontinuity points of \(\frac{\varphi}{\psi}\), around this point in a small circle corresponding to the motion of the hand of a clock, and back to \(0\). At \(0\) it must begin with a definite value of \(s_x\), in order that near the discontinuity value of \(x\) it should at the same time approach the discontinuity value of \(s_x\), and accordingly it is inserted into a definite circuit. Again these loops must not cut one another, nor the loops of the branch-points, in the forbidden manner described above.

Since the periodicity moduli of \(\frac{\varphi}{\psi}\) are all zero, because the function, after traversing a cyclic path, always recovers its original value, in the integral \(\int \frac{\varphi}{\psi}\,dS_{\alpha\beta}\) one has to investigate only those parts which come from the loops laid around the discontinuity points. The loop belonging to \(\alpha,\beta\) gives, similarly as above,
\[
2\pi\sqrt{-1}\int_{\beta}^{\alpha} d\frac{\varphi}{\psi}
=2\pi\sqrt{-1}\left\{\left(\frac{\varphi}{\psi}\right)_{\alpha}-\left(\frac{\varphi}{\psi}\right)_{\beta}\right\}.
\]
If further \(\xi\) is one of the discontinuity points of \(\frac{\varphi}{\psi}\), and if the loop corresponding to it is contracted to an infinitely small circle described about the discontinuity point, then the part of the integral arising from this loop becomes
\[
-2\pi\sqrt{-1}\left(\frac{\varphi}{\frac{\partial\psi}{\partial x}}\right)_{\xi}\cdot
\left(\frac{\partial S_{\alpha\beta}}{\partial x}\right)_{\xi},
\]
so that one obtains the equation
\[
\left(\frac{\varphi}{\psi}\right)_{\alpha}-\left(\frac{\varphi}{\psi}\right)_{\beta}
=\sum\left(\frac{\varphi}{\frac{\partial\psi}{\partial x}}\right)_{\xi}
\left(\frac{\partial S_{\alpha\beta}}{\partial x}\right)_{\xi},
\]
or
\[
\left(\frac{\varphi}{\psi}\right)_{\alpha}-\left(\frac{\varphi}{\psi}\right)_{\beta}
=-\sum\left(\frac{\varphi}{\frac{\partial\psi}{\partial x}}\right)_{\xi}
\frac{\partial}{\partial\xi}\int_{\xi}^{\zeta}dS_{\alpha\beta}.
\]
Here one may now put the function \(\Pi\) for \(S\), and then, applying the theorem on the interchange of parameter and argument, finds
\[
\mathrm{II}\ldots\quad
\left(\frac{\varphi}{\psi}\right)_{\alpha}-\left(\frac{\varphi}{\psi}\right)_{\beta}
=-\sum\left(\frac{\varphi}{\frac{\partial\psi}{\partial x}}\right)_{\xi}
\int_{\alpha}^{\beta}dZ_{\xi}.
\]
This formula contains, if \(\alpha\) is regarded as the variable, the representation of an arbitrary rational function of \(\alpha\) and \(s_{\alpha}\) as an aggregate of integrals of the second kind. If several points \(\xi\) move infinitely close to one another, differential quotients of integrals \(Z\) with respect to the parameter occur.

\section*{Sixth Section.}
\subsection*{The inversion problem.}

\subsubsection*{\S\,38. Auxiliary theorem for judging periodic functions.}

The natural question whether, just as from the integral
\[u=\int_0^x\frac{dx}{\sqrt{1-x^2}}
\]
the upper limit appears as the function \(x=\sin u\) of the integral, or from \(u=\int\frac{dx}{x}\) the upper limit appears as the function \(x=e^u\) of the integral, one can also, for a general algebraic integral of the first kind, treat the upper limit as a function of the integral, was considered by \textsc{Jacobi} in a famous paper, to which we owe the notion of Abelian functions, and he fixed the character of the functions thereby arising.\footnote{Crelle's Journal, vol. XIII.} Although Jacobi's investigation refers initially only to the case \(p=2\), the generalization is immediate. One first needs an auxiliary theorem, which Jacobi states for that case and which, generalized, reads as follows:

Let \(q-1\) rows of \(q\) real quantities each be given:
\[
(1)\ldots\quad
\begin{array}{cccc}
a_1^{(1)},&a_2^{(1)},&\ldots&a_q^{(1)}\\
a_1^{(2)},&a_2^{(2)},&\ldots&a_q^{(2)}\\[-0.15em]
\multicolumn{4}{c}{\cdots\cdots\cdots}\\
a_1^{(q-1)},&a_2^{(q-1)},&\ldots&a_q^{(q-1)}.
\end{array}
\]
Then one can always determine \(q\) positive or negative integers \(m_1,m_2\ldots m_q\) so that the \(q-1\) expressions
\[
(2)\ldots\quad
\begin{array}{rcl}
Q^{(1)}&=&m_1a_1^{(1)}+m_2a_2^{(1)}+\cdots+m_qa_q^{(1)}\\
Q^{(2)}&=&m_1a_1^{(2)}+m_2a_2^{(2)}+\cdots+m_qa_q^{(2)}\\[-0.15em]
\multicolumn{3}{c}{\cdots\cdots\cdots}\\
Q^{(q-1)}&=&m_1a_1^{(q-1)}+m_2a_2^{(q-1)}+\cdots+m_qa_q^{(q-1)}
\end{array}
\]
become respectively smaller than \(q-1\) arbitrarily chosen, however small, quantities \(\varepsilon^{(1)},\varepsilon^{(2)}\ldots\varepsilon^{(q-1)}\).

Assume that the theorem has been proved when, for \(q\), one puts the numbers \(1,2\ldots q-1\). We shall then show that it holds for \(q\) itself.

Among the quantities \(a_q^{(1)},a_q^{(2)}\ldots a_q^{(q-1)}\) at least one must be different from zero. Let it be, for instance, \(a_q^{(q-1)}\). We now form the expressions
\[
(3)\ldots\quad
\begin{array}{c}
Q^{(1)}-Q^{(q-1)}\dfrac{a_q^{(1)}}{a_q^{(q-1)}}\\[0.45em]
Q^{(2)}-Q^{(q-1)}\dfrac{a_q^{(2)}}{a_q^{(q-1)}}\\[-0.1em]
\cdots\cdots\cdots\\[0.25em]
Q^{(q-2)}-Q^{(q-1)}\dfrac{a_q^{(q-2)}}{a_q^{(q-1)}}.
\end{array}
\]
The number of these equations is \(q-2\), the number of unknown integers contained in them is \(q-1\), since \(m_q\) is eliminated. Hence, by the hypothesis, we can determine the numbers \(m_1,m_2\ldots m_{q-1}\) contained in these expressions so that these \(q-2\) expressions become absolutely smaller than arbitrarily small preassigned quantities. Then one can determine \(m_q\) so that \(\frac{Q^{(q-1)}}{a_q^{(q-1)}}\) becomes absolutely smaller than \(\frac12\). The absolute values of the expressions (3) were smaller than arbitrarily small given numbers, and therefore also the expressions \(Q^{(1)},Q^{(2)}\ldots Q^{(q-1)}\) themselves must approach the values \(\frac12a_q^{(1)},\frac12a_q^{(2)}\ldots\frac12a_q^{(q-2)}\) arbitrarily closely. In particular, if the absolute value of \(\frac{Q^{(q-1)}}{a_q^{(q-1)}}\) is equal to \(\frac12-\varepsilon\), and if, apart from sign,
\[
Q^{(i)}-Q^{(q-1)}\frac{a_q^{(i)}}{a_q^{(q-1)}}=\varepsilon^{(i)},
\]
then the absolute value of \(Q^{(i)}\) is not greater than that of \(\varepsilon_i\) increased by that of \(a_q^{(i)}\left(\frac12-\varepsilon\right)\), or not greater than \(\varepsilon^{(i)}+\frac12\sqrt{\bigl(a_q^{(i)}\bigr)^2}\). Thus, putting \(\varepsilon^{(i)}=\frac{\sqrt{\bigl(a_q^{(i)}\bigr)^2}}{m}\), where \(m\) is an arbitrarily large number, one finally has \(Q^{(i)}\) equal to or smaller than the absolute value of \(\left(\frac12+\frac1m\right)a_q^{(i)}\); hence, for all values of \(i\), one may write, apart from sign,
\[
Q^{(i)}<\left(\frac12+\frac1m\right)a_q^{(i)},
\]
with the sole proviso that if one of the quantities \(a_q^{(1)},a_q^{(2)}\ldots a_q^{(q-2)}\) becomes vanishingly small, the corresponding \(Q^{(i)}\) is to be made smaller than an arbitrarily preassigned bound. If we denote the integral combinations \(Q^{(1)},Q^{(2)}\ldots\) so formed by \(a_{q+1}^{(1)},a_{q+1}^{(2)}\ldots a_{q+1}^{(q-1)}\), so that, apart from sign,
\[
a_{q+1}^{(i)}<\frac12a_q^{(i)}\qquad(i=1,2\ldots q-1),
\]
then, if only one of the quantities \(a_{q+1}^{(i)}\) does not become vanishingly small, one can treat the expressions
\[
\begin{array}{rcl}
R^{(1)}&=&n_2a_2^{(1)}+n_3a_3^{(1)}+\cdots+n_{q+1}a_{q+1}^{(1)}\\
R^{(2)}&=&n_2a_2^{(2)}+n_3a_3^{(2)}+\cdots+n_{q+1}a_{q+1}^{(2)}\\[-0.15em]
\multicolumn{3}{c}{\cdots\cdots\cdots}\\
R^{(q-1)}&=&n_2a_2^{(q-1)}+n_3a_3^{(q-1)}+\cdots+n_{q+1}a_{q+1}^{(q-1)}
\end{array}
\]
just as above the expressions \(Q\); that is, one can determine the integers \(n\) so that
\[
R^{(i)}<\frac12a_{q+1}^{(i)}\qquad(i=1,2\ldots q-1),
\]
or, if these linear integral combinations are denoted by \(a_{q+2}^{(i)}\), so that
\[
a_{q+2}^{(i)}<\frac12a_{q+1}^{(i)}\qquad(i=1,2\ldots q-1),
\]
with the sole proviso that, as soon as one of the quantities \(a_{q+1}^{(i)}\) is vanishingly small, \(a_{q+2}^{(i)}\) is to be made smaller than any desired bound.

Proceeding in this way, one forms rows of quantities,
\[
\begin{array}{cccc}
a_{q+1}^{(1)},&a_{q+2}^{(1)},&\ldots&a_{q+k}^{(1)}\\
a_{q+1}^{(2)},&a_{q+2}^{(2)},&\ldots&a_{q+k}^{(2)}\\[-0.15em]
\multicolumn{4}{c}{\cdots\cdots\cdots}\\
a_{q+1}^{(q-1)},&a_{q+2}^{(q-1)},&\ldots&a_{q+k}^{(q-1)}
\end{array}
\]
each of which is either vanishingly small or smaller than half the quantity immediately standing to its left. Thus this process may be continued until all the \(a\)'s in the last vertical row are smaller than arbitrarily preassigned quantities.

But each \(a\) is an integral linear function of the \(q\) quantities \(a\) which stand to its left (where one has to imagine the above scheme completed to the left by the system of the given \(a\)'s), and the coefficients are the same for all \(a\)'s of the same vertical row. Therefore the \(a\)'s of the last vertical row can also be expressed linearly and integrally through the given \(a\)'s, so that
\[
\begin{array}{rcl}
a_{q+k}^{(1)}&=&\mu_1a_1^{(1)}+\mu_2a_2^{(1)}+\cdots+\mu_qa_q^{(1)}\\
a_{q+k}^{(2)}&=&\mu_1a_1^{(2)}+\mu_2a_2^{(2)}+\cdots+\mu_qa_q^{(2)}\\[-0.15em]
\multicolumn{3}{c}{\cdots\cdots\cdots}\\
a_{q+k}^{(q-1)}&=&\mu_1a_1^{(q-1)}+\mu_2a_2^{(q-1)}+\cdots+\mu_qa_q^{(q-1)}.
\end{array}
\]
These expressions therefore supply, as required in the theorem, integral linear combinations of the given \(a\)'s which have been made arbitrarily small.

The algorithm breaks off in the exceptional case mentioned above, when all quantities \(a\) of a vertical row become vanishingly small. One then already has combinations of the required kind and need not continue the operations.

We have thus proved that our theorem is true for an arbitrary number \(q\) if it holds for \(q-1\). It remains only to prove it for \(q=2\), where it reads as follows: if \(a_1,a_2\) are real quantities, one can always determine two numbers \(m_1,m_2\) so that \(m_1a_1+m_2a_2\) becomes smaller than any preassigned, however small, quantity.

How to find the numbers \(m_1,m_2\) in this case is taught by the theory of continued fractions. If \(\left(\frac{a_1}{a_2}\right)\) is the absolute value of \(\frac{a_1}{a_2}\), and if \(\frac{\mu}{\nu},\frac{\mu'}{\nu'}\) are two convergents of the continued-fraction development of \(\left(\frac{a_1}{a_2}\right)\), then the absolute value of \(\left(\frac{a_1}{a_2}\right)-\frac{\mu}{\nu}\) is known to be smaller than \(\frac{1}{\nu\nu'}\), and consequently the absolute value of \(\nu(a_1)-\mu(a_2)\) is smaller than \(\frac{(a_2)}{\nu'}\). One need therefore continue the continued-fraction development only far enough until \(\frac{(a_2)}{\nu'}\) is sufficiently small, in order to obtain the required determination of a linear expression formed from \((a_1),(a_2)\) and smaller than a preassigned quantity. Thus the theorem stated at the beginning is proved in full generality.

We add, for what follows, that in general no integral linear relation exists between quantities \(a\) with the same upper index. For this reason, in general none of the linear combinations which are smaller than given bounds will become absolutely zero, since this would presuppose the existence of such an integral linear relation.

\subsubsection*{\S\,39. Functions with several periods.}

Let us now consider an integral of the first kind and the most general value that it can assume for fixed given limits. If the periodicity moduli of the integral are
\[
a_1+b_1\sqrt{-1},\quad a_2+b_2\sqrt{-1},\quad\ldots,\quad a_{2p}+b_{2p}\sqrt{-1},
\]
where the \(a,b\) now denote real numbers, and if \(I\) is one value of the integral, then the most general value which it can obtain is
\[
I+m_1(a_1+b_1\sqrt{-1})+m_2(a_2+b_2\sqrt{-1})\cdots+m_{2p}(a_{2p}+b_{2p}\sqrt{-1}).
\]
If now \(p>1\), one can always determine the numbers \(m_1,m_2\ldots m_{2p}\), according to the theorem just proved, so that the expressions
\[
\begin{array}{c}
m_1a_1+m_2a_2+\cdots+m_{2p}a_{2p}\\
m_1b_1+m_2b_2+\cdots+m_{2p}b_{2p}
\end{array}
\]
become arbitrarily small. Thus, by a suitable alteration of the path, one can also alter the value of \(I\) arbitrarily; in other words:

An integral of the first kind, for \(p>1\), can acquire any arbitrary value for arbitrarily given limits by suitable choice of the path.

This theorem could have an exception only if the change in the value of \(I\) were not infinitely small but absolutely zero; that is, if there were integers \(m\) such that
\[
\begin{array}{rcl}
m_1a_1+m_2a_2+\cdots+m_{2p}a_{2p}&=&0,\\
m_1b_1+m_2b_2+\cdots+m_{2p}b_{2p}&=&0,
\end{array}
\]
or such that between the periodicity moduli themselves the relation
\[
m_1(a_1+b_1\sqrt{-1})+m_2(a_2+b_2\sqrt{-1})\cdots+m_{2p}(a_{2p}+b_{2p}\sqrt{-1})=0
\]
existed. In that case one could express these \(2p\) periodicity moduli integrally by \(2p-1\) others, so that the integral would in truth possess not \(2p\), but only \(2p-1\), mutually independent periodicity moduli.

To see this, put, for brevity,
\[
a_i+b_i\sqrt{-1}=P_i,
\]
so that the equation
\[
(1)\ldots\quad m_1P_1+m_2P_2+\cdots+m_{2p}P_{2p}=0
\]
is supposed. If \(m_2'\) is the greatest common divisor of \(m_1,m_2\), one can choose two numbers \(n_1,n_2\) so that
\[
\frac{m_1}{m_2'}n_1-\frac{m_2}{m_2'}n_2=1,
\]
and, putting
\[
\begin{array}{rcl}
n_1P_1+n_2P_2&=&P_1',\\
\dfrac{m_2}{m_2'}P_1+\dfrac{m_1}{m_2'}P_2&=&P_2',
\end{array}
\]
\(P_1',P_2'\) become two new periodicity moduli through which \(P_1,P_2\) can be integrally expressed. Equation (1) now passes into
\[
m_2'P_2'+m_3P_3+m_4P_4+\cdots+m_{2p}P_{2p}=0,
\]
which contains one modulus fewer.

If further \(m_3''\) is the greatest common divisor of \(m_2'\) and \(m_3\), and if one chooses \(n_2',n_3'\) so that
\[
\frac{m_2'}{m_3''}n_2'-\frac{m_3}{m_3''}n_3'=1,
\]
then in place of \(P_2',P_3\) one can again introduce the periodicity moduli
\[
\begin{array}{rcl}
n_2'P_2'+n_3'P_3&=&P_2'',\\
\dfrac{m_3}{m_3''}P_2'+\dfrac{m_2'}{m_3''}P_3&=&P_3'',
\end{array}
\]
through which \(P_2',P_3\) can be integrally expressed. Thus \(P_1,P_2,P_3\) have been expressed linearly and integrally through \(P_1',P_2'',P_3''\), while equation (1) passes into
\[
m_3''P_3''+m_4P_4+\cdots+m_{2p}P_{2p}=0,
\]
which contains only \(2p-2\) moduli.

Continuing in this way, one arrives at expressing \(P_1,P_2\ldots P_{2p}\) integrally through new moduli \(P_1',P_2'',P_3''\ldots P_{2p-1}^{(2p-1)}\), while equation (1) passes into
\[
P_{2p}^{(2p-1)}=0.
\]
Thus one can indeed represent the \(2p\) periodicity moduli as integral linear combinations of only \(2p-1\) moduli, and one has an integral with only \(2p-1\) independent periodicity moduli.

Since an integral of the first kind, for every upper limit, can assume every value, conversely the upper limit, considered as a function of the integral, has every arbitrary value for every arbitrary value of the latter. Thus one has an infinitely many-valued function for which it is not so much the value of the arguments that matters as the manner of their formation, that is, the path which the integral must have traversed in order to attain a certain value; for it is solely through this path that the upper limit finds its determination.

\subsubsection*{\S\,40. Concept of Jacobi's inversion problem.}

Jacobi taught how to remove the difficulties arising from this by setting, instead of a single integral, the sum of \(p\) integrals of the same kind with different upper limits equal to a new variable \(v\). We now form \(p\) equations, in each of which a sum of \(p\) integrals of the same kind is set equal to a new variable, so that in each equation the same upper limits recur, while each equation is characterized by the use of a particular integral of the first kind. These equations connect two systems of variables: the system of upper limits (or rather of the corresponding pairs of values \(x,s_x\)) and the \(p\) quantities \(v\). The task of establishing the upper limits from these equations as functions of the quantities \(v\) is Jacobi's inversion problem.

Let us first see how the difficulty discussed above ceases to exist here. The most general values of the \(v\)'s which arise by varying the paths of integration (where of course we assume that all integrals attached to the same upper limit are also to be taken along the same path) arise from one system of values by adding linear integral combinations of the periodicity moduli; the corresponding periodicity moduli of the different integrals are always to be affected by the same undetermined integer. If we now divide the \(p\) increments which the different functions \(v\) have thus received into their \(2p\) real and imaginary parts, then by suitable determination of the integers we can indeed make \(2p-1\) of these quantities arbitrarily small, but not the \((2p)^{\mathrm{th}}\); rather, this last quantity generally becomes larger the smaller one makes the former quantities, and hence the larger the values of the integers are made by continued approximation. Thus in two systems of values of the \(v\)'s corresponding to the same upper limits, the two values differ, for at least one of the variables, by a finite quantity. One can therefore no longer say that every system of upper limits corresponds to every system of the \(v\)'s. But this would still be the case if one wished to introduce fewer than \(p\) integrals in an equation and fewer than \(p\) equations in order to found an inversion problem upon them.

One may also connect Jacobi's inversion problem, as Jacobi himself did,\footnote{Crelle's Journal, vol. IX.} with Abel's theorem. This theorem teaches how to replace a sum of arbitrarily many integrals of the first kind of the same sort by a sum of \(p\) integrals of the same sort, where the upper limits become algebraic functions of the former ones. It is enough to draw an algebraic curve through the points of the curve \(F(x,s_x)=0\) [or \(f(x_1,x_2,x_3)=0\)] which correspond to the upper limits of the given sum, and through a suitable number of constant points of the curve, so that its system of intersection points is thereby completely determined. The \(p\) points which then complete the system of intersections have the property that the sum of the corresponding integrals of the first kind is equal, up to a constant, to the negative sum of the integrals corresponding to the given points. Applying this consideration to all \(p\) mutually independent integrals of the first kind, one has before oneself a special case of the inversion problem, namely one in which the quantities \(v\) are given as negative integral sums, and in which the inversion problem is resolved by an algebraic equation.

\subsubsection*{\S\,41. Formulation of the inversion problem. Abelian functions and transcendentals.}

Let \(x^{(1)},x^{(2)}\ldots x^{(p)}\) denote the points and variables corresponding to the upper limits, and \(c^{(1)},c^{(2)}\ldots c^{(p)}\) those corresponding to the lower limits. Further, let us take the integrals of the first kind in normal form; then the inversion problem is given by the following equations:
\[
(1)\ldots\quad
\left\{\begin{array}{rcl}
\displaystyle\sum_{i=1}^{i=p}\int_{c^{(i)}}^{x^{(i)}}du_1&=&v_1,\\[0.35em]
\displaystyle\sum_{i=1}^{i=p}\int_{c^{(i)}}^{x^{(i)}}du_2&=&v_2,\\[-0.1em]
\multicolumn{3}{c}{\cdots\cdots\cdots}\\[-0.1em]
\displaystyle\sum_{i=1}^{i=p}\int_{c^{(i)}}^{x^{(i)}}du_p&=&v_p.
\end{array}\right.
\]
If now \(\frac{\varphi}{\psi}\) is any homogeneous function of \(x,s_x\), or of \(x_1,x_2,x_3\), whose numerator and denominator are of equally high order, then the values which this function assumes at the different \(p\) points may be regarded as the roots of an algebraic equation
\[
(2)\ldots\quad
\left(\frac{\varphi}{\psi}\right)^p+M_1\left(\frac{\varphi}{\psi}\right)^{p-1}+M_2\left(\frac{\varphi}{\psi}\right)^{p-2}\cdots+M_p=0,
\]
whose coefficients are functions of the \(v\)'s. These functions, or more generally symmetric functions of the \(p\) values which a homogeneous function of zeroth dimension in \(x,s_x\), or in \(x_1,x_2,x_3\), assumes at \(p\) different points, shall be called Abelian functions of the corresponding quantities \(v\).

The restriction involved in considering here only symmetric functions of the \(p\) values of a function \(\frac{\varphi}{\psi}\) is only apparent. For every rational symmetric function, that is, one unchanged by interchange of two points, of the coordinate pairs \(x^{(1)},s_{x^{(1)}};x^{(2)},s_{x^{(2)}};\ldots x^{(p)},s_{x^{(p)}}\), is known to be composed of functions of the kind considered by us. Hence all investigations to be undertaken for these special functions also have validity for those more general ones, and we can extend the definition of Abelian functions to the latter without broadening or narrowing the concept already established.

The most general values which the \(v\)'s can acquire by alteration of the integration paths with unchanged upper limits are
\[
(3)\ldots\quad
\begin{array}{rcl}
v_1'&=&v_1+2m_1\pi\sqrt{-1}+a_{11}q_1+a_{12}q_2\cdots+a_{1p}q_p\\
v_2'&=&v_2+2m_2\pi\sqrt{-1}+a_{21}q_1+a_{22}q_2\cdots+a_{2p}q_p\\[-0.15em]
\multicolumn{3}{c}{\cdots\cdots\cdots}\\
v_p'&=&v_p+2m_p\pi\sqrt{-1}+a_{p1}q_1+a_{p2}q_2\cdots+a_{pp}q_p,
\end{array}
\]
where the \(m,q\) denote arbitrary integers. The Abelian functions obviously remain unchanged when one replaces the \(v\)'s by the \(v'\)'s; they therefore remain unchanged when one simultaneously increases the \(v\)'s by one of the following \(2p\) systems:
\[
(4)\ldots\quad
\begin{array}{rcccccc}
v_1)&2\pi\sqrt{-1}&0&\cdots&0&a_{11}&a_{12}\cdots a_{1p}\\
v_2)&0&2\pi\sqrt{-1}&\cdots&0&a_{21}&a_{22}\cdots a_{2p}\\[-0.15em]
\multicolumn{7}{c}{\cdots\cdots\cdots}\\
v_p)&0&0&\cdots&2\pi\sqrt{-1}&a_{p1}&a_{p2}\cdots a_{pp}.
\end{array}
\]
These systems of quantities are therefore called the systems of periods of the Abelian functions, and the latter themselves, since there are \(2p\) such systems, are called \(2p\)-fold periodic functions.

Further, consider the sum of \(p\) integrals of the third kind \(\Pi_{\xi\eta}\), taken between the same points \(c,x\), and having two given discontinuity points \(\xi,\eta\). This sum
\[
(5)\ldots\quad
\sum_{i=1}^{i=p}\int_{c^{(i)}}^{x^{(i)}}d\Pi_{\xi\eta}
=T_{\xi\eta}\begin{pmatrix}x^{(1)},x^{(2)}\ldots x^{(p)}\\ c^{(1)},c^{(2)}\ldots c^{(p)}\end{pmatrix}
\]
can, since it is a symmetric function of the upper limits,


% --- Continuation: source pp. 162--173 ---

\small

% Source: Clebsch--Gordan, Theorie der Abelschen Functionen, Google source pp. 162--173.
% English translation mirrors the German current slice.

occurs, while the symmetric functions of the latter have already been regarded as functions of the \(v\), we may study it as a function of the \(v\) themselves. In this sense it is to bear the name of an Abelian transcendent, just as the function arising by differentiation with respect to \(\xi\),
\[
(6)\ldots\quad
\sum_{i=1}^{i=p}\int_{c^{(i)}}^{x^{(i)}}dZ_\xi
=T_\xi\begin{pmatrix}x^{(1)}x^{(2)}\cdots x^{(p)}\\ c^{(1)}c^{(2)}\cdots c^{(p)}\end{pmatrix}.
\]
When all arguments in them are of the same kind, we shall denote both transcendents more briefly by
\[
T_{\xi\eta}\begin{pmatrix}x\\c\end{pmatrix},\qquad
T_\xi\begin{pmatrix}x\\c\end{pmatrix}.
\]

As properties of the function \(T_{\xi\eta}\), we know first the one already belonging to the integral \(\Pi\) itself, namely that \(T_{\xi\eta}\) is the difference of two functions of the same kind, one of which depends only on \(\xi\), while the other depends only on \(\eta\). With regard to the limits, therefore, one can compose \(T_{\xi\eta}\) from four functions of the same kind, each with \(p+1\) arguments, in such a way that
\[
\begin{aligned}
T_{\xi\eta}={}&\varphi(x^{(1)},x^{(2)}\ldots x^{(p)},\xi)
-\varphi(x^{(1)},x^{(2)}\ldots x^{(p)},\eta)\\
&-\varphi(c^{(1)},c^{(2)}\ldots c^{(p)},\xi)
+\varphi(c^{(1)},c^{(2)}\ldots c^{(p)},\eta);
\end{aligned}
\]
a representation that will be treated more fully in the following section.

If, moreover, in equations (5), (6) we put for the paths of integration those which transform the integral sums \(v\) into the integral sums \(v'\) (3), then \(\Pi,Z\) change by the corresponding periodicity moduli; and if these are taken from the theorems of the preceding sections, one obtains the following determination of the periodicity moduli of \(T,T\):

If the paths of integration change in such a way that the integrals of the first kind respectively increase by
\[
2m_h\pi\sqrt{-1}+a_{1h}q_1+a_{2h}q_2+\cdots+a_{ph}q_p,
\]
then
\[
T_{\xi\eta}\begin{pmatrix}x\\c\end{pmatrix}\text{ increases by }
q_1\int_\eta^\xi du_1+q_2\int_\eta^\xi du_2+\cdots+q_p\int_\eta^\xi du_p,
\]
and
\[
T_\xi\begin{pmatrix}x\\c\end{pmatrix}\text{ increases by }
q_1\frac{\partial u_1(\xi)}{\partial\xi}+q_2\frac{\partial u_2(\xi)}{\partial\xi}+\cdots+q_p\frac{\partial u_p(\xi)}{\partial\xi}.
\]

\subsubsection*{\S\,42. Reduction of the inversion problem to the investigation of the transcendent \(T_{\xi\eta}\).}

The determination of the coefficients in equation (2) as functions of the \(v\), and thus the inversion problem generally, can be reduced to the investigation of the transcendent \(T\) as a function of the \(v\), by first expressing those coefficients by functions \(T\).

Namely, in the equation of Abel's theorem (Eq. I, p. 127), let \(\alpha\) be taken successively as \(x^{(1)},x^{(2)}\ldots x^{(p)}\), and at the same time let \(\beta\) be taken successively as \(c^{(1)},c^{(2)}\ldots c^{(p)}\). Further, in place of the curve \(\varphi=0\), let us put any curve \(\varphi-\lambda\psi=0\) of the pencil arising from \(\varphi=0\), \(\psi=0\); and let \(\eta\), as there, denote the intersection points of \(f=0\) with \(\psi=0\) which do not at the same time lie on \(\varphi=0\), while \(\xi\) denotes the intersection points of \(\varphi-\lambda\psi=0\) with \(f=0\) which do not lie on \(\varphi\) or \(\psi\). Then Abel's theorem gives us the following series of equations:
\[
\begin{aligned}
\log\!\biggl[\Bigl(\frac{\varphi}{\psi}\Bigr)_{x^{(1)}}-\lambda\biggr]
-\log\!\biggl[\Bigl(\frac{\varphi}{\psi}\Bigr)_{c^{(1)}}-\lambda\biggr]
&=-\sum\int_{c^{(1)}}^{x^{(1)}}d\Pi_{\xi\eta},\\
(7)\ldots\quad
\log\!\biggl[\Bigl(\frac{\varphi}{\psi}\Bigr)_{x^{(2)}}-\lambda\biggr]
-\log\!\biggl[\Bigl(\frac{\varphi}{\psi}\Bigr)_{c^{(2)}}-\lambda\biggr]
&=-\sum\int_{c^{(2)}}^{x^{(2)}}d\Pi_{\xi\eta},\\[-0.3em]
&\hspace{-1.0em}\cdots\cdots\cdots\\[-0.3em]
\log\!\biggl[\Bigl(\frac{\varphi}{\psi}\Bigr)_{x^{(p)}}-\lambda\biggr]
-\log\!\biggl[\Bigl(\frac{\varphi}{\psi}\Bigr)_{c^{(p)}}-\lambda\biggr]
&=-\sum\int_{c^{(p)}}^{x^{(p)}}d\Pi_{\xi\eta}.
\end{aligned}
\]
The pairs \(\xi\eta\) are assigned to one another as Abel's theorem prescribes at the place cited, and the sums on the right refer to the pairs of intersection points so formed. But the paths of integration are here likewise determined by Abel's theorem; they are therefore possibly not the same as those occurring in the integral sums \(v\), but would give integral sums \(v'\) differing from them by periodicity moduli. In truth, however, a change in the paths of integration of the \(x\) has no influence on the right-hand parts of the sums (7); for if, for example, \(x^{(1)}\) describes a period path, the right-hand side of the first equation (7) changes only by an expression of the form
\[
\sum_\eta\int_\eta^\xi dI,
\]
where \(I\) denotes an integral of the first kind, and the sum extends over all point pairs \(\xi,\eta\). But this sum vanishes by Abel's theorem, since \(\xi,\eta\) are intersection systems that can be transformed into one another. Thus this need not be taken into account at all, and, adding equations (7) and at the same time passing from logarithms to numbers, one obtains
\[
\prod_{i=1}^{i=p}
\frac{\bigl(\frac{\varphi}{\psi}\bigr)_{x^{(i)}}-\lambda}
{\bigl(\frac{\varphi}{\psi}\bigr)_{c^{(i)}}-\lambda}
=e^{-\sum T_{\xi\eta}\begin{smallmatrix}x^{(1)}x^{(2)}\cdots x^{(p)}\\c^{(1)}c^{(2)}\cdots c^{(p)}\end{smallmatrix}}.
\]
Put
\[
M=\biggl[\lambda-\Bigl(\frac{\varphi}{\psi}\Bigr)_{c^{(1)}}\biggr]
\biggl[\lambda-\Bigl(\frac{\varphi}{\psi}\Bigr)_{c^{(2)}}\biggr]
\cdots e^{-\sum T_{\xi\eta}\binom{x}{c}},
\]
so that, apart from constants, \(M\) contains only functions \(T\). Then, with regard to the coefficients denoted in \S\,41 (2) by \(M_1,M_2\ldots\), the equation can also be written in the form
\[
\lambda^p+M_1\lambda^{p-1}+\cdots+M_p=N.
\]
One need only form this equation \(p\) times in succession for as many different values of \(\lambda\), in order to have as many equations as unknowns \(M\); these can then be expressed through functions \(T\) by means of this system of linear equations. Thus one has the theorem:

\begin{quote}
The coefficients of the equation of degree \(p\), whose roots are the \(p\) functions \(\bigl(\frac{\varphi}{\psi}\bigr)_{x^{(i)}}\), become rational combinations of various transcendents \(e^T\).
\end{quote}

With respect to the reduction of the inversion problem to such an equation of degree \(p\), it remains to note that solving a single such equation is sufficient for expressing all possible functions \(\frac{\varphi}{\psi}\). For if we know the roots of the one equation
\[
\Bigl(\frac{\varphi}{\psi}\Bigr)^p+M_1\Bigl(\frac{\varphi}{\psi}\Bigr)^{p-1}+\cdots+M_p=0,
\]
then the determination of the corresponding \(p\) values of another function \(\frac{\Phi}{\Psi}\) can always be reduced to linear equations. Namely, if we form successively the equations of degree \(p\) whose roots are the \(p\) values of the functions
\[
\frac{\Phi}{\Psi},\quad \frac{\Phi}{\Psi}\frac{\varphi}{\psi},\quad
\frac{\Phi}{\Psi}\left(\frac{\varphi}{\psi}\right)^2,
\ldots,\quad
\frac{\Phi}{\Psi}\left(\frac{\varphi}{\psi}\right)^{p-1},
\]
and take from each of these equations the expression for the sums of the roots, then we have equations of the form
\[
\begin{aligned}
\Bigl(\frac{\Phi}{\Psi}\Bigr)_{x^{(1)}}+\Bigl(\frac{\Phi}{\Psi}\Bigr)_{x^{(2)}}+\cdots+\Bigl(\frac{\Phi}{\Psi}\Bigr)_{x^{(p)}}&=A,\\
\Bigl(\frac{\Phi}{\Psi}\Bigr)_{x^{(1)}}\Bigl(\frac{\varphi}{\psi}\Bigr)_{x^{(1)}}+\Bigl(\frac{\Phi}{\Psi}\Bigr)_{x^{(2)}}\Bigl(\frac{\varphi}{\psi}\Bigr)_{x^{(2)}}+\cdots+\Bigl(\frac{\Phi}{\Psi}\Bigr)_{x^{(p)}}\Bigl(\frac{\varphi}{\psi}\Bigr)_{x^{(p)}}&=A',\\
\Bigl(\frac{\Phi}{\Psi}\Bigr)_{x^{(1)}}\Bigl(\frac{\varphi}{\psi}\Bigr)_{x^{(1)}}^{2}+\Bigl(\frac{\Phi}{\Psi}\Bigr)_{x^{(2)}}\Bigl(\frac{\varphi}{\psi}\Bigr)_{x^{(2)}}^{2}+\cdots+\Bigl(\frac{\Phi}{\Psi}\Bigr)_{x^{(p)}}\Bigl(\frac{\varphi}{\psi}\Bigr)_{x^{(p)}}^{2}&=A'',\\[-0.3em]
\cdots\cdots\cdots&\\[-0.3em]
\sum_{i=1}^{p}\Bigl(\frac{\Phi}{\Psi}\Bigr)_{x^{(i)}}\Bigl(\frac{\varphi}{\psi}\Bigr)_{x^{(i)}}^{p-1}&=A^{(p-1)}.
\end{aligned}
\]
On the left stand rational combinations of various functions \(e^T\). Since the coefficients of the equations are assumed known, they suffice to determine the functions \(\bigl(\frac{\Phi}{\Psi}\bigr)_{x^{(i)}}\) linearly from them.

\subsubsection*{\S\,43. Extended inversion problem.}

The Jacobi inversion problem may be extended in the following way. Suppose that \(p\) sums, each of \(p+q\) integrals of the first kind of the same type, with upper limit points \(x^{(1)},x^{(2)}\ldots x^{(p+q)}\) and lower points \(c^{(1)},c^{(2)}\ldots c^{(p+q)}\), are set equal to given quantities \(v_1,v_2\ldots v_p\); and suppose further that \(q\) sums, each of \(p+q\) integrals of the third kind of the same type with the same limits, are set equal to given quantities \(w_1,w_2\ldots w_q\). Let the discontinuity points of the integrals of the third kind be \(\xi^{(1)},\eta^{(1)};\xi^{(2)},\eta^{(2)};\ldots;\xi^{(q)},\eta^{(q)}\). Then the equations of the extended inversion problem are
\[
(1)\ldots\quad
\left\{\begin{array}{rcl}
\displaystyle\sum_{h=1}^{h=p+q}\int_{c^{(h)}}^{x^{(h)}}du_1&=&v_1,\\
\displaystyle\sum_{h=1}^{h=p+q}\int_{c^{(h)}}^{x^{(h)}}du_2&=&v_2,\\[-0.2em]
\multicolumn{3}{c}{\cdots\cdots\cdots}\\[-0.2em]
\displaystyle\sum_{h=1}^{h=p+q}\int_{c^{(h)}}^{x^{(h)}}du_p&=&v_p,\\
\displaystyle\sum_{h=1}^{h=p+q}\int_{c^{(h)}}^{x^{(h)}}d\Pi_{\xi^{(1)}\eta^{(1)}}&=&w_1,\\
\displaystyle\sum_{h=1}^{h=p+q}\int_{c^{(h)}}^{x^{(h)}}d\Pi_{\xi^{(2)}\eta^{(2)}}&=&w_2,\\[-0.2em]
\multicolumn{3}{c}{\cdots\cdots\cdots}\\[-0.2em]
\displaystyle\sum_{h=1}^{h=p+q}\int_{c^{(h)}}^{x^{(h)}}d\Pi_{\xi^{(q)}\eta^{(q)}}&=&w_q.
\end{array}\right.
\]
The task here is to represent the value-pairs \(x,s_x\), or the coordinates of the upper limit points, as functions of the quantities \(v,w\), or equivalently to represent as such the coefficients of an equation
\[
\left(\frac{\varphi}{\psi}\right)^{p+q}+N_1\left(\frac{\varphi}{\psi}\right)^{p+q-1}+N_2\left(\frac{\varphi}{\psi}\right)^{p+q-2}\cdots+N_{p+q}=0,
\]
whose roots are the values, corresponding to the upper limits, of a homogeneous function \(\frac{\varphi}{\psi}\), whose numerator and denominator have equally high order.

This inversion problem can be reduced to the Jacobi problem with the aid of Abel's theorem. First determine, for arbitrarily chosen lower limit points \(c^{(p+q+1)},c^{(p+q+2)}\ldots c^{(2p+q)}\), upper limit points \(x^{(p+q+1)},x^{(p+q+2)}\ldots x^{(2p+q)}\) so that, for the corresponding integrals of the first kind, the equations hold:
\[
(2)\ldots\quad
\left\{\begin{array}{rcl}
\displaystyle\sum_{i=1}^{i=p}\int_{c^{(p+q+i)}}^{x^{(p+q+i)}}du_1&=&-v_1,\\
\displaystyle\sum_{i=1}^{i=p}\int_{c^{(p+q+i)}}^{x^{(p+q+i)}}du_2&=&-v_2,\\[-0.2em]
\multicolumn{3}{c}{\cdots\cdots\cdots}\\[-0.2em]
\displaystyle\sum_{i=1}^{i=p}\int_{c^{(p+q+i)}}^{x^{(p+q+i)}}du_p&=&-v_p.
\end{array}\right.
\]
Once this is done with the aid of Jacobi's inversion problem, the first \(p\) equations of system (1) become
\[
\sum_{i=1}^{i=2p+q}\int_{c^{(i)}}^{x^{(i)}}du_k=0,\qquad(k=1,2\ldots p),
\]
equations which, by Abel's theorem, assert that as soon as the lower limit points, together with fixed points of the curve \(f=0\), lie on a curve of order \(m\), the upper limit points also lie with the same points on such a curve. We therefore take such fixed points in sufficient number, and pass through them and the \(2p+q\) points \(c\) a curve \(\Psi=0\); through them and the already known points \(x^{(p+q+1)},x^{(p+q+2)}\ldots x^{(2p+q)}\) we pass a curve \(\Phi=0\), which is also to contain the sought points \(x^{(1)},x^{(2)}\ldots x^{(p+q)}\), and in which one may moreover assume as many coefficients equal to zero as can be destroyed by applying \(f=0\). Then, by adding to the remaining equations of system (1) on both sides the integrals corresponding to the points \(x^{(p+q+1)},x^{(p+q+2)}\ldots x^{(2p+q)}\), and then applying Abel's theorem on the left, we obtain
\[
\log\left(\frac{\Phi}{\Psi}\right)_{\eta^{(h)}}-
\log\left(\frac{\Phi}{\Psi}\right)_{\xi^{(h)}}
=w_h+T_{\xi^{(h)}\eta^{(h)}}\begin{pmatrix}x^{(p+q+1)}\cdots x^{(2p+q)}\\c^{(p+q+1)}\cdots c^{(2p+q)}\end{pmatrix},
\]
\[
h=1,2\ldots q.
\]
For brevity put
\[
(3)\ldots\quad
l_h=\frac{\Psi_{\eta^{(h)}}}{\Psi_{\xi^{(h)}}}
\exp\!\left(w_h+T_{\xi^{(h)}\eta^{(h)}}\begin{pmatrix}x^{(p+q+1)}\cdots\\c^{(p+q+1)}\cdots\end{pmatrix}\right).
\]
Thus \(l_1,l_2\ldots l_q\) are completely given quantities, and for the determination of the function \(\Phi\) one has the equations
\[
(4)\ldots\quad
\begin{array}{rcl}
\Phi_{\eta^{(1)}}-l_1\Phi_{\xi^{(1)}}&=&0,\\
\Phi_{\eta^{(2)}}-l_2\Phi_{\xi^{(2)}}&=&0,\\[-0.3em]
\multicolumn{3}{c}{\cdots\cdots\cdots}\\[-0.3em]
\Phi_{\eta^{(q)}}-l_q\Phi_{\xi^{(q)}}&=&0.
\end{array}
\]
These equations, which are linear in the coefficients of \(\Phi\), are wholly sufficient for determining them. For in \(\Phi\), first, as many coefficients had been made zero as could be made to disappear by adding \(f=0\), and then so many intersection points of \(\Phi=0\) with \(f=0\) had been given that only \(p+q\) of them remained. Since \(p\) intersection points are always determined by the rest, \(q\) equations must suffice to determine the function \(\Phi\) completely, which is achieved by equations (4).

In this way, therefore, we determine the points \(x^{(1)},x^{(2)}\ldots x^{(p+q)}\) sought, as the still unknown \(p+q\) intersection points of the curves \(f=0\), \(\Phi=0\), whose complete determination has been shown above.

It might indeed appear that, for the complete determination of the function \(\Phi\), the complete solution of the equation of degree \(p\) would also be required, on which the determination of the auxiliary points \(x^{(p+q+1)},x^{(p+q+2)}\ldots x^{(2p+q)}\) depends. But in reality it is already enough to know the symmetric functions of the coordinates of these points, which, as shown above, are representable as rational combinations of various functions \(e^{T_{\xi\eta}}\). Namely, if we form the equations
\[
\Phi_{x^{(p+q+1)}}=0,\quad
\Phi_{x^{(p+q+2)}}=0,\ldots,
\Phi_{x^{(2p+q)}}=0,
\]
which say that those points are to lie on the curve \(\Phi=0\), then in place of these equations we may put the \(p\) equations
\[
(5)\ldots\quad
\begin{array}{rcl}
\displaystyle\sum_{i=1}^{i=p}(\Theta_1\cdot\Phi)_{x^{(p+q+i)}}&=&0,\\
\displaystyle\sum_{i=1}^{i=p}(\Theta_2\cdot\Phi)_{x^{(p+q+i)}}&=&0,\\[-0.3em]
\multicolumn{3}{c}{\cdots\cdots\cdots}\\[-0.3em]
\displaystyle\sum_{i=1}^{i=p}(\Theta_p\cdot\Phi)_{x^{(p+q+i)}}&=&0,
\end{array}
\]
where \(\Theta_1,\Theta_2\ldots\Theta_p\) are independent fractional rational functions. If homogeneous coordinates are used, these are to be written so that numerator and denominator of the expressions \(\Theta\cdot\Phi\) become of equal order. These equations say exactly the same as the earlier ones, since their determinant does not vanish by assumption, and they contain only symmetric functions of the \(p\) points which above were expressed by the functions \(e^{T_{\xi\eta}}\). If equations (4), (5) are used to determine the function \(\Phi\), and if one perhaps adds the conditions which say that \(\Phi=0\) is to pass through the double and return points, then one has a sufficient number of equations to determine the coefficients of \(\Phi\) linearly. Moreover all coefficients of \(\Phi\) are thereby expressed rationally through the quantities \(e^w\) and through various transcendents \(e^{T_{\xi\eta}}\). Thus the extended inversion problem has the latter in common with the original one.

One can now also, as was done there, represent the values which a given function \(\frac{\varphi}{\psi}\) assumes at the sought points \(x^{(1)},x^{(2)}\ldots x^{(p)}\) as the roots of an equation whose coefficients become rational functions of the \(e^w\) and of various transcendents \(e^{T_{\xi\eta}}\). We need only write down the initial formulae used above for all \(p+q\) sought point pairs \(c,x\):
\begingroup\footnotesize
\[
\begin{aligned}
\log\!\biggl[\Bigl(\frac{\varphi}{\psi}\Bigr)_{x^{(1)}}-\lambda\biggr]
-\log\!\biggl[\Bigl(\frac{\varphi}{\psi}\Bigr)_{c^{(1)}}-\lambda\biggr]
&=-\sum\int_{c^{(1)}}^{x^{(1)}}d\Pi_{\xi\eta},\\
(6)\ldots\quad
\log\!\biggl[\Bigl(\frac{\varphi}{\psi}\Bigr)_{x^{(2)}}-\lambda\biggr]
-\log\!\biggl[\Bigl(\frac{\varphi}{\psi}\Bigr)_{c^{(2)}}-\lambda\biggr]
&=-\sum\int_{c^{(2)}}^{x^{(2)}}d\Pi_{\xi\eta},\\[-0.4em]
&\hspace{-1.0em}\cdots\cdots\cdots\\[-0.4em]
\log\!\biggl[\Bigl(\frac{\varphi}{\psi}\Bigr)_{x^{(p+q)}}-\lambda\biggr]
-\log\!\biggl[\Bigl(\frac{\varphi}{\psi}\Bigr)_{c^{(p+q)}}-\lambda\biggr]
&=-\sum\int_{c^{(p+q)}}^{x^{(p+q)}}d\Pi_{\xi\eta}.
\end{aligned}
\]
\endgroup
Here again the \(\xi\) denote the intersection points of \(f=0\) with \(\varphi-\lambda\psi=0\), and the \(\eta\) denote the intersection points of \(f=0\) with \(\psi=0\). Summing again, one obtains
\begingroup\footnotesize
\[
\begin{aligned}
\log\Biggl\{&
\frac{\bigl[(\frac{\varphi}{\psi})_{x^{(1)}}-\lambda\bigr]
\bigl[(\frac{\varphi}{\psi})_{x^{(2)}}-\lambda\bigr]\cdots}
{\bigl[(\frac{\varphi}{\psi})_{c^{(1)}}-\lambda\bigr]
\bigl[(\frac{\varphi}{\psi})_{c^{(2)}}-\lambda\bigr]\cdots}
\cdot
\frac{\bigl[(\frac{\varphi}{\psi})_{x^{(p+q)}}-\lambda\bigr]}
{\bigl[(\frac{\varphi}{\psi})_{c^{(p+q)}}-\lambda\bigr]}
\Biggr\}\\
&={} -\sum_{\xi\eta}\left\{\int_{c^{(1)}}^{x^{(1)}}d\Pi_{\xi\eta}+
\int_{c^{(2)}}^{x^{(2)}}d\Pi_{\xi\eta}+\cdots+
\int_{c^{(p+q)}}^{x^{(p+q)}}d\Pi_{\xi\eta}\right\}.
\end{aligned}
\]
\endgroup
The right-hand side of this equation, increased by the sum
\[
-T_{\xi\eta}\begin{pmatrix}x^{(p+q+1)}\cdots x^{(2p+q)}\\c^{(p+q+1)}\cdots c^{(2p+q)}\end{pmatrix},
\]
is, by Abel's theorem, equal to
\[
\log\left[\left(\frac{\Phi_\xi}{\Phi_\eta}\right)\cdot\left(\frac{\Psi_\eta}{\Psi_\xi}\right)\right],
\]
where \(\Phi,\Psi\) denote the expressions defined above, so that that side becomes equal to this logarithm increased by
\[
-T_{\xi\eta}\begin{pmatrix}x^{(p+q+1)}\cdots x^{(2p+q)}\\c^{(p+q+1)}\cdots c^{(2p+q)}\end{pmatrix}.
\]
Substituting this and passing from the logarithms to the numbers gives
\begingroup\footnotesize
\[
\begin{aligned}
(7)\ldots\quad
&\frac{\bigl[(\frac{\varphi}{\psi})_{x^{(1)}}-\lambda\bigr]
\bigl[(\frac{\varphi}{\psi})_{x^{(2)}}-\lambda\bigr]\cdots
\bigl[(\frac{\varphi}{\psi})_{x^{(p+q)}}-\lambda\bigr]}
{\bigl[(\frac{\varphi}{\psi})_{c^{(1)}}-\lambda\bigr]
\bigl[(\frac{\varphi}{\psi})_{c^{(2)}}-\lambda\bigr]\cdots
\bigl[(\frac{\varphi}{\psi})_{c^{(p+q)}}-\lambda\bigr]}\\
&\qquad=\frac{\Phi_{\xi^{(1)}}\Psi_{\eta^{(1)}}}{\Phi_{\eta^{(1)}}\Psi_{\xi^{(1)}}}\cdot
\frac{\Phi_{\xi^{(2)}}\Psi_{\eta^{(2)}}}{\Phi_{\eta^{(2)}}\Psi_{\xi^{(2)}}}\cdots
 e^{-\sum_{\xi\eta}T_{\xi\eta}\begin{smallmatrix}x^{(p+q+1)}\cdots\\c^{(p+q+1)}\cdots\end{smallmatrix}}.
\end{aligned}
\]
\endgroup
Putting in (7), for \(\lambda\), once more successively \(\lambda_1,\lambda_2\ldots\lambda_{p+q}\), we have linear equations for forming the symmetric functions of the quantities \(\bigl(\frac{\varphi}{\psi}\bigr)_x\), and hence for representing the coefficients of the equation whose roots are these quantities. These expressions too are, as one sees, rational functions of the \(e^w\) and of various transcendents \(e^{T_{\xi\eta}}\), as was the case with the earlier problems.

The remark of \S\,42, that the knowledge of one system of functions \(\frac{\varphi}{\psi}\) suffices to determine all others linearly from it, naturally remains completely valid here.

\subsubsection*{\S\,44. Periods of the functions of the extended inversion problem. Modification when integrals of the second kind occur.}

If we also call these symmetric functions Abelian functions, we see that they likewise are periodic functions of their arguments \(v_1,v_2\ldots v_p,w_1,w_2\ldots w_q\). The most general system of values that can arise from a given system \(v_1,v_2\ldots v_p,w_1,w_2\ldots w_q\) by changing the paths of integration is the following:
\[
\begin{array}{rcl}
v_1'&=&v_1+2m_1\pi\sqrt{-1}+\displaystyle\sum_{h=1}^{h=p}a_{1h}q_h,\\
v_2'&=&v_2+2m_2\pi\sqrt{-1}+\displaystyle\sum_{h=1}^{h=p}a_{2h}q_h,\\[-0.3em]
\multicolumn{3}{c}{\cdots\cdots\cdots}\\[-0.3em]
v_p'&=&v_p+2m_p\pi\sqrt{-1}+\displaystyle\sum_{h=1}^{h=p}a_{ph}q_h,\\
w_1'&=&w_1+2n_1\pi\sqrt{-1}+\displaystyle\sum_{h=1}^{h=p}q_h\int_{\eta^{(1)}}^{\xi^{(1)}}du_h,\\
w_2'&=&w_2+2n_2\pi\sqrt{-1}+\displaystyle\sum_{h=1}^{h=p}q_h\int_{\eta^{(2)}}^{\xi^{(2)}}du_h,\\[-0.3em]
\multicolumn{3}{c}{\cdots\cdots\cdots}\\[-0.3em]
w_q'&=&w_q+2n_q\pi\sqrt{-1}+\displaystyle\sum_{h=1}^{h=p}q_h\int_{\eta^{(q)}}^{\xi^{(q)}}du_h.
\end{array}
\]
Thus these functions are \(2p+q\)-fold periodic; here, in particular, \(p+q\) periods \(2\pi\sqrt{-1}\) appear, belonging independently to the different arguments.

We can also form transcendents corresponding to this extended inversion problem. But, as shown above, a sum of \(p+q\) integrals of the third kind of the same type differs from the negative of the transcendent
\[
T_{\xi\eta}\begin{pmatrix}x^{(p+q+1)}\cdots\\c^{(p+q+1)}\cdots\end{pmatrix}
\]
only by the logarithm of a function which, as shown above, is composed rationally from the \(e^w\) and the various expressions \(e^{T_{\xi\eta}}\). Namely, let \(T^{(q)}_{\xi\eta}\) denote the transcendent of the extended inversion problem
\[
T^{(q)}_{\xi\eta}\begin{pmatrix}x^{(1)},x^{(2)}\cdots x^{(p+q)}\\c^{(1)},c^{(2)}\cdots c^{(p+q)}\end{pmatrix}
=\sum_{i=1}^{i=p+q}\int_{c^{(i)}}^{x^{(i)}}d\Pi_{\xi\eta}.
\]
If the paths of integration agree with those of equations (10), then with the choice already made above for the paths of integration in \(T\) we have
\begingroup\footnotesize
\[
T^{(q)}_{\xi\eta}\begin{pmatrix}x^{(1)},x^{(2)}\cdots x^{(p+q)}\\c^{(1)},c^{(2)}\cdots c^{(p+q)}\end{pmatrix}
=-T_{\xi\eta}\begin{pmatrix}x^{(p+q+1)}\cdots x^{(2p+q)}\\c^{(p+q+1)}\cdots c^{(2p+q)}\end{pmatrix}
-\log\frac{\Phi_\xi\Psi_\eta}{\Phi_\eta\Psi_\xi}.
\]
\endgroup
A second treatment of the extended inversion problem is attached to this transcendent, in which the symmetric functions of the values of \(\frac{\varphi}{\psi}\) are reduced to the transcendent \(T_{\xi\eta}^{(q)}\) in exactly the same way that, above, in the Jacobi inversion problem, the corresponding functions were reduced to \(T_{\xi\eta}\).

By differentiating the preceding equation for \(T_{\xi\eta}^{(q)}\), one obtains another equation which reduces the transcendent
\[
T_{\xi}^{(q)}\begin{pmatrix}x^{(1)},x^{(2)}\cdots x^{(p+q)}\\c^{(1)},c^{(2)}\cdots c^{(p+q)}\end{pmatrix}
=\sum_{i=1}^{i=p+q}\int_{c^{(i)}}^{x^{(i)}}dZ_\xi
\]
to \(T_\xi\). For with the same choice of paths of integration in \(T_\xi\) one has
\[
T_\xi^{(q)}\begin{pmatrix}x^{(1)},x^{(2)}\cdots x^{(p+q)}\\c^{(1)},c^{(2)}\cdots c^{(p+q)}\end{pmatrix}
=-T_\xi\begin{pmatrix}x^{(p+q+1)}\cdots x^{(2p+q)}\\c^{(p+q+1)}\cdots c^{(2p+q)}\end{pmatrix}
-\frac{\partial}{\partial\xi}\log\frac{\Phi_\xi}{\Psi_\xi}.
\]
It is further important to note that, among the integrals of the third kind given in the extended inversion problem, some can pass into integrals of the second kind by coalescence of the corresponding discontinuity points \(\xi,\eta\). For those quantities \(w\) for which this occurs, the logarithmic period \(2n\pi\sqrt{-1}\) arising from the integrals of the third kind disappears; hence, if \(r\) sums of integrals of the third kind pass into sums of integrals of the second kind, the corresponding Abelian functions are only \((2p+q-r)\)-fold periodic. In the treatment of the inversion problem, however, a modification would then occur.

One of the equations (1) belonging here is no longer
\[
\sum_{h=1}^{h=p+q}\int_{c^{(h)}}^{x^{(h)}}d\Pi_{\xi^{(i)}\eta^{(i)}}=w_i,
\]
but rather
\[
\sum_{h=1}^{h=p+q}\int_{c^{(h)}}^{x^{(h)}}dZ_{\xi^{(i)}}=t_i.
\]
Thus, adding on the right and left the integrals corresponding to the \(p\) auxiliary quantities \(x\), we find
\[
\sum_{h=1}^{h=2p+q}\int_{c^{(h)}}^{x^{(h)}}dZ_{\xi^{(i)}}=t_i+T_{\xi^{(i)}}\begin{pmatrix}x^{(p+q+1)}\cdots x^{(2p+q)}\\c^{(p+q+1)}\cdots c^{(2p+q)}\end{pmatrix},
\]
and therefore, by applying Abel's theorem,
\[
-\frac{\partial}{\partial\xi^{(i)}}\log\frac{\Phi_{\xi^{(i)}}}{\Psi_{\xi^{(i)}}}
=t_i+T_{\xi^{(i)}}\begin{pmatrix}x^{(p+q+1)}\cdots x^{(2p+q)}\\c^{(p+q+1)}\cdots c^{(2p+q)}\end{pmatrix},
\]
or
\begingroup\footnotesize
\[
\Phi_{\xi^{(i)}}\frac{\partial\Psi_{\xi^{(i)}}}{\partial\xi^{(i)}}
-\Psi_{\xi^{(i)}}\frac{\partial\Phi_{\xi^{(i)}}}{\partial\xi^{(i)}}
=\Phi_{\xi^{(i)}}\Psi_{\xi^{(i)}}
\left[t_i+T_{\xi^{(i)}}\begin{pmatrix}x^{(p+q+1)}\cdots x^{(2p+q)}\\c^{(p+q+1)}\cdots c^{(2p+q)}\end{pmatrix}\right].
\]
\endgroup
This too is a linear equation for the coefficients of \(\Phi\), and it must be put in place of the corresponding equation (4). The coefficients of \(\Phi\) therefore contain the \(t\) rationally, and besides various functions \(e^{T_{\xi\eta}}\), also functions \(T_\xi=\frac{\partial T_{\xi\eta}}{\partial\xi}\), which correspond to the integral sums of the second kind.

In fact, however, this is the whole change. For if one otherwise retains the treatment of the extended inversion problem, only a change in the expressions of the \(T_{\xi\eta}\) results, insofar as these depend on \(\Phi\) and are therefore modified by the altered manner of forming this function.

Of special interest is the case of the extended inversion problem in which the integrals of the third and second kind are such as become discontinuous in the double and return points of \(f=0\).

% --- continuation source pp. 174--181 ---

% Source: Clebsch--Gordan, Theorie der Abelschen Functionen, Google source pp. 174--181.
% English translation mirrors the German current slice.

\begin{center}
{\large\bfseries Seventh Section.}\par\vspace{0.8em}
{\Large\bfseries The function \(T_{\xi\eta}\).}
\end{center}

\subsubsection*{\S\,45. Decomposition of the function \(T_{\xi\eta}\) into the difference of two functions, each of which depends only on \(p+2\) points.}

The transcendent
\[
T_{\xi\eta}\begin{pmatrix}x^{(1)},x^{(2)}\ldots x^{(p)}\\ \alpha^{(1)},\alpha^{(2)}\ldots \alpha^{(p)}\end{pmatrix}
=\int_{\alpha^{(1)}}^{x^{(1)}}d\Pi_{\xi\eta}
+\int_{\alpha^{(2)}}^{x^{(2)}}d\Pi_{\xi\eta}
+\cdots+\int_{\alpha^{(p)}}^{x^{(p)}}d\Pi_{\xi\eta}
\]
still depends in all on \(2p+2\) points, namely on \(\xi\), on \(\eta\), and on the \(2p\) limits of the integrals used. But we have already observed above (\S\,41) that this transcendent can be composed from four functions, each of which depends only on \(p+1\) arguments, in such a way that
\[
\begin{aligned}
T_{\xi\eta}\binom{x}{\alpha}={}&\varphi(x^{(1)},x^{(2)}\ldots;\xi)
-\varphi(x^{(1)},x^{(2)}\ldots;\eta)\\
&-\varphi(\alpha^{(1)},\alpha^{(2)}\ldots;\xi)
+\varphi(\alpha^{(1)},\alpha^{(2)}\ldots;\eta).
\end{aligned}
\]
The function \(\varphi\) is by no means determined; rather \(\varphi(x^{(1)},x^{(2)}\ldots;\xi)\) may be changed by an arbitrary function of the \(x\) and also by an arbitrary function of \(\xi\). If the other functions then undergo the corresponding changes, the expression for \(T_{\xi\eta}\) remains wholly unchanged.

A division of \(T_{\xi\eta}\) can be brought about very easily by decomposing \(T_{\xi\eta}\) into the difference of two integral sums with arbitrarily chosen lower limits, and then, by introducing a new constant point \(\zeta\) and using the formula
\[
d\Pi_{\xi\eta}=d\Pi_{\xi\zeta}-d\Pi_{\eta\zeta},
\]
representing each of these sums again as a difference. But the progress made in this way is only apparent, since \(p+1\) new quantities have been introduced and hence each of the four parts again depends on \(2p+2\) quantities. Therefore, if one wishes to decompose the function \(T_{\xi\eta}\) genuinely into functions depending on a smaller number of quantities, one must take another path.

We first observe the following. If one lays down two curves of order \((n-2)\) which pass through the remaining points of \(f=0\) lying on the line \(\xi,\eta\), and through the double and return points of \(f=0\), and of which one moreover passes through the points \(x\) and the other through the points \(\alpha\), then both curves are exactly determined, and each of them cuts \(f=0\) in \(p\) further points. Let the points lying with the \(x\) on one curve be \(y^{(1)},y^{(2)}\ldots y^{(p)}\), and those lying with the \(\alpha\) on the other curve be \(\beta^{(1)},\beta^{(2)}\ldots\beta^{(p)}\). Then the \(y\) depend only on the \(x\) and on \(\xi,\eta\), while the \(\beta\) depend only on the \(\alpha\) and \(\xi,\eta\). But the two curves stand in the special relation to one another that, together with \(f=0\), they intersect in \(n-2\) points lying on a straight line. Thus the case of the theorem (\S\,13.6) occurs here, according to which, in forming Abel's theorem, the sum of the corresponding integrals of the third kind does not become equal to a logarithmic function, but vanishes. Hence one has
\[
\sum_{h=1}^{h=p}\int_{\alpha^{(h)}}^{x^{(h)}}d\Pi_{\xi\eta}
+\sum_{h=1}^{h=p}\int_{\beta^{(h)}}^{y^{(h)}}d\Pi_{\xi\eta}=0;
\]
here the paths of integration of the second sum are to be chosen so that, when in the first sum the \(x\) are brought near to the \(\alpha\) and the paths of integration are made to vanish, the \(y\) also approach the \(\beta\) so that the paths of integration vanish; and, of course, one and the same path from \(\xi\) to \(\eta\) is not cut by any of the paths of integration.

This equation can also be clothed in the form
\[
2\sum_{h=1}^{h=p}\int_{\alpha^{(h)}}^{x^{(h)}}d\Pi_{\xi\eta}
=\sum_{h=1}^{h=p}\int_{y^{(h)}}^{x^{(h)}}d\Pi_{\xi\eta}
-\sum_{h=1}^{h=p}\int_{\beta^{(h)}}^{\alpha^{(h)}}d\Pi_{\xi\eta};
\]
here too the paths of integration must be taken so that this formula agrees with the preceding one. Thus the path of integration from \(y^{(h)}\) to \(x^{(h)}\) must be composed of the earlier path laid from \(\alpha^{(h)}\) to \(x^{(h)}\) and a path laid from \(y^{(h)}\) to \(\alpha^{(h)}\), while the latter, together with the earlier path from \(\beta^{(h)}\) to \(y^{(h)}\), gives the new path from \(\beta^{(h)}\) to \(\alpha^{(h)}\). If we now introduce the functions \(T\) and note in each of them the upper and lower limit values that occur here, we have
\[
\begin{aligned}
(1)\ldots\quad
T_{\xi\eta}\begin{pmatrix}x^{(1)},x^{(2)}\ldots x^{(p)}\\ \alpha^{(1)},\alpha^{(2)}\ldots \alpha^{(p)}\end{pmatrix}
&=\frac12 T_{\xi\eta}\begin{pmatrix}x^{(1)},x^{(2)}\ldots x^{(p)}\\ y^{(1)},y^{(2)}\ldots y^{(p)}\end{pmatrix}\\
&\quad-\frac12 T_{\xi\eta}\begin{pmatrix}\alpha^{(1)},\alpha^{(2)}\ldots\alpha^{(p)}\\ \beta^{(1)},\beta^{(2)}\ldots\beta^{(p)}\end{pmatrix};
\end{aligned}
\]
an equation which, since the lower limits on the right are determined by the upper ones, we shall for simplicity write in the form
\[
T_{\xi\eta}\binom{x}{\alpha}=\frac12T_{\xi\eta}(x)-\frac12T_{\xi\eta}(\alpha),
\]
in which the lower limits on the right are not indicated separately. We thus see the general function \(T_{\xi\eta}\), without the introduction of new constants, transformed into the half-difference of two other functions of the same kind, one containing no trace of the \(\alpha\), the other no trace of the \(x\); hence each depends only on \(p+2\) quantities. Henceforth, therefore, we shall have to investigate the simpler function
\[
\frac12T_{\xi\eta}\begin{pmatrix}x^{(1)},x^{(2)}\ldots x^{(p)}\\ y^{(1)},y^{(2)}\ldots y^{(p)}\end{pmatrix}=\frac12T_{\xi\eta}(x).
\]

The function \(\frac12T_{\xi\eta}(x)\) is the difference of two functions of the same kind, one of which contains \(\xi\) but not \(\eta\), while the other differs from the former by having \(\eta\) in place of \(\xi\). This property is concealed in the above expression, since the lower limits \(y\) still depend on \(\xi,\eta\); it is therefore necessary to prove it.

\subsubsection*{\S\,46. Proof of the decomposability of the function \(\frac12T_{\xi\eta}(x)\) into two functions depending only on \(p+1\) points.}

For this purpose we note that, by formula (1), our expression can be represented by
\[
\sum\int_{\alpha^{(h)}}^{x^{(h)}}d\Pi_{\xi\eta}
-\frac12\sum\int_{\beta^{(h)}}^{\alpha^{(h)}}d\Pi_{\xi\eta},
\]
where the \(\alpha\) denote any \(p\) constant points. We already know that the first part of this expression has the required property, whatever constant points \(\alpha\) we choose; therefore only the second expression remains to be considered, in which the \(\beta\) still depend on \(\xi,\eta\). Let the \(p\) constant points \(\alpha\) now approach a position in which, together with the double and return points of \(f=0\), they lie on a curve of order \((n-3)\). Then the curve of order \((n-2)\), on which the \(\alpha\) and \(\beta\), together with those exceptional points and the \(n-2\) remaining points of the line \(\xi\eta\), are to lie, passes into the union of this curve of order \((n-3)\) and the line \(\xi\eta\) itself. Of the points \(\beta\), two therefore fall in \(\xi,\eta\) themselves, while the remaining \(p-2\) fall in the other intersection points of the curve of order \((n-3)\) with the curve \(f=0\). Of the \(p\) terms of the sum
\[
\sum\int_{\beta^{(h)}}^{\alpha^{(h)}}d\Pi_{\xi\eta}
\]
therefore \(p-2\) acquire constant limits, and consequently likewise possess the property of decomposing into the difference of two functions, one of which no longer depends on \(\eta\), the other no longer on \(\xi\). Thus only the limit remains to be investigated to which the two remaining integrals tend; their lower limits we temporarily assume to differ little from \(\xi,\eta\), and accordingly denote them by \(\xi+\delta x,\eta+\delta y\). The integral sum arising from this passage to the limit requires a more precise investigation, because in it the lower limits coincide with discontinuity points. Instead of the integrals
\[
\int_{\xi+\delta x}^{\alpha^{(1)}}d\Pi_{\xi\eta}
+\int_{\eta+\delta y}^{\alpha^{(2)}}d\Pi_{\xi\eta},
\]
we can first, by an often used theorem, put
\[
\begin{aligned}
&\int_{\xi+\delta x}^{\alpha^{(1)}}d\Pi_{\xi\zeta}
-\int_{\xi+\delta x}^{\alpha^{(1)}}d\Pi_{\eta\zeta}
+\int_{\eta+\delta y}^{\alpha^{(2)}}d\Pi_{\xi\zeta}
-\int_{\eta+\delta y}^{\alpha^{(2)}}d\Pi_{\eta\zeta}.
\end{aligned}
\]
In the second and third of these integrals the passage to the limit can be carried out immediately, and one obtains
\[
\int_{\eta}^{\alpha^{(2)}}d\Pi_{\xi\zeta}-\int_{\xi}^{\alpha^{(1)}}d\Pi_{\eta\zeta}.
\]
In this difference, by the theorem on interchange of parameter and argument, the parts depending simultaneously on \(\xi,\eta\) cancel; for one has at once
\[
\begin{aligned}
\int_{\eta}^{\alpha^{(2)}}d\Pi_{\xi\zeta}-\int_{\xi}^{\alpha^{(1)}}d\Pi_{\eta\zeta}
={}&\int_{\eta}^{\delta}d\Pi_{\xi\varepsilon}
+\int_{\eta}^{\delta}d\Pi_{\varepsilon\zeta}
+\int_{\delta}^{\alpha^{(2)}}d\Pi_{\xi\zeta}\\
&-\int_{\xi}^{\varepsilon}d\Pi_{\eta\delta}
-\int_{\xi}^{\delta}d\Pi_{\delta\zeta}
-\int_{\varepsilon}^{\alpha^{(1)}}d\Pi_{\eta\zeta},
\end{aligned}
\]
where indeed the first positive term cancels against the first negative one. Thus finally the expression remains to be considered:
\[
\lim\left\{\int_{\xi+\delta x}^{\alpha^{(1)}}d\Pi_{\xi\zeta}
-\int_{\eta+\delta y}^{\alpha^{(2)}}d\Pi_{\eta\zeta}\right\}.
\]
If under the first integral sign one adds \(d\log\Sigma\pm x_1\xi_2\zeta_3\), then the integral
\[
\int_{\xi+\delta x}^{\alpha^{(1)}}d\bigl(\Pi_{\xi\zeta}(x)+\log\Sigma\pm x_1\xi_2\zeta_3\bigr)
\]
arises. It is no longer discontinuous at the point \(\xi\), and hence can be transformed directly into the integral from \(\xi\) to \(\alpha^{(1)}\), which depends only on \(\xi\). In the same way the last of the above integrals, by adding \(d\log\Sigma\pm y_1\eta_2\zeta_3\), is transformed into
\[
\int_{\eta+\delta y}^{\alpha^{(2)}}d\bigl(\Pi_{\eta\zeta}(y)+\log\Sigma\pm y_1\eta_2\zeta_3\bigr),
\]
which is continuous at \(\eta\), hence can be replaced by the integral from \(\eta\) to \(\alpha^{(2)}\), and therefore depends only on \(\eta\).

Thus the consideration of our expression has been reduced to the investigation of the sum of the integrals added above, which must again be subtracted and hence give the terms
\[
-\int_{\alpha^{(1)}}^{\xi+\delta x}d\log\Sigma\pm x_1\xi_2\zeta_3
+\int_{\alpha^{(2)}}^{\eta+\delta y}d\log\Sigma\pm y_1\eta_2\zeta_3
\]
\[
=-\log\left(\frac{\Sigma\pm\delta x_1\xi_2\zeta_3}{\Sigma\pm\alpha_1^{(1)}\xi_2\zeta_3}\cdot
\frac{\Sigma\pm\alpha_1^{(2)}\eta_2\zeta_3}{\Sigma\pm\delta y_1\eta_2\zeta_3}\right).
\]
Here, however, we need attend only to the expression
\[
\log\frac{\Sigma\pm\delta x_1\xi_2\zeta_3}{\Sigma\pm\delta y_1\eta_2\zeta_3},
\]
since the other parts already possess the required property.

The differentials \(\delta x,\delta y\) are bound to one another by the condition that, by hypothesis, the points \(\xi+\delta x,\eta+\delta y\) always have to lie, with the \(n-2\) other points of the line \(\xi\eta\), the double and return points, and \(p\) constant points \(\alpha\), on a curve of order \((n-2)\). The condition thereby expressed is to be formed exactly as the function \(\Omega\) of order \((n-2)\) that stands in the numerator under the integral sign of the integral of the third kind, and whose formation we taught in detail in the first section. We obtain a determinant of \(\frac{n(n-1)}{2}\) rows, which is the result of an elimination. Let \(w=0\) denote the equation of a curve of order \((n-2)\), and suppose that, when \(\xi_1+\lambda\eta_1,\xi_2+\lambda\eta_2,\xi_3+\lambda\eta_3\) are substituted for the variables in it,
\[
w(\xi+\lambda\eta)=v+\lambda v_1+\lambda^2v_2+\cdots+\lambda^{n-2}v_{n-2};
\]
and suppose further that
\[
\frac{f(\xi+\lambda\eta)}{\lambda}=u_1+\lambda u_2+\lambda^2u_3+\cdots+\lambda^{n-2}u_{n-1}.
\]
Then first write the equation \(w=0\) with the variables \(\xi+\delta x\), then with \(\eta+\delta y\), and further with the constant arguments corresponding to the double and return points and to the points \(\alpha\). To these equations add the system
\[
\begin{array}{rcl}
v&=&\mu u_1,\\
v_1&=&\mu u_2,\\[-0.3em]
\multicolumn{3}{c}{\cdots\cdots\cdots}\\[-0.3em]
v_{n-2}&=&\mu u_{n-1}.
\end{array}
\]
Eliminate the coefficients of \(w\), which occur linearly also in the \(v\), and \(\mu\); by \S\,6 the result is the desired equation.

But in place of the equations
\[
w(\xi+\delta x)=0,\qquad w(\eta+\delta y)=0
\]
one can put the equations
\[
\begin{aligned}
w(\xi+\delta x)-v&=-\mu u_1
=\delta x_1\frac{\partial w(\xi)}{\partial\xi_1}
+\delta x_2\frac{\partial w(\xi)}{\partial\xi_2}
+\delta x_3\frac{\partial w(\xi)}{\partial\xi_3},\\
w(\eta+\delta y)-v_{n-2}&=-\mu u_{n-1}
=\delta y_1\frac{\partial w(\eta)}{\partial\eta_1}
+\delta y_2\frac{\partial w(\eta)}{\partial\eta_2}
+\delta y_3\frac{\partial w(\eta)}{\partial\eta_3}.
\end{aligned}
\]
When these equations take the place of the first two, the first rows of the elimination determinant contain only infinitely small terms, except for the last terms, which respectively are \(-u_1\) and \(-u_{n-1}\). Thus, in order to obtain only terms of the first order, we must restrict the determinant to the terms carrying these elements \(u_1\) and \(u_{n-1}\). The developed determinant therefore has the form
\[
Au_{n-1}-Bu_1=0,
\]
where \(A\) and \(B\) are determinants that differ only by one row. Namely \(A=0\) arises from eliminating the coefficients of \(w\) from the equations
\[
\begin{array}{c}
w(\xi+\delta x)=0,\\
w(\alpha^{(1)})=0,\\
w(\alpha^{(2)})=0,\\[-0.3em]
\cdots\cdots\cdots\\[-0.3em]
v=0,\\
v_1=0,\\[-0.3em]
\cdots\cdots\cdots\\[-0.3em]
v_{n-2}=0,
\end{array}
\]
whereas for \(B\) only the first equation is to be replaced by
\[
w(\eta+\delta y)=0.
\]
Therefore \(A=0\), when the quantities \(\xi+\delta x\) are regarded as the variables, is a curve of order \((n-2)\) which is identically satisfied by the coordinates \(\xi_i+\lambda\eta_i\), for which
\[
w(\xi+\lambda\eta)=v+\lambda v_1+\cdots+\lambda^{n-2}v_{n-2},
\]
and which therefore contains the line \(\xi\eta\) as a factor; consequently it can consist besides only of the curve of order \((n-3)\) determined by the points \(\alpha\) and by the double and return points. If the equation of the latter is \(\Phi(\xi+\delta x)=0\), then therefore
\[
A=M\cdot\Phi(\xi+\delta x)\cdot\Sigma\pm(\xi+\delta x)_1\xi_2\eta_3,
\]
where \(M\) no longer depends on \(\xi+\delta x\). Hence also
\[
B=M\cdot\Phi(\eta+\delta y)\cdot\Sigma\pm(\eta_1+\delta y_1)\xi_2\eta_3.
\]
Omitting the factor \(M\), and replacing \(u_{n-1},u_1\) by their values, our condition equation therefore becomes
\[
\begin{aligned}
&\Phi(\xi+\delta x)\,[\Sigma\pm(\xi_1+\delta x_1)\xi_2\eta_3]
\left(\eta_1\frac{\partial f(\xi)}{\partial\xi_1}+\eta_2\frac{\partial f(\xi)}{\partial\xi_2}+\eta_3\frac{\partial f(\xi)}{\partial\xi_3}\right)\\
&=\Phi(\eta+\delta y)\,[\Sigma\pm(\eta_1+\delta y_1)\xi_2\eta_3]
\left(\xi_1\frac{\partial f(\eta)}{\partial\eta_1}+\xi_2\frac{\partial f(\eta)}{\partial\eta_2}+\xi_3\frac{\partial f(\eta)}{\partial\eta_3}\right),
\end{aligned}
\]
or, if we now pass to the limit,
\[
\begin{aligned}
&\Phi(\xi)\Sigma\pm\delta x_1\xi_2\eta_3
\left(\eta_1\frac{\partial f(\xi)}{\partial\xi_1}+\eta_2\frac{\partial f(\xi)}{\partial\xi_2}+\eta_3\frac{\partial f(\xi)}{\partial\xi_3}\right)\\
&=\Phi(\eta)\Sigma\pm\delta y_1\xi_2\eta_3
\left(\xi_1\frac{\partial f(\eta)}{\partial\eta_1}+\xi_2\frac{\partial f(\eta)}{\partial\eta_2}+\xi_3\frac{\partial f(\eta)}{\partial\eta_3}\right).
\end{aligned}
\]
Since \(\xi+\delta x\) must lie on the curve \(f=0\), we have
\[
\frac{\partial f(\xi)}{\partial\xi_1}\delta x_1
+\frac{\partial f(\xi)}{\partial\xi_2}\delta x_2
+\frac{\partial f(\xi)}{\partial\xi_3}\delta x_3=0;
\]
and the identity
\[
\left|
\begin{array}{cccc}
\delta x_1&\delta x_2&\delta x_3&\displaystyle\sum\frac{\partial f(\xi)}{\partial\xi_i}\delta x_i\\
\xi_1&\xi_2&\xi_3&\displaystyle\sum\frac{\partial f(\xi)}{\partial\xi_i}\xi_i\\
\eta_1&\eta_2&\eta_3&\displaystyle\sum\frac{\partial f(\xi)}{\partial\xi_i}\eta_i\\
\zeta_1&\zeta_2&\zeta_3&\displaystyle\sum\frac{\partial f(\xi)}{\partial\xi_i}\zeta_i
\end{array}
\right|=0
\]
therefore gives
\[
\Sigma\pm\delta x_1\xi_2\eta_3=
\Sigma\pm\delta x_1\xi_2\zeta_3\cdot
\frac{\displaystyle\sum\frac{\partial f(\xi)}{\partial\xi_i}\eta_i}
{\displaystyle\sum\frac{\partial f(\xi)}{\partial\xi_i}\zeta_i}.
\]
Likewise one has
\[
\Sigma\pm\delta y_1\xi_2\eta_3=
-\Sigma\pm\delta y_1\eta_2\zeta_3\cdot
\frac{\displaystyle\sum\frac{\partial f(\eta)}{\partial\eta_i}\xi_i}
{\displaystyle\sum\frac{\partial f(\eta)}{\partial\eta_i}\zeta_i}.
\]

% --- continuation source pp. 182--193 ---

If one substitutes these values in the above conditional equation and removes superfluous factors, one has
\[
0=\frac{\Phi'\xi}{\sum_i \frac{\partial f(\xi)}{\partial \xi_i}\zeta_i}
\sum\pm\delta x_1\xi_2\zeta_3
+\frac{\Phi'\eta}{\sum_i \frac{\partial f(\eta)}{\partial \eta_i}\zeta_i}
\sum\pm\delta y_1\eta_2\zeta_3 .
\]
The logarithm to be investigated therefore becomes
\[
\log\left(-\frac{\Phi'\eta}{\Phi'\xi}\cdot
\frac{\sum_i \frac{\partial f(\xi)}{\partial \xi_i}\zeta_i}
{\sum_i \frac{\partial f(\eta)}{\partial \eta_i}\zeta_i}\right),
\]
which is manifestly the difference of a function of \(\xi\) and a function of \(\eta\), as was required above.

Thus we have shown that our transcendent, depending on only \(p+2\) quantities,
\[
\frac12 T_{\xi\eta}\begin{pmatrix}x^{(1)}&x^{(2)}&\cdots\\ y^{(1)}&y^{(2)}&\cdots\end{pmatrix}
\]
still possesses the essential property that it separates into a part depending only on \(\eta\) and a part depending only on \(\xi\), as belonged to the original transcendent \(T_{\xi\eta}\) with \(2p+2\) arguments.

It follows from this that, if one forms three different transcendents \(T_{\xi\eta}\), \(T_{\eta\zeta}\), \(T_{\zeta\xi}\), whose upper limits are always the same points \(x\), while their lower limits arise by the construction given above and may respectively be denoted by \(y,z,t\), the sum of all three vanishes for suitably placed paths; hence
\[
T_{\xi\eta}(x)+T_{\eta\zeta}(x)+T_{\zeta\xi}(x)=0.
\]

\subsubsection*{\S\,47. On the character of the special transcendent \(T_{\xi\eta}(x)\).}

In order to make fully evident the character and the manner of determination of the newly introduced transcendent, we add the following explanations.

We begin with the \(p\) fixed points \(c\), which are the lower limit-points of the inversion problem and which are only not to lie with the double and return points on a curve of order \((n-3)\).\footnote{This assumption is made here for brevity. It presents no essential difficulty also to treat the opposite one.} Let these points be the upper limits of such a transcendent, and let two fixed points \(\lambda,\mu\) be its parameters. They determine \(p\) other points \(b\) as lower limits; if we assign these to the \(c\) in an arbitrary way and connect the corresponding pairs by definite integration paths which do not cut an arbitrarily prescribed path from \(\lambda\) to \(\mu\), then we obtain a perfectly definite function \(T_{\lambda\mu}(c)\), which is a constant.

In this constant, however, we may alter the assignment of the \(c,b\) if we observe the following precautions. We connect all points \(c\), in the imaginary plane, by definite paths, not cutting the parameter curve \(\lambda,\mu\), with a given point \(0\), which likewise corresponds to a definite point of \(f=0\). We then stipulate that, when the assignment of the \(b\) to the \(c\) is changed, the new paths of integration shall be composed of the original paths and of these new lines; then the value of \(T_{\lambda\mu}(c)\) remains wholly unchanged.

Now let one of the parameters move from \(\lambda\) towards \(\xi\), while the points \(c,\mu\) are to remain perfectly fixed. At the same time the points \(b\) move and finally pass into certain other points \(b'\). The differential of the transcendent \(T_{\xi\mu}\) is completely determined in every position (compare also \S\,51); but it may happen that the value of \(T_{\xi\mu}\) becomes infinite. By the properties of the integral of the third kind this could occur if one of the points \(\xi,\mu\) coincided with one of the points \(c,b'\); thus, since \(\mu\) and the \(c\) are considered constant, in the following cases:
\[
\begin{array}{ll}
1) & \text{if }\xi\text{ coincides with one of the points }c,\quad\hfill\\
2) & \text{if }\xi\text{ coincides with one of the points }b',\quad\hfill\\
3) & \text{if }\mu\text{ coincides with one of the points }b'.
\end{array}
\]
But in fact only the first case causes an infinity, since the other two cases cannot occur at all. For if one of the parameters, say \(\xi\), were to coincide with one of the \(b'\), then the curve of order \((n-2)\) which determines the \(b'\) would have \(n-1\) points in common with the line \(\xi\mu\), and would therefore contain that line wholly; hence a second of the points \(b'\) would have to fall in \(\mu\), and the remaining \(p-2\) points \(b'\) would lie, together with the points \(c\), on a curve of order \((n-3)\) passing through the double and return points, which is impossible by the assumption made about the \(c\).

Thus, if the point \(\xi\) is allowed to move in such a way that it never coincides with one of the points \(c\), the function \(T_{\xi\mu}(c)\) remains finite and determined at all times. Since during the whole motion neither parameter coincides with one of the limit-points, the condition contained in the definition of the integral of the third kind also causes no difficulty: namely, that a certain curve connecting the parameters should never be cut by the integration paths. None of these curves is absolutely fixed; rather, within the limits allowed for integration paths in general, it is movable in itself. If all these curves move with \(\xi\), one may at every moment give to the integration paths, or to the curve \(\mu\xi\), such evasions that no intersection occurs. And it is not even necessary for these evasions to take place continuously; all these curves may be placed at every moment as necessity requires, provided only that each of them can be led from one position into another without crossing the points of discontinuity or branch-points. The integration curves may themselves be changed so that both points of discontinuity are crossed at the same time; and the parameter curve may be displaced so that two corresponding limit-points are simultaneously crossed. For the value of the integral of the third kind is not changed if an integration path is replaced by another in such a way that the curve running between the parameters lies wholly between the two; and likewise this value is not changed if the parameter curve is replaced by another so that the whole integration path lies between the two.

Since here the transcendent \(T_{\xi\mu}(c)\) is determined not merely up to periods, but absolutely, \(\frac12T_{\xi\mu}(c)\) is also completely determined. We can further reduce first the transcendent \(T_{\xi\eta}(c)\) to this simple series of values, if we define it by the equation
\[
(1)\ldots\quad T_{\xi\eta}(c)=T_{\xi\mu}(c)-T_{\eta\mu}(c).
\]
If \(\xi\) and \(\eta\) have arisen from \(\lambda\) along definite paths, then the value of \(T_{\xi\eta}(c)\) is completely fixed by this formula.

Finally, every other transcendent \(T_{\xi\eta}(x)\) can be composed from this transcendent \(T_{\xi\eta}(c)\) and from the general transcendent \(T_{\xi\eta}\binom{x}{c}\), if the equation proved admissible above is taken as a definition:
\[
(2)\ldots\quad T_{\xi\eta}(x)=2T_{\xi\eta}\binom{x}{c}+T_{\xi\eta}(c).
\]
In order to fix in this formula the assumptions contained in the transcendent \(T_{\xi\eta}\binom{x}{c}\), let us suppose the \(x\) assigned to the \(c\) in some way and to have arisen from them along definite paths. During the description of these paths only the parameter curve \(\xi\eta\) (which is composed of the original curves \(\xi\mu\) and \(\eta\mu\)) must avoid them in such a way that finally the paths leading from the \(x\) to the \(c\) do not cut the final position of this curve. Thereby \(T_{\xi\eta}\binom{x}{c}\), and hence the right side of the above equation, obtains a wholly definite meaning, so that \(T_{\xi\eta}(x)\) too is perfectly fixed. If in particular the parameter curve is displaced so that the lines \(0c\) also are not cut by it, then the assignment of the \(x\) to the \(c\) may be changed arbitrarily, provided only that the new paths from the \(c\) to the \(x\) are again supposed to be composed of the former paths and of the lines \(0c\); the value of \(T_{\xi\eta}\binom{x}{c}\) is thereby not changed.

If now we denote by \(y\) the lower limit-points corresponding to the \(x\), then conversely the \(x\) also correspond to the \(y\) as upper limit-points and as lower ones; and we obtain two transcendents
\[
T_{\xi\eta}(x),\qquad T_{\xi\eta}(y),
\]
which have the same limit-points, so that their sum must be a periodicity modulus of an integral of the third kind. We wish to investigate under which assumptions, in particular, this sum is equal to zero.

For this purpose we now imagine the function \(T_{\lambda\mu}(c)\) itself to be derived continuously from another one. We start from a contact curve of order \((n-2)\), which passes through the double and return points of \(f=0\), as well as through the \(n-2\) other intersection points of the line \(\lambda\mu\), and which moreover touches the curve \(f=0\) in \(p\) points \(\varepsilon^{(1)},\varepsilon^{(2)},\ldots,\varepsilon^{(p)}\). Suppose now that a curve very near this curve is drawn, which still passes through the double and return points, through the \(n-2\) other points of the line \(\lambda\mu\), and through points \(\xi^{(1)},\xi^{(2)}\ldots\xi^{(p)}\) lying very near the \(\varepsilon\). It will then cut the curve \(f=0\) in a further \(p\) points \(\eta^{(1)},\eta^{(2)},\ldots,\eta^{(p)}\), which lie very near the points \(\varepsilon\) and are connected pairwise with the \(\xi\) by very small paths. If one now takes the \(\varepsilon\) in place of the above fundamental points \(c\), and the \(\eta\) in place of the \(b\), then \(T_{\lambda\mu}(\xi)\) is an infinitely small quantity; but also \(T_{\lambda\mu}(\eta)\), when it arises from \(T_{\lambda\mu}(\xi)\) by the infinitely small motions of the \(\xi\) to the \(\eta\), can only be infinitely small, and the \(\eta\) too can have described only very small paths. Hence the sum \(T_{\lambda\mu}(\xi)+T_{\lambda\mu}(\eta)\) cannot be an aggregate of the finite periodicity moduli, but must vanish.

If now we pass from the \(\varepsilon\) to the \(c\), and at the same time let the \(\eta\) pass to the \(b\), so that the paths leading to the latter are thereby fixed, then the sum
\[
T_{\lambda\mu}(\xi)+T_{\lambda\mu}(\eta)
\]
changes continuously, and therefore the sum can never cease to be zero. And the same holds when the \(c,b\) pass into the \(x,y\), and \(\lambda\) passes into \(\xi\) or \(\eta\), and so on. Under this assumption, therefore, the sum
\[
T_{\xi\eta}(x)+T_{\xi\eta}(y)
\]
is always equal to zero.

We add the following further remark. When \(\lambda\) is moved, the contact curve of the \(\varepsilon\) always passes continuously into other contact curves. If, in particular, we move the point \(\lambda\) toward \(\mu\) on the connecting line given between them, then in the end, when the two have approached one another infinitely closely, we obtain a contact curve which passes through the \(n-2\) remaining points of the tangent drawn at \(\mu\) to \(f=0\), and touches in \(p\) points \(\varepsilon\). We shall later make use of this curve.

\subsubsection*{\S\,48. Theory of the adjoint transcendents.}

A genuine decomposition of the special transcendent \(T_{\xi\eta}(x)\) into the difference of two functions of the same kind, each depending on only \(p+1\) points, is based on the theory of the \emph{adjoint transcendents}.

Through any \(p-1\) of the points \(x\) and through the double and return points of \(f=0\) there is determined a curve of order \((n-3)\), which cuts \(f=0\) in a further \(p-1\) points. If in \(T_{\xi\mu}(x)\), as it was defined in \S\,47, we put these in place of \(p-1\) of the points \(x\), and interchange the remaining point \(x\) with \(\xi\), there arises a series of \(p\) different transcendents, which we call adjoint to the given one. In this we pay no regard to the absolutely constant point \(\mu\).

If the \(p-1\) points lying on the same curve with
\[
x^{(1)},\ x^{(2)},\ldots,x^{(i-1)},\ x^{(i+1)},\ldots,x^{(p)}
\]
are denoted by
\[
x^{1i},\ x^{2i},\ldots,x^{i-1,i},\ x^{i+1,i},\ldots,x^{pi},
\]
then we have the following system of adjoint transcendents:\footnote{It may be noted here as interesting that the lower limit-points of all adjoint transcendents are the same, as is easily proved geometrically from Abel's theorem.}
\[
\begin{aligned}
(1)\ldots\quad T^{(1)}&=T_{x^{(1)}\mu}(\xi,x^{21}\ldots x^{p1}),\\
T^{(2)}&=T_{x^{(2)}\mu}(x^{12},\xi\ldots x^{p2}),\\[-0.35em]
&\cdots\\[-0.35em]
T^{(p)}&=T_{x^{(p)}\mu}(x^{1p},x^{2p}\ldots \xi).
\end{aligned}
\]
To define these transcendents completely exactly in the sense of \S\,47, we must know the paths by which here the \(x^{(i)}\) are connected with \(\mu\), and the points \(\xi,x^{ik}\) with the constant points \(c\). To carry this out, we also connect the point \(0\) (\S\,47) with the point \(\lambda\) by an arbitrarily chosen line. The definite definition of \(T_{\xi\eta}(x)\) requires that the origin of \(\xi\) from \(\lambda\) be given. We now assume that, in the adjoint transcendents, \(x^{(i)}\) is connected with \(\lambda\) by a curve which arises by allowed displacements from the path \(\lambda\xi0c^{(i)}x^{(i)}\), and likewise that \(\xi\) is connected with \(c^{(i)}\) by the path \(\xi0c^{(i)}\), or by one arising from it by allowed displacements.

The determination of the paths along which the points \(x^{ik}\) are to have arisen is led by the following remark. If one lets the points \(y^{(i)}\) belonging to the \(x^{(i)}\) arise on the path indicated in \S\,47, then they are connected with \(0\) by definite paths. If now, in place of \(x^{(1)}\), one chooses a point arbitrarily close to the point \(\xi\), and otherwise carries out the same process, so that the remaining points move from the \(c\) to the remaining \(x\), while this one point moves from \(c^{(1)}\) toward \(\xi\) on the prescribed path through \(0\), then one of the points \(y\) finally falls together with \(\mu\), since the curve of order \((n-2)\), drawn through the \(n-2\) remaining points of \(\xi\mu\) to determine the \(y\), resolves into a curve of order \((n-3)\) and the straight line \(\xi\mu\) itself. The remaining points \(y\) are then, in arbitrary order, the points \(x^{21},x^{31}\ldots x^{p1}\), and they are connected with the \(c\) by definite curves. These shall be used for the first adjoint transcendent, and the corresponding ones for the remaining cases, for the production of the upper limit-points from the \(c\).

It follows that, with the addition of the single line \(0\xi\), all adjoint transcendents are completely and uniquely determined by the given transcendent.

We now wish to transform the \(p\) adjoint transcendents by means of Abel's theorem so that the limit-points occurring in them agree as much as possible with those used in \(T\).

For this purpose we first set up a system of transcendents similar to these \(p\) transcendents, in which the points \(x\) are replaced by arbitrary constant points \(\alpha^{(1)},\alpha^{(2)}\ldots\alpha^{(p)}\), and in which at the same time the point \(\xi\) occurring above as argument is replaced respectively by \(\alpha^{(1)},\alpha^{(2)}\ldots\alpha^{(p)}\). This system, which is by no means a system of adjoint transcendents, is, with analogous notation,
\[
\begin{aligned}
(2)\ldots\quad X^{(1)}&=T_{x^{(1)}\mu}(\alpha^{(1)},\alpha^{21}\ldots\alpha^{p1}),\\
X^{(2)}&=T_{x^{(2)}\mu}(\alpha^{12},\alpha^{(2)}\ldots\alpha^{p2}),\\[-0.35em]
&\cdots\\[-0.35em]
X^{(p)}&=T_{x^{(p)}\mu}(\alpha^{1p},\alpha^{2p}\ldots\alpha^{(p)}).
\end{aligned}
\]
It has the property that each of the functions \(X\) depends on only \emph{one} variable point \(x\), which occurs in the corresponding \(X\) both as parameter and also in the lower limits. The points \(\alpha^{(1)},\alpha^{(2)}\ldots\alpha^{(p)}\) are here thought to be connected with the points \(c\) by arbitrary lines not cutting the lines \(x^{(i)}\mu\); while the \(\alpha^{ik}\) and the paths leading to them are derived from them just as above the \(x^{ik}\) were derived from the \(x^{(i)}\).

Since the same parameters respectively occur in the transcendents \(T^{(1)},T^{(2)}\ldots T^{(p)}\) and \(X^{(1)},X^{(2)}\ldots X^{(p)}\), half the difference of corresponding ones always forms a general transcendent \(T\), namely:
\begingroup\footnotesize
\[
\left\{
\begin{aligned}
\frac12(T^{(1)}-X^{(1)})
&=T_{x^{(1)}\mu}\begin{pmatrix}\xi,x^{21}\ldots x^{p1}\\ \alpha^{(1)},\alpha^{21}\ldots\alpha^{p1}\end{pmatrix} \\
&=\int_{\alpha^{(1)}}^{\xi}d\Pi_{x^{(1)}\mu}
+\int_{\alpha^{21}}^{x^{21}}d\Pi_{x^{(1)}\mu}\cdots
+\int_{\alpha^{p1}}^{x^{p1}}d\Pi_{x^{(1)}\mu},\\[0.35em]
\frac12(T^{(2)}-X^{(2)})
&=T_{x^{(2)}\mu}\begin{pmatrix}x^{12},\xi\ldots x^{p2}\\ \alpha^{12},\alpha^{(2)}\ldots\alpha^{p2}\end{pmatrix} \\
&=\int_{\alpha^{12}}^{x^{12}}d\Pi_{x^{(2)}\mu}
+\int_{\alpha^{(2)}}^{\xi}d\Pi_{x^{(2)}\mu}\cdots
+\int_{\alpha^{p2}}^{x^{p2}}d\Pi_{x^{(2)}\mu},\\[-0.2em]
&\cdots\\[-0.2em]
\frac12(T^{(p)}-X^{(p)})
&=T_{x^{(p)}\mu}\begin{pmatrix}x^{1p},x^{2p}\ldots\xi\\ \alpha^{1p},\alpha^{2p}\ldots\alpha^{(p)}\end{pmatrix} \\
&=\int_{\alpha^{1p}}^{x^{1p}}d\Pi_{x^{(p)}\mu}
+\int_{\alpha^{2p}}^{x^{2p}}d\Pi_{x^{(p)}\mu}\cdots
+\int_{\alpha^{(p)}}^{\xi}d\Pi_{x^{(p)}\mu}.
\end{aligned}
\right.
\tag{3}
\]
\endgroup
According to Abel's theorem, now,
\begingroup\footnotesize
\[
\left\{
\begin{aligned}
&\int_{\alpha^{1i}}^{x^{1i}}d\Pi_{x^{(i)}\mu}
+\int_{\alpha^{2i}}^{x^{2i}}d\Pi_{x^{(i)}\mu}\cdots
+\int_{\alpha^{i-1,i}}^{x^{i-1,i}}d\Pi_{x^{(i)}\mu} \\
&\quad +\int_{\alpha^{i+1,i}}^{x^{i+1,i}}d\Pi_{x^{(i)}\mu}\cdots
+\int_{\alpha^{pi}}^{x^{pi}}d\Pi_{x^{(i)}\mu} \\
&= -\left\{
\int_{\alpha^{(1)}}^{x^{(1)}}d\Pi_{x^{(i)}\mu}
+\int_{\alpha^{(2)}}^{x^{(2)}}d\Pi_{x^{(i)}\mu}\cdots
+\int_{\alpha^{(i-1)}}^{x^{(i-1)}}d\Pi_{x^{(i)}\mu}
\right.\\
&\qquad\left.
+\int_{\alpha^{(i+1)}}^{x^{(i+1)}}d\Pi_{x^{(i)}\mu}\cdots
+\int_{\alpha^{(p)}}^{x^{(p)}}d\Pi_{x^{(i)}\mu}\right\}
+\log\frac{\varphi(\mu)\psi(x^{(i)})}{\psi(\mu)\varphi(x^{(i)})},
\end{aligned}
\right.
\tag{4}
\]
\endgroup
where \(\varphi=0\) is the curve of order \((n-3)\) drawn through \(x^{(1)},x^{(2)}\ldots x^{(i-1)},x^{(i+1)}\ldots x^{(p)}\) and determining the points \(x^{1i},x^{2i}\ldots x^{i-1,i},x^{i+1,i}\ldots x^{pi}\), whereas \(\psi=0\) is the curve of order \((n-3)\) drawn through \(\alpha^{(1)},\alpha^{(2)}\ldots\alpha^{(i-1)},\alpha^{(i+1)}\ldots\alpha^{(p)}\) and determining the points \(\alpha^{1i},\alpha^{2i}\ldots\alpha^{i-1,i},\alpha^{i+1,i}\ldots\alpha^{pi}\).

In particular, denote by
\[
\varphi(y,y^{(1)},y^{(2)}\ldots y^{(p-1)})=0
\]
the equation of a curve of order \((n-3)\), which goes through the double and return points and through \(y^{(1)},y^{(2)}\ldots y^{(p-1)}\), while \(y\) is regarded as a variable point. The expression \(\varphi\) is a determinant which changes only its sign when any two of the points \(y\) are interchanged; hence the above equation may just as well be regarded as the equation of a curve drawn through \(y^{(1)},y^{(2)}\ldots y^{(i-1)},y^{(i+1)}\ldots y^{(p-1)}\), in which \(y^{(i)}\) is the variable point.

With this notation the function under the logarithm sign in (4) becomes
\begingroup\scriptsize
\[
\frac{\varphi(x^{(1)},x^{(2)}\ldots x^{(i-1)},\mu,x^{(i+1)},\ldots,x^{(p)})\,
\varphi(\alpha^{(1)},\alpha^{(2)}\ldots\alpha^{(i-1)},x^{(i)},\alpha^{(i+1)},\ldots,\alpha^{(p)})}
{\varphi(\alpha^{(1)},\alpha^{(2)}\ldots\alpha^{(i-1)},\mu,\alpha^{(i+1)},\ldots,\alpha^{(p)})\,
\varphi(x^{(1)},x^{(2)}\ldots x^{(i-1)},x^{(i)},x^{(i+1)},\ldots,x^{(p)})}.
\]
\endgroup
Hence equations (3) may be replaced by the following:
\begingroup\scriptsize
\[
\begin{aligned}
\frac12(T^{(1)}-X^{(1)})={}&\int_{\alpha^{(1)}}^{\xi}d\Pi_{x^{(1)}\mu}
-\int_{\alpha^{(2)}}^{x^{(2)}}d\Pi_{x^{(1)}\mu}\cdots
-\int_{\alpha^{(p)}}^{x^{(p)}}d\Pi_{x^{(1)}\mu} \\
&+\log\frac{\varphi(\mu,x^{(2)}\ldots x^{(p)})\,\varphi(x^{(1)},\alpha^{(2)}\ldots\alpha^{(p)})}
{\varphi(\mu,\alpha^{(2)}\ldots\alpha^{(p)})\,\varphi(x^{(1)},x^{(2)}\ldots x^{(p)})},\\[0.4em]
\frac12(T^{(2)}-X^{(2)})={}&-\int_{\alpha^{(1)}}^{x^{(1)}}d\Pi_{x^{(2)}\mu}
+\int_{\alpha^{(2)}}^{\xi}d\Pi_{x^{(2)}\mu}\cdots
-\int_{\alpha^{(p)}}^{x^{(p)}}d\Pi_{x^{(2)}\mu} \\
&+\log\frac{\varphi(x^{(1)},\mu\ldots x^{(p)})\,\varphi(\alpha^{(1)},x^{(2)}\ldots\alpha^{(p)})}
{\varphi(\alpha^{(1)},\mu\ldots\alpha^{(p)})\,\varphi(x^{(1)},x^{(2)}\ldots x^{(p)})},\\[-0.2em]
&\cdots\\[-0.2em]
\frac12(T^{(p)}-X^{(p)})={}&-\int_{\alpha^{(1)}}^{x^{(1)}}d\Pi_{x^{(p)}\mu}
-\int_{\alpha^{(2)}}^{x^{(2)}}d\Pi_{x^{(p)}\mu}\cdots
+\int_{\alpha^{(p)}}^{\xi}d\Pi_{x^{(p)}\mu} \\
&+\log\frac{\varphi(x^{(1)},x^{(2)}\ldots\mu)\,\varphi(\alpha^{(1)},\alpha^{(2)}\ldots x^{(p)})}
{\varphi(\alpha^{(1)},\alpha^{(2)}\ldots\mu)\,\varphi(x^{(1)},x^{(2)}\ldots x^{(p)})}.
\end{aligned}
\]
\endgroup
If one adds to these equations the definition of \(T\),
\[
\frac12T=\int_{\alpha^{(1)}}^{x^{(1)}}d\Pi_{\xi\mu}
+\int_{\alpha^{(2)}}^{x^{(2)}}d\Pi_{\xi\mu}\cdots
+\int_{\alpha^{(p)}}^{x^{(p)}}d\Pi_{\xi\mu}+\Xi,
\]
where \(\Xi\) denotes the function formed analogously to the transcendents \(X^{(1)},X^{(2)}\ldots X^{(p)}\),
\[
\Xi=T_{\xi\mu}(\alpha^{(1)},\alpha^{(2)},\ldots,\alpha^{(p)}),
\]
then one sees that in the whole system of transcendents only the limits \(\alpha,x,\xi\) occur in integrals of the third kind, apart from the functions \(X^{(1)},X^{(2)}\ldots X^{(p)}\), each of which depends on only one of the variable points.

We now examine those terms in the difference \(\frac12(T-T^{(i)})\) which depend at the same time on \(\xi\) and on \(x^{(i)}\). These terms can arise only from the difference
\[
\int_{\alpha^{(i)}}^{x^{(i)}}d\Pi_{\xi\mu}-\int_{\alpha^{(i)}}^{\xi}d\Pi_{x^{(i)}\mu};
\]
but this difference is equal to
\begingroup\footnotesize
\[
\int_{\alpha^{(i)}}^{\delta}d\Pi_{\xi\varepsilon}
+\int_{\delta}^{x^{(i)}}d\Pi_{\xi\varepsilon}
+\int_{\alpha^{(i)}}^{x^{(i)}}d\Pi_{\varepsilon\mu}
-\int_{\alpha^{(i)}}^{\varepsilon}d\Pi_{x^{(i)}\delta}
-\int_{\varepsilon}^{\xi}d\Pi_{x^{(i)}\delta}
-\int_{\alpha^{(i)}}^{\xi}d\Pi_{\delta\mu},
\]
\endgroup
where the terms depending simultaneously on the two points \(\xi,x^{(i)}\) cancel. Such terms therefore do not exist at all, and hence
\[
\frac{\partial}{\partial x^{(i)}}\left(\frac{\partial T}{\partial\xi}\right)
=\frac{\partial}{\partial\xi}\left(\frac{\partial T^{(i)}}{\partial x^{(i)}}\right).
\]
In exactly the same way the difference
\[
\frac12(T^{(i)}-T^{(k)})
\]
contains no term depending simultaneously on \(x^{(i)}\) and \(x^{(k)}\). Therefore also
\[
\frac{\partial}{\partial x^{(i)}}\left(\frac{\partial T^{(k)}}{\partial x^{(k)}}\right)
=\frac{\partial}{\partial x^{(k)}}\left(\frac{\partial T^{(i)}}{\partial x^{(i)}}\right).
\]
These equations show that the following remarkable theorem, fundamental for all subsequent investigations, holds:

The differential quotients of the adjoint transcendents
\[
\frac12T,\quad \frac12T^{(1)},\quad \frac12T^{(2)},\ldots,\frac12T^{(p)}
\]
with respect to their corresponding variables
\[
\xi,\quad x^{(1)},\quad x^{(2)},\ldots,x^{(p)}
\]
are the coefficients of a complete differential:
\[
dU=\frac12\left(\frac{\partial T}{\partial\xi}d\xi
+\frac{\partial T^{(1)}}{\partial x^{(1)}}dx^{(1)}\cdots
+\frac{\partial T^{(p)}}{\partial x^{(p)}}dx^{(p)}\right).
\]
It deserves to be noted, moreover, that these differential quotients no longer contain the point \(\mu\) in any way.

\subsubsection*{\S\,49. The function \(U\). The transcendent \(T_{\xi\eta}\) expressed by it.}

We shall define the function \(U\), which is determined by the above equation up to a constant, completely by the stipulation that \(U\) shall vanish when \(\xi,x^{(1)},x^{(2)}\ldots x^{(p)}\) coincide with an arbitrarily given constant system of points \(\zeta,\gamma^{(1)},\gamma^{(2)}\ldots\gamma^{(p)}\).

Thus let us also suppose the point-system \(\zeta,\gamma\), as above the point-system \(\xi,x\), to be connected with the \(\lambda,c\) by definite given curves.

If we now introduce the differential quotients of the \(T\)'s into the expression for \(U\), and integrate between the systems of values \(\zeta,\gamma^{(1)},\gamma^{(2)}\ldots\gamma^{(p)}\) and \(\xi,x^{(1)},x^{(2)}\ldots x^{(p)}\) along the paths prescribed by the graphical representation, or paths derived from them by allowed displacements, we obtain for \(U\) the following completely definite expression:
\begingroup\scriptsize
\[
\begin{aligned}
U={}&\sum_{i=1}^{i=p}\int_{\zeta,\gamma^{(i)}}^{\xi,x^{(i)}}
\left(\frac{\partial\Pi_{x^{(i)}\alpha^{(i)}}(\xi)}{\partial\xi}d\xi
+\frac{\partial\Pi_{\xi\alpha^{(i)}}(x^{(i)})}{\partial x^{(i)}}dx^{(i)}\right)\\
&-\frac12\sum\sum_{i\le k}\int_{\gamma^{(i)},\gamma^{(k)}}^{x^{(i)},x^{(k)}}
\left(\frac{\partial\Pi_{x^{(i)}\alpha^{(i)}}(x^{(k)})}{\partial x^{(k)}}dx^{(k)}
+\frac{\partial\Pi_{x^{(k)}\alpha^{(k)}}(x^{(i)})}{\partial x^{(i)}}dx^{(i)}\right)\\
&+\frac12\int_{\zeta}^{\xi}\frac{\partial\Xi}{\partial\xi}d\xi
+\frac12\sum_{i=1}^{i=p}\int_{\gamma^{(i)}}^{x^{(i)}}\frac{\partial X^{(i)}}{\partial x^{(i)}}dx^{(i)}\\
&-\log\frac{\varphi(x^{(1)},x^{(2)}\ldots x^{(p)})\,\varphi(\gamma^{(1)},\alpha^{(2)}\ldots\alpha^{(p)})\,\varphi(\alpha^{(1)},\gamma^{(2)}\ldots\alpha^{(p)})\cdots\varphi(\alpha^{(1)},\alpha^{(2)}\ldots\gamma^{(p)})}
{\varphi(\gamma^{(1)},\gamma^{(2)}\ldots\gamma^{(p)})\,\varphi(x^{(1)},\alpha^{(2)}\ldots\alpha^{(p)})\,\varphi(\alpha^{(1)},x^{(2)}\ldots\alpha^{(p)})\cdots\varphi(\alpha^{(1)},\alpha^{(2)}\ldots x^{(p)})}.
\end{aligned}
\]
\endgroup
The second sum is divided by \(2\), so that \(i\) and \(k\), independently of one another, may run through all pairs of values except the equal ones, in which case two equal terms always occur. The first parts of this expression are easily reduced to integrals of the third kind. For by familiar rules one has
\begingroup\footnotesize
\[
\begin{aligned}
&\int_{\zeta,\gamma^{(i)}}^{\xi,x^{(i)}}
\left(\frac{\partial\Pi_{x^{(i)}\alpha^{(i)}}(\xi)}{\partial\xi}d\xi
+\frac{\partial\Pi_{\xi\alpha^{(i)}}(x^{(i)})}{\partial x^{(i)}}dx^{(i)}\right)\\
&\quad=\int_{\zeta}^{\xi}\frac{\partial\Pi_{x^{(i)}\alpha^{(i)}}(\xi)}{\partial\xi}d\xi
+\int_{\gamma^{(i)}}^{x^{(i)}}\frac{\partial\Pi_{\xi\alpha^{(i)}}(x^{(i)})}{\partial x^{(i)}}dx^{(i)},
\end{aligned}
\]
\endgroup
or, in the notation always used previously,
\[
=\int_{\zeta}^{\xi}d\Pi_{x^{(i)}\alpha^{(i)}}+
\int_{\gamma^{(i)}}^{x^{(i)}}d\Pi_{\xi\alpha^{(i)}}.
\]
Exactly similarly one has
\begingroup\footnotesize
\[
\begin{aligned}
&\int_{\gamma^{(i)},\gamma^{(k)}}^{x^{(i)},x^{(k)}}
\left(\frac{\partial\Pi_{x^{(i)}\alpha^{(i)}}(x^{(k)})}{\partial x^{(k)}}dx^{(k)}
+\frac{\partial\Pi_{x^{(k)}\alpha^{(k)}}(x^{(i)})}{\partial x^{(i)}}dx^{(i)}\right)\\
&\quad=\int_{\gamma^{(k)}}^{x^{(k)}}d\Pi_{x^{(i)}\alpha^{(i)}}+
\int_{\gamma^{(i)}}^{x^{(i)}}d\Pi_{\gamma^{(k)}\alpha^{(k)}}.
\end{aligned}
\]
\endgroup
Moreover, by the theorem given at the end of \S\,46,
\[
\int_{\zeta}^{\xi}\frac{\partial\Xi}{\partial\xi}d\xi
=T_{\xi\zeta}(\alpha^{(1)},\alpha^{(2)}\ldots\alpha^{(p)}),
\]
\[
\int_{\gamma^{(i)}}^{x^{(i)}}\frac{\partial X^{(i)}}{\partial x^{(i)}}dx^{(i)}
=T_{x^{(i)}\gamma^{(i)}}(\alpha^{1i},\alpha^{2i}\ldots\alpha^{i-1,i},\alpha^{(i)},\alpha^{i+1,i}\ldots\alpha^{pi}).
\]
Consequently the expression for \(U\) finally becomes
\begingroup\scriptsize
\[
\begin{aligned}
U={}&\sum_{i=1}^{i=p}\int_{\zeta}^{\xi}d\Pi_{x^{(i)}\alpha^{(i)}}
+\sum_{i=1}^{i=p}\int_{\gamma^{(i)}}^{x^{(i)}}d\Pi_{\xi\alpha^{(i)}}\\
&-\frac12\sum\sum_{i\le k}\int_{\gamma^{(k)}}^{x^{(k)}}d\Pi_{x^{(i)}\alpha^{(i)}}
-\frac12\sum\sum_{i\le k}\int_{\gamma^{(i)}}^{x^{(i)}}d\Pi_{\gamma^{(k)}\alpha^{(k)}}\\
&+\frac12T_{\xi\zeta}(\alpha^{(1)},\alpha^{(2)}\ldots\alpha^{(p)})
+\frac12\sum_{i=1}^{i=p}T_{x^{(i)}\gamma^{(i)}}(\alpha^{1i},\alpha^{2i}\ldots\alpha^{i-1,i},\alpha^{(i)},\alpha^{i+1,i}\ldots\alpha^{pi})\\
&-\log\frac{\varphi(x^{(1)},x^{(2)}\ldots x^{(p)})\,\varphi(\gamma^{(1)},\alpha^{(2)}\ldots\alpha^{(p)})\,\varphi(\alpha^{(1)},\gamma^{(2)}\ldots\alpha^{(p)})\cdots\varphi(\alpha^{(1)},\alpha^{(2)}\ldots\gamma^{(p)})}
{\varphi(\gamma^{(1)},\gamma^{(2)}\ldots\gamma^{(p)})\,\varphi(x^{(1)},\alpha^{(2)}\ldots\alpha^{(p)})\,\varphi(\alpha^{(1)},x^{(2)}\ldots\alpha^{(p)})\cdots\varphi(\alpha^{(1)},\alpha^{(2)}\ldots x^{(p)})}.
\end{aligned}
\]
\endgroup
That in this expression the \(\xi,\gamma\) influence only an additive constant follows from the manner in which the expression is formed. But it is noteworthy that the constant points \(\alpha^{(i)}\), and hence also the constant points \(\alpha^{ik}\) depending on them, have no influence at all on the value of \(U\), and so occur in it only formally. In

in fact, if we form a second function \(U'\), which differs from the former in that the \(\alpha\) are replaced by other points \(\beta\), one finds \(U-U'=0\). This also follows immediately from the fact that the conjugate transcendents \(T,T^{(1)},\ldots,T^{(p)}\) were already, by their origin, independent of the \(\alpha\), and therefore so are \(dU\) and consequently \(U\) itself.

Since \(\mu\) has entirely disappeared from the expression for \(U\), one sees that \(U\) must have the form
\[
U=\Phi(\xi,x^{(1)},x^{(2)},\ldots,x^{(p)})-\Phi(\zeta,\gamma^{(1)},\gamma^{(2)},\ldots,\gamma^{(p)}),
\]
that is, apart from an additive constant, \(U\) is exclusively a function of the \(p+1\) points \(\xi,x^{(1)},x^{(2)},\ldots x^{(p)}\). We shall accordingly denote it by
\[
U(x^{(1)},x^{(2)},\ldots,x^{(p)};\xi).
\]
We shall now use this function in order to represent by it the special and the general transcendent \(T\).

Indeed, the difference
\[
U(x^{(1)},x^{(2)}\ldots x^{(p)};\xi)-U(x^{(1)},x^{(2)}\ldots x^{(p)};\mu)
\]
is nothing other than the integral of \(dU\), taken between the systems of values
\[
x^{(1)},x^{(2)}\ldots\mu\quad \text{to}\quad x^{(1)},x^{(2)}\ldots\xi.
\]
We shall connect \(\mu\xi\) by the curve already mentioned above, and integrate along this curve. Therefore, with regard to the expression for \(dU\),
\[
\begin{aligned}
&U(x^{(1)},x^{(2)}\ldots x^{(p)};\xi)-U(x^{(1)},x^{(2)}\ldots x^{(p)};\mu)\\
&\quad=\frac12\int_{\mu}^{\xi}\frac{\partial T}{\partial\xi}\,d\xi\\
&\quad=\frac12T_{\xi\mu}(x^{(1)},x^{(2)}\ldots x^{(p)})
-\frac12T_{\mu\mu}(x^{(1)},x^{(2)}\ldots x^{(p)}),
\end{aligned}
\]
where in the first term on the right that curve connects the parameter-points, but in the second term the connecting line becomes infinitely small. Thus the transcendent \(T_{\mu\mu}\) so defined vanishes, and one has the formula
\[
\frac12T_{\xi\mu}(x^{(1)},x^{(2)}\ldots x^{(p)})=U(x^{(1)},x^{(2)}\ldots x^{(p)};\xi)-U(x^{(1)},x^{(2)}\ldots x^{(p)};\mu);
\]
and therefore also, according to the definition of \(T_{\xi\eta}(x)\),
\[
\frac12T_{\xi\eta}(x^{(1)},x^{(2)}\ldots x^{(p)})=U(x^{(1)},x^{(2)}\ldots x^{(p)};\xi)-U(x^{(1)},x^{(2)}\ldots x^{(p)};\eta).
\]
By the function \(U\), therefore, the genuine decomposition of \(\frac12T_{\xi\eta}(x)\) into two functions with \(p+1\) arguments each has been achieved. One immediately also has, for the general transcendent \(T_{\xi\eta}\binom{x}{z}\), the expression
\[
T_{\xi\eta}\binom{x^{(1)},x^{(2)}\ldots x^{(p)}}{z^{(1)},z^{(2)}\ldots z^{(p)}}
=U(x^{(1)},\ldots;\xi)-U(x^{(1)},\ldots;\eta)-U(z^{(1)},\ldots;\xi)+U(z^{(1)},\ldots;\eta).
\]

\subsubsection*{\S\,50. Discontinuities of the function \(U\). The function \(V=e^{-U}\).}

We now wish to investigate the discontinuities of \(U\). This function is composed of integrals of the third kind, of transcendents \(\frac12T_{\xi\eta}(x)\), and of logarithms of algebraic functions. As concerns the discontinuities of the transcendent \(\frac12T_{\xi\eta}(x)\), it is enough to know those of \(T_{\xi\mu}(c)\), where \(c,\mu\) are constant points, the former of which do not lie with the double and cusp points on a curve of order \((n-3)\). For by \S\,45 one can always express the transcendent \(T_{\xi\eta}(x)\) by integrals of the third kind with independent limits and by \(T_{\xi\mu}(c)\); but the discontinuities of the latter function have been set forth in \S\,47.

If we consider the expression \(U\), as it lies explicitly before us in the preceding paragraph, we must first observe that the \(\alpha\) in reality do not affect the value of \(U\) at all, and that \(\xi\) and the \(\gamma\) affect only an additive constant adjoined to it. Therefore one need not take into account those cases in which a term of \(U\) becomes discontinuous because \(\xi\) or one of the \(x\) approaches one of these points; then there necessarily always exists another term which cancels this discontinuity. Thus the following cases remain:

1) The point \(\xi\) coincides with one of the points \(x\), say with \(x^{(i)}\). In such a case \(U\) behaves like the term
\[
\int_{\zeta}^{\xi}d\Pi_{x^{(i)}\alpha^{(i)}},
\]
that is, like \(-\log\Sigma\pm\xi_1x_2^{(i)}\alpha_3^{(i)}\).

2) A point \(x^{(i)}\) coincides with a point \(x^{(k)}\). Then \(U\) behaves like
\[
-\int_{\gamma^{(k)}}^{x^{(k)}}d\Pi_{x^{(i)}\alpha^{(i)}}-\log\varphi(x^{(1)},x^{(2)}\ldots x^{(p)}),
\]
that is, like \(\log\Sigma\pm x_1^{(k)}x_2^{(i)}\alpha_3^{(i)}-\log\varphi(x^{(1)},x^{(2)}\ldots x^{(p)})\); this remains finite as the logarithm of the quotient of two simply vanishing functions.

3) The same is true if, say, \(x^{(i)}\) coincides with one of the double or cusp points of \(f=0\). In this case \(U\) behaves like
\[
-\log\frac{\varphi(x^{(1)},x^{(2)}\ldots x^{(p)})}{\varphi(\alpha^{(1)},\ldots x^{(i)}\ldots\alpha^{(p)})},
\]
and so remains finite.

4) On the other hand \(U\) actually becomes still infinitely large when the \(p\) points \(x\) lie on a curve of order \((n-3)\),
\[
\varphi(x^{(1)},x^{(2)}\ldots x^{(p)})=0,
\]
which passes through the double and cusp points; and in this case \(U\) behaves like
\[
-\log\varphi(x^{(1)},x^{(2)}\ldots x^{(p)}).
\]
Thus one can state the following theorem:

The function \(U(x^{(1)},x^{(2)}\ldots x^{(p)};\xi)\) becomes infinite only when \(\xi\) coincides with one of the \(x\), or when the latter lie on a curve of order \((n-3)\) passing through the double and cusp points of \(f=0\); and in all these cases \(U\) behaves like the negative logarithm of a vanishing quantity.

It is now natural to investigate, instead of the function \(U\) itself, the function
\[
V(x^{(1)},x^{(2)}\ldots x^{(p)};\xi)=e^{-U(x^{(1)},x^{(2)}\ldots x^{(p)};\xi)}.
\]
The function contains the constants \(\zeta,\gamma\) only in a factor. Its most essential property is the following, which follows immediately from the concluding theorem above:

The function \(V=e^{-U}\) is a function of the points \(x^{(1)},x^{(2)}\ldots x^{(p)},\xi\), which never becomes infinitely large (except perhaps as a consequence of infinitely often repeated paths of integration), and which vanishes only when \(\xi\) coincides with one of the points \(x\), or when the points \(x\), together with the double and cusp points, lie on a curve of order \((n-3)\).

With the help of this function one obtains, for the special transcendent \(T_{\xi\eta}(x)\), the expression
\[
\frac12T_{\xi\eta}(x^{(1)},x^{(2)}\ldots x^{(p)})
=\log\frac{V(x^{(1)},x^{(2)}\ldots x^{(p)};\eta)}{V(x^{(1)},x^{(2)}\ldots x^{(p)};\xi)};
\]
for the general transcendent \(T_{\xi\eta}\binom{x}{z}\) one finds
\[
T_{\xi\eta}\binom{x^{(1)},x^{(2)}\ldots x^{(p)}}{z^{(1)},z^{(2)}\ldots z^{(p)}}
=\log\frac{V(x^{(1)},x^{(2)}\ldots x^{(p)};\eta)V(z^{(1)},z^{(2)}\ldots z^{(p)};\xi)}{V(z^{(1)},z^{(2)}\ldots z^{(p)};\eta)V(x^{(1)},x^{(2)}\ldots x^{(p)};\xi)}.
\]
And finally, for the fundamental formula of the inversion problem, with suitably chosen integration paths, one finds
\begingroup\small
\[
\prod_{i=1}^{i=p}\frac{\left(\frac{\varphi}{\psi}\right)_{x^{(i)}}-\lambda}{\left(\frac{\varphi}{\psi}\right)_{c^{(i)}}-\lambda}
=\prod_{\xi\eta}\frac{V(x^{(1)},x^{(2)}\ldots x^{(p)};\xi)V(c^{(1)},c^{(2)}\ldots c^{(p)};\eta)}{V(x^{(1)},x^{(2)}\ldots x^{(p)};\eta)V(c^{(1)},c^{(2)}\ldots c^{(p)};\xi)},
\]
\endgroup
whereby everything is reduced to the investigation of the everywhere finite function \(V\).

\subsubsection*{\S\,51. Properties of the function \(V=e^{-U}\).}

The difference of two functions \(U\), which differ only by the values of \(\xi\), led to the special transcendent \(\frac12T_{\xi\eta}(x)\). In exactly the same way one obtains the sequence of the adjoint transcendents by considering the difference of two functions \(U\) which differ from one another only by one of the points \(x\). Indeed, from the differential expression for \(U\) one obtains
\[
\begin{aligned}
&U(x^{(1)},x^{(2)}\ldots x^{(p)};\xi)-U(y^{(1)},x^{(2)}\ldots x^{(p)};\xi)\nobreak\hspace{0.25em}\nobreak=\frac12\int_{y^{(1)}}^{x^{(1)}}\frac{\partial T^{(1)}}{\partial x^{(1)}}\,dx^{(1)}\\
&\quad=\frac12T_{x^{(1)}\eta}(\xi,x^{21},x^{31}\ldots x^{p1})-\frac12T_{y^{(1)}\eta}(\xi,x^{21},x^{31}\ldots x^{p1})\\
&\quad=\frac12T_{x^{(1)}y^{(1)}}(\xi,x^{21},x^{31}\ldots x^{p1}).
\end{aligned}
\]
This transcendent is to be defined so that it becomes the difference of the two adjoint transcendents
\[
\frac12T_{x^{(1)}\mu}(\xi,x^{21}\ldots)-\frac12T_{y^{(1)}\mu}(\xi,x^{21}\ldots).
\]
Thus one must connect \(y^{(1)}\) by a curve with the curve-system mentioned above, and then obtains both transcendents, from \(\frac12T_{\lambda\mu}(c)\), uniquely according to the preceding discussion.

If one now expresses this adjoint transcendent directly by functions \(U\), one obtains the relation
\[
\tag{1}\begin{aligned}
&U(x^{(1)},x^{(2)}\ldots x^{(p)};\xi)-U(y^{(1)},x^{(2)}\ldots x^{(p)};\xi)\\
&\quad=U(\xi,x^{21},x^{31}\ldots x^{p1};x^{(1)})-U(\xi,x^{21},x^{31}\ldots x^{p1};y^{(1)}),
\end{aligned}
\]
or, if we introduce the function \(V\),
\[
\tag{2}\frac{V(x^{(1)},x^{(2)}\ldots x^{(p)};\xi)}{V(y^{(1)},x^{(2)}\ldots x^{(p)};\xi)}
=\frac{V(\xi,x^{21},x^{31}\ldots x^{p1};x^{(1)})}{V(\xi,x^{21},x^{31}\ldots x^{p1};y^{(1)})}.
\]
This formula teaches how to reduce the quotient of two functions \(V\), which differ by an argument, to the quotient of two \(V\)'s which differ by a parameter, and conversely.

The expression \(\frac12T_{x^{(1)}y^{(1)}}(\xi,x^{21},x^{31}\ldots x^{p1})\), however, can under certain circumstances be reduced to functions \(U\) in still another way: namely when \(y^{(1)}\) lies with the \(x\) on a curve of order \((n-2)\), which passes through the double and cusp points of \(f=0\), and in addition through the remaining \(n-2\) intersection-points of a line \(\xi\eta\).

Let the further intersection-points of this curve with \(f=0\) be denoted by \(y^{(2)},y^{(3)}\ldots y^{(p)}\), and let \(y^{21},y^{31}\ldots y^{p1}\) be the points of \(f=0\) which, with the double and cusp points and with \(y^{(2)},y^{(3)}\ldots y^{(p)}\), lie on a curve of order \((n-3)\). Then the following curves:
\begin{center}
\begin{tabular}{@{}p{5.15in}@{}}
the line \(\xi\eta\),\\
the line \(x^{(1)}y^{(1)}\),\\
the curve of order \((n-3)\) through \(x^{(2)},x^{(3)},\ldots,x^{21},x^{31}\ldots\) and the double and cusp points,\\
the curve of order \((n-3)\) through \(y^{(2)},y^{(3)},\ldots,y^{31},y^{32}\ldots\) and the double and cusp points,
\end{tabular}
\end{center}
together form an improper curve of order \(2n-4\), which has the double and cusp points of \(f=0\) themselves as double points. But among the intersection-points of this curve with \(f=0\), the assumption says that the remaining \(n-2\) points of the line \(\xi\eta\), the \(x\), and the \(y\), lie on a curve of order \((n-2)\) passing through the double and cusp points. The same must therefore be true of the remainder; thus the remaining points of the line \(x^{(1)},y^{(1)}\) lie on such a curve with \(\xi,\eta,x^{21},x^{31}\ldots,y^{21},y^{31}\ldots\). According to the assumptions made above concerning the formation of the transcendent \(\frac12T_{\xi\eta}\), one therefore has the equation
\[
T_{x^{(1)}y^{(1)}}(\xi,x^{21},x^{31}\ldots)+T_{x^{(1)}y^{(1)}}(\eta,y^{21},y^{31}\ldots)=0.
\]
For the functions \(U\) this gives the equation
\[
\begin{aligned}
&U(\xi,x^{21},x^{31}\ldots x^{p1};x^{(1)})-U(\xi,x^{21},x^{31}\ldots x^{p1};y^{(1)})\\
&\quad=-U(\eta,y^{21},y^{31}\ldots y^{p1};x^{(1)})+U(\eta,y^{21},y^{31}\ldots y^{p1};y^{(1)}),
\end{aligned}
\]
or, applying equation (1),
\[
\begin{aligned}
&U(x^{(1)},x^{(2)}\ldots x^{(p)};\xi)-U(y^{(1)},x^{(2)}\ldots x^{(p)};\xi)\\
&
=-U(x^{(1)},y^{(2)}\ldots y^{(p)};\eta)+U(y^{(1)},y^{(2)}\ldots y^{(p)};\eta),
\end{aligned}
\]
or finally, if we introduce the function \(V\),
\[
\tag{3}\frac{V(x^{(1)},x^{(2)}\ldots x^{(p)};\xi)}{V(y^{(1)},x^{(2)}\ldots x^{(p)};\xi)}
=\frac{V(y^{(1)},y^{(2)}\ldots y^{(p)};\eta)}{V(x^{(1)},y^{(2)}\ldots y^{(p)};\eta)}.
\]
Writing this in the form
\[
\frac{V(x^{(1)},x^{(2)}\ldots x^{(p)};\xi)}{V(y^{(1)},y^{(2)}\ldots y^{(p)};\eta)}
=\frac{V(y^{(1)},x^{(2)}\ldots x^{(p)};\xi)}{V(x^{(1)},y^{(2)}\ldots y^{(p)};\eta)},
\]
one sees that the expression on the left is not changed by interchanging a pair \(x^{(1)},y^{(1)}\); here it is assumed only that the points occurring as arguments in numerator and denominator, together with the remaining points of the line \(\xi\eta\), lie on a curve of order \((n-2)\) passing through the double and cusp points. Since this condition is still likewise satisfied on the right, we can continue to interchange further pairs \(x,y\) without changing the value of the quotient, and thus finally arrive at the equation
\[
\tag{4}\frac{V(x^{(1)},x^{(2)}\ldots x^{(p)};\xi)}{V(y^{(1)},y^{(2)}\ldots y^{(p)};\eta)}
=\frac{V(y^{(1)},y^{(2)}\ldots y^{(p)};\xi)}{V(x^{(1)},x^{(2)}\ldots x^{(p)};\eta)}.
\]
Now consider secondly the sum
\[
T_{\xi\eta}(x^{(1)}x^{(2)}\ldots x^{(p)})+T_{\xi\eta}(y^{(1)}y^{(2)}\ldots y^{(p)}).
\]
Again the assumptions of \S\,47 show that this sum vanishes, and therefore, if one transfers the equation so obtained to the functions \(V\), one also has
\[
\tag{5}\frac{V(x^{(1)},x^{(2)}\ldots x^{(p)};\xi)}{V(x^{(1)},x^{(2)}\ldots x^{(p)};\eta)}
=\frac{V(y^{(1)},y^{(2)}\ldots y^{(p)};\eta)}{V(y^{(1)},y^{(2)}\ldots y^{(p)};\xi)}.
\]
The combination of equations (4) and (5) now immediately yields the relation
\[
V(x^{(1)},x^{(2)}\ldots x^{(p)};\xi)=\pm V(y^{(1)},y^{(2)}\ldots y^{(p)};\eta).
\]
Which sign is to be chosen here depends on the manner in which the function \(V\) is defined. For continuity it is sufficient that either the upper sign or always the lower sign be chosen. Therefore we have the following theorem:

I. If one forms the function \(V\) for \(x^{(1)},x^{(2)}\ldots x^{(p)},\xi\) and for \(y^{(1)},y^{(2)}\ldots y^{(p)},\eta\), where the \(x,y\) lie on a curve of order \((n-2)\) which passes through the double and cusp points, and also through the remaining \(n-2\) intersection-points of \(f=0\) with the line \(\xi\eta\), then, with a suitable choice of the integration paths, in both cases there results the same value or the opposite value.

Now let the points \(y,\eta\) be fixed and move the \(x,\xi\); then by the foregoing the value of \(V\) remains unchanged, and one has the following theorem:

II. Let a ray rotate about a point of the curve \(f=0\), and in each position of it draw a curve of order \((n-2)\) through \(n-2\) of its other intersection-points, through the double and cusp points, and through \(p\) fixed points of \(f=0\). Then this curve cuts the latter in a further \(p\) points \(x\), and the function \(V\), formed for these \(p\) points and with a parameter which corresponds to the last intersection-point of the ray, always has the same value.

A special case of theorem I deserves particular emphasis. Namely, if \(\eta\) coincides with one of the points \(x\), say with \(x^{(1)}\), then the curve of order \((n-2)\) passes into the line \(\xi x^{(1)}\) and into a curve of order \((n-3)\), which passes through \(x^{(2)},x^{(3)}\ldots x^{(p)}\) and the double and cusp points and is thereby determined, and whose remaining intersection-points with \(f=0\) are the points \(x^{21},x^{31}\ldots x^{p1}\). Thus one has to replace the \(y\) by \(\xi,x^{21},x^{31}\ldots x^{p1}\), and so obtains the theorem:

III. The function \(V\) does not change, or changes at most its sign, if one interchanges the parameter with one of its arguments and replaces the remaining arguments by those points which lie with them on a curve of order \((n-3)\) passing through the double and cusp points.

Theorem II is equivalent to the setting up of a partial differential equation which the function \(V\) satisfies. Namely, if \(\eta\) is again the fixed point of the ray, and if \(\xi,\xi^{(1)},\xi^{(2)}\ldots\xi^{(n-2)}\) are its remaining intersection-points; and finally if the \(x\), together with the \(\xi\) and fixed points, lie on a curve of order \((n-2)\), then from Abel's theorem for the corresponding integrals of the first kind there follow the differential equations
\[
\frac{\partial u_i(\xi)}{\partial\xi}\,d\xi+
\sum_{h=1}^{h=n-2}\frac{\partial u_i(\xi^{(h)})}{\partial\xi^{(h)}}\,d\xi^{(h)}=0,
\]
\[
\sum_{h=1}^{h=p}\frac{\partial u_i(x^{(h)})}{\partial x^{(h)}}\,dx^{(h)}+
\sum_{h=1}^{h=n-2}\frac{\partial u_i(\xi^{(h)})}{\partial\xi^{(h)}}\,d\xi^{(h)}=0,
\]
or, after elimination of the \(\xi\),
\[
\tag{6}\sum_{h=1}^{h=p}\frac{\partial u_i(x^{(h)})}{\partial x^{(h)}}\,dx^{(h)}-\frac{\partial u_i(\xi)}{\partial\xi}\,d\xi=0.
\quad (i=1,2\ldots p).
\]
These \(p\) differential equations form a system with \(p+1\) variables, whose complete integrals are given by the \(p\) equations
\[
\tag{7}\sum_{h=1}^{h=p}\int_{c^{(h)}}^{x^{(h)}}du_i-\int_{\mu}^{\xi}du_i=\mathrm{Const.}
\quad (i=1,2\ldots p).
\]
At the same time, however, by theorem II,
\[
dV=\frac{\partial V}{\partial\xi}\,d\xi+
\sum_{h=1}^{h=p}\frac{\partial V}{\partial x^{(h)}}\,dx^{(h)}=0.
\]
Combining this with the differential equations (6), and eliminating the \(dx\), one obtains the partial differential equation which \(V\) satisfies:
\[
\det\begin{pmatrix}
\dfrac{\partial V}{\partial x^{(1)}} & \dfrac{\partial V}{\partial x^{(2)}} & \cdots & \dfrac{\partial V}{\partial\xi}\\[0.7em]
\dfrac{\partial u_1(x^{(1)})}{\partial x^{(1)}} & \dfrac{\partial u_1(x^{(2)})}{\partial x^{(2)}} & \cdots & -\dfrac{\partial u_1(\xi)}{\partial\xi}\\
\vdots & \vdots & & \vdots\\
\dfrac{\partial u_p(x^{(1)})}{\partial x^{(1)}} & \dfrac{\partial u_p(x^{(2)})}{\partial x^{(2)}} & \cdots & -\dfrac{\partial u_p(\xi)}{\partial\xi}
\end{pmatrix}=0.
\]
Integrating this, one finds \(V\) as a function of the integrals of the system (6), that is, of the left-hand sides of (7). Thus we have the theorem:

The function \(V\) is in reality dependent on only \(p\) quantities, namely on the \(p\) sums
\[
\sum_{h=1}^{h=p}\int_{c^{(h)}}^{x^{(h)}}du_i-\int_{\mu}^{\xi}du_i,
\quad (i=1,2\ldots p),
\]
or, in the sense of the inversion problem, it is a function of the \(p\) quantities
\[
v_i-\int_{\mu}^{\xi}du_i.
\]
The determination of this function is the solution of the inversion problem. We shall treat this in the following section.

\section*{Eighth Section.}
\section*{Monodromy.}

\subsubsection*{\S\,52. Differential of an Abelian function. Value-systems of the \(v\) for which the system of the \(x\) is indeterminate.}

We have seen above that the inversion problem consists in considering symmetric functions of the coordinates of \(p\) points as functions of the integral sums \(v_1,v_2\ldots v_p\). We now wish to investigate the character of these functions more closely. We shall prove that they are monodromic functions of their arguments.

We can express the connection of the \(x\) with the \(v\) by saying that, between these two systems of quantities, the differential equations
\[
\tag{1}\left\{\begin{aligned}
dv_1&=\frac{\partial u_1(x^{(1)})}{\partial x^{(1)}}dx^{(1)}+\frac{\partial u_1(x^{(2)})}{\partial x^{(2)}}dx^{(2)}\cdots+\frac{\partial u_1(x^{(p)})}{\partial x^{(p)}}dx^{(p)},\\
dv_2&=\frac{\partial u_2(x^{(1)})}{\partial x^{(1)}}dx^{(1)}+\frac{\partial u_2(x^{(2)})}{\partial x^{(2)}}dx^{(2)}\cdots+\frac{\partial u_2(x^{(p)})}{\partial x^{(p)}}dx^{(p)},\\[-0.35em]
&\hspace{4em}\cdots\qquad\cdots\qquad\cdots\\[-0.35em]
dv_p&=\frac{\partial u_p(x^{(1)})}{\partial x^{(1)}}dx^{(1)}+\frac{\partial u_p(x^{(2)})}{\partial x^{(2)}}dx^{(2)}\cdots+\frac{\partial u_p(x^{(p)})}{\partial x^{(p)}}dx^{(p)}
\end{aligned}\right.
\]
hold; while at the same time certain initial values \(x^{(1)}=c^{(1)},x^{(2)}=c^{(2)}\ldots x^{(p)}=c^{(p)}\), corresponding to the initial values \(v_1=0,v_2=0\ldots v_p=0\), together with the associated \(s_x\), are given, so that, while the \(v\) grow from the values 0, the \(x\) gradually move away from the \(c\). If we solve equations (1) for the \(dx\), we first obtain the first differential quotients of the \(x\) with respect to the \(u\), and then, by continued differentiation, the higher differential quotients as functions of the \(x\). If in these expressions we substitute for the \(x\) the values \(c\), there result the coefficients for an expansion, carried out by Taylor's series, in ascending powers of the \(v\); this expansion always converges within certain limits and really represents the functions \(x\), provided only that, as is always possible, the \(c\) have been so chosen that all coefficients of the series remain finite; and indeed, as long (cf. Briot and Bouquet, \emph{Théorie des ff. ell.}) as the differential quotients of the \(x\) with respect to the \(u\) remain synectic. But one may also set up an expansion beyond such a domain, only then taking as starting-point one of the value-systems of the \(x\) already found; in this way, for every value-system of the \(v\), one arrives at a convergent expansion, which of course can be different for different value-systems, and in which those value-systems are always avoided for which the coefficients of the solved equations (1) are no longer finite.

The denominator which occurs in solving equations (1) vanishes when two points \(x\) coincide in the course of the expansion. Hence the expansion of the \(x\) by Taylor's theorem ceases to converge as soon as a value-system is reached for which two points \(x\) coincide. This is avoided if we expand a symmetric function of the \(x\). Such a function
\[
\frac{\varphi(x^{(1)},x^{(2)}\ldots x^{(p)})}{\psi(x^{(1)},x^{(2)}\ldots x^{(p)})}
\]
we shall imagine, in order to fix the ideas, as an algebraic function which is simultaneously rational in the \(x\) and in the \(s_x\), and which therefore does not change if two of the \(x\), and at the same time two of the \(s_x\), are interchanged. Then namely
\[
\tag{2}\frac{\partial\frac{\varphi}{\psi}}{\partial v_i}
=\frac{\displaystyle\sum \pm\frac{\partial u_1(x^{(1)})}{\partial x^{(1)}}\frac{\partial u_2(x^{(2)})}{\partial x^{(2)}}\cdots\frac{\partial\frac{\varphi}{\psi}}{\partial x^{(i)}}\cdots\frac{\partial u_p(x^{(p)})}{\partial x^{(p)}}}{\displaystyle\sum \pm\frac{\partial u_1(x^{(1)})}{\partial x^{(1)}}\frac{\partial u_2(x^{(2)})}{\partial x^{(2)}}\cdots\frac{\partial u_i(x^{(i)})}{\partial x^{(i)}}\cdots\frac{\partial u_p(x^{(p)})}{\partial x^{(p)}}}
\]
is the quotient of two algebraic functional determinants taken with respect to the \(x\); each of these generally vanishes only simply at the coincidence of two \(x\), and their quotient therefore remains finite.

The functional determinants in (2), moreover, admit a remarkable transformation. Namely one has
\[
\frac{\partial u_k(x)}{\partial x}=\frac{\Theta_k(x)}{f'(s_x)},
\]
where \(\Theta_k=0\) is the equation of a curve of order \((n-3)\) passing through the double and return points of the curve \(f=0\). If one also sets
\[
\frac{\partial\frac{\varphi}{\psi}}{\partial x^{(i)}}=\frac{N_i}{f'(s_x^{(i)})},
\]
then formula (2) becomes
\[
\tag{3}\frac{\partial\frac{\varphi}{\psi}}{\partial v_i}
=\frac{\sum\pm\Theta_1(x^{(1)})\Theta_2(x^{(2)})\cdots N_i\cdots\Theta_p(x^{(p)})}{\sum\pm\Theta_1(x^{(1)})\Theta_2(x^{(2)})\cdots\Theta_i(x^{(i)})\cdots\Theta_p(x^{(p)})},
\]
or, if the determinant in the denominator is now denoted by \(H\), and its subdeterminants by \(H_i^{(k)}\),
\[
\tag{4}\frac{\partial\frac{\varphi}{\psi}}{\partial v_i}=\frac{H_i^{(1)}N_1+H_i^{(2)}N_2\cdots+H_i^{(p)}N_p}{H}.
\]
The increment of the function \(\frac{\varphi}{\psi}\) is therefore at all times determined by the formula
\[
\tag{5}d\frac{\varphi}{\psi}=\frac{\sum_i\sum_k H_i^{(k)}N_k\,dv_i}{H},
\]
where the coefficients on the right assume, in general, completely determined values for every system of the \(x\).

An interruption of monodromy can thus occur only for such value-systems for which one of the coefficients of \(d\varphi\) becomes infinitely large or indeterminate. This can happen only if either \(H=0\), or if the denominator of the expression
\[
N_k=f'(s_x^{(k)})\cdot\frac{\psi\frac{\partial\varphi}{\partial x^{(k)}}-\varphi\frac{\partial\psi}{\partial x^{(k)}}}{\psi^2}
\]
vanishes, that is, if \(\psi\) vanishes. But in the latter case one need only consider the reciprocal function \(\frac{\psi}{\varphi}\) to recognize that the monodromy is not disturbed. For one has
\[
d\frac{\psi}{\varphi}=\frac{\sum_i\sum_k H_i^{(k)}P_k\,dv_i}{H},
\]
where
\[
P_k=f'(s_x^{(k)})\cdot\frac{\varphi\frac{\partial\psi}{\partial x^{(k)}}-\psi\frac{\partial\varphi}{\partial x^{(k)}}}{\varphi^2};
\]
and the differential of \(\frac{\psi}{\varphi}\) therefore no longer becomes infinite for \(\psi=0\); rather the individual coefficients take determinate finite values. Thus, for these value-systems, \(\frac{\varphi}{\psi}\) becomes only infinite, without there being any question of branching.

It remains only to investigate the behavior of the function for those value-systems of the \(v\) to which there correspond points \(x\) lying, together with the double and return points of \(f=0\), on a curve of order \((n-3)\). For this is in fact the meaning of the equation \(H=0\). The appearance of the determinant denoted by \(H\) teaches that \(H\) is, in the coordinates of each of the \(p\) points, of order \((n-3)\); that it vanishes when one of the \(x\) coincides with another or with a double or return point. If therefore \(p-1\) points \(x\) are regarded as constant and the \(p\)th as variable, then \(H=0\) is the equation of a curve of the kind just indicated, of order \((n-3)\), or \(H=0\) is the condition that all \(p\) points lie on such a curve.

Here there first occurs the noteworthy circumstance that this condition, although it is only one condition for the \(x\), nevertheless involves two conditions for the \(v\). Abel's theorem teaches this. Let \(\alpha^{(1)},\alpha^{(2)},\ldots,\alpha^{(2p-2)}\) denote any \(2p-2\) points of \(f=0\) which, with the double and return points, lie on a curve of order \((n-3)\), and let \(x^{(p+1)},x^{(p+2)}\ldots x^{(2p-2)}\) be the points in which, in such a case, the curve of order \((n-3)\) determined by the remaining \(x\) still cuts our curve \(f=0\). Then the \(\alpha,x\) can be paired and the paths of integration chosen so that
\[
\sum_{i=1}^{i=2p-2}\int_{\alpha^{(i)}}^{x^{(i)}}du_k=0\qquad (k=1,2\ldots p).
\]
And therefore, by a suitable choice of the paths of integration, one may also establish the equations
\[
\sum_{i=1}^{i=2p-2}\int_{\alpha^{(i)}}^{x^{(p+i)}}du_k=-v_k+\sum_{i=1}^{i=p}\int_{c^{(i)}}^{\alpha^{(i)}}du_k\qquad(k=1,2\ldots p).
\]
Since on the right next to \(-v_k\) there stands a constant, and on the left the lower limits are likewise constant, the \(p\) quantities \(v_k\) depend only on the \(p-2\) points \(x^{(p+1)},x^{(p+2)},\ldots,x^{(2p-2)}\); they are therefore functions of \(p-2\) variables, and consequently two equations must hold among the \(v\) in order that the one condition among the \(x\) may occur, namely that they lie on a curve of order \((n-3)\) drawn through the double and return points.

The same fact appears in another way through the following consideration. Suppose we have found a system of the \(x\) for which \(H=0\), and know the corresponding system of the \(v\). From the equations
\[
dv_k=\frac{\partial u_k(x^{(1)})}{\partial x^{(1)}}dx^{(1)}+\frac{\partial u_k(x^{(2)})}{\partial x^{(2)}}dx^{(2)}\cdots+\frac{\partial u_k(x^{(p)})}{\partial x^{(p)}}dx^{(p)}
\]
it follows, by multiplication with the subdeterminants \(H_i^{(k)}\) and addition, that
\[
H_i^{(1)}dv_1+H_i^{(2)}dv_2\cdots+H_i^{(p)}dv_p=0\qquad(i=1,2\ldots p).
\]
In truth, however, this equation represents not \(p\) different equations, but only a single one; for since \(H\) vanishes, one can always determine such rows of quantities \(a,b\) that the subdeterminants assume the values
\[
H_i^{(k)}=a_i b_k,
\]
whereby the above system of equations reduces to the single equation
\[
\tag{6} b_1dv_1+b_2dv_2\cdots+b_pdv_p=0.
\]
Thus one sees that infinitely many systems of the \(x\) must satisfy a system of the \(v\) which satisfies the two condition equations remarked above, in such a way that the choice of one of these systems of the \(x\), by which the quantities \(b\) acquire determinate values, involves a condition (6) for the increments which the \(v\) are to receive.

The occurrence of such value-systems is in fact a general property of functions of several variables. If such a function is set equal to zero, one obtains one condition among the arguments; if it is set equal to infinity, one obtains a second; the concurrence of both conditions yields value-systems for which the function is indeterminate, where therefore infinitely many function-values correspond to one system of arguments.

If \(v_1,v_2,\ldots,v_p\) are the arguments and \(f=\frac{\varphi}{\psi}\) is the function to be investigated, then the equations \(\varphi=0,\psi=0\) together give the values for which \(f\) becomes indeterminate. If, however, in \(\varphi,\psi\) one introduces for the \(v\) such values \(v+\delta v\) as differ infinitely little from those which make \(\varphi,\psi\) vanish, there results
\[
f=\frac{\sum\frac{\partial\varphi}{\partial v_i}\delta v_i}{\sum\frac{\partial\psi}{\partial v_i}\delta v_i},
\]
so that to each determinate one among the infinitely many function-values belonging to the system of the \(v\), there corresponds a linear relation of the increments \(\delta v\). If we call the entire system of quantities through which the \(p\) quantities \(v\) can run independently of one another a place of \(p\)th dimension, then those value-systems for which \(f\) is indeterminate form a place of \((p-2)\)th dimension. This shows that the monodromy of a function is not necessarily abolished because it assumes infinitely many values for value-systems which form a place of dimension \((p-2)\); that therefore a function otherwise completely and uniquely determined can become infinitely many-valued in a place of dimension \((p-2)\), without one being justified in denying it uniqueness. Thus already the function \(z=\frac{x}{y}\) is completely indeterminate for \(x=0,y=0\), although it is entirely monodromic. But the different values of \(z\) which enter for \(x=0,y=0\) arrange themselves so that to each special value there corresponds a determinate progression
\[
\delta x=z\delta y,
\]
and hence one arrives at a wholly determinate value of \(z\) when \(x\) and \(y\) are made to approach the common value zero in a determinate way.

One must therefore define the concept of monodromy for a function of several variables with the necessary caution. One must not say that a monodromic function is one which possesses only one value for all value-systems of the variables; rather, it already suffices in general that the function be single-valued everywhere, with the exception of a place of \((p-2)\)th dimension in which it can assume every value.

That the symmetric function \(\frac{\varphi}{\psi}\) in our case really assumes the form \(\frac00\) follows from the fact that, by what precedes, it can be represented as a quotient of functions \(V\) which all vanish as soon as the \(x\) lie on one of the curves of order \((n-3)\) under consideration.

Let us therefore be content for the moment with having shown that \(\frac{\varphi}{\psi}\), outside this place of \((p-2)\)th dimension, is a monodromic function of the \(v\), and let us now pass to the investigation of the function \(V\) itself, from which we shall also learn to characterize the function \(\frac{\varphi}{\psi}\) at these exceptional places.

\subsubsection*{\S\,53. Proof that the function \(V\) is a synectic function of the \(v\).}

The function \(V(x^{(1)},x^{(2)}\ldots x^{(p)};\xi)\) has, as proved above, the property of always remaining finite. The differential quotients of its logarithm with respect to the \(x\) (and therefore also with respect to the \(v\)) are by no means algebraic functions; for since, in the integrals of the third kind occurring in \(-U\), parameter and argument are variable, the differential quotients of such an integral are composed of algebraic parts and of the periodicity moduli \(S^{(k)}\) of the incompletely normalized integrals of the third kind. But all these terms have the property, partly of assuming completely determinate values for given points \(x\), and partly at least of being fully determinate when the variable points \(x\), as is the case in the inversion problem, have arisen on determinate paths from given constant points. They have in addition the property of approaching the logarithm of a simply vanishing function, without further factor, where \(\log V\) becomes infinitely large. Hence the function \(V\) is a monodromic function of the \(v\) to the same extent as this holds for the \(x\) themselves. Therefore a disturbance of the monodromy of \(V\) can appear possible only in the following cases:

\begin{center}
\begin{tabular}{@{}p{5.15in}@{}}
1) When two points \(x\) coincide;\\
2) When the points \(x\), with the double and return points of \(f=0\), lie on a curve of order \((n-3)\).
\end{tabular}
\end{center}

If, in considering case 1), we replace the function \(V\), as is permitted by \S\,51, by the function identical with it,
\[
\varepsilon V(y^{(1)},y^{(2)}\ldots y^{(p)};\eta)\qquad(\varepsilon=\pm1),
\]
where the \(y,x\), together with the remaining points of the line \(\xi\eta\) and with the double and return points, lie on a curve of order \((n-2)\), then the \(y\) correspond to another inversion problem in which, up to additive constants, the \(v\) are replaced by \(-v\). The \(y\) are in no position subject to cases 1) or 2); at the place in question they are monodromic functions of the \(v\), and consequently so is \(V\) itself; and case 1) therefore does not ground an interruption of monodromy.

The same auxiliary device also serves for treating the second case. In this case the \(x\) are in fact no longer determinate functions of the \(v\), but the \(y\) still are; the curve of order \((n-2)\) decomposes into the line \(\xi\eta\) and into a curve of order \((n-3)\) which contains the \(x\) and \(p-2\) other points. These, together with the points \(\xi,\eta\) themselves, form the system of the \(y\); hence in the second form of \(V\) one argument coincides with the parameter. In this case, as we know, \(\log V\) behaves like the logarithm of a simply vanishing algebraic function, and the function remains monodromic in the neighborhood of such a value-system. Thus case 2) also cannot ground a disturbance of monodromy. We therefore have the theorem:

\textbf{The function \(V(x^{(1)},x^{(2)}\ldots x^{(p)};\xi)\) is always finite and monodromic, hence a synectic function of the \(v\).}

Thus the inversion problem has been reduced to the consideration of synectic functions. The point will now be to give the actual representation of these functions.

Already at the end of the seventh section it was proved that the function \(V\), depending on \(p+1\) points, contains them only in the \(p\) combinations
\[
\sum_{i=1}^{i=p}\int_{c^{(i)}}^{x^{(i)}}du_h-\int_{\mu}^{\xi}du_h.
\]
We therefore wish to regard it as a function of the \(p\) variables
\[
w_h=\sum_{i=1}^{i=p}\int_{c^{(i)}}^{x^{(i)}}du_h-\int_{\mu}^{\xi}du_h+K_h,
\]
where the \(K_h\) are constants to be determined presently.

If we now replace the \(x\) by the points \(y\), and \(\xi\) by \(\eta\), then \(V\), apart from the sign, remains unchanged; that is, it is, apart from the sign, the same function of the quantities
\[
w'_h=\sum_{i=1}^{i=p}\int_{c^{(i)}}^{y^{(i)}}du_h-\int_{\mu}^{\eta}du_h+K_h.
\]
Moreover between \(w,w'\) there is the relation that \(w_h+w'_h\) is a constant, as Abel's theorem immediately teaches. We can therefore always determine the constants \(K_h\) so that this sum becomes zero. With such a determination, the \(w'\) become equal and opposite to the \(w\), and one therefore has the theorem:

\textbf{The function \(V\), regarded as a function of the \(w\), does not change, or at most changes its sign, if all the \(w\) are simultaneously allowed to change their signs.}

The constants \(K\) are determined from the equations
\[
2K_h+\sum_{i=1}^{i=p}\int_{c^{(i)}}^{x^{(i)}}du_h+\sum_{i=1}^{i=p}\int_{c^{(i)}}^{y^{(i)}}du_h-\int_{\mu}^{\xi}du_h-\int_{\mu}^{\eta}du_h=0,
\]
if one therein sets for \(\xi,\eta\) any two arbitrarily chosen points, and for the \(x,y\) any \(2p\) constant points which, together with the remaining points of the line joining the former two and with the double and return points of \(f=0\), lie on a curve of order \((n-2)\). We shall have to deal with this subject in greater detail below.

\subsubsection*{\S\,54. Periodicity moduli of \(-U=\log V\).}

We now investigate in what way the function \(V\), which depends only on the \(w\), changes when the \(w\) increase by the most general periodicity moduli of integrals of the first kind, thus when \(w_i\) passes into
\[
w_i'=w_i+2m_i\pi\sqrt{-1}+a_{1i}q_1+a_{2i}q_2\cdots+a_{pi}q_p,
\]
where the \(m,q\) are arbitrary positive or negative integers. The function which arises from \(V\) depends only on the \(w_i'\); but these themselves differ from the \(w\) only by constants, and hence the function \(V'\) arising from \(V\) is also only a function of the \(w\). The same is consequently true of the function
\[
\psi(w_1,w_2\ldots w_p)=\log V'-\log V,
\]
which gives the simultaneous increment of \(\log V\), or of the function \(-U\). We shall now determine this function \(\psi\).

We start from the formula (\S\,50):
\[
\tag{1}T_{\xi\eta}\binom{x^{(1)},x^{(2)}\ldots x^{(p)}}{z^{(1)},z^{(2)}\ldots z^{(p)}}
=\log\frac{V(x^{(1)},x^{(2)}\ldots x^{(p)};\eta)V(z^{(1)},z^{(2)}\ldots z^{(p)};\xi)}{V(x^{(1)},x^{(2)}\ldots x^{(p)};\xi)V(z^{(1)},z^{(2)}\ldots z^{(p)};\eta)},
\]
in which four functions \(V\), and therefore correspondingly four systems of \(w\), occur, namely
\[
\begin{aligned}
w_h&=\sum\int_{c^{(i)}}^{x^{(i)}}du_h-\int_{\mu}^{\xi}du_h+K_h,\\
w_h^{\circ}&=\sum\int_{c^{(i)}}^{z^{(i)}}du_h-\int_{\mu}^{\xi}du_h+K_h,\\
W_h&=\sum\int_{c^{(i)}}^{x^{(i)}}du_h-\int_{\mu}^{\eta}du_h+K_h,\\
W_h^{\circ}&=\sum\int_{c^{(i)}}^{z^{(i)}}du_h-\int_{\mu}^{\eta}du_h+K_h.
\end{aligned}
\]
If now a point \(x\) traverses such a period path whose addition makes \(w_h\) pass into \(w'_h\), then at the same time \(W_h\) passes into \(W'_h\), and \(T_{\xi\eta}\binom{x}{z}\) increases by (cf. \S\,41)
\[
\int_{\eta}^{\xi}dI=\int_{\eta}^{\xi}(q_1du_1+q_2du_2\cdots+q_pdu_p),
\]
while the left side of equation (1) increases by
\[
\psi(W_1,W_2\ldots)-\psi(w_1,w_2\ldots).
\]
We therefore have the equation
\[
\tag{2}\psi(W_1,W_2\ldots)-\psi(w_1,w_2\ldots)=\int_{\eta}^{\xi}dI.
\]
Second, if \(\eta\) describes the same period path, but in the opposite direction, then the \(W_h\) again pass into the \(W'_h\), while at the same time the \(W_h^{\circ}\) pass into \(W_h^{\circ\prime}\); the right side of (1) thus increases by
\[
\psi(W_1,W_2,
\ldots)-\psi(W_1^{\circ},W_2^{\circ},\ldots),
\]
while the left side increases by
\[
\sum_{i=1}^{i=p}\int_{z^{(i)}}^{x^{(i)}}dI.
\]
Thus one obtains, secondly, the equation
\[
\tag{3}\psi(W_1,W_2\ldots)-\psi(W_1^{\circ},W_2^{\circ}\ldots)=\sum_{i=1}^{i=p}\int_{z^{(i)}}^{x^{(i)}}dI.
\]
Subtracting (2) from (3), one finds
\[
\psi(w_1,w_2\ldots)-\psi(W_1^{\circ},W_2^{\circ}\ldots)=\sum_{i=1}^{i=p}\int_{z^{(i)}}^{x^{(i)}}dI-\int_{\eta}^{\xi}dI.
\]
If in place of \(z,\eta\) one sets here the constant points \(c,\mu\), then \(\psi(W_1^{\circ},W_2^{\circ}\ldots)\) becomes a constant quantity, and the expression for \(\psi(w_1,w_2\ldots)\) can be written in the form
\[
\tag{4}\psi(w_1,w_2\ldots)=q_1w_1+q_2w_2\cdots+q_pw_p+C.
\]
The constant \(C\), which can now depend only on the numbers \(m,q\), remains to be determined.

If, with the aid of equation (4), we express the function \(V'\) by \(V\), we have
\[
\tag{5}V'=e^{w_1q_1+w_2q_2\cdots+w_pq_p+C}\cdot V.
\]
Since \(V'\) is the same function of the \(w'\) as \(V\) is of the \(w\), \(V'\) passes into \(\varepsilon V\) \((\varepsilon=\pm1)\) when, in place of the quantities \(w'_h\), one sets the quantities
\[
-\{w_h+2m_h\pi\sqrt{-1}+q_1a_{1h}+q_2a_{2h}\cdots+q_pa_{ph}\};
\]
for in this way the \(w'\) are transformed into the \(-w\), and hence \(V'\) into \(\varepsilon V\), since this function remains unchanged up to its sign under a change of sign of all \(w\). But for the same reason \(V\) at the same time passes into \(\varepsilon V'\), and from (5) one therefore obtains the second equation:
\[
\tag{6}V=e^{-\sum q_hw_h-\sum\sum q_hq_ka_{hk}+C}\cdot V';
\]
hence, if (5) is multiplied by (6), and the equal factors are cancelled on both sides,
\[
1=e^{-\sum\sum q_hq_ka_{hk}+2C}.
\]
Consequently
\[
\tag{7}C=\frac12\sum\sum q_hq_ka_{hk}+L\pi\sqrt{-1},
\]
where \(L\) is an integer. It remains finally to determine this integer.

For this purpose let us imagine that to the above increment of the \(w\) there is added a second increment
\[
2m'_h\pi\sqrt{-1}+q'_1a_{1h}+q'_2a_{2h}\cdots+q'_pa_{ph}.
\]
The function
\[
V'=e^{\sum w_hq_h+\frac12\sum\sum q_hq_ka_{hk}+L\pi\sqrt{-1}}\cdot V
\]
is thereby transformed into
\[
e^{\sum w_h(q_h+q'_h)+\frac12\sum(q_h+q'_h)(q_k+q'_k)a_{hk}+(L+L')\pi\sqrt{-1}}\cdot V,
\]
where \(L'\) depends on the numbers \(q',m'\) just as \(L\) depends on the numbers \(q,m\). The same result must be obtained if in \(V\) one directly increases the \(w_h\) by the quantities
\[
2(m_h+m'_h)\pi\sqrt{-1}+a_{1h}(q_1+q'_1)+a_{2h}(q_2+q'_2)\cdots.
\]
In this way, however, \(V\) passes into
\[
e^{\sum w_h(q_h+q'_h)+\frac12\sum(q_h+q'_h)(q_k+q'_k)a_{hk}+L''\pi\sqrt{-1}}\cdot V,
\]
where \(L''\) depends on the sums \(m+m',q+q'\) as \(L,L'\) depend on the individual number series.

The comparison gives, since even multiples of \(\pi\sqrt{-1}\) may be disregarded,
\[
L+L'=L''.
\]
We can therefore in general set the number \(L\) belonging to a sum of arbitrarily many number series equal to the sum of all numbers \(L\) which belong to the individual number series. Let us now start from those number series in which in each case only one number \(m\), or only one number \(q\), is different from zero and is equal to 1. From these we can compose the number series \(m_1,m_2,
\ldots,q_1,q_2\ldots\), by repeating the first \(m_1\) times, the second \(m_2\) times, etc. If we denote by
\[
\varkappa_1,\varkappa_2\ldots\varkappa_p,\lambda_1,\lambda_2\ldots\lambda_p
\]
the values of \(L\) which correspond to the first \(2p\) number series (and indeed we need assign to these numbers only the values 0 or 1, as is also the case for the \(L\)), then for \(L\) we have the expression
\[
L=\varkappa_1m_1+\varkappa_2m_2\cdots+\varkappa_pm_p+\lambda_1q_1+\lambda_2q_2\cdots+\lambda_pq_p.
\]
And so we finally have the theorem:

\textbf{If the \(w_h\) increase by the expressions}
\[
2m_h\pi\sqrt{-1}+\sum_k a_{hk}q_k,
\]
\textbf{then the function \(V\) passes into}
\[
V\cdot e^{\sum w_hq_h+\frac12\sum a_{hk}q_hq_k+(\sum\varkappa_hm_h+\sum\lambda_hq_h)\pi\sqrt{-1}},
\]
\textbf{where the \(\varkappa,\lambda\) are integers to which one may assign the values 0 or 1.}

Thus we have carried the determination of the periodicity modulus of \(\log V\) so far that only the \(2p\) integers \(\varkappa,\lambda\) still remain undetermined. These are defined by definite integrals; they are the periodicity moduli of \(\log V\), diminished by the components of the same already found otherwise.

\subsubsection*{\S\,55. Series expansion.}

In the formula which gives the change of \(V\) under a change of the \(w\) by periodicity moduli, let us put in place of the \(w\) the quantities enlarged by half-periodicity moduli,
\[
\tag{1}\omega_h=w_h-\lambda_h\pi\sqrt{-1}-\frac12(\varkappa_1a_{1h}+\varkappa_2a_{2h}\cdots+\varkappa_pa_{ph}).
\]
By multiplying the equation
\[
V'=V\cdot e^{\sum w_hq_h+\frac12\sum q_hq_ka_{hk}+\sum(\varkappa_hm_h+\lambda_hq_h)\pi\sqrt{-1}}
\]
by
\[
e^{-\frac12\sum\varkappa_h\omega_h-\sum\varkappa_hm_h\pi\sqrt{-1}-\frac12\sum\varkappa_hq_ka_{hk}},
\]
we find the equation
\[
\begin{aligned}
&e^{-\frac12\sum\varkappa_h(\omega_h+2m_h\pi\sqrt{-1}+a_{1h}q_1+a_{2h}q_2\cdots)}\cdot V'\\
&\qquad=e^{-\frac12\sum\varkappa_h\omega_h}\cdot V\cdot e^{\sum q_h\omega_h+\frac12\sum q_hq_ka_{hk}}.
\end{aligned}
\]
The function
\[
e^{-\frac12\sum\varkappa_h\omega_h}\cdot V
\]
therefore does not change if the \(\omega\) are increased by multiples of \(2\pi\sqrt{-1}\), and changes only by the factor
\[
e^{\sum q_h\omega_h+\frac12\sum q_hq_ka_{hk}}
\]
if the \(\omega_h\) are increased by the expressions \(q_1a_{1h}+q_2a_{2h}\cdots\).

From the first circumstance it follows that this function can be expanded in an always convergent series in positive and negative powers of the \(p\) quantities \(e^{\omega_h}\), so that one may set
\[
e^{-\frac12\sum\varkappa_h\omega_h}\cdot V=\sum A_{r_1r_2\ldots r_p}e^{r_1\omega_1+r_2\omega_2\cdots+r_p\omega_p},
\]
where the sum on the right must be extended over all values of the \(r\) from \(-\infty\) to \(+\infty\).

The application of the second part of the preceding theorem then leads to the formula
\[
\sum A_{r_1r_2\ldots r_p}e^{\sum r_h\omega_h+\sum r_hq_ka_{hk}}
=e^{\sum q_h\omega_h+\frac12\sum q_hq_ka_{hk}}\sum A_{r_1r_2\ldots r_p}e^{\sum r_h\omega_h}.
\]
If on the left one substitutes \(r_1+q_1,r_2+q_2\ldots\) for \(r_1,r_2\ldots\), which changes nothing since the sum is to be taken for the \(r\) from \(-\infty\) to \(+\infty\), then one obtains
\[
\begin{aligned}
&\sum A_{r_1+q_1,r_2+q_2\ldots}e^{\sum(r_h+q_h)\omega_h+\sum r_hq_ka_{hk}+\sum q_hq_ka_{hk}}\\
&\qquad=e^{\frac12\sum q_hq_ka_{hk}}\sum A_{r_1r_2\ldots}e^{\sum(r_h+q_h)\omega_h}.
\end{aligned}
\]
If one now compares the coefficients of equal exponentials, one finds
\[
A_{r_1+q_1,r_2+q_2\ldots}=e^{-\sum r_hq_ka_{hk}-\frac12\sum q_hq_ka_{hk}}\cdot A_{r_1r_2\ldots},
\]
or, setting all \(r\) equal to zero,
\[
A_{q_1q_2\ldots q_p}=e^{-\frac12\sum q_hq_ka_{hk}}A_{000\ldots}.
\]
Thus the preceding periodicity conditions determine all coefficients of the series expansion completely, except for one factor common to all, and one has
\[
e^{-\frac12\sum\varkappa_h\omega_h}\cdot V=A_{00\ldots}\sum e^{-\frac12\sum r_hr_ka_{hk}+\sum r_h\omega_h}.
\]
We now denote by \(\Theta(\omega_1,\omega_2\ldots\omega_p)\) the function
\[
\tag{2}\Theta(\omega_1,\omega_2\ldots\omega_p)=\sum e^{-\frac12\sum r_hr_ka_{hk}+\sum r_h\omega_h},
\]
so that
\[
V=A_{00\ldots}\cdot e^{\frac12\sum\varkappa_h\omega_h}\cdot\Theta(\omega_1,\omega_2\ldots\omega_p).
\]
The constant \(A_{00\ldots}\) is easy to determine. For by its origin the function \(U\) became equal to zero, and thus \(V=1\), when the \(x,\xi\) passed into arbitrary initial points \(\gamma,\xi\). If \(\omega_h^{\circ}\) denotes what becomes of \(\omega_h\) under this passage, then
\[
\tag{3}V=e^{\frac12\sum\varkappa_h(\omega_h-\omega_h^{\circ})}\cdot\frac{\Theta(\omega_1,\omega_2\ldots\omega_p)}{\Theta(\omega_1^{\circ},\omega_2^{\circ}\ldots\omega_p^{\circ})}.
\]
This series for \(V\) is, by what precedes, always convergent. The function \(\Theta\), therefore, must likewise possess the property of converging for all finite values of the \(\omega\), in consequence of the values of the \(a_{hk}\) which occur in these developments.

\subsubsection*{\S\,56. The arguments of the function \(\Theta\).}

We now turn to the closer consideration of those quantities \(\omega\) which form the arguments of the function \(\Theta\). They arose by subtracting half-periods from the \(w\), so that we can write
\[
\tag{1}\omega_h=w_h-\frac12P_h^{\varkappa,\lambda}.
\]


when
\[
P_h^{\varkappa,\lambda}=2\lambda_h\pi\sqrt{-1}+\varkappa_1a_{1h}+\varkappa_2a_{2h}\cdots+\varkappa_pa_{ph}
\]
is put. But the arguments \(w\) arose from the \(v\) with the aid of the equation
\[
w_h=v_h-\int_{\mu}^{\xi}du_h+K_h
=\sum_{i=1}^{i=p}\int_{c^{(i)}}^{x^{(i)}}du_h-\int_{\mu}^{\xi}du_h+K_h.
\]
Finally, the constants \(K_h\) are given by
\[
2K=\int_{\mu}^{\xi}du_h+\int_{\mu}^{\eta}du_h
-\sum_{i=1}^{i=p}\int_{c^{(i)}}^{x^{(i)}}du_h
-\sum_{i=1}^{i=p}\int_{c^{(i)}}^{y^{(i)}}du_h,
\]
where the \(x,y\), with the double and return points and with the remaining points of the line \(\xi\eta\), were to lie on a curve of order \((n-2)\). Let this curve now pass into the contact curve mentioned in \(\S\,47\), in which \(\xi,\eta\) fall together into a point \(\mu\), and the \(x\) together with the \(y\) fall pairwise into points \(\varepsilon\); points connected with the points \(\xi,\eta,c,\mu,x,y\) by paths which can be transformed into the integration paths occurring here. Then we have
\[
K_h=-\sum_{i=1}^{i=p}\int_{c^{(i)}}^{\varepsilon^{(i)}}du_h,
\]
and we can therefore write
\[
\tag{2}w_h=\sum_{i=1}^{i=p}\int_{\varepsilon^{(i)}}^{x^{(i)}}du_h-\int_{\mu}^{\xi}du_h.
\]

If we hold the point \(\mu\) fixed, the problem of determining the points \(\varepsilon\) is by no means uniquely soluble. The points \(\varepsilon\) are the contact points of a curve of order \((n-2)\), which passes through the double and return points, and through the other intersection points of the tangent to \(f=0\) drawn at the point \(\mu\). We shall investigate how many such contact curves there are. The \(\varepsilon\) belonged to an arbitrarily chosen contact curve; the contact points of another of these curves we denote by the points \(\alpha\). Then, by Abel's theorem, the equations hold
\[
2\sum_{i=1}^{i=p}\int_{\varepsilon^{(i)}}^{\alpha^{(i)}}du_h=0,
\]
where, if the paths from the \(\varepsilon\) to the \(\alpha\) are thought of arbitrarily, an arbitrary period modulus of an integral of the first kind may be added on the right. Thus one can define the \(\alpha\) from the equation
\[
\tag{3}\sum_{i=1}^{i=p}\int_{\varepsilon^{(i)}}^{\alpha^{(i)}}du_h=\frac12P_h^{\varkappa,\lambda}.
\]
To each system of numbers \(\varkappa,\lambda\) there corresponds a system of the \(\alpha\), hence a contact curve, and also only a single one. Among these different systems, one is the system of the \(\varepsilon\) themselves, namely the one for which all \(\varkappa,\lambda\) vanish. As often as two half-periods \(\frac12P_h^{\varkappa,\lambda}\) do not differ by whole periods, the corresponding systems of the \(\alpha\) must also be distinct from one another; there are therefore as many contact curves as systems of numbers \(\varkappa,\lambda\) can be formed from the numbers 0 and 1, hence \(2^{2p}\). One of these curves is the one originally considered; a second is that whose numbers \(\varkappa,\lambda\) agree with those occurring in the \(\omega\). We shall now preferably denote the contact points of the latter as the points \(\alpha\).

From equations (1), (2), (3) together the values of the \(\omega\) immediately result in the form
\[
\omega_h=\sum_{i=1}^{i=p}\int_{\alpha^{(i)}}^{x^{(i)}}du_h-\int_{\mu}^{\xi}du_h.
\]
And one has the theorem:

The arguments \(\omega\) of the function \(\Theta\) are the integral sums, differing from the \(v\) by constants,
\[
\omega_h=\sum_{i=1}^{i=p}\int_{\alpha^{(i)}}^{x^{(i)}}du_h-\int_{\mu}^{\xi}du_h,
\]
where the \(\alpha\) are the contact points of a definite one among the \(2^{2p}\) curves of order \((n-2)\), which pass through the double and return points, as well as through the \(n-2\) other intersection points of the tangent drawn at \(\mu\) to \(f=0\), and which also touch the curve \(f=0\) at \(p\) points.

It is to be observed that the \(2^{2p}\) curves corresponding to different points \(\mu\) are assigned to one another in a definite way, so that, when one moves \(\mu\) and at the same time lets the \(\alpha\) vary continuously, the value of \(\omega\) undergoes no modification. For this purpose, at another point \(\nu\), instead of the \(\alpha\) one need only choose those points \(\beta\) for which
\[
\sum_{i=1}^{i=p}\int_{\alpha^{(i)}}^{\beta^{(i)}}du_h-\int_{\mu}^{\nu}du_h=0.
\]
If one chooses another system of the \(\beta\), then half a period modulus is added on the right.

\subsubsection*{\S\,57. Solution of the inversion problem by \(\Theta\)-functions.}

We shall now investigate how the solution of the inversion problem is represented with the aid of the synectic function \(\Theta\). From the equation
\[
V(x^{(1)},x^{(2)}\ldots x^{(p)};\xi)
=e^{\frac12\sum\varkappa_h(\omega_h-\omega_h^\circ)}
\cdot\frac{\Theta(\omega_1,\omega_2,\ldots)}
{\Theta(\omega_1^\circ,\omega_2^\circ,\ldots)}
\]
there first results, for the special transcendent \(\frac12T_{\xi\eta}(x)\), the expression
\[
\begin{aligned}
\frac12T_{\xi\eta}(x)
&=\log\frac{V(x^{(1)},x^{(2)}\ldots x^{(p)};\eta)}
{V(x^{(1)},x^{(2)}\ldots x^{(p)};\xi)}\\
&=\frac12\sum\varkappa_h\int_{\eta}^{\xi}du_h
+\log\frac{\Theta\!\left(\sum_{i=1}^{i=p}\int_{\alpha^{(i)}}^{x^{(i)}}du_h-\int_{\mu}^{\eta}du_h\right)}
{\Theta\!\left(\sum_{i=1}^{i=p}\int_{\alpha^{(i)}}^{x^{(i)}}du_h-\int_{\mu}^{\xi}du_h\right)},
\end{aligned}
\]
where, in the \(\Theta\)-functions, only one argument is always put, but written out in full. Then, however, for the general transcendent one finds the formula
\[
\tag{1}
\frac12T_{\xi\eta}\binom{x}{z}
=\frac12T_{\xi\eta}(x)-\frac12T_{\xi\eta}(z)
=\log\frac{
\Theta\!\left(\sum_{i=1}^{i=p}\int_{\alpha^{(i)}}^{x^{(i)}}du_h-\int_{\mu}^{\eta}du_h\right)
\Theta\!\left(\sum_{i=1}^{i=p}\int_{\alpha^{(i)}}^{z^{(i)}}du_h-\int_{\mu}^{\xi}du_h\right)}
{\Theta\!\left(\sum_{i=1}^{i=p}\int_{\alpha^{(i)}}^{x^{(i)}}du_h-\int_{\mu}^{\xi}du_h\right)
\Theta\!\left(\sum_{i=1}^{i=p}\int_{\alpha^{(i)}}^{z^{(i)}}du_h-\int_{\mu}^{\eta}du_h\right)}.
\]
And for the inversion problem one has the formula
\[
\tag{2}
\prod_{i=1}^{i=p}
\frac{\left(\frac{\varphi}{\psi}\right)_{x^{(i)}}-\lambda}
{\left(\frac{\varphi}{\psi}\right)_{c^{(i)}}-\lambda}
=
\prod_{\xi\eta}
\frac{
\Theta\!\left(\sum_{i=1}^{i=p}\int_{\alpha^{(i)}}^{x^{(i)}}du_h-\int_{\mu}^{\xi}du_h\right)
\Theta\!\left(\sum_{i=1}^{i=p}\int_{\alpha^{(i)}}^{c^{(i)}}du_h-\int_{\mu}^{\eta}du_h\right)}
{\Theta\!\left(\sum_{i=1}^{i=p}\int_{\alpha^{(i)}}^{x^{(i)}}du_h-\int_{\mu}^{\eta}du_h\right)
\Theta\!\left(\sum_{i=1}^{i=p}\int_{\alpha^{(i)}}^{c^{(i)}}du_h-\int_{\mu}^{\xi}du_h\right)}.
\]
One sees that the two last formulas, which belong essentially to the inversion problem, have become completely free of the numbers \(\varkappa\), which arose from the accidental arrangement of the transcendent \(\frac12T_{\xi\eta}(x)\) and the function \(V\). But these equations lead naturally to reducing the inversion problem, which was set out in the sixth section, to the more special inversion problem in which the \(c\) are not arbitrarily given points, but precisely the point system of the \(\alpha\) which belongs to a given point \(\mu\), while we can regard the latter as arbitrarily and numerically given.

We then only have, in order to return to the earlier inversion problem, to increase the new values of the \(v\) by constant quantities.

Thus put
\[
\sum_{i=1}^{i=p}\int_{\alpha^{(i)}}^{x^{(i)}}du_h=v_h\qquad(h=1,2\ldots p),
\]
then the solution of the inversion problem is now given in the formula
\[
\tag{3}
\prod_{i=1}^{i=p}\frac{\left(\frac{\varphi}{\psi}\right)_{x^{(i)}}-\lambda}
{\left(\frac{\varphi}{\psi}\right)_{\alpha^{(i)}}-\lambda}
=
\prod_{\xi\eta}
\frac{\Theta\!\left(v_h-\int_{\mu}^{\xi}du_h\right)\Theta\!\left(\int_{\mu}^{\eta}du_h\right)}
{\Theta\!\left(v_h-\int_{\mu}^{\eta}du_h\right)\Theta\!\left(\int_{\mu}^{\xi}du_h\right)}.
\]
At the same time, for the transcendent \(T_{\xi\eta}\) we obtain the expression
\[
T_{\xi\eta}\binom{x}{\alpha}
=\log
\frac{\Theta\!\left(v_h-\int_{\mu}^{\eta}du_h\right)\Theta\!\left(\int_{\mu}^{\xi}du_h\right)}
{\Theta\!\left(v_h-\int_{\mu}^{\xi}du_h\right)\Theta\!\left(\int_{\mu}^{\eta}du_h\right)}.
\]
We can now denote this transcendent as a function of the \(v\) by \(T_{\xi\eta}(v)\), so that one has the defining equation
\[
\tag{4}
T_{\xi\eta}(v)=\log
\frac{\Theta\!\left(v_h-\int_{\mu}^{\eta}du_h\right)\Theta\!\left(\int_{\mu}^{\xi}du_h\right)}
{\Theta\!\left(v_h-\int_{\mu}^{\xi}du_h\right)\Theta\!\left(\int_{\mu}^{\eta}du_h\right)}.
\]

From the properties of the function \(\Theta\) one immediately sees that the functions represented by equations (3), (4) really possess those periodicity properties which were assigned to them above in the sixth section. From the property that, if
\[
v'_h=v_h+2m_h\pi\sqrt{-1}+a_{1h}q_1+a_{2h}q_2\cdots+a_{ph}q_h
\qquad(h=1,2\ldots p)
\]
is put,
\[
\Theta(v')=e^{\sum q_hv_h+\frac12\sum q_hq_ka_{hk}}\Theta(v),
\]
it follows, first, that the functions (3) do not change in any way when the integrals occurring in them are increased by period moduli, and hence that they are indeed \(2p\)-fold periodic functions. For the function \(T_{\xi\eta}(v)\), however, one obtains the equation
\[
T_{\xi\eta}(v')=T_{\xi\eta}(v)+\sum q_h\int_{\eta}^{\xi}du_h,
\]
which is the periodicity property stated in the sixth section.

\subsubsection*{\S\,58. An integral of the third kind and the cross-ratio \(\left(\frac{\varphi}{\psi}\right)_{\xi}:\left(\frac{\varphi}{\psi}\right)_{\eta}\) represented by \(\Theta\)-functions.}

From formulas (1), (2) one can derive some other interesting equations. If in (1) the \(x,z\) are made to coincide pairwise, so that only a single \(x\) differs from the corresponding \(z\), for instance \(x^{(1)}\) from \(z^{(1)}\). If we then put
\[
\int_{\alpha^{(1)}}^{z^{(1)}}du_h+\int_{\alpha^{(2)}}^{x^{(2)}}du_h+\int_{\alpha^{(3)}}^{x^{(3)}}du_h\cdots+\int_{\alpha^{(p)}}^{x^{(p)}}du_h=m_h,
\]
then one finds a single integral of the third kind expressed by \(\Theta\)-functions with the aid of the formula
\[
\tag{5}
\int_{z^{(1)}}^{x^{(1)}}d\Pi_{\xi\eta}
=\log\frac{
\Theta\!\left(\int_{\eta}^{x^{(1)}}du_h+m_h\right)
\Theta\!\left(\int_{\xi}^{z^{(1)}}du_h+m_h\right)}
{\Theta\!\left(\int_{\xi}^{x^{(1)}}du_h+m_h\right)
\Theta\!\left(\int_{\eta}^{z^{(1)}}du_h+m_h\right)}.
\]
And likewise, from formula (2), if in it the \(c\) are replaced by \(z\), but \(p-1\) of these by the \(x\), one obtains quotients of two algebraic functions expressed by \(\Theta\)-functions:
\[
\tag{6}
\frac{\left(\frac{\varphi}{\psi}\right)_{x^{(1)}}}
{\left(\frac{\varphi}{\psi}\right)_{z^{(1)}}}
=
\prod_{\xi\eta}
\frac{
\Theta\!\left(\int_{\xi}^{x^{(1)}}du_h+m_h\right)
\Theta\!\left(\int_{\eta}^{z^{(1)}}du_h+m_h\right)}
{\Theta\!\left(\int_{\eta}^{x^{(1)}}du_h+m_h\right)
\Theta\!\left(\int_{\xi}^{z^{(1)}}du_h+m_h\right)},
\]
where the points \(\xi\) mean the intersection points of \(\varphi=0\) with \(f=0\), and the points \(\eta\) the intersection points of \(\psi=0\) with \(f=0\).

It is noteworthy that these expressions are necessarily independent of the \(p-1\) points \(x^{(2)},x^{(3)}\ldots x^{(p)}\) which occur in the \(p\) quantities \(m\). It follows from this that these \(p\) quantities themselves can influence the right-hand sides of the above equations only in one combination, since only one combination of the \(m\) can exist which is independent of the \(x^{(2)},x^{(3)}\ldots x^{(p)}\).

If one substitutes the expressions (5) in the right-hand side of equation (6), one is led back to Abel's theorem. Conversely, starting from this theorem, one can develop equations analogous to equation (6), but containing a smaller number of factors on the right.

For this purpose consider the equation of Abel's theorem:
\[
\log\left(\frac{\varphi}{\psi}\right)_{\xi}-\log\left(\frac{\varphi}{\psi}\right)_{\eta}
=\sum\int_{x^{(i)}}^{y^{(i)}}d\Pi_{\xi\eta},
\]
where now the \(y\) and \(x\) denote pairs of intersection points of the curves \(\psi=0\) and \(\varphi=0\) with \(f=0\). Let the number of the \(x\) and of the \(y\) be \(r\). Add on the right as many integrals with coincident upper and lower limits as are needed to make the number of all integrals divisible by \(p\) (\(q\cdot p\)); then we can at will combine every \(p\) integrals into a transcendent \(T\), and obtain the equation
\[
\log\left(\frac{\varphi}{\psi}\right)_{\xi}-\log\left(\frac{\varphi}{\psi}\right)_{\eta}
=\sum_{i=0}^{i=q-1}
T_{\xi\eta}\binom{y^{(ip+1)},y^{(ip+2)}\ldots y^{(i+1)p}}
{x^{(ip+1)},x^{(ip+2)}\ldots x^{(i+1)p}}.
\]
Thus put
\[
v_h^{(i)}=\int_{\alpha^{(1)}}^{x^{(ip+1)}}du_h+\int_{\alpha^{(2)}}^{x^{(ip+2)}}du_h\cdots+\int_{\alpha^{(p)}}^{x^{(i+1)p}}du_h,
\]
\[
V_h^{(i)}=\int_{\alpha^{(1)}}^{y^{(ip+1)}}du_h+\int_{\alpha^{(2)}}^{y^{(ip+2)}}du_h\cdots+\int_{\alpha^{(p)}}^{y^{(i+1)p}}du_h,
\]
then, using formula (1), by passing from the logarithms to the numbers we obtain
\[
\frac{\left(\frac{\varphi}{\psi}\right)_{\xi}}
{\left(\frac{\varphi}{\psi}\right)_{\eta}}
=
\prod_{i=0}^{i=q-1}
\frac{
\Theta\!\left(v_h^{(i)}-\int_{\mu}^{\xi}du_h\right)
\Theta\!\left(V_h^{(i)}-\int_{\mu}^{\eta}du_h\right)}
{\Theta\!\left(v_h^{(i)}-\int_{\mu}^{\eta}du_h\right)
\Theta\!\left(V_h^{(i)}-\int_{\mu}^{\xi}du_h\right)}.
\]
The formation of the expression on the right can be carried out in infinitely many different ways, since the added points \(y^{(r+1)},x^{(r+1)},\ldots,y^{q p},x^{q p}\) are arbitrary, and moreover the combining of the \(q p\) integrals into groups of \(p\) can still be altered in many ways.

\subsubsection*{\S\,59. Quotients of alternating functions expressed by \(\Theta\)-functions.}

The equation developed in the preceding section from Abel's theorem can also be used, in the sense of the inversion problem, to express directly by \(\Theta\)-functions certain symmetric functions of \(p\) points \(x\) which no longer have the simple form assumed above.

For the determination of the system of intersections of a curve of order \(m\) \((m>n-3)\), passing through the double and return points, with the curve \(f=0\), suppose that so many intersection points are given that the addition of \(p\) further intersection points would determine the rest and hence fix the entire system of intersections. The equation of such a curve is composed linearly from the equations of \(p+1\) such curves, so that
\[
\varphi=k_1\varphi_1+k_2\varphi_2\cdots+k_{p+1}\varphi_{p+1}=0
\]
becomes the equation of a curve of this kind if the \(k\) are arbitrary constants, while \(\varphi_1=0,\varphi_2=0\ldots\varphi_{p+1}=0\) are special curves satisfying the same conditions, with no linear relation among their left-hand sides. If the curve \(\varphi=0\) is now to pass through \(p\) variable points \(x^{(1)},x^{(2)}\ldots x^{(p)}\), then the factors \(k\) are determined by linear equations, and, if \(X\) denotes the running point, one obtains
\[
\tag{1}
\varphi(X)=
\begin{vmatrix}
\varphi_1(X) & \varphi_2(X) & \cdots & \varphi_{p+1}(X)\\
\varphi_1(x^{(1)}) & \varphi_2(x^{(1)}) & \cdots & \varphi_{p+1}(x^{(1)})\\
\varphi_1(x^{(2)}) & \varphi_2(x^{(2)}) & \cdots & \varphi_{p+1}(x^{(2)})\\
\vdots & \vdots & & \vdots\\
\varphi_1(x^{(p)}) & \varphi_2(x^{(p)}) & \cdots & \varphi_{p+1}(x^{(p)})
\end{vmatrix}.
\]
We put this expression in the last formula of the preceding section in place of \(\varphi\); in place of \(\psi\), however, we put the same function in which only constant points \(a\) have been substituted in place of the \(x\). Each of these curves cuts \(f=0\), apart from the common intersection points, first respectively in the points \(x,a\), and then also in \(p\) further points \(X,A\), which are connected algebraically with the \(x,a\), and between which there also exist the \(p\) transcendent equations
\[
\tag{2}
\sum_{i=1}^{i=p}\int_{a^{(i)}}^{x^{(i)}}du_h+
\sum_{i=1}^{i=p}\int_{A^{(i)}}^{X^{(i)}}du_h=0
\qquad(h=1,2\ldots p).
\]
The groups of the \(x\) in the preceding section are here represented by the group of the \(x\) and that of the \(X\), and likewise the groups of the \(y\) by the group of the \(a\) and that of the \(A\). Thus if one puts
\[
\sum_{i=1}^{i=p}\int_{a^{(i)}}^{x^{(i)}}du_h=v_h,\qquad
\sum_{i=1}^{i=p}\int_{a^{(i)}}^{a^{(i)}}du_h=v_h^\circ,\qquad
\sum_{i=1}^{i=p}\int_{a^{(i)}}^{A^{(i)}}du_h=v'_h,
\]
then from (2) it results that
\[
\sum_{i=1}^{i=p}\int_{a^{(i)}}^{X^{(i)}}du_h=v_h^\circ+v'_h-v_h,
\]
and hence, from the last equation of the preceding section, one obtains
\[
\tag{3}
\begin{aligned}
&\frac{\Sigma\pm\varphi_1(\xi)\varphi_2(x^{(1)})\cdots\varphi_{p+1}(x^{(p)})}
{\Sigma\pm\varphi_1(\eta)\varphi_2(x^{(1)})\cdots\varphi_{p+1}(x^{(p)})}
\cdot
\frac{\Sigma\pm\varphi_1(\eta)\varphi_2(a^{(1)})\cdots\varphi_{p+1}(a^{(p)})}
{\Sigma\pm\varphi_1(\xi)\varphi_2(a^{(1)})\cdots\varphi_{p+1}(a^{(p)})}\\
&\quad=
\frac{
\Theta\!\left(v_h-\int_{\mu}^{\xi}du_h\right)
\Theta\!\left(v_h^\circ+v'_h-v_h-\int_{\mu}^{\xi}du_h\right)
\Theta\!\left(v_h^\circ-\int_{\mu}^{\eta}du_h\right)
\Theta\!\left(v'_h-\int_{\mu}^{\eta}du_h\right)}
{\Theta\!\left(v_h-\int_{\mu}^{\eta}du_h\right)
\Theta\!\left(v_h^\circ+v'_h-v_h-\int_{\mu}^{\eta}du_h\right)
\Theta\!\left(v_h^\circ-\int_{\mu}^{\xi}du_h\right)
\Theta\!\left(v'_h-\int_{\mu}^{\xi}du_h\right)}.
\end{aligned}
\]
We write this equation in the form
\[
\tag{4}
\frac{\Sigma\pm\varphi_1(\xi)\varphi_2(x^{(1)})\cdots\varphi_{p+1}(x^{(p)})}
{\Sigma\pm\varphi_1(\xi)\varphi_2(a^{(1)})\cdots\varphi_{p+1}(a^{(p)})}
=
M\cdot
\frac{
\Theta\!\left(v_h-\int_{\mu}^{\xi}du_h\right)
\Theta\!\left(v_h^\circ+v'_h-v_h-\int_{\mu}^{\xi}du_h\right)}
{\Theta\!\left(v_h^\circ-\int_{\mu}^{\xi}du_h\right)
\Theta\!\left(v'_h-\int_{\mu}^{\xi}du_h\right)}.
\]
Then \(M\) is evidently independent of \(\xi\), since by the preceding equation it does not change when \(\xi\) is replaced by \(\eta\). Let us now, for the moment, denote by \(R_1,R_2\ldots R_{p+1},R_1^\circ,R_2^\circ\ldots R_{p+1}^\circ\) the subdeterminants composed from the \(\varphi_k(x^{(i)})\) and \(\varphi_k(a^{(i)})\), so that
\[
\varphi=R_1\varphi_1(\xi)+R_2\varphi_2(\xi)\cdots+R_{p+1}\varphi_{p+1}(\xi),
\]
\[
\psi=R_1^\circ\varphi_1(\xi)+R_2^\circ\varphi_2(\xi)\cdots+R_{p+1}^\circ\varphi_{p+1}(\xi).
\]
Then one can regard equation (4) as a linear equation with the unknowns \(R_1,R_2\ldots R_{p+1},M\). We form \(p+1\) different equations of this kind by successively introducing \(p+1\) different constant points \(\xi^{(1)},\xi^{(2)}\ldots\xi^{(p+1)}\) in place of \(\xi\). If then \(\Delta_i^{(h)}\) is the subdeterminant of the determinant
\[
\Sigma\pm\varphi_1(\xi^{(1)})\varphi_2(\xi^{(2)})\cdots\varphi_{p+1}(\xi^{(p+1)})
\]
corresponding to the element \(\varphi_i(\xi^{(h)})\), then from the solution of that system there results the equation
\[
mR_i=
\sum_{h=1}^{h=p+1}
\left[R_1^\circ\varphi_1(\xi^{(h)})+R_2^\circ\varphi_2(\xi^{(h)})\cdots+
R_{p+1}^\circ\varphi_{p+1}(\xi^{(h)})\right]
\]
\[
\cdot\Delta_i^{(h)}
\frac{
\Theta\!\left(v-\int_{\mu}^{\xi^{(h)}}du\right)
\Theta\!\left(v^\circ+v'-v-\int_{\mu}^{\xi^{(h)}}du\right)}
{\Theta\!\left(v^\circ-\int_{\mu}^{\xi^{(h)}}du\right)
\Theta\!\left(v'-\int_{\mu}^{\xi^{(h)}}du\right)},
\]
where \(m\) is an undetermined factor. Thus the ratios of the \(p+1\) alternating functions \(R_1,R_2\ldots R_{p+1}\) have thereby been expressed by \(\Theta\)-functions whose arguments differ from the variables of the inversion problem only by constants.
% --- Translation continuation: source PDF pp. 228--243, printed pp. 206--221. ---
\clearpage

\begin{center}
{\bfseries Ninth Section.}\par\vspace{0.6em}
{\Large\bfseries The Function \(\Theta\).}
\end{center}

\subsection*{\S\,60. Vanishing of the function \(\Theta\). Indeterminate case of the inversion problem.}

The function \(\Theta\), which was derived above from the function \(V\), possesses the properties of the latter; it does not become infinite for finite values of \(\omega\), and it vanishes when the \(\omega\) satisfy certain conditions, but in no other case.

Put, as above,
\[
\tag{1}
\omega_h=\sum_{i=1}^{i=p}\int_{\alpha^{(i)}}^{x^{(i)}}du_h-\int_{\mu}^{\xi}du_h .
\]
If the \(\omega\), the arguments of \(\Theta\), are given, then in these equations \(\xi\) is still arbitrary; one can let it coincide with \(\mu\), or conversely move \(\mu\) toward the variable point \(\xi\), whereby the last integral is made to vanish. When one has disposed of \(\mu\), \(\xi\), and hence also of the \(\alpha\), then the \(x\) are in general completely and uniquely determined.

The function \(\Theta\), like \(V\), now always vanishes in the following two cases:
\[
\begin{array}{ll}
1) & \text{when a point }x\text{ coincides with }\xi;\\
2) & \text{when the points }x\text{ lie on a curve of order }(n-3)\text{ passing through}\\
   & \text{the double points and return points;}
\end{array}
\]
and \(\Theta\) can vanish in no other case.

In both cases the \(\omega\) then contain only \(p-1\) variables; for in the first case \(\xi\) cancels directly against the \(x\) coinciding with it; in the second case one can replace the integral sums which have the points \(x^{(1)},x^{(2)}\ldots x^{(p)}\) as upper limits by others which have as upper limiting points the points \(x^{(p+1)},x^{(p+2)}\ldots x^{(2p-2)}\), lying with them on a curve of order \((n-3)\), so that these together with \(\xi\) remain as variables. In both cases, therefore, a condition holds among the \(\omega\).

That the two cases come back to one another, and hence are not essentially different, has already been used above. For since \(\Theta(-\omega)=\Theta(\omega)\), the following mode of expression is justified just as equation (1) is:
\[
\tag{2}
-\omega_h=\sum_{i=1}^{i=p}\int_{\alpha^{(i)}}^{y^{(i)}}du_h-\int_{\mu}^{\eta}du_h .
\]
From the equation
\[
\sum_{i=1}^{i=p}\left(\int_{\alpha^{(i)}}^{x^{(i)}}du_h+
\int_{\alpha^{(i)}}^{y^{(i)}}du_h\right)-
\left(\int_{\mu}^{\xi}du_h+\int_{\mu}^{\eta}du_h\right)=0
\]
it now follows that the \(x,y\) are the intersection points of a curve of order \((n-2)\) with \(f=0\), which passes through the double and return points, as well as through the \(p-2\) points in which the line \(\xi\eta\) still meets it. If the \(x\) lie on a curve of order \((n-3)\), then this curve, together with the line \(\xi\eta\), must take the place of the curve of order \((n-2)\); hence one point \(y\) must fall at \(\xi\), and another at \(\eta\). Thus, while for the \(x\) case (2) occurs, for the \(y\) case (1) occurs, and conversely, if for the \(y\) case (2) occurs, then the \(x\) satisfy the conditions of case (1).

In the applications made above of Abel's theorem, in which the expression
\[
\frac{\left(\dfrac{\varphi}{\psi}\right)_{\xi}}
     {\left(\dfrac{\varphi}{\psi}\right)_{\eta}}
\]
was represented as a quotient of \(\Theta\)-products, we can always avoid the case in which some of the \(\Theta\) might vanish. For we transformed the logarithm of the above algebraic expression into a sum of integrals of the third kind; to these we can always assign others with coinciding upper and lower limits in such a way that groups of \(p\) integrals are formed whose corresponding \(\Theta\)-quotients do not become indeterminate.

In the inversion problem, however, the genuinely indeterminate case does occur: the \(x\) become indeterminate as soon as the integral sums of the first kind are subject to a certain condition, namely the condition saying that the \(x\) lie on a curve of order \((n-3)\) through the double and return points. In the formulae of the inversion problem this shows itself by the simultaneous vanishing of all \(\Theta\)-functions in numerator and denominator.

For this case, however, the given inversion problem can always be reduced easily to another one which is determinate. For although the \(x\) themselves necessarily become indeterminate here, the \(p-2\) points \(y\) in which \(f=0\) is still cut by the curve of order \((n-3)\) through the \(x\) are by no means indeterminate. To these \(p-2\) unknown points \(y\) add two other known points \(y^{(p-1)},y^{(p)}\), arbitrarily chosen on \(f=0\), which only do not lie with the first on a curve of order \((n-3)\). Then
\[
\sum_{i=1}^{i=p}\int_{\alpha^{(i)}}^{x^{(i)}}du_h=v_h,
\]
but also
\[
\tag{3}
\sum_{i=1}^{i=p}\int_{\alpha^{(i)}}^{y^{(i)}}du_h=k_h-v_h,
\]
where \(k_h\) denotes the constant to be determined from Abel's theorem,
\[
k_h=\sum_{i=1}^{i=p}\left(\int_{\alpha^{(i)}}^{x^{(i)}}du_h+
\int_{\alpha^{(i)}}^{y^{(i)}}du_h\right).
\]
If one takes equations (3) as the basis of a new inversion problem, one obtains (cf. \S\,57. 3.)
\[
\prod_{i=1}^{i=p}
\frac{\left(\dfrac{\varphi}{\psi}\right)_{y^{(i)}}-\lambda}
     {\left(\dfrac{\varphi}{\psi}\right)_{\alpha^{(i)}}-\lambda}
=
\prod_{\xi\eta}
\frac{\Theta\left(k_h-v_h-\int_{\mu}^{\xi}du_h\right)\Theta\left(\int_{\mu}^{\eta}du_h\right)}
     {\Theta\left(k_h-v_h-\int_{\mu}^{\eta}du_h\right)\Theta\left(\int_{\mu}^{\xi}du_h\right)},
\]
where in the product on the right two factors, corresponding to the given \(y\), are completely known. Separating these factors and assigning different values to \(\lambda\), one can form an equation of degree \((p-2)\) whose roots are the unknown values of \(\left(\dfrac{\varphi}{\psi}\right)_{y^{(i)}}\).

Finally, one can ask for a direct method of reducing the quotient of two vanishing \(\Theta\)-functions to its true value. Since case (2) above could always be reduced to case (1), we restrict ourselves to the assumption that both \(\Theta\)-functions, in which we may suppose the arbitrary point \(\xi\) to be common, vanish because this point coincides with one of the limiting points. We therefore consider the expression
\[
L=\lim_{y=\xi}
\frac{\Theta\left(\sum_{i=1}^{i=p-1}\int_{\alpha^{(i)}}^{x^{(i)}}du_h+
\int_{\alpha^{(p)}}^{y}du_h-\int_{\mu}^{\xi}du_h\right)}
{\Theta\left(\sum_{i=1}^{i=p-1}\int_{\alpha^{(i)}}^{z^{(i)}}du_h+
\int_{\alpha^{(p)}}^{y}du_h-\int_{\mu}^{\xi}du_h\right)}.
\]
Multiplying this by a similar, non-indeterminate \(\Theta\)-quotient, we have by \S\,57. 1.
\[
\log L+
\log\frac{\Theta\left(\sum_{i=1}^{i=p-1}\int_{\alpha^{(i)}}^{z^{(i)}}du_h+
\int_{\alpha^{(p)}}^{y}du_h-\int_{\mu}^{\eta}du_h\right)}
{\Theta\left(\sum_{i=1}^{i=p-1}\int_{\alpha^{(i)}}^{x^{(i)}}du_h+
\int_{\alpha^{(p)}}^{y}du_h-\int_{\mu}^{\eta}du_h\right)}
= -\sum_{i=1}^{i=p-1}\int_{z^{(i)}}^{x^{(i)}}d\Pi_{\xi\eta}.
\]
As the right-hand side shows, this expression depends on the position of the point \(y\); for the left-hand side we can therefore substitute the expression which arises when \(y\) is replaced by a point \(t\) chosen arbitrarily on \(f=0\). Accordingly, for \(L\) we obtain the quotient of two products of three \(\Theta\)-functions each; the arguments of the functions are
\[
\begin{array}{c@{\quad}c}
\text{in the numerator:} & \text{in the denominator:}\\[0.4em]
\displaystyle \sum_{i=1}^{i=p-1}\int_{\alpha^{(i)}}^{x^{(i)}}du_h+\int_{\alpha^{(p)}}^{t}du_h-\int_{\mu}^{\xi}du_h
&
\displaystyle \sum_{i=1}^{i=p-1}\int_{\alpha^{(i)}}^{z^{(i)}}du_h+\int_{\alpha^{(p)}}^{t}du_h-\int_{\mu}^{\eta}du_h
\\[1.0em]
\displaystyle \sum_{i=1}^{i=p-1}\int_{\alpha^{(i)}}^{z^{(i)}}du_h+\int_{\alpha^{(p)}}^{t}du_h-\int_{\mu}^{\eta}du_h
&
\displaystyle \sum_{i=1}^{i=p-1}\int_{\alpha^{(i)}}^{z^{(i)}}du_h+\int_{\alpha^{(p)}}^{t}du_h-\int_{\mu}^{\xi}du_h
\\[1.0em]
\displaystyle \sum_{i=1}^{i=p-1}\int_{\alpha^{(i)}}^{x^{(i)}}du_h+\int_{\alpha^{(p)}}^{y}du_h-\int_{\mu}^{\eta}du_h
&
\displaystyle \sum_{i=1}^{i=p-1}\int_{\alpha^{(i)}}^{x^{(i)}}du_h+\int_{\alpha^{(p)}}^{y}du_h-\int_{\mu}^{\eta}du_h .
\end{array}
\]
Thus the indeterminate \(\Theta\)-quotient has been reduced to the quotient of non-vanishing products, and the new expression contains the two arbitrarily introduced points \(t,\eta\).

If not just one but several point-pairs \(x,z\) coincide with the point \(\xi\), then the expression is again indeterminate; but by the same procedure one can continually reduce the number of point-pairs causing the indeterminacy, and finally arrives at determinate expressions. If \(k\) point-pairs \(x,z\) have coalesced with \(\xi\), this procedure has to be applied \(k\) times successively. From this we may conjecture that in this case the function \(\Theta\) does not vanish simply, but \(k\)-fold. This is in fact confirmed by the appearance of \(\log\Theta\), as this logarithm is present in the function \(-U\). If, in this function (\S\,49), \(x^{(1)},x^{(2)}\ldots x^{(k)}\) approach the common endpoint \(\xi\), then the discontinuous part of \(-U\), by the properties of the integral of the third kind, becomes
\[
-\sum_{i=1}^{i=k}\log(\xi-x^{(i)})+
\sum_{i<h}\log(x^{(i)}-x^{(h)})-
\log\varphi(x^{(1)},x^{(2)}\ldots x^{(p)}).
\]
But the last two terms remain finite together even when the \(x\) approach one another; hence \(U\) approaches the value
\[
-\log\left[(\xi-x^{(1)})(\xi-x^{(2)})\cdots(\xi-x^{(k)})\right],
\]
or \(\Theta\) approaches the expression
\[
\operatorname{Const.}\,(\xi-x^{(1)})(\xi-x^{(2)})\cdots(\xi-x^{(k)}),
\]
that is, it vanishes \(k\)-fold, and with it vanish its first \(k-1\) differential quotients with respect to \(\xi\).

From the above limiting value of \(-U\) one sees, moreover, that the function \(\Theta\) also vanishes to higher order when not only some of the \(x\) fall at \(\xi\), but the points \(x\), during this approach, also continue to lie on a curve of order \((n-3)\). To the order of the vanishing product \(\xi-x^{(1)}\cdot\xi-x^{(2)}\ldots\) one must then add the order of the vanishing expression
\[
\frac{\varphi}{\Pi(x^{(i)}-x^{(h)})},
\]
which is in general the first.

\subsection*{\S\,61. Special cases in which the function \(\Theta\) vanishes.}

Conversely, a \(\Theta\)-quotient which for any reasons assumes the form \(0/0\) can always be reduced to such \(\Theta\)-quotients as become indeterminate only because the upper limiting points \((x,z)\) lie on curves of order \((n-3)\) passing through the double and return points. In doing this, however, several special circumstances can arise, and these deserve a more exact discussion.

In general, \(mn-p\) intersection points of a curve of order \(m\) through the double and return points of \(f=0\) determine the remaining \(p\). Only when \(m=n-3\), or smaller, does an exception occur; indeed we can restrict ourselves to the case \(m=n-3\), since for smaller values of \(m\) improper curves of order \((n-3)\) can be produced by adjoining irrelevant curves.

Assume arbitrarily any \(p-1\) points \(x^{(1)},x^{(2)}\ldots x^{(p-1)}\) on the curve \(f=0\), and draw through them and through the double and return points a curve of order \((n-3)\). Then this curve is determined and gives \(p-1\) further intersection points on \(f=0\). If \(x^{(p-1)}\) is moved, then in general all the remaining \(p-1\) intersection points will move as well. But it can happen that, for a special position of the points \(x^{(1)},x^{(2)}\ldots x^{(p-2)}\), the motion of \(x^{(p-1)}\) has no influence on one of the remaining \(p-1\) intersection points; that is, those first \(p-2\) points already determine a further intersection point \(x^{(p)}\) without adding the \((p-1)\)-st point, so that a special position of \(x^{(p-1)}\) then fixes only the missing \(p-2\) intersection points. Since this contains only one condition, we may choose the points \(x^{(1)},x^{(2)}\ldots x^{(p-3)}\) arbitrarily and ask how \(x^{(p-2)}\) must be determined so that all curves of the stated kind through \(x^{(1)},x^{(2)}\ldots x^{(p-2)}\) also pass through another point \(x^{(p)}\). If now we draw any three such curves through \(x^{(1)},x^{(2)}\ldots x^{(p-3)}\), and their equations are \(\varphi_1=0,\varphi_2=0,\varphi_3=0\), then, provided there is no linear relation among \(\varphi_1,\varphi_2,\varphi_3\), every fourth curve can be represented in the form
\[
\alpha_1\varphi_1+\alpha_2\varphi_2+\alpha_3\varphi_3=0.
\]
The \(\alpha\) denote variable parameters. They are reduced to one ratio if one considers those curves which pass through the required point \(x^{(p-2)}\), for which therefore
\[
\tag{1}
\alpha_1\varphi_1(x^{(p-2)})+\alpha_2\varphi_2(x^{(p-2)})+\alpha_3\varphi_3(x^{(p-2)})=0.
\]
But it was required that all these curves should pass through a definite other point \(x^{(p)}\), that is, that equation (1) should always imply the other equation
\[
\tag{2}
\alpha_1\varphi_1(x^{(p)})+\alpha_2\varphi_2(x^{(p)})+\alpha_3\varphi_3(x^{(p)})=0.
\]
This is possible only when the three equations hold:
\[
\tag{3}
\begin{cases}
\varphi_1(x^{(p)})=\lambda\varphi_1(x^{(p-2)}),\\
\varphi_2(x^{(p)})=\lambda\varphi_2(x^{(p-2)}),\\
\varphi_3(x^{(p)})=\lambda\varphi_3(x^{(p-2)}).
\end{cases}
\]
From these equations \(\lambda\) and the sought points \(x^{(p)},x^{(p-2)}\) are determined.

Equations (3) show a reciprocity between the points \(x^{(p)},x^{(p-2)}\). Thus there are a number of point-pairs with the property that, together with given \(p-3\) points, they form the common intersection points of a system of curves of order \((n-3)\) with \(f=0\), and therefore, together with those points, exceptionally do not determine a curve of order \((n-3)\).

The number of such point-pairs belonging to a system of \(p-3\) given points is very easily determined by means of the investigations carried out in the third section. The equations
\[
\tag{4}
\begin{cases}
\mu y_1=\varphi_1(x_1,x_2,x_3),\\
\mu y_2=\varphi_2(x_1,x_2,x_3),\\
\mu y_3=\varphi_3(x_1,x_2,x_3)
\end{cases}
\]
are the equations for a one-to-one transformation of order \((n-3)\). The curves \(\varphi=0\) have all double and return points in common, as well as the \(p-3\) arbitrarily chosen points \(x^{(1)},x^{(2)}\ldots x^{(p-3)}\). The point-pairs \(x^{(p)},x^{(p-2)}\) are therefore none other than those which, under this one-to-one substitution, are united to new double or return points. Hence the number of such point-pairs is also easily determined; it is the number of double and return points of the new curve. Since (cf. \S\,18) the order of the new curve is \(p+1\), the number of its double and return points is
\[
d^{(1)}+r^{(1)}=\frac{p\cdot p-3}{2}.
\]
Thus there are \(\dfrac{p\cdot p-3}{2}\) point-pairs which, together with \(p-3\) given points, yield systems of such \(p-1\) points as do not determine a curve of order \((n-3)\) through the double and return points of \(f=0\).

Just as two other points were here sought for \(p-3\) given points, so that with them they determine not one but infinitely many curves of order \((n-3)\), distinguished by the values of an arbitrary parameter, one can likewise treat the question under what circumstances a smaller number of given points already has assigned to it a number of fixed points which together determine a curve-system of order \((n-3)\) with a larger number of parameters. First we can ask how \(p-3\) points \(x^{(1)},x^{(2)}\ldots x^{(p-3)}\) must lie so that every curve of order \((n-3)\) through them (and the double and return points) also passes through a common point. For this purpose imagine adding to the points \(x^{(1)},x^{(2)}\ldots x^{(p-3)}\) two movable points \(x^{(p-2)},x^{(p-1)}\), and seek the condition under which, during the motion of these two new points, one of the remaining intersection points of the curve with \(f=0\), namely \(x^{(p)}\), remains fixed. This requirement involves two conditions. Thus we regard only \(x^{(1)},x^{(2)}\ldots x^{(p-5)}\) as arbitrarily given, and ask how \(x^{(p-4)}\) and \(x^{(p-3)}\) must lie so that all curves through them still pass through a common point \(x^{(p)}\). Since four fewer points are given here than are required to determine a curve, one can draw five curves \(\varphi_1=0,\varphi_2=0\ldots\varphi_5=0\), from which any curve of order \((n-3)\) through those \(p-5\) points and the double and return points is composed. The general curve of the system determined by these points has the form
\[
\alpha_1\varphi_1+\alpha_2\varphi_2\cdots+\alpha_5\varphi_5=0.
\]
The requirement is that every curve of this system through \(x^{(p-4)}\) and \(x^{(p-3)}\) also pass through \(x^{(p)}\), that is, that from the equations
\[
\begin{aligned}
&\alpha_1\varphi_1(x^{(p-4)})+\alpha_2\varphi_2(x^{(p-4)})\cdots+\alpha_5\varphi_5(x^{(p-4)})=0,\\
&\alpha_1\varphi_1(x^{(p-3)})+\alpha_2\varphi_2(x^{(p-3)})\cdots+\alpha_5\varphi_5(x^{(p-3)})=0
\end{aligned}
\]
the following equation follows of itself:
\[
\alpha_1\varphi_1(x^{(p)})+\alpha_2\varphi_2(x^{(p)})\cdots+\alpha_5\varphi_5(x^{(p)})=0.
\]
For this it is necessary and sufficient that the function values of the last equation be linearly composed from those of the first two, that is, that the equations
\[
\varphi_i(x^{(p)})=\lambda\varphi_i(x^{(p-4)})+\mu\varphi_i(x^{(p-3)}),
\qquad i=1,2\ldots 5,
\]
hold. These are five equations which determine \(\lambda,\mu\) and the three points \(x^{(p-4)},x^{(p-3)},x^{(p)}\). Determining the number of their solutions requires an extension of the considerations carried out in the third section.

In general, one can ask when all curves of the kind under consideration that pass through \(x^{(1)},x^{(2)}\ldots x^{(k)}\) also pass through a common point \(x^{(p)}\). To investigate this, add to the fixed points \(x^{(1)},x^{(2)}\ldots x^{(k)}\) movable points \(x^{(k+1)},x^{(k+2)}\ldots x^{(p-1)}\), \(p-k-1\) in number, and examine when their motion leaves unchanged a further intersection point \(x^{(p)}\) of the curves through them with \(f=0\). This involves \(p-k-1\) conditions; hence first only the points \(x^{(1)},x^{(2)}\ldots x^{(2k-p+1)}\) are regarded as arbitrarily given, and the remaining points \(x^{(2k-p+2)},x^{(2k-p+3)}\ldots x^{(k)}\) are determined so that these conditions are satisfied. Through those \(2k-p+1\) points one can now draw \(2p-2k-1\) special curves of order \((n-3)\), \(\varphi_1=0,\varphi_2=0\ldots,\varphi_{2p-2k-1}=0\), from which the most general curve is composed in the form
\[
\alpha_1\varphi_1+\alpha_2\varphi_2\cdots+\alpha_{2p-2k-1}\varphi_{2p-2k-1}=0.
\]
If now a curve of this system passes through the points \(x^{(2k-p+2)},x^{(2k-p+3)}\ldots x^{(k)}\), it is to pass of itself through another point \(x^{(p)}\). Thus from the equations
\[
\begin{aligned}
&\alpha_1\varphi_1(x^{(2k-p+2)})+\alpha_2\varphi_2(x^{(2k-p+2)})\cdots+
\alpha_{2p-2k-1}\varphi_{2p-2k-1}(x^{(2k-p+2)})=0,\\
&\alpha_1\varphi_1(x^{(2k-p+3)})+\alpha_2\varphi_2(x^{(2k-p+3)})\cdots+
\alpha_{2p-2k-1}\varphi_{2p-2k-1}(x^{(2k-p+3)})=0,\\
&\hfill\cdots\hfill,\\
&\alpha_1\varphi_1(x^{(k)})+\alpha_2\varphi_2(x^{(k)})\cdots+
\alpha_{2p-2k-1}\varphi_{2p-2k-1}(x^{(k)})=0
\end{aligned}
\]
the additional equation
\[
\alpha_1\varphi_1(x^{(p)})+\alpha_2\varphi_2(x^{(p)})\cdots+
\alpha_{2p-2k-1}\varphi_{2p-2k-1}(x^{(p)})=0
\]
must follow. Therefore one must have the equations
\[
\varphi_i(x^{(p)})=\lambda\varphi_i(x^{(2k-p+2)})+
\mu\varphi_i(x^{(2k-p+3)})\cdots+
\sigma\varphi_i(x^{(k)}),
\]
\[
i=1,2\ldots 2p-2k-1,
\]
which just suffice to determine \(\lambda,\mu\ldots\sigma\), and in addition the points \(x^{(2k-p+2)},x^{(2k-p+3)}\ldots x^{(k)}\).

If the upper limiting points \(x^{(1)},x^{(2)}\ldots x^{(p)}\) in the \(\Theta\)-function are such that one of the special cases treated here occurs, then the \(\Theta\)-function always vanishes, namely in accordance with case (2) of \S\,60, independently of \(\xi\). But it also vanishes independently of those points \(x\) whose position remained indeterminate above; thus, if every curve of order \((n-3)\) through \(x^{(1)},x^{(2)}\ldots x^{(p-2)}\) also passes through \(x^{(p)}\), it vanishes independently of the special position of \(x^{(p-1)}\); if every curve through \(x^{(1)},x^{(2)}\ldots x^{(p-3)}\) passes through \(x^{(p)}\), it vanishes independently of \(x^{(p-1)}\) and \(x^{(p-2)}\); in general, if every curve through \(x^{(1)},x^{(2)}\ldots x^{(k)}\) passes through \(x^{(p)}\), then it vanishes independently of the position of the points \(x^{(k+1)},x^{(k+2)}\ldots x^{(p-1)}\). The function \(U\) then approaches the function
\[
-\log\varphi(x^{(1)},x^{(2)}\ldots x^{(p)}),
\]
hence \(\Theta\), up to a constant, approaches the value
\[
\varphi(x^{(1)},x^{(2)}\ldots x^{(p)});
\]
and this expression vanishes, but in such a way that not all points \(x\) participate, but rather a series of subdeterminants of the determinant \(\varphi\) vanish which contain only some of the points \(x\).

It is moreover easy to follow the special curves of order \((n-3)\) considered above more exactly in the simplest cases. Thus for a curve of order 5 with two double points, which is the most general example for \(p=4\). If one chooses a point on the curve in any way, one can specify two point-pairs which, together with that point and the double points, become base points of a system of curves of second order; namely the point-pairs in which the connecting line of the given point with one of the double points still cuts the given curve. These lines, combined with arbitrary other lines, form the two systems of conic sections.

\subsection*{\S\,62. The converse of Abel's theorem.}

We now come to those considerations which, conversely, teach how to express combinations of \(\Theta\)-functions by algebraic functions, and which consequently also yield algebraic relations between \(\Theta\)-functions.

Abel's theorem, as expressed at the end of \S\,58, teaches how to represent the double ratio of the four function values
\[
\varphi_\xi,\quad \varphi_\eta,\quad \psi_\xi,\quad \psi_\eta
\]
by \(\Theta\)-functions with the aid of the formula
\[
\tag{1}
\frac{\left(\dfrac{\varphi}{\psi}\right)_{\xi}}
     {\left(\dfrac{\varphi}{\psi}\right)_{\eta}}
=
\prod_{i=0}^{i=q-1}
\frac{\Theta\left(v_h^{(i)}-\int_{\mu}^{\xi}du_h\right)
      \Theta\left(V_h^{(i)}-\int_{\mu}^{\eta}du_h\right)}
     {\Theta\left(v_h^{(i)}-\int_{\mu}^{\eta}du_h\right)
      \Theta\left(V_h^{(i)}-\int_{\mu}^{\xi}du_h\right)},
\]
where
\[
\tag{2}
\begin{aligned}
v_h^{(i)}&=\int_{\alpha^{(1)}}^{x^{(ip+1)}}du_h+
\int_{\alpha^{(2)}}^{x^{(ip+2)}}du_h\cdots+
\int_{\alpha^{(p)}}^{x^{(i-1)p}}du_h,\\
V_h^{(i)}&=\int_{\alpha^{(1)}}^{y^{(ip+1)}}du_h+
\int_{\alpha^{(2)}}^{y^{(ip+2)}}du_h\cdots+
\int_{\alpha^{(p)}}^{y^{(i-1)p}}du_h.
\end{aligned}
\]
The \(q p\) points \(x\) on one side, and the \(q p\) points \(y\) on the other, have no further properties except that, apart from point-pairs \(x,y\) which coincide and could be wholly arbitrary, the points \(y\) are intersections of \(\psi=0\) with \(f=0\), the points \(x\) are intersections of \(\varphi=0\) with \(f=0\), while all remaining intersections of the curves \(\varphi=0\), \(f=0\) also belong to the curves \(\psi=0\), \(f=0\); moreover, that the \(y\) gradually pass into the \(x\) when one lets the curve \(\psi=0\) approach the curve \(\varphi=0\). If we further assume in particular that the curves \(\varphi=0\), \(\psi=0\) are to pass through the double and return points of \(f=0\), then this connection is completely expressed, by Abel's theorem, by the \(p\) equations
\[
\tag{3}
 v_h^0+v_h'\cdots+v_h^{(q-1)}=V_h^0+V_h'\cdots+V_h^{(q-1)}
 \qquad(h=1,2\ldots p).
\]
Thus equation (1) always holds as soon as equation (3) is satisfied. One can therefore reverse Abel's theorem and state it as the following theorem on \(\Theta\)-products:
\[
\begin{gathered}
\text{The product of }q\text{ double ratios}\quad
\prod_{i=0}^{i=q-1}
\frac{\dfrac{\Theta\left(v_h^{(i)}-\int_{\mu}^{\xi}du_h\right)}
              {\Theta\left(V_h^{(i)}-\int_{\mu}^{\xi}du_h\right)}}
     {\dfrac{\Theta\left(v_h^{(i)}-\int_{\mu}^{\eta}du_h\right)}
              {\Theta\left(V_h^{(i)}-\int_{\mu}^{\eta}du_h\right)}}
\end{gathered}
\]
in which only the otherwise arbitrary systems of the \(v,V\) are to be connected by equations (3), is a rational function of the coordinates of \(\xi,\eta\) and of the \(q p\) point-pairs determined from equations (2); and it is equal to
\[
\frac{\left(\dfrac{\varphi}{\psi}\right)_{\xi}}
     {\left(\dfrac{\varphi}{\psi}\right)_{\eta}},
\]
where the curve \(\varphi=0\) vanishes for those \(x\), the curve \(\psi=0\) for those \(y\), which belong to non-coinciding point-pairs, and where both curves pass through the double and return points of \(f=0\).

But one must note that this theorem says nothing when \(q=1\), when only a single double ratio of \(\Theta\)-functions is present. For in this case equations (3) say either that all \(x\) and \(y\) coincide pairwise, or that the \(x\) as well as the \(y\) belong to curves of order \((n-3)\), in which case all \(\Theta\)-functions vanish.

One can also understand the theorem just stated as teaching how to reduce the product
\[
\prod_{i=0}^{i=q-1}
\frac{\Theta\left(v_h^{(i)}-\int_{\mu}^{\xi}du_h\right)}
     {\Theta\left(V_h^{(i)}-\int_{\mu}^{\xi}du_h\right)}
\]
to an algebraic function and to a \(\Theta\)-quotient of the same kind in which only \(\xi\) is replaced by an arbitrary other point \(\eta\). In this sense our theorem gives that of Mr. Riemann (Crelle's Journal vol. 54, p. 153), according to which that \(\Theta\)-quotient, considered as a function of \(\xi\) alone, is an algebraic function.

\subsection*{\S\,63. Products of \(\Theta\)-functions reduced to other and to algebraic functions.}

In the theorem stated above, put the expressions
\[
V_h^{(i)}-\int_{\mu}^{\xi}du_h
\]
with different upper indices \(i\) equal to one another and equal to \(w_h\), and put further
\[
v_h^{(i)}=V_h^{(i)}+m_h^{(i)},
\]
where by the quantities \(m_h^{(i)}\) we understand constants. Then we obtain a \(\Theta\)-quotient which we can denote by
\[
\tag{1}
\frac{\Theta(w+m^{(1)}),\ \Theta(w+m^{(2)})\cdots\Theta(w+m^{(q)})}{\Theta^{q}(w)}.
\]
The arguments \(w\), considered as variable, are here entirely arbitrary, for the equations (3) of the preceding paragraph here pass into
\[
\tag{2}
m_h^{(1)}+m_h^{(2)}+\cdots+m_h^{(q)}=0
\qquad(h=1,2\ldots p).
\]
Under the assumption of these equations among the constants, one can, by the preceding, reduce expression (1) to a similar expression and to an algebraic function. Namely, if one puts
\[
\tag{3}
w_h=\int_{\alpha^{(1)}}^{x^{(1)}}du_h+
\int_{\alpha^{(2)}}^{x^{(2)}}du_h\cdots+
\int_{\alpha^{(p)}}^{x^{(p)}}du_h-
\int_{\mu}^{\xi}du_h,
\]
then quotient (1) can be reduced to a similar one in which only \(\xi\) is replaced by \(\eta\), and to a rational function of the coordinates of \(x^{(1)},x^{(2)}\ldots x^{(p)},\xi,\eta\). At first, the theorem of the preceding paragraph would require that one also introduce the points \(y\), which are defined by the equations
\[
\tag{4}
w_h+m_h^{(i)}=\int_{\alpha^{(1)}}^{y^{(i,1)}}du_h+
\int_{\alpha^{(2)}}^{y^{(i,2)}}du_h\cdots+
\int_{\alpha^{(p)}}^{y^{(i,p)}}du_h-
\int_{\mu}^{\xi}du_h.
\]
If then one lays the curve \(\varphi=0\) through all the points \(y\), and the curve \(\psi=0\) through all the points \(x\), in such a way that the latter has \(q\) intersections with \(f=0\) in each of the points \(x\) (that is, touches with multiplicity \(q\)), while the remaining intersections of \(\varphi=0\), \(f=0\) coincide with the remaining intersections of \(\psi=0\), \(f=0\), and so that finally both curves pass through the double and return points, then
\[
\frac{\left(\dfrac{\varphi}{\psi}\right)_{\xi}}
     {\left(\dfrac{\varphi}{\psi}\right)_{\eta}}
\]
is the desired algebraic function.

Meanwhile, in the present case the points \(y\) can always be reduced to constant points and to the points \(x\). Let \(c^{(1)},c^{(2)}\ldots c^{(p)}\) denote arbitrary points on \(f=0\), and then determine \(q p\) points \(b^{(i,k)}\) so that
\[
\tag{5}
m_h^{(i)}=\int_{c^{(1)}}^{b^{(i,1)}}du_h+
\int_{c^{(2)}}^{b^{(i,2)}}du_h\cdots+
\int_{c^{(p)}}^{b^{(i,p)}}du_h,
\qquad (i=1,2\ldots q,\ h=1,2\ldots p).
\]
Equations (2) then pass into Abel's theorem; if one draws a curve which touches at each of the points \(c\) with multiplicity \(q\), and which passes through the double and return points, then it cuts in further points; if one then draws through these, through the double and return points, and through the first \(q-1\) systems of the \(b\), a curve of equally high order, then this curve cuts in the \(q\)-th system of the \(b\).

At the same time, from (3), (4), (5) one obtains
\[
\int_{x^{(1)}}^{y^{(i,1)}}du_h+\int_{x^{(2)}}^{y^{(i,2)}}du_h\cdots+
\int_{x^{(p)}}^{y^{(i,p)}}du_h
=
\int_{c^{(1)}}^{b^{(i,1)}}du_h+\int_{c^{(2)}}^{b^{(i,2)}}du_h\cdots+
\int_{c^{(p)}}^{b^{(i,p)}}du_h.
\]
This is again Abel's theorem. If one lays a curve through the \(x\), through the double and return points, and through one system of the \(b\), it still cuts in other points; if through these, the double and return points, and the points \(c\), one lays a new curve of the same order, then it cuts \(f=0\) in the corresponding system of the \(y\). These are therefore determined by the \(x\) themselves and by the constant points \(c,b\). Hence we can represent the coefficients of \(\varphi\) as rational functions of the coordinates of \(x^{(1)},x^{(2)}\ldots x^{(p)}\), as will be developed more exactly below.

We may therefore state the theorem as follows:

\medskip
\noindent I. \emph{The expression}
\[
\frac{\Theta(w+m^{(1)}),\ \Theta(w+m^{(2)})\cdots\Theta(w+m^{(q)})}{\Theta^{q}(w)},
\]
in which
\[
m_h^{(1)}+m_h^{(2)}+\cdots+m_h^{(q)}=0\qquad(h=1,2\ldots p),
\]
is, if
\[
w_h=\int_{\alpha^{(1)}}^{x^{(1)}}du_h+
\int_{\alpha^{(2)}}^{x^{(2)}}du_h\cdots+
\int_{\alpha^{(p)}}^{x^{(p)}}du_h-
\int_{\mu}^{\xi}du_h
\]
is put, equal to a rational function of the coordinates of \(x^{(1)},x^{(2)}\ldots x^{(p)},\xi,\eta\), multiplied by the \(\Theta\)-quotient which arises from the given one when \(\xi\) is replaced in it by \(\eta\).

This theorem finds its necessary complement in the following theorem:

\medskip
\noindent II. \emph{The above expression is also equal to a rational function of the coordinates of \(x^{(1)},x^{(2)}\ldots x^{(p)},\xi,a\), multiplied by the \(\Theta\)-quotient which arises from the given one when one of the points \(x\) is replaced by the point \(a\), chosen arbitrarily on \(f=0\).}

This second theorem is easily reduced to the first. For we write all \(\Theta\)-functions in the given expression with negative argument, and so consider the expression
\[
\frac{\Theta(-w-m^{(1)}),\ \Theta(-w-m^{(2)})\cdots\Theta(-w-m^{(q)})}{\Theta^{q}(-w)}.
\]
The quantities \(-m\) still satisfy the same conditions as the \(m\) themselves; if therefore
\[
-w_h=\int_{\alpha^{(1)}}^{z^{(1)}}du_h+
\int_{\alpha^{(2)}}^{z^{(2)}}du_h\cdots+
\int_{\alpha^{(p)}}^{z^{(p)}}du_h-
\int_{\mu}^{\xi}du_h,
\]
then by theorem I the given expression becomes equal to a similar one in which only \(\xi\) is replaced by another point, \(a\), multiplied by a rational function of the coordinates of \(z^{(1)},z^{(2)}\ldots z^{(p)},\xi,a\). But the \(z,\xi\) are connected with the \(x,\xi\) by the equation
\[
\sum_{i=1}^{i=p}\left(\int_{\alpha^{(i)}}^{z^{(i)}}du_h+
\int_{\alpha^{(i)}}^{x^{(i)}}du_h\right)-
\int_{\mu}^{\xi}du_h-\int_{\mu}^{\xi}du_h=0,
\]
from which it follows that the \(x,z\) are the intersection points of \(f=0\) with a curve of order \((n-2)\) passing through the double and return points and through the \(n-2\) remaining points of the line \(\xi\eta\). If \(\xi\) is now allowed to coincide with one of the \(x\), then one of the \(z\) passes into \(\xi\), and the remaining points \(z\) lie with the remaining \(x\) and with the double and return points on a curve of order \((n-3)\). One sees that these become independent of the \(x\) to be transformed into \(a\); hence the \(\Theta\)-quotient to which the given one returns is also wholly independent of this point, and instead always contains the point \(a\) in its place. The algebraic function, however, becomes a rational function of the coordinates of this point \(x\), of the points \(a,\xi\), and of the points lying with the remaining points on that curve of order \((n-3)\). Since the latter occur in \(\varphi,\psi\) only symmetrically, these functions can also be represented rationally by the coordinates of the remaining points \(x\) which determine those points uniquely; and the algebraic function is therefore a rational function of the coordinates of \(x^{(1)},x^{(2)}\ldots x^{(p)},\xi,a\), as was required.

% --- Continuation: source PDF pp. 244--251, printed pp. 222--229. ---
\clearpage

\subsection*{\S\,64. On the representation of the algebraic expressions contained in the foregoing.}

In order to complete the proofs given above, one must show that the curves \(\varphi,\psi\) really depend rationally on the coordinates of the given points. This comes down to proving the following algebraic theorem:

\begin{quote}
If a curve \(\Psi=0\) is to pass through given points \(x\) on \(f=0\), and a second curve \(\Phi=0\) through the remaining intersection-points of \(\Psi=0\) with \(f=0\), then this involves, for the coefficients of \(\Phi=0\), a series of linear conditions which contain the coordinates of the points \(x\) only in a rational manner.
\end{quote}

One may moreover fix \(\Phi=0\) completely by other conditions linear in its coefficients, and then lay a new curve \(\Phi_1=0\) through the rest of the system of intersection-points; for this curve too one then obtains conditions which are linear for the coefficients of the curve and contain the \(x\) rationally, and so on. The whole series of theorems contained here can be summarized by proving a theorem of more general form.

Namely, let, as in Section I, \(x\) and \(s_x\) be the parameters of two pencils of rays. If we express the equations of two curves \(f=0\), \(\Psi=0\) in \(x\) and \(s_x\), and then eliminate \(s_x\), we obtain an equation
\[
X=0,
\]
which gives the parameters \(x\) of those rays of the first pencil which are directed toward the intersection-points of \(f=0\), \(\Psi=0\). At the same time, in the course of the elimination process one finds \(s_x\) expressed by a rational function of \(x\),
\[
s_x=S(x).
\]
This could assume the form \(\frac{0}{0}\) only if two intersection-points lay on the same ray; this can always be avoided by moving the vertex.

If now a part of the intersection-points is given directly, either singly or as a whole, by any determinations whatever, one can always directly give a part of the pencil \(X=0\). Let this part be \(X^{(1)}=0\). The remainder
\[
X^{(2)}=0=\frac{X}{X^{(1)}}
\]
then arises by mere division, and is therefore dependent on the coefficients of \(X^{(1)}\) only in a rational manner.

If now a curve \(\Phi=0\) is to pass through the intersection-points corresponding to the pencil \(X^{(2)}=0\), one first transforms the equation \(\Phi=0\) into an equation between \(x\) and \(s_x\), and then substitutes for \(s_x\) its value \(S(x)\). The equation \(\Phi=0\) then goes over into an equation
\[
\Omega(x)=0,
\]
which must hold for the rays contained in \(X^{(2)}=0\), and hence must contain \(X^{(2)}=0\) as a factor. This is the condition that \(\Phi=0\) pass through the intersection-points belonging to \(X^{(2)}=0\) of \(f=0\), \(\Psi=0\). If one sets
\[
\Omega(x)=X^{(2)}\cdot X^{(3)},
\]
then by equating the coefficients of equal powers of \(x\) one obtains a series of equations which are linear for the unknown coefficients of \(\Phi\) and \(X^{(3)}\), and therefore also for the former alone when the latter are eliminated. These conditions evidently depend in a rational manner on the coefficients of \(X^{(1)}\). Thus one has the theorem which contains the one stated above as a special case:

\begin{quote}
Suppose a part of the intersection-points of \(f=0\), \(\Psi=0\) is given, so that the parameters of the rays \(x\) of a pencil directed toward them can be represented as roots of an equation
\[
X^{(1)}=0,
\]
and the parameters of the rays of a second pencil directed toward the same points can be expressed rationally through \(x\),
\[
s_x=S(x).
\]
Then the condition that a curve \(\Phi=0\) pass through the remaining intersection-points of \(f=0\), \(\Psi=0\) leads to a series of equations which are linear for the coefficients of \(\Phi\) and rational for the coefficients of \(X^{(1)}\), \(S(x)\).
\end{quote}

A special application of this theorem is already contained in \S\,6, where the determination of the function \(\Omega\) led to a similar problem, which was essentially simplified only by the fact that \(\Psi=0\) was there a straight line.

\subsection*{\S\,65. Representation of quotients of \(\Theta\)-products by algebraic functions.}

For brevity we shall call a product
\[
\Theta(w+m^{(1)}),\quad \Theta(w+m^{(2)})\cdots \Theta(w+m^{(q)})
\]
a complete \(\Theta\)-product as soon as the quantities \(m\) satisfy the \(p\) equations
\[
m_h^{(1)}+m_h^{(2)}+\cdots+m_h^{(q)}=0.
\]
The power \(\Theta^q(w)\) is itself such a complete \(\Theta\)-product. Put
\[
w_h=\int_{a^{(1)}}^{x^{(1)}}du_h+
\int_{a^{(2)}}^{x^{(2)}}du_h+\cdots+
\int_{a^{(p)}}^{x^{(p)}}du_h-
\int_{\mu}^{\xi}du_h\qquad(h=1,2\ldots p),
\]
and denote by \(w_h^\circ,w_h^{(1)}\ldots w_h^{(p)}\) the systems which arise successively from this system when first \(\xi\) is replaced by \(\eta\), then also \(x^{(1)}\) by \(a^{(1)}\), then in addition \(x^{(2)}\) by \(a^{(2)}\), and so on. From the quotient
\[
F(w)=\frac{\Theta(w+m^{(1)})\,\Theta(w+m^{(2)})\cdots\Theta(w+m^{(q)})}{\Theta^q(w)}
\]
there arise in this way the following functions, in whose arguments the adjoining points occur as upper limits:
\[
\begin{array}{r@{\ :\ }l}
F(w)&x^{(1)},x^{(2)}\ldots x^{(p)},\xi,\\
F(w^\circ)&x^{(1)},x^{(2)}\ldots x^{(p)},\eta,\\
F(w^{(1)})&a^{(1)},x^{(2)}\ldots x^{(p)},\eta,\\
\vdots&\vdots\\
F(w^{(p)})&a^{(1)},a^{(2)}\ldots a^{(p)},\eta.
\end{array}
\]
By Theorems I and II of \S\,63, the quotient of each two successive functions \(F\) is now an algebraic function of its arguments; specifically one has
\[
\tag{1}
\begin{aligned}
\frac{F(w)}{F(w^\circ)}&=\frac{\varphi_0(\xi)\psi_0(\eta)}{\varphi_0(\eta)\psi_0(\xi)},\\
\frac{F(w^\circ)}{F(w^{(1)})}&=\frac{\varphi_1(x^{(1)})\psi_1(a^{(1)})}{\varphi_1(a^{(1)})\psi_1(x^{(1)})},\\
&\hspace{2em}\cdots\\
\frac{F(w^{(p-1)})}{F(w^{(p)})}&=\frac{\varphi_p(x^{(p)})\psi_p(a^{(p)})}{\varphi_p(a^{(p)})\psi_p(x^{(p)})}.
\end{aligned}
\]
The mode of formation of the curves \(\varphi_i=0\), \(\psi_i=0\) is given in \S\,63; let us here only note that the coefficients of the individual functions depend rationally on the coordinates of the points listed below:
\[
\begin{array}{r@{,\ }l@{\quad}l}
\varphi_0&\psi_0&\text{on }x^{(1)},x^{(2)}\ldots x^{(p)},\\
\varphi_1&\psi_1&\text{on }a^{(1)},x^{(2)}\ldots x^{(p)},\\
\varphi_2&\psi_2&\text{on }a^{(1)},a^{(2)}\ldots x^{(p)},\\
\multicolumn{3}{c}{\cdots}\\
\varphi_p&\psi_p&\text{on }a^{(1)},a^{(2)}\ldots a^{(p)}.
\end{array}
\]
Multiplying equations (1) together, we obtain the equation
\[
\tag{2}
\frac{F(w)}{F(w^{(p)})}=\frac{
\dfrac{\varphi_0(\xi)\varphi_1(x^{(1)})\cdots\varphi_p(x^{(p)})}{\psi_0(\xi)\psi_1(x^{(1)})\cdots\psi_p(x^{(p)})}}
{\dfrac{\varphi_0(\eta)\varphi_1(a^{(1)})\cdots\varphi_p(a^{(p)})}{\psi_0(\eta)\psi_1(a^{(1)})\cdots\psi_p(a^{(p)})}}.
\]
If now we regard \(\eta,a^{(1)},a^{(2)}\ldots a^{(p)}\) as constant points, we obtain \(F(w)\) as a rational function of the coordinates of the points \(x^{(1)},x^{(2)}\ldots\xi\). Finally, if we divide two functions \(F\) which differ by different values of the \(m\), we obtain the following general theorem:

\begin{quote}
The quotient of two complete \(\Theta\)-products with the same number of factors and with the same variables \(w\) is a rational function of the coordinates of the points \(x^{(1)},x^{(2)}\ldots x^{(p)},\xi\), which are connected by the equations
\[
w_h=\int_{\alpha^{(1)}}^{x^{(1)}}du_h+
\int_{\alpha^{(2)}}^{x^{(2)}}du_h+\cdots+
\int_{\alpha^{(p)}}^{x^{(p)}}du_h-
\int_{\mu}^{\xi}du_h
\]
\[(h=1,2\ldots p).
\]
\end{quote}

This is the extension to Abelian functions of the theorem that Mr. Hermite introduced into the theory of elliptic functions as the basis for a large series of further theorems.

The algebraic function is necessarily symmetric in the points \(x^{(1)},x^{(2)}\ldots x^{(p)}\), although it is not so in its form and construction. But, for example, one can give such a symmetric form to the left-hand side of equation (2) by permuting the \(x\), adding the expressions arising in this way, and dividing by their number.

From this one can conclude that also when
\[
w_h=\int_{\gamma^{(1)}}^{z^{(1)}}du_h+
\int_{\gamma^{(2)}}^{z^{(2)}}du_h+\cdots+
\int_{\gamma^{(s)}}^{z^{(s)}}du_h\qquad(h=1,2\ldots p)
\]
is given, this function is still a rational function of the coordinates of \(z^{(1)},z^{(2)}\ldots z^{(s)}\). For if one determines the \(x\) from the equations
\[
\sum_{i=1}^{i=s}\int_{\gamma^{(i)}}^{z^{(i)}}du_h=\sum_{i=1}^{i=p}\int_{\alpha^{(i)}}^{x^{(i)}}du_h-\int_{\mu}^{\xi}du_h,
\]
which again represent Abel's theorem, then by \S\,64 one can express the symmetric functions of the coordinates of the \(x\) as rational functions of the coordinates of the \(z,\xi\), and hence also express the right-hand side of (2) rationally by the coordinates of the \(z\), whereupon \(\xi\) must drop out entirely.

\subsection*{\S\,66. Relations between \(\Theta\)-functions.}

According to the preceding, we can represent the ratios of complete \(\Theta\)-products as rational symmetric functions of the coordinates of \(p\) points \(x\). For the moment we disregard the point \(\xi\). It is arbitrary; one can, for instance, let it coincide with \(\mu\), and thus make it disappear completely from the expressions for the \(w\).

Let \(P_1(w),P_2(w)\ldots P_{p+2}(w)\) denote \(p+2\) complete \(\Theta\)-products which differ from one another only by the values of the \(m\). We then have formulae of the following kind:
\[
\tag{1}
\begin{aligned}
M P_1(w)&=\Phi_1(x^{(1)},x^{(2)}\ldots x^{(p)}),\\
M P_2(w)&=\Phi_2(x^{(1)},x^{(2)}\ldots x^{(p)}),\\
&\hspace{2em}\cdots\\
M P_{p+2}(w)&=\Phi_{p+2}(x^{(1)},x^{(2)}\ldots x^{(p)}),
\end{aligned}
\]
where \(M\) is a common factor, whose nature is not to be further investigated here, while the \(\Phi\) denote rational and symmetric functions of the coordinates of the \(p\) points \(x\). If one combines with these equations the equations
\[
\tag{2}
f(x^{(1)})=0,\quad f(x^{(2)})=0\ldots f(x^{(p)})=0,
\]
or the symmetric equations which can be composed from them, one can eliminate from (1), (2) the \(x\) and the proportionality factor \(M\). One therefore obtains the theorem:

\begin{quote}
Between any \(p+2\) complete \(\Theta\)-products with the same number of factors there exists a homogeneous algebraic equation.
\end{quote}

At the same time, in general, during the elimination process one obtains the symmetric functions of the coordinates of the \(x\) as homogeneous rational functions of degree zero of the given \(p+2\) \(\Theta\)-products. Hence one can conceive the inversion problem in a new way, by saying:

\begin{quote}
The symmetric functions of the coordinates of the \(x\), and thus also the coefficients of an equation whose roots are the function-values \(\varphi(x^{(1)}),\varphi(x^{(2)})\ldots\varphi(x^{(p)})\), can be represented as homogeneous rational functions of \(p+2\) complete \(\Theta\)-products, in general to be chosen arbitrarily, with the same number of factors, among which an algebraic equation itself takes place.
\end{quote}

It is noteworthy that here a single irrationality occurs in the expressions for the \(x\) through the \(p+1\) ratios of the complete \(\Theta\)-products, namely that which is brought about by the relation mentioned. If one introduces these \(p+1\) ratios as new algebraic quantities, then one has expressed the symmetric functions of \(p\) arbitrary points of the curve algebraically and rationally through these \(p+1\) new variables, while they themselves are connected by an equation of higher degree.

Among the simplest relations between \(\Theta\)-products are those which arise from the formulae of the inversion problem itself, as they were presented in \S\,57 and \S\,59 of the preceding section. Let \(\varphi-\lambda\psi=0\) denote curves of a pencil, and let \(\xi^{(i)}\) be the intersection-points of \(f=0\) with a curve of the pencil whose parameter is \(\lambda^{(i)}\). Then, by \S\,57, we have the formulae
\[
\tag{3}
\begin{aligned}
&M\,[\varphi(x^{(1)})-\lambda^{(i)}\psi(x^{(1)})]\cdot[\varphi(x^{(2)})-\lambda^{(i)}\psi(x^{(2)})]\cdots[\varphi(x^{(p)})-\lambda^{(i)}\psi(x^{(p)})]\\
&\qquad = C^{(i)}\prod_{\xi^{(i)}}\Theta\!\biggl(v_h-\int_{\mu}^{\xi^{(i)}}du_h\biggr),
\end{aligned}
\]
where \(M\) is an undetermined factor, and furthermore
\[
v_h=\sum_{i=1}^{i=p}\int_{a^{(i)}}^{x^{(i)}}du_h,
\]
and
\[
C^{(i)}=\frac{[\varphi(a^{(1)})-\lambda^{(i)}\psi(a^{(1)})]\cdot[\varphi(a^{(2)})-\lambda^{(i)}\psi(a^{(2)})]\cdots[\varphi(a^{(p)})-\lambda^{(i)}\psi(a^{(p)})]}{\displaystyle\prod_{\xi^{(i)}}\Theta\!\biggl(\int_{\mu}^{\xi^{(i)}}du_h\biggr)}.
\]
The \(\Theta\)'s occurring in (3) form a system of complete \(\Theta\)-products if the \(v\) are increased by suitable constants, for the sum of the arguments with equal \(h\) has, by Abel's theorem, the same value in each product. If now in (3) one substitutes \(p+2\) different values for \(\lambda^{(i)}\), one can eliminate \(\varphi(x^{(1)})\), \(\psi(x^{(1)})\), and obtains between the \(\Theta\)-products the linear relation
\[
\sum_{i=1}^{i=p+2}
\frac{C^{(i)}\prod_{\xi^{(i)}}\Theta\!\bigl(v_h-\int_{\mu}^{\xi^{(i)}}du_h\bigr)}
{(\lambda^{(i)}-\lambda^{(1)})(\lambda^{(i)}-\lambda^{(2)})\cdots(\lambda^{(i)}-\lambda^{(p+2)})}=0.
\]
In a similar way, relations between \(\Theta\)-products with two factors each arise from the formulae of \S\,59. Keeping the notation chosen there, put
\[
K^{(i)}=\frac{\sum \pm \varphi_1(\xi^{(i)})\varphi_2(a^{(1)})\cdots\varphi_{p+1}(a^{(p)})}
{\Theta\!\bigl(v_h^\circ-\int_{\mu}^{\xi^{(i)}}du_h\bigr)\Theta\!\bigl(v'_h-\int_{\mu}^{\xi^{(i)}}du'_h\bigr)}.
\]
Then one has the equation
\[
\begin{aligned}
R_1\varphi_1(\xi^{(i)})+R_2\varphi_2(\xi^{(i)})+\cdots+R_{p+1}\varphi_{p+1}(\xi^{(i)})
= {}& M K^{(i)}\Theta\!\biggl(v_h-\int_{\mu}^{\xi^{(i)}}du_h\biggr)\\
&\cdot\Theta\!\biggl(v_h^\circ+v'_h-v_h-\int_{\mu}^{\xi^{(i)}}du_h\biggr).
\end{aligned}
\]
Hence, if one substitutes \(p+2\) different points for \(\xi^{(i)}\), and then eliminates the \(R\),
\[
0=\sum_{i=1}^{i=p+2}K^{(i)}\Delta^{(i)}\Theta\!\biggl(v_h-\int_{\mu}^{\xi^{(i)}}du_h\biggr)\Theta\!\biggl(v_h^\circ+v'_h-v_h-\int_{\mu}^{\xi^{(i)}}du_h\biggr),
\]
where \(\Delta^{(1)},\Delta^{(2)}\ldots\Delta^{(p+2)}\) are the various determinants to be formed from the system of quantities \(\varphi_k(\xi^{(i)})\).
% --- Continuation: source PDF pp. 252--270, printed pp. 230--248. ---
\clearpage

\section*{Tenth Section.}
\subsection*{Division.}

\subsection*{\S\,67. Multiplication.}

Suppose that a contact curve is to be drawn which touches the curve \(f=0\) at \(p\) given points with contact of multiplicity \(m\), which passes through the double and cusp points of \(f=0\), and for whose intersection system all points except \(p\) are otherwise given. Then, as was explained in the second section, the determination of the remaining points leads to an equation of degree \(p\). Let us consider a fixed curve and a moving curve of this kind; they are to differ only by the \(p\) points of contact and by \(p\) other points. For the fixed curve, moreover, both sets of points are to coincide with the points denoted above by \(\alpha\), belonging to a given point \(\mu\), so that the fixed curve in fact touches at the \(p\) points \(\alpha\) with multiplicity \(m+1\). Let the points of contact of the moving curve be \(y^{(1)},y^{(2)}\ldots y^{(p)}\), and let the remaining intersection points determined by them be \(x^{(1)},x^{(2)}\ldots x^{(p)}\). Put
\[
v_h=\sum_{i=1}^{i=p}\int_{\alpha^{(i)}}^{x^{(i)}}du_h,\qquad
w_h=\sum_{i=1}^{i=p}\int_{\alpha^{(i)}}^{y^{(i)}}du_h,\qquad
(h=1,2\ldots p).
\]
Then the equations of Abel's theorem pass into
\[
v_h+mw_h=0.\qquad(h=1,2\ldots p)
\]

By what precedes, the quotients of complete \(\Theta\)-products of the same order with the arguments \(v\) can be expressed rationally and symmetrically by the coordinates of the points \(x\), then again, by the procedure of \S\,64, rationally and symmetrically by the coordinates of the points \(y\), and finally, by the inversion problem, the expression so transformed can be expressed by complete \(\Theta\)-products with the arguments \(w\). In this way one obtains the following theorem on the multiplication of Abelian functions:

\begin{quote}
The quotient of two complete \(\Theta\)-products of the same order with the arguments \(mv_1,mv_2\ldots mv_p\) can be represented rationally and homogeneously by complete \(\Theta\)-products with the arguments \(v_1,v_2\ldots v_p\).
\end{quote}

If, in particular, according to \S\,66, we express the ratios of \(p+1\) complete \(\Theta\)-products with the arguments \(mv\) by \(p+2\) complete \(\Theta\)-products with the arguments \(v\), among which latter products an algebraic equation exists, then conversely the ratios of the \(\Theta\)-products with the arguments \(v\) can also be expressed by those with the arguments \(mv\). Geometrically, this algebraic problem leads to the solution of the contact problem which we shall treat in what follows, and to which we give the name of the general division problem, since it is always a matter of passing from functions with the arguments \(mv\) to functions with the arguments \(v\), or from functions with the arguments \(v\) to such functions with the arguments \(\frac{v}{m}\).

\subsection*{\S\,68. The problem of general division.}

Let \(c^{(1)},c^{(2)}\ldots c^{(p)}\) be points at which one lets the curve \(f=0\) be touched with multiplicity \(m\) by another curve. This curve cuts \(f=0\) in the double and cusp points and in certain other points \(a^{(1)},a^{(2)}\ldots a^{(q)},b^{(1)},b^{(2)}\ldots b^{(r)}\). Through the points \(b\) and through given points \(x^{(1)},x^{(2)}\ldots x^{(q)}\), as well as through the double and cusp points, we are to draw a curve of the order of the preceding one which touches \(f=0\) at \(p\) points with multiplicity \(m\); the points of contact \(y^{(1)},y^{(2)}\ldots y^{(p)}\) are sought. By Abel's theorem the \(y\) are then given by the transcendental equations
\[
(1)\ldots\quad
m\sum_{i=1}^{i=p}\int_{c^{(i)}}^{y^{(i)}}du_h+
\sum_{i=1}^{i=q}\int_{a^{(i)}}^{x^{(i)}}du_h=0
\qquad(h=1,2\ldots p).
\]

If in particular we understand by the integrals of the second sum the values which they take on fixed, arbitrarily chosen paths of integration, then on the right-hand side of (1) the zero may be replaced by the most general system of periodicity moduli \(P_h\) of the integrals of the first kind. Thus the determination of the \(y\) is given by the following equations, which have the form of the inversion problem:
\[
(2)\ldots\quad
\sum_{i=1}^{i=p}\int_{c^{(i)}}^{y^{(i)}}du_h
=
\frac{1}{m}\left[
P_h-\sum_{i=1}^{i=q}\int_{a^{(i)}}^{x^{(i)}}du_h
\right]\qquad(h=1,2\ldots p).
\]
From this one sees that, for a fixed choice of the periodicity moduli \(P\), the system of the \(y\), and hence the contact curve, is uniquely determined. At the same time, however, one observes that this system does not change if the \(P_h\) are simultaneously increased or diminished by \(m\) times a system of periodicity moduli. Therefore, in order to obtain all possible contact curves, it is only necessary to assign to the integers \(m_1,m_2\ldots m_p,q_1,q_2\ldots q_p\) occurring in the \(P_h\) the values \(0,1\ldots m-1\). In all one therefore obtains \(m^{2p}\) systems of periodicity moduli, and hence just as many contact curves. The proposed problem therefore has \(m^{2p}\) solutions.\footnote{On the geometrical questions connected with this, compare Crelle's Journal, vol. 63, p. 189 ff.} We shall now see how to develop the solutions algebraically.

Let us imagine that the equation of the sought contact curve, whose order may be the \(\nu^{\text{th}}\), has been written with undetermined coefficients, and let us denote this equation by
\[
\varphi=0.
\]
Together with \(f=0\) it gives the \(\nu n\) intersection points of the two curves. If we draw the rays of an arbitrary pencil \(\alpha+\lambda\beta=0\) which are directed toward these points, then by eliminating the \(x\) from the simultaneous equations
\[
\varphi=0,\qquad f=0,\qquad \alpha+\lambda\beta=0
\]
we find for \(\lambda\) an equation of degree \(\nu n\).

The resulting equation in \(\lambda\), which may be denoted by \(L=0\), contains a series of known factors. For the rays \(\lambda\) which are directed toward the given points \(a,b\), and the double rays which are directed toward the double and cusp points, are known beforehand. If one removes from \(L=0\) the factors corresponding to them, there remains an equation \(L'=0\) of degree \(mp\), which gives the remaining intersection points. These are to combine in groups of \(m\) into \(p\) points; the expression \(L'\) must therefore be the \(m^{\text{th}}\) power of an expression \(\Lambda\) of degree \(p\), so that
\[
(3)\ldots\quad L'=\Lambda^m,
\]
and so that finally the equation \(\Lambda=0\) gives the \(p\) points of contact. For in the course of the elimination just indicated one finds the ratios of the coordinates of the intersection points as rational functions of \(\lambda\); hence these coordinates themselves are given once the corresponding values of \(\lambda\) are known.

Now suppose that the equation \(\varphi=0\) is formed in such a way that one draws \((m-1)p+1\) special curves \(\varphi_1=0,\varphi_2=0\ldots\) of order \(\nu\) through the points \(a,b\) and through the double and cusp points, and then forms the combination
\[
\varphi=\varkappa_1\varphi_1+\varkappa_2\varphi_2+\cdots+
\varkappa_{(m-1)p+1}\varphi_{(m-1)p+1}.
\]
Then \(L\) contains only the \((m-1)p+1\) undetermined constants \(\varkappa\), or rather the \((m-1)p\) ratios among them. The equation \(L=0\) splits of itself into two equations, one of which is independent of the \(\varkappa\) and gives the system of known rays, while the other is the equation \(L'=0\), which is to become the \(m^{\text{th}}\) power of the equation \(\Lambda=0\), and does become this through a suitable determination of the quantities \(\varkappa\).

The left-hand side of equation (3) therefore contains in a rational way the \(\varkappa\) and the coordinates of the \(x\). In order that \(L'\) become the \(m^{\text{th}}\) power of an expression of order \(p\), \((m-1)p\) conditions must be satisfied. In fact, \(\Lambda\) still contains \(p+1\) arbitrary coefficients; if one now equates the corresponding coefficients of \(L'\) and \(\Lambda^m\), one finds \(mp+1\) equations in which the \(p+1\) coefficients of \(\Lambda\) and the \((m-1)p+1\) quantities \(\varkappa\) occur as unknowns. If one eliminates the former, \((m-1)p\) equations remain, which contain only the ratios of the \(\varkappa\) and serve to determine them. If from the equations just considered one also eliminates \(\varkappa_3,\varkappa_4\ldots\varkappa_{(m-1)p+1}\), a single equation finally remains which contains only the quotient of \(\varkappa_1\) and \(\varkappa_2\) as unknown.

This equation necessarily decomposes into factors, according to the character of its roots. For even though it was shown above that, for every one of the contact curves sought, equation (3) must be identically satisfied, the converse by no means follows. Rather, \(m\) rays can also coincide because the \(m\) intersection points corresponding to them distribute themselves in some grouped manner among the points which such a ray has in common with the curve \(f=0\). But those curves \(\varphi=0\) for which something of this kind occurs plainly depend on the position of the vertex of the pencil \(a=0,b=0\) and change with it, whereas the contact curves sought are completely independent of the position of this vertex. Hence one sees that the different equations obtained for \(\varkappa_1:\varkappa_2\), when different pencils \(a,b\) are used, must agree in part of their roots and not in another part, and that only the first part is used for the question treated here. Thus one need only form the equation described above for two different pencils \(a,b\), and seek the common factor of the two; this factor gives the equation for \(\varkappa_1:\varkappa_2\), which supplies the contact curves sought.

The new equation so obtained is necessarily of degree \(m^{2p}\), as follows from the number of contact curves found above. We shall denote this equation by \(K=0\).

The ratios can moreover be arranged so that the unknown \(\varkappa_1:\varkappa_2\) of this equation acquires a simple geometrical meaning. For this purpose it is only necessary to choose the curves \(\varphi_3=0,\varphi_4=0\ldots\varphi_{(m-1)p+1}=0\) so that they have double points at a fixed one of the common intersection points \(a,b\). Since this involves two further conditions, there are in fact \((m-1)p-1\) such curves among whose equations no linear relation exists. In this case, at the point in question the differential quotients of \(\varphi_3,\varphi_4\), etc., vanish, and the equation of the tangent of \(\varphi=0\) at this point becomes
\[
\begin{aligned}
(4)\ldots\quad&
\varkappa_1\left(
\frac{\partial\varphi_1}{\partial x_1}X_1+
\frac{\partial\varphi_1}{\partial x_2}X_2+
\frac{\partial\varphi_1}{\partial x_3}X_3
\right)\\
&+\varkappa_2\left(
\frac{\partial\varphi_2}{\partial x_1}X_1+
\frac{\partial\varphi_2}{\partial x_2}X_2+
\frac{\partial\varphi_2}{\partial x_3}X_3
\right)=0.
\end{aligned}
\]
Thus the values of \(\varkappa_1:\varkappa_2\) corresponding to the various contact curves are the parameters of the tangents drawn at that point to these curves, tangents which belong to the pencil of rays represented by equation (4).

The search for the points \(y\) therefore leads to two algebraic equations. The first, with the unknown \(\varkappa_1:\varkappa_2\), is of degree \(m^{2p}\); its coefficients contain the coordinates of the points \(x\) in a rational and symmetric way. If a root of this equation is known, then \(\varkappa_3,\varkappa_4\), etc., are rational functions of it and of the coordinates of the \(x\); hence, in general, the equation of the corresponding contact curve is represented rationally by this root and by the coordinates of the points \(x\). Secondly, one must solve the equation \(\Lambda=0\) of degree \(p\), which gives the \(p\) corresponding values of \(\lambda\), and hence also, expressed rationally by each such value, the coordinates of the contact points. The coefficients of this equation likewise contain the coefficients of the points \(x\) in a rational and symmetric way.

\subsection*{\S\,69. The problem of special division. Reduction of general division to special division.}

Now, as Jacobi has shown for the elliptic functions (Crelle's Journal, vol. 3, p. 86), and Hermite for the hyperelliptic functions (ibid. vol. 32, p. 293), the solution of the first equation can be accomplished by extraction of roots, as soon as a special division problem is regarded as solved. This special problem arises when the points \(x\) are allowed to coincide with the points \(a\); that is, when, starting from the curve tangent at the points \(c\), one seeks the remaining \(m^{2p}-1\) contact curves which pass through the other intersection points of the first-mentioned curve with \(f=0\), in particular also through the double and cusp points. If \(z^{(1)},z^{(2)}\ldots z^{(p)}\) denote the contact points of such a curve, then corresponding to equation (2) of \S\,68 one has, for their definition, the equations
\[
(1)\ldots\quad
\sum_{i=1}^{i=p}\int_{c^{(i)}}^{z^{(i)}}du_h=\frac{P_h}{m}
\qquad(h=1,2\ldots p),
\]
which, when \(P_h=0\), give the self-evident solution \(z^{(1)}=c^{(1)},z^{(2)}=c^{(2)}\ldots z^{(p)}=c^{(p)}\), while for every other system of the \(P\) they give a definite contact curve. We shall first see how to reduce general division to this special division, and then how the problem of special division itself is to be treated.

For the moment let \(y\) denote the contact points of a fixed, arbitrarily chosen contact curve of the general problem; let \(\eta\) denote those of another curve corresponding to the period system \(P_h'\); and finally let \(\xi\) denote the contact points of a curve of the special problem which belongs to the period system \(P_h-P_h'\). From Abel's theorem one then has
\[
\sum_{i=1}^{i=p}\left(
\int_{c^{(i)}}^{y^{(i)}}du_h-
\int_{c^{(i)}}^{\eta^{(i)}}du_h-
\int_{c^{(i)}}^{\xi^{(i)}}du_h
\right)=0,
\]
or
\[
\sum_{i=1}^{i=p}\left(
\int_{y^{(i)}}^{\eta^{(i)}}du_h+
\int_{c^{(i)}}^{\xi^{(i)}}du_h
\right)=0.
\]
Thus, when the \(y\) are given, one finds the points \(\eta\) as follows. Through the \(y\) and \(c\) one draws a curve which also contains the double and cusp points; through the latter and the remaining intersection points of this curve with \(f=0\) one draws a curve of the same order, which also passes through the \(\xi\); it cuts \(f=0\) in the points \(\eta\). By using in succession all systems of the \(\xi\), one thus obtains from one system \(y\) all the others, that is, from an arbitrarily chosen contact curve of the general problem all the remaining ones. The symmetric functions of the coordinates of the \(\eta\) can therefore be represented rationally by the symmetric functions of the coordinates of the \(y\) and by the symmetric functions of the coordinates of the \(\xi\).

Now to each system of the \(y\), that is, to each contact curve of the general problem, there corresponded above a definite root
\[
\varkappa=\frac{\varkappa_2}{\varkappa_1}
\]
of the equation \(K=0\). The symmetric functions of the coordinates of the \(y\) could be expressed rationally through \(\varkappa\), and therefore also conversely, in infinitely many ways, \(\varkappa\) could be represented as a rational and symmetric function of the coordinates of the \(y\). In a similar manner the \(\eta\) stand in connection with a second root \(\varkappa'\) of that equation.

But if in the equations of the preceding paragraphs we let the \(a\) coincide everywhere with the \(x\), the problem of general division passes into the special one. The equation \(K=0\) is transformed into another equation \(M=0\), of which we know one root, since the tangent direction of the one, originally given, contact curve is known. If under \(M=0\) we understand the equation of degree \(m^{2p}-1\) which remains after removal of the known linear factor, then to every other root \(\mu\) of this equation there corresponds a system of the \(\xi\) such that the symmetric functions of the coordinates of these points are rational functions of \(\mu\). Thus, if we found above that the symmetric functions of the \(y\) can be rationally expressed by those of the \(\eta\) and those of the \(\xi\), we can now state this as the following property of the equation \(K=0\):

\begin{quote}
By means of any root \(\varkappa\) of the equation \(K=0\) we can express all the other roots of this equation rationally, namely with the help of the formula
\[
\varkappa'=\Theta(\varkappa,\mu),
\]
where \(\Theta\) is a certain rational function, \(\mu\) is a root of the equation \(M=0\) of degree \(m^{2p}-1\), and all roots \(\varkappa'\) are obtained in order by putting for \(\mu\) all roots of this equation in succession.
\end{quote}

All roots of the equation \(K=0\) are therefore rational functions of one of them. Let \(\varkappa'\), \(\varkappa''\) be two roots obtained by replacing \(\mu\) by two roots \(\mu'\), \(\mu''\) of \(M=0\), so that
\[
(2)\ldots\quad
\varkappa'=\Theta(\varkappa,\mu'),\qquad \varkappa''=\Theta(\varkappa,\mu'').
\]
Then, by putting \(\varkappa'\) or \(\varkappa''\) in place of \(\varkappa\), one can form other roots of \(K=0\). It is, however, especially easy to see that the two roots
\[
\Theta(\varkappa'',\mu')=\Theta[\Theta(\varkappa,\mu''),\mu'],
\qquad
\Theta(\varkappa',\mu'')=\Theta[\Theta(\varkappa,\mu'),\mu'']
\]
are identical; thus, when different roots of \(M=0\) are used alternately, the order in which they are used makes no difference.

Indeed, from the root \(\varkappa\) belonging to the period system \(P_h\), the use of a root \(\mu'\), belonging to the period system \(P_h-P_h'=Q_h'\), produces a root \(\varkappa'\), which corresponds to the period system \(P_h'=P_h-Q_h'\). Similarly, from \(\varkappa\), the use of the root \(\mu''\) belonging to \(Q_h''\) produces the root \(\varkappa''\), which corresponds to the period system \(P_h-Q_h''\). If in the expression for \(\varkappa'\) one replaces \(\varkappa\) by the root \(\varkappa''\), the period system \(P_h\) is to be replaced by \(P_h-Q_h''\), and one finds the root belonging to the period system \(P_h-Q_h'-Q_h''\):
\[
\varkappa'''=\Theta(\varkappa'',\mu')=\Theta[\Theta(\varkappa,\mu''),\mu'].
\]
Conversely, if in the expression for \(\varkappa''\) one puts the root \(\varkappa'\) in place of \(\varkappa\), one must replace \(P_h\) by \(P_h-Q_h'\), and one arrives at the root belonging to the same period system \(P_h-Q_h'-Q_h''\), but now with the expression
\[
\varkappa'''=\Theta(\varkappa',\mu'')=\Theta[\Theta(\varkappa,\mu'),\mu''].
\]
This proves the theorem just stated.

Secondly, if one always uses the same root \(\mu\) to form the successive roots of \(K=0\),
\[
\varkappa'=\Theta(\varkappa,\mu),\quad
\varkappa''=\Theta(\varkappa',\mu),\quad
\varkappa'''=\Theta(\varkappa'',\mu)\ldots,
\]
then one finds that this cycle closes after \(m\) repetitions, but not earlier, so that
\[
\varkappa^{(m)}=\varkappa.
\]
For if we again denote by \(P_h\) the period system belonging to \(\varkappa\), and by \(Q_h\) that belonging to \(\mu\), then to \(\varkappa'\) belongs the period system \(P_h-Q_h\), to \(\varkappa''\) the system \(P_h-2Q_h\), and finally to \(\varkappa^{(m)}\) the period system \(P_h-mQ_h\), whose \(m^{\text{th}}\) part differs from \(\frac{P}{m}\) only by integral periodicity moduli, which may be omitted. The period system belonging to \(\varkappa^{(m)}\) is therefore identical with that belonging to \(\varkappa\), and consequently \(\varkappa^{(m)}=\varkappa\), as was to be proved.

\subsection*{\S\,70. Solution of the general division equation.}

The equation \(K=0\) therefore belongs to the class of those which, as Abel (Crelle's Jornal, vol. 4, p. 131) has shown, are algebraically solvable. To present the solution more precisely, we introduce a notation for the roots \(\varkappa,\mu\) which represents them in their dependence on the systems of periodicity moduli. If
\[
P_h=2m_h\pi\sqrt{-1}+a_{1h}q_1+a_{2h}q_2\cdots+a_{ph}q_p,
\]
then we denote the roots \(\varkappa,\mu\) belonging to this period system by
\[
{}^{\varkappa}_{m_1,m_2\ldots,q_1,q_2\ldots},\qquad
{}^{\mu}_{m_1,m_2\ldots,q_1,q_2\ldots}.
\]
Equation (2) then passes into
\[
\Theta\!\left(
\varkappa_{m_1,m_2\ldots,q_1,q_2\ldots},
\mu_{n_1,n_2\ldots,r_1,r_2\ldots}
\right)
=
\varkappa_{m_1-n_1,m_2-n_2\ldots,q_1-r_1,q_2-r_2\ldots}.
\]
Now consider the sum
\[
\begin{aligned}
(3)\ldots\quad S_{a_1a_2\ldots b_1b_2\ldots}
&=\sum_{m_i,q_i}
e^{2(a_1m_1+a_2m_2\ldots+b_1q_1+b_2q_2\ldots)
\frac{\pi\sqrt{-1}}{m}}\\
&\qquad\cdot\varkappa_{m_1,m_2\ldots,q_1,q_2\ldots},
\end{aligned}
\]
in which the \(a,b\) are integers, and the \(m_i,q_i\) are each independently to assume all values from \(0\) to \(m-1\). This sum can also be written in the form
\[
\begin{aligned}
(4)\ldots\quad S_{a_1a_2\ldots b_1b_2\ldots}
&=\sum_{n_i,r_i}
e^{2\left[a_1(m_1-n_1)\ldots+b_1(q_1-r_1)\ldots\right]
\frac{\pi\sqrt{-1}}{m}}\\
&\qquad\cdot\varkappa_{m_1-n_1,\ldots,q_1-r_1,\ldots}\\
&=e^{2\left[a_1m_1\ldots+b_1q_1\ldots\right]\frac{\pi\sqrt{-1}}{m}}
\sum_{n_i,r_i}
e^{-2(a_1n_1\ldots+b_1r_1\ldots)\frac{\pi\sqrt{-1}}{m}}\\
&\qquad\cdot
\Theta\!\left(
\varkappa_{m_1\ldots q_1\ldots},
\mu_{n_1\ldots r_1\ldots}
\right),
\end{aligned}
\]
where the \(n_i,r_i\) must now independently take all values from \(0\) to \(m-1\). The last expression is a rational function of \(\varkappa_{m_1,m_2\ldots q_1q_2\ldots}\), and indeed its \(m^{\text{th}}\) power depends on the numbers \(m_i,q_i\) only insofar as these designate the root preferred over the others. But the expression \(S_{a_1a_2\ldots b_1b_2\ldots}\) is unchanged if this root is replaced by another one. Thus, if one successively substitutes for \(\varkappa_{m_1,m_2\ldots q_1q_2\ldots}\) all roots of \(K=0\), and adds the \(m^{\text{th}}\) powers of the expressions so obtained from the above, one obtains \(m^{2p}\) multiplied by the \(m^{\text{th}}\) power of \(S_{a_1a_2\ldots b_1b_2\ldots}\) as a symmetric function of these roots, hence as a rational function of the coefficients of \(K=0\). Therefore
\[
(5)\ldots\quad
S_{a_1a_2\ldots b_1b_2\ldots}
=
\sqrt[m]{T_{a_1a_2\ldots b_1b_2\ldots}},
\]
where \(T_{a_1a_2\ldots b_1b_2\ldots}\) is a function which contains the roots \(\mu\) of the equation \(M=0\), but depends on the coordinates of the points \(x\) only in a rational and symmetric way.

If now equation (5) is formed for all systems of values of the \(a,b\), in which these numbers independently take the values \(0,1,\ldots,m-1\), then the \(m^{2p}\) equations so obtained become just as many linear equations for determining the \(\varkappa\), and by solving them one obtains
\[
\begin{aligned}
m^{2p}\varkappa_{m_1,m_2\ldots q_1q_2\ldots}
&=\sum_{a_i,b_i}
e^{-2(a_1m_1+a_2m_2\ldots+b_1q_1+b_2q_2\ldots)
\frac{\pi\sqrt{-1}}{m}}\\
&\qquad\cdot
\sqrt[m]{T_{a_1a_2\ldots b_1b_2\ldots}}.
\end{aligned}
\]
Only the term \(\sqrt[m]{T_{00\ldots00\ldots}}\) is to be replaced by a rational function; for \(S_{00\ldots00\ldots}\) is nothing other than the sum of all \(\varkappa\), and therefore a coefficient of the equation \(K=0\) itself.

Finally, the \(m^{2p}-1\) root quantities which occur here can be reduced to only \(2p\), which then are ultimately independent of one another and, in their combinations, give the different roots \(\varkappa\). Namely, let us start from the \(2p\) expressions \(S\) in which only one index at a time is different from zero:
\[
S_{10\ldots00\ldots},\quad
S_{01\ldots00\ldots},\quad\ldots,\quad
S_{00\ldots10\ldots},\quad
S_{00\ldots01\ldots},\quad\ldots.
\]
Then
\[
\frac{
S_{a_1a_2\ldots b_1b_2\ldots}
}{
S_{10\ldots00\ldots}^{a_1}
S_{01\ldots00\ldots}^{a_2}\cdots
S_{00\ldots10\ldots}^{b_1}
S_{00\ldots01\ldots}^{b_2}\cdots
}
=U_{a_1a_2\ldots b_1b_2\ldots}
\]
is a rational function of \(\varkappa_{m_1,m_2\ldots q_1q_2\ldots}\), which, formed according to (4), immediately shows that it is independent of the choice of the root \(\varkappa\), and hence can be expressed rationally by the coefficients of the equation \(K=0\). With the help of this equation one also obtains
\[
\begin{aligned}
\sqrt[m]{T_{a_1a_2\ldots b_1b_2\ldots}}
&=
U_{a_1a_2\ldots b_1b_2\ldots}\\
&\quad\cdot
\sqrt[m]{
T_{10\ldots00\ldots}^{a_1}
T_{01\ldots00\ldots}^{a_2}
\cdots
T_{00\ldots10\ldots}^{b_1}
T_{00\ldots01\ldots}^{b_2}
\cdots},
\end{aligned}
\]
so that in the finite expressions only the \(m^{\text{th}}\) roots from the \(2p\) expressions
\[
T_{10\ldots00\ldots},\quad
T_{01\ldots00\ldots},\quad\ldots,\quad
T_{00\ldots10\ldots},\quad
T_{00\ldots01\ldots},\quad\ldots
\]
actually remain; each of them can then be chosen in \(m\) different ways, and thus gives rise to the \(m^{2p}\) different roots \(\varkappa\).

\subsection*{\S\,71. Reduction of the special division equation to the case in which \(m\) is a power of a prime.}

We now come to the problem of special division, that is, to the investigation of the equation \(M=0\), which in solving the equation \(K=0\) we presupposed as solved. The roots of this equation \(M=0\) also stand in remarkable relations to one another. First one sees that the period systems can always be composed linearly and integrally from \(2p\) systems, provided only that the determinant of the integers occurring in these \(2p\) systems is equal to 1. If we now replace these \(2p\) period systems by the integral sums corresponding to them, equation (1) passes into a form in which the integral sums corresponding to any special contact curve are composed linearly and integrally from those \(2p\) systems of integral sums. By Abel's theorem, therefore, the symmetric functions of the coordinates of the contact points of any one of the curves can be expressed by the symmetric functions of the coordinates of the contact points of such \(2p\) curves chosen in various ways; or, all roots of the equation \(M=0\) can always be expressed rationally by such \(2p\) roots of it whose associated period systems have determinant 1.

It should also be mentioned that only one special division problem has to be solved for each number \(m\). If one has determined the points \(z\) from the equations
\[
\sum_{i=1}^{i=p}\int_{c^{(i)}}^{z^{(i)}}du_h=\frac{P_h}{m},
\]
then for other starting points \(\gamma\) one immediately finds the point systems \(\xi\), satisfying the equations
\[
\sum_{i=1}^{i=p}\int_{\gamma^{(i)}}^{\xi^{(i)}}du_h=\frac{P_h}{m},
\]
by Abel's theorem, by subtracting this equation from the previous one.

But the character of the equation \(M=0\) depends entirely on the character of the number \(m\). If \(\nu_1,\nu_2\ldots\) are primes and
\[
m=\nu_1^{\alpha_1}\nu_2^{\alpha_2}\cdots,
\]
then the treatment of the given equation \(M=0\), which supplies the special division for the number \(m\), can always be reduced to the equations which give the division for the individual prime powers \(\nu_1^{\alpha_1},\nu_2^{\alpha_2}\ldots\). For one can always give the characteristic period system belonging to a root of the equation \(M=0\) the form
\[
(1)\ldots\quad
\frac{P'_h}{m}
=
\frac{P_h^{(1)}}{\nu_1^{\alpha_1}}
+
\frac{P_h^{(2)}}{\nu_2^{\alpha_2}}
+\cdots,
\]
where the integers occurring in \(P^{(1)},P^{(2)}\) are always smaller than the corresponding denominators; and one obtains all period systems \(\frac{P_h}{m}\) if, independently of one another, one assigns to the integers contained in \(P_h^{(1)}\) all values from \(0\) to \(\nu_1^{\alpha_1}-1\), to those contained in \(P_h^{(2)}\) all values from \(0\) to \(\nu_2^{\alpha_2}-1\), etc. But if in (1) the different period systems are replaced by the integral sums corresponding to the contact curves, then (1) gives a form of Abel's theorem which teaches how to compose the contact curve of the special problem for the number \(m\) from those for the numbers \(\nu_1^{\alpha_1},\nu_2^{\alpha_2}\ldots\), or which gives the roots of the equation \(M=0\) as rational functions of those roots obtained from the corresponding equations for \(\nu_1^{\alpha_1},\nu_2^{\alpha_2}\ldots\).

In particular, one observes that in equation (1) the general denominator always becomes smaller than \(m\) as soon as, in one of the period systems \(P_h^{(1)},P_h^{(2)}\ldots\), all integers become zero or contain a power of the corresponding denominator as a factor. Such cases, in which \(P_h\) too contains a factor of \(m\) in all its integers, do not give genuine contact curves of the required order. For if in such a case one divides the original equation
\[
m\sum_{i=1}^{i=p}\int_{c^{(i)}}^{z^{(i)}}du_h=P_h
\]
by this common factor \(\delta\), one obtains the equation
\[
\frac{m}{\delta}\sum_{i=1}^{i=p}\int_{c^{(i)}}^{z^{(i)}}du_h=P_h^{(1)},
\]
which shows that the curve corresponding to the \(z\) does not touch with multiplicity \(m\), but only with multiplicity \(\frac{m}{\delta}\), and therefore can enter the present problem only insofar as, counted \(\delta\)-fold, it represents an improper curve touching with multiplicity \(m\). The number of these cases to be excluded, for which the corresponding roots \(\mu\) already occur as roots of lower division equations, is obtained by assigning, in one of the period systems \(P_h^{(1)},P_h^{(2)}\ldots\), to the integers only values divisible by the corresponding prime factor, while the integers of the remaining period systems remain arbitrary.

Now in a period system \(P_h^{(i)}\) the number of systems of integers whose all entries are not divisible by \(\nu_i\) is easy to determine. For the first of the integers occurring in it one chooses any integer not divisible by \(\nu_i\) and smaller than \(\nu_i^{\alpha_i}\), and takes the other \(2p-1\) integers arbitrarily; then one takes for the first number one divisible by \(\nu_i\), for the second one not divisible by \(\nu_i\), and the other \(2p-2\) arbitrarily; then for the first two numbers ones divisible by \(\nu_i\), for the third one not divisible by it, and the remaining \(2p-3\) arbitrarily, and so on. In this way one obtains the following number of combinations which belong to genuine contact curves:
\[
\begin{aligned}
&\nu_i^{\alpha_i-1}(\nu_i-1)\nu_i^{\alpha_i(2p-1)}
+\nu_i^{2(\alpha_i-1)}(\nu_i-1)\nu_i^{\alpha_i(2p-2)}\\
&\quad+\nu_i^{3(\alpha_i-1)}(\nu_i-1)\nu_i^{\alpha_i(2p-3)}
+\cdots+\nu_i^{2p(\alpha_i-1)}(\nu_i-1)\\
&\qquad=\nu_i^{2p(\alpha_i-1)}(\nu_i^{2p}-1).
\end{aligned}
\]
Therefore the total number of genuine contact curves occurring in the special division problem is
\[
\nu_1^{2p(\alpha_1-1)}\nu_2^{2p(\alpha_2-1)}\cdots
(\nu_1^{2p}-1)(\nu_2^{2p}-1)\cdots.
\]

\subsection*{\S\,72. Solution of the special division equation for a prime and powers of one.}

It remains to solve the equation \(M=0\) under the hypothesis that \(m\) is a power of a prime, \(m=\nu^\alpha\). This problem can always be reduced to the special division for \(m=\nu\) and to the general division treated above. Namely, suppose that the special division for \(m=\nu\) has been carried out, that is, suppose that the symmetric functions of the \(z\) have been determined so that
\[
\nu\sum_{i=1}^{i=p}\int_{c^{(i)}}^{z^{(i)}}du_h=P_h,
\]
and let us further determine, by general division for \(m=\nu\), another \(\alpha-1\) point systems such that
\[
\left\{
\begin{aligned}
\nu\sum_{i=1}^{i=p}\int_{c^{(i)}}^{z_1^{(i)}}du_h
&=
\sum_{i=1}^{i=p}\int_{c^{(i)}}^{z^{(i)}}du_h+P_h^{(1)},\\
\nu\sum_{i=1}^{i=p}\int_{c^{(i)}}^{z_2^{(i)}}du_h
&=
\sum_{i=1}^{i=p}\int_{c^{(i)}}^{z_1^{(i)}}du_h+P_h^{(2)},\\
&\ \vdots\\
\nu\sum_{i=1}^{i=p}\int_{c^{(i)}}^{z_{\alpha-1}^{(i)}}du_h
&=
\sum_{i=1}^{i=p}\int_{c^{(i)}}^{z_{\alpha-2}^{(i)}}du_h+P_h^{(\alpha-1)}.
\end{aligned}
\right.
\tag{2}
\]
By combining all these equations one obtains
\[
\nu^\alpha\sum_{i=1}^{i=p}\int_{c^{(i)}}^{z_{\alpha-1}^{(i)}}du_h
=
P_h+\nu P_h^{(1)}+\nu^2P_h^{(2)}+\cdots+\nu^{\alpha-1}P_h^{(\alpha-1)}.
\]
The system of the \(z_{\alpha-1}\) therefore gives the most general contact curve of the special problem for \(m=\nu^\alpha\). The equations (2), as equations of Abel's theorem, also show that the roots of the special division equation for \(m=\nu^\alpha\) are found by repeated extraction of roots from rational functions of the roots of the division equation for \(m=\nu\). Thus this last equation remains to be treated.

The special division for a prime \(\nu\) leads to an equation in \(\mu\) of degree \(\nu^{2p}-1\). The period systems corresponding to its roots are formed by assigning to the integers occurring in the \(P_h\), independently of one another, only the values \(0,1,\ldots,\nu-1\), while excluding only the one case in which all integers would have the value 0; this case corresponds to the given contact curve from which one starts. If the quantities \(P_h\) denote any period system, then
\[
P_h,\quad 2P_h,\quad 3P_h,\ldots,\quad(\nu-1)P_h
\]
form a group of period systems which does not change if, instead of the system \(P_h\), one takes as basis any other system of the group. By the repeatedly used application of Abel's theorem, every root belonging to this group can be expressed rationally by every other one, and a symmetric function \(A\) of the roots of this group is expressible as the same rational function of each of the \(\nu-1\) roots of the group.

The number of groups, which can never have a period system or a root in common, is
\[
\frac{\nu^{2p}-1}{\nu-1}.
\]
If one forms the symmetric function \(A\) for all groups and examines the sum of equal powers of these \(A\), multiplied by \((\nu-1)\), by putting for \((\nu-1)A\) the sum of the \(\nu-1\) representations of \(A\) by the individual roots of the group, then these power sums are found to be symmetric functions of all \(\mu\), and hence rational functions of the coefficients of the given equation. One can therefore immediately set up an equation of degree \(\frac{\nu^{2p}-1}{\nu-1}\), whose roots are the values taken by the symmetric function \(A\) for the different groups. If \(B\) is another symmetric function of the roots of a group, then the sum of all \(B\), the sum of all \(BA\), the sum of all \(BA^2\), etc., can likewise be expressed rationally by the coefficients of the given equation. In this way, once the \(A\) are known, one obtains for every other symmetric function of the roots of a group a linear determination; and it is easy to see that in solving these equations of the first degree every \(B\) can be transformed so that it contains only the one \(A\) corresponding to it.

Thus, if the equation of degree
\[
\frac{\nu^{2p}-1}{\nu-1}
\]
on which the \(A\) depend has been solved, then by means of one root \(A\) of this equation one can represent the remaining symmetric functions of the roots of the corresponding group, and one can therefore at once write down an equation of degree \(\nu-1\), whose roots are the roots of this group and whose coefficients depend only on one root \(A\) of the first-mentioned equation.

The equation for the \(A\) is the only one in these problems which is not algebraically solvable. For the equations of degree \(\nu-1\) corresponding to the individual groups can again be solved according to Abel by means of root quantities. For this purpose let \(g\) be a primitive root of \(\nu\); then the period system of a group can be represented by
\[
P_h,\quad gP_h,\quad g^2P_h,\ldots,\quad g^{\nu-2}P_h.
\]
Let us call the corresponding roots of the group
\[
\mu_0,\mu_1,\mu_2,\ldots,\mu_{\nu-2}.
\]
Then, by Abel's theorem,
\[
\mu_1=\Theta(\mu_0),\quad
\mu_2=\Theta(\mu_1),\ldots,\quad
\mu_{\nu-2}=\Theta(\mu_{\nu-3});
\]
but since \(g^{\nu-1}\equiv 1\pmod{\nu}\), the period system \(g^{\nu-1}P_h\) corresponding to the root
\[
\mu_{\nu-1}=\Theta(\mu_{\nu-2})
\]
returns to \(P_h\) itself, and consequently \(\mu_{\nu-1}=\mu_0\), or
\[
\mu_0=\Theta(\mu_{\nu-2}).
\]
The roots \(\mu_0,\mu_1,\mu_2\ldots\mu_{\nu-2}\), thus ordered, therefore form a cycle, so that, if another root of the group is taken as the basis instead of \(\mu\), these roots are merely cyclically permuted. Now form the expression
\[
F_\lambda=\sum_{i=1}^{i=\nu-2}
e^{\frac{2\lambda i\pi\sqrt{-1}}{\nu-1}}\mu_i.
\]
Then \(F_\lambda\) can be represented as a rational function of any one of the \(\nu-1\) roots \(\mu\). If we put
\[
\Theta[\Theta(\mu)]=\Theta_1(\mu),\quad
\Theta[\Theta_1(\mu)]=\Theta_2(\mu)\quad\text{etc.},
\]
then generally, between two roots \(\mu_k,\mu_h\), one has the relation
\[
\mu_k=\Theta_{k-h-1}(\mu_h),
\]
where, if \(k-h-1\) becomes negative, this index is to be increased by \(\nu-1\). In order, therefore, to express \(F_\lambda\) rationally through \(\mu_h\), one puts
\[
F_\lambda=
\sum_{i=1}^{i=\nu-2}
e^{\frac{2\lambda i\pi\sqrt{-1}}{\nu-1}}\,
\Theta_{i-h-1}(\mu_h),
\]
or, putting \(i-h-1=k\),
\[
F_\lambda=
e^{\frac{2\lambda(h+1)\pi\sqrt{-1}}{\nu-1}}
\sum_{k}
e^{\frac{2\lambda k\pi\sqrt{-1}}{\nu-1}}\Theta_k(\mu_h).
\]
Since here the number \(h\), except in \(\mu_h\), occurs only in a \((\nu-1)^{\text{th}}\) root of unity appearing as a factor, \(F_\lambda^{\nu-1}\) is an expression which has exactly the same form no matter through which root \(\mu_h\) it is represented. Thus, if one forms the \(\nu-1\) expressions of \(F_\lambda^{\nu-1}\) through the different roots \(\mu\) and adds them all, one obtains \((\nu-1)F_\lambda^{\nu-1}\) equal to a symmetric function of these roots, and therefore \(F_\lambda^{\nu-1}\) expressed by the coefficients of our equation of degree \(\nu-1\). Let therefore
\[
F_\lambda^{\nu-1}=\Phi_\lambda.
\]
Then root extraction gives the \(\nu-1\) equations
\[
F_1=\sqrt[\nu-1]{\Phi_1},\qquad
F_2=\sqrt[\nu-1]{\Phi_2},\qquad
\ldots,\qquad
F_{\nu-2}=\sqrt[\nu-1]{\Phi_{\nu-2}},
\]
while \(F_0\) is from the outset a symmetric function of the roots and a coefficient of the equation of degree \(\nu-1\) itself. Thus one has \(\nu-1\) equations of the first degree, from which the \(\nu-1\) roots of the group are linearly determined.

Among the different \((\nu-1)^{\text{th}}\) roots which occur in the expressions for the \(F\), relations still exist. For one sees, just as the corresponding property of \(F_\lambda^{\nu-1}\) was proved above, that the product
\[
\frac{F_\lambda}{F_1^\lambda}
\]
can be expressed rationally by the coefficients of the equation of degree \(\nu-1\). Thus if
\[
\frac{F_\lambda}{F_1^\lambda}=U_\lambda,
\]
where \(U_\lambda\) now denotes a rational function of the given coefficients, one also has
\[
F_\lambda=U_\lambda\left(\sqrt[\nu-1]{\Phi_1}\right)^\lambda.
\]
In the expressions of the roots \(\mu\), therefore, only powers of a single root quantity \(\sqrt[\nu-1]{\Phi_1}\) still occur, and one obtains all roots \(\mu\) from one of them by assigning to this root quantity its different values.

\subsection*{\S\,73. Algebraically representable \(\Theta\)-quotients.}

The consideration of the simple algebraically representable \(\Theta\)-quotients stands in exact connection with the preceding investigations.



% --- Continuation: source PDF pp. 271--278, printed pp. 249--256. ---
% English translation. Formula authority: rendered Google source pages at 300dpi.
% Continuation of §73 from previous tranche; source page 278 also begins §74.

Let \(\frac{\Theta(w+A)}{\Theta(w)}\) be a \(\Theta\)-quotient which can be represented as an algebraic function of the coordinates of the points \(x^{(1)},x^{(2)}\ldots x^{(p)},\xi\), which are connected with the \(w\) by the equations
\[
w_h=\sum_{i=1}^{i=p}\int_{\alpha^{(i)}}^{x^{(i)}}du_h-\int_{\mu}^{\xi}du_h
\qquad(h=1,2\ldots p).
\]
Then, when the paths of integration are changed, that is, when the arguments are increased by periodicity moduli, that quotient must be capable of assuming only a finite number of different values. But if one puts
\[
w_h'=w_h+2m_h\pi\sqrt{-1}+a_{1h}q_1+a_{2h}q_2\cdots+a_{ph}q_p
\]
in place of \(w_h\) (cf. above p.~200), then
\[
\Theta(w'+A)=e^{(w_1+A_1)q_1+(w_2+A_2)q_2\cdots}\Theta(w+A),
\]
\[
\Theta(w')=e^{w_1q_1+w_2q_2\cdots}\Theta(w),
\]
and hence
\[
(1)\ldots\quad
\frac{\Theta(w'+A)}{\Theta(w')}
=
e^{A_1q_1+A_2q_2\cdots+A_pq_p}\cdot
\frac{\Theta(w+A)}{\Theta(w)}.
\]
Therefore, if all \(q\)'s except one remain unchanged, the exponential factor must again assume its original value as soon as this one number \(q_i\) has increased by a certain integer \(r_i\). Thus one must have
\[
e^{A_i(q_i+r_i)}=e^{A_iq_i},\qquad
A_i=\frac{2\pi s_i\sqrt{-1}}{r_i},
\]
where \(s_i\) is again an integer.

Suppose one requires only that \(\frac{\Theta(w+A)}{\Theta(w)}\), multiplied by an exponential factor
\[
e^{B_1w_1+B_2w_2\cdots+B_pw_p},
\]
have a finite series of values when the arguments are increased by periodicity moduli. Then, in place of (1), there enters the equation
\[
\frac{\Theta(w'+A)}{\Theta(w')}\cdot e^{\sum B_iw_i'}=
e^{\sum A_iq_i+2\pi\sqrt{-1}\sum m_iB_i+\sum\sum B_iq_ka_{ik}}
\cdot
\frac{\Theta(w+A)}{\Theta(w)}\cdot e^{\sum B_iw_i},
\]
and therefore the factor
\[
e^{\sum A_iq_i+2\pi\sqrt{-1}\sum m_iB_i+\sum\sum B_iq_ka_{ik}}
\]
must have a finite series of values. If it remains unchanged when \(q_i\) is increased by \(r_i\), or \(m_i\) by \(\varrho_i\), it follows that
\[
e^{2\pi\sqrt{-1}B_i\varrho_i}=1,\qquad
e^{(A_i+\sum_kB_ka_{ik})r_i}=1,
\]
and hence
\[
B_i=-\frac{\sigma_i}{\varrho_i},\qquad
A_i=\frac{2s_i\pi\sqrt{-1}}{r_i}+\frac{1}{\varrho_i}\sum_k a_{ik}\sigma_k,
\]
where again the \(\sigma\)'s and \(s\)'s are integers. These expressions can be brought into a form in which, instead of the \(r_i,\varrho_i\), the greatest common divisor \(m\) of these numbers occurs, and one may therefore put
\[
B_i=-\frac{\sigma_i}{m},\qquad
A_i=\frac1m\left[2s_i\pi\sqrt{-1}+\sum_k a_{ik}\sigma_k\right].
\]
Thus the system of the \(A\)'s is the \(m^{\text{th}}\) part of a general system of periodicity moduli, and one therefore has the theorem (cf. Riemann, l.~c. p.~154):

\begin{quote}
In order that
\[
\frac{\Theta(w+A)}{\Theta w}\,e^{\sum w_iB_i}
\]
be algebraically expressible by the \(\xi,x\) which are determined from the equations
\[
w_h=\sum_{i=1}^{i=p}\int_{c^{(i)}}^{x^{(i)}}du_h-\int_{\mu}^{\xi}du_h,
\]
it is necessary that the \(A\)'s be \(r^{\text{th}}\) parts of a system of periodicity moduli, and if
\[
A_i=\frac1m\left[2s_i\pi\sqrt{-1}+\sum_k a_{ik}\sigma_k\right],
\]
then
\[
B_i=-\frac{\sigma_i}{m}.
\]
\end{quote}

The extended theorem of Hermite now teaches how to carry out this representation as follows.

Put
\[
m_h^{(1)}=m_h^{(2)}\cdots=m_h^{(m-1)}=\frac{P_h}{m},
\]
where \(P_h\) denotes the periodicity modulus
\[
P=2s_h\pi\sqrt{-1}+\sum_k a_{hk}\sigma_k.
\]
Then, in order that the theorem of \S\,65 may be applied, one must have
\[
m_h^{(m)}=-P_h+\frac{P_h}{m},
\]
and hence
\[
\Theta(w+m^{(m)})=
e^{-\sum_h\left(w_h+\frac1mP_h\right)\sigma_h}
\cdot\Theta\left(w+\frac Pm\right).
\]
The algebraically representable quotient is therefore
\[
\left[\frac{\Theta\left(w+\frac Pm\right)}{\Theta(w)}\right]^m
\cdot
e^{-\sum_h\left(w_h+\frac{P_h}{m}\right)\sigma_h}.
\]
If we transform analogously the constant \(\Theta\)-quotient on the right-hand side of the theorem, we obtain for it the expression
\[
\left[\frac{\Theta\left(w^{(p)}+\frac Pm\right)}{\Theta(w^{(p)})}\right]^m
\cdot
e^{-\sum_h\left(w_h^{(p)}+\frac{P_h}{m}\right)\sigma_h},
\]
where the quantities
\[
w_h^{(p)}=\sum_{i=1}^{i=p}\int_{\alpha^{(i)}}^{a^{(i)}}du_h-\int_{\mu}^{\eta}du_h
\]
can assume arbitrary constant values because the points \(a\) may be chosen arbitrarily on \(f=0\). Thus one can finally say that the expression
\[
F=
\frac{\Theta\left(w+\frac Pm\right)}{\Theta(w)}
\cdot
\frac{\Theta(w^{(p)})}{\Theta\left(w^{(p)}+\frac Pm\right)}
\cdot
e^{\sum_h\frac{w_h^{(p)}-w_h}{m}\sigma_h}
\]
is the \(m^{\text{th}}\) root of a rational function of the coordinates of \(x^{(1)},x^{(2)}\ldots x^{(p)},\xi\) (and of the \(a\)'s), as Riemann l.~c. has shown. A further consideration, however, leads to a remarkable transformation of this algebraic function.

According to the representation of \S\,65,
\[
e^{-\sum_h\frac{w_h\sigma_h}{m}}
\frac{\Theta\left(w+\frac Pm\right)}{\Theta(w)}
=
e^{-\sum_h\frac{w_h^{(p)}\sigma_h}{m}}
\frac{\Theta\left(w^{(p)}+\frac Pm\right)}{\Theta(w^{(p)})}
\sqrt[m]{
\frac{
\dfrac{\varphi_0(\xi)\varphi_1(x^{(1)})\cdots\varphi_p(x^{(p)})}
      {\psi_0(\xi)\psi_1(x^{(1)})\cdots\psi_p(x^{(p)})}
}{
\dfrac{\varphi_0(\eta)\varphi_1(a^{(1)})\cdots\varphi_p(a^{(p)})}
      {\psi_0(\eta)\psi_1(a^{(1)})\cdots\psi_p(a^{(p)})}
}
}.
\]
This formula is composed from the individual equations
\[
(1)\ldots
\begin{cases}
e^{-\sum_h\frac{w_h\sigma_h}{m}}
\dfrac{\Theta\left(w+\frac Pm\right)}{\Theta(w)}
=
e^{-\sum_h\frac{w_h^\circ\sigma_h^\circ}{m}}
\dfrac{\Theta\left(w^\circ+\frac Pm\right)}{\Theta(w^\circ)}
\sqrt[m]{\dfrac{\dfrac{\varphi_0(\xi)}{\psi_0(\xi)}}{\dfrac{\varphi_0(\eta)}{\psi_0(\eta)}}},\\[1.0em]
e^{-\sum_h\frac{w_h^\circ\sigma_h^\circ}{m}}
\dfrac{\Theta\left(w^\circ+\frac Pm\right)}{\Theta(w)}
=
e^{-\sum_h\frac{w_h^{(1)}\sigma_h^{(1)}}{m}}
\dfrac{\Theta\left(w^{(1)}+\frac Pm\right)}{\Theta(w^{(1)})}
\sqrt[m]{\dfrac{\dfrac{\varphi_1(x^{(1)})}{\psi_1(x^{(1)})}}{\dfrac{\varphi_1(a^{(1)})}{\psi_1(a^{(1)})}}},\\
\cdots.
\end{cases}
\]

For the definition of the different \(w\)'s see \S\,65, where it is also stated on which of the points \(x,\xi\) the coefficients of the different functions \(\varphi,\psi\) depend, namely:
\[
\begin{array}{ll}
\varphi_0,\psi_0 & \text{on }x^{(1)},x^{(2)}\ldots x^{(p)},\\
\varphi_1,\psi_1 & \text{on }a^{(1)},x^{(2)}\ldots x^{(p)},\\
\varphi_2,\psi_2 & \text{on }a^{(1)},a^{(2)}\ldots x^{(p)},\\
\multicolumn{2}{c}{\cdots}\\
\varphi_p,\psi_p & \text{on }a^{(1)},a^{(2)}\ldots a^{(p)},
\end{array}
\]
while at the same time
\[
(2)\ldots
\begin{cases}
w_h^\circ=\displaystyle\int_{\alpha^{(1)}}^{x^{(1)}}du_h+\int_{\alpha^{(2)}}^{x^{(2)}}du_h\cdots+\int_{\alpha^{(p)}}^{x^{(p)}}du_h-\int_{\mu}^{\eta}du_h,\\[0.55em]
w_h^{(1)}=\displaystyle\int_{\alpha^{(1)}}^{a^{(1)}}du_h+\int_{\alpha^{(2)}}^{x^{(2)}}du_h\cdots+\int_{\alpha^{(p)}}^{x^{(p)}}du_h-\int_{\mu}^{\eta}du_h,\\[0.55em]
w_h^{(2)}=\displaystyle\int_{\alpha^{(1)}}^{a^{(1)}}du_h+\int_{\alpha^{(2)}}^{a^{(2)}}du_h\cdots+\int_{\alpha^{(p)}}^{x^{(p)}}du_h-\int_{\mu}^{\eta}du_h,\\
\cdots,\\
w_h^{(p)}=\displaystyle\int_{\alpha^{(1)}}^{a^{(1)}}du_h+\int_{\alpha^{(2)}}^{a^{(2)}}du_h\cdots+\int_{\alpha^{(p)}}^{a^{(p)}}du_h-\int_{\mu}^{\eta}du_h.
\end{cases}
\]
If now in the equations (1), in the first one, we replace \(x^{(1)},x^{(2)}\ldots x^{(p)}\) by \(a^{(1)},a^{(2)}\ldots a^{(p)}\), in the second \(x^{(2)},x^{(3)}\ldots x^{(p)}\) by \(a^{(2)},a^{(3)}\ldots a^{(p)}\), in the third \(x^{(3)},x^{(4)}\ldots x^{(p)}\) by \(a^{(3)},a^{(4)}\ldots a^{(p)}\), and so on, then the \(w_h^{(i)}\) contained respectively in the arguments of the \(\Theta\)-functions on the left of equation (1) pass into
\[
\begin{aligned}
W_h&=\int_{\alpha^{(1)}}^{a^{(1)}}du_h+\int_{\alpha^{(2)}}^{a^{(2)}}du_h\cdots+\int_{\alpha^{(p)}}^{a^{(p)}}du_h-\int_{\mu}^{\xi}du_h
=w_h^{(p)}-\int_{\eta}^{\xi}du_h,\\
W_h^\circ&=\int_{\alpha^{(1)}}^{x^{(1)}}du_h+\int_{\alpha^{(2)}}^{a^{(2)}}du_h\cdots+\int_{\alpha^{(p)}}^{a^{(p)}}du_h-\int_{\mu}^{\eta}du_h
=w_h^{(p)}+\int_{a^{(1)}}^{x^{(1)}}du_h,\\
W_h^{(1)}&=\int_{\alpha^{(1)}}^{a^{(1)}}du_h+\int_{\alpha^{(2)}}^{x^{(2)}}du_h\cdots+\int_{\alpha^{(p)}}^{a^{(p)}}du_h-\int_{\mu}^{\eta}du_h
=w_h^{(p)}+\int_{a^{(2)}}^{x^{(2)}}du_h,\\
&\hspace{11em}\cdots;
\end{aligned}
\]
at the same time \(\varphi_0,\psi_0\) pass into two expressions \(\Phi,\Psi\), whose coefficients now depend only on constant points, and \(\varphi_1,\varphi_2\ldots\varphi_p\), which are all formed in the same way, pass into \(\Phi_1\), and \(\psi_1,\psi_2\ldots\psi_p\) into \(\Psi_1\), whose coefficients likewise depend only on constant points. The equations (1) therefore transform into the following:
\[
\begin{aligned}
e^{-\sum_h\frac{W_h\sigma_h}{m}}
\frac{\Theta\left(W+\frac Pm\right)}{\Theta(W)}
&=
e^{-\sum_h\frac{w_h^{(p)}\sigma_h}{m}}
\frac{\Theta\left(w^{(p)}+\frac Pm\right)}{\Theta(w^{(p)})}
\sqrt[m]{\dfrac{\dfrac{\Phi(\xi)}{\Psi(\xi)}}{\dfrac{\Phi(\eta)}{\Psi(\eta)}}},\\
e^{-\sum_h\frac{W_h^\circ\sigma_h}{m}}
\frac{\Theta\left(W^\circ+\frac Pm\right)}{\Theta(W^\circ)}
&=
e^{-\sum_h\frac{w_h^{(p)}\sigma_h}{m}}
\frac{\Theta\left(w^{(p)}+\frac Pm\right)}{\Theta(w^{(p)})}
\sqrt[m]{\dfrac{\dfrac{\Phi_1(x^{(1)})}{\Psi_1(x^{(1)})}}{\dfrac{\Phi_1(a^{(1)})}{\Psi_1(a^{(1)})}}},\\
e^{-\sum_h\frac{W_h^{(1)}\sigma_h}{m}}
\frac{\Theta\left(W^{(1)}+\frac Pm\right)}{\Theta(W^{(1)})}
&=
e^{-\sum_h\frac{w_h^{(p)}\sigma_h}{m}}
\frac{\Theta\left(w^{(p)}+\frac Pm\right)}{\Theta(w^{(p)})}
\sqrt[m]{\dfrac{\dfrac{\Phi_1(x^{(2)})}{\Psi_1(x^{(2)})}}{\dfrac{\Phi_1(a^{(2)})}{\Psi_1(a^{(2)})}}},\\
&\hspace{9em}\cdots,
\end{aligned}
\]
in which the algebraic functions still depend on only one variable point each.

If now we divide each of these equations by the corresponding equation (1), we find
\[
(3)\ldots
\begin{aligned}
\frac{\Theta\left(w+\frac Pm\right)\Theta(W)\Theta(w^\circ)\Theta\left(w^{(p)}+\frac Pm\right)}
{\Theta(w)\Theta\left(W+\frac Pm\right)\Theta\left(w^\circ+\frac Pm\right)\Theta(w^{(p)})}
&=
\sqrt[m]{\dfrac{\dfrac{\varphi_0(\xi)}{\psi_0(\xi)}}{\dfrac{\varphi_0(\eta)}{\psi_0(\eta)}}}
\cdot
\sqrt[m]{\dfrac{\dfrac{\Phi(\eta)}{\Psi(\eta)}}{\dfrac{\Phi(\xi)}{\Psi(\xi)}}},\\[0.7em]
\frac{\Theta\left(w^\circ+\frac Pm\right)\Theta(W^\circ)\Theta(w^{(1)})\Theta\left(w^{(p)}+\frac Pm\right)}
{\Theta(w^\circ)\Theta\left(W^\circ+\frac Pm\right)\Theta\left(w^{(1)}+\frac Pm\right)\Theta(w^{(p)})}
&=
\sqrt[m]{\dfrac{\dfrac{\varphi_1(x^{(1)})}{\psi_1(x^{(1)})}}{\dfrac{\varphi_1(a^{(1)})}{\psi_1(a^{(1)})}}}
\cdot
\sqrt[m]{\dfrac{\dfrac{\Phi_1(a^{(1)})}{\Psi_1(a^{(1)})}}{\dfrac{\Phi_1(x^{(1)})}{\Psi_1(x^{(1)})}}},\\[0.7em]
\frac{\Theta\left(w^{(1)}+\frac Pm\right)\Theta(W^{(1)})\Theta(w^{(2)})\Theta\left(w^{(p)}+\frac Pm\right)}
{\Theta(w^{(1)})\Theta\left(W^{(1)}+\frac Pm\right)\Theta\left(w^{(2)}+\frac Pm\right)\Theta(w^{(p)})}
&=
\sqrt[m]{\dfrac{\dfrac{\varphi_2(x^{(1)})}{\psi_2(x^{(1)})}}{\dfrac{\varphi_2(a^{(1)})}{\psi_2(a^{(1)})}}}
\cdot
\sqrt[m]{\dfrac{\dfrac{\Phi_1(a^{(2)})}{\Psi_1(a^{(2)})}}{\dfrac{\Phi_1(x^{(2)})}{\Psi_1(x^{(2)})}}},\\
&\hspace{9em}\cdots .
\end{aligned}
\]
But the expressions on the left can be represented directly as rational functions of the coordinates of the \(x,a,\xi,\eta\), by going back to the theorem of \S\,62. To the \(\Theta\)-quotient of the first equation that theorem is immediately applicable in the form given there. The \(\Theta\)-quotient becomes equal to
\[
\frac{\left(\frac{U}{\overline V}\right)_\xi}{\left(\frac{U}{\overline V}\right)_\eta},
\]
if the curve \(V=0\) passes through the points \(x\) and through the points determined by the equations
\[
\sum_{i=1}^{i=p}\int_{\alpha^{(i)}}^{z^{(i)}}du_h
=
\sum_{i=1}^{i=p}\int_{\alpha^{(i)}}^{a^{(i)}}du_h+\frac{P_h}{m},
\]
or
\[
(4)\ldots\quad
\sum_{i=1}^{i=p}\int_{a^{(i)}}^{z^{(i)}}du_h=\frac{P_h}{m},
\]
and the curve \(U=0\) passes through the points determined by the equations
\[
\sum_{i=1}^{i=p}\int_{\alpha^{(i)}}^{y^{(i)}}du_h
=
\sum_{i=1}^{i=p}\int_{\alpha^{(i)}}^{x^{(i)}}du_h+\frac{P_h}{m},
\]
or
\[
(5)\ldots\quad
\sum_{i=1}^{i=p}\int_{x^{(i)}}^{y^{(i)}}du_h=\frac{P_h}{m},
\]
and through the points \(a\), while all remaining intersection points of \(U=0\) and \(V=0\) with \(f=0\), including the double and cusp points, are common.

Observe that equation (5) introduces no irrationality with respect to the \(x\)'s. For although at first the \(y\)'s belong to a special division problem in which the \(x\)'s are used as the basis, one can nevertheless, by Abel's theorem, if one puts for example the expressions (4) in place of \(\frac Pm\), reduce this irrationality to one taken from a special division problem with constant fundamental points.

The other \(\Theta\)-quotients must first be brought to an analogous form, by using the methods employed at the end of \S\,63, so that, by application of Abel's theorem, the arguments of the \(\Theta\)-functions are transformed into such arguments as differ, in numerator and denominator, only by the negatively taken integral. In this way the left-hand sides of the remaining equations (3) pass into rational functions
\[
\frac{\left(\frac{U_1}{V_1}\right)_{x^{(1)}}}{\left(\frac{U_1}{V_1}\right)_{a^{(1)}}},\quad
\frac{\left(\frac{U_2}{V_2}\right)_{x^{(2)}}}{\left(\frac{U_2}{V_2}\right)_{a^{(2)}}},\ldots
\]
Substituting these into equations (3), one obtains
\[
\begin{aligned}
\sqrt[m]{\dfrac{\dfrac{\varphi_0(\xi)}{\psi_0(\xi)}}{\dfrac{\varphi_0(\eta)}{\psi_0(\eta)}}}
&=
\frac{\left(\frac{U}{V}\right)_\xi}{\left(\frac{U}{V}\right)_\eta}
\cdot
\sqrt[m]{\dfrac{\dfrac{\Phi(\xi)}{\Psi(\xi)}}{\dfrac{\Phi(\eta)}{\Psi(\eta)}}},\\
\sqrt[m]{\dfrac{\dfrac{\varphi_1(x^{(1)})}{\psi_1(x^{(1)})}}{\dfrac{\varphi_1(a^{(1)})}{\psi_1(a^{(1)})}}}
&=
\frac{\left(\frac{U_1}{V_1}\right)_{x^{(1)}}}{\left(\frac{U_1}{V_1}\right)_{a^{(1)}}}
\cdot
\sqrt[m]{\dfrac{\dfrac{\Phi_1(x^{(1)})}{\Psi_1(x^{(1)})}}{\dfrac{\Phi_1(a^{(1)})}{\Psi_1(a^{(1)})}}},\\
\sqrt[m]{\dfrac{\dfrac{\varphi_2(x^{(2)})}{\psi_2(x^{(2)})}}{\dfrac{\varphi_2(a^{(2)})}{\psi_2(a^{(2)})}}}
&=
\frac{\left(\frac{U_2}{V_2}\right)_{x^{(2)}}}{\left(\frac{U_2}{V_2}\right)_{a^{(2)}}}
\cdot
\sqrt[m]{\dfrac{\dfrac{\Phi_1(x^{(2)})}{\Psi_1(x^{(2)})}}{\dfrac{\Phi_1(a^{(2)})}{\Psi_1(a^{(2)})}}},\\
&\hspace{9em}\cdots
\end{aligned}
\]
and thus for our \(\Theta\)-quotient one obtains the following expression:
\[
\begin{aligned}
(6)\ldots\quad
e^{-\sum_h\frac{w_h\sigma_h}{m}}
\frac{\Theta\left(w+\frac Pm\right)}{\Theta(w)}
&=
e^{-\sum_h\frac{w_h^{(p)}\sigma_h}{m}}
\frac{\Theta\left(w^{(p)}+\frac Pm\right)}{\Theta(w^{(p)})}\\
&\quad{}\cdot
\frac{\left(\frac UV\right)_\xi\left(\frac{U_1}{V_1}\right)_{x^{(1)}}\cdots}
     {\left(\frac UV\right)_\eta\left(\frac{U_1}{V_1}\right)_{a^{(1)}}\cdots}
\sqrt[m]{\frac{
\dfrac{\Phi(\xi)}{\Psi(\xi)}\dfrac{\Phi_1(x^{(1)})}{\Psi_1(x^{(1)})}\cdots
}{
\dfrac{\Phi(\eta)}{\Psi(\eta)}\dfrac{\Phi_1(a^{(1)})}{\Psi_1(a^{(1)})}\cdots
}}.
\end{aligned}
\]
The required \(\Theta\)-quotient therefore finally expresses itself as a rational function of the coordinates of \(\xi,x^{(1)},x^{(2)}\ldots\), multiplied by the \(m^{\text{th}}\) root of a product of rational functions, each of which contains only one of the variable points.

This formula contains, besides the \(m^{\text{th}}\) root, no other irrationalities than a root \(\mu\) of the equation \(M=0\) for the special division problem.

That the quotient
\[
e^{-\sum_h\frac{(w_h-w_h^{(p)})\sigma_h}{m}}
\frac{\dfrac{\Theta\left(w+\frac Pm\right)}{\Theta(w)}}
     {\dfrac{\Theta\left(w^{(p)}+\frac Pm\right)}{\Theta(w^{(p)})}}
\]
is expressed by the \(m^{\text{th}}\) root of an algebraic function has its reason in the fact that that quotient changes by an \(m^{\text{th}}\) root of unity when the numbers \(\varrho,\sigma\) are changed by integral multiples of \(m\), while the algebraic function remains wholly unchanged.

\subsection*{\S\,74. Algebraic determination of \(\dfrac{\Theta\left(\frac Pm\right)}{\Theta(0)}\).}

The transcendental term contained in equation (6),
\[
e^{-\sum_h\frac{w_h^{(p)}\sigma_h}{m}}
\frac{\Theta\left(w^{(p)}+\frac Pm\right)}{\Theta(w^{(p)})},
\]
can itself be reduced to algebraic expressions, and thus the given \(\Theta\)-quotient can be expressed purely algebraically,


% --- Source pages 279--291; printed pages 257--269. ---

by choosing the previously arbitrary constant points \(a\) in a suitable manner. We had
\[
w_h^{(p)}=\sum_{i=1}^{i=p}\int_{\alpha^{(i)}}^{a^{(i)}}du_h-\int_{\mu}^{\eta}du_h.
\]
For simplicity let the point \(\eta\) fall into the point \(\mu\), and then put
\[
\sum_{i=1}^{i=p}\int_{\alpha^{(i)}}^{a^{(i)}}du_h=\frac{Q_h}{m},
\]
where the \(Q_h\) denote a system of periodicity moduli. Thus in this case the \(a\) are a solution of a special division problem. In particular one may set all \(Q_h\) equal to zero, and let the \(a\) coincide with the \(\alpha\). Then \(w_h^{p}=0\), and the transcendental factor is reduced to
\[
\frac{\Theta\left(\frac Pm\right)}{\Theta(o)}.
\]
The question now is to express these quantities algebraically, which is done as follows.

In formula (6) set all \(w_h\) equal to zero; let \(\xi,\eta\) coincide with the fixed point \(\mu\), and hence the \(x\) with the \(\alpha\). The \(w_h^{(p)}\), on the other hand, we determine by the equations
\[
w_h^{(p)}=\sum_{i=1}^{i=p}\int_{\alpha^{(i)}}^{a^{(i)}}du_h
=\frac{m-1}{m}P_h.
\]
The symmetric functions of the coordinates of the \(a\) then become rational functions of the root of the equation \(M=0\) belonging to \(\frac{m-1}{m}P_h\). But this root stands, with the root belonging to \(\frac{P_h}{m}\) itself, which has already been used, in the special relation that each is the same rational function of the other, as the theorem of §64 or Abel's theorem teaches. Thus one may also say that the symmetric functions of the coordinates of the \(a\) are rational functions of the root of \(M=0\) belonging to \(\frac{P_h}{m}\).

By means of the indicated values the two quotients
\[
e^{-\sum\frac{w_h\sigma_h}{m}}
\frac{\Theta\left(w+\frac Pm\right)}{\Theta(w)},\qquad
e^{-\sum\frac{w_h^{(p)}\sigma_h}{m}}
\frac{\Theta\left(w^{(p)}+\frac Pm\right)}{\Theta(w^{(p)})}
\]
pass into
\[
\frac{\Theta\left(\frac Pm\right)}{\Theta(o)},\qquad
e^{-\frac{m-1}{m^2}\sum P_h\sigma_h}\cdot
\frac{\Theta(P)}{\Theta\left(\frac{m-1}{m}P\right)}.
\]
Of these expressions the latter, by use of the periodicity equation valid for the \(\Theta\)-function (p. 200), also assumes the form
\[
e^{\sum\frac{P_h\sigma_h}{m^2}}\cdot
\frac{\Theta(o)}{\Theta\left(\frac Pm\right)}.
\]
Equation (6) can therefore be written in the form
\[
\begin{aligned}
e^{-\sum\frac{P_h\sigma_h}{m^2}}
\left[\frac{\Theta\left(\frac Pm\right)}{\Theta(o)}\right]^2
&=
\frac{\left(\frac{U_1}{V_1}\right)_{\alpha^{(1)}}
      \left(\frac{U_2}{V_2}\right)_{\alpha^{(2)}}\cdots}
     {\left(\frac{U_1}{V_1}\right)_{a^{(1)}}
      \left(\frac{U_2}{V_2}\right)_{a^{(2)}}\cdots} \\
&\quad{}\cdot
\sqrt[m]{\frac{
\dfrac{\Phi_1(\alpha^{(1)})\Phi_1(\alpha^{(2)})\cdots}
      {\Psi_1(\alpha^{(1)})\Psi_1(\alpha^{(2)})\cdots}
}{
\dfrac{\Phi_1(a^{(1)})\Phi_1(a^{(2)})\cdots}
      {\Psi_1(a^{(1)})\Psi_1(a^{(2)})\cdots}
}}.
\end{aligned}
\]
The expression on the right is algebraic. The expression under the radical is already a symmetric function of the coordinates of the several points \(\alpha\) on the one hand and of the several points \(a\) on the other, hence a rational function of the root of the special division equation \(M=0\) belonging to \(\frac Pm\). The rational part before the radical is formed unsymmetrically, but it does not change under permutation of the \(\alpha\) among themselves or of the \(a\) among themselves, and can therefore also be brought into symmetric form. Thus, if \(\mu\) is the corresponding root of the division equation, one has the theorem:

The expressions
\[
e^{-\sum\frac{P_h\sigma_h}{2m^2}}
\frac{\Theta\left(\frac Pm\right)}{\Theta(o)}
\]
can be represented in the form
\[
\sqrt{\,A(\mu)\sqrt[m]{B(\mu)}\,},
\]
where \(A,B\) are rational functions, and \(\mu\) is the corresponding root of the equation \(M=0\) of the special division problem.

That this expression is represented by a \((2m)^{\text{th}}\) root is explained by the fact that it changes by a \((2m)^{\text{th}}\) root of unity when the numbers \(s,\sigma\) are changed by whole multiples of \(m\), while the algebraic expressions \(A,B\) remain entirely unchanged.

To the periodicity moduli \(P_h\) and \(-P_h\) there correspond two different roots of \(M=0\), \(\mu,\mu'\), which, when inserted in the above formula, plainly yield the same result. If the above equation is raised to the power \(2m\), so that
\[
A^2(\mu)\cdot B(\mu)
=
e^{-\sum\frac{P_h\sigma_h}{m}}
\left[\frac{\Theta\left(\frac Pm\right)}{\Theta(o)}\right]^{2m},
\]
then, with the aid of this relation, the equation \(M=0\) can be transformed into an equation \(N=0\), whose roots are the quantities
\[
e^{-\sum\frac{P_h\sigma_h}{m}}
\left[\frac{\Theta\left(\frac Pm\right)}{\Theta(o)}\right]^{2m}.
\]
Once this equation \(N=0\) has been solved, in general two associated roots \(\mu,\mu'\) of \(M=0\) are found from each of its roots, and their symmetric functions are rationally expressible, during the formation of the new equation, by means of the new unknown. According to the preceding, only the cases have to be distinguished and treated in which \(m=2\), and in which \(m\) is an odd prime number. In the first case the new equation is of degree
\[
\frac{m^{2p}-1}{2},
\]
or each of the algebraically solvable equations which remain after solution of the auxiliary equation of degree
\[
\frac{m^{2p}-1}{m-1}
\]
is transformed into an equation of degree
\[
\frac{m-1}{2}.
\]
The other case will be discussed more precisely below.

\subsection*{\S\,75. Determination of \(\dfrac{\Theta\left(\frac P2\right)}{\Theta(o)}\). Curves of order \((n-3)\) which touch in \(p\) points and pass through the double and return points.}

The case \(m=2\) is distinguished by the fact that one series among the expressions \(\Theta^{P/2}\) vanishes identically. For, by the periodicity property of the function \(\Theta\) (p. 200),
\[
\Theta\left(\frac P2\right)
=
\Theta\left(-\frac P2\right)\cdot
e^{-\sum\frac{P_h\sigma_h}{2}+\frac12\sum\sum\sigma_h\sigma_k a_{hk}}.
\]
Now since
\[
P_h=2s_h\pi\sqrt{-1}+\sum_k a_{hk}\sigma_k,
\]
this equation becomes
\[
\Theta\left(\frac P2\right)=
\Theta\left(-\frac P2\right)\cdot(-1)^{\sum s_h\sigma_h}.
\]
But since, moreover,
\[
\Theta\left(\frac P2\right)=\Theta\left(-\frac P2\right),
\]
\(\Theta\left(\frac P2\right)\) always vanishes as soon as \(\sum s_h\sigma_h\) is odd, and one has the theorem:

The function \(\Theta\) always vanishes when its arguments are one of the systems of half-periods
\[
\frac12P_h=s_h\pi\sqrt{-1}+\frac12\sum_k a_{hk}\sigma_k
\]
for which the sum \(s_1\sigma_1+s_2\sigma_2+\cdots+s_p\sigma_p\) is an odd number.

Among the functions \(\Theta\left(\frac P2\right)\), therefore, just as many always vanish as the equation
\[
s_1\sigma_1+s_2\sigma_2+\cdots+s_p\sigma_p=2H+1
\]
has solutions, when the numbers \(s,\sigma\) are to be assigned only the values \(0\) or \(1\). To find the number of these solutions, suppose that \(k\) of the numbers \(s\) are equal to \(0\), the remaining \(p-k\) equal to \(1\); the \(\sigma\)'s corresponding to the first numbers \(s\) can then be \(0\) or \(1\), which gives \(2^k\) combinations. The sum of the other \(\sigma\)'s, however, must be odd; hence the last is always determined by the remaining \(p-k-1\). These last numbers \(\sigma\) can therefore still be chosen in \(2^{p-k-1}\) different ways, or in this case there are altogether \(2^{p-1}\) systems of all the \(\sigma\). According as one selects from the numbers \(s_1,s_2\ldots s_p\) such \(k\) which are to be equal to zero, one has
\[
\frac{p\cdot(p-1)\cdots(p-k+1)}{1\cdot2\cdots k}
\]
different cases, and therefore altogether
\[
\frac{p\cdot(p-1)\cdots(p-k+1)}{1\cdot2\cdots k}\cdot2^{p-1}
\]
systems of \(s,\sigma\) for a given \(k\). Finally \(k\) can assume the values \(0,1,\ldots,p-1\); hence the number of all possible solutions is
\[
2^{p-1}\left\{1+\frac p1+\frac{p(p-1)}{1\cdot2}\cdots+p\right\}
=2^{p-1}(2^p-1).
\]
This is therefore the number of systems of half-periodicity moduli for which the \(\Theta\)-function vanishes identically.

If now in the general equations
\[
w_h=\sum_{i=1}^{i=p}\int_{\alpha^{(i)}}^{x^{(i)}}du_h-\int_{\mu}^{\xi}du_h
\]
one again lets the point \(\xi\) coincide with \(\mu\), and replaces the \(w\) by a system of half-periodicity moduli, so that
\[
\frac12P_h=\sum_{i=1}^{i=p}\int_{\alpha^{(i)}}^{x^{(i)}}du_h,
\]
then the function \(\Theta\left(\frac P2\right)\) is zero only and always when either a point \(x^{(i)}\) coincides with \(\mu\), or when the points \(x\) lie on a curve of order \((n-3)\) which passes through the double and return points. Recall the origin of the \(\alpha\): they were the points of contact of a curve of order \((n-2)\) which passes through the double and return points of \(f=0\), as well as through the \((n-2)\) points in which the tangent to \(f=0\) at \(\mu\) further cuts this curve, and which, finally, touches the curve \(f=0\) in \(p\) points. Through those fixed points there pass other contact curves of the same kind; one of them has the \(x\) as its points of contact. If now one of the points \(x\) falls into \(\mu\), then its contact curve decomposes into a contact curve of order \((n-3)\), which passes through the double and return points and touches at the remaining \(p-1\) points \(x\), and into the tangent of the point \(\mu\) itself.

Let the points of contact of such a curve of order \((n-3)\) be \(\beta^{(1)},\beta^{(2)}\ldots\beta^{(p-1)}\). Then there is a certain system \(Q_h\) of periodicity moduli for which
\[
(1)\ldots\quad
\frac{Q_h}{2}
=\int_{\alpha^{(1)}}^{\beta^{(1)}}du_h
+\int_{\alpha^{(2)}}^{\beta^{(2)}}du_h
+\cdots+
\int_{\alpha^{(p-1)}}^{\beta^{(p-1)}}du_h
+\int_{\alpha^{(p)}}^{\mu}du_h.
\]
Let \(R_h\) be such a period system for which the equations
\[
(2)\ldots\quad
\frac{R_h}{2}
=\int_{\alpha^{(1)}}^{\gamma^{(1)}}du_h
+\int_{\alpha^{(2)}}^{\gamma^{(2)}}du_h
+\cdots+
\int_{\alpha^{(p)}}^{\gamma^{(p)}}du_h
\]
hold, where the \(\gamma\)'s lie on a curve of order \((n-3)\) passing through the double and return points. This curve then further cuts in, say, \(\delta^{(1)},\delta^{(2)},\ldots,\delta^{(p-2)}\), and by Abel's theorem one has
\[
\begin{aligned}
(3)\ldots\quad 0
&=\int_{\beta^{(1)}}^{\gamma^{(1)}}du_h
+\int_{\beta^{(2)}}^{\gamma^{(2)}}du_h
+\cdots+
\int_{\beta^{(p-1)}}^{\gamma^{(p-1)}}du_h\\
&\quad{}+\int_{\beta^{(1)}}^{\delta^{(1)}}du_h
+\int_{\beta^{(2)}}^{\delta^{(2)}}du_h
+\cdots+
\int_{\beta^{(p-2)}}^{\delta^{(p-2)}}du_h
+\int_{\beta^{(p-1)}}^{\gamma^{(p)}}du_h.
\end{aligned}
\]
From a combination of equations (1), (2), (3) one now obtains
\[
\frac{Q_h-R_h}{2}
=\int_{\beta^{(1)}}^{\delta^{(1)}}du_h
+\int_{\beta^{(2)}}^{\delta^{(2)}}du_h
+\cdots+
\int_{\beta^{(p-2)}}^{\delta^{(p-2)}}du_h
+\int_{\beta^{(p-1)}}^{\mu}du_h.
\]
Here, if one multiplies by \(2\), the left-hand side is a periodicity modulus; and since the \(\beta\) are points of contact of a contact curve of order \((n-3)\), the points \(\delta^{(1)},\delta^{(2)},\ldots,\delta^{(p-2)},\mu\) must also be points of contact of such a curve. But there are only finitely many such contact curves; hence the point \(\mu\) can always be chosen so that it is not one of their points of contact, and so that the above equation, as well as the assumption that the \(\gamma\) lie on a curve of order \((n-3)\), becomes impossible. Thus we may conclude that whenever \(\Theta\left(\frac P2\right)\) vanishes identically in the indicated way, the contact curve of order \((n-2)\) decomposes into the tangent at \(\mu\) and a contact curve of order \((n-3)\), by one of the points \(x\) coinciding with the point \(\mu\).

As soon as \(s_1\sigma_1+s_2\sigma_2\cdots+s_p\sigma_p=2H+1\), the general periodicity equation of the \(\Theta\)-function gives
\[
\Theta\left(w+\frac P2\right)
=-\Theta\left(-w+\frac P2\right)\cdot e^{\sum\sigma_h w_h},
\]
which can also be brought into the form
\[
e^{-\frac12\sum\sigma_h w_h}\Theta\left(w+\frac P2\right)
=
-e^{\frac12\sum\sigma_h w_h}\Theta\left(-w+\frac P2\right).
\]
The \(2^{p-1}(2^p-1)\) functions
\[
e^{-\frac12\sum\sigma_h w_h}\Theta\left(w+\frac P2\right),
\]
for which \(s_1\sigma_1+s_2\sigma_2\cdots=2H+1\), therefore have the character of odd functions, since they only change sign when all \(w\)'s are made to change their sign simultaneously. Conversely, the remaining \(2^{p-1}(2^p+1)\) functions
\[
e^{-\frac12\sum\sigma_h w_h}\Theta\left(w+\frac P2\right)
\]
have the character of even functions, since they do not change when all \(w\)'s change sign simultaneously. From these latter functions, however, odd functions arise by differentiation with respect to the several arguments \(w\); and if one denotes one of these functions by \(\varphi\), then the expressions
\[
\frac{\partial\varphi}{\partial w_1},\quad
\frac{\partial\varphi}{\partial w_2},\ \ldots,\ 
\frac{\partial\varphi}{\partial w_p}
\]
vanish when all \(w\)'s are simultaneously set equal to zero. Hence, if one sets
\[
w_h+\frac{P_h}{2}
=\sum_{i=1}^{i=p}\int_{\alpha^{(i)}}^{x^{(i)}}du_h-\int_{\mu}^{\xi}du_h
\]
and lets \(\xi\) coincide with \(\mu\), the expressions
\[
\frac{\partial\varphi}{\partial x^{(i)}}
=
\frac{\partial\varphi}{\partial w_1}\frac{\partial w_1}{\partial x^{(i)}}
+
\frac{\partial\varphi}{\partial w_2}\frac{\partial w_2}{\partial x^{(i)}}+\cdots
\]
also vanish. It was proved above that, if one of the functions considered here vanishes while all the \(w\)'s vanish, then one of the points \(x\) falls into \(\mu\), while the other points \(x\) pass into the points of contact of a curve of order \((n-3)\) passing through the double and return points. Since \(\mu\) is arbitrary, whereas the number of contact curves and of their points of contact is finite, one can always suppose \(\mu\) to be different from these contact points. Thus if the function considered is to vanish together with its differential quotients, hence itself at least to second order, the reason cannot be that several of the points \(x\) unite at the point \(\xi\) or \(\mu\), but only that, while one of the \(x\), say \(x^{(p)}\), falls into \(\mu\), all the points \(x\) also lie on a curve of order \((n-3)\) passing through the double and return points.

It follows from this that, among the functions considered above, those which have the character of even functions do not vanish as soon as all \(w\)'s become zero. For in such a case there would have to correspond to the system of half-periodicity moduli
\[
\frac{P_1}{2},\quad \frac{P_2}{2},\ldots,\frac{P_p}{2}
\]
a contact curve of order \((n-3)\) whose points of contact, with each point \(\mu\), that is, with every other point of the curve, lay on a curve of order \((n-3)\) passing through the double and return points. These contact points would therefore have to possess simultaneously two quite different and, in general, by no means coincident properties, defined by quite different algebraic equations: to be points of contact of a contact curve of order \((n-3)\), and to form a system of \(p-1\) points which do not determine a curve of order \((n-3)\) passing through the double and return points.

Thus, in general, only those functions \(\Theta\left(\frac P2\right)\) vanish for which \(\sum s_i\sigma_i\) is odd, and at the same time one has the theorem:

There are \(2^{p-1}(2^p-1)\) curves of order \((n-3)\) which pass through the double and return points and touch the curve \(f=0\) in \(p-1\) points.

The curves are found from the equations
\[
\sum_{i=1}^{i=p-1}\int_{\alpha^{(i)}}^{x^{(i)}}du_h
+\int_{\alpha^{(p)}}^{\mu}du_h
=\frac12P_h
\qquad (h=1,2\ldots p),
\]
which can hold together and provide in the \(x\) the points of contact of the curves, as soon as the integers occurring in the \(P\) satisfy the equation
\[
s_1\sigma_1+s_2\sigma_2\cdots+s_p\sigma_p=2H+1.
\]
The solution of the special class of inversion problems to which the present one belongs was given in §60.

One notes that the case in which all \(P\) vanish is not included among the cases considered here. Hence this case is one of those which are summarized in the following theorem, which supplements the preceding one:

There are \(2^{p-1}(2^p+1)\) curves of order \((n-2)\) which pass through the double and return points, as well as through the \(n-2\) points of intersection of the tangent drawn to \(f=0\) at an arbitrary point \(\mu\), and touch the curve \(f=0\) in \(p\) points.

The points of contact of these curves are found from the equations
\[
\sum_{i=1}^{i=p}\int_{\alpha^{(i)}}^{x^{(i)}}du_h=\frac12P_h,
\qquad (h=1,2\ldots p),
\]
if one takes for the \(P_h\) all those systems of periodicity moduli for which
\[
s_1\sigma_1+s_2\sigma_2\cdots+s_p\sigma_p=2H.
\]

\subsection*{\S\,76. Reduction of the division problems for \(m=2\).}

If one seeks directly by algebraic means to determine a curve \(\varphi=0\) of order \((n-3)\), which touches the curve \(f=0\) in \(p-1\) points and passes through the double and return points, one will set one coefficient of \(\varphi\) equal to \(1\), and from the conditions to be set up will be able to express all the remaining ones rationally by one of them, while this one itself \((\varrho)\) depends on a higher equation \(R=0\). The degree of this equation is, by the preceding, \(2^{p-1}(2^p-1)\). If one knows one of its roots, then the determination of the associated points of contact requires the solution of an equation of degree \(p-1\).

The question of the contact curves of order \((n-2)\), which pass through the double and return points as well as through \(n-2\) fixed points, and touch \(f=0\) in \(p\) points, is by §68 a problem of general division and leads to an equation \(K=0\) of degree \(2^{2p}\). Its solution is obtained, according to §70, by reduction to the special division problem, which consists in seeking the remaining solutions of the above contact problem when one of them is given, and which depends on an equation of degree \(2^{2p}-1\).

But in the present case other circumstances arise than in division by an odd prime number. The special division equation is not the simplest one to which the solution of the general division problem returns; rather, both problems are more simply settled by a special case of the general division problem. This special problem consists in taking the \(n-2\) fixed points of the general problem as the remaining intersection points of the tangent drawn to \(f=0\) at some point \(\mu\). As shown above, in this case only \(2^{p-1}(2^p+1)\) actual contact curves are present, while the others pass into the tangent itself in connection with the various contact curves of order \((n-3)\) which are given by the equation \(R=0\) considered above. If then we form the equation \(K=0\) of the general division problem for this special case, and note that the unknown in it was the parameter of a pencil of rays drawn in one of the \(n-2\) fixed points, namely the parameter of that ray which became tangent to one of the sought contact curves, then \(2^{p-1}(2^p-1)\) of these coincide here with the tangent of the point \(\mu\), which formed a component of just as many decomposed contact curves. The equation \(K=0\) therefore has for this special case a known \(2^{p-1}(2^p-1)\)-fold root, and after its separation there remains an equation \(N=0\) of degree \(2^{p-1}(2^p+1)\), which is to be solved and which supplies the proper contact curves of order \((n-2)\). The roots of this equation of degree \(2^{p-1}(2^p+1)\) are rationally expressible by the roots of the equation \(R=0\) of degree \(2^{p-1}(2^p-1)\).

This rests on the following. For brevity, denote systems of periodicity moduli of the form
\[
2s_h\pi\sqrt{-1}+a_{h1}\sigma_1+a_{h2}\sigma_2+\cdots+a_{hp}\sigma_p,
\qquad (h=1,2\ldots p),
\]
in which the \(s,\sigma\) are to have the values \(0\) or \(1\), by \(G_h\) or \(U_h\), according as \(\sum s_i\sigma_i\) is even or odd. One easily observes that \(2p\) such systems \(U_h\) can be set up in many ways for which the determinant formed from the \(4p^2\) numbers \(s,\sigma\) is equal to \(\pm1\); for example, by using as the different number systems of \(s,\sigma\) the horizontal rows of the schema:
\[
\scriptsize
\begin{array}{c|ccccc|ccccc}
& s_1&s_2&s_3\cdots&s_{p-1}&s_p
& \sigma_1&\sigma_2&\sigma_3\cdots&\sigma_{p-1}&\sigma_p\\ \hline
1)   &1&0&0\cdots&0&0&1&0&0\cdots&0&0\\
2)   &0&1&0\cdots&0&0&0&1&0\cdots&0&0\\
3)   &0&0&1\cdots&0&0&0&0&1\cdots&0&0\\
     & & &\cdots& & & & &\cdots& &\\
p-1) &0&0&0\cdots&1&0&0&0&0\cdots&1&0\\
p)   &0&0&0\cdots&0&1&0&0&0\cdots&0&1\\
p+1) &1&0&0\cdots&0&0&1&1&0\cdots&0&0\\
p+2) &0&1&0\cdots&0&0&0&1&1\cdots&0&0\\
p+3) &0&0&1\cdots&0&0&0&0&1\cdots&0&0\\
     & & &\cdots& & & & &\cdots& &\\
2p-1)&0&0&0\cdots&1&0&0&0&0\cdots&1&1\\
2p)  &0&0&0\cdots&0&1&1&0&0\cdots&0&1
\end{array}
\]
In three of the four squares into which this schema decomposes, only the diagonal elements are different from \(0\); in the last, also those standing immediately to the right of these. If one regards the schema as a determinant and always subtracts the \(i^{\text{th}}\) row from the \((p+i)^{\text{th}}\), there remains a determinant which is plainly equal to \(\pm1\).

It follows that by such \(2p\) systems \(U_h\) one can express the original systems of periodicity moduli
\[
\left(
\begin{array}{cccc|cccc}
2\pi\sqrt{-1}&0&\cdots&0&a_{11}&a_{12}&\cdots&a_{1p}\\
0&2\pi\sqrt{-1}&\cdots&0&a_{21}&a_{22}&\cdots&a_{2p}\\
\vdots&\vdots&\ddots&\vdots&\vdots&\vdots&\ddots&\vdots\\
0&0&\cdots&2\pi\sqrt{-1}&a_{p1}&a_{p2}&\cdots&a_{pp}
\end{array}
\right)
\]
linearly and integrally; and therefore the systems \(G_h\), as well as the remaining systems \(U_h\), can also be represented linearly and integrally by those \(2p\) systems \(U_h\).

Let therefore any system \(G_h\) be connected with those \(2p\) systems \(U_h\) by the equations
\[
G_h=\tau^{(1)}U_h^{(1)}+\tau^{(2)}U_h^{(2)}
+\cdots+\tau^{(2p)}U_h^{(2p)};
\qquad (h=1,2\ldots p),
\]
where the \(\tau\) are whole numbers. Divide this equation by \(2\), and replace the half-periods throughout by the integral sums to which they are equal in the problem of special division. In place of \(\frac12G_h\) there then appears the sum
\[
\sum_{i=1}^{i=p}\int_{\alpha^{(i)}}^{\beta^{(i)}}du_h,
\]
in which the \(\alpha,\beta\) are points of contact of two curves of order \((n-2)\), one of which (the one belonging to the \(\beta\)) depends on the chosen system \(G_h\). The \(\frac12U_h\), however, are to be replaced by sums of the form
\[
\sum_{i=1}^{i=p-1}\int_{\alpha^{(i)}}^{\gamma^{(i)}}du_h
+\int_{\alpha^{(p)}}^{\mu}du_h,
\]
in which the \(\gamma\)'s are points of contact of a curve of order \((n-3)\). If now the roots \(\nu,\nu'\) of \(N=0\) correspond to the \(\alpha,\beta\), and the roots \(\varrho^{(1)},\varrho^{(2)}\ldots\varrho^{(2p)}\) of \(R=0\) correspond to the different systems of \(\gamma\), then by repeated application of Abel's theorem \(\nu'\) is a rational function of \(\nu,\varrho^{(1)},\varrho^{(2)}\ldots\varrho^{(2p)}\). Finally, for every system \(U_h\) not contained among the above \(2p\) systems,
\[
U_h=v^{(1)}U_h^{(1)}+v^{(2)}U_h^{(2)}
+\cdots+v^{(2p)}U_h^{(2p)},\qquad (h=1,2\ldots p),
\]
where the \(v\)'s are whole numbers. It is easy to show that the sum of the numbers \(v\) is always equal to the sum of the numbers \(s\) which occur in the original expression of the system \(U_h\). One can therefore always set up such a system \(U_h\) for which the sum of the \(v\)'s is equal to \(2\). With respect to the \(U^{(1)},U^{(2)}\ldots U^{(2p)}\) given in the above schema, this condition is satisfied, for example, by the system \(U_h\) characterized by the numbers
\[
1\quad1\quad0\cdots0\quad0\ .\quad1\quad0\quad0\cdots0\quad0
\]
and for which
\[
U_h=2U_h^{(1)}+U_h^{(2)}-U_h^{(p+1)}.
\]
If one divides this equation by \(2\) and again inserts the integral sums for the half-periodicity moduli, with systems of \(\gamma\) and \(\mu\) in the upper limits and only the points \(\alpha\) in the lower limits, then in the resulting equation the limits can plainly be brought together so that the lower limits \(\alpha\) occur only in one integral sum, which is moreover multiplied by \(1\). Abel's theorem, applied to the equation thus formed, then shows that the root \(\nu\) of \(N=0\), corresponding to the curve of the \(\alpha\), is rationally expressible by the roots \(\varrho,\varrho^{(1)},\varrho^{(2)},\varrho^{(p+1)}\) which belong to the contact curves of order \((n-3)\) used.

In this way \(\nu\), and therefore also every other root \(\nu'\) of \(N=0\), is expressed rationally by roots of \(R=0\), as was to be proved. These considerations of course suffer an exception for \(p=0\) and \(p=1\).

It is now finally clear how the special division for \(m=2\) is reduced to the problems just treated. If one seeks the solution of the special division problem corresponding to the equation
\[
\sum_{i=1}^{i=p}\int_{c^{(i)}}^{x^{(i)}}du_h=\frac{P_h}{2}
\qquad (h=1,2\ldots p),
\]
one replaces \(\frac{P_h}{2}\) by the integral sum of the corresponding case of that problem, and then finds the symmetric functions of the \(x\), and therefore also a root of the special division equation, rationally expressed by the roots of \(R=0\) and \(N=0\). It deserves to be noted that, accordingly, the solutions of the special division problem likewise split into two classes according as \(P_h\) belongs to the class of the \(U_h\) or to the class of the \(G_h\). The former contains \(2^{p-1}(2^p-1)\), the latter \(2^{p-1}(2^p+1)-1\) period systems. The distinction between the two classes of solutions by geometrical means follows at once from the geometrical interpretation of the case of Abel's theorem which was used for the algebraic solution.

% --- Source pages 292--299; printed pages 270--277. ---

\section*{Eleventh Section.\\The extended inversion problem.}

\subsection*{\S\,77. Reduction of the extended inversion problem to a function \(U^{(q)}\) depending on only \(p+1\) points.}

In the sixth section we gave in detail a solution of the extended inversion problem by which it was reduced to Jacobi's problem. But when the problem is to reduce that problem to \(\Theta\)-functions, it is more suitable to take the other path, indicated briefly in \(\S\,44\). This path treats the extended inversion problem wholly by analogy with the Jacobi problem, by reducing it to the transcendental function
\[
T_{\xi\eta}^{(q)}
\begin{pmatrix}
x^{(1)},&x^{(2)}&\ldots\\
c^{(1)},&c^{(2)}&\ldots
\end{pmatrix}
=\sum_{i=1}^{i=p+q}\int_{c^{(i)}}^{x^{(i)}}d\Pi_{\xi\eta}.
\]
Since this can be treated quite as the transcendental function \(T_{\xi\eta}\) above, it will suffice to emphasize the analogies and briefly set out the special features that enter.

First, however, in order to be able to state the theorems arising in these extensions more simply, we shall reduce the extended inversion problem to a more special formulation, one which can always be given to it without impairing generality. By a one-to-one transformation we can always transform the given curve \(f=0\) into another in which given pairs of points of \(f=0\) coalesce into double points. Other double points will also appear, necessarily conjugate to these, but they are at first of no concern here. We suppose this one-to-one transformation arranged so that the pairs of points \(\xi^{(1)},\eta^{(1)};\ \xi^{(2)},\eta^{(2)},\ldots;\ \xi^{(q)},\eta^{(q)}\) occurring in the extended inversion problem of the sixth section coalesce into double points. In other words, we shall assume that the integrals of the third kind given in the extended inversion problem become discontinuous only at \emph{double points} of the given curve; some of these points may also pass over into cusps, a case to be considered later. Thus assume first that the \(q\) fundamental integrals of the third kind are discontinuous at \(q\) double points of \(f=0\), while besides these there may be \(\delta\) double points and \(r\) cusps, so that
\[
p=\frac{(n-1)\cdot(n-2)}{2}-q-\delta-r.
\]
Since the fundamental integrals of the third kind now outwardly have the form of integrals of the first kind, we shall denote them by
\[
u_{p+1},\ u_{p+2}\ldots u_{p+q},
\]
and so obtain the inversion problem in the form
\[
v_h=\sum_{i=1}^{i=p+q}\int_{c^{(i)}}^{x^{(i)}}du_h,
\qquad (h=1,2\ldots p+q).
\]
If, in analogy with \(\S\,42\), one applies Abel's theorem to a pencil of curves \(\varphi-\lambda\psi=0\) passing through the \(\delta+r\) exceptional points, then one finds the solution of the extended inversion problem, as before that of the Jacobi problem, in the form
\[
\prod_{i=1}^{i=p+q}
\frac{\left(\dfrac{\varphi}{\psi}\right)_{x^{(i)}}-\lambda}
     {\left(\dfrac{\varphi}{\psi}\right)_{c^{(i)}}-\lambda}
=
e^{-\sum T_{\xi\eta}^{(q)}
\begin{pmatrix}
x^{(1)},x^{(2)}\ldots x^{(p+q)}\\
c^{(1)},c^{(2)}\ldots c^{(p+q)}
\end{pmatrix}},
\]
where the sum extends over all pairs of points \(\xi,\eta\), and where the \(\xi\) are the intersection points of \(f=0\) with \(\varphi-\lambda\psi=0\), while the \(\eta\) are the intersection points of \(f=0\) with \(\psi=0\).

In applying Abel's theorem one must here always use curves which pass through the remaining \(\delta+r\) exceptional points, but not through those which are the discontinuity points of the integrals of the third kind given in the inversion problem. In the intersection system of such a curve, when its order exceeds the \((n-3)\)-rd, \(p+q\) points are always determined by the others (\(\S\,9\)); but only \(p+q-1\) are so determined when the curve is of order \(n-3\). In particular, lay a curve of order \(n-2\) through the \(\delta+r\) exceptional points, through the \(n-2\) points in which \(f=0\) is further cut by a line \(\xi\eta\), and through \(p+q\) given points \(x\). It cuts \(f=0\) in \(p+q\) further points \(y\). A transcendental function \(T_{\xi\eta}^{(q)}\) having the \(x\) as upper limits and the corresponding \(y\) as lower limits will be called the special transcendental function \(T_{\xi\eta}^{(q)}(x^{(1)},x^{(2)}\ldots x^{(p+q)})\). Then, exactly as in \(\S\,45\), one can express the general transcendental function \(T_{\xi\eta}^{(q)}\) by two special ones with the aid of the formula
\[
T_{\xi\eta}^{(q)}
\begin{pmatrix}
x^{(1)},x^{(2)}\ldots\\
c^{(1)},c^{(2)}\ldots
\end{pmatrix}
=\frac12\left[
T_{\xi\eta}^{(q)}(x^{(1)},x^{(2)}\ldots)
-
T_{\xi\eta}^{(q)}(c^{(1)},c^{(2)}\ldots)
\right].
\]
The theory of adjoint transcendental functions can also be extended to this case. If one lays a curve of order \(n-3\) through the \(\delta+r\) exceptional points and through \(p+q-1\) of the \(x\), it cuts in \(p+q-1\) further points, which according to the omitted point \(x\) are to be denoted as follows:
\[
\begin{array}{c@{\qquad}l}
x^{(1)}&x^{21},\ x^{31}\ldots\\
x^{(2)}&x^{12},\ x^{32}\ldots\\
x^{(3)}&x^{13},\ x^{23}\ldots\\
&\ldots
\end{array}
\]
If one then forms the adjoint transcendental functions
\[
T^{(1)}=T_{x^{(1)}\eta}(\xi,x^{21},x^{31}\ldots),
\qquad
T^{(2)}=T_{x^{(2)}\eta}(x^{12},\xi,x^{32}\ldots),
\]
and moreover denotes by
\[
T^0=T_{\xi\eta}^{(q)}(x^{(1)},x^{(2)},x^{(3)}\ldots)
\]
the original transcendental function, then again
\[
dU^{(q)}=\frac12\left[
\frac{\partial T^0}{\partial\xi}\,d\xi
+\frac{\partial T^{(1)}}{\partial x^{(1)}}\,dx^{(1)}
+\frac{\partial T^{(2)}}{\partial x^{(2)}}\,dx^{(2)}
+\cdots+
\frac{\partial T^{(p+q)}}{\partial x^{(p+q)}}\,dx^{(p+q)}
\right]
\]
is a complete differential. Introducing constant auxiliary points \(\alpha\), and then constant lower limits \(\zeta,\gamma\), one obtains, exactly as in the seventh section, an expression for the function \(U^{(q)}\) itself, namely
\[
\begin{aligned}
U^{(q)}(x^{(1)},x^{(2)},\ldots;\xi)
={}&
\sum_{i=1}^{i=p+q}\int_{\zeta}^{\xi}d\Pi_{x^{(i)}\alpha^{(i)}}
+\sum_{i=1}^{i=p+q}\int_{\gamma^{(i)}}^{x^{(i)}}d\Pi_{\xi\alpha^{(i)}}\\
&-\frac12\sum_{i>k}\sum\int_{\gamma^{(k)}}^{x^{(k)}}d\Pi_{x^{(i)}\alpha^{(i)}}
-\frac12\sum_{i>k}\sum\int_{\gamma^{(i)}}^{x^{(i)}}d\Pi_{\gamma^{(k)}\alpha^{(k)}}\\
&+\frac12 T_{\xi\zeta}^{(q)}(\alpha^{(1)},\alpha^{(2)}\ldots\alpha^{(p+q)})\\
&+\frac12\sum_{i=1}^{i=p+q}
T_{x^{(i)}\gamma^{(i)}}^{(q)}
(\alpha^{1i},\alpha^{2i}\ldots\alpha^{i-1,i},\alpha^{(i)},\alpha^{i+1,i}\ldots\alpha^{p+q,i})\\
&+\log\frac{%
\begin{gathered}\scriptstyle
\varphi(x^{(1)},x^{(2)}\ldots x^{(p+q)})\,
\varphi(\gamma^{(1)},\alpha^{(2)}\ldots\alpha^{(p+q)})\\[-0.15em]
\scriptstyle{}\cdot
\varphi(\alpha^{(1)},\gamma^{(2)}\ldots\alpha^{(p+q)})\cdots
\varphi(\alpha^{(1)},\alpha^{(2)}\ldots\gamma^{(p+q)})
\end{gathered}}{%
\begin{gathered}\scriptstyle
\varphi(\gamma^{(1)},\gamma^{(2)}\ldots\gamma^{(p+q)})\,
\varphi(x^{(1)},\alpha^{(2)}\ldots\alpha^{(p+q)})\\[-0.15em]
\scriptstyle{}\cdot
\varphi(\alpha^{(1)},x^{(2)}\ldots\alpha^{(p+q)})\cdots
\varphi(\alpha^{(1)},\alpha^{(2)}\ldots x^{(p+q)})
\end{gathered}}.
\end{aligned}
\]
Here the function \(\varphi\), when set equal to zero, gives the condition that the points appended to it as arguments lie with the \(\delta+r\) exceptional points on a curve of order \(n-3\).

One now proves, as in \(\S\,51\), the following two theorems.

I. If the function \(U^{(q)}\) is formed for \(x^{(1)},x^{(2)}\ldots x^{(p+q)}, \xi\), and for \(y^{(1)},y^{(2)}\ldots y^{(p+q)}, \eta\), where the \(x,y\) lie on a curve of order \(n-2\) which passes through the separated \(\delta+r\) exceptional points and through the \(n-2\) remaining intersection points of \(f=0\) with the line \(\xi\eta\), then, with a suitable choice of integration paths, the same value or the opposite value is obtained both times.

II. Let a ray rotate about a point of the curve \(f=0\), and in each position of it let a curve of order \(n-2\) be laid through \(n-2\) of its other intersection points, through the \(\delta+r\) exceptional points of \(f=0\), and through \(p+q\) fixed points of \(f=0\). It then cuts the curve in another \(p+q\) points \(x\), and the function \(U^{(q)}\), formed for these \(p+q\) points and with a parameter corresponding to the last intersection point of the ray, always has the same value.

The second theorem again replaces a partial differential equation satisfied by the function \(U^{(q)}\); it says that this function depends only on the \(p+q\) sums
\[
\sum_{h=1}^{h=p+q}\int_{c^{(h)}}^{x^{(h)}}du_k-\int_{\mu}^{\xi}du_k,
\qquad (k=1,2\ldots p+q),
\]
or on the \(p+q\) arguments
\[
v_k-\int_{\mu}^{\xi}du_k,
\]
where \(\mu\) may be any constant point.

\subsection*{\S\,78. Discontinuities of \(U^{(q)}\). The synectic function \(V^{(q)}=e^{-U^{(q)}}\).}

Up to this point the analogy with the investigations of the seventh section is so complete that one can take the results from it without more. But this changes somewhat when we pass to the investigation of the discontinuities of \(U^{(q)}\). This function becomes infinite first, like the function \(U\) of the seventh section,

\[
\begin{array}{ll}
1)&\text{when }\xi\text{ coincides with one of the points }x;\\
2)&\text{when the points }x\text{ lie on a curve of order }(n-3)\\
&\text{passing through the remaining }\delta+r\text{ exceptional points,}
\end{array}
\]
and in both cases it approaches the negative logarithm of an algebraic function vanishing simply. But the function \(U^{(q)}\) becomes infinite in yet a third case, namely when one of the points \(\xi,x\) approaches one of the \(q\) double points distinguished in the extended inversion problem. To determine what value the function \(U^{(q)}\) approaches in this case, suppose, for instance, that \(x^{(1)}\) approaches infinitely near one of those double points. Then in \(U^{(q)}\) the term
\[
\frac12 T_{x^{(1)}\gamma^{(1)}}^{(q)}(\alpha^{(1)},\alpha^{21},\alpha^{31}\ldots\alpha^{p+q,1})
\]
becomes infinite. For the curve of order \(n-2\) by which one constructs the lower limiting points passes through the \(n-2\) points in which \(f=0\) is further cut by the line \(x^{(1)}\gamma^{(1)}\). If now \(x^{(1)}\) lies near a double point on one branch of \(f=0\), then this line cuts once more near the same point, and indeed the other branch of \(f=0\); hence the curve of order \(n-2\) must pass through a point of the line \(x^{(1)}\gamma^{(1)}\) infinitely near the double point, and therefore cuts near the double point a second further time, again on the first branch of \(f=0\). Thus one further intersection point, that is, a lower limiting point, must approach infinitely near the double point. Let this point be \(\beta^{(1)}\); then the discontinuous part of the above transcendental function will be the term
\[
\frac12\int_{\beta^{(1)}}^{\alpha^{(1)}}d\Pi_{x^{(1)}\gamma^{(1)}},
\]
in which the argument \(\beta^{(1)}\) and the parameter \(x^{(1)}\) approach each other infinitely closely, and indeed on the same one of the two branches of \(f=0\) which intersect at the double point in question.

How this integral behaves in the limit we see from formula (5) of \(\S\,58\), which expresses a single integral of the third kind by \(\Theta\)-functions. According to it, the above integral behaves, as \(\beta^{(1)}\) approaches \(x^{(1)}\), like the logarithm of a \(\Theta\)-function which vanishes by the coincidence of the parameter \(x^{(1)}\) with the argument \(\beta^{(1)}\), thus like \(\log(\beta^{(1)}-x^{(1)})\). Hence \(T^{(q)}\) approaches the value \(\frac12\log(\beta^{(1)}-x^{(1)})\).

Meanwhile, by the approach of \(x^{(1)}\) to one of the \(q\) distinguished double points, the corresponding quantity \(v_{p+i}\) also becomes infinite because of the integral
\[
\int_{c^{(1)}}^{x^{(1)}}du_{p+i},
\]
in which the argument \(x^{(1)}\) approaches one of the parameters, namely that one of the two points thought of as coalesced in the double point on whose branch \(x^{(1)}\) moves. If we denote this double point by \(\delta\), the above integral, and hence also \(v_i\), behaves in the limit like \(\pm\log(\delta-x^{(1)})\); the upper or lower sign is to be chosen according as \(x^{(1)}\) moves on one or the other branch, that is, according as \(x^{(1)}\) approaches the one or the other parameter.

During the mentioned motion of the points \(\beta\) and of the point \(x^{(1)}\), by Abel's theorem the \(p+q\) expressions mentioned above as arguments of a function \(U^{(q)}\),
\[
\sum_{h=1}^{h=p+q}\int_{c^{(h)}}^{\beta^{(h)}}du_k-\int_{\mu}^{x^{(1)}}du_k,
\qquad (k=1,2\ldots p+q),
\]
retain constant finite values. But when \(x^{(1)}\), and with it one of the \(\beta\), moves toward the double point \(\delta\), then in the expression
\[
\sum_{h=1}^{h=p+q}\int_{c^{(h)}}^{\beta^{(h)}}du_{p+i}-\int_{\mu}^{x^{(1)}}du_{p+1}
\]
the two integrals
\[
\int_{c^{(1)}}^{\beta^{(1)}}du_{p+i}-\int_{\mu}^{x^{(1)}}du_{p+i}
\]
become infinite and approach the expressions
\[
\pm\left[\log(\beta^{(1)}-\delta)-\log(x^{(1)}-\delta)\right].
\]
Thus the quantity
\[
\log\frac{\beta^{(1)}-\delta}{x^{(1)}-\delta}
\]
remains finite, or the sum
\[
T_{x^{(1)}\gamma^{(1)}}^{(q)}(\alpha^{(1)},\alpha^{21}\ldots)\pm v_{p+i}
\]
retains a finite value. Therefore one may say that the function \(U^{(q)}\) approaches the infinitely large value \(\pm\frac12v_{p+i}\) when one of the points \(x\) approaches the \(i\)-th of the \(q\) distinguished double points.

Noting then that the function \(U^{(q)}\) can become infinite only when one of the quantities \(v\) itself assumes an infinitely large value, or in the cases cited above under (1), (2), in which it approaches the negative logarithm of a quantity vanishing simply, one concludes, as in \(\S\,53\), that the function
\[
V^{(q)}=e^{-U^{(q)}}
\]
is a synectic function of the \(p+q\) arguments
\[
\sum_{h=1}^{h=p+q}\int_{c^{(h)}}^{x^{(h)}}du_k-\int_{\mu}^{\xi}du_k,
\qquad (k=1,2\ldots p+q).
\]
And this function vanishes only, and always, when either \(\xi\) coincides with one of the \(x\), or when the points \(x\) lie on a curve of order \(n-3\) passing through the separated \(\delta+r\) exceptional points.

The theorem stated above under I. now gives still the property of the function \(V^{(q)}\) (cf. \(\S\,53\)) that, considered as a function of the \(p+q\) arguments
\[
w_h=\sum_{i=1}^{i=p+q}\int_{c^{(i)}}^{x^{(i)}}du_h-\int_{\mu}^{\xi}du_h+K_h,
\]
with a suitable choice of the constants \(K\), it remains unchanged up to sign when all \(w\) simultaneously change sign. The \(K\) must be so chosen that, if the \(y\) denote those \(p+q\) points which, together with the \(x\), with the \(n-2\) remaining intersection points of a line \(\xi\eta\), and with the \(\delta+r\) exceptional points, lie on a curve of order \(n-2\), then
\[
(1)\quad
2K_h+
\sum_{i=1}^{i=p+q}\int_{c^{(i)}}^{x^{(i)}}du_h+
\sum_{i=1}^{i=p+q}\int_{c^{(i)}}^{y^{(i)}}du_h
-\int_{\mu}^{\xi}du_h-\int_{\mu}^{\eta}du_h=0,
\]
whereby Abel's theorem gives constant values for the \(K_h\).

\subsection*{\S\,79. Series expansion of the function \(V^{(q)}\) from its periodicity properties.}

We now determine the expression of the synectic function \(V^{(q)}\) completely by its periodicity moduli. First form the transcendental functions
\[
\frac12T_{\xi\eta}^{(q)}(x^{(1)},x^{(2)}\ldots x^{(p+q)})
=
\log
\frac{V^{(q)}(x^{(1)},x^{(2)}\ldots x^{(p+q)};\eta)}
     {V^{(q)}(x^{(1)},x^{(2)}\ldots x^{(p+q)};\xi)},
\]
\[
T_{\xi\eta}^{(q)}
\begin{pmatrix}
x^{(1)},x^{(2)}\ldots x^{(p+q)}\\
z^{(1)},z^{(2)}\ldots z^{(p+q)}
\end{pmatrix}
=
\log
\frac{
V^{(q)}(x^{(1)},x^{(2)}\ldots;\eta)
V^{(q)}(z^{(1)},z^{(2)}\ldots;\xi)
}{
V^{(q)}(x^{(1)},x^{(2)}\ldots;\xi)
V^{(q)}(z^{(1)},z^{(2)}\ldots;\eta)
}.
\]
From the last formula we obtain, exactly as in \(\S\,54\), the change undergone by \(V^{(q)}\) when the \(w\) on which \(V^{(q)}(x^{(1)},x^{(2)}\ldots;\xi)\) depends increase by periodicity moduli. The most general change of the quantities \(w_1,w_2\ldots w_p\) consists in \(w_i\) passing into
\[
w_i'=w_i+2m_i\pi\sqrt{-1}+a_{1i}q_1+a_{2i}q_2+\cdots+a_{pi}q_p,
\qquad (i=1,2\ldots p).
\]

% --- Source pages 300--308; printed pages 278--286. ---

passes over. Then the remaining \(w\) change by the periodicity moduli which belong to the integrals of the third kind occurring in them; and, in addition, one may add mutually independent periodicity moduli which are integral multiples of \(2\pi\sqrt{-1}\). Let \(u_{ki}\) denote the quantity by which the integral \(u_{p+i}\) increases when \(w_1,w_2\ldots w_p\) increase by the periodicity moduli \(a_{1k},a_{2k}\ldots a_{pk}\). Then the most general values which the quantities \(w_{p+1},w_{p+2}\ldots\) can assume at the same time as the passage of \(w_1,w_2\ldots w_p\) into \(w'_1,w'_2\ldots w'_p\) are the following:
\[
w'_{p+i}=w_{p+i}+2n_i\pi\sqrt{-1}+u_{1i}q_1+u_{2i}q_2+\cdots+u_{pi}q_p,
\qquad (i=1,2\ldots q).
\]
Concerning the \(u_{ki}\) we remark that, for each integral \(u_{p+i}\), in the imaginary plane we must think of the parameter points united in the double point as being joined by a curve which starts from the point representing the double point and returns to it, but in such a way that the beginning of this curve corresponds to the one branch of \(f=0\) through that double point, and the end to the other. The integral \(u_{ki}\), according to \(\S\,34\), (6), is then the integral \(u_k\) carried over this curve, where it must always be fixed beforehand which branch of \(f=0\) is to correspond to the first parameter and which to the second, that is, with which this path of integration begins and with which it ends.

Thus, while the \(p+q\) quantities \(w\) transform into the \(p+q\) quantities \(w'\), \(V^{(q)}\) passes into
\[
(2)\quad
V'^{(q)}=e^{w_1q_1+w_2q_2+\cdots+w_pq_p+C}\cdot V^{(q)},
\]
where \(C\) is a constant. Exactly as before, one finds for \(C\) the expression
\[
C=\frac12\sum\sum q_hq_ka_{hk}+L\pi\sqrt{-1},
\]
where \(L\) is an integer.

If we form \(L\) for the number system \(m,n,q\) used here, and denote by \(L'\), \(L''\) the numbers into which \(L\) passes when \(m,n,q\) are respectively replaced by other numbers \(m',n',q'\), or by the sums \(m+m',n+n',q+q'\), then, as before, one has
\[
L+L'=L'',
\]
and therefore concludes
\[
\begin{aligned}
L={}&\varkappa_1m_1+\varkappa_2m_2+\cdots+\varkappa_pm_p
+\lambda_1q_1+\lambda_2q_2+\cdots+\lambda_pq_p\\
&+\nu_1n_1+\nu_2n_2+\cdots+\nu_qn_q,
\end{aligned}
\]
where the \(\varkappa,\lambda,\nu\) denote fixed integers depending only on the original definition of the transcendental function \(\frac12T_{\xi\eta}^{(q)}(x^{(1)},x^{(2)}\ldots)\), and for which it suffices to assign the values 0 or 1.

Now introduce, instead of the \(w\), the arguments
\[
\begin{aligned}
\omega_h&=w_h-\lambda_h\pi\sqrt{-1}-\frac12(\varkappa_1a_{1h}+\varkappa_2a_{2h}+\cdots+\varkappa_pa_{ph}),
&& (h=1,2\ldots p),\\
\omega_{p+h}&=w_{p+h}-\nu_h\pi\sqrt{-1}-\frac12(\varkappa_1u_{1h}+\varkappa_2u_{2h}+\cdots+\varkappa_pu_{ph}),
&& (h=1,2\ldots q).
\end{aligned}
\]
These differ from the former quantities by half-periodicity moduli. Thus one may give equation (2), or the equation arising from it,
\[
(3)\quad
V'^{(q)}=
 e^{\sum n_hq_h+\frac12\sum\sum q_hq_ka_{hk}+
 \pi\sqrt{-1}(\sum\varkappa_hm_h+\sum\lambda_hq_h+\sum\nu_hn_h)}\cdot V^{(q)},
\]
the form
\[
\begin{aligned}
&e^{-\frac12(\varkappa_1\omega'_1+\varkappa_2\omega'_2+\cdots+\nu_1\omega'_{p+1}+\nu_2\omega'_{p+2}+\cdots)}\,V'^{(q)}\\
&=e^{-\frac12(\varkappa_1\omega_1+\varkappa_2\omega_2+\cdots+\nu_1\omega_{p+1}+\nu_2\omega_{p+2}+\cdots)}\,V^{(q)}\\
&\qquad\cdot e^{q_1\omega_1+q_2\omega_2+\cdots+\frac12\sum\sum q_hq_ka_{hk}-\frac12\sum\sum q_h\nu_ku_{hk}}.
\end{aligned}
\]
Here the \(\omega'\) are thought of as having arisen from the \(\omega\) by the same additions of periodicity moduli by which the \(w'\) arose from the \(w\). The function
\[
e^{-\frac12(\varkappa_1\omega_1+\varkappa_2\omega_2+\cdots+\nu_1\omega_{p+1}+\nu_2\omega_{p+2}+\cdots)}\,V^{(q)}
\]
therefore does not change when the \(\omega\) independently increase by multiples of \(2\pi\sqrt{-1}\), and it changes by the factor
\[
e^{q_1\omega_1+q_2\omega_2+\cdots+q_p\omega_p+\frac12\sum\sum q_hq_ka_{hk}-\frac12\sum\sum q_h\nu_ku_{hk}}
\]
when the \(\omega\) increase respectively by the systems
\[
\begin{array}{rcl}
\omega_i&:&a_{1i}q_1+a_{2i}q_2+\cdots+a_{pi}q_p,\\
\omega_{p+i}&:&u_{1i}q_1+u_{2i}q_2+\cdots+u_{pi}q_p.
\end{array}
\]

By the first property the above function can be expanded into an always convergent series of the form
\[
\sum_{r,s} A_{r_1r_2\ldots r_ps_1s_2\ldots s_q}
 e^{r_1\omega_1+r_2\omega_2+\cdots+r_p\omega_p+s_1\omega_{p+1}+s_2\omega_{p+2}+\cdots+s_q\omega_{p+q}},
\]
where the numbers \(r,s\), independently of one another, are to receive all integral values from \(-\infty\) to \(+\infty\). The last property then gives the identical relation
\[
\begin{aligned}
&\sum_{r,s} A_{r_1r_2\ldots s_1s_2\ldots}
 e^{\sum r_i\omega_i+\sum s_i\omega_{p+i}+\sum\sum q_kr_ia_{ki}+\sum\sum q_ks_iu_{ki}}\\
&=e^{q_1\omega_1+q_2\omega_2+\cdots+q_p\omega_p+\frac12\sum\sum q_hq_ka_{hk}-\frac12\sum\sum q_h\nu_ku_{hk}}
\sum_{r,s} A_{r_1r_2\ldots s_1s_2\ldots}e^{\sum r_i\omega_i+\sum s_i\omega_{p+i}}.
\end{aligned}
\]
Putting on the left \(r_1+q_1,r_2+q_2\ldots\) in place of \(r_1,r_2\ldots\), and then comparing, one obtains the relation between the coefficients:
\[
\begin{aligned}
&A_{r_1+q_1,r_2+q_2,\ldots,s_1,s_2,\ldots}
 e^{\sum\sum q_kr_ia_{ki}+\sum\sum q_ks_iu_{ki}+\sum\sum q_kq_ia_{ki}}\\
&=A_{r_1r_2\ldots s_1s_2\ldots}
 e^{\frac12\sum\sum q_hq_ka_{hk}-\frac12\sum\sum q_h\nu_ku_{hk}}.
\end{aligned}
\]
Or, when all \(r\) are made to vanish and \(A_{s_1s_2\ldots}\) is written for \(A_{00\ldots s_1s_2\ldots}\),
\[
A_{q_1q_2\ldots s_1s_2\ldots}=A_{s_1s_2\ldots}
 e^{-\frac12\sum\sum q_hq_ka_{hk}-\sum\sum q_ks_iu_{ki}-\frac12\sum\sum q_h\nu_ku_{hk}}.
\]
The series to be determined thus transforms into
\[
\begin{aligned}
&\sum_s\left\{\sum_r e^{-\frac12\sum\sum r_hr_ka_{hk}+\sum r_h\omega_h-
\sum\sum r_h\left(s_i+\frac12\nu_i\right)u_{hi}}\right\}
 A_{s_1s_2\ldots}e^{\sum s_i\omega_{p+i}}\\
&=\sum_s A_{s_1s_2\ldots s_q}e^{\sum s_i\omega_{p+i}}
 \Theta\left[\omega_h-\left(s_1+\frac{\nu_1}{2}\right)u_{h1}-\left(s_2+\frac{\nu_2}{2}\right)u_{h2}\cdots\right],
\end{aligned}
\]
where only one argument is written for the \(\Theta\)-function instead of all. Thus, for the function \(V^{(q)}\), one has the expression
\[
\begin{aligned}
V={}&e^{\frac12(\sum\varkappa_h\omega_h+\sum\nu_h\omega_{p+h})}\,\cdot
\sum_s A_{s_1s_2\ldots s_q}e^{\sum s_i\omega_{p+i}}\\
&\cdot\Theta\left[\omega_h-\left(s_1+\frac{\nu_1}{2}\right)u_{h1}-\left(s_2+\frac{\nu_2}{2}\right)u_{h2}\cdots\right].
\end{aligned}
\]
This expression simplifies very considerably if one recalls (cf. \(\S\,78\)) that, for an infinitely large value of \(v_{p+i}\), and hence also of \(w_{p+i}\), the function behaves like \(e^{\pm\frac12w_{p+i}}\). It follows that all the numbers \(\nu\) must be put equal to 1, and that then only those terms in the sum may be retained in which the \(s\) take the values 0 or \(-1\), so that in the expansion of \(V^{(q)}\) no higher power of \(e^{\pm\frac12w_{p+i}}\) than the first occurs. If one then again denotes \(2s_i+\nu_i\) by \(s_i\), the expression of \(V^{(q)}\) becomes
\[
\begin{aligned}
(1)\quad V^{(q)}={}&e^{\frac12\sum\varkappa_h\omega_h}
\sum_{s_1,s_2,\ldots=\pm1} A_{s_1s_2\ldots s_q}
 e^{\frac12\sum s_i\omega_{p+i}}\\
&\cdot\Theta\left(\omega_h-\frac{s_1u_{h1}+s_2u_{h2}+\cdots}{2}\right),
\end{aligned}
\]
where the sum now still consists of \(2^q\) terms, whose coefficients \(A_{s_1s_2\ldots}\) remain to be determined.

\subsection*{\S\,80. Determination of the coefficients \(A\) of the function \(V^{(q)}\).}

Before passing to the determination of the coefficients \(A\), we shall examine more closely the arguments of the \(\Theta\)-functions occurring in \(V\). The quantities \(\omega\) can be represented in the form
\[
\omega_h=\sum_{i=1}^{i=p+q}\int_{\beta^{(i)}}^{x^{(i)}}du_h-\int_{\mu}^{\xi}du_h,
\qquad (h=1,2\ldots p+q),
\]
provided the lower limiting points \(\beta\) are suitably determined. It is easy to prove that the \(\beta\) are the points of contact of a curve of order \(n-2\), which passes through the separated \(\delta+r\) double points and cusps and through the \(n-2\) remaining intersection points of the tangent drawn to \(f=0\) at the point \(\mu\), and which touches in \(p+q\) points. Indeed, in equation (1) the \(y\) pass into the \(\beta\) if the \(\beta\) are put in place of the \(x\) and \(\xi,\eta\) are made to coincide with \(\mu\). Then from (1) one has
\[
K_h+\sum_{i=1}^{i=p+q}\int_{c^{(i)}}^{\beta^{(i)}}du_h
\]
equal to a half-periodicity modulus, and this changes according to the chosen system of the \(\beta\). Thus
\[
w_h=v_h+K_h-\int_{\mu}^{\xi}du_h
\]
differs from the expression
\[
\sum_{i=1}^{i=p+q}\int_{\beta^{(i)}}^{x^{(i)}}du_h-\int_{\mu}^{\xi}du_h
\]
only by the half of a system of periodicity moduli; the same holds for the \(\omega_h\), however the system of \(\beta\) is chosen. Hence this system can finally be chosen so that the above expression agrees exactly with \(\omega_h\).

Now suppose that a curve of order \(n-2\) is moved so that it always passes through the \(n-2\) remaining intersection points of the tangent at \(\mu\), and also through the separated \(\delta+r\) exceptional points. Let the initial curve be the one which touches at the points \(\beta^{(1)},\beta^{(2)}\ldots\beta^{(p+q)}\); let the final curve be the curve used in the eighth section, which passes through all double points and cusps and touches at the \(p\) points \(\alpha\) which appeared as lower limiting points in the Jacobi inversion problem. If one forms, for the transition from one curve to the other, the transcendental equations of Abel's theorem, then the \(q\) equations corresponding to the \(q\) distinguished double points become illusory; but for the integrals of the first kind one obtains equations of the form
\[
\begin{aligned}
(1)\quad
2\sum_{i=1}^{i=p}\int_{\beta^{(i)}}^{\alpha^{(i)}}du_h
+\sum_{i=1}^{i=q}\left(
\int_{\beta^{(p+i)}}^{\varepsilon^{(+i)}}du_h
+\int_{\beta^{(p+i)}}^{\varepsilon^{(-i)}}du_h\right)=0,\\
(h=1,2\ldots p),
\end{aligned}
\]
where \(\varepsilon^{(+i)},\varepsilon^{(-i)}\) denote the two points of \(f=0\) united in one of the preferred double points, distinguished according to the branch of the curve on which one approaches the double point. In consequence of this equation one can give the quantities \(\omega_1,\omega_2\ldots\omega_p\) the form
\[
\begin{aligned}
\omega_h={}&\sum_{i=1}^{i=p}\int_{\alpha^{(i)}}^{x^{(i)}}du_h
+\frac12\sum_{i=1}^{i=q}\left(\int_{\varepsilon^{(+i)}}^{x^{(p+i)}}du_h+
\int_{\varepsilon^{(-i)}}^{x^{(p+i)}}du_h\right)-\int_{\mu}^{\xi}du_h,\\
&\hfill (h=1,2\ldots p),
\end{aligned}
\]
or also
\[
\omega_h=\sum_{i=1}^{i=p}\int_{\alpha^{(i)}}^{x^{(i)}}du_h
+\sum_{i=1}^{i=q}\int_{\varepsilon^{(is_i)}}^{x^{(p+i)}}du_h
+\frac12\sum_{i=1}^{i=q}\int_{\varepsilon^{(-is_i)}}^{\varepsilon^{(is_i)}}du_h
-\int_{\mu}^{\xi}du_h,
\]
where the numbers \(s_1,s_2\ldots s_q\), as above, have the meaning \(\pm1\).

Now we can evidently choose the integrals \(u_{hi}\), which above were determined only up to periodicity moduli, so that they agree with those integrals which occur in the last form, beginning and ending at the same double point, and which are exactly determined by the transition of curves carried out above. Thus, if \(\varepsilon^{(i)}\) is thought of as corresponding to the first parameter of the integral \(u_{p+i}\), and \(\varepsilon^{(-i)}\) to the second, then
\[
\int_{\varepsilon^{(-is_i)}}^{\varepsilon^{(is_i)}}du_h=s_i\cdot u_{hi},
\]
and one can put the above equation for \(\omega_h\) in the form
\[
\omega_h-\frac12(s_1u_{h1}+s_2u_{h2}+\cdots)
=\sum_{i=1}^{i=p}\int_{\alpha^{(i)}}^{x^{(i)}}du_h
+\sum_{i=1}^{i=q}\int_{\varepsilon^{(is_i)}}^{x^{(p+i)}}du_h-\int_{\mu}^{\xi}du_h.
\]
Here one has before oneself the arguments of the different \(\Theta\)-functions occurring in \(V^{(q)}\). They have the form of the \(\omega_h\) themselves; only here the lower limits are replaced by the points \(\alpha\) occurring in the original inversion problem and by the preferred double points themselves, for which the appropriate direction of continuation must always be chosen.

By \(\S\,78\), the function \(V^{(q)}\) vanishes when one of the points \(x\) coincides with \(\xi\). Let \(x^{(p+i)}\) coincide with \(\xi\), and let \(x^{(p+1)},x^{(p+2)}\ldots x^{(p+i-1)},x^{(p+i+1)}\ldots x^{(p+q)}\) coincide with \(\varepsilon^{(1\cdot s_1)},\varepsilon^{(2\cdot s_2)}\ldots\varepsilon^{((i-1)s_{i-1})},\varepsilon^{((i+1)s_{i+1})}\ldots\varepsilon^{(qs_q)}\). Thus these last \(x\) approach, on predetermined branches, \(q-1\) of the preferred double points. Then \(V^{(q)}\) vanishes, while at the same time the quantities \(\omega_{p+1},\omega_{p+2}\ldots\omega_{p+i-1},\omega_{p+i+1},\ldots\omega_{p+q}\) acquire infinitely large values; and indeed, since \(\varepsilon^{(k)}\) represents the first parameter and \(\varepsilon^{(-k)}\) the second, \(e^{s_k\omega_k}\) becomes infinitely large as soon as one of the points \(x\) approaches the point \(\varepsilon^{(ks_k)}\). If one therefore restricts the function \(V^{(q)}\) to those terms with respect to which all the others become infinitely small, one obtains the equation
\[
\begin{aligned}
(2)\quad 0={}&A_{s_1s_2\ldots s_{i-1},1,s_{i+1}\ldots s_q}
 e^{\frac12\omega_{p+i}}
 \Theta\left(\sum_{k=1}^{k=p}\int_{\alpha^{(k)}}^{x^{(k)}}du_h+
 \int_{\varepsilon^{(i)}}^{\mu}du_h\right)\\
&+A_{s_1s_2\ldots s_{i-1},-1,s_{i+1}\ldots s_q}
 e^{-\frac12\omega_{p+i}}
 \Theta\left(\sum_{k=1}^{k=p}\int_{\alpha^{(k)}}^{x^{(k)}}du_h+
 \int_{\varepsilon^{(-i)}}^{\mu}du_h\right),
\end{aligned}
\]
which must hold for all arbitrary points \(x^{(1)},x^{(2)}\ldots x^{(p)}\), provided that at the same time the quantity \(\omega_{p+i}\) is given its value:
\[
\omega_{p+i}=\sum_{k=1}^{k=p}\int_{\beta^{(k)}}^{x^{(k)}}du_{p+i}
+\sum_{k=1}^{k=i-1}\int_{\beta^{(p+k)}}^{\varepsilon^{(ks_k)}}du_{p+i}
+\int_{\beta^{(p+i)}}^{\mu}du_{p+i}
+\sum_{k=i+1}^{k=q}\int_{\beta^{(p+k)}}^{\varepsilon^{(ks_k)}}du_{p+i}.
\]
Now in particular we can give the points \(x\) such positions \(\gamma\) that
\[
(3)\quad
2\sum_{k=1}^{k=p}\int_{\alpha^{(k)}}^{\gamma^{(k)}}du_h+
\int_{\varepsilon^{(i)}}^{\mu}du_h+
\int_{\varepsilon^{(-i)}}^{\mu}du_h=0,
\qquad (h=1,2\ldots p).
\]
In this case the two \(\Theta\)-functions in equation (2) become equal, and one obtains
\[
(4)\quad
A_{s_1s_2\ldots s_{i-1},-1,s_{i+1}\ldots s_q}
=-e^{\omega_{p+i}}A_{s_1s_2\ldots s_{i-1},1,s_{i+1}\ldots s_q},
\]
where now
\[
\omega_{p+i}=\sum_{k=1}^{k=p}\int_{\beta^{(k)}}^{\gamma^{(k)}}du_{p+i}
+\sum_{k=1}^{k=i-1}\int_{\beta^{(p+k)}}^{\varepsilon^{(ks_k)}}du_{p+i}
+\int_{\beta^{(p+i)}}^{\mu}du_{p+i}
+\sum_{k=i+1}^{k=q}\int_{\beta^{(p+k)}}^{\varepsilon^{(ks_k)}}du_{p+i}.
\tag{5}
\]
Equation (3) defines the \(\gamma\) as the points of contact of a curve of order \(n-2\) passing through all double points and cusps. Think of the tangent at \(\mu\) and an arbitrary ray through the double point \(\varepsilon^{(i)},\varepsilon^{(-i)}\) as drawn. Then the integral sum
\[
\int_{\varepsilon^{(i)}}^{\mu}du_h+
\int_{\varepsilon^{(-i)}}^{\mu}du_h
\]
can be replaced by the negative sum of the \(n-2\) integrals whose upper limits are the intersection points of the tangent and whose lower limits are the intersection points of the ray. If these integrals are put into equation (3), then, according to Abel's theorem, it says that, just as the \(\alpha\) are the points of contact of a curve of order \(n-2\) laid through all double points and cusps and the \(n-2\) intersection points of the tangent, so the \(\gamma\) become the points of contact of such a curve passing through all double points and cusps and the \(n-2\) intersection points of the ray; but the latter curve is thereby cut by the ray in \(n+2\) points, and therefore must contain this whole ray. Hence the \(\gamma\) are the contact points of a curve of order \(n-3\), which passes through all double points and cusps, one excepted, and touches \(f=0\) in \(p\) points.

Since these points \(\gamma\) are assigned to the double point \(\varepsilon^{(i)},\varepsilon^{(-i)}\), we shall denote them by
\[
\gamma^{i,1},\gamma^{i,2},\ldots,\gamma^{i,p},
\]
and further put
\[
\begin{aligned}
a_i={}&\sum_{k=1}^{k=p}\int_{\beta^{(k)}}^{\gamma^{i,k}}du_{p+i}
+\int_{\beta^{(p+i)}}^{\mu}du_{p+i}
+\frac12\sum_{k=1}^{k=i-1}\left(\int_{\beta^{(p+k)}}^{\varepsilon^{(k)}}du_{p+i}+\int_{\beta^{(p+k)}}^{\varepsilon^{(-k)}}du_{p+i}\right)\\
&+\frac12\sum_{k=i+1}^{k=q}\left(\int_{\beta^{(p+k)}}^{\varepsilon^{(k)}}du_{p+i}+\int_{\beta^{(p+k)}}^{\varepsilon^{(-k)}}du_{p+i}\right),\\
b_{ik}=b_{ki}={}&\int_{\varepsilon^{(-i)}}^{\varepsilon^{(i)}}du_{p+k}=\int_{\varepsilon^{(-k)}}^{\varepsilon^{(k)}}du_{p+i}.
\end{aligned}
\]
The various equations derivable from (4) then evidently give
\[
A_{s_1s_2\ldots s_q}=(-1)^{\frac{q-s_1-s_2-\cdots-s_q}{2}}
 e^{-\frac12\sum_i s_ia_i-\frac14\sum_{i\ge k}\sum s_is_kb_{ik}}\cdot C,
\]
and consequently one obtains
\[
V^{(q)}=e^{\frac12\sum\varkappa_h\omega_h}\cdot C\cdot \Theta^{(q)}(\omega_1,\omega_2\ldots\omega_{p+q}),
\]
where \(\Theta^{(q)}\) is understood to mean the function
\[
\begin{aligned}
(6)\quad \Theta^{(q)}(\omega)={}&
\sum_{s_1s_2\ldots s_q=\pm1}
(-1)^{\frac{-s_1-s_2-\cdots-s_q}{2}}
 e^{\frac12\sum_i s_i(\omega_{p+i}-a_i)-\frac14\sum\sum s_is_kb_{ik}}\\
&\cdot\Theta\left(\omega_h-\frac{s_1u_{h1}+s_2u_{h2}+\cdots}{2}\right).
\end{aligned}
\]
The arguments occurring in the exponential quantities here have a form quite similar to that of the arguments of the \(\Theta\)-function; for
\[
\begin{aligned}
\omega_{p+i}-a_i-\sum_{k\ge i}s_kb_{ik}
={}&\sum_{k=1}^{k=p}\int_{\gamma^{i,k}}^{x^{(k)}}du_{p+i}
+\sum_{k=1}^{k=i-1}\int_{\varepsilon^{(ks_k)}}^{x^{(p+k)}}du_{p+i}
+\int_{\xi}^{x^{(p+i)}}du_{p+i}\\
&+\sum_{k=i+1}^{k=q}\int_{\varepsilon^{(ks_k)}}^{x^{(p+k)}}du_{p+i}.
\end{aligned}
\]
The constant \(C\) is determined by the condition that \(U^{(q)}=-\log V^{(q)}\) must vanish for \(\xi=\zeta\), \(x^{(i)}=\gamma^{(i)}\). If, therefore, one writes \(\omega_1^\circ,\omega_2^\circ\ldots\omega_{p+q}^\circ\) for the values assumed by the \(w\) when \(\xi\) falls into \(\zeta\) and the \(x\) into the \(\gamma\), then one has
\[
V^{(q)}=e^{\frac12\sum\varkappa_h(\omega_h-\omega_h^\circ)}\cdot
\frac{\Theta^{(q)}(\omega)}{\Theta^{(q)}(\omega^\circ)}.
\]
With the aid of the transcendental function \(\Theta^{(q)}\), the solution of the extended inversion problem proceeds as that of the Jacobi problem. If, as before, one writes only one argument of the function \(\Theta^{(q)}\), one has
\[
\begin{aligned}
\frac12T_{\xi\eta}^{(q)}(x^{(1)},x^{(2)}\ldots x^{(p+q)})
&=\log\frac{V^{(q)}(x^{(1)}x^{(2)}\ldots x^{(p+q)};\eta)}{V^{(q)}(x^{(1)}x^{(2)}\ldots x^{(p+q)};\xi)}\\
&=\frac12\sum\varkappa_h\int_{\eta}^{\xi}du_h+
\log\frac{\Theta^{(q)}\left(\sum_{i=1}^{i=p+q}\int_{\beta^{(i)}}^{x^{(i)}}du_h-\int_{\mu}^{\eta}du_h\right)}{\Theta^{(q)}\left(\sum_{i=1}^{i=p+q}\int_{\beta^{(i)}}^{x^{(i)}}du_h-\int_{\mu}^{\xi}du_h\right)}.
\end{aligned}
\]
Furthermore,
\[
\begin{aligned}
&T_{\xi\eta}^{(q)}
\begin{pmatrix}
x^{(1)},x^{(2)}\ldots x^{(p+q)}\\
z^{(1)},z^{(2)}\ldots z^{(p+q)}
\end{pmatrix}\\
&=\log\frac{
\Theta^{(q)}\left(\sum_{i=1}^{i=p+q}\int_{\beta^{(i)}}^{x^{(i)}}du_h-\int_{\mu}^{\eta}du_h\right)
\Theta^{(q)}\left(\sum_{i=1}^{i=p+q}\int_{\beta^{(i)}}^{z^{(i)}}du_h-\int_{\mu}^{\xi}du_h\right)
}{
\Theta^{(q)}\left(\sum_{i=1}^{i=p+q}\int_{\beta^{(i)}}^{z^{(i)}}du_h-\int_{\mu}^{\eta}du_h\right)
\Theta^{(q)}\left(\sum_{i=1}^{i=p+q}\int_{\beta^{(i)}}^{x^{(i)}}du_h-\int_{\mu}^{\xi}du_h\right)
}.
\end{aligned}
\]
If \(\varphi=0\), \(\psi=0\) are the equations of two curves of the same high order, and if the \(\xi\) denote the intersection points of \(f=0\) with \(\varphi-\lambda\psi=0\), while the \(\eta\) denote the intersection points of \(f=0\) with \(\psi=0\), then the symmetric functions of the values which \(\frac{\varphi}{\psi}\) assumes at the points \(x^{(1)},x^{(2)}\ldots x^{(p+q)}\) are found by means of the formula
\[
\prod_{i=1}^{i=p+q}
\frac{\left(\dfrac{\varphi}{\psi}\right)_{x^{(i)}}-\lambda}{\left(\dfrac{\varphi}{\psi}\right)_{\beta^{(i)}}-\lambda}
=\prod_{\xi\eta}
\frac{
\Theta^{(q)}\left(v_h-\int_{\mu}^{\xi}du_h\right)
\Theta^{(q)}\left(\int_{\mu}^{\eta}du_h\right)
}{
\Theta^{(q)}\left(v_h-\int_{\mu}^{\eta}du_h\right)
\Theta^{(q)}\left(\int_{\mu}^{\xi}du_h\right)
},
\]
where the \(v\) denote the integral sums
\[
v_h=\sum_{i=1}^{i=p+q}\int_{\beta^{(i)}}^{x^{(i)}}du_h,
\qquad (h=1,2\ldots p+q).
\]

% --- Source pages 309--321; printed pages 287--299. ---

\subsection*{\S\,81. With a suitable arrangement the function \(\Theta^{(q)}\) is a function of its arguments that is even or odd with \(q\).}

We now wish to investigate what becomes of the function \(\Theta^{(q)}(\omega)\) when all its arguments change sign. If in the formula for \(\Theta^{(q)}\) one puts \(-\omega\) in place of \(\omega\), and at the same time \(-s_1,-s_2\ldots\) in place of \(s_1,s_2\ldots\) (the latter changes absolutely nothing), then one has
\[
\begin{aligned}
\Theta^{(q)}(-\omega)={}&
\sum_{s_1,s_2\ldots s_q=\pm1}(-1)^{\frac{s_1+s_2+\cdots+s_q}{2}}
\exp\!\left
\{\frac12\sum_i s_i(\omega_{p+i}+a_i)-\frac14\sum_{i>k}\sum s_i s_k b_{ik}\right\}\\
&\cdot \Theta\!\left(\omega_h-\frac{s_1u_{h1}+s_2u_{h2}+\cdots}{2}\right).
\end{aligned}
\]
The corresponding terms of \(\Theta^{(q)}(\omega)\) and \(\Theta^{(q)}(-\omega)\) therefore differ only by the coefficients
\[
(-1)^{-\frac{s_1+s_2+\cdots}{2}}e^{-\frac12\sum_i s_ia_i},\qquad
(-1)^{\frac{s_1+s_2+\cdots}{2}}e^{\frac12\sum_i s_ia_i},
\]
whose quotient has the value
\[
e^{-\sum_i s_i a_i}\cdot(-1)^q.
\]
Let us examine more closely the part standing in the exponent of \(e\). With insertion of the values of the \(a\), it becomes
\[
\begin{aligned}
-\frac12\sum_i s_i\Bigg[
&2\sum_{k=1}^{k=p}\int_{\beta^{(k)}}^{\gamma^{i,k}}du_{p+i}
+\sum_{k=1}^{k=i-1}\left(\int_{\beta^{(p+k)}}^{\varepsilon^{(k)}}du_{p+i}
+\int_{\beta^{(p+k)}}^{\varepsilon^{(-k)}}du_{p+i}\right)\\
&+2\int_{\beta^{(p+i)}}^{\mu}du_{p+i}
+\sum_{k=i+1}^{k=q}\left(\int_{\beta^{(p+k)}}^{\varepsilon^{(k)}}du_{p+i}
+\int_{\beta^{(p+k)}}^{\varepsilon^{(-k)}}du_{p+i}\right)
\Bigg].
\end{aligned}
\]
If in the bracketed part we complete it by integrals with coincident upper and lower limits, representing the points in which \(f=0\) is further cut by the tangent at \(\mu\), then the lower limits form the complete intersection system of the curve of order \((n-2)\) by which the \(\beta\) were defined, while the upper limits form the aggregate of the intersections of the tangent at \(\mu\) with those of the curve of order \((n-3)\) which defines the \(\gamma\). By Abel's theorem, therefore, the bracketed expression is equal to a periodicity modulus of the integral \(u_{p+i}\), increased by the logarithm of an algebraic function. The latter vanishes because the coordinates of the parameter points of \(u_{p+i}\) coincide, and yet the two curves do not pass through the relevant double point, so that the left sides of their equations assume equal and non-vanishing values for \(\varepsilon^{(i)}\) and \(\varepsilon^{(-i)}\). Thus only the periodicity modulus remains to be determined. It consists of a multiple of \(2\pi\sqrt{-1}\), belonging to the integral of the third kind \(u_{p+i}\) as such, and of a further part which would require corresponding terms also for the integrals of the first kind, and therefore vanishes with those terms. This second part can therefore be investigated for the integrals of the first kind, namely for the integrals
\[
\begin{aligned}
&2\sum_{k=1}^{k=p}\int_{\beta^{(k)}}^{\gamma^{i,k}}du_h
+\sum_{k=1}^{k=i-1}\left(\int_{\beta^{(p+k)}}^{\varepsilon^{(k)}}du_h
+\int_{\beta^{(p+k)}}^{\varepsilon^{(-k)}}du_h\right)
+2\int_{\beta^{(p+i)}}^{\mu}du_h\\
&\hspace{3em}
+\sum_{k=i+1}^{k=q}\left(\int_{\beta^{(p+k)}}^{\varepsilon^{(k)}}du_h
+\int_{\beta^{(p+k)}}^{\varepsilon^{(-k)}}du_h\right),
\qquad (h=1,2\ldots p).
\end{aligned}
\]
Comparing now equations \(\S\,80\) (1) and (3), which refer to integrals of the first kind and to the same paths of integration, one immediately finds that this sum vanishes. Hence it follows that the periodicity moduli of the integrals of the third kind can consist only of a term \(2\varrho_i\pi\sqrt{-1}\), and that the quotient under investigation therefore has the value
\[
(-1)^{q+\sum\varrho_i s_i},
\]
or, since all \(s_i\) are equal to \(\pm1\),
\[
(-1)^{q+\sum\varrho_i},
\]
so that finally
\[
\Theta^{(q)}(-\omega)=(-1)^{q+\sum\varrho_i}\cdot\Theta^{(q)}(\omega).
\]
As for the numbers \(\varrho\), in them there is a real arbitrariness. Namely, suppose the system of the \(\beta\) determined in the way indicated on p. 281. Then one can find other systems \(\beta'\), other contact curves, related to the original ones in such a way that
\[
\sum_{k=1}^{k=p+q}\int_{\beta^{(k)}}^{\beta'^{(k)}}du_h=0,
\qquad (h=1,2\ldots p),
\]
\[
\sum_{k=1}^{k=p+q}\int_{\beta^{(k)}}^{\beta'^{(k)}}du_{p+h}=\sigma_h\pi\sqrt{-1},
\qquad (h=1,2\ldots q),
\]
and which, by Abel's theorem, prove to be contact points of such curves, as one sees at once by multiplying the preceding equations by 2. If one now introduces the \(\beta'\) instead of the \(\beta\), then \(\omega_1,\omega_2\ldots\omega_p\), that is, the arguments of the \(\Theta\)-functions occurring in \(\Theta^{(q)}\), do not change; but each of the quantities \(\omega_{p+i},a_{i,s_1,s_2\ldots}\) changes by \(\sigma_i\pi\sqrt{-1}\). One sees that the function \(\Theta^{(q)}\) is thereby changed only insofar as the variables \(\omega_{p+i}\) acquire a changed meaning. But at the same time the quotient investigated above changes by
\[
-\sum s_i\sigma_i,
\]
or by \(\sum\sigma_i\). Henceforth we shall always suppose that instead of the \(\beta\) that system of the \(\beta'\) is put for which \(\sigma_i=\varrho_i\), \((i=1,2\ldots q)\). Then one has
\[
\Theta^{(q)}(-\omega)=(-1)^q\Theta^{(q)}(\omega),
\]
and the function \(\Theta^{(q)}\) so defined is therefore an even or odd function of its arguments, according as \(q\) is even or odd.

For this function so defined one always has
\[
a_i=0,
\qquad (i=1,2\ldots q),
\]
and hence finally for \(\Theta^{(q)}\) the expression
\[
\begin{aligned}
\Theta^{(q)}(\omega)={}&
\sum_{s_1s_2\ldots s_q=\pm1}
\exp\!\left\{\frac12\sum_i s_i\omega_{p+i}
-\frac14\sum_{i>k}\sum s_i s_k b_{ik}
-\frac{\pi\sqrt{-1}}{2}\sum_i s_i\right\}\\
&\cdot\Theta\!\left(\omega_h-\frac{s_1u_{h1}+s_2u_{h2}+\cdots}{2}\right),
\end{aligned}
\]
while
\[
\omega_h=\sum_{i=1}^{i=p+q}\int_{c^{(i)}}^{x^{(i)}}du_h-\int_{\mu}^{\xi}du_h,
\qquad (h=1,2\ldots p+q).
\]

\subsection*{\S\,82. The division for the extended inversion problem.}

We shall not enter into the extensions of those theorems which constitute the contents of the ninth section, and which are easily extended to the problems of the extended inversion problem and to the function \(\Theta^{(q)}\). On the other hand, it seems appropriate to devote a few words to the \emph{division}, specifically to the special division, since the extension of the general division problem is solved exactly as the general division treated in the ninth section. The problem of the general division will here be to determine the \(x\) from the equations
\[
m\sum_{i=1}^{i=p+q}\int_{c^{(i)}}^{x^{(i)}}du_h
+\sum_{i=1}^{i=r}\int_{a^{(i)}}^{y^{(i)}}du_h=0,
\qquad (h=1,2\ldots p+q),
\]
or, through given points \(y\), to lay a curve which passes through the separated \(\delta+r\) exceptional points, has \(m\)-point contact in \(p+q\) points, and has its remaining intersections with \(f=0\) in common with a similar contact curve, touching at the points \(c\) and passing through the points \(a\). The problem admits \(m^{2p+q}\) solutions; the exponent of \(m\) gives the number of mutually independent systems of periodicity moduli of the extended inversion problem. The roots of the algebraic equation corresponding to those solutions are found, as in \(\S\,70\), algebraically, assuming the roots of a special division equation.

The latter, arising from the equation
\[
(1)\quad m\sum_{i=1}^{i=p+q}\int_{c^{(i)}}^{x^{(i)}}du_h=P_h,
\qquad (h=1,2\ldots p+q),
\]
has, besides the self-evident solution \(x^{(1)}=c^{(1)}\), \(x^{(2)}=c^{(2)}\), etc., a further \(m^{2p+q}-1\). We shall reduce the corresponding algebraic equation, that is, the special division equation of the extended inversion problem, to the corresponding one of Jacobi's problem.

Let a problem of special division for Jacobi's inversion problem be given by the equation
\[
(2)\quad m\sum_{i=1}^{i=p}\int_{\gamma^{(i)}}^{z^{(i)}}du_h=P_h,
\qquad (h=1,2\ldots p).
\]
Now compare two solutions of problems (1) and (2) with one another, for which respectively
\[
(3)\quad
\left\{\begin{aligned}
\sum_{i=1}^{i=p+q}\int_{c^{(i)}}^{x^{(i)}}du_h
&=\frac{2\lambda_h\pi\sqrt{-1}}{m}
+\frac1m(\varkappa_1a_{1h}+\varkappa_2a_{2h}+\cdots),
&& (h=1,2\ldots p),\\
\sum_{i=1}^{i=p+q}\int_{c^{(i)}}^{x^{(i)}}du_{p+h}
&=\frac{2\nu_h\pi\sqrt{-1}}{m}
+\frac1m(\varkappa_1u_{1h}+\varkappa_2u_{2h}+\cdots),
&& (h=1,2\ldots q),
\end{aligned}\right.
\]
and
\[
(4)\quad
\sum_{i=1}^{i=p}\int_{\gamma^{(i)}}^{z^{(i)}}du_h
=\frac{2\lambda_h\pi\sqrt{-1}}{m}
+\frac1m(\varkappa_1a_{1h}+\varkappa_2a_{2h}+\cdots),
\qquad (h=1,2\ldots p).
\]
From the combination of the first equation (3) with (4) follows
\[
(5)\quad
\sum_{i=1}^{i=p+q}\int_{c^{(i)}}^{x^{(i)}}du_h
+\sum_{i=1}^{i=p}\int_{z^{(i)}}^{\gamma^{(i)}}du_h=0,
\qquad (h=1,2\ldots p).
\]
If the points \(z\), or at least the symmetric functions of their coordinates, that is, a root of the special division equation for Jacobi's inversion problem, are therefore assumed known, then one can lay a curve \(\psi=0\) through them, through the \(c\), and through all double and cusp points; this curve cuts in further points \(b\). Through these, as well as through the \(\gamma\) and the double and cusp points, there then passes a second curve \(\varphi=0\) of the same order as the preceding one, whose remaining intersections with \(f=0\) are the \(x\). But this curve, or even its intersection system, is not yet completely determined thereby; rather, initially \(q\) of the \(x\), or as many coefficients of \(\varphi=0\) not destroyed by \(f=0\), remain undetermined.

Now first suppose the parameter points of the integrals of the third kind \(u_{p+h}\) to be slightly different from the \(q\) preferred double points, and located near each double point on the different branches of \(f=0\) in \(\varepsilon^{(h)}\), \(\varepsilon^{(-h)}\). Then, because of equation (5), by Abel's theorem one also has
\[
(6)\quad
\sum_{i=1}^{i=p+q}\int_{c^{(i)}}^{x^{(i)}}du_{p+h}
+\sum_{i=1}^{i=p}\int_{z^{(i)}}^{\gamma^{(i)}}du_{p+h}
=-\log
\frac{\dfrac{\varphi(\varepsilon^{(h)})}{\psi(\varepsilon^{(h)})}}
{\dfrac{\varphi(\varepsilon^{(-h)})}{\psi(\varepsilon^{(-h)})}}.
\]
If now the points \(\varepsilon^{(h)}\) and \(\varepsilon^{(-h)}\) approach the corresponding double point more and more, then at the same time \(\varphi(\varepsilon^{(h)})\), \(\psi(\varepsilon^{(h)})\), \(\varphi(\varepsilon^{(-h)})\), \(\psi(\varepsilon^{(-h)})\) vanish. Thus, if
\[
\begin{array}{lll}
\varepsilon_1^{(h)}=\vartheta_1^{(h)}+\delta\xi_1^{(h)},&
\varepsilon_2^{(h)}=\vartheta_2^{(h)}+\delta\xi_2^{(h)},&
\varepsilon_3^{(h)}=\vartheta_3^{(h)}+\delta\xi_3^{(h)},\\
\varepsilon_1^{(-h)}=\vartheta_1^{(h)}+\delta\eta_1^{(h)},&
\varepsilon_2^{(-h)}=\vartheta_2^{(h)}+\delta\eta_2^{(h)},&
\varepsilon_3^{(-h)}=\vartheta_3^{(h)}+\delta\eta_3^{(h)},
\end{array}
\]
where the \(\delta\xi\), \(\delta\eta\) denote the infinitely small increments of the coordinates of a point moving from the double point on one or the other branch of \(f=0\), one has
\[
\lim
\frac{\dfrac{\varphi(\varepsilon^{(h)})}{\psi(\varepsilon^{(h)})}}
{\dfrac{\varphi(\varepsilon^{(-h)})}{\psi(\varepsilon^{(-h)})}}
=
\frac{\dfrac{\varphi'(\vartheta_1^{(h)})\delta\xi_1^{(h)}+\varphi'(\vartheta_2^{(h)})\delta\xi_2^{(h)}+\varphi'(\vartheta_3^{(h)})\delta\xi_3^{(h)}}
{\psi'(\vartheta_1^{(h)})\delta\xi_1^{(h)}+\psi'(\vartheta_2^{(h)})\delta\xi_2^{(h)}+\psi'(\vartheta_3^{(h)})\delta\xi_3^{(h)}}}
{\dfrac{\varphi'(\vartheta_1^{(h)})\delta\eta_1^{(h)}+\varphi'(\vartheta_2^{(h)})\delta\eta_2^{(h)}+\varphi'(\vartheta_3^{(h)})\delta\eta_3^{(h)}}
{\psi'(\vartheta_1^{(h)})\delta\eta_1^{(h)}+\psi'(\vartheta_2^{(h)})\delta\eta_2^{(h)}+\psi'(\vartheta_3^{(h)})\delta\eta_3^{(h)}}}.
\]
Of the three \(d\xi\) as well as of the three \(d\eta\), one may set one arbitrarily equal to zero; the ratio of the others is then found from the quadratic equation furnished by the lowest terms of \(f=0\) near the double point:
\[
\sum_i\sum_k
\frac{\partial^2f(\vartheta^{(h)})}{\partial\vartheta_i^{(h)}\partial\vartheta_k^{(h)}}
\,\delta\xi_i^{(h)}\delta\xi_k^{(h)}=0,
\]
and this equation must also be satisfied if the \(d\xi\) are replaced by the \(d\eta\).

Furthermore, note that equation (4) can be regarded as a form of Abel's theorem arising from the equation
\[
m\sum_{i=1}^{i=p}\int_{\gamma^{(i)}}^{z^{(i)}}du_h
=\lambda\int_t^{\delta}du_h+\varkappa_1\int_{t^{(1)}}^{\delta^{(1)}}du_h
+\varkappa_2\int_{t^{(2)}}^{\delta^{(2)}}du_h+\cdots+\varkappa_p\int_{t^{(p)}}^{\delta^{(p)}}du_h,
\qquad (h=1,2\ldots p),
\]
by letting the variable points \(t\) run through period paths until they coincide with the points \(\delta\), thereby producing the periodicity moduli on the right side of (5). Thus lay a curve \(\Phi=0\), passing through double points and cusps, which touches with multiplicity \(m\) at the points \(\gamma\), and at the points \(\delta,\delta^{(1)},\delta^{(2)}\ldots\) touches with multiplicities \(\lambda\), \(\varkappa_1\), etc.; and lay through its remaining intersections with \(f=0\), as well as through the double points and cusps, another curve \(\Psi=0\) of the same order, touching at the points \(t,t^{(1)},t^{(2)}\ldots\) with multiplicities \(\lambda\), \(\varkappa_1\), etc. If one further thinks of the second curve as arising from the first by continuous motion in such a way that it never loses its contact properties, then the second curve will touch at the points \(z\), and these will form the corresponding solution of the special division problem of Jacobi's inversion problem when the \(t\) have returned to the \(\delta\) along the indicated period paths. Thus for any integral of the third kind one first has
\[
m\sum_{i=1}^{i=p}\int_{\gamma^{(i)}}^{z^{(i)}}d\Pi_{\xi\eta}
+\lambda\int_{\delta}^{t}d\Pi_{\xi\eta}
+\varkappa_1\int_{\delta^{(1)}}^{t^{(1)}}d\Pi_{\xi\eta}+\cdots
=-\log\frac{\dfrac{\Phi(\xi)}{\Psi(\xi)}}{\dfrac{\Phi(\eta)}{\Psi(\eta)}}.
\]
But when the \(t\) have traversed their period paths, and when at the same time the integral \(\Pi_{\xi\eta}\) passes into the integral \(u_{p+h}\), this equation passes, similarly to (5), into the form
\[
\begin{aligned}
\sum_{i=1}^{i=p}\int_{\gamma^{(i)}}^{z^{(i)}}du_{p+h}
={}&\frac1m(\varkappa_1u_{1h}+\varkappa_2u_{2h}+\cdots)\\
&-\frac1m\log
\frac{\dfrac{\Phi'(\vartheta_1^{(h)})\delta\xi_1^{(h)}+\Phi'(\vartheta_2^{(h)})\delta\xi_2^{(h)}+\cdots}
{\Psi'(\vartheta_1^{(h)})\delta\xi_1^{(h)}+\Psi'(\vartheta_2^{(h)})\delta\xi_2^{(h)}+\cdots}}
{\dfrac{\Phi'(\vartheta_1^{(h)})\delta\eta_1^{(h)}+\Phi'(\vartheta_2^{(h)})\delta\eta_2^{(h)}+\cdots}
{\Psi'(\vartheta_1^{(h)})\delta\eta_1^{(h)}+\Psi'(\vartheta_2^{(h)})\delta\eta_2^{(h)}+\cdots}},\\
&\hspace{13em}(h=1,2\ldots q).
\end{aligned}
\]
The combination of this equation with the second equation (3) and with (6) now gives at once that the expression
\[
\log\frac{\dfrac{\varphi'(\vartheta_1^{(h)})\delta\xi_1^{(h)}+\cdots}
{\psi'(\vartheta_1^{(h)})\delta\xi_1^{(h)}+\cdots}}
{\dfrac{\varphi'(\vartheta_1^{(h)})\delta\eta_1^{(h)}+\cdots}
{\psi'(\vartheta_1^{(h)})\delta\eta_1^{(h)}+\cdots}}
-\frac1m
\log\frac{\dfrac{\Phi'(\vartheta_1^{(h)})\delta\xi_1^{(h)}+\cdots}
{\Psi'(\vartheta_1^{(h)})\delta\xi_1^{(h)}+\cdots}}
{\dfrac{\Phi'(\vartheta_1^{(h)})\delta\eta_1^{(h)}+\cdots}
{\Psi'(\vartheta_1^{(h)})\delta\eta_1^{(h)}+\cdots}}
\]
is equal to the logarithm of an \(m\)-th root of unity, or that
\[
(7)\quad
\frac{\dfrac{\varphi'(\vartheta_1^{(h)})\delta\xi_1^{(h)}+\cdots}
{\psi'(\vartheta_1^{(h)})\delta\xi_1^{(h)}+\cdots}}
{\dfrac{\varphi'(\vartheta_1^{(h)})\delta\eta_1^{(h)}+\cdots}
{\psi'(\vartheta_1^{(h)})\delta\eta_1^{(h)}+\cdots}}
=
\sqrt[m]{
\frac{\dfrac{\Phi'(\vartheta_1^{(h)})\delta\xi_1^{(h)}+\cdots}
{\Psi'(\vartheta_1^{(h)})\delta\xi_1^{(h)}+\cdots}}
{\dfrac{\Phi'(\vartheta_1^{(h)})\delta\eta_1^{(h)}+\cdots}
{\Psi'(\vartheta_1^{(h)})\delta\eta_1^{(h)}+\cdots}}
},
\qquad (h=1,2\ldots q).
\]
In this equation, if the relevant solution of the special division problem for Jacobi's inversion problem is known, everything is given except the coefficients of the function \(\varphi\); and hence the present formula gives \(q\) further linear equations which serve to determine these coefficients completely. They are obtained expressed by rational functions and \(m\)-th roots of rational functions of a root of the assumed known division problem, and one can therefore also express in this way a root of the extended inversion problem through it. Thus to each root of the Jacobi division problem there belong \(m^q\) roots of the extended one, which are obtained by giving independently to the \(q\) different \(m\)-th roots all possible values. The \(m^q\) curves of such a group stand in the relation mentioned above to one another. If through all double points and cusps, through the \(c\) and the \(\gamma\), one lays a curve \(\psi=0\), then every system of the \(x\) belonging to a curve of the group lies on a curve \(\varphi=0\) which is of the same order as \(\psi=0\), with the double points and cusps, the points \(z\), and the remaining intersections of \(\psi=0\) with \(f=0\).

Among the groups, however, there is one associated with the identical root of the Jacobi division equation, for which the \(z\) coincide with the \(\gamma\). For this group the function from which in (7) the \(m\)-th root was drawn becomes equal to unity; and one of the curves \(\varphi\) coincides with \(\psi\), namely that one for which the \(m\)-th root itself is set equal to 1. One observes that the remaining \(m^q-1\) curves of this group are found without adding the division equation for the Jacobi inversion problem, so that they require no irrationalities other than the \(m\)-th roots of unity.

\subsection*{\S\,83. Division for \(m=2\).}

As for the special division of the extended inversion problem for \(m=2\), we note that from the theorem for \(\Theta^{(q)}\) given on p. 279 the periodicity formula follows:
\[
\Theta^{(q)}(\omega+P)=
\exp\!\left\{\sum q_i\omega_i+\frac12\sum\sum q_hq_ka_{hk}
+\pi\sqrt{-1}\sum n_i\right\}\Theta^{(q)}(\omega),
\]
provided that
\[
\begin{array}{ll}
P_h=2m_h\pi\sqrt{-1}+q_1a_{1h}+q_2a_{2h}+\cdots,&(h=1,2\ldots p),\\
P_{p+h}=2n_h\pi\sqrt{-1}+q_1u_{1h}+q_2u_{2h}+\cdots,&(h=1,2\ldots q).
\end{array}
\]
Now put \(\omega_h=-\frac{P_h}{2}\). Then we have
\[
\Theta^{(q)}\!\left(\frac P2\right)
=e^{\pi\sqrt{-1}\left(\sum n_i-\sum m_iq_i\right)}
\Theta^{(q)}\!\left(-\frac P2\right),
\]
or, taking into account the equation
\[
\Theta^{(q)}(-\omega)=(-1)^q\Theta^{(q)}(\omega),
\]
also
\[
0=\Theta^{(q)}\!\left(\frac P2\right)
\left(1-(-1)^{\sum m_iq_i+\sum n_i+q}\right).
\]
Thus the expression \(\Theta^{(q)}\!\left(\frac P2\right)\) always vanishes when
\[
(1)\quad \sum m_iq_i+\sum n_i+q=2H+1,
\]
and one can show as in \(\S\,75\) that in every other case it generally does not vanish. In formula (1) all \(m_i,q_i,n_i\) can evidently be chosen arbitrarily, except for one of the latter, which is determined by the others. Thus one sees that among the expressions \(\Theta^{(q)}\!\left(\frac P2\right)\) a number \(2^{2p+q-1}\) vanish, and that an equal number do not vanish; hence also among the \(2^{2p+q}\) functions \(\Theta^{(q)}\!\left(\omega+\frac P2\right)\) just as many are even as odd.

As in the tenth section, one shows that the non-vanishing functions \(\Theta^{(q)}\!\left(\frac P2\right)\) correspond to the actual contact curves which pass through the separated \(\delta+r\) double points and cusps and through the \(n-2\) remaining intersections of the tangent drawn at the point \(\mu\), and in addition touch the given curve in \(p+q\) points. The other functions \(\Theta^{(q)}\!\left(\frac P2\right)\), by contrast, correspond to such contact curves as decompose into the tangent at \(\mu\) itself and curves of order \((n-3)\). One further concludes that there are \(2^{2p+q-1}\) curves of order \((n-3)\) passing through \(\delta+r=d+r-q\) double points and cusps of the curve and touching the curve \(f=0\) in \(p+q\) points. To find these curves, however, according to the preceding, no equation other than the one giving the curves of order \((n-3)\) laid through all double points and cusps and touching in \(p\) points is required.

\subsection*{\S\,84. Influence of cusps.}

It remains finally to mention the modifications undergone by the preceding investigations when some of the \(q\) preferred double points pass over into cusps. In the preceding put \(q+r\) in place of \(r\), and suppose that in the equation \(f=0\) there are \(r\) variable parameters which independently approach the limit zero, while at the same time the discontinuity points of the last \(r\) integrals \(u_{p+q+1},u_{p+q+2},\ldots,u_{p+q+r}\) pass from double points to cusps. Let the \(r\) variable parameters be \(\tau_1,\tau_2\ldots\tau_r\); geometrically we may think of them as defined by placing pencils of rays at the corresponding double points of the curve \(f=0\), and denoting by \(\tau_i\) the difference of the parameters of those rays in each pencil which are the tangents of the double point. While the \(\tau\) approach zero, and hence the double points approach cusps, those last \(r\) integrals become infinitely small through the approach of the parameter points, and one may put
\[
u_{p+q+i}=\tau_iU_i,
\qquad (i=1,2\ldots r),
\]
where the \(U\) are now finite integrals of the second kind. In the equations of the inversion problem one must then replace the last \(r\) equations
\[
\sum_{h=1}^{h=p+q+r}\int_{c^{(h)}}^{x^{(h)}}du_{p+q+i}=v_{p+q+i},
\qquad (i=1,2\ldots r)
\]
by
\[
(1)\quad \sum_{h=1}^{h=p+q+r}\int_{c^{(h)}}^{x^{(h)}}dU_i=V_i,
\]
where
\[
v_{p+q+i}=\tau_iV_i,
\]
and the \(V_i\) now denote finite variables.

The integrals (1), as integrals of the second kind, have no periodicity modulus \(2\pi\sqrt{-1}\) any longer; however \(V_i\) may pass into
\[
V_i'=V_i+U_{1i}q_1+U_{2i}q_2+\cdots+U_{pi}q_p,
\qquad (i=1,2\ldots r),
\]
while the integral sums of the first and third kind transform into
\[
\begin{array}{ll}
v_i'=v_i+2m_i\pi\sqrt{-1}+a_{1i}q_1+a_{2i}q_2+\cdots+a_{pi}q_p,&(i=1,2\ldots p),\\
v_{p+i}'=v_{p+i}+2n_i\pi\sqrt{-1}+u_{1i}q_1+u_{2i}q_2+\cdots+u_{pi}q_p,&(i=1,2\ldots q).
\end{array}
\]
The transcendental functions \(U^{(q)}\), \(V^{(q)}\) correspond in this limiting transition to other transcendental functions \(U^{(q,r)}\), \(V^{(q,r)}\), whose form can be read off at once from \(\S\,77\), and which behave, with respect to their continuity and their vanishing, quite similarly to those functions.

It is different with the function \(\Theta^{(q,r)}\), which arises from \(\Theta^{(q+r)}\) by the limiting transition. It is easy to see that the function \(\Theta^{(q)}\) vanishes as soon as any one of the \(q\) double points occurring in it is transformed into a cusp. If, for example, the \(q\)-th double point passes into such a point, the points \(\varepsilon^{(q)},\varepsilon^{(-q)}\), previously to be regarded as separate, coalesce into two infinitely near points of the curve; \(\omega_{p+q}\) vanishes together with the \(u_{hq}\) and the \(a_q,b_{hq}\); in \(\Theta^{(q)}\), the terms multiplied by
\[
(-1)^{\frac{-s_1-s_2-\cdots-1}{2}}
\quad\text{and}\quad
(-1)^{\frac{-s_1-s_2-\cdots+1}{2}}
\]
become equal to one another, and therefore destroy each other, since these powers of \(\sqrt{-1}\) always have opposite values.

It follows that when \(r\) double points of \(\Theta^{(q+r)}\) pass into cusps, this function approaches the limit
\[
\lim \Theta^{(q+r)}=\tau_1\tau_2\ldots\tau_r\Theta^{(q,r)},
\]
where only now \(\Theta^{(q,r)}\) is a finite function. Since only quotients of the \(\Theta^{(q)}\) occur throughout, one obtains the formulas of the new inversion problem by using this function everywhere in place of \(\Theta^{(q+r)}\), a function which, according to the preceding, must be synectic and must retain the property of vanishing only under the circumstances specified for \(\Theta^{(q)}\).

As for the actual formation of the function \(\Theta^{(q,r)}\), set
\[
\begin{array}{lll}
\omega_{p+q+i}=\tau_iW_i,&(i=1,2\ldots r),\\
b_{h,q+i}=\tau_iB_{hi},&(i=1,2\ldots r,\ h=1,2\ldots q),\\
b_{q+k,q+i}=\tau_i\tau_kC_{ki},&(i,k=1,2\ldots r),\\
u_{h,q+i}=\tau_iU_{hi},&(i=1,2\ldots r,\ h=1,2\ldots p),
\end{array}
\]
and then one finds
\[
\begin{aligned}
&\Theta^{(q,r)}(w_1,w_2\ldots w_{p+q},W_1,W_2\ldots W_r)\\
&=\frac{\partial^q}{\partial\tau_1\partial\tau_2\cdots}
\Bigg\{\sum_{s_1s_2\ldots s_r=\pm1}
\exp\!\left(\frac12\sum_{i=1}^{i=r}s_i\tau_iW_{p+i}
-\frac{\pi\sqrt{-1}}2\sum_{i=1}^{i=r}s_i
-\frac14\sum_{i\ge k}\sum s_is_k\tau_i\tau_kC_{ik}\right)\\
&\qquad\cdot
\Theta^{(q)}\!\left(\omega_h-\frac{\tau_1s_1U_{h1}+\tau_2s_2U_{h2}+\cdots}{2},\
\omega_{p+h}-\frac{\tau_1s_1B_{h1}+\tau_2s_2B_{h2}+\cdots}{2}\right)\Bigg\},
\end{aligned}
\]
where on the right in \(\Theta^{(q)}\) only two arguments are given instead of \(p+q\), and where, after carrying out the differentiation, all \(\tau\) are to be set equal to zero.

In the transition from double points to cusps the problem of division is modified by the fact that integrals of the second kind no longer have any peculiar periodicity moduli. Hence the problem of general division for the present inversion problem with \(p+q+r\) variables has only \(m^{2p+q}\) solutions, and that of special division has \(m^{2p+q}-1\).

On the other hand, both the property of being an even or odd function, according as \(q+r\) is even or odd, is retained by \(\Theta^{(q,r)}\), and also the property of acquiring, by increasing the variables by a system \(P_h\) of periodicity moduli, the factor
\[
\exp\!\left\{\sum_{i=1}^{i=p}q_i\omega_i
+\frac12\sum_{h,k=1}^{h,k=p}q_hq_ka_{hk}
+\pi\sqrt{-1}\sum_{i=1}^{i=q}n_i\right\}.
\]
Therefore, as in \(\S\,83\), one also has \(\Theta^{(q,r)}\left(\frac P2\right)=0\) as soon as
\[
\sum_{i=1}^{i=p}m_iq_i+\sum_{i=1}^{i=q}n_i+q+r=2H+1.
\]
Thus here as there, there are \(2^{2p+q-1}\) proper contact curves of order \((n-2)\), which pass through the exceptional points present besides the \(q+r\) preferred ones and through the intersections of the tangent of \(\mu\), and which moreover touch in \(p+q+r\) points. Equally large is the number of curves of order \((n-3)\) which pass through those exceptional points and touch in \(p+q+r-1\) points. The only exception is the case \(q=0\), that is, the case in which all preferred double points become cusps. In this case one can no longer choose the \(m_i,q_i\) arbitrarily; instead they must satisfy the equation
\[
\sum_{i=1}^{i=p}m_iq_i+r=2H+1,
\]
so that the corresponding system of half periodicity moduli leads to a contact curve of order \((n-3)\). Hence in this case, if \(r\) is even, there are \(2^{p-1}(2^p-1)\) contact curves of order \((n-3)\) which pass through the double points and cusps, except for \(r\) of the latter, and touch in \(p+r-1\) points; but if \(r\) is odd there are \(2^{p-1}(2^p+1)\) such curves. The exact reverse holds for the proper contact curves of order \((n-2)\) which pass through the double points and cusps just indicated, as well as through the \(n-2\) intersections of the tangent of a point \(\mu\), and which touch in \(p+r\) points.


% --- Source pages 322--331; printed pages 300--309. ---

\section*{Twelfth Section.\\The infinitely many forms of the function \(\Theta\).}

\subsection*{\S\,85. Canonical determinants.}

In the fourth section it was shown how one can arrive at a system of normal periodicity moduli. This was done by means of the linear integral relations
\[
\begin{aligned}
(1)\quad
M^{(h)}&=\alpha_1^{(h)}(a_1)+\alpha_2^{(h)}(a_2)+\cdots+\alpha_{2p}^{(h)}(a_{2p}),\\
N^{(h)}&=\alpha_1^{(h)}(b_1)+\alpha_2^{(h)}(b_2)+\cdots+\alpha_{2p}^{(h)}(b_{2p})
\qquad (h=1,2\ldots2p)
\end{aligned}
\]
which produced the identical transformation equation
\[
(2)\quad
\sum\sum c_{ik}(a_i)(b_k)
=\sum_{h=1}^{h=p}\left(M^{(2h-1)}N^{(2h)}-N^{(2h-1)}M^{(2h)}\right).
\]
Here the \(c\) denoted integers, \(c_{ik}=-c_{ki}\), and the determinant of the \(c\) was equal to 1. If then the \((a),(b)\) are the periodicity moduli of any two integrals of the first kind, the expressions (1) are normal periodicity moduli of them, as was necessary for setting up the theory carried out above.

The transformation formula (2) can be fulfilled in infinitely many ways. Suppose that it has been fulfilled in one way; thus a system of the \(\alpha,M,N\) is given. We seek the conditions which all other solutions of the transformation problem must satisfy.

Thus set, in order to obtain a new system of normal periodicity moduli,
\[
\begin{aligned}
(3)\quad
R_k&=\beta_{1,k}M^{(1)}+\beta_{2,k}M^{(2)}+\cdots+\beta_{2p,k}M^{(2p)},\\
S_k&=\beta_{1,k}N^{(1)}+\beta_{2,k}N^{(2)}+\cdots+\beta_{2p,k}N^{(2p)}
\qquad (k=1,2\ldots2p).
\end{aligned}
\]
The integers \(\beta\) are then to be determined so that identically
\[
(4)\quad
\sum_{k=1}^{k=p}(R_{2k-1}S_{2k}-S_{2k-1}R_{2k})
=\sum_{h=1}^{h=p}\left(M^{(2h-1)}N^{(2h)}-N^{(2h-1)}M^{(2h)}\right),
\]
or, when on the left the quantities (3) are replaced by their values, so that
\[
\begin{aligned}
&\sum_{k=1}^{k=p}\sum_{s,r=1}^{s,r=2p}
(\beta_{s,2k-1}\beta_{r,2k}-\beta_{r,2k-1}\beta_{s,2k})M^{(s)}N^{(r)}\\
&\hspace{2em}=\sum_{h=1}^{h=p}\left(M^{(2h-1)}N^{(2h)}-N^{(2h-1)}M^{(2h)}\right).
\end{aligned}
\]
Comparison of the coefficients therefore shows that the integers \(\beta\) must satisfy the following conditions, in which, in order not to obtain any twice, we may suppose \(s<r\):
\[
(5)\quad
\sum_{k=1}^{k=p}\left(\beta_{s,2k-1}\beta_{r,2k}-\beta_{r,2k-1}\beta_{s,2k}\right)=0,
\]
except when \(s\) is an odd number and \(r\) the following even number, in which case
\[
(6)\quad
\sum_{k=1}^{k=p}\left(\beta_{2h-1,2k-1}\beta_{2h,2k}-\beta_{2h,2k-1}\beta_{2h-1,2k}\right)=1.
\]
Determinants of the kind composed here from the integers \(\beta\)\footnote{Such determinants were first studied by Hermite in the Comptes Rendus of 1855, namely in connection with general transformations of the function \(\Theta\) for \(p=2\).} will be called \emph{canonical}. Thus, by defining as elements of a canonical determinant such integers as satisfy equations (5), (6), we can state the theorem contained in this definition:

\begin{center}
\textbf{The transition from one system of normal periodicity moduli to another is effected with the aid of a canonical transformation determinant, and every such determinant gives one of the required transitions.}
\end{center}

The canonical determinant has the value \(\pm1\); for its square, multiplied by the determinant
\[
\left|
\begin{array}{rrrrrrr}
0&1&0&0&\cdots&0&0\\
-1&0&0&0&\cdots&0&0\\
0&0&0&1&\cdots&0&0\\
0&0&-1&0&\cdots&0&0\\
\multicolumn{7}{c}{\cdots}\\
0&0&0&0&\cdots&0&1\\
0&0&0&0&\cdots&-1&0
\end{array}
\right|=1
\]
of the bilinear function which forms the right-hand side of (4), gives the determinant, wholly identical with the preceding one, of the bilinear function on the left-hand side of (4); thus the square of the canonical determinant is 1. Denote the value of the canonical determinant itself by \(\varepsilon=\pm1\), and by \(\gamma_{hk}\) the minor corresponding to the element \(\beta_{hk}\). The solution of equations (3) then gives
\[
\begin{aligned}
(7)\quad
M^{(h)}&=\varepsilon[\gamma_{h1}R_1+\gamma_{h2}R_2+\cdots+\gamma_{h,2p}R_{2p}],\\
N^{(h)}&=\varepsilon[\gamma_{h1}S_1+\gamma_{h2}S_2+\cdots+\gamma_{h,2p}S_{2p}].
\end{aligned}
\]
If on the left-hand side of equation (4) one replaces the \(R\) by their values from (3), and on the right-hand side the \(N\) by their values from (7), one obtains the equation
\[
\begin{aligned}
&\sum_{k=1}^{k=p}\sum_{r=1}^{r=2p}(\beta_{r,2k-1}S_{2k}-\beta_{r,2k}S_{2k-1})M^{(r)}\\
&\hspace{2em}=\varepsilon\sum_{h=1}^{h=p}\sum_{s=1}^{s=2p}
(\gamma_{2h,s}M^{(2h-1)}-\gamma_{2h-1,s}M^{(2h)})S_s,
\end{aligned}
\]
and, by comparing the coefficients on both sides, one obtains the following values for the minors \(\gamma\):
\[
\begin{array}{rlrl}
(8)\quad
\gamma_{2h,2k}&=\varepsilon\beta_{2h-1,2k-1},&
\gamma_{2h,2k-1}&=-\varepsilon\beta_{2h-1,2k},\\
\gamma_{2h-1,2k}&=-\varepsilon\beta_{2h,2k-1},&
\gamma_{2h-1,2k-1}&=\varepsilon\beta_{2h,2k}.
\end{array}
\]
Therefore equations (7), the solutions of (4), can also be written in the form
\[
\begin{aligned}
(9)\quad
M^{(2h-1)}&=\beta_{2h,2}R_1-\beta_{2h,1}R_2+\beta_{2h,4}R_3-\beta_{2h,3}R_4+\cdots,\\
M^{(2h)}&=-\beta_{2h-1,2}R_1+\beta_{2h-1,1}R_2-\beta_{2h-1,4}R_3+\beta_{2h-1,3}R_4+\cdots,\\
N^{(2h-1)}&=\beta_{2h,2}S_1-\beta_{2h,1}S_2+\beta_{2h,4}S_3-\beta_{2h,3}S_4+\cdots,\\
N^{(2h)}&=-\beta_{2h-1,2}S_1+\beta_{2h-1,1}S_2-\beta_{2h-1,4}S_3+\beta_{2h-1,3}S_4+\cdots.
\end{aligned}
\]
By definition every transformation determinant connecting two systems of normal periodicity moduli is canonical; in particular, this is so for the determinant of equations (7) or (9). The same also holds for the determinant of the \(\gamma\) themselves; for since equations (5), (6) are fulfilled when the \(\beta\) are replaced by the \(\varepsilon\gamma\), they are also fulfilled when the \(\gamma\) alone are put in place of the \(\beta\).

If in equations (3) the \(M,N\) are expressed by any two other systems of normal periodicity moduli,
\[
\begin{aligned}
M^{(h)}&=\delta_{1,h}V_1+\delta_{2,h}V_2+\cdots+\delta_{2p,h}V_{2p},\\
N^{(h)}&=\delta_{1,h}W_1+\delta_{2,h}W_2+\cdots+\delta_{2p,h}W_{2p},
\end{aligned}
\]
then formulas (3) pass into the following:
\[
\begin{aligned}
R_k&=\vartheta_{1,k}V_1+\vartheta_{2,k}V_2+\cdots+\vartheta_{2p,k}V_{2p},\\
S_k&=\vartheta_{1,k}W_1+\vartheta_{2,k}W_2+\cdots+\vartheta_{2p,k}W_{2p},
\end{aligned}
\]
where
\[
\vartheta_{i,k}=\beta_{1,k}\delta_{i,1}+\beta_{2,k}\delta_{i,2}+\cdots+\beta_{2p,k}\delta_{i,2p}.
\]
The determinants of the \(\delta\) and of the \(\vartheta\) are canonical by definition. Hence one has the theorem:

\begin{center}
\textbf{By multiplying two canonical determinants}
\end{center}
\[
\left|\begin{array}{cccc}
\beta_{11}&\beta_{21}&\cdots&\beta_{2p,1}\\
\beta_{12}&\beta_{22}&\cdots&\beta_{2p,2}\\
\multicolumn{4}{c}{\cdots}\\
\beta_{1,2p}&\beta_{2,2p}&\cdots&\beta_{2p,2p}
\end{array}\right|,
\qquad
\left|\begin{array}{cccc}
\delta_{11}&\delta_{21}&\cdots&\delta_{2p,1}\\
\delta_{12}&\delta_{22}&\cdots&\delta_{2p,2}\\
\multicolumn{4}{c}{\cdots}\\
\delta_{1,2p}&\delta_{2,2p}&\cdots&\delta_{2p,2p}
\end{array}\right|,
\]
one obtains, if in one of them horizontal and vertical rows are interchanged, the determinant, likewise canonical,
\[
\left|\begin{array}{cccc}
\vartheta_{11}&\vartheta_{21}&\cdots&\vartheta_{2p,1}\\
\vartheta_{12}&\vartheta_{22}&\cdots&\vartheta_{2p,2}\\
\multicolumn{4}{c}{\cdots}\\
\vartheta_{1,2p}&\vartheta_{2,2p}&\cdots&\vartheta_{2p,2p}
\end{array}\right|,
\]
in which
\[
\vartheta_{ik}=\beta_{1,k}\delta_{i,1}+\beta_{2,k}\delta_{i,2}+\cdots+\beta_{2p,k}\delta_{i,2p}.
\]
Since, moreover, for the validity of this theorem it must be immaterial in which of the two determinants horizontal and vertical rows are interchanged, and since the two determinants thereby arising differ only by interchange of horizontal and vertical rows, the further theorem follows:

\begin{center}
\textbf{A canonical determinant does not lose this property when its horizontal rows are interchanged with the vertical rows.}
\end{center}

The above theorem on multiplication therefore remains correct even if the proviso about interchange of rows is omitted; but then the elements of the resulting determinant no longer give the coefficients of the substitution composed from the substitutions corresponding to the individual multiplied determinants.

The question of the most general canonical determinant, or of the most general system of normal periodicity moduli, can now be understood as that of finding the simplest canonical determinants from which all others can be derived by multiplication. The corresponding linear substitutions then lead from certain most original systems of normal periodicity moduli to all possible ones.

Instead, however, of composing the given determinant from determinants as simple as possible, one may, provided only that the adjoints of the simplest canonical determinants again belong to the simplest ones, proceed so as to reduce the given determinant to a simple one by multiplication with simple determinants. The last, multiplied by the adjoints of the others, is then equal to the given one, which is thereby resolved into simple factors. Thus there arise the two questions: first, which canonical determinants we shall regard as \emph{simple}; second, how, by multiplication with simple determinants, one reduces a given determinant itself to a simple one.

\subsection*{\S\,86. Reduction of canonical determinants to simple ones.}

It is enough to list the following as the \emph{simplest} canonical determinants:

I. The diagonal terms are 1, except for two, \(\beta_{2h-1,2h-1}\) and \(\beta_{2h,2h}\), which are zero. In their place \(\beta_{2h,2h-1}=\pm1\), \(\beta_{2h-1,2h}=\mp1\); all other elements are zero. Apart from the sign \(\pm\), there are \(p\) such determinants. For \(p=2\) they are
\[
\left|\begin{array}{rrrr}0&\pm1&0&0\\ \mp1&0&0&0\\ 0&0&1&0\\ 0&0&0&1\end{array}\right|,
\qquad
\left|\begin{array}{rrrr}1&0&0&0\\ 0&1&0&0\\ 0&0&0&\pm1\\ 0&0&\mp1&0\end{array}\right|.
\]

II. The diagonal terms are 1, one element \(\beta_{2h-1,2h}=\pm1\), and all other entries zero. Apart from the sign \(\pm\), there are again \(p\) such determinants. For \(p=2\):
\[
\left|\begin{array}{rrrr}1&0&0&0\\ \pm1&1&0&0\\ 0&0&1&0\\ 0&0&0&1\end{array}\right|,
\qquad
\left|\begin{array}{rrrr}1&0&0&0\\ 0&1&0&0\\ 0&0&1&0\\ 0&0&\pm1&1\end{array}\right|.
\]

III. The diagonal terms are 1, one element \(\beta_{2h-1,2k-1}\ (k\gtrless h)=\pm1\), \(\beta_{2k,2h}=\mp1\), and all other entries zero. Apart from the sign \(\pm\), there are \(p(p-1)\) such determinants. For \(p=2\):
\[
\left|\begin{array}{rrrr}1&0&0&0\\ 0&1&0&\mp1\\ \pm1&0&1&0\\ 0&0&0&1\end{array}\right|,
\qquad
\left|\begin{array}{rrrr}1&0&\pm1&0\\ 0&1&0&0\\ 0&0&1&0\\ 0&\mp1&0&1\end{array}\right|.
\]
That these determinants are all canonical follows at once by considering the conditions to be satisfied. But the second condition is also fulfilled, namely that the determinants formed from the minors are again of the same kind as the original ones. Indeed, when one forms the adjoint of any one of the above determinants, the original one always reappears, except that \(\pm\) is changed into \(\mp\), and \(\mp\) into \(\pm\).

We now examine the effect of multiplying a given canonical determinant
\[
\left|\begin{array}{cccc}
\beta_{11}&\beta_{21}&\cdots&\beta_{2p,1}\\
\beta_{12}&\beta_{22}&\cdots&\beta_{2p,2}\\
\multicolumn{4}{c}{\cdots}\\
\beta_{1,2p}&\beta_{2,2p}&\cdots&\beta_{2p,2p}
\end{array}\right|
\]
by one of the simplest canonical determinants listed under I, II, III. In order that the substitution corresponding to the resulting determinant be a combination of the two substitutions corresponding to the given determinant and the simple determinant, we interchange the horizontal rows with the vertical rows in the latter before multiplication.

The effect of multiplying the determinant of the \(\beta\) by a simple determinant I is evidently (cf. \(\S\,86\)) that two vertical rows, the \((2h-1)\)-st and the \((2h)\)-th, are interchanged, and that one of them changes sign in the process.

The effect of multiplication by a simple determinant II is no other than that the \((2h)\)-th vertical row is added to or subtracted from the preceding odd \((2h-1)\)-st row.

The effect, finally, of multiplication by a determinant III is that the \((2h-1)\)-st vertical row is increased or decreased by the \((2k-1)\)-st, and at the same time, conversely, the \((2k)\)-th is decreased or increased by the \((2h)\)-th.

Thus with the aid of multiplication by simple canonical determinants we can always perform the following operations on the given determinant:

I. Interchange an even vertical row with the preceding odd one, changing the sign of one of them.

II. Subtract an even vertical row from the preceding odd one, or add it to that row.

III. Subtract an odd vertical row from another odd row or add it to it, and at the same time add the following even row to the even row following the former, or subtract it from that row.

We now want to see whether, by applying these processes, we can transform the given determinant into a simple one, or into the identity substitution determinant, in which all terms not on the diagonal are replaced by zeros.

First we can add an allowed fourth process to the above three. Apply process I to the vertical rows \((2h-1),(2h)\). Then in the \((2h-1)\)-st position we may put the row \(-(2h)\), and in the \((2h)\)-th the row \((2h-1)\). With the aid of process II the row \(-(2h)\) now in the \((2h-1)\)-st position passes into the aggregate of the row \(\pm(2h-1)-(2h)\). Applying process I once more, we then have in the \((2h-1)\)-st position the unchanged row \((2h-1)\), but in the \((2h)\)-th position the difference of the rows \((2h)\mp(2h-1)\). Thus it has been found that the following process can also always be performed:

IV. Subtract an odd vertical row from the following even one, or add it to that row.

By alternating processes II and IV one can evidently destroy all even terms in some horizontal row, for example the first. When all even terms of this row are zero, the action of process III, so far as this horizontal row is concerned, reduces to adding and subtracting the remaining odd terms. Thus these can again be destroyed by one another, so that finally a single term remains, and this can be made the first term of the row. Since the determinant has value \(\pm1\), this term too is necessarily \(\pm1\).

The second horizontal row is now treated similarly, only with the first term of it always excluded. One destroys all the remaining odd terms by processes II, IV, and then the even terms by III up to the second, which then becomes a factor of the determinant and therefore has the value \(\pm1\). By this one finally destroys the first term according to II.

Equation (5) of \(\S\,85\) now teaches that if in two successive horizontal rows \((2h-1),(2h)\) all elements vanish except the diagonal terms, the same must occur for the corresponding vertical terms. Since the nature of the determinant is not changed by interchanging horizontal and vertical rows, it is enough to prove the converse statement. Put in (5) \(r=2h-1\) or \(r=2h\), and assume that all \(\beta_{2h-1,i}\), \(\beta_{2h,i}\) vanish except \(\beta_{2h-1,2h-1}\), \(\beta_{2h,2h}\); then equation (5) immediately gives \(\beta_{s,2h}=0\), \(\beta_{s,2h-1}=0\), as was to be proved.

The reduced determinant therefore contains in the first horizontal and vertical rows only the diagonal terms. Thus one may operate with the other rows of the determinant as though the first two rows were not present at all. By the same procedure one then transforms also the third and fourth horizontal and vertical rows into rows containing only diagonal terms, then the fifth and sixth, and so on, until in general only diagonal terms remain.

Thus every canonical determinant can be composed with the aid of determinants I, II, III; hence every canonical substitution can also be composed from the \(2p(p+1)\) simple substitutions corresponding to those \(2p(p+1)\) simple determinants\footnote{The idea of composing integral substitution determinants of value 1 from simple determinants in which only the digits 1 and 0 occur belongs, as far as we know, to Mr. Kronecker.}.

\subsection*{\S\,87. Transformation of integrals of the first and third kind when the system of normal periods is changed.}

We now wish to investigate how the normal integrals of the first and third kind take shape when an altered system of normal periodicity moduli is taken as basis. For this purpose we return first to the formulas of \(\S\,30\). In these, put in place of the \(H\) the normal integrals of the first kind \(u_1,u_2,u_3\), corresponding to a given period system, and denote by \(v_1,v_2,v_3\) the new normal integrals of the first kind whose system of normal periodicity moduli is connected with the former by the canonical determinant
\[
\left|\begin{array}{cccc}
\beta_{11}&\beta_{21}&\cdots&\beta_{2p,1}\\
\beta_{12}&\beta_{22}&\cdots&\beta_{2p,2}\\
\multicolumn{4}{c}{\cdots}\\
\beta_{1,2p}&\beta_{2,2p}&\cdots&\beta_{2p,2p}
\end{array}\right|.
\]
If one takes as the original system of normal periodicity moduli the following,
\[
\begin{array}{cccccc}
2\pi\sqrt{-1}&a_{11}&0&a_{12}&0&a_{13}\cdots\\
0&a_{21}&2\pi\sqrt{-1}&a_{22}&0&a_{23}\cdots\\
0&a_{31}&0&a_{32}&2\pi\sqrt{-1}&a_{33}\cdots\\
\multicolumn{6}{c}{\cdots}
\end{array}
\]
then the new system of periodicity moduli,
\[
\begin{array}{cccc}
u_1^{(1)}&u_1^{(2)}&\cdots&u_1^{(2p)}\\
u_2^{(1)}&u_2^{(2)}&\cdots&u_2^{(2p)}\\
\multicolumn{4}{c}{\cdots}\\
u_p^{(1)}&u_p^{(2)}&\cdots&u_p^{(2p)}
\end{array}
\]
is given by the formula
\[
\left\{
\begin{aligned}
(1)\quad u_1^{(h)}&=\beta_{1,h}\cdot2\pi\sqrt{-1}+a_{11}\beta_{2,h}+a_{21}\beta_{4,h}+a_{31}\beta_{6,h}+\cdots,\\
u_2^{(h)}&=\beta_{3,h}\cdot2\pi\sqrt{-1}+a_{12}\beta_{2,h}+a_{22}\beta_{4,h}+a_{32}\beta_{6,h}+\cdots,\\
u_3^{(h)}&=\beta_{5,h}\cdot2\pi\sqrt{-1}+a_{13}\beta_{2,h}+a_{23}\beta_{4,h}+a_{33}\beta_{6,h}+\cdots,\\
&\hspace{5em}\cdots\qquad (h=1,2\ldots2p).
\end{aligned}
\right.
\]
Now set
\[
\begin{aligned}
(2)\quad
v_1&=m_{11}u_1+m_{12}u_2+\cdots+m_{1p}u_p,\\
v_2&=m_{21}u_1+m_{22}u_2+\cdots+m_{2p}u_p,\\
&\hspace{5em}\cdots\\
v_p&=m_{p1}u_1+m_{p2}u_2+\cdots+m_{pp}u_p.
\end{aligned}
\]
Then the \(m\) are determined, according to \(\S\,30\), by the formula
\[
(3)\quad
m_{hk}=2\pi\sqrt{-1}\,\frac{\partial\log\Delta}{\partial u_k^{(2h-1)}},
\]
where
\[
(4)\quad
\Delta=\Sigma\pm u_1^{(1)}u_2^{(3)}u_3^{(5)}\cdots u_p^{(2p-1)}.
\]
As a consequence of this determination, the system of periodicity moduli of the new normal integrals \(v\) becomes the following:
\[
\begin{array}{cccccc}
2\pi\sqrt{-1}&b_{11}&0&b_{12}&0&b_{13}\cdots\\
0&b_{21}&2\pi\sqrt{-1}&b_{22}&0&b_{23}\cdots\\
0&b_{31}&0&b_{32}&2\pi\sqrt{-1}&b_{33}\cdots\\
\multicolumn{6}{c}{\cdots,}
\end{array}
\]
where
\[
(5)\quad
b_{hk}=m_{h1}u_1^{(2k)}+m_{h2}u_2^{(2k)}+\cdots
=m_{k1}u_1^{(2h)}+m_{k2}u_2^{(2h)}+\cdots.
\]
After the new integrals of the first kind have been normalized in this way, it is easy to normalize correspondingly the integrals of the third kind. From the integral \(S_{\xi\eta}\) the normal integral \(\Pi_{\xi\eta}\) arose by making vanish, through addition of integrals of the first kind, those periodicity moduli which correspond to the periodicity moduli \(2\pi\sqrt{-1}\).

% --- Source pages 332--342; printed pages 310--320. ---

Thus, from the integral \(\int_z^x d\Pi_{\xi\eta}\) normalized for the \(u\), one obtains the integral \(\int_z^x dP_{\xi\eta}\) normalized for the \(v\) by means of the formula
\[
\int_z^x dP_{\xi\eta}
=
\int_z^x d\Pi_{\xi\eta}
-\frac{1}{2\pi\sqrt{-1}}
 \int_z^x(K_1\,dv_1+K_2\,dv_2+\cdots+K_p\,dv_p),
\]
by making vanish those periodicity moduli which correspond to the periodicity moduli \(2\pi\sqrt{-1}\) of the \(v\), that is, to the periodicity moduli
\[
u_k^{(1)},\ u_k^{(3)},\ldots,u_k^{(2p-1)}
\]
of the \(u\). Therefore \(K_h\) must be set equal to the periodicity modulus of \(\Pi\) which belongs to \(u_k^{(2h-1)}\), or one must put
\[
K_h=\int_{\eta}^{\xi}(\beta_{2,2h-1}du_1+\beta_{4,2h-1}du_2+\cdots+\beta_{2p,2h-1}du_p),
\]
so that for \(P_{\xi\eta}\) one obtains the defining equation
\[
(6)\quad
\int_z^x dP_{\xi\eta}
=
\int_z^x d\Pi_{\xi\eta}
-\frac{1}{2\pi\sqrt{-1}}
\sum_{h=1}^{h=p}\sum_{k=1}^{k=p}
\beta_{2k,2h-1}\int_{\eta}^{\xi}du_k\int_z^x dv_h.
\]
The last part of this expression, which still contains \(u,v\) mixed, can advantageously be represented by the latter alone, and then one also sees that, as is indeed necessary by the theorem on interchange of parameter and argument, it remains unchanged when \(x,z\) are interchanged with \(\xi,\eta\). For this purpose we first form the solutions of equations (2), which are very simple. Since namely the \(u_k\) reduce to \(u_k^{(2p-1)}\) when all \(v\) vanish except \(v_p\), and when this last is equal to \(2\pi\sqrt{-1}\), one necessarily has
\[
(7)\quad
u_k=\frac{1}{2\pi\sqrt{-1}}
\left(u_k^{(1)}v_1+u_k^{(3)}v_2+\cdots+u_k^{(2p-1)}v_p\right).
\]
Hence from (6), using this formula, one obtains
\[
(8)\quad
\int_z^x dP_{\xi\eta}
=
\int_z^x d\Pi_{\xi\eta}
+\frac{1}{4\pi^2}
\sum_{h=1}^{h=p}\sum_{m=1}^{m=p}
c_{hm}\int_z^x dv_h\int_{\eta}^{\xi}dv_m,
\]
where
\[
(9)\quad
c_{hm}=\sum_{k=1}^{k=p}\beta_{2k,2h-1}u_k^{(2m-1)}.
\]
If now parameter and argument are interchanged in (8), the last sum is found to be changed only in so far as \(c_{mh}\) takes the place of \(c_{hm}\). Thus one must have
\[
c_{hm}=c_{mh},
\]
or, substituting the values of the \(u_k^{(2m-1)}\) from (1),
\[
\begin{aligned}
0={}&2\pi\sqrt{-1}\sum_{k=1}^{k=p}
(\beta_{2k,2h-1}\beta_{2k-1,2m-1}
-\beta_{2k,2m-1}\beta_{2k-1,2h-1})\\
&+\sum_{s=1}^{s=p}\sum_{k=1}^{k=p}a_{sk}
(\beta_{2k,2h-1}\beta_{2s,2m-1}
-\beta_{2k,2m-1}\beta_{2s,2h-1}).
\end{aligned}
\]
This equation is in fact fulfilled. For clearly \(a_{sk}\) and \(a_{ks}\) are multiplied by equal and opposite numbers; hence the terms of the last sum cancel pairwise, and the coefficient of \(2\pi\sqrt{-1}\) vanishes because of the properties of a canonical determinant.

\subsection*{\S\,88. Transformation of the \(\Theta\)-function by canonical substitutions.}

If equation (8) is summed over \(p\) pairs of points \(x,z\), one reaches the new form of the transcendental function \(T\); and by treating this further as was done in Sections VII and VIII, one reaches the new forms of the functions \(U,V,\Theta\). Since in the end only the last matters, we shall pass at once from equation (8) to it. Denote the transcendental function \(\Theta\), now that the constants in the exponent are also changed, by \(\Theta(\omega,a)\), and the new transcendental function by \(\Theta(w,b)\).

From the formula
\[
\sum_{i=1}^{i=p}\int_{z^{(i)}}^{x^{(i)}}d\Pi_{\xi\eta}
=
-\log
\frac{
\Theta\!\left(\sum\int_{\alpha^{(i)}}^{x^{(i)}}du_h-\int_{\mu}^{\xi}du_h,a\right)
\Theta\!\left(\sum\int_{\alpha^{(i)}}^{z^{(i)}}du_h-\int_{\mu}^{\eta}du_h,a\right)
}{
\Theta\!\left(\sum\int_{\alpha^{(i)}}^{x^{(i)}}du_h-\int_{\mu}^{\eta}du_h,a\right)
\Theta\!\left(\sum\int_{\alpha^{(i)}}^{z^{(i)}}du_h-\int_{\mu}^{\xi}du_h,a\right)
}
\]
there follows at once the analogous formula
\[
\sum_{i=1}^{i=p}\int_{z^{(i)}}^{x^{(i)}}dP_{\xi\eta}
=
-\log
\frac{
\Theta\!\left(\sum\int_{\gamma^{(i)}}^{x^{(i)}}dv_h-\int_{\mu}^{\xi}dv_h,b\right)
\Theta\!\left(\sum\int_{\gamma^{(i)}}^{z^{(i)}}dv_h-\int_{\mu}^{\eta}dv_h,b\right)
}{
\Theta\!\left(\sum\int_{\gamma^{(i)}}^{x^{(i)}}dv_h-\int_{\mu}^{\eta}dv_h,b\right)
\Theta\!\left(\sum\int_{\gamma^{(i)}}^{z^{(i)}}dv_h-\int_{\mu}^{\xi}dv_h,b\right)
}.
\]
In this last formula the lower limits \(\alpha\) of the first have been replaced by new limits \(\gamma\); for one cannot decide in advance whether the contact curve of the \(\alpha\) must here be replaced by another one, whose points of contact we denote by the points \(\gamma\).

Now summing equation (8) over the \(p\) pairs of points \(x,z\), and inserting the above values for the integral sums, we obtain a relation between the different kinds of \(\Theta\), namely
\[
\resizebox{0.98\linewidth}{!}{$
\begin{aligned}
\log
\frac{
\frac{
\Theta\!\left(\sum\int_{\gamma^{(i)}}^{x^{(i)}}dv_h-\int_{\mu}^{\xi}dv_h,b\right)
\Theta\!\left(\sum\int_{\gamma^{(i)}}^{z^{(i)}}dv_h-\int_{\mu}^{\eta}dv_h,b\right)}
{\Theta\!\left(\sum\int_{\alpha^{(i)}}^{x^{(i)}}du_h-\int_{\mu}^{\xi}du_h,a\right)
\Theta\!\left(\sum\int_{\alpha^{(i)}}^{z^{(i)}}du_h-\int_{\mu}^{\eta}du_h,a\right)}
}{
\frac{
\Theta\!\left(\sum\int_{\gamma^{(i)}}^{x^{(i)}}dv_h-\int_{\mu}^{\eta}dv_h,b\right)
\Theta\!\left(\sum\int_{\gamma^{(i)}}^{z^{(i)}}dv_h-\int_{\mu}^{\xi}dv_h,b\right)}
{\Theta\!\left(\sum\int_{\alpha^{(i)}}^{x^{(i)}}du_h-\int_{\mu}^{\eta}du_h,a\right)
\Theta\!\left(\sum\int_{\alpha^{(i)}}^{z^{(i)}}du_h-\int_{\mu}^{\xi}du_h,a\right)}
}
={}&-\frac{1}{4\pi^2}
\sum_{h=1}^{h=p}\sum_{m=1}^{m=p}
c_{hm}
\left(\sum\int_{z^{(i)}}^{x^{(i)}}dv_h\right)
\int_{\eta}^{\xi}dv_m .
\end{aligned}
$}
\]
Now the identity evidently holds:
\[
\begin{aligned}
&\sum\sum c_{hm}\left(\sum\int_{z^{(i)}}^{x^{(i)}}dv_h\right)\int_{\eta}^{\xi}dv_m\\
={}&\frac12\sum\sum c_{hm}
\left(\sum\int_{\gamma^{(i)}}^{x^{(i)}}dv_h-\int_{\mu}^{\eta}dv_h\right)
\left(\sum\int_{\gamma^{(i)}}^{x^{(i)}}dv_m-\int_{\mu}^{\eta}dv_m\right)\\
&+\frac12\sum\sum c_{hm}
\left(\sum\int_{\gamma^{(i)}}^{z^{(i)}}dv_h-\int_{\mu}^{\xi}dv_h\right)
\left(\sum\int_{\gamma^{(i)}}^{z^{(i)}}dv_m-\int_{\mu}^{\xi}dv_m\right)\\
&-\frac12\sum\sum c_{hm}
\left(\sum\int_{\gamma^{(i)}}^{x^{(i)}}dv_h-\int_{\mu}^{\xi}dv_h\right)
\left(\sum\int_{\gamma^{(i)}}^{x^{(i)}}dv_m-\int_{\mu}^{\xi}dv_m\right)\\
&-\frac12\sum\sum c_{hm}
\left(\sum\int_{\gamma^{(i)}}^{z^{(i)}}dv_h-\int_{\mu}^{\eta}dv_h\right)
\left(\sum\int_{\gamma^{(i)}}^{z^{(i)}}dv_m-\int_{\mu}^{\eta}dv_m\right).
\end{aligned}
\]
Thus if, for brevity, one puts
\[
\resizebox{0.98\linewidth}{!}{$
\begin{aligned}
&\log
\frac{
\Theta\!\left(\sum\int_{\gamma^{(i)}}^{x^{(i)}}dv_h-\int_{\mu}^{\xi}dv_h,b\right)}
{\Theta\!\left(\sum\int_{\alpha^{(i)}}^{x^{(i)}}du_h-\int_{\mu}^{\xi}du_h,a\right)}
-\frac{1}{8\pi^2}\sum\sum c_{hm}
\left(\sum\int_{\gamma^{(i)}}^{x^{(i)}}dv_h-\int_{\mu}^{\xi}dv_h\right)
\left(\sum\int_{\gamma^{(i)}}^{x^{(i)}}dv_m-\int_{\mu}^{\xi}dv_m\right)\\
&\hspace{8em}=
\varphi\!\left(\sum\int_{\alpha^{(i)}}^{x^{(i)}}du-\int_{\mu}^{\xi}du\right),
\end{aligned}
$}
\]
then between the four functions \(\varphi\), each depending on only \(p\) arguments, one has the equation
\[
\begin{aligned}
&\varphi\!\left(\sum\int_{\alpha^{(i)}}^{x^{(i)}}du-\int_{\mu}^{\xi}du\right)
+\varphi\!\left(\sum\int_{\alpha^{(i)}}^{z^{(i)}}du-\int_{\mu}^{\eta}du\right)\\
&-\varphi\!\left(\sum\int_{\alpha^{(i)}}^{x^{(i)}}du-\int_{\mu}^{\eta}du\right)
-\varphi\!\left(\sum\int_{\alpha^{(i)}}^{z^{(i)}}du-\int_{\mu}^{\xi}du\right)=0.
\end{aligned}
\]
This equation is evidently satisfiable, as is also easily proved by differentiating with respect to the arguments, only if the function \(\varphi\) is a \emph{linear function} of its arguments. Thus put
\[
\left.
\begin{aligned}
\sum\int_{\alpha^{(i)}}^{x^{(i)}}du_h-\int_{\mu}^{\xi}du_h&=\omega_h,\\
\sum\int_{\gamma^{(i)}}^{x^{(i)}}dv_h-\int_{\mu}^{\xi}dv_h&=w_h,
\end{aligned}
\right\}
\qquad (h=1,2\ldots p);
\]
then
\[
\varphi(\omega)=\log C+c_1\omega_1+c_2\omega_2+\cdots+c_p\omega_p,
\]
where \(c,C\) denote constants still to be determined, and between the two kinds of \(\Theta\)-functions there is the equation
\[
(10)\quad
e^{-\frac{1}{8\pi^2}\sum\sum c_{hm}w_hw_m}
\Theta(w,b)=C\cdot e^{\sum c_h\omega_h}\cdot\Theta(\omega,a).
\]
Between \(\omega,w\), apart from additive half-periodicity moduli which must be added in order to make the limits equal, there exist the same relations that connected \(u\) with \(v\), namely
\[
(11)\quad
\omega_k-\frac12P_k
=
\frac{1}{2\pi\sqrt{-1}}
\left(u_k^{(1)}w_1+u_k^{(3)}w_2+\cdots+u_k^{(2p-1)}w_p\right);
\]
the half-periodicity moduli, however, serve to define the \(\gamma\) from the \(\alpha\), since one has
\[
(12)\quad
\sum\int_{\alpha^{(i)}}^{\gamma^{(i)}}du_k
=\frac12P_k
=\varrho_k\pi\sqrt{-1}
+\frac12(a_{1k}\sigma_1+a_{2k}\sigma_2+\cdots+a_{pk}\sigma_p).
\]
In relation (10) the constants \(c,C\), and the integers \(\varrho,\sigma\), still remain to be determined.

To carry out these determinations, let all \(w_h\) pass into \(-w_h\), whereby \(\omega_k\) passes into
\[
\frac12P_k-\left(\omega_k-\frac12P_k\right)=P_k-\omega_k.
\]
By this process the left-hand side of equation (10) remains unchanged; the right-hand side must therefore also remain unchanged, and hence
\[
\begin{aligned}
e^{\sum c_h\omega_h}\Theta(\omega,a)
&=e^{\sum c_h(P_h-\omega_h)}\Theta(P-\omega,a)\\
&=e^{\sum c_h(P_h-\omega_h)-\sum\sigma_h\omega_h+\frac12\sum\sum\sigma_h\sigma_ka_{hk}}\Theta(\omega,a).
\end{aligned}
\]
Comparison of the exponents first gives
\[
c_h=-\frac12\sigma_h;
\]
and then
\[
1=(-1)^{\sum\sigma_h\varrho_h}.
\]
The determination of the \(c\) is thus reduced to that of the \(\sigma\); the last equation teaches that \(\sum\sigma_h\varrho_h\) is always an even number, which gives the following theorem for the contact curves:

\begin{center}
\textbf{In transforming \(\Theta\)-functions into others by canonical substitutions, the lower limiting points that appear are always points of contact of proper contact curves of order \(n-2\) (\(\S\,75\)).}
\end{center}

\subsection*{\S\,89. Contact curves which occur in the lower limits.}

The determination of \(\varrho,\sigma\) themselves is obtained by investigating the periodicity properties of the two sides of (10). If we let the system of \(w_h\) increase by the periodicity moduli
\[
Q_h=2m_h\pi\sqrt{-1}+b_{1h}q_1+b_{2h}q_2+\cdots,
\]
then the \(\omega_h\) increase by the corresponding periodicity moduli
\[
\begin{aligned}
R_h&=m_1u_h^{(1)}+m_2u_h^{(3)}+m_3u_h^{(5)}+\cdots
+q_1u_h^{(2)}+q_2u_h^{(4)}+\cdots\\
&=2\mu_h\pi\sqrt{-1}+a_{1h}\nu_1+a_{2h}\nu_2+\cdots,
\end{aligned}
\]
where, with regard to equations (1), the numbers \(\mu,\nu\) receive the values
\[
\begin{aligned}
\mu_h={}&m_1\beta_{2h-1,1}+m_2\beta_{2h-1,3}+m_3\beta_{2h-1,5}+\cdots\\
&+q_1\beta_{2h-1,2}+q_2\beta_{2h-1,4}+q_3\beta_{2h-1,6}+\cdots,\\
\nu_h={}&m_1\beta_{2h,1}+m_2\beta_{2h,3}+m_3\beta_{2h,5}+\cdots\\
&+q_1\beta_{2h,2}+q_2\beta_{2h,4}+q_3\beta_{2h,6}+\cdots .
\end{aligned}
\tag{13}
\]
Now put in (10) simultaneously \(w_h+Q_h\) in place of \(w_h\), and \(\omega_h+R_h\) in place of \(\omega_h\), and compare the changes undergone thereby by the two sides of equation (10). We obtain
\[
\resizebox{0.98\linewidth}{!}{$
\begin{aligned}
e^{-\frac{1}{4\pi^2}\sum\sum c_{hk}w_hQ_k
-\frac{1}{8\pi^2}\sum\sum c_{hk}Q_hQ_k
+\sum w_hq_h+\frac12\sum\sum q_hq_kb_{hk}}
=
e^{-\frac12\sum\sigma_hR_h+\sum\nu_k\omega_k+\frac12\sum\sum\nu_h\nu_ka_{hk}},
\end{aligned}
$}
\]
or, expressing \(\omega_h\) by the \(w\) from (11),
\[
=e^{-\frac12\sum\sigma_hR_h+\frac12\sum\nu_hP_h
+\frac{1}{2\pi\sqrt{-1}}\sum\sum\nu_kw_hu_k^{(2h-1)}
+\frac12\sum\sum\nu_h\nu_ka_{hk}}.
\]
In this equation put \(w_h=-\frac12Q_h\); then
\[
\begin{aligned}
(-1)^{\sum q_hm_h}
&=(-1)^{\sum(\nu_h\varrho_h-\sigma_h\mu_h)}
e^{-\frac{1}{4\pi\sqrt{-1}}\sum\sum\nu_kQ_hu_k^{(2h-1)}
+\frac12\sum\sum\nu_h\nu_ka_{hk}}.
\end{aligned}
\]
Now observe that, when as in (7) the \(v_h\) increase by the periodicity moduli \(Q_h\), \(u_k\) must increase by the corresponding periodicity modulus \(R_k\). Thus
\[
R_k=\frac{1}{2\pi\sqrt{-1}}\sum_h Q_hu_k^{(2h-1)},
\]
and the above equation becomes
\[
\begin{aligned}
(-1)^{\sum q_hm_h}
&=(-1)^{\sum(\nu_h\varrho_h-\sigma_h\mu_h)}
e^{-\frac12\sum\nu_kR_k+\frac12\sum\sum\nu_h\nu_ka_{hk}}\\
&=(-1)^{\sum(\nu_h\varrho_h-\sigma_h\mu_h-\nu_h\mu_h)}.
\end{aligned}
\]
Therefore, for all integral systems of values of the \(m,q\), the equation
\[
(14)\quad
\sum q_hm_h-\sum\nu_h\mu_h+\sum\nu_h\varrho_h-\sum\sigma_h\mu_h
\equiv0 \pmod{2}
\]
must be fulfilled. If in this equation we put
\[
m_k=\beta_{2h,2k},\qquad q_k=-\beta_{2h,2k-1}\qquad (k=1,2\ldots p),
\]
then equations (13) show that, because of the properties of the elements of a canonical determinant, all \(\mu\) and \(\nu\) vanish except \(\mu_h\), which becomes equal to 1. Equation (14) therefore gives
\[
(15)\quad
\sigma_h\equiv\sum_{k=1}^{k=p}\beta_{2h,2k}\beta_{2h,2k-1}\pmod{2}.
\]
If on the other hand one puts
\[
m_k=-\beta_{2h-1,2k},\qquad q_k=\beta_{2h-1,2k-1}\qquad (k=1,2\ldots p),
\]
then by (13) all \(\mu,\nu\) vanish except \(\nu_h\), which becomes equal to 1, and (14) gives
\[
(16)\quad
\varrho_h\equiv\sum_{k=1}^{k=p}\beta_{2h-1,2k}\beta_{2h-1,2k-1}\pmod{2}.
\]
The numbers \(\varrho,\sigma\) are thereby determined as far as the nature of the matter permits; for if they are increased by even numbers, the right-hand side of equation (10) changes only by a constant. The formula therefore remains unchanged, and only the constant \(C\) receives a different meaning. Indeed, if, while keeping the \(w\) unchanged, we change \(\varrho_h,\sigma_h\) by \(2r_h,2s_h\), then \(\omega_k\) receives the increment
\[
M_k=2r_k\pi\sqrt{-1}+a_{1k}s_1+a_{2k}s_2+\cdots .
\]
Hence
\[
e^{-\frac12\sum\sigma_h\omega_h}\Theta(\omega,a)
\]
is transformed into
\[
e^{-\frac12\sum(\sigma_h+2s_h)(\omega_h+M_h)}
e^{\sum s_h\omega_h+\frac12\sum\sum s_hs_ka_{hk}}\Theta(\omega,a),
\]
which in fact differs from the preceding expression only by a constant.

The connection between the two \(\Theta\)-functions arising from different systems of normal periodicity moduli is thus finally expressed by the equation
\[
(17)\quad
e^{-\frac{1}{8\pi^2}\sum\sum c_{hm}w_hw_m}\Theta(w,b)
=
C\,e^{-\frac12\sum\sigma_h\omega_h}\Theta(\omega,a),
\]
where \(w,\omega\) are connected by the equations
\[
\omega_k-\frac12P_k=
\frac{1}{2\pi\sqrt{-1}}
\left(u_k^{(1)}w_1+u_k^{(3)}w_2+\cdots+u_k^{(2p-1)}w_p\right),
\]
\[
P_k=\varrho_k\pi\sqrt{-1}+\frac12(a_{1k}\sigma_1+a_{2k}\sigma_2+\cdots+a_{pk}\sigma_p),
\]
and \(\varrho,\sigma\) are defined by equations (15), (16) by means of the elements of the canonical determinant. To determine them completely, one may assume that they take only the values 0 or 1.

It remains to show that the numbers \(\sigma,\varrho\), as given by (15), (16), can indeed be even or odd independently of one another, with the sole restriction
\[
\sum\sigma_h\varrho_h\equiv0\pmod{2};
\]
in other words, to show that for different choices of the normal periodicity moduli all proper contact curves \(\alpha,\gamma\) can occur in the lower limits.

This is done simply by writing down canonical determinants which actually furnish all contact curves. First consider the canonical determinant
\[
\left|
\begin{array}{cccccc}
a_1&b_1&0&0&0&0\cdots\\
c_1&d_1&0&0&0&0\cdots\\
0&0&a_2&b_2&0&0\cdots\\
0&0&c_2&d_2&0&0\cdots\\
0&0&0&0&a_3&b_3\cdots\\
0&0&0&0&c_3&d_3\cdots\\
&&\cdots
\end{array}
\right|.
\]
The conditions under which it is canonical are
\[
a_1d_1-b_1c_1=\pm1,\qquad
a_2d_2-b_2c_2=\pm1,\qquad
a_3d_3-b_3c_3=\pm1\quad\text{etc.}
\]
At the same time one has
\[
\varrho_1=a_1c_1,\quad\varrho_2=a_2c_2,\quad\varrho_3=a_3c_3\cdots,
\]
\[
\sigma_1=b_1d_1,\quad\sigma_2=b_2d_2,\quad\sigma_3=b_3d_3\cdots .
\]
It is then immediately seen that, with a suitable determination of \(a,b,c,d\), the \(\varrho,\sigma\) can be even or odd independently of one another; the only restriction is that none of the pairs \(\varrho_i,\sigma_i\) may consist of two \emph{odd} numbers. Because of the equation \(\sum\sigma_h\varrho_h\equiv0\pmod{2}\), the latter case can occur only for an even number of pairs. If \(\sigma_i,\varrho_i\) and \(\sigma_k,\varrho_k\) are to become two such pairs, then one need only replace the elements
\[
\begin{array}{ccccc}
a_i&b_i&\cdots&0&0\\
c_i&d_i&\cdots&0&0\\
&&\cdots\\
0&0&\cdots&a_k&b_k\\
0&0&\cdots&c_k&d_k
\end{array}
\]
of the above determinant by
\[
\begin{array}{ccccc}
1&1&\cdots&0&-1\\
0&1&\cdots&0&-1\\
&&\cdots\\
1&0&\cdots&1&0\\
1&0&\cdots&1&1
\end{array}.
\]
As often as this is done, one has two pairs \(\sigma,\varrho\) which consist of two odd numbers; the remaining pairs, as before, are independently composed of even and odd, or of two even, numbers.

Thus determinants have in fact been formed which furnish all required cases, and with this prove the above theorem. At the same time one sees that all canonical determinants divide into \(2^{p-1}(2^p+1)\) groups, according as the relevant systems of \(\varrho,\sigma\) lead to one or another contact curve.

\subsection*{\S\,90. Reduction of the constant occurring in the transformation formula for \(\Theta\) to a numerical one.}

We now come to finding the value of the constant \(C\) in the formula
\[
(1)\quad
e^{-\frac{1}{8\pi^2}\sum\sum c_{hm}w_hw_m+\frac12\sum\sigma_h\omega_h}
\Theta(w,b)
=C\cdot\Theta(\omega,a).
\]
For this purpose we first note that the function
\[
\Theta(\omega,a)=\sum e^{-\frac12\sum\sum a_{ik}m_im_k+\sum m_i\omega_i}
\]
plainly satisfies the following partial differential equations\footnote{In these equations the \(a_{ik}\) are treated as independent of one another. But in fact formula (1) also holds, as one easily sees, for the fully general case in which the \(a_{ik}\) are arbitrary quantities, endowed only with the property of making the \(\Theta\)-series converge. To prove this, one need only define the synectic function on the left by its periodicity properties; it can then be expanded in a series which agrees with the right-hand side of (1).}:
\[
(2)\quad
\frac{\partial\Theta(\omega,a)}{\partial a_{ii}}
=-\frac12\frac{\partial^2\Theta(\omega,a)}{\partial\omega_i^2},
\qquad
\frac{\partial\Theta(\omega,a)}{\partial a_{ik}}
=-\frac{\partial^2\Theta(\omega,a)}{\partial\omega_i\partial\omega_k}.
\]
Dividing both sides of (1) by \(C\), the left-hand side of the equation must also satisfy these partial differential equations; here \(c,C\) are functions only of the \(a_{ik}\), whereas the \(w\) depend both on the \(a_{ik}\) and on the \(\omega_h\). And since the left-hand side of equation (1) is a sum of exponential quantities of different kinds, each term must satisfy those differential equations for itself. In particular consider the term arising from the term 1 of the function \(\Theta(w,b)\); it is \(\frac{e^\psi}{C}\), if for brevity one sets
\[
\psi=-\frac{1}{8\pi^2}\sum\sum c_{hm}w_hw_m+\frac12\sum\sigma_h\omega_h.
\]
The function \(\frac{e^\psi}{C}\) must therefore satisfy equations (2), or the single total differential equation into which those can be combined:
\[
\sum\frac{\partial\frac{e^\psi}{C}}{\partial a_{ii}}\,da_{ii}
+\sum\frac{\partial\frac{e^\psi}{C}}{\partial a_{ik}}\,da_{ik}
=
-\frac12\sum\frac{\partial^2\frac{e^\psi}{C}}{\partial\omega_i^2}\,da_{ii}
-\sum\frac{\partial^2\frac{e^\psi}{C}}{\partial\omega_i\partial\omega_k}\,da_{ik},
\]
where in the sums each combination \(ik\) is to be used only once, not multiplied by the factor 2. One can write instead
\[
\begin{aligned}
-\frac{dC}{C}
+\sum\frac{\partial\psi}{\partial a_{ii}}\,da_{ii}
+\sum\frac{\partial\psi}{\partial a_{ik}}\,da_{ik}
={}&-\frac12\sum\frac{\partial^2\psi}{\partial\omega_i^2}\,da_{ii}
-\sum\frac{\partial^2\psi}{\partial\omega_i\partial\omega_k}\,da_{ik}\\
&-\frac12\sum\left(\frac{\partial\psi}{\partial\omega_i}\right)^2da_{ii}
-\sum\frac{\partial\psi}{\partial\omega_i}\frac{\partial\psi}{\partial\omega_k}da_{ik},
\end{aligned}
\]
or
\[
\begin{aligned}
(3)\quad d\log C={}&
\sum\frac{\partial\psi}{\partial a_{ii}}\,da_{ii}
+\sum\frac{\partial\psi}{\partial a_{ik}}\,da_{ik}
+\frac12\sum\frac{\partial^2\psi}{\partial\omega_i^2}\,da_{ii}\\
&+\sum\frac{\partial^2\psi}{\partial\omega_i\partial\omega_k}\,da_{ik}
+\frac12\sum\left(\frac{\partial\psi}{\partial\omega_i}\right)^2da_{ii}
+\sum\frac{\partial\psi}{\partial\omega_i}\frac{\partial\psi}{\partial\omega_k}da_{ik}.
\end{aligned}
\]
Now form the differential quotients of \(\psi\), and then substitute for the \(\omega_h\) those constant values for which the \(w_h\) vanish. One finds
\[
(4)\quad
\frac{\partial\psi}{\partial a_{ii}}=0,\qquad
\frac{\partial\psi}{\partial a_{ik}}=0,\qquad
\frac{\partial\psi}{\partial\omega_i}=\frac12\sigma_i,
\]
\[
\frac{\partial^2\psi}{\partial\omega_i\partial\omega_k}
=
-\frac{1}{4\pi^2}\sum\sum c_{hm}
\frac{\partial w_h}{\partial\omega_i}
\frac{\partial w_m}{\partial\omega_k}.
\]
Further, substituting the value
\[
c_{hm}=\sum_r\beta_{2r,2h-1}u_r^{(2m-1)},
\]
one obtains
\[
\frac{\partial^2\psi}{\partial\omega_i\partial\omega_k}
=
\sum_h\sum_m\sum_r
\beta_{2r,2h-1}u_r^{(2m-1)}
\frac{\partial}{\partial\omega_i}\frac{w_h}{2\pi\sqrt{-1}}\,
\frac{\partial}{\partial\omega_k}\frac{w_m}{2\pi\sqrt{-1}},
\]
or, with regard to (11) of \(\S\,88\),
\[
=
\sum_h\sum_r\beta_{2r,2h-1}
\frac{\partial}{\partial\omega_i}\frac{w_h}{2\pi\sqrt{-1}}\,
\frac{\partial\omega_r}{\partial\omega_k}
=
\sum_h\beta_{2k,2h-1}
\frac{\partial}{\partial\omega_i}\frac{w_h}{2\pi\sqrt{-1}}.
\]
Finally, if we introduce the function \(\Delta\) of \(\S\,87\), then
\[
\frac{\partial}{\partial\omega_i}\frac{w_h}{2\pi\sqrt{-1}}
=
\frac{\partial\log\Delta}{\partial u_i^{(2h-1)}},
\]
and since also
\[
\beta_{2k,2h-1}=\frac{\partial u_i^{(2h-1)}}{\partial a_{ik}},
\]
it finally follows that
\[
(5)\quad
\frac{\partial^2\psi}{\partial\omega_i\partial\omega_k}
=
\frac{\partial\log\Delta}{\partial a_{ik}},
\]
if the last differential quotient is understood so that differentiation with respect to \(a_{ik}\) is taken only in so far as it occurs in the \(u\) with lower index \(i\). Substituting the values (4), (5) in (3), we get
\[
d\log C=\frac12d\log\Delta+\frac18d\!\left(\sum\sum\sigma_i\sigma_ka_{ik}\right),
\]
or
\[
C=\gamma\cdot\sqrt{\Delta}\cdot e^{\frac18\sum\sum\sigma_i\sigma_ka_{ik}},
\]
where now the constant \(\gamma\) depends neither on the \(\omega\), nor on the \(a_{ik}\), but must be purely numerical. The formula for the connection between the two functions \(\Theta\) is transformed into the following:
\[
e^{-\frac{1}{8\pi^2}\sum\sum c_{hm}w_hw_m}
\Theta(w,b)
=
\gamma\sqrt{\Delta}\,
e^{-\frac12\sum\sigma_h\omega_h+\frac18\sum\sum\sigma_i\sigma_ka_{ik}}
\Theta(\omega,a).
\]


% --- Source pages 343--355; printed pages 321--333. ---

\subsection*{\S\,91. The numerical constant in the composition of transformations.}

The quantity \(\gamma\) is a purely numerical constant, depending on the properties of the numbers \(\beta\). To determine it, we must go back to the composition of the canonical determinant from simple ones. Suppose two successive substitutions with canonical determinant are composed, as was carried out in \(\S\,85\). The first substitution, in which we denote the elements of the canonical determinants by \(\beta\), and the corresponding combinations of them by \(\varrho,\sigma\), carries the variables \(\omega\) to the variables \(v\) with the aid of the equations
\[
(1)\quad
\omega_k-\frac12P_k
=\frac{1}{2\pi\sqrt{-1}}\left(u_k^{(1)}v_1+u_k^{(3)}v_2+\cdots+u_k^{(2p-1)}v_p\right).
\]
The second substitution carries the \(v\) to new variables \(W\). The elements of the canonical determinant used here are denoted by \(\beta'\), and the corresponding combinations of them by \(\varrho',\sigma'\); the fundamental periodicity moduli, which here are integral combinations of the \(b_{ik}\) and \(2\pi\sqrt{-1}\), and which correspond to the \(u_k^{(i)}\), shall be represented by \(v_k^{(i)}\). Finally, let the periods \(P\) here be replaced by periods \(P'\), which are likewise integral combinations of the \(b_{ik}\) and \(2\pi\sqrt{-1}\). The formula for the relation among the variables is then
\[
(2)\quad
v_k-\frac12P_k'
=\frac{1}{2\pi\sqrt{-1}}\left(v_k^{(1)}W_1+v_k^{(3)}W_2+\cdots+v_k^{(2p-1)}W_p\right).
\]
Finally, from these two there is composed a substitution carrying the \(\omega\) to the \(W\):
\[
(3)\quad
\omega_k-\frac12P_k''
=\frac{1}{2\pi\sqrt{-1}}\left(w_k^{(1)}W_1+w_k^{(3)}W_2+\cdots+w_k^{(2p-1)}W_p\right).
\]
In this substitution we denote by \(\beta'',\varrho'',\sigma''\) the quantities corresponding to \(\beta,\varrho,\sigma\); the periods \(w_k^{(h)},P_k''\) are again composed of the \(a_{ik}\) and \(2\pi\sqrt{-1}\). Substituting the values of the \(v_k\) from (2) into (1), and comparing the resulting formula with (3), gives the following relations:
\[
(4)\quad
\left\{\begin{aligned}
2\pi\sqrt{-1}\,w_k^{(1)}&=u_k^{(1)}v_1^{(1)}+u_k^{(3)}v_2^{(1)}+\cdots+u_k^{(2p-1)}v_p^{(1)},\\
2\pi\sqrt{-1}\,w_k^{(3)}&=u_k^{(1)}v_1^{(3)}+u_k^{(3)}v_2^{(3)}+\cdots+u_k^{(2p-1)}v_p^{(3)},\\
&\hspace{4em}\cdots\\
2\pi\sqrt{-1}\,w_k^{(2p-1)}&=u_k^{(1)}v_1^{(2p-1)}+u_k^{(3)}v_2^{(2p-1)}+\cdots+u_k^{(2p-1)}v_p^{(2p-1)}.
\end{aligned}\right.
\]
\[
(5)\quad
2\pi\sqrt{-1}\,(P_k''-P_k)
=u_k^{(1)}P_1'+u_k^{(3)}P_2'+\cdots+u_k^{(2p-1)}P_{2p}'.
\]
The first \(p\) of these equations give nothing other than some of the relations
\[
(6)\quad
\beta_{k,r}''=\beta_{k1}\beta_{1r}'+\beta_{k2}\beta_{2r}'+\cdots+\beta_{k,2p}\beta_{2p,r}',
\]
by means of which the elements of the composed canonical determinant are formed. The last equation, however, teaches that the numbers \(\varrho'',\sigma''\) are not merely determined, like \(\varrho,\sigma,\varrho',\sigma'\), by the elements \(\beta,\beta'\) up to even numbers, but that they are composed from these and from \(\varrho,\sigma,\varrho',\sigma'\) in a completely definite way. Substituting for \(u_k^{(h)},P_k,P_k',P_k''\) their values
\[
(7)\quad
u_k^{(h)}=\beta_{2k-1,h}\,2\pi\sqrt{-1}+\sum_m\beta_{2m,h}a_{mk}
\]
\[
(8)\quad
\left\{\begin{aligned}
P_k&=\varrho_k\cdot2\pi\sqrt{-1}+\sigma_1a_{1k}+\sigma_2a_{2k}+\cdots+\sigma_pa_{pk},\\
P_k''&=\varrho_k''\cdot2\pi\sqrt{-1}+\sigma_1''a_{1k}+\sigma_2''a_{2k}+\cdots+\sigma_p''a_{pk},\\
P_k'&=\varrho_k'\cdot2\pi\sqrt{-1}+\sigma_1'b_{1k}+\sigma_2'b_{2k}+\cdots+\sigma_p'b_{pk},
\end{aligned}\right.
\]
one obtains
\[
(\varrho_k''-\varrho_k)2\pi\sqrt{-1}+\sum_h(\sigma_h''-\sigma_h)a_{hk}
=\sum_h u_k^{(2h-1)}\varrho_h'
+\frac{1}{2\pi\sqrt{-1}}\sum_h\sum_i u_k^{(2h-1)}b_{hi}\sigma_i'.
\]
In the meantime we know that an increase of the \(v_k\) by the periodicity moduli \(b_{ki}\) corresponds to an increase of the \(\omega_k\) by \(u_k^{(2i)}\), so that
\[
(9)\quad
u_k^{(2i)}=\frac{1}{2\pi\sqrt{-1}}\sum_hu_k^{(2h-1)}b_{hi}.
\]
Thus the above formula passes into
\[
(\varrho_k''-\varrho_k)2\pi\sqrt{-1}+\sum_h(\sigma_h''-\sigma_h)a_{hk}
=\sum_hu_k^{(2h-1)}\varrho_h'+\sum_hu_k^{(2h)}\sigma_h'.
\]
Thus, to determine \(\varrho'',\sigma''\), one obtains the formulae:
\[
(10)\quad
\left\{\begin{aligned}
\varrho_k''-\varrho_k&=\sum_h(\varrho_h'\beta_{2k-1,2h-1}+\sigma_h'\beta_{2k-1,2h}),\\
\sigma_k''-\sigma_k&=\sum_h(\varrho_h'\beta_{2k,2h-1}+\sigma_h'\beta_{2k,2h}).
\end{aligned}\right.
\]
The three formulae for the \(\Theta\)-functions are obtained from the earlier formula by replacing, for the second and third substitutions, \(c,a,b,\Delta\) respectively by \(c',b,A,\Delta'\) and by \(c'',a,A,\Delta''\), so that we obtain
\[
\begin{aligned}
e^{-\frac{1}{8\pi^2}\sum\sum c_{hm}v_hv_m}\Theta(v,b)
&=\gamma\sqrt{\Delta}\,e^{-\frac12\sum\sigma_h\omega_h+\frac18\sum\sum\sigma_i\sigma_ka_{ik}}\Theta(\omega,a),\\
e^{-\frac{1}{8\pi^2}\sum\sum c_{hm}'W_hW_m}\Theta(W,A)
&=\gamma'\sqrt{\Delta'}\,e^{-\frac12\sum\sigma_h'v_h+\frac18\sum\sum\sigma_i'\sigma_k'b_{ik}}\Theta(v,b),\\
e^{-\frac{1}{8\pi^2}\sum\sum c_{hm}''W_hW_m}\Theta(W,A)
&=\gamma''\sqrt{\Delta''}\,e^{-\frac12\sum\sigma_h''\omega_h+\frac18\sum\sum\sigma_i''\sigma_k''a_{ik}}\Theta(\omega,a).
\end{aligned}
\]
We now divide the product of the first two equations by the third and at the same time put the \(W\) equal to zero, so that \(\omega_h=\frac12P_h''\) and \(v_h=\frac12P_h'\). We then find
\[
\begin{aligned}
(11)\quad
e^{-\frac{1}{32\pi^2}\sum\sum c_{hm}P_h'P_m'}
={}&\frac{\gamma\gamma'}{\gamma''}\cdot\frac{\sqrt\Delta\sqrt{\Delta'}}{\sqrt{\Delta''}}\\
&\cdot e^{\frac14\sum[(\sigma_h''-\sigma_h)P_h''-\sigma_h'P_h']
-\frac18\sum\sum[(\sigma_i''\sigma_k''-\sigma_i\sigma_k)a_{ik}-\sigma_i'\sigma_k'b_{ik}]}.
\end{aligned}
\]
Since \(\Delta\) is the determinant of the \(u_k^{(2h-1)}\), \(\Delta'\) that of the \(v_k^{(2h-1)}\), and \(\Delta''\) that of the \(w_k^{(2h-1)}\), we have, by (4),
\[
(12)\quad
\Delta\Delta'=\Delta''\cdot(2\pi\sqrt{-1})^{\frac p4}.
\]
Taking \((2\pi\sqrt{-1})^{\frac p2}\) in an arbitrarily fixed sense once for all, we shall always choose the sign of \(\sqrt{\Delta''}\) so that also
\[
(13)\quad
\sqrt\Delta\cdot\sqrt{\Delta'}=\sqrt{\Delta''}\cdot(2\pi\sqrt{-1})^{\frac p2}.
\]
We further put the \(a_{ik}\) equal to zero in (11), since only the purely numerical ratio \(\frac{\gamma\gamma'}{\gamma''}\) is to be determined. Then \(P_h\) immediately reduces to \(\varrho_h\cdot2\pi\sqrt{-1}\), \(P_h''\) to \(\varrho_h''\cdot2\pi\sqrt{-1}\), and the \(b_{ik}\) are determined by the equations into which (9) now pass:
\[
(14)\quad
\beta_{2k-1,2i}\cdot2\pi\sqrt{-1}=\sum_h\beta_{2k-1,2h-1}b_{ih}.
\]
Substituting the values of the \(c_{hm}\), one has
\[
\begin{aligned}
(15)\quad
\sum\sum c_{hm}P_h'P_m'
&=2\pi\sqrt{-1}\sum\sum\sum\beta_{2k,2h-1}u_k^{(2m-1)}P_m'P_h'\\
&=2\pi\sqrt{-1}\sum\sum\beta_{2k,2h-1}(P_k''-P_k)P_h'.
\end{aligned}
\]
Since the \(a_{ik}\) are supposed to vanish, this reduces to
\[
\begin{aligned}
&(2\pi\sqrt{-1})^2\sum\sum\beta_{2k,2h-1}(\varrho_k''-\varrho_k)P_h'\\
&=(2\pi\sqrt{-1})^3\sum\sum\beta_{2k,2h-1}(\varrho_k''-\varrho_k)\varrho_h'\\
&\quad+(2\pi\sqrt{-1})^2\sum\sum\sum\beta_{2k,2h-1}(\varrho_k''-\varrho_k)\sigma_m'b_{hm}.
\end{aligned}
\]
Substituting in the second part of this expression for \(\varrho_k''-\varrho_k\) its value from (10), we get
\[
\begin{aligned}
(16)\quad
\sum_k\sum_h\beta_{2k,2h-1}(\varrho_k''-\varrho_k)b_{hm}
={}&\sum_k\sum_h\sum_i\beta_{2k,2h-1}\bigl(\varrho_i'\beta_{2k-1,2i-1}\nobreak+\sigma_i'\beta_{2k-1,2i}\bigr)b_{hm}.
\end{aligned}
\]
Now from the equations holding among the elements of the canonical determinant we have
\[
(17)\quad
\left\{\begin{aligned}
\sum_k\beta_{2k,2h-1}\beta_{2k-1,2i-1}
&=\sum_k\beta_{2k-1,2h-1}\beta_{2k,2i-1},\\
\sum_k\beta_{2k,2h-1}\beta_{2k-1,2i}
&=\sum_k\beta_{2k-1,2h-1}\beta_{2k,2i},
\end{aligned}\right.
\]
and only when \(h=i\) is the latter equation to be replaced by
\[
(18)\quad
\sum_k\beta_{2k,2i-1}\beta_{2k-1,2i}
=-1+\sum_k\beta_{2k-1,2i-1}\beta_{2k,2i}.
\]
Expression (16) is therefore transformed into
\[
\sum_k\sum_h\sum_i\beta_{2k-1,2h-1}(\varrho_i'\beta_{2k,2i-1}+\sigma_i'\beta_{2k,2i})b_{hm}
-
\sum_h\sigma_h'b_{hm},
\]
and summing with respect to \(h\), we obtain from (14):
\[
2\pi\sqrt{-1}\sum_k\sum_i\sum_m\beta_{2k-1,2m}(\varrho_i'\beta_{2k,2i-1}+\sigma_i'\beta_{2k,2i})\sigma_m'
-
\sum_h\sum_m\sigma_h'\sigma_m'b_{hm}.
\]
Thus expression (15) becomes
\[
\begin{aligned}
&(2\pi\sqrt{-1})^3\sum_k\sum_h\beta_{2k,2h-1}(\varrho_k''-\varrho_k)\varrho_h'\\
&\quad+(2\pi\sqrt{-1})^3\sum_k\sum_i\sum_m\beta_{2k-1,2m}(\varrho_i'\beta_{2k,2i-1}+\sigma_i'\beta_{2k,2i})\sigma_m'\\
&\quad-(2\pi\sqrt{-1})^2\sum_h\sum_m\sigma_h'\sigma_m'b_{hm}.
\end{aligned}
\]
Collecting everything, one finally obtains from (11) the formula
\[
\begin{aligned}
(19)\quad
\frac{\gamma\gamma'}{\gamma''}
={}&\frac{1}{(2\pi\sqrt{-1})^{p/2}}
\exp\!\left(\frac{2\pi\sqrt{-1}}{8}\sum_k\sum_h\beta_{2k,2h-1}(\varrho_k''-\varrho_k)\varrho_h'\right)\\
&\cdot\exp\!\left(\frac{2\pi\sqrt{-1}}{8}\sum_k\sum_m\sum_i\beta_{2k-1,2m}\sigma_m'(\varrho_i'\beta_{2k,2i-1}+\sigma_i'\beta_{2k,2i})\right)\\
&\cdot\exp\!\left(-\frac{2\pi\sqrt{-1}}{4}\bigl[(\sigma_h''-\sigma_h)\varrho_h''-\sigma_h'\varrho_h'\bigr]\right).
\end{aligned}
\]
Thus, apart from the factor \((2\pi\sqrt{-1})^{p/2}\), which may also be regarded as an integrating part of each factor \(\gamma\), \(\gamma''\) differs from the product \(\gamma\gamma'\) only by an eighth root of unity. If we put
\[
\gamma=\delta(2\pi\sqrt{-1})^{-p/2},\qquad
\gamma'=\delta'(2\pi\sqrt{-1})^{-p/2},\qquad
\gamma''=\delta''(2\pi\sqrt{-1})^{-p/2},
\]
then the \(\delta\) are connected by the equation
\[
\begin{aligned}
(20)\quad
\frac{\delta\delta'}{\delta''}
={}&\exp\!\left(\frac{2\pi\sqrt{-1}}{8}\sum\sum\beta_{2k,2h-1}(\varrho_k''-\varrho_k)\varrho_h'\right)\\
&\cdot\exp\!\left(
\frac{2\pi\sqrt{-1}}{8}\sum\sum\sum\beta_{2k-1,2m}\sigma_m'(\varrho_i'\beta_{2k,2i-1}+\sigma_i'\beta_{2k,2i})
\right.\\
&\qquad\left.
-\frac{2\pi\sqrt{-1}}{4}\bigl[(\sigma_h''-\sigma_h)\varrho_h''-\sigma_h'\varrho_h'\bigr]\right),
\end{aligned}
\]
and the formula for \(\Theta\) is
\[
\begin{aligned}
(21)\quad
e^{-\frac{1}{8\pi^2}\sum\sum c_{hm}v_hv_m}\Theta(v,b)
={}&\delta\sqrt{\frac{\Delta}{(2\pi\sqrt{-1})^p}}\,
 e^{-\frac12\sum\sigma_h\omega_h+\frac18\sum\sum\sigma_i\sigma_ka_{ik}}\Theta(\omega,a).
\end{aligned}
\]
Formula (20) can moreover be given in a substantially simpler form. From (10) one has
\[
\sum\sum\beta_{2k,2h-1}(\varrho_k''-\varrho_k)\varrho_h'
=\sum\sum\sum\beta_{2k,2h-1}\varrho_i'(\beta_{2k-1,2i-1}\varrho_i'+\beta_{2k-1,2i}\sigma_i').
\]
If we now apply formulas (17) to the sums to be taken over \(k\), this becomes
\[
\sum\sum\sum\beta_{2k-1,2h-1}\varrho_h'(\beta_{2k,2i-1}\varrho_i'+\beta_{2k,2i}\sigma_i'),
\]
and the expression multiplied by \(\frac{2\pi\sqrt{-1}}8\) in (20) is therefore
\[
\begin{aligned}
&\sum_k\left[\sum_h(\beta_{2k-1,2h-1}\varrho_h'+\beta_{2k-1,2h}\sigma_h')\right]
\left[\sum_i(\beta_{2k,2i-1}\varrho_i'+\beta_{2k,2i}\sigma_i')\right]\\
&=\sum_k(\varrho_k''-\varrho_k)(\sigma_k''-\sigma_k),
\end{aligned}
\]
so that (20) takes the form
\[
(22)\quad
\frac{\delta\delta'}{\delta''}
=\exp\!\left(\frac{2\pi\sqrt{-1}}8\sum_k[2\sigma_k'\varrho_k'-(\varrho_k''+\varrho_k)(\sigma_k''-\sigma_k)]\right).
\]
Finally this formula is still further simplified by putting
\[
\delta=\varepsilon e^{\frac{2\pi\sqrt{-1}}4\sum_k\varrho_k\sigma_k},
\qquad
\delta'=\varepsilon' e^{\frac{2\pi\sqrt{-1}}4\sum_k\varrho_k'\sigma_k'},
\qquad
\delta''=\varepsilon'' e^{\frac{2\pi\sqrt{-1}}4\sum_k\varrho_k''\sigma_k''}.
\]
The formula for the \(\Theta\)-functions then becomes
\[
\begin{aligned}
(23)\quad
e^{-\frac{1}{8\pi^2}\sum\sum c_{hm}v_hv_m}\Theta(v,b)
={}&\varepsilon\sqrt{\frac{\Delta}{(2\pi\sqrt{-1})^p}}
 e^{\begin{subarray}{l}\frac{2\pi\sqrt{-1}}4\sum\varrho_k\sigma_k-\frac12\sum\sigma_h\omega_h\\[0.1em]
 +\frac18\sum\sum\sigma_i\sigma_ka_{ik}\end{subarray}}\Theta(\omega,a),
\end{aligned}
\]
and the formula for the connection among the numbers \(\varepsilon\) is
\[
(24)\quad
\frac{\varepsilon\varepsilon'}{\varepsilon''}
=\exp\!\left(\frac{2\pi\sqrt{-1}}8\sum_k(\varrho_k''-\varrho_k)(\sigma_k''+\sigma_k)\right).
\]

\subsection*{\S\,92. Determination of the numerical constant for the simple transformations.}

It remains only to set up the general formula between the various \(\Theta\)'s for the case in which the canonical determinant is a \emph{simple} one, and in this case actually to determine the coefficients \(\varepsilon\). Since every substitution can be composed from simple ones, formula (24) of the preceding section gives the means of determining \(\varepsilon\) for every substitution from these. We consider in order the various simple substitutions, putting the non-vanishing elements of the simple determinant in definite positions for simplicity of presentation. It is immediately clear what changes when the \(\pm1\) are put in other positions.

\begin{center}I.\end{center}

Determinant of the substitution:
\[
\left|\begin{array}{ccccc}
0&\pm1&0&0&\cdots\\
\mp1&0&0&0&\cdots\\
0&0&1&0&\cdots\\
0&0&0&1&\cdots\\
\cdots&\cdots&\cdots&\cdots&
\end{array}\right|.
\]
All \(\varrho,\sigma\) vanish. The \(v\) are given by the formulae
\[
\begin{gathered}
v_1=\pm\frac{2\pi\sqrt{-1}}{a_{11}}\omega_1,\qquad
v_2=\frac{a_{11}\omega_2-a_{21}\omega_1}{a_{11}},\qquad
v_3=\frac{a_{11}\omega_3-a_{31}\omega_1}{a_{11}},\\
\cdots\qquad
v_p=\frac{a_{11}\omega_p-a_{p1}\omega_1}{a_{11}}.
\end{gathered}
\]
Furthermore,
\[
\begin{gathered}
b_{11}=\frac{4\pi^2}{a_{11}},
\quad b_{12}=\pm2\pi\sqrt{-1}\frac{a_{12}}{a_{11}},
\quad b_{13}=\pm2\pi\sqrt{-1}\frac{a_{13}}{a_{11}},\ldots,
\quad b_{1p}=\pm2\pi\sqrt{-1}\frac{a_{1p}}{a_{11}},\\
b_{ik}=\frac{a_{ik}a_{11}-a_{1i}a_{1k}}{a_{11}},\qquad (i,k=2,3\ldots p),\\
\Delta=\pm(2\pi\sqrt{-1})^{p-1}a_{11},
\qquad c_{11}=a_{11},\quad\hbox{all other }c_{ik}\hbox{ vanish.}
\end{gathered}
\]
The transformation formula for the \(\Theta\)-function is therefore
\[
e^{-\frac{a_{11}}{8\pi^2}v_1^2}\Theta(v,b)
=\varepsilon\sqrt{\frac{\pm a_{11}}{2\pi\sqrt{-1}}}\Theta(\omega,a).
\]
The determination of \(\varepsilon\) is made by deriving the formula directly. Namely,
\[
\Theta(\omega,a)=\sum e^{-\frac12\sum\sum a_{ik}m_im_k+\sum m_i\omega_i}
\]
\[
=\sum e^{-\frac12\sum{}'\sum{}'a_{ik}m_im_k+\sum{}'m_i\omega_i}
\Theta(\omega_1-a_{12}m_2-a_{13}m_3\ldots,a_{11}),
\]
where the sign \(!\) indicates that the summation is to be restricted to the indices \(2,3\ldots p\), and where the \(\Theta\)-function under the summation sign is the elliptic one, for which the definition is
\[
\Theta(\omega_1,a_{11})=\sum e^{-\frac12a_{11}m_1^2+m_1\omega_1}.
\]
According to Jacobi, however, at the same time
\[
\Theta(\omega_1,a_{11})=\sqrt{\frac{2\pi}{a_{11}}}
 e^{\frac{\omega_1^2}{2a_{11}}}
\Theta\left(\pm\frac{2\pi\sqrt{-1}}{a_{11}}\omega_1,\frac{4\pi^2}{a_{11}}\right).
\]
As for the sign of the root, one may put \(\omega_1=0\), \(a_{11}=2\pi\), and sees that in this case \(\sqrt{\frac{2\pi}{a_{11}}}\) must pass into 1. Since now the real part of \(a_{11}\) must always be positive if the \(\Theta\)-function is to have any meaning, it follows that \(\sqrt{\frac{2\pi}{a_{11}}}\), whose real part therefore likewise can never vanish, must always be chosen so that its real part is positive, as in the special case of 1. Substituting now in the above formula for \(\omega_1\)
\[
\omega_1-m_2a_{21}-m_3a_{31}\ldots
\]
and introducing the transformed value into \(\Theta(\omega,a)\), one obtains
\[
\Theta(\omega,a)=\sqrt{\frac{2\pi}{a_{11}}}e^{\frac{\omega_1^2}{2a_{11}}}\sum e^{-\frac12\sum\sum b_{ik}m_im_k+\sum m_iv_i},
\]
where the \(v,b\) agree exactly with the definitions given above, and where also
\[
e^{\frac{\omega_1^2}{2a_{11}}}=e^{-\frac{a_{11}}{8\pi^2}v_1^2},
\]
so that agreement with the formula derived from our general theory becomes complete if one puts
\[
\varepsilon\sqrt{\frac{\pm a_{11}}{2\pi\sqrt{-1}}}=\sqrt{\frac{a_{11}}{2\pi}},
\]
and takes the root \(\sqrt{\frac{a_{11}}{2\pi}}\) so that its real part is positive.

\begin{center}II.\end{center}

Determinant of the substitution:
\[
\left|\begin{array}{ccccc}
1&0&0&0&\cdots\\
\pm1&1&0&0&\cdots\\
0&0&1&0&\cdots\\
0&0&0&1&\cdots\\
\cdots&\cdots&\cdots&\cdots&
\end{array}\right|.
\]
The \(\varrho,\sigma\) vanish except for \(\varrho_1\), which becomes congruent to \(\pm1\), and can therefore be set equal to 1. Thus
\[
P_1=\pm2\pi\sqrt{-1},\qquad P_2=0,
\qquad P_3=0\ldots
\]
The variables are connected by the formulae
\[
v_1=\omega_1\mp\pi\sqrt{-1},\qquad
v_2=\omega_2,
\qquad v_3=\omega_3\ldots
\]
Moreover
\[
b_{11}=a_{11}\pm2\pi\sqrt{-1},\quad
\hbox{in all other cases }b_{ik}=a_{ik},
\]
\[
c_{ik}=0\quad(\hbox{for all values of }i,k),
\qquad\Delta=(2\pi\sqrt{-1})^p.
\]
Thus the transformation formula for the \(\Theta\)-function becomes
\[
\Theta(v,b)=\varepsilon\Theta(\omega,a),
\]
or
\[
\sum e^{-\frac12\sum\sum a_{ik}m_im_k-m_1^2\pi\sqrt{-1}+\sum m_i\omega_i\mp m_1\pi\sqrt{-1}}
=\varepsilon\sum e^{-\frac12\sum\sum a_{ik}m_im_k+\sum m_i\omega_i}.
\]
Since \(e^{-m_1(m_1\pm1)\pi\sqrt{-1}}=1\), the \(\Theta\)-functions on both sides agree completely, and hence
\[
\varepsilon=1.
\]

\begin{center}III.\end{center}

Determinant of the substitution:
\[
\left|\begin{array}{cccccc}
1&0&\pm1&0&0&\cdots\\
0&1&0&0&0&\cdots\\
0&0&1&0&0&\cdots\\
0&\mp1&0&1&0&\cdots\\
0&0&0&0&1&\cdots\\
\cdots&\cdots&\cdots&\cdots&\cdots&
\end{array}\right|.
\]
All \(\varrho,\sigma\) vanish. We also have
\[
v_i=\omega_i,
\]
with the single exception of \(v_2\), for which
\[
v_2=\omega_2\mp\omega_1.
\]
All \(c_{ik}\) vanish. The \(b_{ik}\) are obtained from the formulae
\[
\begin{gathered}
b_{11}=a_{11},\qquad b_{12}=a_{12}\mp a_{11},\qquad b_{22}=a_{22}\mp a_{12}+a_{11},\\
b_{2k}=a_{2k}\mp a_{1k}\qquad(k=3,4\ldots),
\end{gathered}
\]
and in all other cases \(b_{ik}=a_{ik}\). Finally \(\Delta=(2\pi\sqrt{-1})^p\). Thus the transformation formula for \(\Theta\) is
\[
\Theta(v,b)=\varepsilon\Theta(\omega,a),
\]
or
\[
\begin{aligned}
&\sum\sum e^{-\frac12\sum\sum a_{ik}m_im_k\pm(a_{11}m_1+a_{12}m_2+a_{13}m_3+\cdots)m_2}
 e^{-\frac12a_{11}m_2^2+\sum m_i\omega_i\mp m_2\omega_1}\\
&\qquad=\varepsilon\sum\sum e^{-\frac12\sum\sum a_{ik}m_im_k+\sum m_i\omega_i}.
\end{aligned}
\]
Here again the two sums agree completely if in the sum on the left one puts \(m_1\pm m_2\) in place of \(m_1\), which can be done without changing the value of the sum. Thus again
\[
\varepsilon=1.
\]
We shall add what form the reduction formula for \(\frac{\varepsilon\varepsilon'}{\varepsilon''}\) takes when, in composing the substitutions, an arbitrary \((\beta)\) is combined with one of the simple \((\beta')\). We shall then have the means of determining the number \(\varepsilon\) for any canonical substitution, once the decomposition of the canonical determinant into simple ones has been found and, conversely, that determinant is compounded from these. The formulae are:

Composition of an arbitrary determinant with I.:
\[
\varepsilon''=\varepsilon\varepsilon',
\]
where \(\varepsilon'\) is given by the formula
\[
\varepsilon'\sqrt{\frac{\pm a_{11}}{2\pi\sqrt{-1}}}=\sqrt{\frac{a_{11}}{2\pi}}.
\]
Composition of an arbitrary determinant with II.:
\[
\varepsilon''=\varepsilon\cdot
\exp\!\left(\mp\frac{2\pi\sqrt{-1}}8\sum_k\beta_{2k-1,1}(\beta_{2k,1}+2\sigma_k)\right).
\]
Composition of an arbitrary determinant with III.
\[
\varepsilon''=\varepsilon.
\]
We add to these investigations only a few remarks concerning the determination of the coefficients \(a_{ik}\) and of the lower limiting points \(\alpha\) in any expansion to be made by \(\Theta\)-functions.

\subsection*{\S\,93. Computation of the \(a_{ik}\) for the different kinds of \(\Theta\)-functions. Direct construction of the series expansions.}

In the developments of the earlier sections we nowhere defined the lower limiting points \(\alpha\), that is, the preferred contact curve of order \(n-2\). The considerations of the preceding paragraphs now taught that in fact all \(2^{p-1}(2^p+1)\) proper contact curves of order \(n-2\) can furnish the lower limiting points; and the occurrence of one or another contact curve is determined by the system of fundamental periodicity moduli taken as basis.

Instead, therefore, of determining these first, one can start from an arbitrarily chosen contact curve of order \(n-2\), whose contact points are to appear in the lower limiting points of the arguments of the \(\Theta\)-functions.

We have shown above (\(\S\,75\)) how the \(2^{p-1}(2^p+1)-1\) different non-vanishing functions
\[
(1)\quad
e^{-\frac12\sum P_h\sigma_h}\left(\frac{\Theta\left(\frac P2\right)}{\Theta(0)}\right)^4
\]
can be expressed rationally and algebraically through the roots of the equation on which the finding of that contact curve depends. If one knows the preferred contact curve, which we choose here at will, then each of the expressions (1) can be expressed through the two roots of that equation which belong to the preferred curve and to some other one. To be sure, one does not know a priori how the different expressions (1) correspond to these different roots. But if one forms the symmetric functions of the expressions (1), one obtains completely known expressions depending only on that arbitrarily preferred root.

The equations so formed can be used to compute approximately the quantities \(e^{a_{ik}}\) by algebraically known functions. In fact, nothing more is required for this than first restricting the \(\Theta\)-functions occurring in equations (1) to their first terms, thereby forming equations which allow the solution directly and suffice as first approximation equations. Continuing the approximation then gives the quantities \(e^{a_{ik}}\) with arbitrary accuracy.

Thus, starting from an arbitrarily chosen contact curve, one finds \(a_{ik}\) belonging to a definite class of \(\Theta\)-functions, and hence one obtains the means of setting up one of the expansions by \(\Theta\)-functions. Substituting these found values of the \(a_{ik}\) in the expressions (1), comparison with the corresponding algebraic expressions shows which of the various contact curves are assigned to the various expressions (1).

The method indicated here agrees exactly with that which arises in the theory of elliptic functions for expressing the quantity \(q\) by \(\sqrt{k}\) and \(\sqrt{k'}\) from the formulae given on p. 184 of the Fundamenta. But in general it also gives a direct instruction for reducing the solution of the inversion problem to a not wholly unmanageable number of operations. One sees that, for producing the final formulae, it is \emph{not} necessary to determine the position of the branch points and their assignment to the roots \(s_x\) corresponding to a point \(O\), by means of arbitrarily assumed systems of loops, nor to calculate the periodicity moduli by definite integrals, and so on.

The construction of the final formulae requires first the knowledge of the \(a_{ik}\), of which we have just spoken, but then also the knowledge of those combinations of integrals of the first kind which we called the associated normal integrals. Let \(I_1,I_2\ldots I_p\) be any mutually independent integrals of the first kind, and suppose that the inversion problem is given in the form
\[
(2)\quad
\begin{aligned}
V_1&=\sum_{i=1}^{i=p}\int_{\alpha^{(i)}}^{x^{(i)}}dI_1,\\
V_2&=\sum_{i=1}^{i=p}\int_{\alpha^{(i)}}^{x^{(i)}}dI_2,\\[-0.25em]
&\quad\cdots\\[-0.25em]
V_p&=\sum_{i=1}^{i=p}\int_{\alpha^{(i)}}^{x^{(i)}}dI_p.
\end{aligned}
\]
Then the task is to find those combinations
\[
u_k=\beta_{1k}I_1+\beta_{2k}I_2\ldots+\beta_{pk}I_p
\]
which are the normal integrals of the first kind, or the combinations
\[
v_k=\beta_{1k}V_1+\beta_{2k}V_2\ldots+\beta_{pk}V_p,
\]
which occur as arguments of the \(\Theta\)-functions. The coefficients \(\beta\) are to be determined.

From the considerations of the tenth section we were able to represent the function
\[
\left[\frac{\Theta\left(v+\frac P2\right)}{\Theta(v)}\right]^4
\]
as a precisely determined algebraic function of the \(x\). Differentiating the equation given by this representation, we obtain an equation of the form
\[
A_1dv_1+A_2dv_2\ldots+A_pdv_p=B^{(1)}dx^{(1)}+B^{(2)}dx^{(2)}\ldots+B^{(p)}dx^{(p)},
\]
or also, if we express the differentials of the \(x\) from (2) by the \(dV\),
\[
=C_1dV_1+C_2dV_2\ldots+C_pdV_p.
\]
The \(A,C\) contain both the \(v\) and the \(x\); but when one sets the \(x\) equal to the \(\alpha\), all \(v\) vanish, and the coefficients \(A,C\) thereby become completely known constant quantities. Thus we obtain linear equations with constant coefficients between the \(dv\) and \(dV\). If we form \(p\) such equations, we can express one set of differentials linearly through the other, and the obtained coefficients are the required quantities \(\beta\).

It is therefore possible, from the outset, to reach the construction of the series expansions with arbitrary approximation. But here we have reached the boundary of the investigations that can be undertaken with the aid of the methods developed in this work. The more exact execution of the indicated considerations requires methods remote from those developed here, and by its extent must be reserved for an independent presentation.
\end{document}
